\documentclass[a4paper, openany]{book}
\usepackage[czech]{babel}
\usepackage{graphicx} 
\usepackage{tikz}
\usetikzlibrary{decorations.pathmorphing, calc, babel}
\usepackage{xcolor}
\usepackage{halloweenmath}
\usepackage{amsmath, amssymb}
\usepackage{setspace}
\usepackage[hidelinks]{hyperref}
\usepackage{enumitem}
\usepackage{pdfpages}
\usepackage{fancyhdr}
\usepackage{lipsum}
\usepackage{titlesec}
\usepackage[margin=1in]{geometry} % Set margin to 1 inch
\usepackage{float}
\usepackage{subcaption}
\usepackage{wrapfig}
\usepackage{setspace}
\usepackage{tcolorbox}
\usepackage{siunitx}
\usepackage{datetime}
\usepackage{indentfirst}
\usepackage{amsthm}

\tcbuselibrary{skins, breakable, theorems}

% povoli TeXu trochu vic mezer, aby radky nepretekaly za okraj
\setlength{\emergencystretch}{2em}

% ceske znaceni goniometrickych funkci
\DeclareMathOperator{\tg}{tg}
\DeclareMathOperator{\cotg}{cotg}



\fancyhf{} 
\pagestyle{fancy}
\fancyhead[RE]{\small{\nouppercase{\leftmark}}} 
\fancyhead[LE]{$\mathrightbat$\hspace{0.45em}\thepage\hspace{0.45em}$\mathleftbat$} 
\fancyfoot[RE]{$\mathfrak{evelyn}$}
\fancyfoot[LE]{\texttt{\today, \currenttime}}

\fancyhead[RO]{$\mathrightbat$\hspace{0.45em}\thepage\hspace{0.45em}$\mathleftbat$} 
\fancyhead[LO]{\small{\nouppercase{\rightmark}}} 
\fancyfoot[RO]{\texttt{\today, \currenttime}}
\fancyfoot[LO]{$\mathfrak{evelyn}$}

\fancypagestyle{plain}{% % <-- this is new
  \fancyhf{} 
  
  \fancyfoot[RE]{$\mathfrak{evelyn}$}
  \fancyfoot[LE]{\texttt{\today, \currenttime}}

  \fancyfoot[C]{$\mathrightbat$\hspace{0.45em}\thepage\hspace{0.45em}$\mathleftbat$}

  \fancyfoot[RO]{\texttt{\today, \currenttime}}
  \fancyfoot[LO]{$\mathfrak{evelyn}$}
  % \fancyfoot[]{\thepage} % same placement as with page style "fancy"
  \renewcommand{\headrulewidth}{0pt}}


% Define a command for typical math quotes
\newcommand{\mathquote}[2]{%
    \begin{flushright}
    \begin{minipage}[t]{0.5\textwidth} % Adjust the width as needed
    \textit{#1}

    \raggedleft#2
    \end{minipage}
    \end{flushright}%
    \vspace{0.3cm}%
}


% \definecolor{colorofknowledge}{RGB}{194, 141, 225}
\definecolor{colorofwisdom}{RGB}{147, 17, 17}

\hypersetup{
    colorlinks=true,
    linkcolor=colorofwisdom, % Change this color to your desired one
    filecolor=colorofwisdom,      
    urlcolor=colorofwisdom,
}

% Define a command for examples, theorems, etc.
\newcounter{counter}[section]
\renewcommand{\thecounter}{\thesection.\arabic{counter}}

\newenvironment{example}[1][]{\refstepcounter{counter}\par\medskip
  \noindent \textbf{Příklad~\thecounter#1} \rmfamily\label{ex:\thecounter}}{\medskip}

\newenvironment{problem}[1][]{\refstepcounter{counter}\par\medskip
  \noindent \textbf{Cvičení~\thecounter#1} \rmfamily\label{pr:\thecounter}}{\medskip}

\newenvironment{hint}[1][]{
  \noindent \textbf{Nápověda~\thecounter#1} \rmfamily}{\medskip}

\newenvironment{solution}[1][]{
  \noindent \textbf{Řešení~\thecounter#1} \rmfamily}{\medskip}

\newenvironment{note}[1][]{\refstepcounter{counter}\par\medskip
  \noindent \textbf{Poznámka~\thecounter#1} \rmfamily}{\medskip}

\newenvironment{proofbox}[1][]{
  \begin{tcolorbox}[colback=white, sharp corners]}{\end{tcolorbox}}

  \newenvironment{definitionbox}[1][]{
    \begin{tcolorbox}[enhanced, breakable,
      colback=black!05!white,
      colframe=black,
      fonttitle=\bfseries,
      coltitle=white,
      colbacktitle=black,
      sharp corners,
      boxrule=0.5mm,
      frame style={ultra thick, double distance=1pt, decorate, decoration={random steps, segment length=3mm, amplitude=0.5mm}},
      drop fuzzy shadow,
      overlay={
        \node[anchor=north east, font=\small] at (frame.north east) {$\mathleftbat$};
      }
    ]
  }{
    \end{tcolorbox}
  }

% Ornament na konec kapitoly
\newcommand{\koneckapitoly}{%
  \par\vspace{1.5em}
  \begin{center}
    \begin{tikzpicture}
      \draw[thick, decorate, decoration={random steps, segment length=3mm, amplitude=0.4mm}] (-4,0) -- (-0.7,0);
      \node[scale=1.8] at (0,0) {$\mathbat$};
      \draw[thick, decorate, decoration={random steps, segment length=3mm, amplitude=0.4mm}] (0.7,0) -- (4,0);
    \end{tikzpicture}
  \end{center}
}

% Nadpisy kapitol: "Kapitola N" frakturou (Euler Fraktur nema ceske znaky,
% proto je fraktura jen u cisla kapitoly a nazev zustava v beznem pismu)
\titleformat{\chapter}[display]
  {\normalfont\bfseries\raggedright}
  {\huge\usefont{U}{euf}{b}{n}\chaptertitlename\ \thechapter}
  {20pt}
  {\Huge}

\renewenvironment{proof}[1][]{
  \noindent \textsc{Důkaz.} #1 \rmfamily}{\hfill$\blacksquare$\medskip}

% TODO: Improve, so it actually looks decent (still not satisfied with the result)
\newcommand{\bookcitation}[3]{
  \begin{definitionbox}
      #1, \textit{#2}, #3
  \end{definitionbox}
}



\begin{document}

\begin{titlepage}
  \centering
  \vspace*{2cm}
  
  % Thesis title
  \textbf{\Huge Maturita Formalita}\\[0.5cm]
  
  % Subtitle (if any)
  \textbf{\Large aneb Jak zvládnout matematiku na jedničku?}\\[1.5cm]
  
  % Date
  \textbf{\Large \today}

  \vfill
  Stížnosti a návrhy na vylepšení můžete zasílat na můj email: \\[0.5cm]
  \textbf{\href{mailto: evelyn.vybiral@gmail.com}{evelyn.vybiral@gmail.com}}.
\end{titlepage}

\thispagestyle{empty}

\setcounter{tocdepth}{1}
\tableofcontents

\chapter*{Úvod}
  
\addcontentsline{toc}{chapter}{Úvod}

\mathquote{Sine mathematica vita nulla est (Bez matematiky život nestojí za nic).}{autor neznámý}

Tento text má být snaha o trochu pokročilejší nahled na elementární matematiku potřebnou k maturitě. Měl by být vhodný jak pro studenty, kteří se chtějí připravit na maturitu, tak pro ty, kteří chtějí získat hlubší porozumění matematice. Než se pustíme do vysvětlení jak tato kniha funguje, chci Vás odkázat opět na můj email, kam můžete zasílat veškeré dotazy, nejasnosti a připomínky: 
\begin{center}
  \textbf{\href{mailto: evelyn.vybiral@gmail.com}{evelyn.vybiral@gmail.com}}.
\end{center}

\begin{definitionbox}
  \textbf{Poznámka k autorství}: Kapitoly 1 a 2 jsem napsala sama. Kapitoly 3 až 25 vznikly ve spolupráci s umělou inteligencí (Claude): AI je psala podle mého zadání, stylu a připomínek a já je postupně procházím, opravuji a upravuji. Pokud v nich najdete chybu, je to chyba nás obou a budu moc ráda, když mi ji napíšete na email výše. Vtipy, které nefungují, jsou nejspíš její. Ty dobré jsou moje.
\end{definitionbox}

Bežné matematické knihy jsou psané ve formátu definice-věta-důkaz (co tyto věci znamenají se dozvíte brzy). Toto však v tomto textu naleznete pouze v omezené míře. Místo toho se budeme snažit vysvětlit matematické pojmy a věty tak, aby byly srozumitelné pro každého. Pokusím se avšak přidat nějaký nahled na to, jak by to mohlo vypadat, kdybychom se snažili být formální.

Formát knihy bude většinou následující: nejdříve se seznámíme s nějakým pojmem což bude většinou vyznačeno následujícím rámečkem:
\begin{definitionbox}
  Důležitý pojem nebo pravidlo. Rámečkům věnujte extra pozornost.
\end{definitionbox}
Poté čtenáři představíme nějaký příklad, který by měl ilustrovat, jak daný pojem funguje. Příklady budou vyznačeny následovně:
\begin{example}
  Příklad na procvičení daného pojmu. Tento příklad bude i řešený.
\end{example}
\begin{problem}
  Cvičení na procvičení daného pojmu. 
\end{problem}

\begin{hint}
  Pokud bude cvičení náročnější budeme se snažit vždy poskytnout nápovědu aby čtenáři mohla být cesta k řešení usnadněna.
\end{hint}

\begin{solution}
  Řešení cvičení. Občas bude vynecháno, aby si čtenář mohl cvičení zkusit sám. Vynechávat řešení budeme většinou jen u problémů, které už byly dostatečně procvičeny.
\end{solution}



\chapter{Matematická logika a důkazy}
%chktex-file 44

\mathquote{Only prove something that already seems to be obvious.}{Grant Sanderson}


V této kapitole se seznámíme s některými základními pojmy a pravidly, které nám pomohou v matematickém uvažování. Základními bloky jsou pojem výrok a práce s nimi. Po zavedení těchto pojmů se seznámíme s důkazy, díky kterým můžeme ověřit správnost tvrzení.

\section{Výroková logika}

Zde se budeme pouze věnovat základní výrokové logice a na konci si lehce zmíníme něco o predikátové logice. Výroková logika je základním nástrojem matematiky, který nám umožňuje formulovat veškeré matematické tvrzení.

\begin{definitionbox}
  Výrok je tvrzení u kterého lze rozhodnout, zda je pravdivé nebo nepravdivé. Výroky budeme značit písmeny $p, q, r, \ldots$. Pravdivost výroku budeme značit $1$ a nepravdivost $0$. Těmto číslům se říká pravdivostní hodnoty.
\end{definitionbox}
\begin{example}
  Uvedeme si zde několik příkladů výroků:
  \begin{itemize}
    \item $p$: Venku prší.
    \item $q$: $2 + 2 = 69$.
    \item $r$: Praha je hlavní město České republiky.
    \item $s$: $4$ je prvočíslo.
    \item $t$: $x^2 + 1 = 0$ má reálná řešení.
  \end{itemize}
\end{example}

\begin{problem}
  Koukněte se na příklad~\ref{ex:1.1.1} a rozhodněte, zda jsou jednotlivé výroky pravdivé nebo nepravdivé.
\end{problem}

\begin{solution}
  \begin{itemize}
    \item $p$: Venku prší. --- Špatně zapsaný výrok. Není určeno, kde se nacházíme.
    \item $q$: $2 + 2 = 69$. --- Nepravda. 0.
    \item $r$: Praha je hlavní město České republiky. --- Pravda. 1.
    \item $s$: $4$ je prvočíslo. --- Nepravda. 0.
    \item $t$: $x^2 + 1 = 0$ má reálná řešení. --- Nepravda. 0.
  \end{itemize}
\end{solution}

\pagebreak
\begin{example}
  Důležité je také uvést si, co výrok není. 
  \begin{itemize}
    \item Jak se máš?
    \item Bum!
    \item Pozor na psa!
    \item Černá je nejlepší barva. (Toto je názor, nikoliv výrok. Redakce si však myslí, že to je pravda.)
  \end{itemize}
\end{example}

Ještě než se vrhneme na logické spojky, představíme si co to znamená negace výroku.

\begin{definitionbox}
  \textbf{Negace} nebo-li NOT výroku $p$ je výrok, který tvrdí opak výroku $p$. Negaci výroku $p$ budeme značit $\neg p$.
\end{definitionbox}

\begin{example}
  Vezměme si výrok $p$: Voda v moři je slaná. Negací tohoto výroku je výrok $\neg p$: Voda v moři není slaná. 
\end{example}

\begin{definitionbox}
  Upozorňujeme, že v příkladu~\ref{ex:1.1.4} opravdu měníme jen slovo `je' na `není'. Nepoužíváme například `Voda v moři je sladká'. Důvodem pro to je, že ne všechno je dichotomie pravda/nepravda. Některé výroky mohou být pravdivé, nepravdivé nebo něco mezi tím. Vezměme například výrok \textit{Kámen je zelený}. Negace je \textit{Kámen není zelený}. Ale nevíme, zda je kámen modrý, červený nebo třeba duhový.
\end{definitionbox}

\begin{problem}
  Opět se vraťte k příkladu~\ref{ex:1.1.1} a určete negace jednotlivých výroků.
\end{problem}

\begin{solution}
  \begin{itemize}
    \item $\neg p$: Venku neprší.
    \item $\neg q$: $2 + 2 \neq 69$.
    \item $\neg r$: Praha není hlavní město České republiky.
    \item $\neg s$: $4$ není prvočíslo.
    \item $\neg t$: $x^2 + 1 = 0$ nemá reálná řešení.
  \end{itemize}
\end{solution}

Výroky samotné jsou sice fajn ale co s nimi dělat? K tomu nám poslouží logické spojky. Je to způsob jakým můžeme spojovat výroky do složitějších celků. Proč? Dozvíte se za chvíli.

\begin{definitionbox}
  Rozlišujeme 4 základní logické spojky:

  \textbf{Konjunkce ($\land$):} Konjunkce nebo-li AND výroků $p$ a $q$ je výrok, který je pravdivý právě tehdy, když jsou pravdivé oba výroky $p$ a $q$. 

  \textbf{Disjunkce ($\lor$):} Disjunkce nebo-li OR výroků $p$ a $q$ je výrok, který je pravdivý právě tehdy, když je pravdivý alespoň jeden z výroků $p$ a $q$.

  \textbf{Implikace ($\Rightarrow$):} Implikace nebo-li IF-THEN výroků $p$ a $q$ je výrok, který je nepravdivý právě tehdy, když je výrok $p$ pravdivý a $q$ nepravdivý. V tomto případě lze nad tím přemýšlet jako nad slibem. Když mi někdo něco slíbí, tak to musí splnit, ale pokud to neslíbí, tak je jedno co udělá.

  \textbf{Ekvivalence ($\Leftrightarrow$):} Ekvivalence nebo-li IF-AND-ONLY-IF výroků $p$ a $q$ je výrok, který je pravdivý právě tehdy, když jsou oba výroky $p$ a $q$ pravdivé nebo oba nepravdivé.

  \vspace{0.2cm}

  Pro větší přehlednost si můžete vytvořit tabulku pravdivostních hodnot pro jednotlivé logické spojky.

  \begin{center}
    \begin{tabular}{|c|c|c|c|c|c|c|}
      \hline
      $p$ & $q$ & $p \land q$ & $p \lor q$ & $p \Rightarrow q$ & $p \Leftrightarrow q$ \\
      \hline
      0 & 0 & 0 & 0 & 1 & 1 \\
      0 & 1 & 0 & 1 & 1 & 0 \\
      1 & 0 & 0 & 1 & 0 & 0 \\
      1 & 1 & 1 & 1 & 1 & 1 \\
      \hline
    \end{tabular}
  \end{center}
\end{definitionbox}

\pagebreak

\begin{example}
  Uveďme si několik příkladů na logické spojky:
  \begin{itemize}
    \item Konjunkce: $p \land q$: Venku prší a je teplo.
    \item Disjunkce: $p \lor q$: Venku prší nebo je teplo.
    \item Implikace: $p \Rightarrow q$: Když prší, tak je mokro.
    \item Ekvivalence: $p \Leftrightarrow q$: Venku prší právě tehdy, když je mokro.
  \end{itemize}
\end{example}

\begin{problem}
  Vraťte se k příkladu~\ref{ex:1.1.1} a zkuste si spojit výroky pomocí logických spojek.
\end{problem}

\section{Odbočka --- Predikátová logika}
Predikátová logika staví na výrokové logice a rozšiřuje ji o kvantifikátory. Pracuje se zde také s množinami výroků. Ty nám umožňují mluvit o všech prvcích nějaké množiny. Základními kvantifikátory jsou:
\begin{itemize}
  \item Univerzální kvantifikátor $\forall$: Výrok $\forall x \in M: P(x)$ znamená, že výrok $P(x)$ platí pro všechny prvky množiny $M$.
  \item Existenciální kvantifikátor $\exists$: Výrok $\exists x \in M: P(x)$ znamená, že existuje prvek množiny $M$, pro který platí výrok $P(x)$.
\end{itemize}

Toto se nám bude hodit při důkazech, kde budeme chtít mluvit o všech prvcích nějaké množiny (co je množina se také dozvíte). Negace a spojování kvantifikátorů také funguje, ale je to lehce složitější.

\begin{example}
  V predikátové logice se pracuje v takzvaném univerzu. To si vždy zvolíme. Může to být například množina všech reálných čísel, množina všech lidí, množina všech zvířat, atd. Většinou se značí písmenem $U$. Ukažme si nyní dva příklady. Pokud si za naše univerzum zvolíme množinu všech reálných čísel, můžeme napsat:
  \begin{itemize}
    \item $\forall x \in \mathbb{R}: x^2 \geq 0$. Tento výrok je pravdivý, protože druhá mocnina reálného čísla je vždy nezáporná.
    \item $\exists x \in \mathbb{R}: x^2 = -1$. Tento výrok je nepravdivý, protože druhá mocnina reálného čísla je vždy nezáporná.  
  \end{itemize}
  Nyní zvolme univerzum všech lidí. Zvolme si takzvané predikáty $M(x)$ a $G(x)$, které budou respektive znamenat, že $x$ zvládnul maturitu a $x$ je génius. Nyní můžeme napsat:
  \begin{itemize}
    \item $\forall x \in \text{Lidé}: M(x) \Rightarrow G(x)$. Tento výrok říká, že každý, kdo zvládl maturitu, je génius. O pravdivosti se zde raději nebudeme bavit.
    \item $\exists x \in \text{Lidé}: M(x) \land G(x)$. Tento výrok říká, že existuje někdo, kdo zvládl maturitu a je génius.
  \end{itemize}

  Všimněte si, že v prvním příkladě jsme mluvili o všech reálných číslech a v druhém o všech lidech. To je důležité, protože výrok může být pravdivý v jednom univerzu a nepravdivý v jiném.

  \begin{definitionbox}
    Nyní si ukážeme, jak se neguje kvantifikátor. Negace univerzálního kvantifikátoru je existenciální kvantifikátor a naopak.
  \begin{equation}
    \begin{aligned}
      \neg (\forall x \in M: P(x)) \Leftrightarrow \exists x \in M: \neg P(x) \\
      \neg (\exists x \in M: P(x)) \Leftrightarrow \forall x \in M: \neg P(x)  
    \end{aligned}
  \end{equation}
  Je to celkem intuitivní. Pokud něco platí pro všechny, stačí nám najít jedno, kde to neplatí. Pokud neexistuje ani jedno, kde to platí, tak to pro všechny neplatí, tedy pro každé platí negace.
  \end{definitionbox}
\end{example}

\section{Odbočka --- Přirozená dedukce}
Jeden z formálních dokazovacích systémů je přirozená dedukce. Tento systém je založen na pravidlech, které nám umožňují odvodit nové výroky z již známých.\footnote{Přirozenou dedukci lze zavést i pro predikátovou logiku, ale to je nad rámec tohoto textu.} Většina pravidel je intuitivních a vychází z toho, jak my lidé přirozeně uvažujeme. Tato sekce je uvedena opět jako odbočka, ale doporučuji si ji projít abyste v dalších kapitolách chápali jak se dělají důkazy.

Je zde pár pravidel, ale ty si uvádět nebudeme. Koukneme se na příklady a vysvětlíme si to u nich.

\begin{definitionbox}
  Většinou naším cílem bude odvodit nějaký výrok $q$ z předpokladů $p_1, p_2, \ldots, p_n$. To značíme jako $p_1, p_2, \ldots, p_n \vdash q$.  Předpoklad je výrok, který považujeme za pravdivý. Nutno ještě zmínit co znamená pomocný předpoklad. Pro větší přehlednost si kreslíme box kolem našeho `světa' ve kterém se pohybujeme. Pokud si vytvoříme nový box, tak to znamená, že v něm a jen v něm platí ty předpoklady, které jsme si tam napsali. Pokud z něj odejdeme, tak už ty předpoklady neplatí. Může se to zdát zbytečné, ale většinou nám to pomůže důkaz dokončit viz příklad~\ref{ex:1.3.2}.
\end{definitionbox}

\begin{example}
  Ukažme si nyní, jak bychom mohli dokázat, že $p \land q \vdash q \land p$. Budeme postupovat pomocí přirozené dedukce. Začneme tedy předpokladem $p \land q$ a budeme chtít dokázat $q \land p$. 
  \begin{proofbox}
    1. $p \land q$ (předpoklad) \\
    2. $p$ (eliminace konjunkce z 1) \\
    3. $q$ (eliminace konjunkce z 1) \\
    4. $q \land p$ (zavedení konjunkce z 3 a 2)
  \end{proofbox}
  Tímto jsme dokázali, že $p \land q \vdash q \land p$. Zde jsme použili pravidla eliminace konjunkce a zavedení konjunkce. Tato pravidla vychází z toho, že pokud víme, že platí $p \land q$, tak víme, že platí $p$ a $q$ samostatně a naopak. U ostatních pravidel to tak jednoduché nebude.
\end{example}

Provedeme si nyní nějakou přípravu, která bude působit zbytečně, ale v dalších kapitolách se nám to bude hodit.
\begin{example}
  Ukažme si, jak bychom mohli dokázat, že $p \Rightarrow q \vdash \neg q \Rightarrow \neg p$. Budeme postupovat pomocí přirozené dedukce. Začneme tedy předpokladem $p \Rightarrow q$ a budeme chtít dokázat $\neg q \Rightarrow \neg p$. 
  \begin{proofbox}
    1. $p \Rightarrow q$ (předpoklad) 
    \begin{proofbox}
      2. $\neg q$ (pomocný předpoklad) \\
      3. $\neg p$ (modus tollens z 1 a 2)
    \end{proofbox}
    4. $\neg q \Rightarrow \neg p$ (zavedení implikace z 2 a 3)
  \end{proofbox}
  Tímto jsme dokázali, že $p \Rightarrow q \vdash \neg q \Rightarrow \neg p$. Zde jsme použili pravidlo modus tollens. To říká, že pokud víme, že platí $p \Rightarrow q$ a $\neg q$, tak můžeme odvodit $\neg p$. To vychází přímo z definice implikace. Pokud někdo něco nesplnil, ale slib `vyšel' pozitivně, tak to znamená, že to ani neslíbil. Zde je krásně vidět, jak nám pomocný předpoklad pomohl. Protože jsme viděli, že když budeme předpokládat $\neg q$, tak můžeme odvodit $\neg p$ a z toho jsme zavedli chtěnou implikaci.
\end{example}

Pozorného čtenáře jistě napadlo, že místo tohoto postupu si můžeme napsat pravdivostní tabulku a zjistit, že tvrzení platí. To bychom se připravili o všechnu zábavu a pro složitější tvrzení by tabulka rostla exponenciálně.
Uvedeme si zde pro představu jak by vypadala tabulka pro příklad~\ref{ex:1.3.2}
\begin{center}
  \begin{tabular}{|c|c|c|c|c|c|}
    \hline
    $p$ & $q$ & $p \Rightarrow q$ & $\neg q$ & $\neg p$ & $\neg q \Rightarrow \neg p$ \\
    \hline
    0 & 0 & 1 & 1 & 1 & 1 \\
    0 & 1 & 1 & 0 & 1 & 1 \\
    1 & 0 & 0 & 1 & 0 & 0 \\
    1 & 1 & 1 & 0 & 0 & 1 \\
    \hline
  \end{tabular}
\end{center}

\begin{example}
  Podívejme se nyní na důležitý příklad $p\Rightarrow q \dashv\vdash \neg p\lor q$. 
  \begin{proofbox}
    1. $p \Rightarrow q$ (předpoklad) \\
    2. $p \lor \neg p$ (tertium non datur) 
    \begin{proofbox}
      3. $p$ (pomocný předpoklad) \\
      4. $q$ (modus ponens z 1 a 3) \\
      5. $\neg p \lor q$ (zavedení disjunkce z 4)
    \end{proofbox}
    \begin{proofbox}
      6. $\neg p$ (pomocný předpoklad) \\
      7. $\neg p \lor q$ (zavedení disjunkce z 6)
    \end{proofbox}
    8. $\neg p \lor q$ (eliminace disjunkce z 2, 3--5 a 6--7)
  \end{proofbox}
  \begin{proofbox}
    1. $\neg p \lor q$ (předpoklad) 
    \begin{proofbox}
      2. $p$ (pomocný předpoklad) \\
      3. $q$ (eliminace disjunkce z 1 a 2) 
    \end{proofbox}
    4. $p \Rightarrow q$ (zavedení implikace z 2 a 3)
  \end{proofbox}
  Využili jsme zde pravidla \textit{tertium non datur} nebo-li pravidlo vyloučeného třetího. To říká, že v našem logickém systému platí, že platí buď $p$ nebo $\neg p$.
\end{example}

\begin{note}
  Pro zajímavost existují další běžné logické spojky, které se dají odvodit z těchto základních. Máme zde například XOR (exclusive OR) $\oplus$, který říká, že platí buď $p$ nebo $q$, ale ne oba zároveň. Dalšími jsou NAND $\mid$ a NOR $\downarrow$, které jsou negacemi AND a OR.\@ 
\end{note}

\begin{problem}
  Napište si, jak by vypadaly výše zmíněné spojky pomocí základních spojek.
\end{problem}

\begin{example}
  Pojďme si nyní odpočinout od úplně formálních příkladů a podívat na trochu méně formální příklad. Mějme následující tvrzení:
  \begin{itemize}
    \item Na výlet pojede Petr nebo Quido.
    \item Jestliže pojede Quido, pojede Simona a nepojede Renata.
    \item Jestliže pojede Tomáš, pojede i Renata.
    \item Jestliže pojede Simona, pojede Tomáš.
  \end{itemize}
  Chceme zjistit, zda je pravdivé tvrzení `Petr pojede na výlet'.\footnote{Toto lze také převést na problém sémantického důsledku a řešit rezoluční metodou. Tímto spojením se zabývá věta o úplnosti, která je též nad rámec této středoškolské knihy.} Budeme postupovat pomocí přirozené dedukce. Začneme tedy předpoklady, které máme a budeme chtít dokázat, že Petr pojede na výlet. Pojmenujeme si výroky $p$: Petr pojede na výlet, $q$: Quido pojede na výlet, $s$: Simona pojede na výlet, $t$: Tomáš pojede na výlet, $r$: Renata pojede na výlet.

  \begin{proofbox}
    1. $p \lor q$ (předpoklad) \\
    2. $q \Rightarrow (s \land \neg r)$ (předpoklad) \\
    3. $t \Rightarrow r$ (předpoklad) \\
    4. $s \Rightarrow t$ (předpoklad) 
  \begin{proofbox}
    5. $\neg p$ (pomocný předpoklad) \\
    6. $q$ (eliminace disjunkce z 1 a 5) \\
    7. $s \land \neg r$ (modus ponens z 2 a 6) \\
    8. $s$ (eliminace konjunkce z 7) \\
    9. $t$ (modus ponens z 4 a 8) \\
    10. $r$ (modus ponens z 3 a 9) \\
    11. $\neg r$ (eliminace konjunkce z 7) \\
    12. $\perp$ (eliminace negace z 10 a 11)
  \end{proofbox}
    13. $\lnot\lnot p$ (reductio ad absurdum z 5--12) \\
    14. $p$ (eliminace dvojí negace z 13)
  \end{proofbox}
  Tímto jsme dokázali, že Petr pojede na výlet. Zde jsme použili pravidla eliminace disjunkce, modus ponens a reductio ad absurdum. Pojďme si spolu pravidla vysvětlit. 
  \begin{itemize}
    \item Eliminace disjunkce: Pokud víme, že platí $p \lor q$ a víme, že jedno z nich neplatí, tak můžeme odvodit to druhé. 
    \item Modus ponens: Pokud víme, že platí $p \Rightarrow q$ a $p$, tak můžeme odvodit $q$. Někdo mi něco slíbil a já vím, že slib byl úspěšný, takže to i splnil.
    \item Reductio ad absurdum: Pokud předpokládáme něco a dokážeme, že z toho vyplývá spor, tak můžeme odvodit negaci našeho předpokladu.
  \end{itemize}
\end{example}

\begin{problem}
  Ať si taky něco procvičíte, zkuste si dokázat, že $p \Rightarrow q, q \Rightarrow r \vdash p \Rightarrow r$.\footnote{Pokud se vám to povede, tak gratuluji. Nejsou zde vypsaná pravidla, takže se ani nepředpokládá, že byste to měli zvládnout.}
\end{problem}

% \pagebreak

\begin{solution}
  \begin{proofbox}
    1. $p \Rightarrow q$ (předpoklad) \\
    2. $q \Rightarrow r$ (předpoklad) 
  \begin{proofbox}
    3. $p$ (pomocný předpoklad) \\
    4. $q$ (modus ponens z 1 a 3) \\
    5. $r$ (modus ponens z 2 a 4)
  \end{proofbox}
    6. $p \Rightarrow r$ (zavedení implikace z 3--5)
  \end{proofbox}
\end{solution}

\begin{definitionbox}
  Abyste opravdu věřili, že tato metoda je lepší než pravdivostní tabulky, tak se pojďme u příkladu~\ref{ex:1.3.6} zamyslet nad velikostí tabulky. Vzhledem k tomu, že máme 5 proměnných, tak abychom dosáhli všech možných kombinací, musí mít tabulka $2^5 = 32$ řádků! Představte si, že je všechny musíte vyplňovat a pak z toho něco vyčíst. 
\end{definitionbox}

Nebudeme vás nadále trápit, ale ještě si ukážeme odkud se berou vzorečky negací složených výroků.

\begin{example} \textbf{(De Morganovy zákony)}
  De Morganovy zákony jsou zákony, které nám říkají, jak se neguje konjunkce a disjunkce. Jsou to následující vzorečky:
  \begin{equation}
    \begin{aligned}
      \neg (p \land q) &\dashv\vdash \neg p \lor \neg q \\
      \neg (p \lor q) &\dashv\vdash \neg p \land \neg q
    \end{aligned}
  \end{equation}

  Tyto zákony se dají dokázat pomocí pravdivostní tabulky, ale my si je zde formálně dokážeme. Dají se zobecnit matematickou indukcí (už se brzy dozvíte co to je) na libovolný počet výroků následujícím způsobem:
  \begin{equation}
    \begin{aligned}
      \neg \left(\bigwedge_{i=1}^{n} p_i\right) &\dashv\vdash \bigvee_{i=1}^{n} \neg p_i \\
      \neg \left(\bigvee_{i=1}^{n} p_i\right) &\dashv\vdash \bigwedge_{i=1}^{n} \neg p_i.
    \end{aligned}
  \end{equation}

  Uvedeme zde důkaz prvního De Morganova zákona. Zbytek necháme jako cvičení pro čtenáře.

  \begin{proofbox}
    % $\neg (p \land q) \vdash \neg p \lor \neg q$ 
    % \vspace{0.4cm}

    1. $\neg (p \land q)$ (předpoklad)
    \begin{proofbox}
      2. $\neg(\neg p \lor \neg q)$ (pomocný předpoklad)
      \begin{proofbox}
        3. $\neg p$ (pomocný předpoklad) \\
        4. $\neg p \lor \neg q$ (zavedení disjunkce z 3) \\
        5. $\perp$ (eliminace negace z 4 a 2)
      \end{proofbox}
      6. $p$ (reductio ad absurdum z 3--5) 
      \begin{proofbox}
        7. $\neg q$ (pomocný předpoklad) \\
        8. $\neg p \lor \neg q$ (zavedení disjunkce z 7) \\
        9. $\perp$ (eliminace negace z 8 a 2)
      \end{proofbox}
      10. $q$ (reductio ad absurdum z 7--9) \\
      11. $p \land q$ (zavedení konjunkce z 6 a 10) \\
      12. $\perp$ (eliminace negace z 11 a 1)
    \end{proofbox}
    13. $\neg\neg(\neg p \lor \neg q)$ (reductio ad absurdum z 2--12) \\
    14. $\neg p \lor \neg q$ (eliminace dvojí negace z 13)
  \end{proofbox}

  \begin{proofbox}
    % $\neg p \lor \neg q \vdash \neg (p \land q)$ 
    % \vspace{0.4cm}

    1. $\neg p \lor \neg q$ (předpoklad)
    \begin{proofbox}
      2. $p \land q$ (pomocný předpoklad) \\
      3. $p$ (eliminace konjunkce z 2) \\
      4. $q$ (eliminace konjunkce z 2)
      \begin{proofbox}
        5. $\neg p$ (pomocný předpoklad) \\
        6. $\perp$ (eliminace negace z 3 a 5)
      \end{proofbox}
      \begin{proofbox}
        7. $\neg q$ (pomocný předpoklad) \\
        8. $\perp$ (eliminace negace z 4 a 7)
      \end{proofbox}
      9. $\perp$ (eliminace disjunkce z 1, 5--6 a 7--8)
    \end{proofbox}
    10. $\neg(p \land q)$ (reductio ad absurdum z 2--9)
  \end{proofbox}

\end{example}

Mohli bychom pokračovat dál, ale máme tu ještě spoustu dalších zajímavých věcí, které chceme probrat. Pokud vás téma logiky zaujalo, tak si doporučuji přečíst si knihu 
\bookcitation{Jiří Velebil}{\href{https://math.fel.cvut.cz/en/people/velebil/files/y01mlo/logika.pdf}{Velmi jemný úvod do matematické logiky}}{Praha 2007}
Chcete-li jít více do hloubky a dozvědět se i více o těch prapodivných boxech, které jsme používali, tak si přečtěte knihu
\bookcitation{M. Demlová a B. Pondělíček}{Matematická logika}{skriptum FEL ČVUT, Praha, 1997}
\section{Důkazy}
Důkazy jsou základním nástrojem matematiky a bez nich bychom se nedostali daleko. V této kapitole si ukážeme několik základních typů důkazů, které se v matematice používají. V úvodní kapitole byl zmíněn systém definice-věta-důkaz. Než tedy přejdeme k důkazům, tak si ještě jednou připomeneme co tento systém znamená.

\begin{definitionbox}
  \textbf{Definice} je vymezení pojmu, se kterým hodláme pracovat. Definice obsahuje název a charakteristické vlastnosti daného pojmu.

  \textbf{Věta} je výrok, který je pravdivý a musí být dokázán. Většinou rozvíjíme nějaký pojem, který jsme si definovali.

  \textbf{Důkaz} je postup, který nám ukazuje, proč je věta pravdivá. Důkaz se skládá z předpokladů a závěru. Předpoklady jsou věty, které už máme dokázané a závěr je věta, kterou chceme dokázat.
\end{definitionbox}

\begin{example}
  Vyzkoušejme si zde jak by vypadal jednoduchý důkaz. V matematické literatuře by to vypadalo následovně:

  \vspace{0.3cm}

  \noindent\textbf{Věta 1.2.3\footnote{Nenechte se zmást značením. Je zde uvedeno jako příklad, nezapadá však do našeho číslování.} (O součtu dvou sudých čísel)}: Součet dvou sudých čísel je sudé číslo.

  \begin{proof}
    Nechť $a$ a $b$ jsou sudá čísla. Pak existují taková čísla $k, l \in \mathbb{Z}$, že $a = 2k$ a $b = 2l$. Potom
    \begin{equation}
      a + b = 2k + 2l = 2(k + l),
    \end{equation}
    kde $k + l \in \mathbb{Z}$. Tedy $a + b$ je sudé číslo.
  \end{proof}
  
  Důkaz není úplně korektní, ale za předpokladu, že víme, co znamená, že číslo je sudé\footnote{Správná definice je ovšem, že číslo je sudé, pokud lze napsat jako 2k, kde $k\in\mathbb{Z}$. U lichých je to 2k + 1, kde $k\in\mathbb{Z}$}, tak bychom mohli tento důkaz přijmout.
\end{example}

\begin{problem}
  Zkuste si nyní dokázat, že součet dvou lichých čísel je sudé číslo.
\end{problem}

\begin{solution}

  \begin{proof}
    Nechť $a$ a $b$ jsou lichá čísla. Pak existují taková čísla $k, l \in \mathbb{Z}$, že $a = 2k + 1$ a $b = 2l + 1$. Potom
    \begin{equation}
      a + b = 2k + 1 + 2l + 1 = 2(k + l + 1),
    \end{equation}
    kde $k + l + 1 \in \mathbb{Z}$. Tedy $a + b$ je sudé číslo.
  \end{proof}
\end{solution}

\begin{definitionbox}
  Pojďme si nyní ukázat několik základních typů důkazů, které se v matematice používají. Mějme předpoklad $p$ a závěr $q$.

  \textbf{Důkaz přímý}: Důkaz, kde se jde \textit{přímo} z předpokladu k závěru. V jazyku logických spojek to znamená, že se použije implikace $p \Rightarrow q$. Tento typ důkazu jsme si ukázali v příkladu~\ref{ex:1.4.1}.

  \textbf{Důkaz nepřímý}: Důkaz, kde se používá takzvaná obměněná implikace. V jazyku logických spojek to znamená, že se použije implikace $\neg q \Rightarrow \neg p$. To jsme si odvodili v příkladu~\ref{ex:1.3.2}.

  \textbf{Důkaz sporem}: Důkaz, kde se zneguje výraz $p \Rightarrow q$ a dokáže se, že z toho vyplývá nějaký spor. Zde můžeme použít příklad~\ref{ex:1.3.3} ve kterém jsme si dokázali, že $p\Rightarrow q \vdash \neg p \lor q$. Pokud tento výrok znegujeme za pomoci De Morganových zákonů, tak dostaneme $p \land \neg q$. Toto je nyní náš předpoklad, ze kterého se snažíme vyvodit spor. 

  \vspace{0.1cm}

  Poslední typ zde úplně nezapadá, takže si ho uvedeme za chvíli samostatně.
\end{definitionbox}

\begin{example}
  Ukažme si nyní příklad přímého důkazu. Mějme následující větu: \textit{Pokud je číslo $n$ sudé, tak je i $n^2$ sudé.}. Zde opět využijeme naší definice sudých čísel z předchozích příkladů. Máme tedy číslo $n$, které je sudé. To znamená, že lze napsat jako $2k$, kde $k\in\mathbb{Z}$. Nyní si spočítáme $n^2$: 
  \begin{equation}
    n^2 = (2k)^2 = 4k^2 = 2(2k^2).
  \end{equation}
  S trochou důvtipu lze vidět, že $2k^2 \in \mathbb{Z}$.\footnote{Zde bychom potřebovali něco o teorii celých čísel. Nám postačí, když tomu budete věřit.} To znamená, že se nám podařilo $n^2$ napsat jako $2l$, kde $l\in\mathbb{Z}$. Tedy $n^2$ je sudé číslo.
\end{example}
\begin{problem}
  Zkuste si nyní dokázat, že pokud je číslo $n$ liché, tak je i $n^2$ liché. Postupujte podobně jako v předchozím příkladu.
\end{problem}

\begin{example} 
  Představte si, že máte dokázat větu: \textit{Pokud je číslo $n^2$ sudé, tak je i $n$ sudé.} Zde zkusíme použít nepřímý důkaz. Provedeme obměněnou implikaci a dostaneme výrok $\neg q \Rightarrow \neg p$. Tedy pokud je $n$ liché, tak je i $n^2$ liché. Toto jste již udělali ve cvičení~\ref{pr:1.4.4}.
\end{example}

\begin{example}
  Přejděme na důkaz sporem. Oblíbeným příkladem je důkaz iracionality čísla $\sqrt{2}$. Chceme ukázat, že $\sqrt{2}\notin\mathbb{Q}$. Provedením negace dostaneme tedy předpoklad, že $\sqrt{2}$ je racionální číslo. To znamená, že existují čísla $p$ a $q$ taková, že $\sqrt{2} = \frac{p}{q}$.\footnote{Definice racionálních čísel.} Dále si můžeme předpokládat, že $p$ a $q$ jsou nesoudělná čísla. To znamená, že nemají žádný společný dělitel. Nyní si můžeme napsat:
  \begin{equation}
    \sqrt{2} = \frac{p}{q} \Rightarrow 2 = \frac{p^2}{q^2} \Rightarrow 2q^2 = p^2.
  \end{equation}
  To znamená, že $p^2$ je sudé číslo. Z předchozích příkladů víme, že to znamená, že $p$ je sudé číslo. To znamená, že existuje číslo $k$ takové, že $p = 2k$. Dosadíme zpět do našeho vztahu:
  \begin{equation}
    2q^2 = (2k)^2 \Rightarrow 2q^2 = 4k^2 \Rightarrow q^2 = 2k^2.
  \end{equation}
  To znamená, že $q^2$ je sudé číslo. Z předchozích příkladů víme, že to znamená, že $q$ je sudé číslo. To je ale spor s naším předpokladem, že $p$ a $q$ nemají žádný společný dělitel. Tedy jsme dokázali, že $\sqrt{2}$ je iracionální číslo.
\end{example}
\begin{definitionbox}
  K důkazům se určitě ještě vrátíme v dalších kapitolách u konkrétních vět, které stojí za dokázání. Teď si uvedeme ještě jeden běžný typ důkazu, který se v matematice používá.

  \vspace{0.2cm}

  \textbf{Důkaz matematickou indukcí}: Důkaz, který se používá pro dokazování vět, které platí pro všechna přirozená čísla. Důkaz se skládá ze tří kroků:
  \begin{itemize}
    \item \textbf{Základní krok}: Dokážeme, že věta platí pro první přirozené číslo.
    \item \textbf{Indukční krok}: Předpokládáme, že věta platí pro nějaké přirozené číslo $n$. Tomuto se říká indukční předpoklad.
    \item \textbf{Závěr}: Využitím indukčního předpokladu dokážeme, že věta platí i pro $n+1$ čísel.
  \end{itemize}
\end{definitionbox}

\begin{example}
  Ukažme si nyní důkaz matematickou indukcí na příkladu. Mějme následující větu: \textit{Pro každé přirozené číslo $n$ platí, že $1 + 2 + \cdots + n = \frac{n(n+1)}{2}$.}\footnote{Zkuste se nejdřív zamyslet, proč by to mohlo platit. Legenda praví, že toto tvrzení objevil slavný matematik Gauss, když byl ještě dítě.}  

  \begin{proof}

    \textbf{Základní krok}: Pro $n = 1$ platí, že $1 = \frac{1(1+1)}{2} = 1$. Tedy věta platí pro $n = 1$.

    \textbf{Indukční krok}: Předpokládejme, že věta platí pro nějaké přirozené číslo $n$. To znamená, že platí rovnost $1 + 2 + \cdots + n = \frac{n(n+1)}{2}$.

    \textbf{Závěr}: Nyní se pokusíme dokázat, že věta platí i pro $n+1$. Tedy
    \begin{equation}
      1 + 2 + \cdots + n + (n+1) = \frac{n(n+1)}{2} + (n+1) = \frac{n(n+1) + 2(n+1)}{2} = \frac{(n+1)(n+2)}{2}.
    \end{equation}
    Tímto jsme dokázali, že věta platí i pro $n+1$. Tedy věta platí pro všechna přirozená čísla.
  \end{proof}
\end{example}
\section{Příběh --- Logik, který se bál jídla}

V září 1930 se v Královci konala konference o základech matematiky. Při závěrečné diskusi se přihlásil o slovo nenápadný čtyřiadvacetiletý Vídeňák Kurt Gödel a skoro mimochodem oznámil výsledek, který otřásl celou logikou. V každém bezesporném formálním systému, který je dost silný na to, aby popsal aritmetiku přirozených čísel, existují \textbf{pravdivá tvrzení, která v něm nelze dokázat}. Den nato pronesl ve stejném městě David Hilbert slavný projev zakončený slovy `Musíme vědět. Budeme vědět.' (víc v kapitole~\ref{kap:hilbert}). Že některé věci vědět nebudeme, už bylo dokázáno.

Gödelův trik připomíná paradox lháře. Sestrojil aritmetické tvrzení, které o sobě v jistém smyslu říká `já nejsem dokazatelné'. Kdyby šlo dokázat, systém by dokazoval nepravdu, a byl by tedy sporný. Takže dokázat nejde, a proto je pravdivé. Je to úvaha sporem, přesně taková, jakou jsme poznali v této kapitole, jen o samotném dokazování.

V roce 1940 Gödel s manželkou Adele odjel z nacisty obsazeného Rakouska transsibiřskou magistrálou a přes Japonsko do Princetonu. Tam se spřátelil s Albertem Einsteinem, který prý tvrdil, že do institutu chodí už jen proto, aby se mohl s Gödelem procházet domů. Když Gödel v roce 1947 žádal o americké občanství, připravoval se na pohovor tak důkladně, že v americké ústavě našel logickou mezeru, která by podle něj umožnila zákonně nastolit diktaturu. Začal o ní soudci vykládat a Einstein se svědkem Oskarem Morgensternem ho prý museli rychle uklidnit. Občanství dostal.

S přibývajícím věkem u Gödela sílila paranoia. Bál se, že ho někdo otráví, a jedl jen jídlo, které mu připravila a ochutnala Adele. Když byla v roce 1977 na půl roku hospitalizována, přestal jíst úplně. Zemřel v lednu 1978 v princetonské nemocnici. Vážil asi třicet kilogramů. Úmrtní list uvádí `podvýživu a vyčerpání způsobené poruchou osobnosti'.

\koneckapitoly

\chapter{Množiny, relace a čísla}
\mathquote{From the paradise that Cantor created for us no-one shall be able to expel us.}{David Hilbert}

\section{Odbočka --- Axiomatická teorie množin}
Než se pustíme do pro nás více relevantních témat, pojďme se seznámit s rigorózní teorií množin na které je postavena celá matematika. Teorii množin vytvořil Georg Cantor na konci 19. století, axiomy, které si zde ukážeme, ale sepsali až na začátku 20. století Ernst Zermelo a Abraham Fraenkel. Je základem pro všechny matematické koncepty, které se týkají množin. Tato teorie je ale velice složitá a zmiňuji ji zde jen proto, abychom měli představu o tom, jak by mohla matematika vypadat, kdybychom se snažili být formální.

Hlavní roli v této teorii hrají axiomy, což jsou základní pravidla, která musíme přijmout jako pravdivá, abychom mohli dále pracovat. Než se pustíme na samotné axiomy, vysvětleme si rychle následující znaky, které nám usnadní zápis. Množiny značíme velkými písmeny, například $X$, $Y$, $Z$, atd. Prvky množin značíme malými písmeny, například $x$, $y$, $z$, atd. Pokud prvek $x$ patří do množiny $X$, značíme to $x\in X$ a pokud nepatří, tak $x\notin X$. Systém množin je vlastně také množina, která obsahuje množiny a značíme ji $\mathcal{X}$, $\mathcal{Y}$, $\mathcal{Z}$, atd. Projděme si nyní axiomy množinové teorie a pak si ukážeme základní příklady a vysvetlení těchto axiomů. 

\begin{definitionbox}
  \textbf{Axiomy množinové teorie}
  \begin{enumerate}
    \item \textbf{Axiom of extension}: Dvě množiny jsou rovny právě tehdy, když obsahují stejné prvky.
    \item \textbf{Axiom of the empty set}: Existuje množina, která neobsahuje žádné prvky. Tato množina se nazývá prázdná množina a značí se $\emptyset$.
    \item \textbf{Axiom schema of separation}: Pro každou množinu $x$ a každou vlastnost $P$ existuje množina, která obsahuje všechny prvky množiny $x$ splňující vlastnost $P$.
    \item \textbf{Axiom of pairing}: Pro každé dvě množiny $X$ a $Y$ existuje množina $\{X, Y\}$, jejímiž prvky jsou právě množiny $X$ a $Y$. 
    \item \textbf{Axiom of union}: Pro každý systém množin $\mathcal{S}$ existuje množina, která obsahuje všechny prvky všech množin v systému $\mathcal{S}$. Tato množina se nazývá sjednocení množin v systému $\mathcal{S}$ a značí se $\bigcup \mathcal{S}$. Pokud bychom chtěli sjednotit pouze množiny $X$ a $Y$, značili bychom to $X\cup Y$.
    \item \textbf{Axiom of power set}: Pro každou množinu $X$ existuje systém všech podmnožin množiny $X$. Tento systém se nazývá potenční množina množiny $X$ a značí se $\mathcal{P}(X)$. 
    \item \textbf{Axiom of infinity}: Existuje systém množin $\mathcal{S}$ takový, že $\emptyset\in\mathcal{S}$ a pro každou množinu $x\in\mathcal{S}$ platí $x\cup\{x\}\in\mathcal{S}$. 
    \item \textbf{Axiom of choice}: Pro každý systém neprázdných množin $\mathcal{S}$ existuje funkce $f$, která každé množině $X\in\mathcal{S}$ přiřazuje prvek $f(X)$, který je prvkem množiny $X$.
    \item \textbf{Axiom schema of replacement}: Pro každou množinu $X$ a každou funkci $f$ existuje množina, která obsahuje všechny prvky $f(x)$ pro každý prvek $x\in X$.
  \end{enumerate}
  V poslední řadě chci dodat, že na přesných axiomech množinové teorie se zdroje moc neshodují a můžete najít různé formulace. Nám výše uvedené axiomy postačí pro naše účely. Pojďme si nyní podrobněji vysvětlit jednotlivé axiomy a ukázat si nějaké příklady.
\end{definitionbox}

\subsection{Axiom of extension}
K tomuto axiomu moc slov nepotřebujeme. Vezměme například tři množiny definované následovně\footnote{Všimněte si, že množiny zde zapisujeme pomocí závorek a prvky množiny oddělujeme čárkami.}:
\begin{align*}
  X &= \{1, 2, 3\},\\
  Y &= \{1, 2, 3\},\\
  Z &= \{1, 2, 3, 4\}.
\end{align*}
Množiny $X$ a $Y$ jsou rovny, protože obsahují stejné prvky, zatímco množina $Z$ je odlišná, protože obsahuje navíc prvek $4$.
\subsection{Axiom of the empty set}
I tento axiom je poměrně jednoduchý. Prázdná množina je množina, která neobsahuje žádné prvky. Můžeme ji zapsat následovně: $\emptyset = \{\}$.
\subsection{Axiom schema of separation}
Tento axiom nám říká, že pro každou množinu $x$ a každou vlastnost $P$ existuje množina, která obsahuje všechny prvky množiny $x$ splňující vlastnost $P$. Například, pokud máme množinu $X = \{1, 2, 3, 4, 5\}$ a vlastnost $P$, že prvek $x$ je sudý, můžeme vytvořit množinu $Y = \{2, 4\}$.
\subsection{Axiom of pairing}
Tento axiom nám říká, že pro každé dvě množiny $X$ a $Y$ existuje množina, jejímiž prvky jsou právě množiny $X$ a $Y$ (ne jejich prvky, to dělá až axiom sjednocení). Například, pokud máme množiny $X = \{1, 2\}$ a $Y = \{3, 4\}$, můžeme vytvořit množinu $Z = \{\{1, 2\}, \{3, 4\}\}$. Tento axiom využijeme při důležité definici kartézského součinu množin.
\subsection{Axiom of union}
Tento axiom nám říká, že pro každý systém množin $\mathcal{S}$ existuje množina, která obsahuje všechny prvky všech množin v systému $\mathcal{S}$. Například, pokud máme množiny $X = \{1, 2\}$ a $Y = \{2, 3\}$, můžeme vytvořit množinu $Z = X\cup Y = \{1, 2, 3\}$.
\subsection{Axiom of power set}
Tento axiom nám říká, že pro každou množinu $X$ existuje systém všech podmnožin množiny $X$. Například, pokud máme množinu $X = \{1, 2\}$, můžeme vytvořit množinu $\mathcal{P}(X) = \{\emptyset, \{1\}, \{2\}, \{1, 2\}\}$. Této množině se říká potenční množina množiny $X$.
\subsection{Axiom of infinity} \label{axiom_of_infinity}
Tento axiom nám říká, že existuje systém množin $\mathcal{S}$ takový, že $\emptyset\in\mathcal{S}$ a pro každou množinu $x\in\mathcal{S}$ platí $x\cup\{x\}\in\mathcal{S}$. Tento axiom nám zaručuje existenci nekonečných množin. Pomocí tohoto axiomu můžeme definovat například množinu přirozených čísel $\mathbb{N}$ a to následovně:
\begin{align*}
  0 &= \emptyset,\\
  1 &= \{0\} = \{\emptyset\},\\
  2 &= \{0, 1\} = \{\emptyset, \{\emptyset\}\},\\
  3 &= \{0, 1, 2\} = \{\emptyset, \{\emptyset\}, \{\emptyset, \{\emptyset\}\}\},\\
  &\vdots
\end{align*}
\subsection{Axiom of choice}
Tento axiom je nejkontroverznější z celé množinové teorie, protože pokud ho nepřijmeme, spousta zajímavých matematických konceptů se nám zhroutí. Bez něj například nejde dokázat, že každý vektorový prostor má bázi, a i některé věty z analýzy potřebují alespoň jeho slabší verzi (například když z nekonečně mnoha množin potřebujeme vybrat posloupnost prvků). Obyčejné limity, se kterými budeme počítat, se bez něj ale obejdou. K tomuto axiomu se ještě vrátíme, až se budeme zabývat konkrétními matematickými koncepty. Známým paradoxem spojeným s tímto axiomem je Banachův--Tarského paradox, který říká, že kouli můžeme rozložit na konečně mnoho kusů a ty pak poskládat do dvou koulí, z nichž každá má stejný objem jako ta původní. 
\subsection{Axiom schema of replacement}
Tento axiom nám říká, že pro každou množinu $X$ a každou funkci $f$ existuje množina, která obsahuje všechny prvky $f(x)$ pro každý prvek $x\in X$. Například, pokud máme množinu $X = \{1, 2, 3\}$ a funkci $f(x) = x^2$\footnote{Zde předpokládáme, že víte, co dělá základní operace na druhou.}, můžeme vytvořit množinu $Y = \{1, 4, 9\}$.

\section{Množiny}
Nebudeme se zde zabývat, co přímo množina je, protože bychom se pohybovali na hranici mezi matematikou a filozofií. Množinou tedy rozumíme určitou kolekci prvků, které mohou být cokoliv. Pokud jste nečetli předchozí odbočku, zopakujme si nyní základní definice a pojmy spojené s množinami.
\begin{definitionbox}
  \textbf{Značení}: Množiny značíme velkými písmeny, například $X$, $Y$, $Z$. Prvky množin značíme malými písmeny, například $x$, $y$, $z$.\\
  \textbf{Prvek množiny}: Prvek $x$ je prvkem množiny $X$, pokud $x\in X$. Opak platí, pokud $x\notin X$.\\
  \textbf{Prázdná množina}: Množina, která neobsahuje žádné prvky. Značíme ji $\emptyset$. \\
  \textbf{Podmnožina}: Množina $Y$ je podmnožinou množiny $X$, pokud každý prvek množiny $Y$ je také prvkem množiny $X$. Značíme to $Y\subseteq X$. \\
  \textbf{Rovnost množin}: Množiny $X$ a $Y$ jsou rovny, pokud obsahují stejné prvky. Značíme to $X = Y$. \\
  \textbf{Sjednocení množin}: Sjednocení množin $X$ a $Y$ je množina, která obsahuje všechny prvky množin $X$ a $Y$. Značíme to $X\cup Y$. \\
  \textbf{Průnik množin}: Průnik množin $X$ a $Y$ je množina, která obsahuje pouze ty prvky, které jsou zároveň v množinách $X$ a $Y$. Značíme to $X\cap Y$. \\
  \textbf{Rozdíl množin}: Rozdíl množin $X$ a $Y$ je množina, která obsahuje pouze ty prvky, které jsou v množině $X$ a nejsou v množině $Y$. Značíme to $X\setminus Y$. \\
  \textbf{Potenční množina}: Potenční množina množiny $X$ je množina všech podmnožin množiny $X$. Značíme to $\mathcal{P}(X)$. \\
\end{definitionbox}

% TODO obrazky?!

\begin{example}
  Mějme množiny $X = \{1, 2, 3\}$ a $Y = \{2, 3, 4\}$. Potom platí:
  \begin{align*}
    X\cup Y &= \{1, 2, 3, 4\},\\
    X\cap Y &= \{2, 3\},\\
    X\setminus Y &= \{1\},\\
    Y\setminus X &= \{4\}. \\
    \mathcal{P}(X) &= \{\emptyset, \{1\}, \{2\}, \{3\}, \{1, 2\}, \{1, 3\}, \{2, 3\}, \{1, 2, 3\}\}.
  \end{align*}
\end{example}

\begin{problem}
  Mějme množiny $X = \{1, 2, 3, 4\}$ a $Y = \{3, 4, 5\}$. Zjistěte $X\cup Y$, $X\cap Y$, $X\setminus Y$, $Y\setminus X$ a $\mathcal{P}(Y)$.
\end{problem}

\begin{solution}
  Máme:
  \begin{align*}
    X\cup Y &= \{1, 2, 3, 4, 5\},\\
    X\cap Y &= \{3, 4\},\\
    X\setminus Y &= \{1, 2\},\\
    Y\setminus X &= \{5\},\\
    \mathcal{P}(Y) &= \{\emptyset, \{3\}, \{4\}, \{5\}, \{3, 4\}, \{3, 5\}, \{4, 5\}, \{3, 4, 5\}\}.
  \end{align*}
\end{solution}

% TODO jeste sem pls zajeb nejake priklady setu

\section{Relace}
Když už nyní víme, co jsou množiny, je na čase se zabývat vztahy mezi nimi. Takovým vztahům se říká relace a jsou základem pro spoustu dalších matematických koncept. Relace můžeme chápat jako způsob, jakým jsou prvky množin propojeny. Nejprve si zopakujme základní definice a pojmy spojené s relacemi. Než si řekneme, co je to relace, musíme nejdříve definovat kartézský součin množin.

\begin{definitionbox}
  \textbf{Kartézský součin množin}: Kartézský součin množin $X$ a $Y$ je množina všech uspořádaných dvojic $(x, y)$, kde $x\in X$ a $y\in Y$. Značíme to $X\times Y$. 
\end{definitionbox}

\begin{example}
  Mějme množiny $X = \{1, 2\}$ a $Y = \{3, 4\}$. Potom platí:
  \begin{align*}
    X\times Y &= \{(1, 3), (1, 4), (2, 3), (2, 4)\}.
  \end{align*}
\end{example}

\begin{definitionbox}
  \textbf{Relace}: Relace $\mathcal R$ mezi množinami $X$ a $Y$ je podmnožina kartézského součinu $X\times Y$. Pokud máme relaci $\mathcal R$ mezi množinami $X$ a $Y$, značíme to $\mathcal R\subseteq X\times Y$. Prvek $(x, y)\in R$ značíme $x\mathcal Ry$ a říkáme, že prvek $x$ je v relaci $R$ s prvkem $y$. Je to ve zkratce jakási omezující vlastnost na kartézský součin množin. 
\end{definitionbox}

\begin{example}
  Mějme množiny $X = \{1, 2\}$ a $Y = \{3, 4\}$ a relaci $\mathcal R = \{(1, 3), (2, 4)\}$. Potom platí:
  \begin{align*}
    1\mathcal R3,\\
    2\mathcal R4.
  \end{align*}
V tomto případě máme relaci zadanou explicitně a slouží jen na ukázání poněkud zvláštního zápisu\footnote{Tento zápis, jak brzy uvidíte, moc používat nebudeme.}.
\end{example}

\begin{example}
  Nyní si pojďme ukázat příklad, kde je relace zadaná jako nějaká vlastnost, kterou prvky množin splňují. Mějme množiny $X = \{1, 2, 3\}$ a $Y = \{3, 4, 5\}$ a relaci $\mathcal R = \{(x, y)\in X\times Y\ |\ x < y\}$. Což v překladu znamená, že prvek $x$ je v relaci s prvkem $y$, pokud $x$ je menší než $y$. Potom platí:
  \begin{align*}
    1\mathcal R3,\\
    1\mathcal R4,\\
    1\mathcal R5,\\
    2\mathcal R3,\\
    2\mathcal R4,\\
    2\mathcal R5,\\
    3\mathcal R4,\\
    3\mathcal R5.
  \end{align*}

\end{example}

\begin{problem}
  Mějme množiny $X = \{1, 2, 3\}$ a $Y = \{3, 4, 5\}$ a relaci $\mathcal R = \{(x, y)\in X\times Y\ |\ x + y = 5\}$. Zjistěte, které dvojice prvků množin $X$ a $Y$ splňují vlastnost $x + y = 5$.
\end{problem}

\begin{solution}
  Máme:
  \begin{align*}
    1\mathcal R4,\\
    2\mathcal R3.
  \end{align*}
\end{solution}

\begin{note}
  Jak již bylo zmíněno v definice kartézského součinu, prvek $(x, y)$ je uspořádaná dvojice, což znamená, že pořadí prvků má význam. To znamená, že ve většině případů $(x, y)\neq(y, x)$. 
\end{note}

\begin{definitionbox}
  \textbf{Vlastnosti relací}. \\
    Nechť $\mathcal R$ je relace na množině $X$. Potom můžeme definovat následující vlastnosti relací: \vspace{5mm}

    \textbf{Reflexivita}: Relace $\mathcal R$ je reflexivní, pokud pro každý prvek $x\in X$ platí $x\mathcal Rx$.\ \vspace{1mm}

    \textbf{Symetrie}: Relace $\mathcal R$ je symetrická, pokud pro každý prvek $x, y\in X$ platí, že pokud $x\mathcal Ry$, tak $y\mathcal Rx$.\ \vspace{1mm}

    \textbf{Antisymetrie}: Relace $\mathcal R$ je antisymetrická, pokud pro každý prvek $x, y\in X$ platí, že pokud $x\mathcal Ry$ a $y\mathcal Rx$, tak $x = y$.\ \vspace{1mm}

    \textbf{Tranzitivita}: Relace $\mathcal R$ je tranzitivní, pokud pro každý prvek $x, y, z\in X$ platí, že pokud $x\mathcal Ry$ a $y\mathcal Rz$, tak $x\mathcal Rz$.\ \vspace{1mm}

  Pokud je relace reflexivní, symetrická a tranzitivní, říkáme, že je \textbf{relace ekvivalence}.
\end{definitionbox}

Vlastnosti relací se nejlépe ukazující na číselných množinách a proto si nejprve pojďme říct něco o nich.

\section{Odbočka --- Tělesa}
Těleso je ve zkratce množina prvků, která splňuje `hezké' vlastnosti. Jdou mezi sebou `hezky' sčítat a násobit. Pojďme si nyní zadefinovat co přesně těleso je.

\begin{definitionbox}
  \textbf{Těleso} je zadáno množinou $\mathbb{F}$ a obsahuje dvě funkce, které jdou abstraktně zadefinovat, ale pro nás bude stačit, že to je sčítaní a násobení. Těleso musí splňovat následující axiomy:

  \begin{enumerate}
    \item Axiomy pro sčítání skalárů (Skaláry jsou prvky tělesa. Například v případě $\mathbb{R}$ jsou skaláry reálná čísla. To že jsou reálná čísla těleso si řekneme později.):
    \begin{enumerate}
      \item Existence nulového prvku: Existuje prvek $0\in\mathbb{F}$ takový, že pro každý prvek $x\in\mathbb{F}$ platí $x + 0 = x$. Tento prvek se nazývá nulový prvek.
      \item Komutativita sčítání: Pro každé dva prvky $x, y\in\mathbb{F}$ platí, že $x + y = y + x$.
      \item Asociativita sčítání: Pro každé tři prvky $x, y, z\in\mathbb{F}$ platí, že $(x + y) + z = x + (y + z)$.
      \item Existence opačného prvku: Pro každý prvek $x\in\mathbb{F}$ existuje prvek $-x\in\mathbb{F}$ takový, že $x + (-x) = 0$.
    \end{enumerate}
    \item Axiomy pro násobení skalárů:
    \begin{enumerate}
      \item Existence jednotkového prvku: Existuje prvek $1\in\mathbb{F}$, $1\neq 0$, takový, že pro každý prvek $x\in\mathbb{F}$ platí $x\cdot 1 = x$. Tento prvek se nazývá jednotkový prvek.
      \item Komutativita násobení: Pro každé dva prvky $x, y\in\mathbb{F}$ platí, že $x\cdot y = y\cdot x$.
      \item Asociativita násobení: Pro každé tři prvky $x, y, z\in\mathbb{F}$ platí, že $(x\cdot y)\cdot z = x\cdot (y\cdot z)$.
    \end{enumerate}
    \item Distributivní zákon: Pro každé tři prvky $x, y, z\in\mathbb{F}$ platí, že $x\cdot (y + z) = x\cdot y + x\cdot z$.
    \item Existence inverzního prvku: Pro každý nenulový prvek $x\in\mathbb{F}$ existuje prvek $x^{-1}\in\mathbb{F}$ takový, že $x\cdot x^{-1} = 1$. Tento prvek se nazývá inverzní prvek.
  \end{enumerate}
  

\end{definitionbox}

\section{Čísla}
O číslech jste pravděpodobně již slyšeli, ale stejně je nutné si představit základní číselné množiny, se kterými budeme pracovat. 

\subsection{Přirozená čísla}
Přirozená čísla, jak jejich název napovídá, jsou čísla, která se nejčastěji vyskutují v našem životě. Je to také první množina, která se považovala za číselnou množinu po překvapivě dlouhou dobu.\footnote{Nula byla často rozporuplná a někteří matematikové ji neuznávali jako číslo a jejich argument byl, že není potřeba číslo, které značí nic.} 
\begin{definitionbox}
  \textbf{Přirozená čísla} jsou množina čísel, která začínají od jedničky a pokračují dál. Značíme je $\mathbb{N} = \{1, 2, 3, \ldots\}$. Pokud si vzpomínáte na odbočku o axiomech, tak konkrétně v~\ref{axiom_of_infinity} jsme si ukázali, jak můžeme přirozená čísla definovat pomocí množinové teorie.

  V případě, že chceme zahrnout nulu, mluvíme o \textbf{přirozených číslech včetně nuly} a značíme je $\mathbb{N}_0 = \{0, 1, 2, 3, \ldots\}$.
\end{definitionbox}

\begin{problem}
  Jsou přirozená čísla těleso?
\end{problem}

\begin{hint}
  Zkuste si projít axiomy tělesa a zjistit, zda přirozená čísla splňují všechny axiomy.
\end{hint}

\begin{solution}
  Nejsou.

  \begin{proof}
    Přirozená čísla nesplňují axiomy tělesa, protože například neexistuje inverzní prvek pro sčítání. Pokud vezmeme libovolné přirozené číslo $n$, tak jeho inverzní prvek by měl být $-n$, ale toto číslo není přirozené.
  \end{proof}
\end{solution}

\subsection{Celá čísla}
\begin{definitionbox}
  \textbf{Celá čísla} jsou množina čísel, která zahrnují přirozená čísla, ale také záporná čísla a nulu. Značíme je $\mathbb{Z} = \{\ldots, -3, -2, -1, 0, 1, 2, 3, \ldots\}$.
\end{definitionbox}

\begin{problem}
  Jsou celá čísla těleso?
\end{problem}

\begin{solution}
  Nejsou.

  \begin{proof}
    Celá čísla nesplňují axiomy tělesa, protože například neexistuje inverzní prvek pro násobení. Pokud vezmeme libovolné celé číslo $n$, tak jeho inverzní prvek by měl být $\frac{1}{n}$, ale toto číslo není celé.
  \end{proof}
\end{solution}

\subsection{Racionální čísla}
\begin{definitionbox}
  \textbf{Racionální čísla} jsou množina čísel, která zahrnují celá čísla, ale také zlomky. Značíme je $\mathbb{Q} = \left\{\frac{a}{b}\ |\ a, b\in\mathbb{Z}, b\neq 0\right\}$.\ \vspace{2mm}

  Tato definice je trochu složitější, ale znamená to, že racionální čísla jsou všechna čísla, která můžeme zapsat jako zlomek, kde čitatel i jmenovatel jsou celá čísla a jmenovatel není nula. 
\end{definitionbox}

\begin{problem}
  Jsou racionální čísla těleso?
\end{problem}

\begin{solution}
  Ano.

  \begin{proof}
    Racionální čísla splňují axiomy tělesa, protože jsou uzavřená na sčítání, násobení, existuje nulový prvek, existuje opačný prvek, existuje jednotkový prvek, jsou komutativní a asociativní. Pokud tomu nevěříte, zkuste si projít axiomy tělesa a zjistit, zdali racionální čísla splňují všechny axiomy.
  \end{proof}
\end{solution}

\subsection{Reálná čísla}
\begin{definitionbox}
  \textbf{Reálná čísla} jsou množina čísel, která zahrnují racionální čísla, ale také iracionální čísla. Značíme je $\mathbb{R}$. Přesnou definici zde nebudeme uvádět, protože je složitější.
\end{definitionbox}

Proč ale existují reálná čísla? Proč nestačí racionální čísla? Odpověď je jednoduchá, ale zároveň složitá. Ukažme si příklad, který nám ukáže číslo, které není racionální.

\begin{example}
  Mějme čtverec o stranách délky $1$. Jaká je délka úhlopříčky čtverce? Pokud si vzpomenete na Pythagorovu větu\footnote{K této větě se ještě dostaneme, takže pokud nevíte o co jde, tak stačí, když mi budete věřit.}, tak délka úhlopříčky je $\sqrt{1^2 + 1^2} = \sqrt{2}$. Nyní si pojďme ukázat, že číslo $\sqrt{2}$ není racionální.

  \begin{proof}
    Sporem. Předpokládejme, že číslo $\sqrt{2}$ je racionální. To znamená, že existují nesoudělná\footnote{Nemají žádného společného dělitele.} celá čísla $a$ a $b$ taková, že $\sqrt{2} = \frac{a}{b}$. Potom můžeme napsat:
    \begin{align*}
      \sqrt{2} &= \frac{a}{b},\\
      2 &= \frac{a^2}{b^2},\\
      2b^2 &= a^2.
    \end{align*}
    To znamená, že $a^2$ je sudé číslo, což znamená, že $a$ je také sudé číslo. Toto tvrzení víme z příkladu \ref{ex:1.4.5}. Takže existuje celé číslo $c$ takové, že $a = 2c$. Dosadíme zpět do rovnice:
    \begin{align*}
      2b^2 &= {(2c)}^2,\\ 
      2b^2 &= 4c^2,\\
      b^2 &= 2c^2.
    \end{align*}
    To znamená, že $b^2$ je sudé číslo, což znamená, že $b$ je také sudé číslo. To ale znamená, že čísla $a$ a $b$ mají společného dělitele a došli jsme ke sporu. Tedy jsme dokázali, že číslo $\sqrt{2}$ není racionální.
  \end{proof}
\end{example}

Když tedy $\sqrt{2}$ není racionální číslo, co tedy je? Těmto číslům se říká iracionální čísla a je to podmnožina čísel reálných. Velice známým iracionálním číslem je $\pi$ nebo $e$.

\begin{problem}
  Ukažte, že realná čísla jsou těleso.
\end{problem}

\begin{note}
  Reálná čísla jsou takzvaně nespočetná množina. To zjednodušeně znamená, že jich je více než přirozených čísel. Těch je ale také nekonečně mnoho. Ve zkratce, počet prvků množiny reálných čísel je větší nekonečno než počet prvků množiny přirozených čísel. Více se tím zde zabývat nebudeme.
\end{note}

\subsection{Komplexní čísla}
\begin{definitionbox}
  \textbf{Komplexní čísla} jsou množina čísel, která zahrnují reálná čísla, ale také imaginární čísla. Značíme je $\mathbb{C} = \{a + bi\ |\ a, b\in\mathbb{R}\}$, kde $i$ je imaginární jednotka, která je definována vztahem $i^2 = -1$ (neformálně se píše $i = \sqrt{-1}$).
\end{definitionbox}

Proč potřebujeme další číselnou množinu? Odpověď je opět jednoduchá, ale zároveň složitá. Komplexní čísla nám umožňují řešit rovnice, které nemají řešení v reálných číslech. Například rovnice $x^2 + 1 = 0$ nemá řešení v reálných číslech, protože v nich $\sqrt{-1}$ není definované, ale má řešení v komplexních číslech.

% TODO add chapter here
Ke komplexním číslům se ještě vrátíme, protože jsou velice zajímavé a důležité. Nyní si jen pojďme uvědomit množinový vztah mezi výše zmíněnými číselnými množinami.
\begin{equation}
  \mathbb{N} \subset \mathbb{Z} \subset \mathbb{Q} \subset \mathbb{R} \subset \mathbb{C}
\end{equation}
\section{Příběh --- Dopis, který zbořil dům}

Gottlob Frege, profesor matematiky v Jeně, strávil přes dvacet let tím, že stavěl celou aritmetiku na čisté logice. Jeho systém počítal s tím, čemu se dnes říká \textit{naivní teorie množin}: každá vlastnost určuje množinu, totiž množinu všech věcí, které tu vlastnost mají. V červnu 1902, když se tiskl druhý díl jeho životního díla \textit{Základní zákony aritmetiky}, mu přišel dopis od mladého Angličana Bertranda Russella.

Russell se ptal na zdánlivě nevinnou věc. Vezměme vlastnost `nebýt prvkem sebe sama' a množinu
\begin{equation*}
  R = \{x \mid x\notin x\}.
\end{equation*}
Je $R$ prvkem sebe sama? Pokud $R\in R$, pak $R$ tu vlastnost nemá, a tedy $R\notin R$. Pokud $R\notin R$, pak tu vlastnost má, a tedy $R\in R$. Dostáváme $R\in R\Leftrightarrow R\notin R$, což je spor. Lidová verze zní: vesnický holič holí právě ty muže, kteří se neholí sami. Holí se holič sám?

Frege okamžitě pochopil, že spor je přímo v základech jeho systému. Do dodatku ke knize, která už byla v tisku, napsal: `Vědeckému spisovateli se sotva může přihodit něco nepříjemnějšího, než když se mu po dokončení díla zhroutí jeden ze základů jeho stavby. Do takové situace mě dostal dopis pana Bertranda Russella, právě když se tisk tohoto svazku chýlil ke konci.' Svůj program už nikdy nezachránil. Zahořkl a jeho deník z posledních let života, zveřejněný až v roce 1994, obsahuje otevřeně antisemitské a protidemokratické názory. Velký logik, smutný člověk.

Z Russellova paradoxu se matematika poučila. Proto axiomatická teorie množin ze začátku této kapitoly nedovoluje vytvořit množinu z libovolné vlastnosti. Schéma axiomů vydělení dovoluje jen \textit{vybrat} z již existující množiny ty prvky, které danou vlastnost mají. `Množina všech množin' tak vůbec neexistuje a paradox nejde ani zapsat. Základ tohoto systému navrhl v roce 1908 Ernst Zermelo, mimo jiné právě jako odpověď na paradoxy, jako byl ten Russellův.

Russell sám se později stal slavným filozofem a pacifistou. V roce 1918 byl za protiválečný článek odsouzen k šesti měsícům vězení, kde napsal knihu \textit{Úvod do matematické filozofie}. V roce 1950 dostal Nobelovu cenu za literaturu.

\koneckapitoly

\chapter{Mocniny, odmocniny a algebraické výrazy}
\mathquote{Algebra is generous; she often gives more than is asked of her.}{Jean le Rond d'Alembert}

V předchozí kapitole jsme si postavili číselné množiny. Teď s čísly konečně začneme pořádně počítat. Tato kapitola je takovou `posilovnou' celé maturity --- mocniny, odmocniny a úpravy výrazů se objeví úplně všude, od rovnic přes funkce až po derivace. Pokud zde něco nebude sedět, bude to bolet později. Budeme se tedy snažit každé pravidlo nejen uvést, ale i ukázat, \textit{proč} platí. Většina těchto pravidel totiž není nic, co by se muselo biflovat, ale přirozeně vyplývá z definic.

\section{Mocniny s přirozeným exponentem}

Začněme tím nejjednodušším. Pokud chceme číslo vynásobit samo sebou několikrát, bylo by nepraktické to vždy vypisovat. Místo $2\cdot 2\cdot 2\cdot 2\cdot 2$ je mnohem příjemnější napsat $2^5$.

\begin{definitionbox}
  \textbf{Mocnina s přirozeným exponentem}: Pro každé $a\in\mathbb{R}$ a $n\in\mathbb{N}$ definujeme
  \begin{equation*}
    a^n = \underbrace{a\cdot a\cdot\ \cdots\ \cdot a}_{n\text{-krát}}.
  \end{equation*}
  Číslo $a$ nazýváme \textbf{základ} (mocněnec) a číslo $n$ nazýváme \textbf{exponent} (mocnitel).

  \vspace{0.2cm}

  Pro ty, kterým tečky v definici nejsou dost formální\footnote{Tečky $\cdots$ jsou totiž v matematice trochu podvod. Předpokládají, že čtenář `ví, jak to pokračuje'.}, můžeme mocninu zadefinovat i \textbf{rekurzivně}:
  \begin{align*}
    a^1 &= a,\\
    a^{n+1} &= a^n\cdot a \quad\text{pro } n\in\mathbb{N}.
  \end{align*}
\end{definitionbox}

Rekurzivní definice je ve skutečnosti velice podobná matematické indukci z první kapitoly. Řekneme, co se stane pro $n = 1$, a pak řekneme, jak se z $n$ dostat na $n+1$. To nám zaručí, že mocnina je definovaná pro všechna přirozená čísla.

\begin{example}
  Spočítejme několik mocnin:
  \begin{align*}
    3^4 &= 3\cdot 3\cdot 3\cdot 3 = 81,\\
    {(-2)}^3 &= (-2)\cdot(-2)\cdot(-2) = -8,\\
    {(-2)}^4 &= (-2)\cdot(-2)\cdot(-2)\cdot(-2) = 16,\\
    -2^4 &= -(2\cdot 2\cdot 2\cdot 2) = -16,\\
    {\left(\tfrac{1}{2}\right)}^3 &= \tfrac{1}{2}\cdot\tfrac{1}{2}\cdot\tfrac{1}{2} = \tfrac{1}{8}.
  \end{align*}
  Všimněte si rozdílu mezi ${(-2)}^4$ a $-2^4$. Mocnina má přednost před znaménkem mínus, takže $-2^4$ znamená `minus (dva na čtvrtou)'. Tohle je jedna z nejčastějších chyb u maturity, takže si na ni dávejte extra pozor.
\end{example}

\begin{note}
  Z příkladu výše je vidět, že záporné číslo umocněné na sudý exponent je kladné a na lichý exponent záporné. Je to proto, že každá dvojice mínusů se `vyruší'. Zkuste si to dokázat matematickou indukcí.
\end{note}

Teď přijdou pravidla pro počítání s mocninami. Ta vypadají na první pohled jako něco, co se musí naučit nazpaměť, ale ve skutečnosti jsou to jen důsledky toho, že počítáme, kolik máme kde činitelů.

\begin{definitionbox}
  \textbf{Pravidla pro počítání s mocninami}: Pro každé $a, b\in\mathbb{R}$ a $m, n\in\mathbb{N}$ platí:
  \begin{enumerate}
    \item $a^m\cdot a^n = a^{m+n}$,
    \item $\dfrac{a^m}{a^n} = a^{m-n}$, pokud $a\neq 0$ a $m > n$,
    \item ${(a^m)}^n = a^{m\cdot n}$,
    \item ${(a\cdot b)}^n = a^n\cdot b^n$,
    \item ${\left(\dfrac{a}{b}\right)}^n = \dfrac{a^n}{b^n}$, pokud $b\neq 0$.
  \end{enumerate}
\end{definitionbox}

\begin{example}\label{ex:mocniny-soucin}
  Dokažme si první pravidlo. Nejdříve neformálně, tak jak by to udělal každý rozumný člověk:
  \begin{equation*}
    a^m\cdot a^n = \underbrace{a\cdot\ \cdots\ \cdot a}_{m\text{-krát}}\cdot\underbrace{a\cdot\ \cdots\ \cdot a}_{n\text{-krát}} = \underbrace{a\cdot\ \cdots\ \cdot a}_{(m+n)\text{-krát}} = a^{m+n}.
  \end{equation*}
  A nyní formálně, pomocí matematické indukce a rekurzivní definice.

  \begin{proof}
    Zvolme pevně $m\in\mathbb{N}$ a indukcí podle $n$ dokážeme, že $a^m\cdot a^n = a^{m+n}$.

    \textbf{Základní krok}: Pro $n = 1$ máme $a^m\cdot a^1 = a^m\cdot a = a^{m+1}$, což je přesně rekurzivní definice.

    \textbf{Indukční krok}: Předpokládejme, že pro nějaké $n$ platí $a^m\cdot a^n = a^{m+n}$.

    \textbf{Závěr}: Potom
    \begin{equation*}
      a^m\cdot a^{n+1} = a^m\cdot(a^n\cdot a) = (a^m\cdot a^n)\cdot a = a^{m+n}\cdot a = a^{m+n+1},
    \end{equation*}
    kde jsme použili postupně definici, asociativitu násobení (jeden z axiomů tělesa!), indukční předpoklad a opět definici.
  \end{proof}
\end{example}

\begin{problem}\label{pr:mocniny-mocnina}
  Dokažte (alespoň neformálně) pravidla 3 a 4.
\end{problem}

\begin{hint}
  U pravidla 3 si rozmyslete, že ${(a^m)}^n$ znamená, že máme $n$ `balíčků' a v každém z nich je $m$ činitelů $a$. U pravidla 4 si pomozte komutativitou násobení a přeskládejte činitele.
\end{hint}

\begin{solution}
  Pravidlo 3:
  \begin{equation*}
    {(a^m)}^n = \underbrace{a^m\cdot a^m\cdot\ \cdots\ \cdot a^m}_{n\text{-krát}} = a^{\overbrace{m + m + \cdots + m}^{n\text{-krát}}} = a^{m\cdot n},
  \end{equation*}
  kde jsme opakovaně použili pravidlo 1.

  Pravidlo 4:
  \begin{equation*}
    {(ab)}^n = \underbrace{(ab)\cdot(ab)\cdot\ \cdots\ \cdot(ab)}_{n\text{-krát}} = \underbrace{a\cdot\ \cdots\ \cdot a}_{n\text{-krát}}\cdot\underbrace{b\cdot\ \cdots\ \cdot b}_{n\text{-krát}} = a^n b^n.
  \end{equation*}
\end{solution}

\begin{problem}
  Zjednodušte následující výrazy:
  \begin{enumerate}[label=(\alph*)]
    \item $x^3\cdot x^5\cdot x$,
    \item ${(2a^2b)}^3$,
    \item $\dfrac{{(x^2y)}^4}{x^5y^2}$, kde $x, y\neq 0$,
    \item $\dfrac{3^{10}\cdot 9^2}{27^4}$.
  \end{enumerate}
\end{problem}

\begin{solution}
  \begin{enumerate}[label=(\alph*)]
    \item $x^{3+5+1} = x^9$,
    \item $2^3\cdot{(a^2)}^3\cdot b^3 = 8a^6b^3$,
    \item $\dfrac{x^8y^4}{x^5y^2} = x^3y^2$,
    \item Vše převedeme na základ $3$: $\dfrac{3^{10}\cdot{(3^2)}^2}{{(3^3)}^4} = \dfrac{3^{10}\cdot 3^4}{3^{12}} = 3^{2} = 9$.
  \end{enumerate}
  Poslední příklad je typický maturitní trik. Pokud vidíte různé základy, které jsou mocninami stejného čísla, vždy se je vyplatí převést na společný základ.
\end{solution}

\begin{definitionbox}
  Upozorňujeme, že pro sčítání \textbf{žádné} podobné pravidlo neexistuje! Obecně
  \begin{equation*}
    {(a + b)}^n \neq a^n + b^n \qquad\text{a}\qquad a^m + a^n \neq a^{m+n}.
  \end{equation*}
  Stačí si vzít $a = b = 1$ a $n = 2$. Potom ${(1 + 1)}^2 = 4$, ale $1^2 + 1^2 = 2$. Tomuto chybnému `pravidlu' se v anglicky mluvících zemích přezdívá \textit{freshman's dream}, tedy sen prváka.
\end{definitionbox}

\section{Mocniny s celým exponentem}

Doposud jsme měli v exponentu jen přirozená čísla. Co by ale mělo znamenat $2^0$ nebo $2^{-3}$? Vynásobit dvojku samu sebou nulakrát nebo dokonce mínus třikrát? To zní jako nesmysl. Zde přichází na řadu typický matematický přístup: chceme, aby naše pravidla platila i nadále, a podle toho si novou definici \textit{vynutíme}.

Podívejme se na pravidlo $\frac{a^m}{a^n} = a^{m-n}$. Pokud bychom chtěli, aby fungovalo i pro $m = n$, musí platit
\begin{equation*}
  1 = \frac{a^n}{a^n} = a^{n-n} = a^0.
\end{equation*}
Podobně pokud by mělo platit pro $m = 0$, musí platit
\begin{equation*}
  \frac{1}{a^n} = \frac{a^0}{a^n} = a^{0-n} = a^{-n}.
\end{equation*}

\begin{definitionbox}
  \textbf{Mocnina s nulovým a záporným exponentem}: Pro každé $a\in\mathbb{R}\setminus\{0\}$ a $n\in\mathbb{N}$ definujeme
  \begin{align*}
    a^0 &= 1,\\
    a^{-n} &= \frac{1}{a^n}.
  \end{align*}
  Při této definici platí všech pět pravidel z předchozí sekce pro všechna $m, n\in\mathbb{Z}$ (pokud nedělíme nulou) a pravidlo pro podíl už nepotřebuje podmínku $m > n$.
\end{definitionbox}

\begin{note}
  Výraz $0^0$ jsme nechali nedefinovaný. V některých oblastech matematiky (například v kombinatorice) se dodefinovává jako $0^0 = 1$, protože je to praktické. V analýze se ale objevuje jako takzvaný neurčitý výraz, protože se k němu můžeme `blížit' různými způsoby a dostávat různé výsledky. U maturity tento výraz považujte za nedefinovaný.
\end{note}

\begin{example}
  Spočítejme:
  \begin{align*}
    5^{-2} &= \frac{1}{5^2} = \frac{1}{25},\\
    {\left(\frac{2}{3}\right)}^{-3} &= \frac{1}{{\left(\frac{2}{3}\right)}^3} = \frac{1}{\frac{8}{27}} = \frac{27}{8} = {\left(\frac{3}{2}\right)}^3,\\
    {(-7)}^0 &= 1,\\
    \frac{a^{-3}b^2}{a^2b^{-4}} &= a^{-3-2}b^{2-(-4)} = a^{-5}b^6 = \frac{b^6}{a^5}.
  \end{align*}
  Z druhého řádku si odneste užitečný trik: záporný exponent u zlomku zlomek `otočí', tedy ${\left(\frac{a}{b}\right)}^{-n} = {\left(\frac{b}{a}\right)}^{n}$.
\end{example}

\begin{problem}
  Zjednodušte a výsledek zapište bez záporných exponentů ($x, y \neq 0$):
  \begin{enumerate}[label=(\alph*)]
    \item ${(x^{-2}y^3)}^{-2}$,
    \item $\dfrac{2^{-1} + 2^{-2}}{2^{-3}}$,
    \item ${\left(\dfrac{x^{-1}}{3y^{2}}\right)}^{-2}\cdot\dfrac{1}{9x^{2}}$.
  \end{enumerate}
\end{problem}

\begin{solution}
  \begin{enumerate}[label=(\alph*)]
    \item $x^{4}y^{-6} = \dfrac{x^4}{y^6}$,
    \item Pozor, zde se jedná o součet, nemůžeme tedy jen sečíst exponenty. $\dfrac{\frac{1}{2} + \frac{1}{4}}{\frac{1}{8}} = \dfrac{\frac{3}{4}}{\frac{1}{8}} = 6$,
    \item ${\left(\dfrac{3y^2}{x^{-1}}\right)}^{2}\cdot\dfrac{1}{9x^2} = {(3xy^2)}^2\cdot\dfrac{1}{9x^2} = \dfrac{9x^2y^4}{9x^2} = y^4$.
  \end{enumerate}
\end{solution}

\subsection{Zápis čísla v desítkovém tvaru}

Mocniny s celým exponentem se nám hodí i v praxi, například ve fyzice a chemii. Velká a malá čísla se totiž zapisují pomocí takzvaného \textbf{vědeckého zápisu}.

\begin{definitionbox}
  Každé kladné reálné číslo $x$ lze jednoznačně zapsat ve tvaru
  \begin{equation*}
    x = a\cdot 10^n,
  \end{equation*}
  kde $1\leq a < 10$ a $n\in\mathbb{Z}$. Tomuto tvaru říkáme \textbf{vědecký} (také semilogaritmický) \textbf{zápis} čísla.
\end{definitionbox}

\begin{example}
  Rychlost světla je přibližně $300\,000\,000\ \si{m/s} = 3\cdot 10^8\ \si{m/s}$. Hmotnost elektronu je přibližně $0{,}000\,000\,000\,000\,000\,000\,000\,000\,000\,911\ \si{kg} = 9{,}11\cdot 10^{-31}\ \si{kg}$. Doufám, že je jasné, proč fyzici dávají přednost druhému zápisu.
\end{example}

\section{Odmocniny}

Odmocňování je opačná operace k mocnění. Pokud víme, že $x^2 = 9$, chceme se nějak dostat k číslu $x$. Zde ale narážíme na první problém: rovnice $x^2 = 9$ má dvě řešení, $3$ a $-3$. Která z nich by měla být `ta' odmocnina? Matematici se dohodli, že ta nezáporná.

\begin{definitionbox}
  \textbf{$n$-tá odmocnina}: Nechť $a \geq 0$ a $n\in\mathbb{N}$. Pak \textbf{$n$-tá odmocnina} z čísla $a$ je to \textbf{nezáporné} číslo $x$, pro které platí $x^n = a$. Značíme ji $\sqrt[n]{a}$. Pro $n = 2$ píšeme zkráceně $\sqrt{a}$ a mluvíme o druhé odmocnině.

  \vspace{0.2cm}

  Symbolicky:
  \begin{equation*}
    \sqrt[n]{a} = x \quad\Leftrightarrow\quad x\geq 0\ \land\ x^n = a.
  \end{equation*}
\end{definitionbox}

Na tom, že takové číslo $x$ vůbec existuje a je jen jedno, stojí kus teorie reálných čísel, který tu dokazovat nebudeme.\footnote{Existence plyne z takzvané úplnosti reálných čísel, jednoznačnost z toho, že funkce $x^n$ je na $[0, \infty)$ rostoucí. Obojí uvidíte v pozdějších kapitolách.} Mimochodem, v racionálních číslech by to neplatilo --- v první kapitole jsme dokázali, že žádné racionální číslo umocněné na druhou nedá $2$.

\begin{definitionbox}
  \textbf{Dvě velmi důležité poznámky}:
  \begin{enumerate}
    \item Odmocnina je \textbf{vždy nezáporná}. Tedy $\sqrt{9} = 3$, a ne $\pm 3$. Znaménko $\pm$ se objeví až při řešení rovnice $x^2 = 9$, ne při počítání odmocniny.
    \item Pro každé $x\in\mathbb{R}$ platí
    \begin{equation*}
      \sqrt{x^2} = |x|.
    \end{equation*}
    Například $\sqrt{{(-5)}^2} = \sqrt{25} = 5 = |-5|$. Pokud byste napsali $\sqrt{x^2} = x$, pro záporná $x$ by vám vyšlo záporné číslo, což je spor s bodem 1.
  \end{enumerate}
\end{definitionbox}

\begin{note}
  Pro liché $n$ se v některých učebnicích definuje i odmocnina ze záporného čísla, například $\sqrt[3]{-8} = -2$, protože ${(-2)}^3 = -8$. Je to rozumné, protože pro liché $n$ je takové číslo opět jediné. V tomto textu (a ve většině českých učebnic) budeme ale pro jednoduchost odmocňovat jen nezáporná čísla a místo $\sqrt[3]{-8}$ budeme psát $-\sqrt[3]{8}$.
\end{note}

\begin{definitionbox}
  \textbf{Pravidla pro počítání s odmocninami}: Pro každé $a, b\geq 0$ a $m, n, k\in\mathbb{N}$ platí:
  \begin{enumerate}
    \item $\sqrt[n]{a}\cdot\sqrt[n]{b} = \sqrt[n]{a\cdot b}$,
    \item $\dfrac{\sqrt[n]{a}}{\sqrt[n]{b}} = \sqrt[n]{\dfrac{a}{b}}$, pokud $b\neq 0$,
    \item ${\left(\sqrt[n]{a}\right)}^m = \sqrt[n]{a^m}$,
    \item $\sqrt[m]{\sqrt[n]{a}} = \sqrt[m\cdot n]{a}$,
    \item $\sqrt[n]{a^m} = \sqrt[n\cdot k]{a^{m\cdot k}}$.
  \end{enumerate}
\end{definitionbox}

\begin{example}\label{ex:odmocniny-soucin}
  Dokažme si první pravidlo. Krásně na něm uvidíme, jak se s definicí odmocniny pracuje.

  \begin{proof}
    Označme $x = \sqrt[n]{a}$ a $y = \sqrt[n]{b}$. Podle definice je $x, y\geq 0$, $x^n = a$ a $y^n = b$. Chceme ukázat, že $x\cdot y$ je $n$-tá odmocnina z $ab$. Podle definice tedy musíme ověřit dvě věci:
    \begin{itemize}
      \item $xy \geq 0$. To platí, protože součin dvou nezáporných čísel je nezáporný.
      \item ${(xy)}^n = ab$. To platí, protože podle pravidel pro mocniny je ${(xy)}^n = x^ny^n = ab$.
    \end{itemize}
    Tedy $xy = \sqrt[n]{ab}$, což jsme chtěli dokázat.
  \end{proof}

  Všimněte si, že jsme nic nepočítali. Jen jsme ověřili, že kandidát na odmocninu splňuje definici. Tento postup je v matematice velice častý.
\end{example}

\begin{problem}
  Stejným způsobem jako v příkladu~\ref{ex:odmocniny-soucin} dokažte pravidlo 4.
\end{problem}

\begin{hint}
  Označte $x = \sqrt[m]{\sqrt[n]{a}}$ a umocněte postupně na $m$ a pak na $n$.
\end{hint}

\begin{solution}
  \begin{proof}
    Nechť $x = \sqrt[m]{\sqrt[n]{a}}$. Podle definice je $x\geq 0$ a $x^m = \sqrt[n]{a}$. Umocněním na $n$ dostaneme ${(x^m)}^n = a$, tedy $x^{mn} = a$. Protože $x \geq 0$, je podle definice $x = \sqrt[mn]{a}$.
  \end{proof}
\end{solution}

\begin{definitionbox}
  Opět pozor, pro sčítání nic takového neplatí:
  \begin{equation*}
    \sqrt{a + b} \neq \sqrt{a} + \sqrt{b}.
  \end{equation*}
  Protipříklad: $\sqrt{9 + 16} = \sqrt{25} = 5$, ale $\sqrt{9} + \sqrt{16} = 3 + 4 = 7$.
\end{definitionbox}

\subsection{Částečné odmocnění a usměrňování}

Dvě techniky, které budete u maturity potřebovat neustále, jsou částečné odmocnění a usměrňování zlomků.

\begin{definitionbox}
  \textbf{Částečné odmocnění}: Pokud pod odmocninou najdeme činitel, který je $n$-tou mocninou, můžeme ho `vytáhnout' před odmocninu:
  \begin{equation*}
    \sqrt[n]{a^n b} = a\sqrt[n]{b}, \qquad a, b\geq 0.
  \end{equation*}

  \textbf{Usměrnění zlomku}: Úprava zlomku tak, aby ve jmenovateli nebyla odmocnina. Zlomek rozšíříme vhodným výrazem:
  \begin{equation*}
    \frac{1}{\sqrt{a}} = \frac{1}{\sqrt{a}}\cdot\frac{\sqrt{a}}{\sqrt{a}} = \frac{\sqrt{a}}{a}, \qquad
    \frac{1}{\sqrt{a} + \sqrt{b}} = \frac{1}{\sqrt{a} + \sqrt{b}}\cdot\frac{\sqrt{a} - \sqrt{b}}{\sqrt{a} - \sqrt{b}} = \frac{\sqrt{a} - \sqrt{b}}{a - b}.
  \end{equation*}
  Ve druhém případě jsme využili vzorec $(x + y)(x - y) = x^2 - y^2$, ke kterému se brzy dostaneme.
\end{definitionbox}

\begin{example}
  Částečně odmocníme:
  \begin{align*}
    \sqrt{72} &= \sqrt{36\cdot 2} = 6\sqrt{2},\\
    \sqrt[3]{54} &= \sqrt[3]{27\cdot 2} = 3\sqrt[3]{2},\\
    \sqrt{50} - \sqrt{18} + \sqrt{8} &= 5\sqrt{2} - 3\sqrt{2} + 2\sqrt{2} = 4\sqrt{2}.
  \end{align*}
  Poslední příklad ukazuje, proč se částečné odmocnění vyplatí. Zdánlivě nesouvisející odmocniny se najednou daly sečíst, protože jsme je převedli na `stejný druh', podobně jako když sčítáme $5x - 3x + 2x$.
\end{example}

\begin{example}
  Usměrníme:
  \begin{align*}
    \frac{6}{\sqrt{3}} &= \frac{6\sqrt{3}}{3} = 2\sqrt{3},\\
    \frac{2}{\sqrt{5} - 1} &= \frac{2(\sqrt{5} + 1)}{(\sqrt{5} - 1)(\sqrt{5} + 1)} = \frac{2(\sqrt{5} + 1)}{5 - 1} = \frac{\sqrt{5} + 1}{2}.
  \end{align*}
  Mimochodem, číslo $\frac{1 + \sqrt{5}}{2}\approx 1{,}618$ je slavný \textbf{zlatý řez}, který se objevuje v umění, architektuře i v přírodě.
\end{example}

\begin{problem}
  Zjednodušte (částečně odmocněte a usměrněte):
  \begin{enumerate}[label=(\alph*)]
    \item $\sqrt{12} + \sqrt{27} - \sqrt{75}$,
    \item $\dfrac{10}{\sqrt{5}}$,
    \item $\dfrac{\sqrt{3}}{2 + \sqrt{3}}$,
    \item $\sqrt{a^5b^2}$, kde $a\geq 0$ a $b < 0$.
  \end{enumerate}
\end{problem}

\begin{solution}
  \begin{enumerate}[label=(\alph*)]
    \item $2\sqrt{3} + 3\sqrt{3} - 5\sqrt{3} = 0$,
    \item $\dfrac{10\sqrt{5}}{5} = 2\sqrt{5}$,
    \item $\dfrac{\sqrt{3}(2 - \sqrt{3})}{(2 + \sqrt{3})(2 - \sqrt{3})} = \dfrac{2\sqrt{3} - 3}{4 - 3} = 2\sqrt{3} - 3$,
    \item $\sqrt{a^4\cdot a\cdot b^2} = a^2\cdot|b|\cdot\sqrt{a} = -a^2b\sqrt{a}$. Pozor na to, že $b$ je záporné, takže $\sqrt{b^2} = |b| = -b$.
  \end{enumerate}
\end{solution}

\section{Mocniny s racionálním exponentem}

Udělejme ještě jeden krok. Stejně jako jsme si vynutili definici záporného exponentu, vynutíme si nyní definici zlomku v exponentu. Pokud by pravidlo ${(a^m)}^n = a^{mn}$ mělo platit i pro zlomky, pak musí platit
\begin{equation*}
  {\left(a^{\frac{1}{n}}\right)}^n = a^{\frac{1}{n}\cdot n} = a^1 = a.
\end{equation*}
Tedy $a^{\frac{1}{n}}$ je číslo, které umocněné na $n$ dá $a$. To už ale známe, je to přece $\sqrt[n]{a}$!

\begin{definitionbox}
  \textbf{Mocnina s racionálním exponentem}: Pro každé $a > 0$, $m\in\mathbb{Z}$ a $n\in\mathbb{N}$ definujeme
  \begin{equation*}
    a^{\frac{m}{n}} = \sqrt[n]{a^m}.
  \end{equation*}
  Pro $a = 0$ a kladný exponent klademe $0^{\frac{m}{n}} = 0$.

  \vspace{0.2cm}

  Pro mocniny s racionálním exponentem opět platí všechna pravidla pro počítání s mocninami (pro kladné základy). Díky tomu můžeme \textbf{odmocniny převádět na mocniny} a počítat s nimi mnohem pohodlněji.
\end{definitionbox}

\begin{note}
  Proč požadujeme $a > 0$? Zkusme definici použít na $a = -8$. Na jednu stranu bychom mohli chtít ${(-8)}^{\frac{1}{3}} = -2$. Na druhou stranu ale $\frac{1}{3} = \frac{2}{6}$, takže by mělo platit ${(-8)}^{\frac{1}{3}} = {(-8)}^{\frac{2}{6}} = \sqrt[6]{{(-8)}^2} = \sqrt[6]{64} = 2$. Dostali jsme dva různé výsledky pro stejné číslo, což je katastrofa. Proto se omezíme na kladné základy.
\end{note}

\begin{example}
  Převod na mocniny je velice silná zbraň. Zjednodušme
  \begin{equation*}
    \frac{\sqrt{a}\cdot\sqrt[3]{a^2}}{\sqrt[6]{a}}, \qquad a > 0.
  \end{equation*}
  Převedeme vše na mocniny a sečteme exponenty:
  \begin{equation*}
    \frac{a^{\frac{1}{2}}\cdot a^{\frac{2}{3}}}{a^{\frac{1}{6}}} = a^{\frac{1}{2} + \frac{2}{3} - \frac{1}{6}} = a^{\frac{3 + 4 - 1}{6}} = a^1 = a.
  \end{equation*}
  Zkuste si to samé spočítat jen s pravidly pro odmocniny. Bude to fungovat, ale bude to mnohem více bolet.
\end{example}

\begin{problem}
  Zjednodušte ($x > 0$):
  \begin{enumerate}[label=(\alph*)]
    \item $8^{\frac{2}{3}}$,
    \item $16^{-\frac{3}{4}}$,
    \item $\sqrt{x\sqrt{x\sqrt{x}}}$,
    \item $\dfrac{\sqrt[4]{x^3}\cdot x^{-\frac{1}{2}}}{\sqrt[8]{x^2}}$.
  \end{enumerate}
\end{problem}

\begin{hint}
  U (c) postupujte zevnitř: $x\sqrt{x} = x^{\frac{3}{2}}$ a tak dále.
\end{hint}

\begin{solution}
  \begin{enumerate}[label=(\alph*)]
    \item ${(\sqrt[3]{8})}^2 = 2^2 = 4$,
    \item $\dfrac{1}{{(\sqrt[4]{16})}^3} = \dfrac{1}{2^3} = \dfrac{1}{8}$,
    \item $x\sqrt{x} = x^{\frac{3}{2}}$, pak $x\cdot\sqrt{x^{\frac{3}{2}}} = x\cdot x^{\frac{3}{4}} = x^{\frac{7}{4}}$ a nakonec $\sqrt{x^{\frac{7}{4}}} = x^{\frac{7}{8}} = \sqrt[8]{x^7}$,
    \item $x^{\frac{3}{4} - \frac{1}{2} - \frac{2}{8}} = x^{0} = 1$.
  \end{enumerate}
\end{solution}

\section{Odbočka --- Iracionální exponent}

Co by mělo znamenat $2^{\sqrt{2}}$? Exponent $\sqrt{2}$ není zlomek, takže žádná z našich definic nefunguje. Myšlenka je ale jednoduchá. Číslo $\sqrt{2}$ můžeme libovolně přesně aproximovat zlomky:
\begin{equation*}
  1{,}4 = \tfrac{14}{10},\quad 1{,}41 = \tfrac{141}{100},\quad 1{,}414 = \tfrac{1414}{1000},\quad \ldots
\end{equation*}
A mocniny $2^{1{,}4}$, $2^{1{,}41}$, $2^{1{,}414}$, \ldots už spočítat umíme. Ukazuje se, že tato čísla se `blíží' k nějakému konkrétnímu číslu a to prohlásíme za $2^{\sqrt{2}}\approx 2{,}665$. Formálně se to dělá pomocí limit, ke kterým se dostaneme v pozdějších kapitolách. Důležité je, že i pro iracionální exponenty platí všechna pravidla pro mocniny s kladným základem.

Na závěr odbočky si ukážeme jeden z nejkrásnějších krátkých důkazů v matematice. Využijeme v něm princip \textit{tertium non datur} z první kapitoly.

\begin{definitionbox}
  \textbf{Věta}: Existují iracionální čísla $a, b$ taková, že $a^b$ je racionální.
\end{definitionbox}

\begin{proof}
  Uvažujme číslo $\sqrt{2}^{\sqrt{2}}$. Podle principu \textit{tertium non datur} je toto číslo buď racionální, nebo iracionální.
  \begin{itemize}
    \item Pokud je $\sqrt{2}^{\sqrt{2}}$ racionální, jsme hotovi. Stačí zvolit $a = b = \sqrt{2}$, o kterém víme, že je iracionální.
    \item Pokud je $\sqrt{2}^{\sqrt{2}}$ iracionální, zvolme $a = \sqrt{2}^{\sqrt{2}}$ a $b = \sqrt{2}$. Obě čísla jsou iracionální a
    \begin{equation*}
      a^b = {\left(\sqrt{2}^{\sqrt{2}}\right)}^{\sqrt{2}} = \sqrt{2}^{\sqrt{2}\cdot\sqrt{2}} = \sqrt{2}^2 = 2,
    \end{equation*}
    což je racionální číslo.
  \end{itemize}
  V obou případech jsme taková čísla našli.
\end{proof}

Na tomto důkazu je zvláštní to, že na jeho konci \textit{nevíme}, která z možností nastala. Víme jen, že jedna z nich určitě funguje. Takovým důkazům se říká \textbf{nekonstruktivní}.\footnote{Pro zvědavé: díky hluboké Gelfondově--Schneiderově větě dnes víme, že $\sqrt{2}^{\sqrt{2}}$ je iracionální, takže nastal druhý případ. Pro důkaz výše to ale vůbec nepotřebujeme.} Někteří matematici (takzvaní intuicionisté) takové důkazy odmítají přijmout, protože odmítají právě princip \textit{tertium non datur}.

\section{Algebraické výrazy}

Nyní už máme v ruce všechny operace a můžeme se podívat na výrazy. Výraz je, zjednodušeně řečeno, `smysluplně' zapsaná kombinace čísel, proměnných, operací a závorek.

\begin{definitionbox}
  \textbf{Algebraický výraz} je zápis složený z čísel, proměnných ($x, y, a, b, \ldots$), znaků operací ($+, -, \cdot, :,$ mocnin a odmocnin) a závorek, který dává smysl.

  \textbf{Hodnota výrazu} je číslo, které dostaneme, když za všechny proměnné dosadíme konkrétní čísla a výraz spočítáme.

  \textbf{Definiční obor výrazu} (také \textbf{podmínky}) je množina všech hodnot proměnných, pro které má výraz smysl. Výraz typicky nemá smysl, pokud:
  \begin{itemize}
    \item dělíme nulou (jmenovatel se nesmí rovnat nule),
    \item odmocňujeme záporné číslo (pod sudou odmocninou musí být nezáporné číslo),
    \item a později také logaritmujeme nekladné číslo.
  \end{itemize}
\end{definitionbox}

\begin{example}
  Určeme definiční obor výrazu
  \begin{equation*}
    V(x) = \frac{\sqrt{x + 3}}{x^2 - 4}.
  \end{equation*}
  Pod odmocninou musí být nezáporné číslo, tedy $x + 3\geq 0$, neboli $x\geq -3$. Jmenovatel nesmí být nula, tedy $x^2\neq 4$, neboli $x\neq 2$ a $x\neq -2$. Celkem dostáváme
  \begin{equation*}
    D(V) = \{x\in\mathbb{R}\ |\ x\geq -3,\ x\neq\pm 2\} = [-3, -2)\cup(-2, 2)\cup(2, \infty).
  \end{equation*}
  Zápis pomocí intervalů si podrobně vysvětlíme v další kapitole. Všimněte si, že definiční obor je \textit{průnik} podmínek, protože všechny musí platit zároveň. Opět se nám zde hodí množiny z předchozí kapitoly.
\end{example}

\begin{note}
  Definiční obor \textbf{vždy} určujeme z původního výrazu, ještě před jakoukoli úpravou! Například výraz $\frac{x^2}{x}$ se dá upravit na $x$, ale zatímco $x$ je definované pro všechna reálná čísla, původní výraz pro $x = 0$ definovaný není. Správně tedy píšeme $\frac{x^2}{x} = x$ pro $x\neq 0$. U ústní maturity se na podmínky zkoušející ptají velice rádi.
\end{note}

\subsection{Mnohočleny}

Nejdůležitějším typem výrazů jsou mnohočleny (také polynomy). Jsou to výrazy, ve kterých proměnnou jen sčítáme, odčítáme, násobíme číslem a umocňujeme na přirozené číslo.

\begin{definitionbox}
  \textbf{Mnohočlen} (polynom) v proměnné $x$ je výraz tvaru
  \begin{equation*}
    a_nx^n + a_{n-1}x^{n-1} + \cdots + a_1x + a_0,
  \end{equation*}
  kde $n\in\mathbb{N}_0$ a $a_0, a_1, \ldots, a_n\in\mathbb{R}$ jsou takzvané \textbf{koeficienty}. Pokud $a_n\neq 0$, říkáme, že mnohočlen má \textbf{stupeň} $n$. Jednotlivé sčítance $a_kx^k$ nazýváme \textbf{členy} mnohočlenu, $a_0$ je \textbf{absolutní člen}.

  \vspace{0.2cm}

  Mnohočleny stupně $1$ nazýváme \textbf{lineární} ($ax + b$), stupně $2$ \textbf{kvadratické} ($ax^2 + bx + c$) a stupně $3$ \textbf{kubické}.
\end{definitionbox}

\begin{example}
  Výraz $3x^4 - x^2 + 7$ je mnohočlen stupně $4$ s koeficienty $a_4 = 3$, $a_3 = 0$, $a_2 = -1$, $a_1 = 0$ a $a_0 = 7$. Výrazy $\frac{1}{x}$ nebo $\sqrt{x}$ mnohočleny nejsou, protože $x^{-1}$ a $x^{\frac{1}{2}}$ nemají přirozený exponent.
\end{example}

Mnohočleny sčítáme tak, že sečteme členy se stejnou mocninou (tzv. \textbf{sčítání podobných členů}). Násobíme je tak, že vynásobíme `každý s každým' (to je vlastně jen opakované použití distributivního zákona z axiomů tělesa) a pak sečteme podobné členy.

\begin{example}
  Mějme mnohočleny $P(x) = 2x^2 - 3x + 1$ a $Q(x) = x - 2$. Potom
  \begin{align*}
    P(x) + Q(x) &= 2x^2 - 3x + 1 + x - 2 = 2x^2 - 2x - 1,\\
    P(x)\cdot Q(x) &= (2x^2 - 3x + 1)(x - 2)\\
    &= 2x^3 - 4x^2 - 3x^2 + 6x + x - 2\\
    &= 2x^3 - 7x^2 + 7x - 2.
  \end{align*}
  Všimněte si, že stupeň součinu je součet stupňů ($2 + 1 = 3$). Proč tomu tak je?
\end{example}

\subsection{Vzorce}

Některé součiny se objevují tak často, že se vyplatí si je pamatovat. Ani ty si ale nemusíte pamatovat slepě. Stačí je jednou odvodit (roznásobit) a pochopit.

\begin{definitionbox}
  \textbf{Základní algebraické vzorce}: Pro všechna $a, b\in\mathbb{R}$ platí
  \begin{align*}
    {(a + b)}^2 &= a^2 + 2ab + b^2,\\
    {(a - b)}^2 &= a^2 - 2ab + b^2,\\
    (a + b)(a - b) &= a^2 - b^2,\\
    {(a + b)}^3 &= a^3 + 3a^2b + 3ab^2 + b^3,\\
    {(a - b)}^3 &= a^3 - 3a^2b + 3ab^2 - b^3,\\
    a^3 + b^3 &= (a + b)(a^2 - ab + b^2),\\
    a^3 - b^3 &= (a - b)(a^2 + ab + b^2).
  \end{align*}
\end{definitionbox}

\begin{example}\label{ex:vzorec-ctverec}
  Vzorec ${(a + b)}^2 = a^2 + 2ab + b^2$ má krásný geometrický význam, který znali už staří Řekové. Představte si čtverec o straně $a + b$. Jeho obsah je ${(a + b)}^2$. Zároveň se ale dá rozdělit na čtverec o straně $a$, čtverec o straně $b$ a dva obdélníky o stranách $a$ a $b$:
  \begin{center}
    \begin{tikzpicture}[scale=0.8]
      \fill[colorofwisdom!25] (0,0) rectangle (3,3);
      \fill[black!15] (3,3) rectangle (5,5);
      \fill[black!5] (3,0) rectangle (5,3);
      \fill[black!5] (0,3) rectangle (3,5);
      \draw[thick] (0,0) rectangle (5,5);
      \draw[thick] (3,0) -- (3,5);
      \draw[thick] (0,3) -- (5,3);
      \node at (1.5,1.5) {$a^2$};
      \node at (4,4) {$b^2$};
      \node at (4,1.5) {$ab$};
      \node at (1.5,4) {$ab$};
      \node[below] at (1.5,0) {$a$};
      \node[below] at (4,0) {$b$};
      \node[left] at (0,1.5) {$a$};
      \node[left] at (0,4) {$b$};
    \end{tikzpicture}
  \end{center}
  Sečtením obsahů dostaneme přesně $a^2 + 2ab + b^2$. Tento obrázek je mimochodem také důkazem toho, že ${(a + b)}^2 \neq a^2 + b^2$ --- ty dva obdélníky tam prostě jsou.
\end{example}

\begin{problem}
  Dokažte roznásobením vzorce pro $(a + b)(a - b)$, ${(a + b)}^3$ a $a^3 - b^3$.
\end{problem}

\begin{solution}
  \begin{align*}
    (a + b)(a - b) &= a^2 - ab + ab - b^2 = a^2 - b^2,\\
    {(a + b)}^3 &= {(a + b)}^2(a + b) = (a^2 + 2ab + b^2)(a + b)\\
    &= a^3 + a^2b + 2a^2b + 2ab^2 + ab^2 + b^3 = a^3 + 3a^2b + 3ab^2 + b^3,\\
    (a - b)(a^2 + ab + b^2) &= a^3 + a^2b + ab^2 - a^2b - ab^2 - b^3 = a^3 - b^3.
  \end{align*}
\end{solution}

\begin{note}
  Koeficienty ve vzorcích $1, 2, 1$ a $1, 3, 3, 1$ nejsou náhodné. Pokud si je napíšete pod sebe, dostanete začátek takzvaného \textbf{Pascalova trojúhelníku}, ve kterém je každé číslo součtem dvou čísel nad ním:
  \begin{center}
    \begin{tabular}{ccccccccc}
      &&&& 1 &&&& \\
      &&& 1 && 1 &&& \\
      && 1 && 2 && 1 && \\
      & 1 && 3 && 3 && 1 & \\
      1 && 4 && 6 && 4 && 1
    \end{tabular}
  \end{center}
  Poslední řádek nám tedy říká, že ${(a + b)}^4 = a^4 + 4a^3b + 6a^2b^2 + 4ab^3 + b^4$. Obecně to popisuje \textbf{binomická věta}, kterou si dokážeme v kapitole o kombinatorice.
\end{note}

\begin{problem}
  Upravte pomocí vzorců:
  \begin{enumerate}[label=(\alph*)]
    \item ${(3x - 2y)}^2$,
    \item $(2a + 5)(2a - 5)$,
    \item ${(x + 1)}^3 - {(x - 1)}^3$,
    \item $101^2 - 99^2$ (zpaměti!).
  \end{enumerate}
\end{problem}

\begin{solution}
  \begin{enumerate}[label=(\alph*)]
    \item $9x^2 - 12xy + 4y^2$,
    \item $4a^2 - 25$,
    \item $(x^3 + 3x^2 + 3x + 1) - (x^3 - 3x^2 + 3x - 1) = 6x^2 + 2$,
    \item $(101 + 99)(101 - 99) = 200\cdot 2 = 400$.
  \end{enumerate}
\end{solution}

\subsection{Rozklad mnohočlenu na součin}

Roznásobování je snadné, to zvládne každý. Opačný směr --- tedy zapsat mnohočlen jako součin jednodušších výrazů --- je mnohem zajímavější a důležitější. Pokud totiž máme mnohočlen zapsaný jako součin, můžeme snadno zjistit, kdy je roven nule, krátit zlomky nebo zjišťovat znaménko. Proč je to tak, vychází z následující jednoduché, ale klíčové vlastnosti reálných čísel.

\begin{definitionbox}
  \textbf{Věta (o nulovém součinu)}: Pro všechna $a, b\in\mathbb{R}$ platí
  \begin{equation*}
    a\cdot b = 0\quad\Leftrightarrow\quad a = 0\ \lor\ b = 0.
  \end{equation*}
\end{definitionbox}

\begin{proof}
  Implikace zprava doleva je jasná, protože $0\cdot b = 0$ a $a\cdot 0 = 0$.\footnote{I toto se dá dokázat z axiomů tělesa: $0\cdot b = (0 + 0)\cdot b = 0\cdot b + 0\cdot b$ a po odečtení $0\cdot b$ od obou stran dostaneme $0 = 0\cdot b$.} Zleva doprava: předpokládejme, že $ab = 0$. Pokud $a = 0$, jsme hotovi. Pokud $a\neq 0$, existuje podle axiomů tělesa inverzní prvek $a^{-1}$ a platí
  \begin{equation*}
    b = 1\cdot b = (a^{-1}a)\cdot b = a^{-1}(ab) = a^{-1}\cdot 0 = 0.
  \end{equation*}
\end{proof}

\begin{definitionbox}
  \textbf{Metody rozkladu na součin}:
  \begin{enumerate}
    \item \textbf{Vytýkání}: $ab + ac = a(b + c)$. Vždy zkuste jako první!
    \item \textbf{Použití vzorců}: zejména $a^2 - b^2 = (a - b)(a + b)$ a ${(a\pm b)}^2$.
    \item \textbf{Postupné vytýkání}: rozdělíme členy do skupin, z každé skupiny vytkneme a pak vytkneme společnou závorku.
    \item \textbf{Rozklad kvadratického trojčlenu}: hledáme čísla $r, s$ taková, že
    \begin{equation*}
      x^2 + px + q = (x - r)(x - s),\qquad\text{kde } r + s = -p,\ r\cdot s = q.
    \end{equation*}
    Obecně tuto metodu rozebereme v kapitole o kvadratických rovnicích.
  \end{enumerate}
\end{definitionbox}

\begin{example}
  Rozložme na součin:
  \begin{align*}
    6x^3 - 9x^2 &= 3x^2(2x - 3),\\
    4x^2 - 25 &= (2x - 5)(2x + 5),\\
    x^2 - 6x + 9 &= {(x - 3)}^2,\\
    x^3 - 2x^2 + 3x - 6 &= x^2(x - 2) + 3(x - 2) = (x - 2)(x^2 + 3),\\
    x^2 - 5x + 6 &= (x - 2)(x - 3),\\
    x^4 - 1 &= (x^2 - 1)(x^2 + 1) = (x - 1)(x + 1)(x^2 + 1).
  \end{align*}
  U pátého řádku jsme hledali dvě čísla, jejichž součet je $5$ a součin $6$. To jsou čísla $2$ a $3$. Na posledním řádku jsme vzorec použili dvakrát. Výraz $x^2 + 1$ už v reálných číslech dále rozložit nejde (proč?).
\end{example}

\begin{problem}
  Rozložte na součin co nejvíce činitelů:
  \begin{enumerate}[label=(\alph*)]
    \item $2a^2b - 8b^3$,
    \item $x^2 + 2x - 15$,
    \item $ax - ay + bx - by$,
    \item $x^3 + 8$,
    \item $3x^3 - 12x^2 + 12x$.
  \end{enumerate}
\end{problem}

\begin{solution}
  \begin{enumerate}[label=(\alph*)]
    \item $2b(a^2 - 4b^2) = 2b(a - 2b)(a + 2b)$,
    \item hledáme čísla se součinem $-15$ a součtem $-2$, tedy $-5$ a $3$: $(x - 3)(x + 5)$,
    \item $a(x - y) + b(x - y) = (x - y)(a + b)$,
    \item $x^3 + 2^3 = (x + 2)(x^2 - 2x + 4)$,
    \item $3x(x^2 - 4x + 4) = 3x{(x - 2)}^2$.
  \end{enumerate}
\end{solution}

\subsection{Lomené výrazy}

Lomený výraz je zlomek, v jehož jmenovateli se vyskytuje proměnná. S lomenými výrazy počítáme úplně stejně jako se zlomky, jen si musíme pamatovat na podmínky.

\begin{definitionbox}
  \textbf{Počítání s lomenými výrazy}: Pro výrazy $A, B, C, D$ (kde jmenovatelé jsou nenulové) platí:
  \begin{align*}
    \frac{A\cdot C}{B\cdot C} &= \frac{A}{B} &&\text{(krácení, resp. rozšiřování),}\\
    \frac{A}{B} + \frac{C}{D} &= \frac{AD + BC}{BD} &&\text{(sčítání),}\\
    \frac{A}{B}\cdot\frac{C}{D} &= \frac{AC}{BD} &&\text{(násobení),}\\
    \frac{A}{B} : \frac{C}{D} &= \frac{A}{B}\cdot\frac{D}{C} = \frac{AD}{BC} &&\text{(dělení, navíc } C\neq 0\text{).}
  \end{align*}
  Při sčítání se vyplatí místo součinu jmenovatelů $BD$ použít \textbf{nejmenší společný násobek} jmenovatelů. Ten najdeme tak, že jmenovatele rozložíme na součin.

  \vspace{0.2cm}

  \textbf{Postup úpravy lomeného výrazu}:
  \begin{enumerate}
    \item Rozložit všechny čitatele a jmenovatele na součin.
    \item Určit podmínky (všechny činitele ve jmenovatelích $\neq 0$, u dělení i čitatel dělitele).
    \item Upravit a krátit. Krátit můžeme jen \textbf{činitele}, nikdy sčítance!
  \end{enumerate}
\end{definitionbox}

\begin{example}
  Upravme a určeme podmínky:
  \begin{equation*}
    \frac{x^2 - 9}{x^2 + 3x}.
  \end{equation*}
  Rozložíme: $\frac{(x - 3)(x + 3)}{x(x + 3)}$. Podmínky: $x\neq 0$ a $x\neq -3$. Nyní můžeme krátit činitelem $(x + 3)$:
  \begin{equation*}
    \frac{x^2 - 9}{x^2 + 3x} = \frac{x - 3}{x}, \qquad x\neq 0,\ x\neq -3.
  \end{equation*}
  Podmínku $x\neq -3$ musíme uvést, i když se ve výsledku už nijak neprojevuje. Původní výraz totiž pro $x = -3$ neexistuje, zatímco upravený ano.
\end{example}

\begin{example}
  Sečtěme:
  \begin{equation*}
    \frac{2}{x - 1} - \frac{3}{x + 1} + \frac{6}{x^2 - 1}.
  \end{equation*}
  Rozložíme poslední jmenovatel: $x^2 - 1 = (x - 1)(x + 1)$. Podmínky: $x\neq\pm 1$. Společný jmenovatel je tedy $(x - 1)(x + 1)$ a
  \begin{equation*}
    \frac{2(x + 1) - 3(x - 1) + 6}{(x - 1)(x + 1)} = \frac{2x + 2 - 3x + 3 + 6}{(x - 1)(x + 1)} = \frac{11 - x}{(x - 1)(x + 1)}, \qquad x\neq\pm 1.
  \end{equation*}
\end{example}

\begin{example}
  Častým postrachem jsou \textbf{složené zlomky}. Zjednodušme
  \begin{equation*}
    \frac{\dfrac{1}{x} - \dfrac{1}{y}}{\dfrac{1}{x^2} - \dfrac{1}{y^2}}.
  \end{equation*}
  Podmínky: $x\neq 0$, $y\neq 0$ a jmenovatel hlavního zlomku nesmí být nula, tedy $x^2\neq y^2$, neboli $x\neq\pm y$. Upravíme čitatele a jmenovatele zvlášť:
  \begin{equation*}
    \frac{\dfrac{y - x}{xy}}{\dfrac{y^2 - x^2}{x^2y^2}} = \frac{y - x}{xy}\cdot\frac{x^2y^2}{(y - x)(y + x)} = \frac{xy}{x + y}.
  \end{equation*}
\end{example}

\begin{problem}
  Zjednodušte a určete podmínky:
  \begin{enumerate}[label=(\alph*)]
    \item $\dfrac{a^2 - 4a + 4}{a^2 - 4}$,
    \item $\dfrac{x}{x - 2} - \dfrac{4}{x^2 - 2x}$,
    \item $\dfrac{x^2 - 1}{x^2 + x}\ :\ \dfrac{x - 1}{x^2}$,
    \item $\left(\dfrac{a}{a - b} - \dfrac{b}{a + b}\right)\cdot\dfrac{a^2 - b^2}{a^2 + b^2}$.
  \end{enumerate}
\end{problem}

\begin{solution}
  \begin{enumerate}[label=(\alph*)]
    \item $\dfrac{{(a - 2)}^2}{(a - 2)(a + 2)} = \dfrac{a - 2}{a + 2}$, $a\neq\pm 2$.
    \item $x^2 - 2x = x(x - 2)$, podmínky $x\neq 0, 2$. Pak $\dfrac{x^2 - 4}{x(x - 2)} = \dfrac{(x - 2)(x + 2)}{x(x - 2)} = \dfrac{x + 2}{x}$.
    \item Podmínky: $x\neq 0$, $x\neq -1$ (jmenovatel prvního zlomku), $x\neq 1$ (čitatel dělitele). Pak
    \begin{equation*}
      \frac{(x - 1)(x + 1)}{x(x + 1)}\cdot\frac{x^2}{x - 1} = x.
    \end{equation*}
    \item Podmínky: $a\neq\pm b$ (výraz $a^2 + b^2$ je nulový jen pro $a = b = 0$, což už je vyloučeno). Závorka je $\dfrac{a(a + b) - b(a - b)}{(a - b)(a + b)} = \dfrac{a^2 + b^2}{(a - b)(a + b)}$ a po vynásobení dostaneme $1$.
  \end{enumerate}
  Výsledek (d) je typický. Pokud vám v maturitním příkladu vyjde nějaké `hezké' číslo, je to dobré znamení.
\end{solution}

\subsection{Dělení mnohočlenu mnohočlenem}

Stejně jako umíme dělit čísla `pod sebou' se zbytkem, umíme dělit i mnohočleny.

\begin{definitionbox}
  \textbf{Věta (o dělení se zbytkem)}: Pro mnohočleny $P$ a $Q$, kde $Q$ není nulový, existují jednoznačně určené mnohočleny $S$ (podíl) a $R$ (zbytek) takové, že
  \begin{equation*}
    P(x) = Q(x)\cdot S(x) + R(x),
  \end{equation*}
  a stupeň $R$ je menší než stupeň $Q$ (nebo je $R$ nulový).
\end{definitionbox}

\begin{example}
  Vydělme $(2x^3 - 3x^2 + 4x - 5) : (x - 1)$. Postupujeme jako u čísel: vždy vydělíme člen s nejvyšší mocninou, výsledek vynásobíme dělitelem a odečteme.
  \shorthandoff{-}% babel czech dela z - aktivni znak, coz rozbiji \cline
  \begin{equation*}
    \begin{array}{rrrrl}
      2x^3 & -\ 3x^2 & +\ 4x & -\ 5 & : (x - 1) = 2x^2 - x + 3\\
      -(2x^3 & -\ 2x^2) & & & \\ \cline{1-2}
      & -x^2 & +\ 4x & & \\
      & -(-x^2 & +\ x) & & \\ \cline{2-3}
      & & 3x & -\ 5 & \\
      & & -(3x & -\ 3) & \\ \cline{3-4}
      & & & -2 &
    \end{array}
  \end{equation*}
  Tedy $2x^3 - 3x^2 + 4x - 5 = (x - 1)(2x^2 - x + 3) - 2$. Zkuste si to roznásobit a ověřit.
\end{example}

\begin{note}
  Všimněte si, že zbytek $-2$ je přesně hodnota $P(1) = 2 - 3 + 4 - 5 = -2$. To není náhoda. Pokud dělíme mnohočlen $P(x)$ dvojčlenem $(x - r)$, je zbytek vždy roven $P(r)$. Stačí do rovnosti $P(x) = (x - r)S(x) + R$ dosadit $x = r$. Speciálně: pokud je $P(r) = 0$, je mnohočlen $P$ dělitelný dvojčlenem $(x - r)$. Toto pozorování nám velice pomůže při řešení rovnic vyšších stupňů.
\end{note}

\begin{problem}
  Vydělte $(x^3 - 6x^2 + 11x - 6) : (x - 1)$ a výsledný podíl dále rozložte na součin.
\end{problem}

\begin{solution}
  Podíl je $x^2 - 5x + 6$ se zbytkem $0$ (ověřte: $P(1) = 1 - 6 + 11 - 6 = 0$). Dále $x^2 - 5x + 6 = (x - 2)(x - 3)$, takže
  \begin{equation*}
    x^3 - 6x^2 + 11x - 6 = (x - 1)(x - 2)(x - 3).
  \end{equation*}
\end{solution}

\section{Shrnutí a otázky k ústní zkoušce}

U ústní maturity se od vás nečeká jen to, že umíte počítat, ale i to, že rozumíte, \textit{proč} věci fungují. Zkuste si zodpovědět následující otázky. Pokud to zvládnete bez nahlížení do textu, máte toto téma v kapse.

\begin{problem}
  Odpovězte vlastními slovy:
  \begin{enumerate}
    \item Jak je definována mocnina s přirozeným, celým a racionálním exponentem? Proč je $a^0 = 1$?
    \item Proč u racionálních exponentů požadujeme kladný základ?
    \item Jaký je rozdíl mezi $\sqrt{9}$ a řešením rovnice $x^2 = 9$? Čemu se rovná $\sqrt{x^2}$?
    \item Co je definiční obor výrazu a proč ho určujeme ještě před úpravami?
    \item Co je mnohočlen a jeho stupeň? Jak se mnohočleny sčítají, násobí a dělí?
    \item Odvoďte vzorec pro ${(a - b)}^2$ a nakreslete jeho geometrický význam.
    \item Jaké znáte metody rozkladu mnohočlenu na součin? Proč je rozklad užitečný?
    \item Proč můžeme krátit jen činitele a ne sčítance? Uveďte protipříklad.
  \end{enumerate}
\end{problem}

\begin{problem}
  Na závěr jeden těžší příklad pro odvážné. Zjednodušte pro $a > 0$, $a\neq 1$:
  \begin{equation*}
    \left(\frac{a^{\frac{3}{2}} - 1}{a^{\frac{1}{2}} - 1} - a^{\frac{1}{2}}\right)\cdot\frac{a - 1}{a + 1}.
  \end{equation*}
\end{problem}

\begin{hint}
  Označte $t = a^{\frac{1}{2}} = \sqrt{a}$. Pak $a^{\frac{3}{2}} = t^3$, $a = t^2$ a můžete použít vzorec pro $t^3 - 1$.
\end{hint}

\begin{solution}
  Se substitucí $t = \sqrt{a}$ ($t > 0$, $t\neq 1$) dostaneme
  \begin{equation*}
    \left(\frac{t^3 - 1}{t - 1} - t\right)\cdot\frac{t^2 - 1}{t^2 + 1} = \left(t^2 + t + 1 - t\right)\cdot\frac{t^2 - 1}{t^2 + 1} = (t^2 + 1)\cdot\frac{t^2 - 1}{t^2 + 1} = t^2 - 1 = a - 1.
  \end{equation*}
  Substituce je velice užitečný trik. Složitý výraz se tím `převlékne' za jednoduchý, který už umíme upravit.
\end{solution}

\koneckapitoly

\chapter{Lineární rovnice, nerovnice, soustavy a matice}
\mathquote{Do not worry about your difficulties in mathematics. I can assure you mine are still greater.}{Albert Einstein}

Rovnice jsou pravděpodobně to první, co si člověk představí pod slovem `matematika'. Najdi $x$. V této kapitole si ale ukážeme, že řešení rovnic není jen mechanické přehazování čísel z jedné strany na druhou. Za každým krokem stojí logika z první kapitoly a pokud ji pochopíme, vyhneme se většině chyb. Začneme lineárními rovnicemi, pak se podíváme na nerovnice a nakonec na soustavy rovnic.

\section{Co je to rovnice?}

V první kapitole jsme si řekli, že výrok je tvrzení, o kterém lze rozhodnout, zda je pravdivé. Je tedy $2x + 1 = 7$ výrok? Není! Dokud nevíme, co je $x$, nemůžeme o pravdivosti rozhodnout. Pro $x = 3$ je to pravda, pro $x = 0$ nepravda. Takovým tvrzením jsme v odbočce o predikátové logice říkali predikáty.

\begin{definitionbox}
  \textbf{Rovnice} s neznámou $x$ je zápis $L(x) = P(x)$, kde $L$ a $P$ jsou výrazy. Výraz $L(x)$ nazýváme \textbf{levá strana} a $P(x)$ \textbf{pravá strana} rovnice.

  \vspace{0.2cm}

  \textbf{Definiční obor rovnice} $D$ je množina všech čísel, pro která mají smysl obě strany rovnice.

  \textbf{Řešení} (také kořen) rovnice je každé číslo $x_0\in D$, pro které je $L(x_0) = P(x_0)$ pravdivý výrok.

  \textbf{Množina všech řešení} (značíme $K$, jako kořeny) je tedy
  \begin{equation*}
    K = \{x\in D\ |\ L(x) = P(x)\}.
  \end{equation*}
  \textbf{Vyřešit rovnici} znamená najít množinu $K$, tedy najít \textit{všechna} řešení a ukázat, že jiná neexistují.
\end{definitionbox}

Poslední věta je důležitá. Pokud uhodnete, že $x = 3$ je řešením $2x + 1 = 7$, ještě jste rovnici nevyřešili. Musíte také ukázat, že žádné jiné řešení neexistuje. K tomu nám slouží ekvivalentní úpravy.

\begin{definitionbox}
  \textbf{Ekvivalentní úprava} rovnice je úprava, po které má nová rovnice \textbf{stejnou} množinu řešení jako původní rovnice. Mezi ekvivalentní úpravy patří:
  \begin{enumerate}
    \item přičtení (odečtení) stejného čísla nebo výrazu, který je definovaný na celém $D$, k oběma stranám,
    \item vynásobení (vydělení) obou stran stejným \textbf{nenulovým} číslem nebo výrazem, který je na $D$ definovaný a nenulový,
    \item záměna stran rovnice,
    \item úprava jedné strany rovnice (roznásobení, sečtení podobných členů, \ldots), pokud se nezmění definiční obor.
  \end{enumerate}

  \textbf{Důsledková úprava} je úprava, po které nová rovnice obsahuje \textbf{všechna} řešení té původní, ale může obsahovat i nějaká navíc. Typicky je to \textbf{umocnění obou stran} nebo vynásobení výrazem, který může být nulový. Po důsledkové úpravě je \textbf{nutné provést zkoušku}.
\end{definitionbox}

Pokud se na to podíváme optikou první kapitoly, ekvivalentní úprava znamená, že mezi původní a novou rovnicí platí ekvivalence $\Leftrightarrow$. Důsledková úprava zaručuje jen implikaci $\Rightarrow$. Proto ty názvy.

\begin{example}
  Proč je násobení nulou zakázané? Vezměme rovnici $x = 2$, která má jediné řešení. Pokud obě strany vynásobíme nulou, dostaneme $0 = 0$, což platí pro všechna $x$. Najednou máme nekonečně mnoho `řešení'.

  Proč umocnění není ekvivalentní? Vezměme opět $x = 2$. Umocněním dostaneme $x^2 = 4$, která má řešení $2$ a $-2$. Řešení $-2$ přibylo. Z $a = b$ plyne $a^2 = b^2$, ale z $a^2 = b^2$ plyne jen $a = b$ \textbf{nebo} $a = -b$.
\end{example}

\begin{definitionbox}
  \textbf{Zkouška} je dosazení nalezeného čísla do \textbf{původní} rovnice (zvlášť levá a pravá strana). Zkouška je \textbf{nutnou} součástí řešení, pokud jsme použili důsledkovou úpravu. Pokud jsme používali jen ekvivalentní úpravy, je zkouška pouze kontrolou, zda jsme neudělali početní chybu. I tak ji doporučujeme dělat vždy.
\end{definitionbox}

\section{Lineární rovnice}

\begin{definitionbox}
  \textbf{Lineární rovnice} s neznámou $x$ je rovnice, kterou lze ekvivalentními úpravami převést na tvar
  \begin{equation*}
    ax + b = 0,\qquad a, b\in\mathbb{R}.
  \end{equation*}
\end{definitionbox}

Lineární rovnici umíme vyřešit úplně obecně. To je vzácná věc, takže si to užijme.

\begin{definitionbox}
  \textbf{Věta (o řešení lineární rovnice)}: Pro rovnici $ax + b = 0$ s neznámou $x\in\mathbb{R}$ platí:
  \begin{enumerate}
    \item Pokud $a\neq 0$, má rovnice právě jedno řešení $x = -\frac{b}{a}$, tedy $K = \left\{-\frac{b}{a}\right\}$.
    \item Pokud $a = 0$ a $b\neq 0$, nemá rovnice žádné řešení, tedy $K = \emptyset$.
    \item Pokud $a = 0$ a $b = 0$, je řešením každé reálné číslo, tedy $K = \mathbb{R}$.
  \end{enumerate}
\end{definitionbox}

\begin{proof}
  \begin{enumerate}
    \item Pokud $a\neq 0$, můžeme k oběma stranám přičíst $-b$ a obě strany vydělit číslem $a$. Obě úpravy jsou ekvivalentní, takže
    \begin{equation*}
      ax + b = 0\quad\Leftrightarrow\quad ax = -b\quad\Leftrightarrow\quad x = -\frac{b}{a}.
    \end{equation*}
    \item Pokud $a = 0$, má rovnice tvar $0\cdot x + b = 0$, neboli $b = 0$. To je ale výrok, který vůbec nezávisí na $x$. Pokud $b\neq 0$, je nepravdivý pro všechna $x$, takže $K = \emptyset$.
    \item Ze stejného důvodu je pro $b = 0$ výrok $0 = 0$ pravdivý pro všechna $x$, takže $K = \mathbb{R}$.
  \end{enumerate}
\end{proof}

\begin{example}
  Vyřešme rovnici $3(x - 2) - 2(x + 1) = 4 - x$.
  \begin{align*}
    3x - 6 - 2x - 2 &= 4 - x,\\
    x - 8 &= 4 - x,\\
    2x &= 12,\\
    x &= 6.
  \end{align*}
  Zkouška: $L(6) = 3\cdot 4 - 2\cdot 7 = -2$, $P(6) = 4 - 6 = -2$. Tedy $L(6) = P(6)$ a $K = \{6\}$.
\end{example}

\begin{example}
  Vyřešme rovnici $2(x + 3) = 2x + 5$. Roznásobením dostaneme $2x + 6 = 2x + 5$ a po odečtení $2x$ dostaneme $6 = 5$. To je nepravdivý výrok, tedy $K = \emptyset$.

  Naopak rovnice $2(x + 3) = 2x + 6$ po stejné úpravě vede na $6 = 6$, a tedy $K = \mathbb{R}$. Takovým rovnicím se někdy říká \textbf{identity}. Například všechny vzorce z předchozí kapitoly jsou identity.
\end{example}

\begin{problem}
  Vyřešte v $\mathbb{R}$:
  \begin{enumerate}[label=(\alph*)]
    \item $5x - 3 = 2x + 9$,
    \item $\dfrac{x - 1}{2} - \dfrac{x + 2}{3} = 1$,
    \item $(x - 2)^2 - (x + 1)^2 = 3$,
    \item $\dfrac{2x + 1}{4} - \dfrac{x}{2} = \dfrac{1}{4}$.
  \end{enumerate}
\end{problem}

\begin{hint}
  U (b) a (d) vynásobte obě strany nejmenším společným násobkem jmenovatelů. U (c) se nelekněte druhých mocnin. Po roznásobení se $x^2$ odečtou.
\end{hint}

\begin{solution}
  \begin{enumerate}[label=(\alph*)]
    \item $3x = 12$, tedy $K = \{4\}$.
    \item Vynásobíme šesti: $3(x - 1) - 2(x + 2) = 6$, tedy $x - 7 = 6$ a $K = \{13\}$.
    \item $x^2 - 4x + 4 - x^2 - 2x - 1 = 3$, tedy $-6x + 3 = 3$ a $K = \{0\}$.
    \item Vynásobíme čtyřmi: $2x + 1 - 2x = 1$, tedy $1 = 1$ a $K = \mathbb{R}$.
  \end{enumerate}
\end{solution}

\subsection{Rovnice s neznámou ve jmenovateli}

Pokud se neznámá objeví ve jmenovateli, musíme nejdříve určit definiční obor. Pak rovnici vynásobíme společným jmenovatelem. Tato úprava je ekvivalentní jen na definičním oboru, protože na něm je jmenovatel nenulový. Proto musíme nakonec zkontrolovat, zda výsledek do definičního oboru patří.

\begin{example}
  Vyřešme rovnici
  \begin{equation*}
    \frac{3}{x - 2} + \frac{1}{x} = \frac{6}{x(x - 2)}.
  \end{equation*}
  Podmínky: $x\neq 0$, $x\neq 2$. Vynásobíme $x(x - 2)$:
  \begin{align*}
    3x + (x - 2) &= 6,\\
    4x &= 8,\\
    x &= 2.
  \end{align*}
  Jenže $x = 2$ nepatří do definičního oboru! Rovnice tedy nemá řešení, $K = \emptyset$. Pokud bychom zapomněli na podmínky, napsali bychom špatný výsledek. Zkouškou bychom chybu odhalili také, protože bychom zjistili, že do původní rovnice $2$ dosadit nejde.
\end{example}

\begin{problem}
  Vyřešte v $\mathbb{R}$:
  \begin{enumerate}[label=(\alph*)]
    \item $\dfrac{x + 1}{x - 1} = 3$,
    \item $\dfrac{2}{x + 3} - \dfrac{1}{x - 3} = \dfrac{x - 15}{x^2 - 9}$,
    \item $\dfrac{x}{x - 1} - \dfrac{1}{x + 1} = \dfrac{x^2 + 1}{x^2 - 1}$.
  \end{enumerate}
\end{problem}

\begin{solution}
  \begin{enumerate}[label=(\alph*)]
    \item $x\neq 1$. Pak $x + 1 = 3x - 3$, tedy $x = 2$ a $K = \{2\}$.
    \item $x\neq\pm 3$. Vynásobíme $(x + 3)(x - 3)$: $2(x - 3) - (x + 3) = x - 15$, tedy $x - 9 = x - 15$, neboli $-9 = -15$. $K = \emptyset$.
    \item $x\neq\pm 1$. Vynásobíme $(x - 1)(x + 1)$: $x(x + 1) - (x - 1) = x^2 + 1$, tedy $x^2 + 1 = x^2 + 1$. To platí vždy, ale musíme vyřadit čísla mimo definiční obor: $K = \mathbb{R}\setminus\{-1, 1\}$.
  \end{enumerate}
  Výsledek (c) je pěkná past. Rovnice je `skoro' identita, ale podmínky z ní vyřazují dva body.
\end{solution}

\subsection{Rovnice s parametrem}

Občas v rovnici kromě neznámé vystupuje ještě nějaké další písmeno, takzvaný \textbf{parametr}. Parametr je číslo, které neznáme, ale je pevně dané. Úkolem je vyřešit rovnici pro \textit{všechny} možné hodnoty parametru. Tomuto rozboru říkáme \textbf{diskuze}. Pro lineární rovnice nám přesně řekne, co dělat, věta o řešení lineární rovnice.

\begin{example}
  Vyřešme rovnici $a(x - 1) = 2x + 3$ s neznámou $x\in\mathbb{R}$ a parametrem $a\in\mathbb{R}$.

  Převedeme na tvar z věty:
  \begin{equation*}
    ax - a = 2x + 3\quad\Leftrightarrow\quad (a - 2)x = a + 3.
  \end{equation*}
  Nyní musíme rozlišit, zda je koeficient u $x$ nulový.
  \begin{itemize}
    \item Pokud $a\neq 2$, můžeme dělit a $x = \frac{a + 3}{a - 2}$.
    \item Pokud $a = 2$, dostáváme $0\cdot x = 5$, což nemá řešení.
  \end{itemize}
  Výsledek se obvykle zapisuje do tabulky:
  \begin{center}
    \begin{tabular}{|c|c|}
      \hline
      Parametr & Množina řešení \\
      \hline
      $a = 2$ & $K = \emptyset$ \\
      $a\in\mathbb{R}\setminus\{2\}$ & $K = \left\{\frac{a + 3}{a - 2}\right\}$ \\
      \hline
    \end{tabular}
  \end{center}
\end{example}

\begin{problem}
  Vyřešte s neznámou $x\in\mathbb{R}$ a parametrem $p\in\mathbb{R}$:
  \begin{enumerate}[label=(\alph*)]
    \item $px + 1 = x + p$,
    \item $p^2x - 4 = 2p + 4x$.
  \end{enumerate}
\end{problem}

\begin{solution}
  \begin{enumerate}[label=(\alph*)]
    \item $(p - 1)x = p - 1$. Pro $p\neq 1$ je $K = \{1\}$. Pro $p = 1$ dostaneme $0 = 0$, tedy $K = \mathbb{R}$.
    \item $(p^2 - 4)x = 2p + 4$, neboli $(p - 2)(p + 2)x = 2(p + 2)$.
    \begin{itemize}
      \item Pro $p = 2$: $0\cdot x = 8$, tedy $K = \emptyset$.
      \item Pro $p = -2$: $0\cdot x = 0$, tedy $K = \mathbb{R}$.
      \item Pro $p\neq\pm 2$: $x = \frac{2(p + 2)}{(p - 2)(p + 2)} = \frac{2}{p - 2}$, tedy $K = \left\{\frac{2}{p - 2}\right\}$.
    \end{itemize}
  \end{enumerate}
\end{solution}

\subsection{Vyjádření neznámé ze vzorce}

Velice praktická dovednost, zejména pro fyziku, je vyjádřit jednu veličinu ze vzorce. Je to vlastně jen řešení rovnice, ve které je `neznámou' jedno písmeno a ostatní písmena se chovají jako parametry. Opět musíme dávat pozor, čím dělíme.

\begin{example}
  Ze vzorce pro obsah lichoběžníku $S = \frac{(a + c)v}{2}$ vyjádřeme $c$ (všechny veličiny jsou kladné):
  \begin{equation*}
    2S = (a + c)v\quad\Rightarrow\quad \frac{2S}{v} = a + c\quad\Rightarrow\quad c = \frac{2S}{v} - a.
  \end{equation*}
\end{example}

\begin{problem}
  Vyjádřete:
  \begin{enumerate}[label=(\alph*)]
    \item ze vzorce $\dfrac{1}{f} = \dfrac{1}{a} + \dfrac{1}{b}$ (zobrazovací rovnice čočky) veličinu $b$,
    \item ze vzorce $v = v_0 + at$ veličinu $t$,
    \item ze vzorce $R = \dfrac{R_1R_2}{R_1 + R_2}$ veličinu $R_1$.
  \end{enumerate}
\end{problem}

\begin{solution}
  \begin{enumerate}[label=(\alph*)]
    \item $\dfrac{1}{b} = \dfrac{1}{f} - \dfrac{1}{a} = \dfrac{a - f}{af}$, tedy $b = \dfrac{af}{a - f}$ (pro $a\neq f$).
    \item $t = \dfrac{v - v_0}{a}$ (pro $a\neq 0$).
    \item $R(R_1 + R_2) = R_1R_2$, tedy $RR_1 - R_1R_2 = -RR_2$, neboli $R_1(R - R_2) = -RR_2$ a $R_1 = \dfrac{RR_2}{R_2 - R}$ (pro $R\neq R_2$).
  \end{enumerate}
\end{solution}

\section{Intervaly}

Než se pustíme do nerovnic, potřebujeme si zavést šikovný způsob zápisu jejich řešení. Řešením nerovnice totiž většinou není jedno číslo, ale celý `kus' reálné osy.

\begin{definitionbox}
  \textbf{Intervaly}: Nechť $a, b\in\mathbb{R}$, $a < b$. Definujeme
  \begin{center}
    \renewcommand{\arraystretch}{1.3}
    \begin{tabular}{|l|l|l|}
      \hline
      Název & Zápis & Množina \\
      \hline
      uzavřený & $[a, b]$ & $\{x\in\mathbb{R}\ |\ a\leq x\leq b\}$ \\
      otevřený & $(a, b)$ & $\{x\in\mathbb{R}\ |\ a < x < b\}$ \\
      zleva uzavřený, zprava otevřený & $[a, b)$ & $\{x\in\mathbb{R}\ |\ a\leq x < b\}$ \\
      zleva otevřený, zprava uzavřený & $(a, b]$ & $\{x\in\mathbb{R}\ |\ a < x\leq b\}$ \\
      \hline
      neomezené & $[a, \infty)$ & $\{x\in\mathbb{R}\ |\ x\geq a\}$ \\
       & $(a, \infty)$ & $\{x\in\mathbb{R}\ |\ x > a\}$ \\
       & $(-\infty, b]$ & $\{x\in\mathbb{R}\ |\ x\leq b\}$ \\
       & $(-\infty, b)$ & $\{x\in\mathbb{R}\ |\ x < b\}$ \\
       & $(-\infty, \infty)$ & $\mathbb{R}$ \\
      \hline
    \end{tabular}
  \end{center}
  Hranatá závorka znamená, že krajní bod do intervalu \textbf{patří}, kulatá, že \textbf{nepatří}. U nekonečna je vždy kulatá závorka, protože $\infty$ není reálné číslo, a tedy do intervalu patřit nemůže.
\end{definitionbox}

\begin{note}
  V mnoha českých učebnicích se uzavřený konec intervalu značí lomenou závorkou, například $\langle a, b)$ místo $[a, b)$. Oba zápisy znamenají totéž. U maturity můžete použít kterýkoli z nich.
\end{note}

Intervaly si můžeme znázornit na číselné ose. Bod, který do intervalu patří, kreslíme jako plné kolečko, bod, který nepatří, jako prázdné kolečko.

\begin{example}
  Interval $[-1, 3)$ na číselné ose:
  \begin{center}
    \begin{tikzpicture}
      \draw[->, thick] (-3,0) -- (5,0) node[right] {$x$};
      \foreach \x in {-2,-1,0,1,2,3,4} \draw (\x,0.1) -- (\x,-0.1) node[below] {$\x$};
      \draw[line width=3pt, colorofwisdom] (-1,0) -- (3,0);
      \filldraw[colorofwisdom] (-1,0) circle (3pt);
      \filldraw[fill=white, draw=colorofwisdom, thick] (3,0) circle (3pt);
    \end{tikzpicture}
  \end{center}
\end{example}

Protože intervaly jsou množiny, můžeme s nimi provádět všechny množinové operace z předchozí kapitoly. Průnik a sjednocení intervalů budeme při řešení nerovnic potřebovat neustále.

\begin{problem}
  Určete:
  \begin{enumerate}[label=(\alph*)]
    \item $[-2, 5)\cap(1, 7]$,
    \item $(-\infty, 3]\cup(0, 10)$,
    \item $(-\infty, 1)\cap[1, \infty)$,
    \item $\mathbb{R}\setminus(-1, 2]$.
  \end{enumerate}
\end{problem}

\begin{solution}
  \begin{enumerate}[label=(\alph*)]
    \item $(1, 5)$,
    \item $(-\infty, 10)$,
    \item $\emptyset$ (bod $1$ do prvního intervalu nepatří),
    \item $(-\infty, -1]\cup(2, \infty)$.
  \end{enumerate}
\end{solution}

\section{Lineární nerovnice}

\begin{definitionbox}
  \textbf{Nerovnice} je zápis $L(x) < P(x)$, případně se znaky $>$, $\leq$, $\geq$. Definiční obor, řešení a množina řešení se definují stejně jako u rovnic.

  \textbf{Lineární nerovnice} je nerovnice, kterou lze ekvivalentními úpravami převést na jeden z tvarů
  \begin{equation*}
    ax + b < 0,\quad ax + b > 0,\quad ax + b\leq 0,\quad ax + b\geq 0.
  \end{equation*}
\end{definitionbox}

S nerovnicemi pracujeme téměř stejně jako s rovnicemi. Je zde ale jeden zásadní rozdíl, který je potřeba si vrýt do paměti.

\begin{definitionbox}
  \textbf{Ekvivalentní úpravy nerovnic}:
  \begin{enumerate}
    \item Přičtení stejného čísla nebo výrazu k oběma stranám \textbf{nemění} znak nerovnosti.
    \item Vynásobení obou stran stejným \textbf{kladným} číslem \textbf{nemění} znak nerovnosti.
    \item Vynásobení obou stran stejným \textbf{záporným} číslem \textbf{obrátí} znak nerovnosti ($<$ na $>$, $\leq$ na $\geq$ a naopak).
  \end{enumerate}
  Z toho plyne, že nerovnici \textbf{nikdy nesmíme násobit výrazem s neznámou}, pokud nevíme, jaké má znaménko!
\end{definitionbox}

Proč se znak otočí? Intuitivně: $2 < 3$, ale $-2 > -3$, protože na záporné straně osy je `větší' číslo blíže nule. Dokažme si to.

\begin{proof}
  Nechť $a < b$ a $c < 0$. Chceme ukázat, že $ac > bc$. Z $a < b$ plyne, že $b - a > 0$. Z $c < 0$ plyne, že $-c > 0$. Součin dvou kladných čísel je kladný, tedy
  \begin{equation*}
    (b - a)(-c) > 0\quad\Leftrightarrow\quad -bc + ac > 0\quad\Leftrightarrow\quad ac > bc.
  \end{equation*}
\end{proof}

\begin{example}
  Vyřešme nerovnici $3 - 2x\geq x - 6$.
  \begin{align*}
    -2x - x &\geq -6 - 3,\\
    -3x &\geq -9 \quad\Big|\ :(-3),\\
    x &\leq 3.
  \end{align*}
  V posledním kroku jsme dělili záporným číslem, a proto se znak otočil. Množina řešení je $K = (-\infty, 3]$.
\end{example}

\begin{problem}
  Vyřešte v $\mathbb{R}$:
  \begin{enumerate}[label=(\alph*)]
    \item $4x + 7 < 2x - 1$,
    \item $\dfrac{x}{3} - \dfrac{x - 1}{2}\geq 1$,
    \item $2(x - 1) > 2x + 3$,
    \item $(x + 2)^2\leq x^2 + 4x + 4$.
  \end{enumerate}
\end{problem}

\begin{solution}
  \begin{enumerate}[label=(\alph*)]
    \item $2x < -8$, tedy $x < -4$ a $K = (-\infty, -4)$.
    \item Vynásobíme šesti: $2x - 3x + 3\geq 6$, tedy $-x\geq 3$ a $x\leq -3$. $K = (-\infty, -3]$.
    \item $-2 > 3$, což neplatí nikdy. $K = \emptyset$.
    \item Levá strana se po roznásobení rovná pravé, tedy $0\leq 0$, což platí vždy. $K = \mathbb{R}$.
  \end{enumerate}
\end{solution}

\subsection{Soustavy nerovnic}

Pokud má neznámá splňovat několik nerovnic zároveň, mluvíme o soustavě nerovnic. Každou nerovnici vyřešíme zvlášť a výsledná množina řešení je \textbf{průnik} jednotlivých množin řešení. Opět se zde krásně potkává logická spojka $\land$ s množinovým průnikem $\cap$.

\begin{example}
  Vyřešme soustavu
  \begin{align*}
    2x - 1 &> -5,\\
    3 - x &\geq 2x - 6.
  \end{align*}
  První nerovnice: $2x > -4$, tedy $x > -2$ a $K_1 = (-2, \infty)$. Druhá nerovnice: $-3x\geq -9$, tedy $x\leq 3$ a $K_2 = (-\infty, 3]$. Celkem
  \begin{equation*}
    K = K_1\cap K_2 = (-2, 3].
  \end{equation*}
\end{example}

\begin{example}
  Zápis $-1 < 2x + 3\leq 7$ je zkrácený zápis soustavy dvou nerovnic $-1 < 2x + 3$ a $2x + 3\leq 7$. Takovou `dvojitou' nerovnici můžeme upravovat najednou, pokud provádíme stejnou úpravu se všemi třemi částmi:
  \begin{equation*}
    -1 < 2x + 3\leq 7\quad\Leftrightarrow\quad -4 < 2x\leq 4\quad\Leftrightarrow\quad -2 < x\leq 2.
  \end{equation*}
  Tedy $K = (-2, 2]$.
\end{example}

\subsection{Nerovnice v součinovém a podílovém tvaru}

Co když potřebujeme vyřešit nerovnici jako $\frac{x - 1}{x + 2}\geq 0$? Vynásobit $(x + 2)$ nemůžeme, protože nevíme, zda je tento výraz kladný nebo záporný. Místo toho využijeme toho, že znaménko součinu (i podílu) závisí jen na znaménkách jednotlivých činitelů.

\begin{definitionbox}
  \textbf{Metoda nulových bodů} (také metoda intervalů) pro nerovnice tvaru $\frac{A(x)}{B(x)}\gtrless 0$ nebo $A(x)\cdot B(x)\gtrless 0$:
  \begin{enumerate}
    \item Převedeme všechno na jednu stranu, aby na druhé byla nula. \textbf{Nikdy nenásobíme jmenovatelem!}
    \item Čitatele i jmenovatele rozložíme na součin lineárních činitelů.
    \item Najdeme \textbf{nulové body}, tedy čísla, ve kterých je některý činitel nulový.
    \item Nulové body rozdělí reálnou osu na intervaly. V každém z nich má každý činitel stálé znaménko. Znaménka zapíšeme do tabulky.
    \item Vybereme intervaly, kde má výraz požadované znaménko. Pozor na nulové body jmenovatele, ty do řešení nepatří nikdy.
  \end{enumerate}
\end{definitionbox}

\begin{example}\label{ex:podilova-nerovnice}
  Vyřešme nerovnici
  \begin{equation*}
    \frac{x - 1}{x + 2}\geq 0.
  \end{equation*}
  Podmínka: $x\neq -2$. Nulové body: $x = 1$ (čitatel) a $x = -2$ (jmenovatel). Sestavíme tabulku znamének. Znaménko výrazu $x - 1$ v každém intervalu zjistíme například dosazením libovolného čísla z daného intervalu.
  \begin{center}
    \renewcommand{\arraystretch}{1.3}
    \begin{tabular}{|c|c|c|c|}
      \hline
       & $(-\infty, -2)$ & $(-2, 1)$ & $(1, \infty)$ \\
      \hline
      $x - 1$ & $-$ & $-$ & $+$ \\
      $x + 2$ & $-$ & $+$ & $+$ \\
      \hline
      $\frac{x - 1}{x + 2}$ & $+$ & $-$ & $+$ \\
      \hline
    \end{tabular}
  \end{center}
  Hledáme, kde je výraz kladný nebo nulový. Kladný je na $(-\infty, -2)$ a $(1, \infty)$, nulový je pro $x = 1$. Bod $x = -2$ do řešení nepatří. Celkem
  \begin{equation*}
    K = (-\infty, -2)\cup[1, \infty).
  \end{equation*}
\end{example}

\begin{note}
  Proč nemůžeme v příkladu~\ref{ex:podilova-nerovnice} jednoduše vynásobit obě strany $(x + 2)$? Dostali bychom $x - 1\geq 0$, tedy $x\geq 1$, a přišli bychom o celý interval $(-\infty, -2)$. Pro $x < -2$ je totiž $x + 2$ záporné a znak nerovnosti by se měl otočit. Pokud bychom chtěli přesto násobit, museli bychom rozlišit dva případy ($x + 2 > 0$ a $x + 2 < 0$). To je také správný postup, jen je pracnější.
\end{note}

\begin{problem}
  Vyřešte v $\mathbb{R}$:
  \begin{enumerate}[label=(\alph*)]
    \item $(x + 3)(2 - x) > 0$,
    \item $\dfrac{2x + 4}{x - 3}\leq 0$,
    \item $\dfrac{x + 1}{x - 2} > 2$,
    \item $\dfrac{x(x - 4)}{x + 1}\geq 0$.
  \end{enumerate}
\end{problem}

\begin{hint}
  U (c) nejdříve odečtěte dvojku a převeďte levou stranu na společný jmenovatel. U (b) a (c) si dejte pozor na nulový bod jmenovatele.
\end{hint}

\begin{solution}
  \begin{enumerate}[label=(\alph*)]
    \item Nulové body $-3$ a $2$. Výraz je kladný mezi nimi, $K = (-3, 2)$.
    \item Nulové body $-2$ a $3$. Výraz je záporný na $(-2, 3)$, nulový v $-2$. $K = [-2, 3)$.
    \item $\dfrac{x + 1 - 2(x - 2)}{x - 2} > 0$, tedy $\dfrac{5 - x}{x - 2} > 0$. Nulové body $2$ a $5$, $K = (2, 5)$.
    \item Nulové body $-1$, $0$ a $4$. Tabulka:
    \begin{center}
      \renewcommand{\arraystretch}{1.3}
      \begin{tabular}{|c|c|c|c|c|}
        \hline
         & $(-\infty, -1)$ & $(-1, 0)$ & $(0, 4)$ & $(4, \infty)$ \\
        \hline
        $x$ & $-$ & $-$ & $+$ & $+$ \\
        $x - 4$ & $-$ & $-$ & $-$ & $+$ \\
        $x + 1$ & $-$ & $+$ & $+$ & $+$ \\
        \hline
        celkem & $-$ & $+$ & $-$ & $+$ \\
        \hline
      \end{tabular}
    \end{center}
    $K = (-1, 0]\cup[4, \infty)$.
  \end{enumerate}
\end{solution}

\section{Soustavy lineárních rovnic}

Často máme víc neznámých a víc podmínek najednou. Například: součet dvou čísel je $10$ a jejich rozdíl je $4$. Která čísla to jsou? Zde potřebujeme vyřešit dvě rovnice zároveň.

\begin{definitionbox}
  \textbf{Soustava dvou lineárních rovnic o dvou neznámých} je soustava
  \begin{align*}
    a_1x + b_1y &= c_1,\\
    a_2x + b_2y &= c_2,
  \end{align*}
  kde $a_1, a_2, b_1, b_2, c_1, c_2\in\mathbb{R}$. \textbf{Řešením} soustavy je každá uspořádaná dvojice $(x, y)\in\mathbb{R}\times\mathbb{R}$, která splňuje \textbf{obě} rovnice zároveň.
\end{definitionbox}

Všimněte si, že řešení je uspořádaná dvojice, tedy prvek kartézského součinu z předchozí kapitoly. Množina řešení soustavy je tedy podmnožinou $\mathbb{R}\times\mathbb{R}$, neboli relace!

\begin{definitionbox}
  \textbf{Metody řešení soustav}:
  \begin{enumerate}
    \item \textbf{Dosazovací metoda}: Z jedné rovnice vyjádříme jednu neznámou a dosadíme ji do druhé rovnice. Tím dostaneme jednu rovnici o jedné neznámé.
    \item \textbf{Sčítací} (eliminační) \textbf{metoda}: Rovnice vynásobíme vhodnými nenulovými čísly tak, aby se po jejich sečtení jedna neznámá vyrušila.
    \item \textbf{Srovnávací metoda}: Z obou rovnic vyjádříme stejnou neznámou a výrazy porovnáme.
    \item \textbf{Grafická metoda}: Každá rovnice popisuje přímku v rovině. Řešením je jejich průsečík.
  \end{enumerate}
\end{definitionbox}

\begin{example}
  Vyřešme soustavu
  \begin{align*}
    x + y &= 10,\\
    x - y &= 4.
  \end{align*}
  \textbf{Dosazovací metoda}: Z první rovnice $y = 10 - x$. Dosadíme do druhé: $x - (10 - x) = 4$, tedy $2x = 14$ a $x = 7$. Pak $y = 10 - 7 = 3$.

  \textbf{Sčítací metoda}: Sečteme obě rovnice: $2x = 14$, tedy $x = 7$. Odečteme je: $2y = 6$, tedy $y = 3$.

  Obě metody dají samozřejmě stejný výsledek, $K = \{(7, 3)\}$. Zkouška: $7 + 3 = 10$ a $7 - 3 = 4$.
\end{example}

\begin{example}
  Sčítací metoda ukáže svou sílu u `ošklivějších' koeficientů. Vyřešme
  \begin{align*}
    3x + 4y &= 2,\\
    5x - 6y &= 16.
  \end{align*}
  Chceme se zbavit $y$. Nejmenší společný násobek $4$ a $6$ je $12$. První rovnici vynásobíme $3$ a druhou $2$:
  \begin{align*}
    9x + 12y &= 6,\\
    10x - 12y &= 32.
  \end{align*}
  Sečtením dostaneme $19x = 38$, tedy $x = 2$. Dosazením do první rovnice $6 + 4y = 2$, tedy $y = -1$. $K = \{(2, -1)\}$.
\end{example}

\subsection{Geometrický význam a počet řešení}

Každá rovnice $ax + by = c$ (kde $a, b$ nejsou obě nulová) popisuje v rovině přímku. Vyřešit soustavu tedy znamená najít body, které leží na obou přímkách zároveň. A jak mohou dvě přímky v rovině ležet? Buď se protínají v jednom bodě, nebo jsou rovnoběžné různé, nebo splývají.

\begin{center}
  \begin{tikzpicture}[scale=0.7]
    % protínající
    \begin{scope}
      \draw[->] (-0.5,0) -- (4,0) node[right] {\small $x$};
      \draw[->] (0,-0.5) -- (0,4) node[above] {\small $y$};
      \draw[thick, colorofwisdom] (-0.3,3.8) -- (3.8,-0.3);
      \draw[thick] (-0.3,0.3) -- (3.5,2.2);
      \filldraw (2.33,1.47) circle (2.5pt);
      \node[below] at (1.75,-0.6) {\small jedno řešení};
    \end{scope}
    % rovnobezne
    \begin{scope}[xshift=6cm]
      \draw[->] (-0.5,0) -- (4,0) node[right] {\small $x$};
      \draw[->] (0,-0.5) -- (0,4) node[above] {\small $y$};
      \draw[thick, colorofwisdom] (-0.3,1) -- (3.8,3.5);
      \draw[thick] (-0.3,-0.2) -- (3.8,2.3);
      \node[below] at (1.75,-0.6) {\small žádné řešení};
    \end{scope}
    % splyvajici
    \begin{scope}[xshift=12cm]
      \draw[->] (-0.5,0) -- (4,0) node[right] {\small $x$};
      \draw[->] (0,-0.5) -- (0,4) node[above] {\small $y$};
      \draw[line width=3pt, colorofwisdom!40] (-0.3,0.5) -- (3.8,3.3);
      \draw[thick] (-0.3,0.5) -- (3.8,3.3);
      \node[below] at (1.75,-0.6) {\small nekonečně mnoho};
    \end{scope}
  \end{tikzpicture}
\end{center}

\begin{definitionbox}
  \textbf{Počet řešení soustavy dvou lineárních rovnic o dvou neznámých}: Soustava má buď
  \begin{itemize}
    \item \textbf{právě jedno} řešení (přímky jsou různoběžné),
    \item \textbf{žádné} řešení (přímky jsou rovnoběžné a různé), nebo
    \item \textbf{nekonečně mnoho} řešení (přímky splývají).
  \end{itemize}
  Nikdy nemůže mít například právě dvě řešení.
\end{definitionbox}

\begin{example}
  Vyřešme soustavu
  \begin{align*}
    2x - y &= 3,\\
    -4x + 2y &= 1.
  \end{align*}
  Vynásobíme první rovnici dvěma a sečteme: $0 = 7$. Spor, soustava nemá řešení, $K = \emptyset$. Geometricky: obě přímky mají tvar $y = 2x + \text{něco}$, takže jsou rovnoběžné.

  Kdyby druhá rovnice byla $-4x + 2y = -6$, dostali bychom $0 = 0$. Druhá rovnice by byla jen dvojnásobkem první, přímky by splývaly a řešením by byly všechny dvojice $(x, 2x - 3)$, kde $x\in\mathbb{R}$. To zapisujeme
  \begin{equation*}
    K = \{(t, 2t - 3)\ |\ t\in\mathbb{R}\}.
  \end{equation*}
  Písmeno $t$ se zde nazývá parametr a každá jeho hodnota nám dá jedno řešení.
\end{example}

\begin{problem}
  Vyřešte soustavy v $\mathbb{R}\times\mathbb{R}$:
  \begin{enumerate}[label=(\alph*)]
    \item $\begin{aligned}[t] 2x + 3y &= 12,\\ x - y &= 1,\end{aligned}$
    \item $\begin{aligned}[t] \tfrac{x}{2} + \tfrac{y}{3} &= 3,\\ \tfrac{x}{3} - \tfrac{y}{4} &= \tfrac{7}{12},\end{aligned}$
    \item $\begin{aligned}[t] 6x - 9y &= 3,\\ 4x - 6y &= 2.\end{aligned}$
  \end{enumerate}
\end{problem}

\begin{solution}
  \begin{enumerate}[label=(\alph*)]
    \item Z druhé $x = 1 + y$, dosazením $2 + 2y + 3y = 12$, tedy $y = 2$ a $x = 3$. $K = \{(3, 2)\}$.
    \item Odstraníme zlomky: první rovnici vynásobíme $6$, druhou $12$. Dostaneme $3x + 2y = 18$ a $4x - 3y = 7$. První vynásobíme $3$, druhou $2$ a sečteme: $17x = 68$, tedy $x = 4$ a z první rovnice $y = \frac{18 - 12}{2} = 3$. $K = \{(4, 3)\}$.
    \item První rovnice je $\frac{3}{2}$-násobkem druhé, přímky splývají. Z druhé rovnice vyjádříme $x = \frac{1 + 3y}{2}$ a za $y$ zvolíme parametr $t$. Tedy $K = \left\{\left(\frac{1 + 3t}{2}, t\right)\ \middle|\ t\in\mathbb{R}\right\}$.
  \end{enumerate}
\end{solution}

\subsection{Soustavy tří rovnic o třech neznámých}

Se třemi neznámými postupujeme stejně: postupně eliminujeme neznámé, až zbude jedna rovnice s jednou neznámou. Nejsystematičtější postup se nazývá \textbf{Gaussova eliminační metoda}. Soustavu upravujeme tak, aby z druhé a třetí rovnice zmizelo $x$ a pak ze třetí rovnice zmizelo $y$. Dostaneme takzvaný \textbf{trojúhelníkový tvar}, ze kterého neznámé dopočítáme odspodu.

\begin{example}
  Vyřešme soustavu
  \begin{align*}
    x + y + z &= 6, \tag{I}\\
    2x - y + z &= 3, \tag{II}\\
    x + 2y - z &= 2. \tag{III}
  \end{align*}
  Eliminujeme $x$: (II) $- 2\cdot$(I) a (III) $-$ (I):
  \begin{align*}
    x + y + z &= 6,\\
    -3y - z &= -9,\\
    y - 2z &= -4.
  \end{align*}
  Eliminujeme $y$ ze třetí rovnice: $3\cdot$(III') $+$ (II'):
  \begin{align*}
    x + y + z &= 6,\\
    -3y - z &= -9,\\
    -7z &= -21.
  \end{align*}
  Odspodu: $z = 3$, pak $-3y - 3 = -9$, tedy $y = 2$, a nakonec $x = 6 - 2 - 3 = 1$. $K = \{(1, 2, 3)\}$. Zkoušku si proveďte sami.
\end{example}

\begin{note}
  Geometricky popisuje každá rovnice $ax + by + cz = d$ rovinu v prostoru. Tři roviny se mohou protnout v jednom bodě (jedno řešení), v přímce (nekonečně mnoho řešení), nebo nemusí mít žádný společný bod. Více o tom v kapitole o analytické geometrii v prostoru.
\end{note}

\begin{problem}
  Vyřešte soustavu
  \begin{align*}
    x + 2y - z &= 3,\\
    3x - y + 2z &= 7,\\
    2x + 3y + z &= 11.
  \end{align*}
\end{problem}

\begin{solution}
  (II) $- 3\cdot$(I): $-7y + 5z = 8$. (III) $- 2\cdot$(I): $-y + 3z = 8$. Z poslední rovnice $y = 3z - 8$, dosazením do předposlední $-21z + 56 + 5z = 8$, tedy $-16z = -48$ a $z = 3$. Pak $y = 1$ a z první rovnice $x = 1 - 2 + 3 = 2$. $K = \{(2, 1, 3)\}$.
\end{solution}

\section{Odbočka --- Determinanty a Cramerovo pravidlo}

Podívejme se na soustavu dvou rovnic úplně obecně. Vyřešme
\begin{align*}
  a_1x + b_1y &= c_1,\\
  a_2x + b_2y &= c_2
\end{align*}
sčítací metodou. První rovnici vynásobíme $b_2$, druhou $-b_1$ a sečteme. Členy s $y$ se vyruší a zůstane
\begin{equation*}
  (a_1b_2 - a_2b_1)x = c_1b_2 - c_2b_1.
\end{equation*}
Podobně vynásobením $-a_2$ a $a_1$ dostaneme
\begin{equation*}
  (a_1b_2 - a_2b_1)y = a_1c_2 - a_2c_1.
\end{equation*}
Ve všech třech výrazech se opakuje stejný vzor: `křížem' vynásobíme a odečteme. To si zaslouží vlastní název.

\begin{definitionbox}
  \textbf{Determinant} matice $2\times 2$ je číslo
  \begin{equation*}
    \begin{vmatrix} a & b\\ c & d \end{vmatrix} = ad - bc.
  \end{equation*}

  \textbf{Cramerovo pravidlo}: Označme
  \begin{equation*}
    D = \begin{vmatrix} a_1 & b_1\\ a_2 & b_2 \end{vmatrix},\quad
    D_x = \begin{vmatrix} c_1 & b_1\\ c_2 & b_2 \end{vmatrix},\quad
    D_y = \begin{vmatrix} a_1 & c_1\\ a_2 & c_2 \end{vmatrix}.
  \end{equation*}
  Pokud $D\neq 0$, má soustava právě jedno řešení
  \begin{equation*}
    x = \frac{D_x}{D},\qquad y = \frac{D_y}{D}.
  \end{equation*}
  Pokud $D = 0$, má soustava buď žádné, nebo nekonečně mnoho řešení.
\end{definitionbox}

Determinant $D_x$ vznikne z $D$ tak, že sloupec koeficientů u $x$ nahradíme pravými stranami. Podobně $D_y$. Výše uvedené odvození je zároveň důkazem Cramerova pravidla. Podmínka $D\neq 0$ geometricky znamená, že přímky nejsou rovnoběžné.\footnote{Determinant má i krásný geometrický význam: $|D|$ je obsah rovnoběžníku, který vytvoří vektory $(a_1, a_2)$ a $(b_1, b_2)$. Pokud je nulový, vektory leží na jedné přímce. Více se dozvíte v kapitole o vektorech.}

\begin{example}
  Vyřešme pomocí Cramerova pravidla soustavu
  \begin{align*}
    3x + 4y &= 2,\\
    5x - 6y &= 16.
  \end{align*}
  \begin{equation*}
    D = \begin{vmatrix} 3 & 4\\ 5 & -6 \end{vmatrix} = -18 - 20 = -38,\quad
    D_x = \begin{vmatrix} 2 & 4\\ 16 & -6 \end{vmatrix} = -12 - 64 = -76,\quad
    D_y = \begin{vmatrix} 3 & 2\\ 5 & 16 \end{vmatrix} = 48 - 10 = 38.
  \end{equation*}
  Tedy $x = \frac{-76}{-38} = 2$ a $y = \frac{38}{-38} = -1$. Stejný výsledek jako dříve sčítací metodou.
\end{example}

\section{Slovní úlohy}

Slovní úlohy jsou místo, kde se matematika potkává se skutečným světem. Postup je vždy stejný:
\begin{definitionbox}
  \textbf{Postup řešení slovní úlohy}:
  \begin{enumerate}
    \item Pozorně si přečíst zadání a určit, co hledáme. To označíme neznámou (a napíšeme si, co neznámá znamená a v jakých jednotkách).
    \item Přeložit vztahy ze zadání do rovnic.
    \item Rovnice vyřešit.
    \item Provést zkoušku \textbf{dosazením do zadání} (ne do rovnic, protože v nich můžeme mít chybu).
    \item Napsat slovní odpověď.
  \end{enumerate}
\end{definitionbox}

\begin{example}[ (pohyb)]
  Z Prahy vyjede v $8{:}00$ nákladní auto rychlostí $60\ \si{km/h}$. V $9{:}00$ za ním vyjede osobní auto rychlostí $90\ \si{km/h}$. Kdy a jak daleko od Prahy osobní auto dohoní nákladní?

  Označme $t$ dobu jízdy osobního auta v hodinách. Nákladní auto jede o hodinu déle, tedy $t + 1$ hodin. V okamžiku setkání ujela obě auta stejnou dráhu:
  \begin{equation*}
    90t = 60(t + 1)\quad\Leftrightarrow\quad 30t = 60\quad\Leftrightarrow\quad t = 2.
  \end{equation*}
  Zkouška: osobní auto ujede za $2$ hodiny $180\ \si{km}$, nákladní za $3$ hodiny také $180\ \si{km}$. Odpověď: Osobní auto dohoní nákladní v $11{:}00$, $180\ \si{km}$ od Prahy.
\end{example}

\begin{example}[ (společná práce)]
  Jeden bazén napustí první čerpadlo za $6$ hodin a druhé za $3$ hodiny. Za jak dlouho ho napustí obě čerpadla dohromady?

  Klíčová myšlenka: první čerpadlo napustí za hodinu $\frac{1}{6}$ bazénu, druhé $\frac{1}{3}$ bazénu. Pokud obě pracují $t$ hodin, napustí dohromady celý bazén, tedy
  \begin{equation*}
    \frac{t}{6} + \frac{t}{3} = 1\quad\Leftrightarrow\quad t + 2t = 6\quad\Leftrightarrow\quad t = 2.
  \end{equation*}
  Odpověď: Obě čerpadla napustí bazén za $2$ hodiny. Častá chyba je výsledky zprůměrovat nebo sečíst. Rozumná kontrola: dohromady to musí být rychlejší než rychlejší z čerpadel, tedy méně než $3$ hodiny.
\end{example}

\begin{example}[ (směsi)]
  Kolik litrů $80\%$ roztoku a kolik litrů $30\%$ roztoku musíme smíchat, abychom dostali $10$ litrů $50\%$ roztoku?

  Označme $x$ litrů $80\%$ roztoku a $y$ litrů $30\%$ roztoku. Celkový objem je $x + y = 10$. Množství čisté látky se nemění: $0{,}8x + 0{,}3y = 0{,}5\cdot 10 = 5$. Z první rovnice $y = 10 - x$, dosazením
  \begin{equation*}
    0{,}8x + 3 - 0{,}3x = 5\quad\Leftrightarrow\quad 0{,}5x = 2\quad\Leftrightarrow\quad x = 4.
  \end{equation*}
  Tedy $y = 6$. Zkouška: $0{,}8\cdot 4 + 0{,}3\cdot 6 = 3{,}2 + 1{,}8 = 5$ litrů čisté látky v $10$ litrech, tedy $50\%$. Odpověď: $4$ litry $80\%$ a $6$ litrů $30\%$ roztoku.
\end{example}

\begin{problem}
  Otci je dnes $42$ let a synovi $12$. Za kolik let bude otec třikrát starší než syn?
\end{problem}

\begin{solution}
  Za $x$ let: $42 + x = 3(12 + x)$, tedy $42 + x = 36 + 3x$ a $x = 3$. Zkouška: za $3$ roky bude otci $45$ a synovi $15$, $45 = 3\cdot 15$. Odpověď: Za $3$ roky.
\end{solution}

\begin{problem}
  Na výletě bylo $35$ dětí. Kdyby přišli ještě $3$ chlapci, bylo by chlapců dvakrát více než dívek. Kolik bylo na výletě chlapců a kolik dívek?
\end{problem}

\begin{solution}
  Označme $c$ počet chlapců a $d$ počet dívek. Pak $c + d = 35$ a $c + 3 = 2d$. Z první $c = 35 - d$, dosazením $38 - d = 2d$, tedy $d = \frac{38}{3}$. To není celé číslo! Úloha tedy nemá řešení. I takový závěr je u slovní úlohy správnou odpovědí. Kdyby přišli $4$ chlapci, dostali bychom $39 = 3d$, tedy $13$ dívek a $22$ chlapců.
\end{solution}

\begin{problem}
  Dva cyklisté vyjedou současně proti sobě z míst vzdálených $72\ \si{km}$. První jede rychlostí $20\ \si{km/h}$, druhý $16\ \si{km/h}$. Za jak dlouho se potkají?
\end{problem}

\begin{solution}
  Za čas $t$ ujedou dohromady $20t + 16t = 72$, tedy $t = 2$. Potkají se za $2$ hodiny.
\end{solution}

\section{Shrnutí a otázky k ústní zkoušce}

\begin{problem}
  Odpovězte vlastními slovy:
  \begin{enumerate}
    \item Co je rovnice a co znamená ji vyřešit? Proč rovnice není výrok?
    \item Jaký je rozdíl mezi ekvivalentní a důsledkovou úpravou? Uveďte příklady. Kdy je zkouška nutná?
    \item Rozeberte všechny případy řešení rovnice $ax + b = 0$.
    \item Jak se mění znak nerovnosti při násobení? Dokažte to.
    \item Proč nesmíme nerovnici násobit výrazem s neznámou? Jak se řeší nerovnice v podílovém tvaru?
    \item Co jsou intervaly a jak se zapisují? Jak souvisí soustava nerovnic s průnikem množin?
    \item Jaké znáte metody řešení soustav lineárních rovnic? Kolik řešení může soustava dvou rovnic o dvou neznámých mít a proč? Nakreslete.
    \item Co je to determinant a k čemu se hodí?
  \end{enumerate}
\end{problem}

\koneckapitoly

\chapter{Kvadratické rovnice a nerovnice}
\label{kap:kvadraticke}
\mathquote{Mathematics is the art of giving the same name to different things.}{Henri Poincaré}

Lineární rovnice jsme v předchozí kapitole vyřešili úplně a definitivně. Teď přidáme neznámou na druhou a najednou se začnou dít zajímavé věci. Rovnice může mít dvě řešení, jedno, nebo žádné. Objeví se diskriminant, Vietovy vzorce a parabola. Kvadratické rovnice jsou pravděpodobně nejčastějším typem rovnic, se kterým se u maturity potkáte, a to nejen v této kapitole, ale i v geometrii, ve fyzice a vlastně všude.

\section{Kvadratická rovnice}

\begin{definitionbox}
  \textbf{Kvadratická rovnice} s neznámou $x$ je rovnice, kterou lze ekvivalentními úpravami převést na tvar
  \begin{equation*}
    ax^2 + bx + c = 0,\qquad a, b, c\in\mathbb{R},\ a\neq 0.
  \end{equation*}
  Výraz $ax^2$ nazýváme \textbf{kvadratický člen}, $bx$ \textbf{lineární člen} a $c$ \textbf{absolutní člen}. Pokud $a = 1$, mluvíme o \textbf{normovaném} tvaru $x^2 + px + q = 0$.
\end{definitionbox}

Podmínka $a\neq 0$ je důležitá. Pokud by $a = 0$, dostali bychom lineární rovnici z předchozí kapitoly a nic by se nedělo.

\subsection{Neúplné kvadratické rovnice}

Než se pustíme do obecného vzorce, podívejme se na dva speciální případy, které se dají vyřešit mnohem rychleji. Používat na ně vzorec je jako jet do vedlejší místnosti autem. Jde to, ale proč.

\begin{definitionbox}
  \textbf{Ryze kvadratická rovnice} ($b = 0$): $ax^2 + c = 0$. Řešíme ji osamostatněním $x^2$:
  \begin{equation*}
    x^2 = -\frac{c}{a}.
  \end{equation*}
  Pokud $-\frac{c}{a} > 0$, má rovnice dvě řešení $x_{1,2} = \pm\sqrt{-\frac{c}{a}}$. Pokud $-\frac{c}{a} = 0$, má jedno řešení $x = 0$. Pokud $-\frac{c}{a} < 0$, nemá v $\mathbb{R}$ řešení.

  \vspace{0.2cm}

  \textbf{Kvadratická rovnice bez absolutního členu} ($c = 0$): $ax^2 + bx = 0$. Řešíme ji vytknutím a větou o nulovém součinu z kapitoly~\ref{kap:vyrazy}:
  \begin{equation*}
    x(ax + b) = 0\quad\Leftrightarrow\quad x = 0\ \lor\ x = -\frac{b}{a}.
  \end{equation*}
\end{definitionbox}

\begin{example}
  Vyřešme:
  \begin{align*}
    4x^2 - 9 = 0 &\quad\Leftrightarrow\quad x^2 = \tfrac{9}{4} &&\Leftrightarrow\quad x = \pm\tfrac{3}{2}, &&K = \left\{-\tfrac{3}{2}, \tfrac{3}{2}\right\},\\
    x^2 + 16 = 0 &\quad\Leftrightarrow\quad x^2 = -16 && &&K = \emptyset,\\
    3x^2 - 12x = 0 &\quad\Leftrightarrow\quad 3x(x - 4) = 0 && &&K = \{0, 4\}.
  \end{align*}
\end{example}

\begin{note}
  U rovnice $3x^2 = 12x$ by vás mohlo napadnout vydělit obě strany číslem $3x$. Nedělejte to! Výraz $3x$ může být nulový a dělením byste přišli o řešení $x = 0$. Tohle je jedna z nejčastějších chyb vůbec. Vytýkejte, nedělte.
\end{note}

\subsection{Doplnění na čtverec a obecný vzorec}

Jak vyřešit rovnici $x^2 + 6x + 5 = 0$? Trik je přepsat levou stranu tak, aby v ní byl `celý čtverec'. Víme, že ${(x + 3)}^2 = x^2 + 6x + 9$. Naše rovnice má místo $9$ jen $5$, takže
\begin{equation*}
  x^2 + 6x + 5 = (x^2 + 6x + 9) - 4 = {(x + 3)}^2 - 4.
\end{equation*}
Rovnice tedy zní ${(x + 3)}^2 = 4$, a to už je v podstatě ryze kvadratická rovnice: $x + 3 = \pm 2$, tedy $x = -1$ nebo $x = -5$. Tomuto postupu se říká \textbf{doplnění na čtverec} a je to jedna z nejužitečnějších technik celé středoškolské matematiky. Uvidíme ji znovu u paraboly, u kružnice i u integrálů.

Pokud stejný postup provedeme s obecnou rovnicí, dostaneme slavný vzorec.

\begin{definitionbox}
  \textbf{Diskriminant} kvadratické rovnice $ax^2 + bx + c = 0$ je číslo
  \begin{equation*}
    D = b^2 - 4ac.
  \end{equation*}

  \textbf{Věta (o řešení kvadratické rovnice)}:
  \begin{itemize}
    \item Pokud $D > 0$, má rovnice dvě různá reálná řešení
    \begin{equation*}
      x_{1,2} = \frac{-b\pm\sqrt{D}}{2a}.
    \end{equation*}
    \item Pokud $D = 0$, má rovnice jedno (tzv. \textbf{dvojnásobné}) řešení $x = -\frac{b}{2a}$.
    \item Pokud $D < 0$, nemá rovnice v $\mathbb{R}$ žádné řešení.
  \end{itemize}
\end{definitionbox}

\begin{proof}
  Protože $a\neq 0$, můžeme rovnici vydělit číslem $a$:
  \begin{equation*}
    x^2 + \frac{b}{a}x + \frac{c}{a} = 0.
  \end{equation*}
  Doplníme na čtverec. Platí ${\left(x + \frac{b}{2a}\right)}^2 = x^2 + \frac{b}{a}x + \frac{b^2}{4a^2}$, a tedy
  \begin{equation*}
    {\left(x + \frac{b}{2a}\right)}^2 - \frac{b^2}{4a^2} + \frac{c}{a} = 0\quad\Leftrightarrow\quad {\left(x + \frac{b}{2a}\right)}^2 = \frac{b^2 - 4ac}{4a^2} = \frac{D}{4a^2}.
  \end{equation*}
  Levá strana je druhá mocnina, a tedy je vždy nezáporná. Jmenovatel $4a^2$ je kladný.
  \begin{itemize}
    \item Pokud $D < 0$, je pravá strana záporná a rovnice nemá řešení.
    \item Pokud $D = 0$, dostaneme ${\left(x + \frac{b}{2a}\right)}^2 = 0$, a tedy $x = -\frac{b}{2a}$.
    \item Pokud $D > 0$, je $x + \frac{b}{2a} = \pm\sqrt{\frac{D}{4a^2}} = \pm\frac{\sqrt{D}}{2|a|}$. Protože před zlomkem stejně stojí $\pm$, můžeme absolutní hodnotu vynechat a dostaneme $x = \frac{-b\pm\sqrt{D}}{2a}$.
  \end{itemize}
\end{proof}

Všimněte si, že vzorec platí i pro $D = 0$, jen obě řešení splynou. Proto se mluví o dvojnásobném kořeni.

\begin{example}\label{ex:kvadr-vzorec}
  Vyřešme rovnici $2x^2 - 3x - 2 = 0$. Máme $a = 2$, $b = -3$, $c = -2$, takže
  \begin{equation*}
    D = {(-3)}^2 - 4\cdot 2\cdot(-2) = 9 + 16 = 25,\qquad x_{1,2} = \frac{3\pm 5}{4}.
  \end{equation*}
  Tedy $x_1 = 2$ a $x_2 = -\frac{1}{2}$, $K = \left\{-\frac{1}{2}, 2\right\}$.

  Rovnice $x^2 - 6x + 9 = 0$ má $D = 36 - 36 = 0$, a tedy jediné řešení $x = 3$. To jsme ale mohli vidět rovnou, protože levá strana je ${(x - 3)}^2$.

  Rovnice $x^2 + 2x + 5 = 0$ má $D = 4 - 20 = -16 < 0$, a tedy $K = \emptyset$. V kapitole o komplexních číslech uvidíme, že řešení přece jen má, jen ne mezi reálnými čísly.
\end{example}

\begin{note}
  Nejčastější chyba ve vzorci je znaménko u $b$. Pokud je $b = -3$, je $-b = 3$ a $b^2 = 9$, ne $-9$. Druhá mocnina záporného čísla je kladná, i když vám kalkulačka po zadání \texttt{-3\^{}2} tvrdí opak.\footnote{Kalkulačka má pravdu, protože $-3^2 = -9$, jak víme z kapitoly~\ref{kap:vyrazy}. Problém je v tom, co jste do ní napsali.}
\end{note}

\begin{problem}
  Vyřešte v $\mathbb{R}$:
  \begin{enumerate}[label=(\alph*)]
    \item $x^2 - 7x + 12 = 0$,
    \item $3x^2 + 5x - 2 = 0$,
    \item $4x^2 - 4x + 1 = 0$,
    \item $x^2 - 2x - 1 = 0$,
    \item $(x + 1)(x - 3) = 2x - 7$.
  \end{enumerate}
\end{problem}

\begin{solution}
  \begin{enumerate}[label=(\alph*)]
    \item $D = 49 - 48 = 1$, $x_{1,2} = \frac{7\pm 1}{2}$, $K = \{3, 4\}$.
    \item $D = 25 + 24 = 49$, $x_{1,2} = \frac{-5\pm 7}{6}$, $K = \left\{-2, \frac{1}{3}\right\}$.
    \item $D = 16 - 16 = 0$, $x = \frac{4}{8} = \frac{1}{2}$. Levá strana je ${(2x - 1)}^2$.
    \item $D = 4 + 4 = 8$, $x_{1,2} = \frac{2\pm\sqrt{8}}{2} = \frac{2\pm 2\sqrt{2}}{2} = 1\pm\sqrt{2}$.
    \item Roznásobíme: $x^2 - 2x - 3 = 2x - 7$, tedy $x^2 - 4x + 4 = 0$, neboli ${(x - 2)}^2 = 0$. $K = \{2\}$.
  \end{enumerate}
\end{solution}

\begin{example}
  Kvadratická rovnice se skrývá i za zlatým řezem z kapitoly~\ref{kap:vyrazy}. Zlatý řez je takové kladné číslo $\varphi$, pro které platí $\varphi^2 = \varphi + 1$. Neboli $\varphi^2 - \varphi - 1 = 0$, a tedy
  \begin{equation*}
    \varphi = \frac{1\pm\sqrt{1 + 4}}{2} = \frac{1\pm\sqrt{5}}{2}.
  \end{equation*}
  Kladné řešení je $\frac{1 + \sqrt{5}}{2}\approx 1{,}618$. Druhé řešení je záporné, takže ho zlatý řez nepotřebuje.
\end{example}

\section{Vietovy vzorce a rozklad na součin}

Podívejme se, co se stane, když kořeny kvadratické rovnice sečteme a vynásobíme.

\begin{definitionbox}
  \textbf{Vietovy vzorce}: Pokud $x_1, x_2$ jsou kořeny rovnice $ax^2 + bx + c = 0$, platí
  \begin{equation*}
    x_1 + x_2 = -\frac{b}{a},\qquad x_1\cdot x_2 = \frac{c}{a}.
  \end{equation*}
  Pro normovanou rovnici $x^2 + px + q = 0$ speciálně $x_1 + x_2 = -p$ a $x_1x_2 = q$.

  \vspace{0.2cm}

  \textbf{Rozklad kvadratického trojčlenu}: Pokud $D\geq 0$, platí pro všechna $x\in\mathbb{R}$
  \begin{equation*}
    ax^2 + bx + c = a(x - x_1)(x - x_2).
  \end{equation*}
  Pokud $D < 0$, nelze trojčlen v $\mathbb{R}$ rozložit na součin lineárních činitelů.
\end{definitionbox}

\begin{proof}
  Sečteme kořeny ze vzorce:
  \begin{equation*}
    x_1 + x_2 = \frac{-b + \sqrt{D}}{2a} + \frac{-b - \sqrt{D}}{2a} = \frac{-2b}{2a} = -\frac{b}{a}.
  \end{equation*}
  Vynásobíme je (a použijeme vzorec $(u + v)(u - v) = u^2 - v^2$):
  \begin{equation*}
    x_1x_2 = \frac{(-b + \sqrt{D})(-b - \sqrt{D})}{4a^2} = \frac{b^2 - D}{4a^2} = \frac{b^2 - b^2 + 4ac}{4a^2} = \frac{c}{a}.
  \end{equation*}
  Rozklad pak plyne z roznásobení pravé strany:
  \begin{equation*}
    a(x - x_1)(x - x_2) = a\left(x^2 - (x_1 + x_2)x + x_1x_2\right) = a\left(x^2 + \frac{b}{a}x + \frac{c}{a}\right) = ax^2 + bx + c.
  \end{equation*}
\end{proof}

François Viète byl francouzský právník, který matematiku dělal ve volném čase. Mimo jiné prý během války rozluštil šifrovanou korespondenci španělského krále tak spolehlivě, že si Španělé stěžovali u papeže, že Francouzi používají černou magii.\footnote{Papež stížnost neřešil. Jeho vlastní luštitelé tu šifru totiž prý také dávno prolomili.}

Vietovy vzorce mají spoustu využití:
\begin{itemize}
  \item Kořeny normované rovnice s `hezkými' kořeny často uhodneme zpaměti. Pro $x^2 - 5x + 6 = 0$ hledáme dvě čísla se součtem $5$ a součinem $6$. To jsou $2$ a $3$.
  \item Můžeme sestavit rovnici se zadanými kořeny.
  \item Můžeme počítat s kořeny, aniž bychom je znali.
  \item Můžeme ověřit, zda jsme rovnici vyřešili správně.
\end{itemize}

\begin{example}
  Sestavme normovanou kvadratickou rovnici s kořeny $3$ a $-2$. Podle Vietových vzorců je $p = -(3 + (-2)) = -1$ a $q = 3\cdot(-2) = -6$. Hledaná rovnice je $x^2 - x - 6 = 0$. Ekvivalentně $(x - 3)(x + 2) = 0$.
\end{example}

\begin{example}
  Bez řešení rovnice $x^2 - 5x + 3 = 0$ spočítejme $x_1^2 + x_2^2$ a $\frac{1}{x_1} + \frac{1}{x_2}$. Víme, že $x_1 + x_2 = 5$ a $x_1x_2 = 3$ (a že kořeny existují, protože $D = 13 > 0$). Pak
  \begin{align*}
    x_1^2 + x_2^2 &= {(x_1 + x_2)}^2 - 2x_1x_2 = 25 - 6 = 19,\\
    \frac{1}{x_1} + \frac{1}{x_2} &= \frac{x_1 + x_2}{x_1x_2} = \frac{5}{3}.
  \end{align*}
  Kořeny jsou přitom $\frac{5\pm\sqrt{13}}{2}$, se kterými by se počítalo podstatně hůř.
\end{example}

\begin{example}
  Rozložme $6x^2 - x - 2$ na součin. $D = 1 + 48 = 49$, kořeny jsou $\frac{1\pm 7}{12}$, tedy $\frac{2}{3}$ a $-\frac{1}{2}$. Potom
  \begin{equation*}
    6x^2 - x - 2 = 6\left(x - \frac{2}{3}\right)\left(x + \frac{1}{2}\right) = (3x - 2)(2x + 1).
  \end{equation*}
  V posledním kroku jsme $6$ rozdělili jako $3\cdot 2$ a každý činitel jsme jedním vynásobili, abychom se zbavili zlomků. Roznásobením si ověřte, že výsledek sedí.
\end{example}

\begin{problem}
  \begin{enumerate}[label=(\alph*)]
    \item Uhodněte kořeny rovnic $x^2 + 7x + 10 = 0$ a $x^2 - x - 20 = 0$.
    \item Sestavte kvadratickou rovnici s celočíselnými koeficienty a kořeny $\frac{1}{2}$ a $-3$.
    \item Rozložte na součin $2x^2 + 5x - 3$ a $x^2 - 6x + 7$.
    \item Zkraťte $\dfrac{x^2 - x - 6}{x^2 - 9}$.
    \item Pro kořeny rovnice $2x^2 - 6x + 1 = 0$ spočítejte $x_1^2x_2 + x_1x_2^2$.
  \end{enumerate}
\end{problem}

\begin{solution}
  \begin{enumerate}[label=(\alph*)]
    \item Součin $10$, součet $-7$: kořeny $-2$ a $-5$. Součin $-20$, součet $1$: kořeny $5$ a $-4$.
    \item $\left(x - \frac{1}{2}\right)(x + 3) = 0$, po vynásobení dvěma $(2x - 1)(x + 3) = 0$, tedy $2x^2 + 5x - 3 = 0$.
    \item Z (b) víme, že $2x^2 + 5x - 3 = (2x - 1)(x + 3)$. Pro druhý trojčlen je $D = 36 - 28 = 8$, kořeny $3\pm\sqrt{2}$, tedy $x^2 - 6x + 7 = (x - 3 - \sqrt{2})(x - 3 + \sqrt{2})$.
    \item $\dfrac{(x - 3)(x + 2)}{(x - 3)(x + 3)} = \dfrac{x + 2}{x + 3}$ pro $x\neq\pm 3$.
    \item $x_1x_2(x_1 + x_2) = \frac{1}{2}\cdot 3 = \frac{3}{2}$.
  \end{enumerate}
\end{solution}

\section{Kvadratické rovnice s parametrem}

U lineárních rovnic s parametrem jsme rozhodovali podle toho, zda je koeficient u $x$ nulový. U kvadratických rovnic rozhoduje hlavně diskriminant. A pozor: pokud parametr stojí u $x^2$, musíme zvlášť vyšetřit i případ, kdy rovnice přestane být kvadratická.

\begin{example}
  Určeme počet řešení rovnice $x^2 + 2x + m = 0$ v závislosti na parametru $m\in\mathbb{R}$. Diskriminant je $D = 4 - 4m$.
  \begin{itemize}
    \item $D > 0\Leftrightarrow m < 1$: dvě řešení $x_{1,2} = -1\pm\sqrt{1 - m}$,
    \item $D = 0\Leftrightarrow m = 1$: jedno řešení $x = -1$,
    \item $D < 0\Leftrightarrow m > 1$: žádné řešení.
  \end{itemize}
\end{example}

\begin{example}
  Vyřešme rovnici $mx^2 - 2x + 1 = 0$ s parametrem $m\in\mathbb{R}$.

  Pro $m = 0$ dostaneme lineární rovnici $-2x + 1 = 0$, tedy $x = \frac{1}{2}$. Tento případ se strašně snadno přehlédne.

  Pro $m\neq 0$ je rovnice kvadratická s $D = 4 - 4m$. Pro $m < 1$ (a $m\neq 0$) má dvě řešení $x_{1,2} = \frac{2\pm\sqrt{4 - 4m}}{2m} = \frac{1\pm\sqrt{1 - m}}{m}$, pro $m = 1$ jedno řešení $x = 1$ a pro $m > 1$ žádné.
  \begin{center}
    \renewcommand{\arraystretch}{1.4}
    \begin{tabular}{|c|c|}
      \hline
      Parametr & Množina řešení \\
      \hline
      $m\in(-\infty, 0)\cup(0, 1)$ & $K = \left\{\frac{1 - \sqrt{1 - m}}{m}, \frac{1 + \sqrt{1 - m}}{m}\right\}$ \\
      $m = 0$ & $K = \left\{\frac{1}{2}\right\}$ \\
      $m = 1$ & $K = \{1\}$ \\
      $m\in(1, \infty)$ & $K = \emptyset$ \\
      \hline
    \end{tabular}
  \end{center}
\end{example}

\begin{problem}
  \begin{enumerate}[label=(\alph*)]
    \item Pro která $p\in\mathbb{R}$ má rovnice $x^2 - px + 4 = 0$ dvojnásobný kořen? Určete ho.
    \item Pro která $k\in\mathbb{R}$ má rovnice $x^2 + kx + 9 = 0$ dva různé reálné kořeny?
    \item Určete $c$ tak, aby jeden kořen rovnice $x^2 - 8x + c = 0$ byl trojnásobkem druhého.
  \end{enumerate}
\end{problem}

\begin{solution}
  \begin{enumerate}[label=(\alph*)]
    \item $D = p^2 - 16 = 0$, tedy $p = \pm 4$. Pro $p = 4$ je kořen $x = 2$, pro $p = -4$ je $x = -2$.
    \item $D = k^2 - 36 > 0$, tedy $|k| > 6$, neboli $k\in(-\infty, -6)\cup(6, \infty)$.
    \item Podle Vieta $x_1 + 3x_1 = 8$, tedy $x_1 = 2$, $x_2 = 6$ a $c = x_1x_2 = 12$.
  \end{enumerate}
\end{solution}

\section{Rovnice vyšších stupňů}

Pro rovnice třetího a čtvrtého stupně existují obecné vzorce, ale jsou tak ošklivé, že se u maturity nepoužívají.\footnote{Vzorec pro kubickou rovnici má na svědomí jeden z nejdramatičtějších příběhů dějin matematiky, viz příběh na konci kapitoly.} Pro pátý a vyšší stupeň obecný vzorec neexistuje, jak jsme zmínili u Galoise v kapitole~\ref{kap:rovnice}. U maturity se proto objevují jen rovnice, které se dají na kvadratické převést nějakým trikem.

\subsection{Substituce}

\begin{definitionbox}
  \textbf{Substituce}: Pokud se v rovnici opakovaně vyskytuje stejný výraz $V(x)$, označíme ho novou proměnnou $t = V(x)$, vyřešíme rovnici pro $t$ a pak se vrátíme k $x$ řešením rovnic $V(x) = t_i$.

  \vspace{0.2cm}

  Typickým příkladem je \textbf{bikvadratická rovnice} $ax^4 + bx^2 + c = 0$, kterou substitucí $t = x^2$ převedeme na kvadratickou rovnici $at^2 + bt + c = 0$. Záporná řešení $t$ zahodíme, protože $x^2$ nemůže být záporné.
\end{definitionbox}

\begin{example}
  Vyřešme $x^4 - 5x^2 + 4 = 0$. Substituce $t = x^2$: $t^2 - 5t + 4 = 0$, tedy $t_1 = 1$, $t_2 = 4$. Zpět: $x^2 = 1$ dává $x = \pm 1$ a $x^2 = 4$ dává $x = \pm 2$. Celkem $K = \{-2, -1, 1, 2\}$.

  Rovnice $x^4 + 3x^2 - 4 = 0$ dává $t^2 + 3t - 4 = 0$, tedy $t_1 = 1$, $t_2 = -4$. Rovnice $x^2 = -4$ nemá řešení, takže $K = \{-1, 1\}$.
\end{example}

\begin{example}
  Substituce funguje i pro složitější výrazy. Vyřešme ${(x^2 - 2x)}^2 - 2(x^2 - 2x) - 3 = 0$. Substituce $t = x^2 - 2x$ dává $t^2 - 2t - 3 = 0$, tedy $t = 3$ nebo $t = -1$.
  \begin{itemize}
    \item $x^2 - 2x = 3\Leftrightarrow x^2 - 2x - 3 = 0\Leftrightarrow x\in\{-1, 3\}$,
    \item $x^2 - 2x = -1\Leftrightarrow{(x - 1)}^2 = 0\Leftrightarrow x = 1$.
  \end{itemize}
  $K = \{-1, 1, 3\}$. Kdo by roznásobil závorku, dostal by rovnici čtvrtého stupně a pravděpodobně by se rozplakal.
\end{example}

\subsection{Celočíselné kořeny}

\begin{definitionbox}
  \textbf{Věta (o celočíselných kořenech)}: Nechť $P(x) = a_nx^n + \cdots + a_1x + a_0$ je mnohočlen s \textbf{celočíselnými} koeficienty. Pokud je celé číslo $r$ jeho kořenem, pak $r$ dělí absolutní člen $a_0$.
\end{definitionbox}

\begin{proof}
  Pokud $P(r) = 0$, pak $a_nr^n + \cdots + a_1r + a_0 = 0$, a tedy
  \begin{equation*}
    a_0 = -r\left(a_nr^{n-1} + \cdots + a_2r + a_1\right).
  \end{equation*}
  Výraz v závorce je celé číslo, takže $a_0$ je celočíselným násobkem $r$.
\end{proof}

Strategie je tedy jasná. Vyzkoušíme dělitele absolutního členu. Když najdeme kořen $r$, vydělíme mnohočlen dvojčlenem $(x - r)$, jak jsme to dělali v kapitole~\ref{kap:vyrazy}, a dostaneme mnohočlen o stupeň nižší.

\begin{example}
  Vyřešme $x^3 - 2x^2 - 5x + 6 = 0$. Dělitelé čísla $6$ jsou $\pm 1, \pm 2, \pm 3, \pm 6$. Zkusíme $x = 1$: $1 - 2 - 5 + 6 = 0$. Máme kořen! Vydělíme:
  \begin{equation*}
    (x^3 - 2x^2 - 5x + 6) : (x - 1) = x^2 - x - 6 = (x - 3)(x + 2).
  \end{equation*}
  Tedy $x^3 - 2x^2 - 5x + 6 = (x - 1)(x - 3)(x + 2)$ a $K = \{-2, 1, 3\}$.
\end{example}

\begin{note}
  Věta nám nezaručuje, že nějaký celočíselný kořen existuje. Například $x^3 - 2 = 0$ má jediný reálný kořen $\sqrt[3]{2}$, který celý není. Věta nám jen říká, kde hledat, pokud tam něco je.
\end{note}

\begin{problem}
  Vyřešte v $\mathbb{R}$:
  \begin{enumerate}[label=(\alph*)]
    \item $x^4 - 10x^2 + 9 = 0$,
    \item $2x^4 - 7x^2 - 4 = 0$,
    \item $x^3 + x^2 - 4x - 4 = 0$,
    \item $x^3 - 3x + 2 = 0$.
  \end{enumerate}
\end{problem}

\begin{hint}
  U (c) zkuste postupné vytýkání z kapitoly~\ref{kap:vyrazy}.
\end{hint}

\begin{solution}
  \begin{enumerate}[label=(\alph*)]
    \item $t^2 - 10t + 9 = 0$, $t\in\{1, 9\}$, $K = \{-3, -1, 1, 3\}$.
    \item $2t^2 - 7t - 4 = 0$, $D = 49 + 32 = 81$, $t = \frac{7\pm 9}{4}$, tedy $t = 4$ nebo $t = -\frac{1}{2}$. Záporné $t$ zahodíme, $K = \{-2, 2\}$.
    \item $x^2(x + 1) - 4(x + 1) = (x + 1)(x^2 - 4) = (x + 1)(x - 2)(x + 2)$, $K = \{-2, -1, 2\}$.
    \item $x = 1$ je kořen ($1 - 3 + 2 = 0$). Dělením $x^2 + x - 2 = (x - 1)(x + 2)$, celkem $(x - 1)^2(x + 2)$ a $K = \{-2, 1\}$. Kořen $1$ je dvojnásobný.
  \end{enumerate}
\end{solution}

\section{Kvadratické nerovnice}

\begin{definitionbox}
  \textbf{Kvadratická nerovnice} je nerovnice, kterou lze převést na tvar $ax^2 + bx + c > 0$ (případně $<$, $\geq$, $\leq$), kde $a\neq 0$.

  \vspace{0.2cm}

  \textbf{Postup řešení}:
  \begin{enumerate}
    \item Převedeme vše na jednu stranu.
    \item Najdeme kořeny $x_1, x_2$ rovnice $ax^2 + bx + c = 0$.
    \item Pokud kořeny existují, rozložíme trojčlen na $a(x - x_1)(x - x_2)$ a použijeme metodu nulových bodů z kapitoly~\ref{kap:rovnice}.
    \item Pokud $D < 0$, má trojčlen všude stejné znaménko jako $a$. Nerovnice tedy platí buď pro všechna $x$, nebo pro žádné.
  \end{enumerate}
\end{definitionbox}

Nejrychlejší je ale představit si graf. Funkce $y = ax^2 + bx + c$ má jako graf parabolu (podrobně ji rozebereme v kapitole o elementárních funkcích). Pokud $a > 0$, je parabola otevřená nahoru, pokud $a < 0$, dolů. Kořeny jsou body, kde parabola protíná osu $x$. Řešit nerovnici $ax^2 + bx + c > 0$ pak znamená zjistit, kde je parabola nad osou $x$.

\begin{center}
  \begin{tikzpicture}[scale=0.6]
    \begin{scope}
      \draw[->] (-3.2,0) -- (3.4,0) node[right] {\small $x$};
      \draw[->] (0,-2.6) -- (0,3.4) node[above] {\small $y$};
      \draw[thick, colorofwisdom, domain=-2.6:2.6, samples=60] plot (\x, {0.5*\x*\x - 2});
      \filldraw (-2,0) circle (2pt) node[below left] {\small $x_1$};
      \filldraw (2,0) circle (2pt) node[below right] {\small $x_2$};
      \node at (0,-3.4) {\small $a > 0,\ D > 0$};
    \end{scope}
    \begin{scope}[xshift=8cm]
      \draw[->] (-3.2,0) -- (3.4,0) node[right] {\small $x$};
      \draw[->] (0,-2.6) -- (0,3.4) node[above] {\small $y$};
      \draw[thick, colorofwisdom, domain=-0.7:2.7, samples=60] plot (\x, {(\x-1)*(\x-1)});
      \filldraw (1,0) circle (2pt) node[below] {\small $x_0$};
      \node at (0,-3.4) {\small $a > 0,\ D = 0$};
    \end{scope}
    \begin{scope}[xshift=16cm]
      \draw[->] (-3.2,0) -- (3.4,0) node[right] {\small $x$};
      \draw[->] (0,-2.6) -- (0,3.4) node[above] {\small $y$};
      \draw[thick, colorofwisdom, domain=-2.4:2.4, samples=60] plot (\x, {0.4*\x*\x + 0.8});
      \node at (0,-3.4) {\small $a > 0,\ D < 0$};
    \end{scope}
  \end{tikzpicture}
\end{center}

\begin{example}
  Vyřešme nerovnice:
  \begin{enumerate}[label=(\alph*)]
    \item $x^2 - x - 6 < 0$. Rozklad $(x - 3)(x + 2) < 0$. Parabola je otevřená nahoru a pod osou je mezi kořeny, $K = (-2, 3)$.
    \item $x^2 - 4x + 4 > 0$. Levá strana je ${(x - 2)}^2$, která je kladná všude kromě $x = 2$. $K = \mathbb{R}\setminus\{2\}$.
    \item $x^2 + x + 1\geq 0$. $D = 1 - 4 < 0$ a $a > 0$, takže je trojčlen vždy kladný. $K = \mathbb{R}$.
    \item $-x^2 + 2x + 3\geq 0$. Vynásobíme $-1$ (a otočíme znak): $x^2 - 2x - 3\leq 0$, tedy $(x - 3)(x + 1)\leq 0$ a $K = [-1, 3]$.
  \end{enumerate}
\end{example}

\begin{note}
  Proč je $x^2 + x + 1$ vždy kladné? Doplněním na čtverec: $x^2 + x + 1 = {\left(x + \frac{1}{2}\right)}^2 + \frac{3}{4}$. Druhá mocnina je nezáporná a přičteme k ní kladné číslo. To je důkaz celého bodu 4 v postupu výše, jen pro konkrétní čísla.
\end{note}

\begin{problem}
  Vyřešte v $\mathbb{R}$:
  \begin{enumerate}[label=(\alph*)]
    \item $x^2 - 5x + 4\geq 0$,
    \item $2x^2 + x - 1 < 0$,
    \item $x^2 + 4 < 4x$,
    \item $-2x^2 + 3x - 5 > 0$,
    \item $x^2\leq 9$.
  \end{enumerate}
\end{problem}

\begin{solution}
  \begin{enumerate}[label=(\alph*)]
    \item $(x - 1)(x - 4)\geq 0$, $K = (-\infty, 1]\cup[4, \infty)$.
    \item Kořeny $\frac{-1\pm 3}{4}$, tedy $-1$ a $\frac{1}{2}$. $K = \left(-1, \frac{1}{2}\right)$.
    \item ${(x - 2)}^2 < 0$, to neplatí nikdy. $K = \emptyset$.
    \item Vynásobíme $-1$: $2x^2 - 3x + 5 < 0$. $D = 9 - 40 < 0$ a $a > 0$, levá strana je vždy kladná. $K = \emptyset$.
    \item $(x - 3)(x + 3)\leq 0$, $K = [-3, 3]$. Častá chyba je odmocnit a napsat $x\leq 3$. Tím se ztratí celá levá polovina.
  \end{enumerate}
\end{solution}

\section{Soustavy s kvadratickou rovnicí}

Pokud je v soustavě alespoň jedna rovnice kvadratická, metody z předchozí kapitoly (sčítací, matice) přestanou fungovat. Zůstane nám ale dosazovací metoda, která funguje vždy.

\begin{example}
  Vyřešme soustavu
  \begin{align*}
    x + y &= 5,\\
    xy &= 6.
  \end{align*}
  Z první rovnice $y = 5 - x$, dosazením $x(5 - x) = 6$, tedy $x^2 - 5x + 6 = 0$ a $x\in\{2, 3\}$. Dopočítáme $y$: $K = \{(2, 3), (3, 2)\}$.

  Alternativně: podle Vietových vzorců jsou $x$ a $y$ kořeny rovnice $t^2 - 5t + 6 = 0$. Vietovy vzorce zde fungují obráceně a ušetří nám práci.
\end{example}

\begin{example}
  Vyřešme soustavu
  \begin{align*}
    x^2 + y^2 &= 25,\\
    x - y &= 1.
  \end{align*}
  Z druhé rovnice $x = y + 1$, dosazením ${(y + 1)}^2 + y^2 = 25$, tedy $2y^2 + 2y - 24 = 0$, neboli $y^2 + y - 12 = 0$, a tedy $y = 3$ nebo $y = -4$. Odpovídající $x$ jsou $4$ a $-3$. $K = \{(4, 3), (-3, -4)\}$.

  Geometricky: první rovnice je kružnice se středem v počátku a poloměrem $5$, druhá je přímka. Přímka protíná kružnici ve dvou bodech. Kdyby ji jen tečně škrtla, měli bychom jedno řešení, a kdyby ji minula, žádné. Více v kapitole o kuželosečkách.
\end{example}

\begin{problem}
  Vyřešte soustavy:
  \begin{enumerate}[label=(\alph*)]
    \item $x - y = 2$, $xy = 15$,
    \item $x^2 - y^2 = 16$, $x + y = 8$,
    \item $y = x^2 - 3$, $y = 2x$.
  \end{enumerate}
\end{problem}

\begin{solution}
  \begin{enumerate}[label=(\alph*)]
    \item $x = y + 2$, $(y + 2)y = 15$, $y^2 + 2y - 15 = 0$, $y\in\{3, -5\}$. $K = \{(5, 3), (-3, -5)\}$.
    \item $x^2 - y^2 = (x - y)(x + y) = 8(x - y) = 16$, tedy $x - y = 2$. Spolu s $x + y = 8$: $K = \{(5, 3)\}$. Zde stačilo všimnout si vzorce.
    \item $x^2 - 3 = 2x$, $x^2 - 2x - 3 = 0$, $x\in\{-1, 3\}$. $K = \{(-1, -2), (3, 6)\}$. Geometricky průsečíky paraboly a přímky.
  \end{enumerate}
\end{solution}

\section{Slovní úlohy}

\begin{example}
  Obdélníková zahrada má obvod $34\ \si{m}$ a obsah $60\ \si{m^2}$. Jaké má rozměry?

  Označme strany $a, b$. Pak $2(a + b) = 34$ a $ab = 60$, tedy $a + b = 17$ a $ab = 60$. Podle Vieta jsou $a, b$ kořeny $t^2 - 17t + 60 = 0$, tedy $t\in\{5, 12\}$. Zahrada má rozměry $5\ \si{m}\times 12\ \si{m}$. Mimochodem, její úhlopříčka je přesně $13\ \si{m}$, protože $5^2 + 12^2 = 169$.
\end{example}

\begin{example}
  Kámen pustíme z věže a za $t$ sekund spadne o $s = 5t^2$ metrů (zaokrouhlili jsme $g\approx 10\ \si{m/s^2}$). Za jak dlouho dopadne kámen z výšky $45\ \si{m}$?

  $5t^2 = 45$, tedy $t^2 = 9$ a $t = \pm 3$. Záporný čas nedává smysl, takže kámen dopadne za $3$ sekundy. U slovních úloh je vždy potřeba zkontrolovat, zda řešení dává smysl ve skutečném světě. Matematika nám klidně nabídne i kámen, který dopadl před třemi sekundami, ještě než jsme ho pustili.
\end{example}

\begin{problem}
  Součin dvou po sobě jdoucích přirozených čísel je $156$. Která jsou to čísla?
\end{problem}

\begin{solution}
  $n(n + 1) = 156$, $n^2 + n - 156 = 0$, $D = 1 + 624 = 625$, $n = \frac{-1\pm 25}{2}$. Přirozené řešení je $n = 12$, čísla jsou $12$ a $13$.
\end{solution}

\begin{problem}
  Skupina kamarádů si pronajala chatu za $12\,000$ Kč a náklady si rozdělila rovným dílem. Kdyby jich jelo o dva víc, zaplatil by každý o $500$ Kč méně. Kolik jich jelo?
\end{problem}

\begin{solution}
  Označme $n$ počet lidí. Pak $\frac{12000}{n} - \frac{12000}{n + 2} = 500$. Vynásobíme $n(n + 2)$: $12000(n + 2) - 12000n = 500n(n + 2)$, tedy $24000 = 500n^2 + 1000n$, neboli $n^2 + 2n - 48 = 0$. Kořeny jsou $6$ a $-8$, jelo jich tedy $6$. Zkouška: $2000 - 1500 = 500$.
\end{solution}

\section{Příběh --- Kubická rovnice a jeden skandál}

Vzorec pro kvadratickou rovnici znali v nějaké podobě už Babyloňané před čtyřmi tisíci lety. Na vzorec pro rovnici třetího stupně se čekalo až do renesanční Itálie, a když konečně přišel, provázel ho skandál.

Kolem roku 1515 našel řešení některých kubických rovnic Scipione del Ferro, ale nikomu ho neřekl. V té době se matematika dělala formou veřejných soubojů, ve kterých si soupeři navzájem zadávali úlohy, a tajný vzorec byl cennější než diplom. Nezávisle na něm vzorec objevil Niccolò Fontana, kterému se říkalo Tartaglia, tedy `koktavec'. Jako dítě mu totiž při francouzském plenění města Brescie voják rozsekl mečem čelist a od té doby špatně mluvil.\footnote{Matka ho prý ošetřovala tak, že mu rány olizovala, jako to dělají psi. Přežil. Vousy si pak nechal narůst, aby zakryl jizvy.}

Tartagliův vzorec chtěl za každou cenu znát Gerolamo Cardano, lékař, astrolog a vášnivý hráč. Tartaglia mu ho nakonec prozradil pod podmínkou, že ho nikdy nezveřejní. Cardano ho v roce 1545 zveřejnil ve své knize \textit{Ars Magna}. Bránil se tím, že si mezitím přečetl starší poznámky del Ferra, takže vlastně publikuje jeho vzorec, ne Tartagliův. Tartaglia to neviděl stejně a zbytek života strávil hádkami. Dnes se vzorci říká \textbf{Cardanův vzorec}, takže si můžete udělat obrázek o tom, jak funguje spravedlnost.

Ani Cardano ale nedopadl dobře. Jeho nejstaršího syna popravili za otrávení manželky, sám byl uvězněn inkvizicí za to, že sestavil horoskop Ježíše Krista, a podle legendy si předpověděl datum vlastní smrti. Když se ten den blížil a on byl stále zdravý, prý přestal jíst, aby jeho předpověď vyšla.\footnote{Pro přesnost: tohle je opravdu jen legenda. Ale stojí za to ji vyprávět.}

Mimochodem, právě při řešení kubických rovnic se poprvé objevily odmocniny ze záporných čísel. Cardanův vzorec totiž někdy vede přes $\sqrt{-1}$, i když je výsledek nakonec úplně obyčejné reálné číslo. Komplexní čísla tedy nevznikla proto, že by je někdo chtěl, ale proto, že se jim nedalo vyhnout. K tomu se vrátíme v kapitole o komplexních číslech.

\section{Shrnutí a otázky k ústní zkoušce}

\begin{problem}
  Odpovězte vlastními slovy:
  \begin{enumerate}
    \item Co je kvadratická rovnice? Jak se řeší neúplné kvadratické rovnice?
    \item Odvoďte vzorec pro kořeny kvadratické rovnice doplněním na čtverec.
    \item Co je diskriminant a co nám říká o počtu řešení?
    \item Vyslovte a dokažte Vietovy vzorce. K čemu se hodí?
    \item Jak rozložíte kvadratický trojčlen na součin? Kdy to nejde?
    \item Jak se řeší kvadratická rovnice s parametrem? Na co si dát pozor, když parametr stojí u $x^2$?
    \item Jak se řeší bikvadratické rovnice a rovnice vyšších stupňů? Vyslovte a dokažte větu o celočíselných kořenech.
    \item Jak se řeší kvadratické nerovnice? Načrtněte, jak souvisí s grafem paraboly.
    \item Jak se řeší soustavy, ve kterých je kvadratická rovnice? Jaký mají geometrický význam?
  \end{enumerate}
\end{problem}

\koneckapitoly

\chapter{Absolutní hodnota a iracionální rovnice}
\label{kap:absolutni}
\mathquote{Mathematics is not about numbers, equations, computations, or algorithms: it is about understanding.}{William Paul Thurston}

V této kapitole se podíváme na dva typy rovnic, které vypadají nevinně, ale umí potrápit. Rovnice s absolutní hodnotou, u kterých musíme rozebírat případy, a iracionální rovnice, tedy rovnice s neznámou pod odmocninou, u kterých se poprvé setkáme s tím, že správně provedený výpočet vyprodukuje špatné řešení. Obojí spojuje jedna věc: musíme hodně dávat pozor na to, co přesně děláme.

\section{Absolutní hodnota}

\begin{definitionbox}
  \textbf{Absolutní hodnota} reálného čísla $a$ je číslo
  \begin{equation*}
    |a| = \begin{cases}
      a & \text{pro } a\geq 0,\\
      -a & \text{pro } a < 0.
    \end{cases}
  \end{equation*}
  \textbf{Geometrický význam}: $|a|$ je vzdálenost čísla $a$ od nuly na číselné ose. Obecněji $|a - b|$ je \textbf{vzdálenost} čísel $a$ a $b$.
\end{definitionbox}

Definice na první pohled děsí tím, že se v ní objevuje $-a$. Absolutní hodnota přece nemůže být záporná! A taky není. Pokud je $a < 0$, je $-a$ kladné číslo. Například $|-5| = -(-5) = 5$. Absolutní hodnota z čísla prostě udělá kladné číslo, bez ohledu na to, čím si prošlo.

\begin{center}
  \begin{tikzpicture}
    \draw[->, thick] (-4.5,0) -- (4.5,0) node[right] {$x$};
    \foreach \x in {-4,...,4} \draw (\x,0.1) -- (\x,-0.1) node[below] {\small $\x$};
    \draw[colorofwisdom, line width=2pt] (-3,0.4) -- (0,0.4);
    \node[above, colorofwisdom] at (-1.5,0.4) {$|-3| = 3$};
    \draw[line width=2pt] (1,-0.8) -- (4,-0.8);
    \node[below] at (2.5,-0.8) {$|4 - 1| = |1 - 4| = 3$};
  \end{tikzpicture}
\end{center}

Jak vypadá graf funkce $y = |x|$? Pro $x\geq 0$ je to přímka $y = x$, pro $x < 0$ přímka $y = -x$. Jinak řečeno: vezmeme přímku $y = x$ a její část pod osou $x$ \textit{překlopíme} nahoru. Vznikne tvar písmene V se špičkou v počátku. Posouvání pak funguje jako u každé jiné funkce. Graf $y = |x - 2|$ má špičku v bodě $2$, tedy tam, kde je výraz uvnitř absolutní hodnoty nulový, a graf $y = |x| - 1$ je celé V posunuté o $1$ dolů:
\begin{center}
  \begin{tikzpicture}[scale=0.55]
    \begin{scope}
      \draw[->] (-2.8,0) -- (3.3,0) node[right] {\small $x$};
      \draw[->] (0,-2) -- (0,3.4) node[above] {\small $y$};
      \draw[dashed, gray, thick] (-1.8,-1.8) -- (0,0);
      \draw[->, gray] (-1.3,-1.1) to[bend left=40] (-1.3,1.1);
      \draw[very thick, colorofwisdom] (-2.6,2.6) -- (0,0) -- (2.9,2.9);
      \node[font=\small] at (0.3,-2.7) {$y = |x|$};
    \end{scope}
    \begin{scope}[xshift=7.3cm]
      \draw[->] (-1.4,0) -- (5,0) node[right] {\small $x$};
      \draw[->] (0,-2) -- (0,3.4) node[above] {\small $y$};
      \draw[very thick, colorofwisdom] (-1,3) -- (2,0) -- (4.8,2.8);
      \filldraw (2,0) circle (3pt) node[below, font=\small] {$2$};
      \node[font=\small] at (1.8,-2.7) {$y = |x - 2|$};
    \end{scope}
    \begin{scope}[xshift=15.5cm]
      \draw[->] (-2.8,0) -- (3.3,0) node[right] {\small $x$};
      \draw[->] (0,-2) -- (0,3.4) node[above] {\small $y$};
      \draw[dashed, gray] (-2.6,2.6) -- (0,0) -- (2.9,2.9);
      \draw[very thick, colorofwisdom] (-2.6,1.6) -- (0,-1) -- (2.9,1.9);
      \filldraw (-1,0) circle (3pt) node[below left, font=\small] {$-1$};
      \filldraw (1,0) circle (3pt) node[below right, font=\small] {$1$};
      \node[font=\small] at (0.3,-2.7) {$y = |x| - 1$};
    \end{scope}
  \end{tikzpicture}
\end{center}
Ostrá špička je typickým poznávacím znamením absolutní hodnoty. V kapitole o derivacích uvidíme, že právě v ní nemá graf tečnu.

\begin{definitionbox}
  \textbf{Vlastnosti absolutní hodnoty}: Pro všechna $a, b\in\mathbb{R}$ platí
  \begin{enumerate}
    \item $|a|\geq 0$ a $|a| = 0\Leftrightarrow a = 0$,
    \item $|-a| = |a|$,
    \item $|a\cdot b| = |a|\cdot|b|$ a $\left|\frac{a}{b}\right| = \frac{|a|}{|b|}$ pro $b\neq 0$,
    \item $|a|^2 = a^2$ a $\sqrt{a^2} = |a|$,
    \item $-|a|\leq a\leq|a|$,
    \item $|a + b|\leq|a| + |b|$ \quad (\textbf{trojúhelníková nerovnost}).
  \end{enumerate}
\end{definitionbox}

Většina vlastností plyne přímo z definice rozebráním případů $a\geq 0$ a $a < 0$. Zkuste si to. Nejzajímavější je ta poslední, a tak si ji dokážeme.

\begin{proof}[(trojúhelníkové nerovnosti)]
  Obě strany jsou nezáporné, takže stačí dokázat nerovnost pro jejich druhé mocniny.\footnote{Pro nezáporná čísla $u, v$ platí $u\leq v\Leftrightarrow u^2\leq v^2$. Pro záporná čísla by to neplatilo, viz $-3\leq 1$, ale $9 > 1$.} Podle vlastnosti 4 je
  \begin{equation*}
    |a + b|^2 = {(a + b)}^2 = a^2 + 2ab + b^2,\qquad {(|a| + |b|)}^2 = a^2 + 2|a||b| + b^2 = a^2 + 2|ab| + b^2.
  \end{equation*}
  Podle vlastnosti 5 je $ab\leq|ab|$, a tedy $|a + b|^2\leq{(|a| + |b|)}^2$.
\end{proof}

Proč se jí říká trojúhelníková? Ve vektorech (kapitola o analytické geometrii) bude mít tvar $|\vec{u} + \vec{v}|\leq|\vec{u}| + |\vec{v}|$ a říká, že žádná strana trojúhelníku není delší než součet zbylých dvou. Neboli: přímá cesta je vždy nejkratší. Rovnost nastává, právě když $a$ a $b$ mají stejné znaménko (nebo je jedno z nich nula).

Už na číselné ose je ale vidět, o co jde. Číslo $a + b$ najdeme tak, že z nuly uděláme krok $a$ a pak krok $b$. Celkem jsme ušli $|a| + |b|$. Pokud oba kroky míří stejným směrem, skončíme přesně v této vzdálenosti od nuly. Pokud míří proti sobě, druhý krok část cesty `vrátí' a skončíme blíž:
\begin{center}
  \begin{tikzpicture}[scale=0.85]
    \begin{scope}
      \draw[->] (-1,0) -- (4.5,0);
      \foreach \x in {0,...,4} \draw (\x,0.08) -- (\x,-0.08) node[below, font=\scriptsize] {$\x$};
      \draw[->, very thick, colorofwisdom] (0,0.3) -- (2,0.3) node[midway, above, font=\small] {$a = 2$};
      \draw[->, very thick] (2,0.3) -- (3.5,0.3) node[midway, above, font=\small] {$b = 1{,}5$};
      \filldraw (3.5,0) circle (2.5pt);
      \node[font=\small] at (1.75,-1) {$|a + b| = 3{,}5 = |a| + |b|$};
    \end{scope}
    \begin{scope}[xshift=7.5cm]
      \draw[->] (-1,0) -- (4.5,0);
      \foreach \x in {0,...,4} \draw (\x,0.08) -- (\x,-0.08) node[below, font=\scriptsize] {$\x$};
      \draw[->, very thick, colorofwisdom] (0,0.3) -- (3,0.3) node[midway, above, font=\small] {$a = 3$};
      \draw[->, very thick] (3,1.05) -- (1,1.05) node[midway, above, font=\small] {$b = -2$};
      \draw[dotted] (3,0.3) -- (3,1.05);
      \filldraw (1,0) circle (2.5pt);
      \node[font=\small] at (1.75,-1) {$|a + b| = 1 < 5 = |a| + |b|$};
    \end{scope}
  \end{tikzpicture}
\end{center}

\begin{problem}
  \begin{enumerate}[label=(\alph*)]
    \item Spočítejte $|3 - \pi|$, $|\sqrt{2} - 1|$ a $|1 - \sqrt{3}|$.
    \item Zjednodušte $|x - 2| + |x + 1|$ pro $x\in(-1, 2)$.
    \item Najděte $a, b$ taková, aby $|a + b| < |a| + |b|$.
  \end{enumerate}
\end{problem}

\begin{solution}
  \begin{enumerate}[label=(\alph*)]
    \item $\pi > 3$, takže $|3 - \pi| = \pi - 3$. Dále $\sqrt{2} > 1$, takže $|\sqrt{2} - 1| = \sqrt{2} - 1$, a $\sqrt{3} > 1$, takže $|1 - \sqrt{3}| = \sqrt{3} - 1$.
    \item Pro $x\in(-1, 2)$ je $x - 2 < 0$ a $x + 1 > 0$, takže výraz je $-(x - 2) + (x + 1) = 3$. Pro všechna taková $x$ vychází stejně. To je geometricky jasné: součet vzdáleností bodu mezi $-1$ a $2$ od těchto dvou bodů je přesně vzdálenost mezi nimi.
    \item Stačí čísla s opačným znaménkem, například $a = 1$, $b = -1$: $0 < 2$.
  \end{enumerate}
\end{solution}

\section{Rovnice a nerovnice s absolutní hodnotou}

\subsection{Jednoduché případy}

Pokud je v rovnici jen jedna absolutní hodnota tvaru $|x - a|$, nejrychlejší je myslet geometricky.

\begin{definitionbox}
  Pro $b > 0$ platí:
  \begin{align*}
    |x - a| = b &\quad\Leftrightarrow\quad x = a - b\ \lor\ x = a + b,\\
    |x - a| < b &\quad\Leftrightarrow\quad a - b < x < a + b,\\
    |x - a| > b &\quad\Leftrightarrow\quad x < a - b\ \lor\ x > a + b.
  \end{align*}
  Slovy: hledáme čísla, jejichž vzdálenost od $a$ je rovna, menší nebo větší než $b$. Pro $b < 0$ nemá rovnice $|x - a| = b$ řešení (vzdálenost nemůže být záporná) a pro $b = 0$ má jediné řešení $x = a$.
\end{definitionbox}

Představte si, že do bodu $a$ zapíchneme kružítko s rozevřením $b$. Body, kam kružítko dosáhne přesně, řeší rovnici. Body uvnitř řeší nerovnici `$<$' a body venku nerovnici `$>$':
\begin{center}
  \begin{tikzpicture}
    \draw[->, thick] (-4.5,0) -- (4.5,0) node[right] {$x$};
    \draw[line width=3pt, colorofwisdom] (-2.5,0) -- (2.5,0);
    \draw[line width=3pt, black!25] (-4.3,0) -- (-2.5,0);
    \draw[line width=3pt, black!25] (2.5,0) -- (4.3,0);
    \filldraw (0,0) circle (2pt);
    \filldraw[fill=white, draw=black, thick] (-2.5,0) circle (3pt);
    \filldraw[fill=white, draw=black, thick] (2.5,0) circle (3pt);
    \node[below] at (0,-0.15) {$a$};
    \node[below] at (-2.5,-0.15) {$a - b$};
    \node[below] at (2.5,-0.15) {$a + b$};
    \draw[decorate, decoration={brace, amplitude=5pt}] (-2.5,0.2) -- (0,0.2) node[midway, above=4pt, font=\small] {$b$};
    \draw[decorate, decoration={brace, amplitude=5pt}] (0,0.2) -- (2.5,0.2) node[midway, above=4pt, font=\small] {$b$};
    \node[colorofwisdom, font=\small] at (0,-0.95) {$|x - a| < b$};
    \node[gray, font=\small] at (-3.6,-0.95) {$|x - a| > b$};
    \node[gray, font=\small] at (3.6,-0.95) {$|x - a| > b$};
    \node[font=\small] at (0,1.25) {krajní body: $|x - a| = b$};
  \end{tikzpicture}
\end{center}

\begin{example}
  Vyřešme:
  \begin{enumerate}[label=(\alph*)]
    \item $|x - 3| = 5$. Čísla ve vzdálenosti $5$ od $3$ jsou $-2$ a $8$. $K = \{-2, 8\}$.
    \item $|x + 2|\leq 1$. Protože $x + 2 = x - (-2)$, hledáme čísla ve vzdálenosti nejvýše $1$ od $-2$. $K = [-3, -1]$.
    \item $|2x - 1| > 3$. Buď $2x - 1 > 3$, tedy $x > 2$, nebo $2x - 1 < -3$, tedy $x < -1$. $K = (-\infty, -1)\cup(2, \infty)$.
    \item $|x - 7| = -2$. Absolutní hodnota nemůže být záporná, $K = \emptyset$. Nic nepočítáme.
  \end{enumerate}
\end{example}

\subsection{Metoda nulových bodů}

Co když je v rovnici absolutních hodnot víc, nebo je kombinovaná s dalšími výrazy? Pak použijeme obecnou metodu. Myšlenka je jednoduchá: absolutní hodnota je `problém' jen proto, že nevíme, jestli je výraz uvnitř kladný, nebo záporný. Rozdělíme tedy osu na intervaly, na kterých to víme.

\begin{definitionbox}
  \textbf{Metoda nulových bodů pro absolutní hodnoty}:
  \begin{enumerate}
    \item Najdeme \textbf{nulové body} všech výrazů v absolutních hodnotách.
    \item Nulové body rozdělí reálnou osu na intervaly. Na každém intervalu má každý výraz stálé znaménko, takže absolutní hodnoty můžeme \textbf{odstranit} (výraz buď opíšeme, nebo vezmeme s opačným znaménkem).
    \item Na každém intervalu zvlášť vyřešíme vzniklou rovnici bez absolutních hodnot.
    \item \textbf{Ověříme}, zda nalezené řešení leží v daném intervalu. Pokud ne, zahodíme ho.
    \item Výsledek je sjednocení řešení ze všech intervalů.
  \end{enumerate}
\end{definitionbox}

Krok 4 je důležitý a zapomíná se na něj nejčastěji. Na každém intervalu řešíme jinou rovnici, která platí \textit{jen} na tom intervalu. Řešení mimo interval je řešením špatné rovnice.

\begin{example}\label{ex:abs-dve}
  Vyřešme $|x + 1| + |x - 2| = 5$. Nulové body jsou $-1$ a $2$. Sestavíme tabulku (do které bodu nulové body přiřadíme, je jedno, jen musíme být konzistentní):
  \begin{center}
    \renewcommand{\arraystretch}{1.4}
    \begin{tabular}{|c|c|c|c|}
      \hline
       & $(-\infty, -1)$ & $[-1, 2)$ & $[2, \infty)$ \\
      \hline
      $|x + 1|$ & $-x - 1$ & $x + 1$ & $x + 1$ \\
      $|x - 2|$ & $-x + 2$ & $-x + 2$ & $x - 2$ \\
      \hline
      rovnice & $-2x + 1 = 5$ & $3 = 5$ & $2x - 1 = 5$ \\
      řešení & $x = -2$ & --- & $x = 3$ \\
      leží v intervalu? & ano & --- & ano \\
      \hline
    \end{tabular}
  \end{center}
  $K = \{-2, 3\}$. Geometricky: hledáme body, jejichž součet vzdáleností od $-1$ a od $2$ je $5$. Body mezi $-1$ a $2$ mají tento součet vždy $3$, takže musíme jít ven. A vyjít musíme o $1$ na každou stranu.
\end{example}

Celý příklad si můžeme nakreslit. Graf funkce $y = |x + 1| + |x - 2|$ se skládá ze tří úseček, na každém intervalu z tabulky jedné. Nulové body jsou přesně místa, kde se graf `láme'. Uprostřed je graf vodorovný ve výšce $3$, po stranách stoupá. Řešení rovnice jsou $x$-ové souřadnice průsečíků s vodorovnou přímkou $y = 5$:
\begin{center}
  \begin{tikzpicture}[xscale=0.85, yscale=0.5]
    \draw[->] (-3.8,0) -- (4.8,0) node[right] {\small $x$};
    \draw[->] (0,-0.5) -- (0,7.8) node[above] {\small $y$};
    \foreach \x in {-3,...,4} \draw (\x,0.12) -- (\x,-0.12) node[below, font=\scriptsize] {$\x$};
    \draw[dashed, gray] (-1,0) -- (-1,3);
    \draw[dashed, gray] (2,0) -- (2,3);
    \draw[very thick, colorofwisdom] (-3,7) -- (-1,3) -- (2,3) -- (4,7);
    \draw[thick] (-3.6,5) -- (4.6,5) node[right, font=\small] {$y = 5$};
    \filldraw (-2,5) circle (3pt);
    \filldraw (3,5) circle (3pt);
    \draw[dotted, thick] (-2,5) -- (-2,0);
    \draw[dotted, thick] (3,5) -- (3,0);
    \node[left, font=\small] at (-3.05,7) {$-2x + 1$};
    \node[above, font=\small] at (0.5,3) {$3$};
    \node[right, font=\small] at (4.05,7) {$2x - 1$};
  \end{tikzpicture}
\end{center}
Z obrázku je hned vidět i to, že rovnice $|x + 1| + |x - 2| = c$ nemá pro $c < 3$ žádné řešení.

\begin{example}
  Vyřešme $|x - 1| = 2x + 4$. Nulový bod je $1$.
  \begin{itemize}
    \item Pro $x < 1$: $-x + 1 = 2x + 4$, tedy $x = -1$. To leží v $(-\infty, 1)$, takže je to řešení.
    \item Pro $x\geq 1$: $x - 1 = 2x + 4$, tedy $x = -5$. To \textbf{neleží} v $[1, \infty)$, zahodíme.
  \end{itemize}
  $K = \{-1\}$. Zkouška: $|-2| = 2 = 2\cdot(-1) + 4$. Kdybychom zapomněli na krok 4, tvrdili bychom, že $x = -5$ je řešení, ale $|-6| = 6\neq -6$.
\end{example}

Odkud se `řešení' $x = -5$ vzalo? Na intervalu $[1, \infty)$ jsme řešili rovnici $x - 1 = 2x + 4$. Tím jsme vlastně hledali průsečík přímky $y = 2x + 4$ s \textit{celou} přímkou $y = x - 1$, a ne jen s její částí napravo od $1$, která tvoří graf absolutní hodnoty. Průsečík vyšel na čárkované části, která do grafu nepatří:
\begin{center}
  \begin{tikzpicture}[scale=0.42]
    \draw[->] (-6.8,0) -- (4.5,0) node[right] {\small $x$};
    \draw[->] (0,-7.6) -- (0,6) node[above] {\small $y$};
    \foreach \x in {-6,...,3} \draw (\x,0.15) -- (\x,-0.15);
    \draw[dashed, gray, thick] (1,0) -- (-5.8,-6.8);
    \draw[very thick, colorofwisdom] (-4.2,5.2) -- (1,0) -- (4,3) node[right, font=\small] {$y = |x - 1|$};
    \draw[thick] (-5.6,-7.2) -- (0.6,5.2) node[right, font=\small] {$y = 2x + 4$};
    \filldraw (-1,2) circle (5pt);
    \draw[gray] (-1,2) -- (-2.6,1.2) node[left, font=\small, black] {$(-1, 2)$};
    \filldraw[fill=white, draw=black, thick] (-5,-6) circle (5pt) node[left=3pt, font=\small] {$(-5, -6)$};
    \node[gray, font=\scriptsize, right] at (-2.6,-3.9) {$y = x - 1$};
    \node[font=\small, below] at (1,-0.1) {$1$};
  \end{tikzpicture}
\end{center}

\begin{example}
  Nerovnice se řeší úplně stejně. Vyřešme $|x - 3| - |x| < 1$. Nulové body $0$ a $3$.
  \begin{itemize}
    \item Pro $x < 0$: $-(x - 3) - (-x) = 3 < 1$. Neplatí, $K_1 = \emptyset$.
    \item Pro $0\leq x < 3$: $-(x - 3) - x = 3 - 2x < 1$, tedy $x > 1$. Průnik s intervalem: $K_2 = (1, 3)$.
    \item Pro $x\geq 3$: $(x - 3) - x = -3 < 1$. Platí vždy, $K_3 = [3, \infty)$.
  \end{itemize}
  $K = K_1\cup K_2\cup K_3 = (1, \infty)$. U nerovnic neděláme `kontrolu, zda řešení leží v intervalu', ale rovnou děláme průnik s intervalem.
\end{example}

\begin{note}
  Někdy se vyplatí oběma stranám umocnit. Pokud jsou obě strany nezáporné, je to ekvivalentní úprava. Například
  \begin{equation*}
    |x - 1| > |x + 3|\quad\Leftrightarrow\quad{(x - 1)}^2 > {(x + 3)}^2\quad\Leftrightarrow\quad -2x + 1 > 6x + 9\quad\Leftrightarrow\quad x < -1.
  \end{equation*}
  Geometricky hledáme body, které jsou dál od $1$ než od $-3$. To jsou všechny body vlevo od středu, tedy od $-1$.
  \begin{center}
    \begin{tikzpicture}[scale=0.8]
      \draw[->, thick] (-5.5,0) -- (3.5,0) node[right] {$x$};
      \foreach \x in {-5,...,3} \draw (\x,0.08) -- (\x,-0.08) node[below, font=\scriptsize] {$\x$};
      \draw[line width=3pt, colorofwisdom] (-5.3,0) -- (-1,0);
      \filldraw[fill=white, draw=colorofwisdom, thick] (-1,0) circle (3pt);
      \filldraw (-3,0) circle (2.5pt);
      \filldraw (1,0) circle (2.5pt);
      \draw[dashed, gray] (-1,-0.7) -- (-1,1.3) node[above, font=\scriptsize] {střed};
      \draw[decorate, decoration={brace, amplitude=4pt}] (-3,0.25) -- (-1,0.25) node[midway, above=3pt, font=\scriptsize] {$2$};
      \draw[decorate, decoration={brace, amplitude=4pt}] (-1,0.25) -- (1,0.25) node[midway, above=3pt, font=\scriptsize] {$2$};
      \node[colorofwisdom, font=\small] at (-3.8,-0.9) {blíž k $-3$ než k $1$};
    \end{tikzpicture}
  \end{center}
\end{note}

\begin{problem}
  Vyřešte v $\mathbb{R}$:
  \begin{enumerate}[label=(\alph*)]
    \item $|3x + 2| = 7$,
    \item $|x - 4| < 2x + 1$,
    \item $|x| + |x - 4| = 4$,
    \item $|x + 2| - |x - 1| = x$,
    \item $|x^2 - 4| = 3x$.
  \end{enumerate}
\end{problem}

\begin{solution}
  \begin{enumerate}[label=(\alph*)]
    \item $3x + 2 = \pm 7$, tedy $x = \frac{5}{3}$ nebo $x = -3$.
    \item Pro $x < 4$: $4 - x < 2x + 1$, tedy $x > 1$, průnik $(1, 4)$. Pro $x\geq 4$: $x - 4 < 2x + 1$, tedy $x > -5$, průnik $[4, \infty)$. $K = (1, \infty)$.
    \item Pro $x < 0$: $-2x + 4 = 4$, tedy $x = 0$, neleží v intervalu. Pro $0\leq x < 4$: $4 = 4$, platí vždy. Pro $x\geq 4$: $2x - 4 = 4$, tedy $x = 4$, leží v intervalu. $K = [0, 4]$. Geometricky: součet vzdáleností od $0$ a $4$ je roven $4$ přesně pro body mezi nimi.
    \item Nulové body $-2$ a $1$. Pro $x < -2$: $-x - 2 + x - 1 = -3 = x$, tedy $x = -3$, leží. Pro $-2\leq x < 1$: $x + 2 + x - 1 = 2x + 1 = x$, tedy $x = -1$, leží. Pro $x\geq 1$: $x + 2 - x + 1 = 3 = x$, leží. $K = \{-3, -1, 3\}$.
    \item Pravá strana musí být nezáporná, tedy $x\geq 0$. Pro $0\leq x < 2$: $4 - x^2 = 3x$, $x^2 + 3x - 4 = 0$, $x\in\{1, -4\}$, v intervalu je jen $1$. Pro $x\geq 2$: $x^2 - 4 = 3x$, $x^2 - 3x - 4 = 0$, $x\in\{4, -1\}$, v intervalu je jen $4$. $K = \{1, 4\}$.
  \end{enumerate}
\end{solution}

\subsection{Absolutní hodnota s parametrem}

\begin{example}
  Určeme počet řešení rovnice $|x - 2| = a$ v závislosti na parametru $a\in\mathbb{R}$. Podle geometrického významu: pro $a < 0$ žádné, pro $a = 0$ jedno ($x = 2$), pro $a > 0$ dvě ($x = 2\pm a$).

  Zajímavější je rovnice $|x - 1| + |x + 1| = a$. Z příkladu~\ref{ex:abs-dve} už tušíme, že součet vzdáleností od $-1$ a $1$ je vždy alespoň $2$ (vzdálenost těch dvou bodů), a to přesně pro body mezi nimi. Tedy pro $a < 2$ žádné řešení, pro $a = 2$ celý interval $[-1, 1]$ a pro $a > 2$ dvě řešení $x = \pm\frac{a}{2}$. Kdo nevěří, ať si to rozepíše metodou nulových bodů.
\end{example}

Opět pomůže obrázek. Graf levé strany je `vana' s vodorovným dnem ve výšce $2$ a pravá strana je vodorovná přímka ve výšce $a$. Pokud je přímka pod dnem, nic neprotne. Pokud leží přesně na dně, splyne s ním na celém intervalu $[-1, 1]$. A pokud je výš, protne stěny vany ve dvou bodech:
\begin{center}
  \begin{tikzpicture}[xscale=0.9, yscale=0.55]
    \draw[->] (-3.2,0) -- (3.4,0) node[right] {\small $x$};
    \draw[->] (0,-0.5) -- (0,5.8) node[above] {\small $y$};
    \foreach \x in {-2,-1,1,2} \draw (\x,0.12) -- (\x,-0.12) node[below, font=\scriptsize] {$\x$};
    \draw[very thick] (-2.5,5) -- (-1,2) -- (1,2) -- (2.5,5);
    \draw[thick, colorofwisdom] (-3,3) -- (3,3) node[right, font=\small] {$a = 3$: dvě řešení};
    \filldraw[colorofwisdom] (-1.5,3) circle (3pt);
    \filldraw[colorofwisdom] (1.5,3) circle (3pt);
    \draw[line width=3pt, colorofwisdom!50] (-1,2) -- (1,2);
    \draw[thick, colorofwisdom] (-3,2) -- (3,2) node[right, font=\small] {$a = 2$: celý interval $[-1, 1]$};
    \draw[thick, gray] (-3,1) -- (3,1) node[right, font=\small, black] {$a = 1$: žádné řešení};
  \end{tikzpicture}
\end{center}

\section{Iracionální rovnice}

\begin{definitionbox}
  \textbf{Iracionální rovnice} je rovnice, ve které se neznámá vyskytuje pod odmocninou.

  \vspace{0.2cm}

  \textbf{Postup řešení}:
  \begin{enumerate}
    \item Určíme definiční obor (výrazy pod sudými odmocninami $\geq 0$).
    \item Odmocninu \textbf{osamostatníme} na jedné straně rovnice.
    \item Obě strany \textbf{umocníme}. To je \textbf{důsledková} úprava!
    \item Pokud v rovnici zbyly odmocniny, opakujeme kroky 2 a 3.
    \item Vyřešíme vzniklou rovnici.
    \item Provedeme \textbf{zkoušku}. Ta je zde \textbf{nutnou součástí řešení}, ne jen kontrolou.
  \end{enumerate}
\end{definitionbox}

Proč je zkouška nutná? V kapitole~\ref{kap:rovnice} jsme si řekli, že umocnění není ekvivalentní úprava. Z $a = b$ sice plyne $a^2 = b^2$, ale z $a^2 = b^2$ plyne jen $a = \pm b$. Umocněním tedy do rovnice `propašujeme' i řešení rovnice $a = -b$, která nás vůbec nezajímala. Takovým řešením se říká \textbf{falešné kořeny} a zkouška je jediný spolehlivý způsob, jak je odhalit.

\begin{example}
  Vyřešme $\sqrt{x + 3} = x + 1$. Definiční obor: $x\geq -3$. Umocníme:
  \begin{equation*}
    x + 3 = x^2 + 2x + 1\quad\Leftrightarrow\quad x^2 + x - 2 = 0\quad\Leftrightarrow\quad x\in\{1, -2\}.
  \end{equation*}
  Zkouška:
  \begin{itemize}
    \item $x = 1$: $L = \sqrt{4} = 2$, $P = 2$. Vyhovuje.
    \item $x = -2$: $L = \sqrt{1} = 1$, $P = -1$. \textbf{Nevyhovuje!}
  \end{itemize}
  $K = \{1\}$. Kořen $-2$ je falešný. Přišel z rovnice $\sqrt{x + 3} = -(x + 1)$, kterou jsme umocněním omylem vyřešili také.
\end{example}

Na obrázku je to krásně vidět. Rovnici $x + 3 = {(x + 1)}^2$, kterou jsme po umocnění řešili, splňují body přímky $y = x + 1$, pro které platí $y^2 = x + 3$. To je ale `ležatá' parabola, která má dvě poloviny: horní $y = \sqrt{x + 3}$ (graf naší levé strany) a dolní $y = -\sqrt{x + 3}$. Umocněním jsme k odmocnině nepozorovaně přibrali i její zrcadlový obraz. Falešný kořen $-2$ je průsečík s touto dolní, čárkovanou polovinou:
\begin{center}
  \begin{tikzpicture}[scale=0.75]
    \draw[->] (-3.8,0) -- (3.6,0) node[right] {\small $x$};
    \draw[->] (0,-2.9) -- (0,3.4) node[above] {\small $y$};
    \draw[very thick, colorofwisdom, domain=-3:3, samples=120] plot (\x, {sqrt(\x+3)}) node[right, font=\small] {$y = \sqrt{x + 3}$};
    \draw[dashed, thick, gray, domain=-3:3, samples=120] plot (\x, {-sqrt(\x+3)}) node[right, font=\small] {$y = -\sqrt{x + 3}$};
    \draw[thick] (-3.6,-2.6) -- (2.3,3.3) node[left, font=\small] {$y = x + 1$};
    \filldraw (1,2) circle (2.5pt) node[below right, font=\small] {$(1, 2)$};
    \draw[gray] (-2,-1) -- (-1.2,-2.3) node[below, font=\small, black] {$(-2, -1)$};
    \filldraw[fill=white, draw=black, thick] (-2,-1) circle (2.5pt);
  \end{tikzpicture}
\end{center}

\begin{note}
  Místo zkoušky můžeme na začátku přidat podmínku, že strana, na které není odmocnina, musí být nezáporná (protože odmocnina je nezáporná). V předchozím příkladu $x + 1\geq 0$, tedy $x\geq -1$, a kořen $-2$ by vypadl rovnou. Za těchto podmínek už je umocnění ekvivalentní úprava. U maturity ale doporučujeme dělat zkoušku tak jako tak. Podmínka se snadno pokazí, zkouška ne.
\end{note}

\begin{example}
  Vyřešme $\sqrt{2x + 3} - \sqrt{x + 1} = 1$. Definiční obor: $x\geq -1$. Osamostatníme jednu odmocninu a umocníme:
  \begin{align*}
    \sqrt{2x + 3} &= 1 + \sqrt{x + 1},\\
    2x + 3 &= 1 + 2\sqrt{x + 1} + x + 1,\\
    x + 1 &= 2\sqrt{x + 1}.
  \end{align*}
  Znovu umocníme: $x^2 + 2x + 1 = 4x + 4$, tedy $x^2 - 2x - 3 = 0$ a $x\in\{3, -1\}$. Zkouška: $x = 3$: $\sqrt{9} - \sqrt{4} = 1$, vyhovuje. $x = -1$: $\sqrt{1} - \sqrt{0} = 1$, vyhovuje. $K = \{-1, 3\}$.

  Tentokrát falešné kořeny nevznikly. Ze zkoušky jsme se to ale dozvěděli až na konci, takže ji rozhodně nebylo zbytečné dělat.
\end{example}

\begin{note}
  Ve třetím řádku předchozího příkladu by vás mohlo lákat vydělit rovnici výrazem $\sqrt{x + 1}$. Pak byste přišli o kořen $x = -1$, pro který je $\sqrt{x + 1} = 0$. Vytýkejte, nedělte. Pořád dokola.
\end{note}

\begin{example}
  Někdy stačí podívat se na definiční obor. Rovnice $\sqrt{x - 5} + \sqrt{2 - x} = 1$ vyžaduje $x\geq 5$ a zároveň $x\leq 2$. To nesplňuje žádné číslo, takže $K = \emptyset$ a nemusíme vůbec nic počítat. Kdo by začal umocňovat, dostal by po chvíli nějaká čísla, zkouškou by je pak pracně vyloučil a měl by pocit dobře vykonané práce. Zbytečně.
\end{example}

\begin{example}
  Substituce funguje i tady. Rovnice $x - 5\sqrt{x} + 6 = 0$ se substitucí $t = \sqrt{x}\geq 0$ změní na $t^2 - 5t + 6 = 0$, tedy $t\in\{2, 3\}$ a $x\in\{4, 9\}$. Zkouška: $4 - 10 + 6 = 0$ a $9 - 15 + 6 = 0$.
\end{example}

\begin{problem}
  Vyřešte v $\mathbb{R}$:
  \begin{enumerate}[label=(\alph*)]
    \item $\sqrt{x + 7} = x - 5$,
    \item $\sqrt{3x - 2} = \sqrt{x + 4}$,
    \item $\sqrt{x + 5} + \sqrt{x} = 5$,
    \item $x + \sqrt{x} = 12$,
    \item $\sqrt{x^2 + 9} = x - 3$.
  \end{enumerate}
\end{problem}

\begin{solution}
  \begin{enumerate}[label=(\alph*)]
    \item $x + 7 = x^2 - 10x + 25$, $x^2 - 11x + 18 = 0$, $x\in\{2, 9\}$. Zkouška: $x = 9$: $4 = 4$. $x = 2$: $3\neq -3$. $K = \{9\}$.
    \item $3x - 2 = x + 4$, $x = 3$. Zkouška: $\sqrt{7} = \sqrt{7}$. $K = \{3\}$.
    \item $\sqrt{x + 5} = 5 - \sqrt{x}$, umocníme: $x + 5 = 25 - 10\sqrt{x} + x$, tedy $10\sqrt{x} = 20$, $\sqrt{x} = 2$, $x = 4$. Zkouška: $3 + 2 = 5$. $K = \{4\}$.
    \item Substituce $t = \sqrt{x}$: $t^2 + t - 12 = 0$, $t\in\{3, -4\}$. Záporné $t$ zahodíme, $x = 9$. $K = \{9\}$.
    \item $x^2 + 9 = x^2 - 6x + 9$, tedy $x = 0$. Zkouška: $L = 3$, $P = -3$. $K = \emptyset$. Podmínka $x - 3\geq 0$ by to řekla rovnou.
  \end{enumerate}
\end{solution}

\section{Iracionální nerovnice}

U nerovnic je situace složitější, protože zkouška nefunguje. Nemůžeme dosadit celý interval. Musíme tedy umocňovat jen tehdy, když víme, že jsou obě strany nezáporné, a ostatní případy rozebrat zvlášť.

\begin{example}
  Vyřešme $\sqrt{x - 1} < 3$. Definiční obor $x\geq 1$. Obě strany jsou nezáporné, takže umocnění je ekvivalentní: $x - 1 < 9$, tedy $x < 10$. S definičním oborem $K = [1, 10)$.
\end{example}

\begin{example}
  Vyřešme $\sqrt{x + 2} > x$. Definiční obor $x\geq -2$. Pravá strana může mít libovolné znaménko, a tak rozebereme dva případy.
  \begin{itemize}
    \item $x\in[-2, 0)$: pravá strana je záporná a levá nezáporná, takže nerovnice platí vždy. $K_1 = [-2, 0)$.
    \item $x\geq 0$: obě strany jsou nezáporné a můžeme umocnit: $x + 2 > x^2$, tedy $x^2 - x - 2 < 0$, neboli $(x - 2)(x + 1) < 0$, a tedy $x\in(-1, 2)$. Průnik s $[0, \infty)$ je $K_2 = [0, 2)$.
  \end{itemize}
  $K = K_1\cup K_2 = [-2, 2)$.
\end{example}

\begin{center}
  \begin{tikzpicture}[scale=0.9]
    \draw[->] (-2.8,0) -- (3.5,0) node[right] {\small $x$};
    \draw[->] (0,-2.5) -- (0,3) node[above] {\small $y$};
    \draw[thick, colorofwisdom, domain=-2:3.2, samples=80] plot (\x, {sqrt(\x+2)}) node[right] {\small $y = \sqrt{x + 2}$};
    \draw[thick, domain=-2.4:2.9] plot (\x, \x) node[above right] {\small $y = x$};
    \filldraw (2,2) circle (2pt);
    \filldraw[colorofwisdom] (-2,0) circle (2pt);
    \draw[line width=3pt, colorofwisdom!50] (-2,-0.05) -- (2,-0.05);
    \node[below] at (2,0) {\small $2$};
    \node[below] at (-2,-0.1) {\small $-2$};
  \end{tikzpicture}
\end{center}

Na obrázku je vidět, že graf odmocniny leží nad přímkou $y = x$ přesně na intervalu $[-2, 2)$. Obrázek je u iracionálních nerovnic velice dobrá kontrola.

\begin{problem}
  Vyřešte v $\mathbb{R}$:
  \begin{enumerate}[label=(\alph*)]
    \item $\sqrt{2x + 1}\geq 3$,
    \item $\sqrt{x + 6} < x$,
    \item $\sqrt{4 - x^2} > -1$.
  \end{enumerate}
\end{problem}

\begin{solution}
  \begin{enumerate}[label=(\alph*)]
    \item $x\geq -\frac{1}{2}$, umocníme: $2x + 1\geq 9$, $x\geq 4$. $K = [4, \infty)$.
    \item Definiční obor $x\geq -6$. Pro $x < 0$ je pravá strana záporná a nerovnice neplatí. Pro $x\geq 0$ umocníme: $x + 6 < x^2$, $(x - 3)(x + 2) > 0$, tedy $x > 3$ (průnik s $x\geq 0$). $K = (3, \infty)$.
    \item Odmocnina je vždy nezáporná, takže je větší než $-1$ všude, kde je definovaná. $K = [-2, 2]$. Opět žádné počítání.
  \end{enumerate}
\end{solution}

\section{Odbočka --- Taxíková geometrie}

Absolutní hodnota měří vzdálenost na přímce. Jak měřit vzdálenost v rovině? Obvykle podle Pythagorovy věty: vzdálenost bodů $A = (a_1, a_2)$ a $B = (b_1, b_2)$ je
\begin{equation*}
  d(A, B) = \sqrt{{(a_1 - b_1)}^2 + {(a_2 - b_2)}^2}.
\end{equation*}
Představte si ale, že jste v městě s pravoúhlou sítí ulic, jako je Manhattan. Diagonálně přes domy projít nemůžete, takže skutečná vzdálenost, kterou ujdete, je
\begin{equation*}
  d_{\text{taxi}}(A, B) = |a_1 - b_1| + |a_2 - b_2|.
\end{equation*}
Tomu se říká \textbf{taxíková} (také manhattanská) \textbf{vzdálenost}.\footnote{V Praze to nefunguje, protože ulice nejsou pravoúhlé a taxikář vás stejně poveze oklikou.} Pro body $A = (0, 0)$ a $B = (4, 3)$ je taxíková vzdálenost $4 + 3 = 7$, zatímco vzdušnou čarou je to jen $5$. Zajímavé je, že \textit{každá} cesta po ulicích, která jde jen doprava a nahoru, je stejně dlouhá:
\begin{center}
  \begin{tikzpicture}[scale=0.8]
    \draw[gray!50] (-0.3,-0.3) grid (4.3,3.3);
    \draw[dashed, thick] (0,0) -- (4,3) node[pos=0.45, above left, font=\small] {$5$};
    \draw[very thick, colorofwisdom] (0,0) -- (4,0) -- (4,3);
    \draw[very thick, black!70] (0,0) -- (1,0) -- (1,1) -- (2,1) -- (2,2) -- (3,2) -- (3,3) -- (4,3);
    \filldraw (0,0) circle (2.5pt) node[below left] {$A$};
    \filldraw (4,3) circle (2.5pt) node[above right] {$B$};
    \node[right, font=\small, align=left] at (4.8,1.5) {obě plné cesty\\ mají délku $7$};
  \end{tikzpicture}
\end{center}
Taxíková vzdálenost splňuje všechny rozumné vlastnosti vzdálenosti, včetně trojúhelníkové nerovnosti (dokažte si to pomocí té naší). Geometrie se ale chová velice zvláštně. Co je třeba `kružnice' se středem v počátku a poloměrem $1$? Je to množina bodů, pro které $|x| + |y| = 1$. Metodou nulových bodů (pro každý kvadrant zvlášť) zjistíte, že to je čtverec postavený na špičku:
\begin{center}
  \begin{tikzpicture}[scale=1.3]
    \draw[->] (-1.6,0) -- (1.6,0) node[right] {\small $x$};
    \draw[->] (0,-1.4) -- (0,1.5) node[above] {\small $y$};
    \draw[thick, colorofwisdom] (1,0) -- (0,1) -- (-1,0) -- (0,-1) -- cycle;
    \draw[dashed] (0,0) circle (1);
    \node[below right] at (1,0) {\small $1$};
    \node[above right] at (0,1) {\small $1$};
    \node[colorofwisdom, anchor=south west] at (0.6,0.6) {\small $|x| + |y| = 1$};
  \end{tikzpicture}
\end{center}
Čárkovaně je pro srovnání obyčejná kružnice. Taxíková geometrie není jen hříčka. Používá se v robotice, při zpracování obrazu i ve strojovém učení, kde se mezi `vzdálenostmi' vybírá podle toho, která se pro danou úlohu hodí víc.

\section{Shrnutí a otázky k ústní zkoušce}

\begin{problem}
  Odpovězte vlastními slovy:
  \begin{enumerate}
    \item Jak je definována absolutní hodnota a jaký má geometrický význam?
    \item Jaké má absolutní hodnota vlastnosti? Dokažte trojúhelníkovou nerovnost.
    \item Jak se řeší rovnice a nerovnice tvaru $|x - a| = b$, $|x - a| < b$ a $|x - a| > b$?
    \item Popište metodu nulových bodů. Proč musíme kontrolovat, zda řešení leží v daném intervalu?
    \item Co je iracionální rovnice a jak se řeší? Proč je zkouška nutnou součástí řešení?
    \item Co jsou falešné kořeny a odkud se berou?
    \item Jak se řeší iracionální nerovnice a proč u nich nemůžeme spoléhat na zkoušku?
  \end{enumerate}
\end{problem}

\koneckapitoly

\chapter{Funkce a jejich vlastnosti}
\label{kap:funkce}
\mathquote{In mathematics you don't understand things. You just get used to them.}{John von Neumann}

Funkce jsou pravděpodobně nejdůležitějším pojmem celé matematiky. Popisují, jak jedna veličina závisí na jiné: dráha na čase, cena na množství, teplota na nadmořské výšce. V této kapitole si funkce pořádně zadefinujeme (pomocí relací z kapitoly 2, takže se nám konečně vyplatí, že jsme je zaváděli), popíšeme jejich vlastnosti a naučíme se s nimi pracovat. Konkrétní funkce (lineární, kvadratické, exponenciální, goniometrické, \ldots) přijdou v dalších kapitolách. Tahle kapitola je o tom, jak o funkcích mluvit.

\section{Pojem funkce}

Intuitivně je funkce předpis, který každému číslu $x$ přiřadí nějaké číslo $y$. Můžeme si ji představit jako automat: dáte do něj $x$ a vypadne jedno $y$. Důležité je slovo \textit{jedno}. Pokud by z automatu pro stejné $x$ jednou vypadlo jedno a podruhé jiné číslo, nebyl by to automat, ale loterie.

\begin{definitionbox}
  \textbf{Funkce} $f$ z množiny $A$ do množiny $B$ je relace $f\subseteq A\times B$ taková, že ke každému $x\in A$ existuje \textbf{právě jedno} $y\in B$ s vlastností $(x, y)\in f$. Toto $y$ značíme $f(x)$ a nazýváme ho \textbf{funkční hodnota} v bodě $x$. Píšeme
  \begin{equation*}
    f: A\to B,\qquad x\mapsto f(x),
  \end{equation*}
  nebo zkráceně $y = f(x)$. Proměnnou $x$ nazýváme \textbf{nezávisle proměnná} (argument), proměnnou $y$ \textbf{závisle proměnná}.

  \vspace{0.2cm}

  \textbf{Definiční obor} $D(f)$ je množina všech $x$, pro která je funkce definovaná. \textbf{Obor hodnot} $H(f)$ je množina všech hodnot, kterých funkce nabývá:
  \begin{equation*}
    H(f) = \{f(x)\ |\ x\in D(f)\}.
  \end{equation*}
  \textbf{Graf funkce} je množina bodů $\{(x, f(x))\ |\ x\in D(f)\}$ v rovině. Z definice je graf funkce vlastně funkce sama.
\end{definitionbox}

V tomto textu budeme pracovat s \textbf{reálnými funkcemi reálné proměnné}, tedy $A, B\subseteq\mathbb{R}$. Pokud je funkce zadaná jen předpisem (například $f(x) = \sqrt{x - 1}$) a definiční obor není výslovně uveden, rozumíme jím největší množinu, pro kterou má předpis smysl. Určujeme ho stejně jako definiční obor výrazu v kapitole~\ref{kap:vyrazy}.

\begin{example}
  Několik funkcí a ne-funkcí:
  \begin{itemize}
    \item $f(x) = x^2$ je funkce, $D(f) = \mathbb{R}$, $H(f) = [0, \infty)$. Různým $x$ může příslušet stejné $y$ (například $f(2) = f(-2) = 4$), to definice dovoluje.
    \item $g(x) = \frac{1}{x - 2}$ je funkce s $D(g) = \mathbb{R}\setminus\{2\}$ a $H(g) = \mathbb{R}\setminus\{0\}$.
    \item Relace $\{(x, y)\in\mathbb{R}^2\ |\ x^2 + y^2 = 1\}$ (kružnice) \textbf{není} funkce, protože například k $x = 0$ patří dvě hodnoty $y = 1$ i $y = -1$.
    \item Relace $\{(1, 2), (2, 3), (3, 2)\}$ je funkce s $D = \{1, 2, 3\}$ a $H = \{2, 3\}$. Funkce nemusí mít předpis, stačí tabulka.
  \end{itemize}
\end{example}

\begin{note}
  Jak poznat z obrázku, zda je křivka grafem funkce? \textbf{Test svislou přímkou}: každá svislá přímka smí křivku protnout nejvýše jednou. Pokud ji někde protne dvakrát, máme pro jedno $x$ dvě hodnoty $y$. Kružnice testem neprojde.
\end{note}

\begin{definitionbox}
  \textbf{Způsoby zadání funkce}:
  \begin{itemize}
    \item \textbf{předpisem} (rovnicí), například $f(x) = 3x^2 - 1$,
    \item \textbf{tabulkou} (u funkcí s konečným definičním oborem nebo u naměřených dat),
    \item \textbf{grafem},
    \item \textbf{slovně}, například `každému přirozenému číslu přiřadíme jeho nejmenšího prvočíselného dělitele'.
  \end{itemize}
  Funkce může být zadána i \textbf{po částech}, tedy různými předpisy na různých intervalech. Takovou funkcí je například absolutní hodnota z kapitoly~\ref{kap:absolutni}.
\end{definitionbox}

\begin{problem}
  Určete definiční obor:
  \begin{enumerate}[label=(\alph*)]
    \item $f(x) = \dfrac{x + 1}{x^2 - 5x + 6}$,
    \item $g(x) = \sqrt{4 - x^2}$,
    \item $h(x) = \dfrac{\sqrt{x + 3}}{|x| - 1}$,
    \item $k(x) = \sqrt{\dfrac{x - 1}{x + 2}}$.
  \end{enumerate}
\end{problem}

\begin{solution}
  \begin{enumerate}[label=(\alph*)]
    \item $x^2 - 5x + 6 = (x - 2)(x - 3)\neq 0$, tedy $D(f) = \mathbb{R}\setminus\{2, 3\}$.
    \item $4 - x^2\geq 0$, tedy $D(g) = [-2, 2]$.
    \item $x\geq -3$ a $|x|\neq 1$, tedy $D(h) = [-3, -1)\cup(-1, 1)\cup(1, \infty)$.
    \item $\frac{x - 1}{x + 2}\geq 0$, podle příkladu v kapitole~\ref{kap:rovnice} je $D(k) = (-\infty, -2)\cup[1, \infty)$.
  \end{enumerate}
\end{solution}

\section{Vlastnosti funkcí}

\subsection{Monotonie}

\begin{definitionbox}
  Nechť $M\subseteq D(f)$. Funkce $f$ je na množině $M$
  \begin{itemize}
    \item \textbf{rostoucí}, pokud pro všechna $x_1, x_2\in M$ platí $x_1 < x_2\Rightarrow f(x_1) < f(x_2)$,
    \item \textbf{klesající}, pokud pro všechna $x_1, x_2\in M$ platí $x_1 < x_2\Rightarrow f(x_1) > f(x_2)$,
    \item \textbf{neklesající}, pokud pro všechna $x_1, x_2\in M$ platí $x_1 < x_2\Rightarrow f(x_1)\leq f(x_2)$,
    \item \textbf{nerostoucí}, pokud pro všechna $x_1, x_2\in M$ platí $x_1 < x_2\Rightarrow f(x_1)\geq f(x_2)$.
  \end{itemize}
  Funkce, které mají některou z těchto vlastností, nazýváme \textbf{monotónní}, rostoucí a klesající funkce \textbf{ryze monotónní}. Pokud množinu $M$ neuvedeme, myslíme celý definiční obor.
\end{definitionbox}

Slovy: funkce je rostoucí, když s rostoucím $x$ roste i $y$. Graf rostoucí funkce `jde do kopce', když ho čteme zleva doprava. V definici si všimněte kvantifikátoru `pro všechna'. Nestačí ověřit pár bodů, musí to platit pro každou dvojici.

\begin{example}
  Dokažme, že $f(x) = 2x + 1$ je rostoucí na $\mathbb{R}$. Nechť $x_1 < x_2$. Pak $2x_1 < 2x_2$ (násobíme kladným číslem) a $2x_1 + 1 < 2x_2 + 1$, tedy $f(x_1) < f(x_2)$.

  Dokažme, že $g(x) = \frac{1}{x}$ je klesající na $(0, \infty)$. Nechť $0 < x_1 < x_2$. Pak
  \begin{equation*}
    g(x_1) - g(x_2) = \frac{1}{x_1} - \frac{1}{x_2} = \frac{x_2 - x_1}{x_1x_2} > 0,
  \end{equation*}
  protože čitatel i jmenovatel jsou kladné. Tedy $g(x_1) > g(x_2)$. Úplně stejně se ukáže, že $g$ je klesající i na $(-\infty, 0)$.
\end{example}

\begin{note}\label{note:1x}
  Pozor, funkce $\frac{1}{x}$ \textbf{není} klesající na celém svém definičním oboru $\mathbb{R}\setminus\{0\}$! Stačí vzít $x_1 = -1 < 1 = x_2$, a $g(-1) = -1 < 1 = g(1)$. Je klesající na každém z intervalů $(-\infty, 0)$ a $(0, \infty)$ zvlášť, ale ne na jejich sjednocení. Tohle je oblíbená chyták otázka u ústní maturity. Správná odpověď je `klesající na $(-\infty, 0)$ a na $(0, \infty)$', nikdy `klesající na $\mathbb{R}\setminus\{0\}$'.
\end{note}

\subsection{Prostá funkce}

\begin{definitionbox}
  Funkce $f$ je \textbf{prostá}, pokud pro všechna $x_1, x_2\in D(f)$ platí
  \begin{equation*}
    x_1\neq x_2\quad\Rightarrow\quad f(x_1)\neq f(x_2).
  \end{equation*}
  Ekvivalentně (obměnou implikace z kapitoly 1): $f(x_1) = f(x_2)\Rightarrow x_1 = x_2$.
\end{definitionbox}

Prostá funkce nikdy nepřiřadí stejnou hodnotu dvěma různým číslům. Graficky: každá \textbf{vodorovná} přímka protne graf nejvýše jednou. Funkce $x^2$ prostá není ($f(-2) = f(2)$), funkce $x^3$ ano.

\begin{definitionbox}
  \textbf{Věta}: Každá ryze monotónní funkce je prostá.
\end{definitionbox}

\begin{proof}
  Nechť $f$ je rostoucí a $x_1\neq x_2$. Bez újmy na obecnosti\footnote{Tato fráze se v matematice používá, když by druhý případ vypadal úplně stejně, jen s prohozenými písmeny. Je to způsob, jak elegantně říct `ten druhý případ se mi nechce psát'.} $x_1 < x_2$. Pak $f(x_1) < f(x_2)$, a tedy $f(x_1)\neq f(x_2)$. Pro klesající funkci stejně.
\end{proof}

Obrácená implikace neplatí! Funkce $\frac{1}{x}$ je prostá (zkuste si to dokázat), ale podle poznámky~\ref{note:1x} na svém definičním oboru monotónní není.

\subsection{Sudost a lichost}

\begin{definitionbox}
  Funkce $f$ je
  \begin{itemize}
    \item \textbf{sudá}, pokud pro každé $x\in D(f)$ je $-x\in D(f)$ a $f(-x) = f(x)$,
    \item \textbf{lichá}, pokud pro každé $x\in D(f)$ je $-x\in D(f)$ a $f(-x) = -f(x)$.
  \end{itemize}
  Graf sudé funkce je \textbf{souměrný podle osy $y$}, graf liché funkce je \textbf{souměrný podle počátku}.
\end{definitionbox}

Názvy pochází z mocnin: $x^2, x^4, x^6, \ldots$ jsou sudé funkce a $x, x^3, x^5, \ldots$ liché. Nenechte se ale zmást. Většina funkcí není ani sudá, ani lichá. Funkce není jako číslo, které musí být jedno nebo druhé.

\begin{center}
  \begin{tikzpicture}[scale=0.7]
    \begin{scope}
      \draw[->] (-2.6,0) -- (2.8,0) node[right] {\small $x$};
      \draw[->] (0,-0.6) -- (0,3.4) node[above] {\small $y$};
      \draw[thick, colorofwisdom, domain=-1.8:1.8, samples=60] plot (\x, {\x*\x});
      \draw[dashed] (-1.5,0) -- (-1.5,2.25) -- (1.5,2.25) -- (1.5,0);
      \node[below] at (-1.5,0) {\small $-x$};
      \node[below] at (1.5,0) {\small $x$};
      \node at (0,-1.3) {\small sudá: $f(x) = x^2$};
    \end{scope}
    \begin{scope}[xshift=8cm, yshift=1.4cm]
      \draw[->] (-2.6,0) -- (2.8,0) node[right] {\small $x$};
      \draw[->] (0,-2) -- (0,2) node[above] {\small $y$};
      \draw[thick, colorofwisdom, domain=-1.55:1.55, samples=60] plot (\x, {0.5*\x*\x*\x});
      \draw[dashed] (-1.4,0) -- (-1.4,-1.372) (1.4,0) -- (1.4,1.372);
      \draw[dashed, gray] (-1.4,-1.372) -- (1.4,1.372);
      \node[above] at (-1.4,0) {\small $-x$};
      \node[below] at (1.4,0) {\small $x$};
      \node at (0,-2.7) {\small lichá: $f(x) = \frac{1}{2}x^3$};
    \end{scope}
  \end{tikzpicture}
\end{center}

\begin{example}
  Rozhodněme o sudosti a lichosti:
  \begin{itemize}
    \item $f(x) = x^4 - x^2$: $f(-x) = {(-x)}^4 - {(-x)}^2 = x^4 - x^2 = f(x)$. Sudá.
    \item $g(x) = x^3 + x$: $g(-x) = -x^3 - x = -g(x)$. Lichá.
    \item $h(x) = x^2 + x$: $h(1) = 2$, $h(-1) = 0$. Protože $h(-1)\neq h(1)$ a $h(-1)\neq -h(1)$, není ani sudá, ani lichá. Pro vyvrácení stačí jeden protipříklad.
    \item $k(x) = \sqrt{x}$: $D(k) = [0, \infty)$ není souměrný podle nuly ($1\in D$, ale $-1\notin D$). Není ani sudá, ani lichá a nemusíme nic počítat.
  \end{itemize}
\end{example}

\begin{problem}
  \begin{enumerate}[label=(\alph*)]
    \item Rozhodněte o sudosti a lichosti funkcí $\frac{x}{x^2 + 1}$, $|x| + 1$, $\frac{1}{x - 1}$.
    \item Existuje funkce, která je zároveň sudá i lichá?
    \item Dokažte, že součin dvou lichých funkcí je sudá funkce.
  \end{enumerate}
\end{problem}

\begin{solution}
  \begin{enumerate}[label=(\alph*)]
    \item První je lichá, druhá sudá. Třetí má definiční obor $\mathbb{R}\setminus\{1\}$, který není souměrný, takže není ani sudá, ani lichá.
    \item Ano, $f(x) = 0$. Z $f(-x) = f(x)$ a $f(-x) = -f(x)$ plyne $f(x) = -f(x)$, tedy $f(x) = 0$. Je to tedy jediná taková funkce (na daném souměrném definičním oboru).
    \item Nechť $f, g$ jsou liché a $h = f\cdot g$. Pak $h(-x) = f(-x)g(-x) = (-f(x))(-g(x)) = f(x)g(x) = h(x)$.
  \end{enumerate}
\end{solution}

\subsection{Omezenost a extrémy}

\begin{definitionbox}
  Funkce $f$ je na množině $M\subseteq D(f)$
  \begin{itemize}
    \item \textbf{shora omezená}, pokud existuje $h\in\mathbb{R}$ takové, že $f(x)\leq h$ pro všechna $x\in M$,
    \item \textbf{zdola omezená}, pokud existuje $d\in\mathbb{R}$ takové, že $f(x)\geq d$ pro všechna $x\in M$,
    \item \textbf{omezená}, pokud je omezená shora i zdola.
  \end{itemize}
  Funkce $f$ má v bodě $a\in M$ \textbf{maximum} na $M$, pokud $f(x)\leq f(a)$ pro všechna $x\in M$, a \textbf{minimum}, pokud $f(x)\geq f(a)$ pro všechna $x\in M$. Maximum a minimum společně nazýváme \textbf{extrémy}.
\end{definitionbox}

\begin{example}
  Funkce $f(x) = x^2$ je zdola omezená (například číslem $0$), shora omezená není. Má minimum v bodě $0$, maximum nemá.

  Funkce $g(x) = \frac{1}{x}$ na intervalu $(0, \infty)$ je zdola omezená číslem $0$, ale minimum nemá! Hodnoty se k nule libovolně blíží, ale nikdy jí nedosáhnou. Omezenost tedy neznamená, že existuje extrém. V analýze se pro `nejlepší' dolní mez zavádí pojem infimum, ale to už je nad rámec maturity.

  Funkce $\sin x$ je omezená ($-1\leq\sin x\leq 1$) a maxima nabývá v nekonečně mnoha bodech.
\end{example}

\subsection{Periodičnost}

\begin{definitionbox}
  Funkce $f$ je \textbf{periodická} s periodou $p > 0$, pokud pro každé $x\in D(f)$ je i $x\pm p\in D(f)$ a
  \begin{equation*}
    f(x + p) = f(x).
  \end{equation*}
  Nejmenší takové $p$ (pokud existuje) nazýváme \textbf{základní perioda}.
\end{definitionbox}

Typickými periodickými funkcemi jsou goniometrické funkce, kterým patří kapitola~\ref{kap:goniometrie}. Funkce $\sin x$ má základní periodu $2\pi$. Pokud je $p$ perioda, je periodou i $2p$, $3p$ a tak dál, proto mluvíme o základní periodě. Zajímavost: konstantní funkce je periodická s libovolnou periodou, a proto základní periodu nemá.

\section{Operace s funkcemi}

\subsection{Skládání funkcí}

\begin{definitionbox}
  \textbf{Složená funkce} $g\circ f$ (čteme `$g$ po $f$') je funkce
  \begin{equation*}
    (g\circ f)(x) = g(f(x)).
  \end{equation*}
  Nejdřív se tedy použije $f$ (\textbf{vnitřní funkce}) a na výsledek $g$ (\textbf{vnější funkce}). Definiční obor tvoří ta $x\in D(f)$, pro která je $f(x)\in D(g)$.
\end{definitionbox}

\begin{example}
  Nechť $f(x) = x + 1$ a $g(x) = x^2$. Pak
  \begin{equation*}
    (g\circ f)(x) = g(x + 1) = {(x + 1)}^2,\qquad (f\circ g)(x) = f(x^2) = x^2 + 1.
  \end{equation*}
  Tedy $g\circ f\neq f\circ g$. Skládání funkcí \textbf{není komutativní}, stejně jako násobení matic.\footnote{To není náhoda. Matice jsou ve skutečnosti zápisem určitého druhu funkcí (lineárních zobrazení) a násobení matic odpovídá jejich skládání. To je důvod, proč je násobení matic definované tak divně.} Obléct si ponožky a pak boty dá jiný výsledek než obráceně.
\end{example}

\begin{example}
  Nechť $f(x) = \sqrt{x}$ a $g(x) = x - 4$. Pak $(f\circ g)(x) = \sqrt{x - 4}$ s definičním oborem $[4, \infty)$, zatímco $(g\circ f)(x) = \sqrt{x} - 4$ s definičním oborem $[0, \infty)$.
\end{example}

\begin{problem}
  \begin{enumerate}[label=(\alph*)]
    \item Pro $f(x) = 2x - 1$ a $g(x) = \frac{1}{x}$ určete $f\circ g$, $g\circ f$ a $f\circ f$ včetně definičních oborů.
    \item Funkci $h(x) = \sqrt{x^2 + 1}$ zapište jako složení dvou jednodušších funkcí.
  \end{enumerate}
\end{problem}

\begin{solution}
  \begin{enumerate}[label=(\alph*)]
    \item $(f\circ g)(x) = \frac{2}{x} - 1$, $D = \mathbb{R}\setminus\{0\}$. $(g\circ f)(x) = \frac{1}{2x - 1}$, $D = \mathbb{R}\setminus\left\{\frac{1}{2}\right\}$. $(f\circ f)(x) = 2(2x - 1) - 1 = 4x - 3$, $D = \mathbb{R}$.
    \item $h = g\circ f$, kde $f(x) = x^2 + 1$ a $g(x) = \sqrt{x}$. Tento rozklad budeme potřebovat u derivací.
  \end{enumerate}
\end{solution}

\subsection{Inverzní funkce}

Funkce přiřazuje číslu $x$ číslo $y$. Můžeme jít obráceně a z $y$ zjistit, z jakého $x$ vzniklo? To jde jen tehdy, když každé $y$ vzniklo z jediného $x$, tedy když je funkce prostá.

\begin{definitionbox}
  Nechť $f$ je \textbf{prostá} funkce. \textbf{Inverzní funkce} $f^{-1}$ je funkce, která každému $y\in H(f)$ přiřadí to (jediné) $x\in D(f)$, pro které $f(x) = y$. Platí:
  \begin{itemize}
    \item $D(f^{-1}) = H(f)$ a $H(f^{-1}) = D(f)$,
    \item $f^{-1}(f(x)) = x$ pro všechna $x\in D(f)$ a $f(f^{-1}(y)) = y$ pro všechna $y\in H(f)$,
    \item graf $f^{-1}$ je souměrný s grafem $f$ podle přímky $y = x$.
  \end{itemize}

  \textbf{Postup hledání}: V rovnici $y = f(x)$ vyjádříme $x$ pomocí $y$ a pak písmena $x$ a $y$ prohodíme (protože nezávislou proměnnou obvykle značíme $x$).
\end{definitionbox}

\begin{note}
  Značení $f^{-1}$ je bohužel matoucí. Neznamená $\frac{1}{f}$! Například pro $f(x) = 2x$ je $f^{-1}(x) = \frac{x}{2}$, ale $\frac{1}{f(x)} = \frac{1}{2x}$. Exponent $-1$ zde znamená `inverze vzhledem ke skládání', podobně jako $A^{-1}$ u matic znamenalo inverzi vzhledem k násobení matic.
\end{note}

\begin{example}
  Najděme inverzní funkci k $f(x) = 2x - 3$. Funkce je rostoucí, tedy prostá. Z $y = 2x - 3$ vyjádříme $x = \frac{y + 3}{2}$ a prohodíme písmena: $f^{-1}(x) = \frac{x + 3}{2}$. Ověření: $f^{-1}(f(x)) = \frac{2x - 3 + 3}{2} = x$.
\end{example}

\begin{example}
  Funkce $f(x) = x^2$ na $\mathbb{R}$ prostá není, takže inverzní funkci nemá. Pokud ale definiční obor omezíme na $[0, \infty)$, prostá je a inverzní funkce je $f^{-1}(x) = \sqrt{x}$. Přesně takhle je ostatně definována odmocnina v kapitole~\ref{kap:vyrazy}: je to inverze k mocnině na nezáporných číslech.
  \begin{center}
    \begin{tikzpicture}[scale=0.85]
      \draw[->] (-0.5,0) -- (3.6,0) node[right] {\small $x$};
      \draw[->] (0,-0.5) -- (0,3.6) node[above] {\small $y$};
      \draw[dashed, gray] (-0.3,-0.3) -- (3.4,3.4) node[above right] {\small $y = x$};
      \draw[thick, colorofwisdom, domain=0:1.85, samples=60] plot (\x, {\x*\x}) node[above] {\small $y = x^2$};
      \draw[thick, domain=0:3.4, samples=80] plot (\x, {sqrt(\x)}) node[right] {\small $y = \sqrt{x}$};
    \end{tikzpicture}
  \end{center}
\end{example}

\begin{problem}
  Najděte inverzní funkci a určete její definiční obor:
  \begin{enumerate}[label=(\alph*)]
    \item $f(x) = 5 - 3x$,
    \item $g(x) = \dfrac{2x + 1}{x - 3}$,
    \item $h(x) = x^2 + 1$ pro $x\leq 0$.
  \end{enumerate}
\end{problem}

\begin{solution}
  \begin{enumerate}[label=(\alph*)]
    \item $x = \frac{5 - y}{3}$, tedy $f^{-1}(x) = \frac{5 - x}{3}$, $D = \mathbb{R}$.
    \item $y(x - 3) = 2x + 1$, tedy $x(y - 2) = 3y + 1$ a $x = \frac{3y + 1}{y - 2}$. $g^{-1}(x) = \frac{3x + 1}{x - 2}$, $D(g^{-1}) = H(g) = \mathbb{R}\setminus\{2\}$.
    \item $x^2 = y - 1$ a protože $x\leq 0$, je $x = -\sqrt{y - 1}$. $h^{-1}(x) = -\sqrt{x - 1}$, $D(h^{-1}) = H(h) = [1, \infty)$. Pozor na znaménko: kdybychom napsali $+\sqrt{\ }$, dostali bychom inverzi k jiné větvi paraboly.
  \end{enumerate}
\end{solution}

\section{Transformace grafů}

Pokud známe graf funkce $f$, můžeme z něj snadno nakreslit grafy mnoha příbuzných funkcí. To je jedna z nejužitečnějších dovedností pro celou analýzu, protože místo bodování (počítání hodnot v mnoha bodech) stačí znát pár základních grafů a vědět, jak se posouvají.

\begin{definitionbox}
  \textbf{Transformace grafu} funkce $y = f(x)$ (pro $c > 0$, $a > 0$):
  \begin{center}
    \renewcommand{\arraystretch}{1.4}
    \begin{tabular}{|l|l|}
      \hline
      Nová funkce & Co se stane s grafem \\
      \hline
      $y = f(x) + c$ & posun o $c$ \textbf{nahoru} \\
      $y = f(x) - c$ & posun o $c$ \textbf{dolů} \\
      $y = f(x - c)$ & posun o $c$ \textbf{doprava} \\
      $y = f(x + c)$ & posun o $c$ \textbf{doleva} \\
      $y = -f(x)$ & překlopení podle osy $x$ \\
      $y = f(-x)$ & překlopení podle osy $y$ \\
      $y = a\cdot f(x)$ & $a$-násobné roztažení ve směru osy $y$ (pro $a < 1$ stlačení) \\
      $y = f(a\cdot x)$ & $a$-násobné stlačení ve směru osy $x$ (pro $a < 1$ roztažení) \\
      $y = |f(x)|$ & části grafu pod osou $x$ se překlopí nahoru \\
      $y = f(|x|)$ & pravá část grafu se zrcadlí doleva (levá se zahodí) \\
      \hline
    \end{tabular}
  \end{center}
\end{definitionbox}

Zdaleka nejčastější chyba je u vodorovných posunů. $f(x - 2)$ posouvá graf \textbf{doprava}, ne doleva. Proč? Funkce $g(x) = f(x - 2)$ nabývá v bodě $x = 2$ hodnoty $f(0)$. Co dřív bylo v nule, je teď ve dvojce. Graf se tedy posunul doprava. Vodorovné transformace prostě dělají `opak toho, co vidíte', a je lepší to pochopit než to biflovat.

\begin{example}
  Nakresleme graf $y = {(x - 2)}^2 - 1$. Vyjdeme z paraboly $y = x^2$, posuneme ji o $2$ doprava a o $1$ dolů. Vrchol se tak přesune z $(0, 0)$ do $(2, -1)$.
  \begin{center}
    \begin{tikzpicture}[scale=0.75]
      \draw[->] (-2.5,0) -- (4.8,0) node[right] {\small $x$};
      \draw[->] (0,-1.6) -- (0,4.4) node[above] {\small $y$};
      \foreach \x in {-2,-1,1,2,3,4} \draw (\x,0.08) -- (\x,-0.08) node[below] {\tiny $\x$};
      \draw[thick, gray, dashed, domain=-2:2, samples=60] plot (\x, {\x*\x}) node[left] {\small $y = x^2$};
      \draw[thick, colorofwisdom, domain=-0.2:4.2, samples=60] plot (\x, {(\x-2)*(\x-2)-1}) node[right] {\small $y = (x - 2)^2 - 1$};
      \filldraw[colorofwisdom] (2,-1) circle (2pt);
      \draw[->, gray] (0.2,0.1) -- (1.8,-0.85);
    \end{tikzpicture}
  \end{center}
\end{example}

\begin{example}
  Graf $y = |x^2 - 4|$ dostaneme z paraboly $y = x^2 - 4$ tak, že část mezi $-2$ a $2$, která je pod osou $x$, překlopíme nahoru. Graf $y = {|x|}^2 - 4$ je naopak stejný jako $y = x^2 - 4$, protože parabola je sudá a pravá strana je zrcadlem levé už teď.
\end{example}

\begin{problem}
  Popište, jak vzniknou z grafu $y = f(x)$ grafy:
  \begin{enumerate}[label=(\alph*)]
    \item $y = f(x + 3) - 2$,
    \item $y = -f(x) + 1$,
    \item $y = 2f(x - 1)$,
    \item $y = |f(x)| - 1$.
  \end{enumerate}
  Výsledky si nakreslete pro $f(x) = x^2$ a $f(x) = \frac{1}{x}$.
\end{problem}

\begin{solution}
  \begin{enumerate}[label=(\alph*)]
    \item Posun o $3$ doleva a o $2$ dolů.
    \item Překlopení podle osy $x$ a pak posun o $1$ nahoru. Na pořadí záleží! Kdybychom nejdřív posunuli a pak překlopili, dostali bychom $-(f(x) + 1) = -f(x) - 1$.
    \item Posun o $1$ doprava a dvojnásobné roztažení ve směru osy $y$.
    \item Části pod osou $x$ se překlopí nahoru a pak se celý graf posune o $1$ dolů.
  \end{enumerate}
\end{solution}

\section{Odbočka --- Monstra}

Když se v 19. století matematici začali vážně zabývat tím, co je funkce, zjistili, že definice `ke každému $x$ právě jedno $y$' dovoluje věci, které nikdo nečekal.

Peter Gustav Lejeune Dirichlet přišel s funkcí
\begin{equation*}
  D(x) = \begin{cases}
    1 & \text{pro } x\in\mathbb{Q},\\
    0 & \text{pro } x\notin\mathbb{Q}.
  \end{cases}
\end{equation*}
Je to naprosto legitimní funkce, splňuje definici. Jenže se nedá nakreslit. Racionální i iracionální čísla jsou na reálné ose `všude' (mezi každými dvěma čísly leží racionální i iracionální číslo), takže graf by vypadal jako dvě husté vodorovné čáry v $y = 0$ a $y = 1$, ze kterých ale ve skutečnosti chybí nekonečně mnoho bodů. Tato funkce není spojitá v žádném bodě.

Karl Weierstrass pak v roce 1872 sestrojil funkci, která je spojitá všude (její graf se dá nakreslit jedním tahem), ale nemá tečnu v žádném bodě. Graf je tak `zubatý', že zuby má v každém sebemenším kousku. Matematici té doby byli zděšení. Henri Poincaré, autor citátu na začátku kapitoly~\ref{kap:kvadraticke}, takové funkce nazval \textbf{monstry}. Charles Hermite napsal kolegovi:
\begin{center}
  \textit{`S hrůzou a odporem se odvracím od této politováníhodné rány spojitých funkcí, které nemají derivaci.'}
\end{center}
Dnes víme, že monstra nejsou výjimka, ale pravidlo. V jistém přesném smyslu je `skoro každá' spojitá funkce taková jako ta Weierstrassova. Hezké hladké funkce, se kterými pracujeme ve škole, jsou ve skutečnosti ty vzácné. A mimochodem, podobně zubaté křivky dnes používáme k popisu pobřeží, horských hřebenů nebo cen akcií. Monstra se ukázala být docela užitečná.

\section{Shrnutí a otázky k ústní zkoušce}

\begin{problem}
  Odpovězte vlastními slovy:
  \begin{enumerate}
    \item Co je funkce? Jak souvisí s relacemi? Co je definiční obor, obor hodnot a graf?
    \item Jak lze funkci zadat? Jak poznáte z obrázku, zda jde o graf funkce?
    \item Definujte rostoucí, klesající, nerostoucí a neklesající funkci. Dokažte, že $\frac{1}{x}$ je klesající na $(0, \infty)$. Proč není klesající na $\mathbb{R}\setminus\{0\}$?
    \item Co je prostá funkce? Jaký je vztah mezi prostou a ryze monotónní funkcí?
    \item Co je sudá a lichá funkce? Jak to poznáte z grafu?
    \item Co je omezená funkce? Co je maximum a minimum? Může být funkce omezená a nemít extrém?
    \item Co je periodická funkce?
    \item Co je složená funkce? Je skládání funkcí komutativní?
    \item Kdy existuje inverzní funkce? Jak ji najdeme a jak vypadá její graf?
    \item Jak se změní graf funkce, když k ní přičteme konstantu, když ji posuneme v argumentu, když ji vynásobíme $-1$ a když na ni použijeme absolutní hodnotu?
  \end{enumerate}
\end{problem}

\koneckapitoly

\chapter{Elementární funkce}
\label{kap:elementarni}
\mathquote{Philosophy is written in this grand book, the universe\ldots\ It is written in the language of mathematics.}{Galileo Galilei}

V minulé kapitole jsme si řekli, jak o funkcích mluvit. Teď se podíváme na konkrétní funkce, se kterými se budete potkávat neustále: lineární, kvadratické, lineární lomené a mocninné. Ke každé z nich budete u ústní maturity potřebovat umět nakreslit graf a popsat vlastnosti z předchozí kapitoly. Dobrá zpráva je, že když pochopíte, \textit{proč} graf vypadá tak, jak vypadá, nemusíte si skoro nic pamatovat.

\section{Lineární funkce}

\begin{definitionbox}
  \textbf{Lineární funkce} je funkce $f: \mathbb{R}\to\mathbb{R}$ daná předpisem
  \begin{equation*}
    f(x) = ax + b,\qquad a, b\in\mathbb{R}.
  \end{equation*}
  Grafem lineární funkce je \textbf{přímka}. Pro $a = 0$ dostáváme \textbf{konstantní funkci} $f(x) = b$, jejímž grafem je vodorovná přímka. Pro $b = 0$ dostáváme \textbf{přímou úměrnost} $f(x) = ax$, jejíž graf prochází počátkem.

  \vspace{0.2cm}

  \textbf{Vlastnosti}:
  \begin{itemize}
    \item Koeficient $a$ se nazývá \textbf{směrnice}. Udává, o kolik se změní $y$, když se $x$ zvětší o $1$. Pokud přímka svírá s kladnou poloosou $x$ úhel $\alpha$, je $a = \tg\alpha$.
    \item Koeficient $b$ je $y$-ová souřadnice průsečíku s osou $y$, tedy $f(0) = b$.
    \item Pro $a > 0$ je funkce \textbf{rostoucí}, pro $a < 0$ \textbf{klesající}, pro $a = 0$ konstantní.
    \item Pro $a\neq 0$ je $H(f) = \mathbb{R}$, funkce je prostá a inverzní funkce je opět lineární.
  \end{itemize}
\end{definitionbox}

\begin{proof}[(monotonie)]
  Nechť $x_1 < x_2$. Pak $f(x_2) - f(x_1) = ax_2 + b - ax_1 - b = a(x_2 - x_1)$. Protože $x_2 - x_1 > 0$, má rozdíl stejné znaménko jako $a$. Pro $a > 0$ je tedy $f(x_1) < f(x_2)$ a funkce roste, pro $a < 0$ klesá.
\end{proof}

\begin{center}
  \begin{tikzpicture}[scale=0.8]
    \draw[->] (-2.5,0) -- (4,0) node[right] {\small $x$};
    \draw[->] (0,-2) -- (0,4) node[above] {\small $y$};
    \draw[thick, colorofwisdom, domain=-1.5:2.4] plot (\x, {1.2*\x + 1}) node[right] {\small $a > 0$};
    \draw[thick, domain=-2:3.2] plot (\x, {-0.6*\x + 2}) node[above right] {\small $a < 0$};
    \draw[thick, gray, domain=-2.3:3.8] plot (\x, {-1.3}) node[right] {\small $a = 0$};
    \draw[dashed] (1,2.2) -- (2,2.2) node[midway, below] {\tiny $1$} -- (2,3.4) node[midway, right] {\tiny $a$};
    \filldraw[colorofwisdom] (0,1) circle (2pt) node[left] {\small $b$};
  \end{tikzpicture}
\end{center}

Představte si, že jdete po grafu zleva doprava. Směrnice $a$ říká, o kolik vystoupáte, když uděláte krok délky $1$ doprava. Je to přesně \textbf{stoupání} silnice: dopravní značka `stoupání $10\,\%$' znamená, že na $100$ metrů vodorovně vystoupáte o $10$ metrů, tedy $a = 0{,}1$. Záporné $a$ znamená klesání. Koeficient $b$ pak jen říká, v jaké výšce přímka protíná osu $y$, tedy kde `startujeme'.

\begin{center}
  \begin{tikzpicture}[scale=0.75]
    \begin{scope}
      \draw[->] (-2.8,0) -- (3,0) node[right] {\small $x$};
      \draw[->] (0,-1.8) -- (0,3.9) node[above] {\small $y$};
      \draw[thick, colorofwisdom] (-1.25,-1.5) -- (1.25,3.5) node[above] {\scriptsize $a = 2$};
      \draw[thick, colorofwisdom!70] (-2.5,-1.5) -- (2.5,3.5) node[right] {\scriptsize $a = 1$};
      \draw[thick, colorofwisdom!45] (-2.5,-0.25) -- (2.5,2.25) node[right] {\scriptsize $a = \frac{1}{2}$};
      \draw[thick, gray] (-2.5,1) -- (2.5,1) node[right] {\scriptsize $a = 0$};
      \draw[thick] (-2.5,3.5) -- (2.5,-1.5) node[right] {\scriptsize $a = -1$};
      \fill (0,1) circle (2pt);
      \node[align=center] at (0.1,-2.6) {\small stejné $b = 1$, různé $a$:\\\small přímka se otáčí};
    \end{scope}
    \begin{scope}[xshift=7.6cm]
      \draw[->] (-2.8,0) -- (3,0) node[below] {\small $x$};
      \draw[->] (0,-2.4) -- (0,3.9) node[above] {\small $y$};
      \foreach \b/\c in {-1/gray, 0/black, 1/colorofwisdom!60, 2/colorofwisdom} {
        \draw[thick, \c] (-2.5,{\b - 1.25}) -- (2.5,{\b + 1.25}) node[right] {\scriptsize $b = \b$};
        \fill[\c] (0,\b) circle (2pt);
      }
      \node[align=center] at (0.1,-2.6) {\small stejné $a = \frac{1}{2}$, různé $b$:\\\small přímka se posouvá};
    \end{scope}
  \end{tikzpicture}
\end{center}

\begin{example}
  Najděme lineární funkci, jejíž graf prochází body $A = (1, 3)$ a $B = (3, 7)$. Směrnice je podíl změny $y$ a změny $x$:
  \begin{equation*}
    a = \frac{7 - 3}{3 - 1} = 2.
  \end{equation*}
  Na obrázku je to `schod': vodorovně jdeme o $2$, svisle vystoupáme o $4$, tedy na každý krok délky $1$ připadá stoupání $2$.
  \begin{center}
    \begin{tikzpicture}[xscale=0.9, yscale=0.45]
      \draw[->] (-0.5,0) -- (4.2,0) node[right] {\small $x$};
      \draw[->] (0,-0.6) -- (0,8.8) node[above] {\small $y$};
      \foreach \x in {1,2,3} \draw (\x,0.15) -- (\x,-0.15) node[below] {\tiny $\x$};
      \foreach \y in {1,3,7} \draw (0.07,\y) -- (-0.07,\y) node[left] {\tiny $\y$};
      \draw[thick, colorofwisdom, domain=-0.3:3.7] plot (\x, {2*\x + 1}) node[right] {\small $y = 2x + 1$};
      \fill[colorofwisdom!15] (1,3) -- (3,3) -- (3,7) -- cycle;
      \draw[thick] (1,3) -- (3,3) node[midway, below] {\scriptsize $\Delta x = 2$} -- (3,7) node[midway, right] {\scriptsize $\Delta y = 4$};
      \fill (1,3) circle (2.5pt) node[above left] {\small $A$};
      \fill (3,7) circle (2.5pt) node[above left] {\small $B$};
      \fill (0,1) circle (2.5pt);
    \end{tikzpicture}
  \end{center}
  Z $f(1) = 3$ dostaneme $2 + b = 3$, tedy $b = 1$ a $f(x) = 2x + 1$. Jinak řečeno, řešíme soustavu dvou lineárních rovnic $a + b = 3$, $3a + b = 7$ z kapitoly~\ref{kap:rovnice}.
\end{example}

\begin{note}
  Grafy dvou lineárních funkcí jsou rovnoběžné, právě když mají stejnou směrnici. Kolmé jsou (pro nenulové směrnice), právě když $a_1\cdot a_2 = -1$. To dokážeme v kapitole o analytické geometrii.
\end{note}

\begin{problem}
  \begin{enumerate}[label=(\alph*)]
    \item Určete průsečíky grafu $f(x) = -\frac{1}{2}x + 3$ s osami a rozhodněte, zda je $f$ rostoucí.
    \item Najděte lineární funkci, jejíž graf prochází bodem $(2, -1)$ a je rovnoběžný s grafem $y = 3x$.
    \item Jaký úhel svírá graf funkce $y = x$ s osou $x$? A graf $y = \sqrt{3}x$?
    \item Taxi účtuje nástupní sazbu $60$ Kč a $28$ Kč za kilometr. Vyjádřete cenu jako funkci ujetých kilometrů. Kolik kilometrů ujedete za $480$ Kč?
  \end{enumerate}
\end{problem}

\begin{solution}
  \begin{enumerate}[label=(\alph*)]
    \item S osou $y$ v $(0, 3)$, s osou $x$: $-\frac{1}{2}x + 3 = 0$, tedy $(6, 0)$. Směrnice je záporná, funkce je klesající.
    \item $a = 3$, $-1 = 6 + b$, $b = -7$. Tedy $y = 3x - 7$.
    \item $\tg\alpha = 1$, $\alpha = 45^\circ$. $\tg\alpha = \sqrt{3}$, $\alpha = 60^\circ$.
    \item $c(k) = 28k + 60$. Z $28k + 60 = 480$ plyne $k = 15$ kilometrů.
  \end{enumerate}
\end{solution}

\section{Kvadratická funkce}

\begin{definitionbox}
  \textbf{Kvadratická funkce} je funkce $f: \mathbb{R}\to\mathbb{R}$ daná předpisem
  \begin{equation*}
    f(x) = ax^2 + bx + c,\qquad a, b, c\in\mathbb{R},\ a\neq 0.
  \end{equation*}
  Grafem kvadratické funkce je \textbf{parabola}. Doplněním na čtverec (kapitola~\ref{kap:kvadraticke}) ji převedeme na \textbf{vrcholový tvar}
  \begin{equation*}
    f(x) = a{(x - m)}^2 + n,\qquad\text{kde } m = -\frac{b}{2a},\ n = f(m) = c - \frac{b^2}{4a}.
  \end{equation*}
  Bod $V = [m, n]$ je \textbf{vrchol} paraboly a přímka $x = m$ je její \textbf{osa souměrnosti}.

  \vspace{0.2cm}

  \textbf{Vlastnosti} (pro $a > 0$; pro $a < 0$ se vše obrátí):
  \begin{itemize}
    \item parabola je otevřená \textbf{nahoru},
    \item funkce je klesající na $(-\infty, m]$ a rostoucí na $[m, \infty)$,
    \item v bodě $m$ má \textbf{minimum} $n$, maximum nemá,
    \item $H(f) = [n, \infty)$, funkce je zdola omezená a shora neomezená,
    \item průsečík s osou $y$ je $[0, c]$, průsečíky s osou $x$ jsou kořeny rovnice $ax^2 + bx + c = 0$.
  \end{itemize}
\end{definitionbox}

Proč tomu tak je? Ve vrcholovém tvaru je výraz ${(x - m)}^2$ vždy nezáporný a nulový jen pro $x = m$. Pro $a > 0$ je tedy $f(x) = a{(x - m)}^2 + n\geq n$ a rovnost nastane jen ve vrcholu. Zbytek je transformace grafu z kapitoly~\ref{kap:funkce}: parabolu $y = x^2$ roztáhneme $a$-krát, posuneme o $m$ doprava a o $n$ nahoru.

Samotný koeficient $a$ určuje `tvar' paraboly. Čím je $|a|$ větší, tím je parabola užší a strmější. Pro $a$ blízké nule je široká a plochá. Záporné $a$ ji obrátí dnem vzhůru.

\begin{center}
  \begin{tikzpicture}[scale=0.8]
    \draw[->] (-3,0) -- (3.2,0) node[right] {\small $x$};
    \draw[->] (0,-2.9) -- (0,4.2) node[above] {\small $y$};
    \foreach \x in {-2,-1,1,2} \draw (\x,0.07) -- (\x,-0.07) node[below] {\tiny $\x$};
    \draw[thick, colorofwisdom, domain=-1.34:1.34, samples=60] plot (\x, {2*\x*\x}) node[above] {\scriptsize $a = 2$};
    \draw[thick, colorofwisdom!60, domain=-1.9:1.9, samples=60] plot (\x, {\x*\x});
    \node[colorofwisdom!60, above] at (-1.9,3.61) {\scriptsize $a = 1$};
    \draw[thick, gray, domain=-2.68:2.68, samples=60] plot (\x, {0.5*\x*\x}) node[above] {\scriptsize $a = \frac{1}{2}$};
    \draw[thick, domain=-1.6:1.6, samples=60] plot (\x, {-\x*\x}) node[right] {\scriptsize $a = -1$};
  \end{tikzpicture}
\end{center}

\begin{example}
  Vyšetřeme funkci $f(x) = x^2 - 4x + 3$. Doplníme na čtverec:
  \begin{equation*}
    f(x) = (x^2 - 4x + 4) - 1 = {(x - 2)}^2 - 1.
  \end{equation*}
  Vrchol je $V = [2, -1]$, $a = 1 > 0$, parabola je otevřená nahoru. Průsečík s osou $y$: $[0, 3]$. Průsečíky s osou $x$: $x^2 - 4x + 3 = (x - 1)(x - 3) = 0$, tedy $[1, 0]$ a $[3, 0]$. Funkce je klesající na $(-\infty, 2]$, rostoucí na $[2, \infty)$, $H(f) = [-1, \infty)$, minimum $-1$ v bodě $2$.

  Druhá funkce $g(x) = -2x^2 + 4x + 6 = -2(x^2 - 2x) + 6 = -2{(x - 1)}^2 + 8$ má vrchol $[1, 8]$, je otevřená dolů, má maximum $8$ a $H(g) = (-\infty, 8]$. Kořeny jsou $-1$ a $3$.
  \begin{center}
    \begin{tikzpicture}[scale=0.55]
      \draw[->] (-2.5,0) -- (5,0) node[right] {\small $x$};
      \draw[->] (0,-2) -- (0,9) node[above] {\small $y$};
      \foreach \x in {-1,1,2,3,4} \draw (\x,0.1) -- (\x,-0.1) node[below] {\tiny $\x$};
      \foreach \y in {2,4,6,8} \draw (0.1,\y) -- (-0.1,\y) node[left] {\tiny $\y$};
      \draw[thick, colorofwisdom, domain=-0.6:4.6, samples=60] plot (\x, {(\x-2)*(\x-2)-1}) node[right] {\small $f$};
      \draw[thick, domain=-1.3:3.3, samples=60] plot (\x, {-2*(\x-1)*(\x-1)+8}) node[right] {\small $g$};
      \draw[dashed, colorofwisdom] (2,-1) -- (2,6.5);
      \draw[dashed] (1,8) -- (1,-1.5);
      \filldraw[colorofwisdom] (2,-1) circle (3pt);
      \filldraw (1,8) circle (3pt);
    \end{tikzpicture}
  \end{center}
\end{example}

\begin{note}
  Vrchol je přesně uprostřed mezi kořeny (pokud existují). To plyne ze souměrnosti paraboly a z Vietových vzorců: $\frac{x_1 + x_2}{2} = -\frac{b}{2a} = m$. Pokud tedy znáte kořeny, vrchol najdete zpaměti.
\end{note}

Kořeny kvadratické rovnice $ax^2 + bx + c = 0$ jsou přesně místa, kde parabola protne osu $x$. Podle toho, jak vysoko parabola `visí', nastanou tři situace. Diskriminant tedy není nic tajemného, jen říká, jestli vrchol leží pod osou, na ní, nebo nad ní.

\begin{center}
  \begin{tikzpicture}[scale=0.75]
    \foreach \sx/\c/\capt in {0/-1/{$D > 0$: dva kořeny}, 5/0/{$D = 0$: jeden (dvojnásobný)}, 10/1/{$D < 0$: žádný reálný kořen}} {
      \begin{scope}[xshift=\sx cm]
        \draw[->] (-2,0) -- (2.2,0) node[right] {\scriptsize $x$};
        \draw[->] (0,-1.4) -- (0,3.9) node[above] {\scriptsize $y$};
        \draw[thick, colorofwisdom, domain=-1.65:1.65, samples=60] plot (\x, {\x*\x + \c});
        \node at (0,-2) {\scriptsize \capt};
      \end{scope}
    }
    \draw[line width=2.5pt, colorofwisdom!50] (-1,0) -- (1,0);
    \fill (-1,0) circle (2.2pt) (1,0) circle (2.2pt) (5,0) circle (2.2pt);
    \node[above, colorofwisdom] at (-1.05,0.05) {\scriptsize $x_1$};
    \node[above, colorofwisdom] at (1.05,0.05) {\scriptsize $x_2$};
  \end{tikzpicture}
\end{center}

Z obrázku vyčteme i řešení kvadratické nerovnice. Na prvním obrázku je parabola pod osou $x$ (tedy $ax^2 + bx + c < 0$) přesně mezi kořeny, na zvýrazněném intervalu $(x_1, x_2)$, a nad osou vně tohoto intervalu. Pro $a < 0$ je to přesně naopak.

\subsection{Úlohy na extrém}

Protože kvadratická funkce má vždy jeden extrém ve vrcholu, umíme řešit jednoduché optimalizační úlohy, aniž bychom potřebovali derivace.

\begin{example}
  Ze $20$ metrů pletiva chceme oplotit obdélníkový záhon co největšího obsahu. Jaké má mít rozměry?

  Označme jednu stranu $x$. Pak druhá strana je $10 - x$ a obsah je
  \begin{equation*}
    S(x) = x(10 - x) = -x^2 + 10x = -{(x - 5)}^2 + 25.
  \end{equation*}
  Parabola je otevřená dolů a má maximum $25$ v bodě $x = 5$. Nejlepší je tedy čtverec $5\times 5$ metrů. Mezi všemi obdélníky se stejným obvodem má největší obsah čtverec. Tahle odpověď se objevuje v matematice pořád dokola, jen v různých převlecích.

  \begin{center}
    \begin{tikzpicture}
      \begin{scope}[scale=0.25]
        \draw[thick, fill=colorofwisdom!15] (0,0) rectangle (2,8);
        \node[below] at (1,0) {\scriptsize $2\times 8$};
        \node[above] at (1,8) {\scriptsize $S = 16$};
        \draw[thick, fill=colorofwisdom!40] (5,0) rectangle (10,5);
        \node[below] at (7.5,0) {\scriptsize $5\times 5$};
        \node[above] at (7.5,5) {\scriptsize $S = 25$};
        \draw[thick, fill=colorofwisdom!15] (13,0) rectangle (21,2);
        \node[below] at (17,0) {\scriptsize $8\times 2$};
        \node[above] at (17,2) {\scriptsize $S = 16$};
      \end{scope}
      \begin{scope}[xshift=6.6cm, xscale=0.42, yscale=0.1]
        \draw[->] (-0.5,0) -- (11,0) node[right] {\small $x$};
        \draw[->] (0,-2) -- (0,30) node[above] {\small $S$};
        \foreach \x in {2,5,8,10} \draw (\x,0.6) -- (\x,-0.6) node[below] {\tiny $\x$};
        \foreach \y in {16,25} \draw (0.1,\y) -- (-0.1,\y) node[left] {\tiny $\y$};
        \draw[thick, colorofwisdom, domain=0:10, samples=60] plot (\x, {\x*(10 - \x)});
        \draw[dashed, gray] (0,25) -- (5,25) -- (5,0);
        \fill (2,16) circle (0.17 and 0.7) (8,16) circle (0.17 and 0.7);
        \fill[colorofwisdom] (5,25) circle (0.2 and 0.85);
      \end{scope}
    \end{tikzpicture}
  \end{center}
\end{example}

\begin{problem}
  \begin{enumerate}[label=(\alph*)]
    \item Vyšetřete funkci $f(x) = 2x^2 + 8x + 5$ (vrchol, monotonie, obor hodnot, průsečíky s osami).
    \item Určete kvadratickou funkci, jejíž graf má vrchol $[1, -4]$ a prochází bodem $[3, 4]$.
    \item Součet dvou čísel je $14$. Pro jaká čísla je jejich součin největší?
    \item Kámen vyhozený svisle vzhůru je v čase $t$ ve výšce $h(t) = 20t - 5t^2$ metrů. Kdy dosáhne nejvyššího bodu a jak vysoko to je?
  \end{enumerate}
\end{problem}

\begin{solution}
  \begin{enumerate}[label=(\alph*)]
    \item $f(x) = 2(x^2 + 4x) + 5 = 2{(x + 2)}^2 - 3$. Vrchol $[-2, -3]$, klesající na $(-\infty, -2]$, rostoucí na $[-2, \infty)$, $H = [-3, \infty)$. Osa $y$: $[0, 5]$. Osa $x$: $2x^2 + 8x + 5 = 0$, $x = \frac{-8\pm\sqrt{24}}{4} = -2\pm\frac{\sqrt{6}}{2}$.
    \item $f(x) = a{(x - 1)}^2 - 4$, z $f(3) = 4$ plyne $4a - 4 = 4$, $a = 2$. Tedy $f(x) = 2{(x - 1)}^2 - 4 = 2x^2 - 4x - 2$.
    \item $x(14 - x) = -{(x - 7)}^2 + 49$, maximum pro $x = 7$, čísla jsou $7$ a $7$.
    \item $h(t) = -5(t^2 - 4t) = -5{(t - 2)}^2 + 20$. Nejvýš je v čase $t = 2\ \si{s}$ ve výšce $20\ \si{m}$.
  \end{enumerate}
\end{solution}

\section{Nepřímá úměrnost a lineární lomená funkce}

\begin{definitionbox}
  \textbf{Nepřímá úměrnost} je funkce $f(x) = \frac{k}{x}$, kde $k\neq 0$. Grafem je \textbf{rovnoosá hyperbola} se \textbf{asymptotami} $x = 0$ a $y = 0$.
  \begin{itemize}
    \item $D(f) = \mathbb{R}\setminus\{0\}$, $H(f) = \mathbb{R}\setminus\{0\}$,
    \item funkce je lichá a prostá,
    \item pro $k > 0$ je klesající na $(-\infty, 0)$ a na $(0, \infty)$ (graf v 1. a 3. kvadrantu), pro $k < 0$ je rostoucí na každém z těchto intervalů (graf ve 2. a 4. kvadrantu),
    \item není omezená a nemá extrémy.
  \end{itemize}

  \textbf{Lineární lomená funkce} je funkce
  \begin{equation*}
    f(x) = \frac{ax + b}{cx + d},\qquad c\neq 0,\ ad - bc\neq 0.
  \end{equation*}
  Úpravou ji vždy převedeme na tvar
  \begin{equation*}
    f(x) = \frac{k}{x - m} + n,
  \end{equation*}
  ze kterého je vidět, že jde o hyperbolu $y = \frac{k}{x}$ posunutou o $m$ doprava a o $n$ nahoru. Asymptoty jsou $x = m$ a $y = n$, $D(f) = \mathbb{R}\setminus\{m\}$ a $H(f) = \mathbb{R}\setminus\{n\}$.
\end{definitionbox}

Podmínky v definici mají svůj smysl. Pro $c = 0$ by šlo o lineární funkci. Pro $ad - bc = 0$ (ano, to je determinant z kapitoly~\ref{kap:rovnice}) by se čitatel a jmenovatel daly zkrátit a dostali bychom konstantní funkci s jedním vyňatým bodem.

Proč má graf nepřímé úměrnosti právě takový tvar? Rovnost $y = \frac{k}{x}$ je totéž jako $x\cdot y = k$. Každý bod grafu je tedy rohem obdélníku (s druhým rohem v počátku), který má obsah přesně $k$. Když obdélník protahujeme do šířky, musí se stejně rychle zplošťovat. Proto se graf přimyká k osám, ale nikdy se jich nedotkne: obdélník s nulovou výškou by měl obsah nula, ne $k$.

\begin{center}
  \begin{tikzpicture}[scale=0.7]
    \begin{scope}
      \fill[colorofwisdom!25] (0,0) rectangle (1,2);
      \fill[colorofwisdom!12] (0,0) rectangle (2.5,0.8);
      \draw (1,0) -- (1,2) -- (0,2);
      \draw (2.5,0) -- (2.5,0.8) -- (0,0.8);
      \draw[->] (-3.6,0) -- (3.8,0) node[right] {\small $x$};
      \draw[->] (0,-3.6) -- (0,3.8) node[above] {\small $y$};
      \draw[thick, colorofwisdom, domain=0.6:3.4, samples=60] plot (\x, {2/\x});
      \draw[thick, colorofwisdom, domain=-3.4:-0.6, samples=60] plot (\x, {2/\x});
      \fill (1,2) circle (2pt) node[right] {\scriptsize $[1, 2]$};
      \fill (2.5,0.8) circle (2pt) node[above right] {\scriptsize $[2{,}5;\, 0{,}8]$};
      \node at (0,-4.3) {\small $y = \frac{2}{x}$: oba obdélníky mají obsah $2$};
    \end{scope}
    \begin{scope}[xshift=9cm]
      \draw[->] (-3.6,0) -- (3.8,0) node[right] {\small $x$};
      \draw[->] (0,-3.6) -- (0,3.8) node[above] {\small $y$};
      \draw[thick, domain=0.6:3.4, samples=60] plot (\x, {-2/\x});
      \draw[thick, domain=-3.4:-0.6, samples=60] plot (\x, {-2/\x});
      \node at (0,-4.3) {\small $y = -\frac{2}{x}$: pro $k < 0$ 2. a 4. kvadrant};
    \end{scope}
  \end{tikzpicture}
\end{center}

A proč jsou u lineární lomené funkce $\frac{k}{x - m} + n$ asymptoty právě $x = m$ a $y = n$? Když je $x$ hodně daleko od $m$, je zlomek $\frac{k}{x - m}$ skoro nula a funkce má hodnotu skoro $n$. Když se naopak $x$ blíží k $m$, dělíme skoro nulou a hodnota utíká do nekonečna.

\begin{example}
  Upravme $f(x) = \frac{2x + 1}{x - 1}$. Trik: v čitateli `vyrobíme' jmenovatel.
  \begin{equation*}
    \frac{2x + 1}{x - 1} = \frac{2(x - 1) + 3}{x - 1} = 2 + \frac{3}{x - 1}.
  \end{equation*}
  Jde tedy o hyperbolu $y = \frac{3}{x}$ posunutou o $1$ doprava a o $2$ nahoru. Asymptoty $x = 1$, $y = 2$, $D(f) = \mathbb{R}\setminus\{1\}$, $H(f) = \mathbb{R}\setminus\{2\}$. Protože $k = 3 > 0$, je funkce klesající na $(-\infty, 1)$ i na $(1, \infty)$. Průsečíky: s osou $y$ v $[0, -1]$, s osou $x$ v $\left[-\frac{1}{2}, 0\right]$.
  \begin{center}
    \begin{tikzpicture}[scale=0.6]
      \draw[->] (-4.5,0) -- (6,0) node[right] {\small $x$};
      \draw[->] (0,-4) -- (0,7) node[above] {\small $y$};
      \draw[dashed, gray] (1,-4) -- (1,7) node[above right] {\small $x = 1$};
      \draw[dashed, gray] (-4.5,2) -- (6,2) node[right] {\small $y = 2$};
      \draw[thick, colorofwisdom, domain=-4.5:0.4, samples=60] plot (\x, {2 + 3/(\x - 1)});
      \draw[thick, colorofwisdom, domain=1.6:6, samples=60] plot (\x, {2 + 3/(\x - 1)});
      \filldraw (0,-1) circle (2.5pt);
      \filldraw (-0.5,0) circle (2.5pt);
    \end{tikzpicture}
  \end{center}
\end{example}

\begin{problem}
  Upravte na tvar $\frac{k}{x - m} + n$, určete asymptoty, definiční obor, obor hodnot a monotonii:
  \begin{enumerate}[label=(\alph*)]
    \item $f(x) = \dfrac{x + 3}{x + 1}$,
    \item $g(x) = \dfrac{3x - 1}{x - 2}$,
    \item $h(x) = \dfrac{-2x + 5}{x - 2}$.
  \end{enumerate}
\end{problem}

\begin{solution}
  \begin{enumerate}[label=(\alph*)]
    \item $\frac{(x + 1) + 2}{x + 1} = 1 + \frac{2}{x + 1}$. Asymptoty $x = -1$, $y = 1$. $D = \mathbb{R}\setminus\{-1\}$, $H = \mathbb{R}\setminus\{1\}$. Klesající na $(-\infty, -1)$ a na $(-1, \infty)$.
    \item $\frac{3(x - 2) + 5}{x - 2} = 3 + \frac{5}{x - 2}$. Asymptoty $x = 2$, $y = 3$. Klesající na obou intervalech.
    \item $\frac{-2(x - 2) + 1}{x - 2} = -2 + \frac{1}{x - 2}$. Asymptoty $x = 2$, $y = -2$. Klesající na $(-\infty, 2)$ a na $(2, \infty)$.
  \end{enumerate}
\end{solution}

\begin{example}
  Nepřímá úměrnost je všude kolem nás. Pokud jedeme $120\ \si{km}$, je doba jízdy $t = \frac{120}{v}$. Dvakrát rychleji znamená dvakrát kratší dobu. Pokud si $n$ lidí dělí pizzu, dostane každý $\frac{1}{n}$ pizzy. Graf ukazuje, že pro velká $n$ se kousek blíží nule, ale nikdy jí nedosáhne. Což z vlastní zkušenosti z oslav nemusí být pravda.
\end{example}

\section{Mocninné funkce}

\begin{definitionbox}
  \textbf{Mocninná funkce} je funkce $f(x) = x^r$, kde $r\in\mathbb{R}$ je konstanta. U maturity nás zajímají hlavně tyto případy:
  \begin{center}
    \renewcommand{\arraystretch}{1.4}
    \begin{tabular}{|l|l|l|l|}
      \hline
      Exponent & Příklad & $D(f)$ & Vlastnosti \\
      \hline
      $n\in\mathbb{N}$ sudé & $x^2, x^4$ & $\mathbb{R}$ & sudá, $H = [0, \infty)$, klesající na $(-\infty, 0]$, rostoucí na $[0, \infty)$ \\
      $n\in\mathbb{N}$ liché & $x, x^3$ & $\mathbb{R}$ & lichá, $H = \mathbb{R}$, rostoucí, prostá \\
      $-n$, $n$ sudé & $x^{-2}$ & $\mathbb{R}\setminus\{0\}$ & sudá, $H = (0, \infty)$, rostoucí na $(-\infty, 0)$, klesající na $(0, \infty)$ \\
      $-n$, $n$ liché & $x^{-1}$ & $\mathbb{R}\setminus\{0\}$ & lichá, $H = \mathbb{R}\setminus\{0\}$, klesající na $(-\infty, 0)$ i na $(0, \infty)$ \\
      $\frac{1}{n}$ & $\sqrt{x}, \sqrt[3]{x}$ & $[0, \infty)$ & rostoucí, $H = [0, \infty)$, inverzní k $x^n$ na $[0, \infty)$ \\
      \hline
    \end{tabular}
  \end{center}
\end{definitionbox}

\begin{center}
  \begin{tikzpicture}[scale=0.85]
    \begin{scope}
      \draw[->] (-2.2,0) -- (2.4,0) node[right] {\small $x$};
      \draw[->] (0,-2.2) -- (0,2.6) node[above] {\small $y$};
      \draw[thick, colorofwisdom, domain=-1.5:1.5, samples=60] plot (\x, {\x*\x}) node[above] {\small $x^2$};
      \draw[thick, domain=-1.25:1.25, samples=60] plot (\x, {\x*\x*\x*\x});
      \node[left] at (-1.25,2.44) {\small $x^4$};
      \draw[thick, dashed, domain=-1.3:1.3, samples=60] plot (\x, {\x*\x*\x});
      \node[left] at (-1.3,-2.197) {\small $x^3$};
      \filldraw (1,1) circle (1.5pt) (-1,1) circle (1.5pt) (-1,-1) circle (1.5pt);
      \node at (0,-2.7) {\small kladné celé exponenty};
    \end{scope}
    \begin{scope}[xshift=6.5cm]
      \draw[->] (-2.2,0) -- (2.4,0) node[right] {\small $x$};
      \draw[->] (0,-2.2) -- (0,2.6) node[above] {\small $y$};
      \draw[thick, colorofwisdom, domain=0.42:2.2, samples=60] plot (\x, {1/\x});
      \node[below left, colorofwisdom] at (-0.9,-1.11) {\small $x^{-1}$};
      \draw[thick, colorofwisdom, domain=-2.2:-0.42, samples=60] plot (\x, {1/\x});
      \draw[thick, domain=0.65:2.2, samples=60] plot (\x, {1/(\x*\x)}) node[right] {\small $x^{-2}$};
      \draw[thick, domain=-2.2:-0.65, samples=60] plot (\x, {1/(\x*\x)});
      \node at (0,-2.7) {\small záporné exponenty};
    \end{scope}
    \begin{scope}[xshift=12cm]
      \draw[->] (-0.5,0) -- (3.2,0) node[right] {\small $x$};
      \draw[->] (0,-0.5) -- (0,2.6) node[above] {\small $y$};
      \draw[thick, colorofwisdom, domain=0:3, samples=80] plot (\x, {sqrt(\x)}) node[above] {\small $\sqrt{x}$};
      \draw[thick, domain=0:3, samples=80] plot (\x, {\x^(1/3)}) node[below right] {\small $\sqrt[3]{x}$};
      \filldraw (1,1) circle (1.5pt);
      \node at (1.4,-1.2) {\small odmocniny};
    \end{scope}
  \end{tikzpicture}
\end{center}

Všimněte si, že všechny grafy procházejí bodem $[1, 1]$, protože $1^r = 1$ pro každé $r$. Pro $x\in(0, 1)$ je $x^4 < x^2 < x$, ale pro $x > 1$ naopak $x^4 > x^2 > x$. Vyšší mocnina je u nuly `plošší' a dál od nuly `strmější'. Totéž platí obráceně pro odmocniny.

\begin{center}
  \begin{tikzpicture}[scale=1.8]
    \fill[colorofwisdom!10] (0,0) rectangle (1,1);
    \draw[dashed, gray] (1,0) -- (1,2.25);
    \draw[dashed, gray] (0,1) -- (2.05,1);
    \draw[->] (-0.1,0) -- (2.3,0) node[right] {\small $x$};
    \draw[->] (0,-0.1) -- (0,2.4) node[above] {\small $y$};
    \node[below] at (1,0) {\scriptsize $1$};
    \node[left] at (0,1) {\scriptsize $1$};
    \draw[thick, gray, domain=0:2] plot (\x, {\x}) node[right] {\scriptsize $x$};
    \draw[thick, colorofwisdom!60, domain=0:1.48, samples=60] plot (\x, {\x*\x}) node[right] {\scriptsize $x^2$};
    \draw[thick, colorofwisdom, domain=0:1.3, samples=60] plot (\x, {\x*\x*\x}) node[left] {\scriptsize $x^3$};
    \fill (1,1) circle (0.8pt);
    \node at (0.5,-0.2) {\scriptsize $x^3 < x^2 < x$};
    \node at (1.65,-0.2) {\scriptsize $x^3 > x^2 > x$};
  \end{tikzpicture}
\end{center}

Proč? Násobit číslem menším než $1$ znamená zmenšovat. Číslo $0{,}5$ dá $0{,}5^2 = 0{,}25$ a $0{,}5^3 = 0{,}125$, každé další násobení výsledek zmenší. Naopak $2$ dá $4$ a $8$, každé násobení zvětšuje. Bod $x = 1$ je hranice, kde se nic nemění.

\begin{note}
  Proč je $x^{-2}$ rostoucí na $(-\infty, 0)$? Pro $x_1 < x_2 < 0$ je $x_1^2 > x_2^2$ (čím dál od nuly, tím větší druhá mocnina) a převrácené hodnoty se otočí, tedy $\frac{1}{x_1^2} < \frac{1}{x_2^2}$. Na grafu: když jdeme zleva k nule, hodnoty rostou do nekonečna.
\end{note}

\begin{problem}
  \begin{enumerate}[label=(\alph*)]
    \item Seřaďte podle velikosti čísla $0{,}5^2$, $0{,}5^3$, $0{,}5^{\frac{1}{2}}$, $0{,}5^{-1}$.
    \item Nakreslete graf funkce $f(x) = {(x + 1)}^3 - 2$ a určete její obor hodnot a průsečík s osou $y$.
    \item Pro která $x$ platí $x^3 > x$?
  \end{enumerate}
\end{problem}

\begin{solution}
  \begin{enumerate}[label=(\alph*)]
    \item $0{,}5^3 = 0{,}125 < 0{,}5^2 = 0{,}25 < 0{,}5^{\frac{1}{2}}\approx 0{,}707 < 0{,}5^{-1} = 2$. Pro základ menší než $1$ s rostoucím exponentem hodnota klesá.
    \item Posun grafu $y = x^3$ o $1$ doleva a o $2$ dolů. $H = \mathbb{R}$, průsečík s osou $y$: $f(0) = 1 - 2 = -1$.
    \item $x^3 - x = x(x - 1)(x + 1) > 0$. Metodou nulových bodů $K = (-1, 0)\cup(1, \infty)$.
  \end{enumerate}
\end{solution}

\section{Příběh --- Hořící zrcadla}

Proč se parabola jmenuje parabola? Název pochází od řeckého matematika Apollónia z Pergy, který zkoumal křivky vzniklé řezem kužele rovinou. Parabola (z řeckého \textit{parabolé}, `přiložení vedle') vznikne, když je řez rovnoběžný s povrchovou přímkou kužele. Víc o tom v kapitole o kuželosečkách.

Parabola má jednu výjimečnou vlastnost. Všechny paprsky, které na ni dopadají rovnoběžně s osou, se po odrazu setkají v jednom bodě, v takzvaném \textbf{ohnisku}. Proto mají satelitní antény, reflektory aut i solární vařiče parabolický tvar. Funguje to i obráceně: zdroj světla v ohnisku vyšle rovnoběžný svazek paprsků.

Podle legendy toho využil Archimédés při obraně Syrakus proti Římanům v roce 212 př.~n.~l. Prý nechal vojáky postavit obrovská zrcadla, kterými soustředil sluneční paprsky na římské lodě a zapálil je. Jestli se to opravdu stalo, je sporné. V roce 2005 to zkusili studenti z MIT s dřevěnou lodí a 127 zrcadly a loď sice začala hořet, ale trvalo to dlouho a muselo být ideální počasí.\footnote{Pořad Bořiči mýtů (Mythbusters) to zkusil několikrát a mýtus prohlásil za vyvrácený. Takže záleží, komu věříte.}

Syrakusy nakonec stejně padly. Podle Plútarcha byl Archimédés v té době zabraný do geometrického problému a kreslil obrazce do písku. Když k němu přišel římský voják, Archimédés ho prý okřikl `Neruš mé kruhy!' a voják ho na místě zabil. Římský velitel Marcellus, který chtěl Archiméda zajmout živého, byl prý jeho smrtí rozezlen. Archimédés si přál, aby mu na hrob vytesali kouli vepsanou do válce, protože objev vztahu mezi jejich objemy považoval za svůj největší úspěch. O 137 let později našel Cicero zarostlý a zapomenutý hrob právě podle tohoto obrazce. Ten vztah si ukážeme v kapitole o stereometrii.

\section{Shrnutí a otázky k ústní zkoušce}

\begin{problem}
  Odpovězte vlastními slovy:
  \begin{enumerate}
    \item Co je lineární funkce? Jaký význam mají koeficienty $a$ a $b$? Dokažte, kdy je lineární funkce rostoucí.
    \item Jak najdete lineární funkci, jejíž graf prochází dvěma danými body?
    \item Co je kvadratická funkce? Jak najdete vrchol paraboly? Popište vlastnosti kvadratické funkce v závislosti na znaménku $a$.
    \item Jak souvisí graf kvadratické funkce s řešením kvadratické rovnice a nerovnice?
    \item Co je nepřímá úměrnost a lineární lomená funkce? Jak najdete asymptoty?
    \item Proč v definici lineární lomené funkce požadujeme $ad - bc\neq 0$?
    \item Popište vlastnosti mocninných funkcí s kladným celým, záporným celým a racionálním exponentem. Nakreslete grafy.
    \item Uveďte příklady využití lineárních, kvadratických a lomených funkcí v praxi.
  \end{enumerate}
\end{problem}

\koneckapitoly

\chapter{Exponenciální a logaritmické funkce}
\label{kap:exponencialni}
\mathquote{[Logarithms,] by shortening the labours, doubled the life of the astronomer.}{Pierre-Simon Laplace}

Představte si, že vám někdo nabídne dvě možnosti. Buď dostanete milion korun hned, nebo dostanete dnes jednu korunu, zítra dvě, pozítří čtyři a každý další den dvojnásobek předchozího, a to po dobu jednoho měsíce. Většina lidí sáhne po milionu. Třicátý den by ale druhá možnost vyplatila $2^{29}$, tedy přes půl miliardy korun, a celkem dohromady přes miliardu. Lidský mozek je prostě na exponenciální růst krátký. V této kapitole se podíváme na funkce, které takový růst popisují, a na logaritmy, které jsou jejich opakem.

\section{Exponenciální funkce}

V kapitole~\ref{kap:vyrazy} jsme si postupně zadefinovali mocniny s přirozeným, celým, racionálním a nakonec i iracionálním exponentem. Pro kladný základ $a$ tedy má výraz $a^x$ smysl pro každé reálné $x$. Mocninná funkce měla proměnnou v základu a pevný exponent. Teď to otočíme.

\begin{definitionbox}
  \textbf{Exponenciální funkce} o základu $a$ je funkce
  \begin{equation*}
    f(x) = a^x,\qquad a > 0,\ a\neq 1.
  \end{equation*}

  \textbf{Vlastnosti}:
  \begin{itemize}
    \item $D(f) = \mathbb{R}$, $H(f) = (0, \infty)$,
    \item pro $a > 1$ je \textbf{rostoucí}, pro $0 < a < 1$ \textbf{klesající},
    \item je prostá, zdola omezená (nulou), shora neomezená, nemá extrémy,
    \item graf prochází bodem $[0, 1]$, protože $a^0 = 1$, a bodem $[1, a]$,
    \item osa $x$ (přímka $y = 0$) je \textbf{asymptota} grafu.
  \end{itemize}
\end{definitionbox}

Proč podmínky na základ? Pro $a = 1$ bychom dostali konstantní funkci $1^x = 1$, což je nuda. Pro $a\leq 0$ by výraz $a^x$ nebyl definovaný pro mnoho $x$ (zkuste ${(-2)}^{\frac{1}{2}}$), jak jsme viděli v kapitole~\ref{kap:vyrazy}.

\begin{center}
  \begin{tikzpicture}[scale=0.85]
    \draw[->] (-3.5,0) -- (3.5,0) node[right] {\small $x$};
    \draw[->] (0,-0.5) -- (0,5) node[above] {\small $y$};
    \foreach \x in {-3,-2,-1,1,2,3} \draw (\x,0.08) -- (\x,-0.08) node[below] {\tiny $\x$};
    \foreach \y in {1,2,3,4} \draw (0.08,\y) -- (-0.08,\y) node[left] {\tiny $\y$};
    \draw[thick, colorofwisdom, domain=-3.4:2.2, samples=60] plot (\x, {pow(2,\x)}) node[right] {\small $y = 2^x$};
    \draw[thick, domain=-2.2:3.4, samples=60] plot (\x, {pow(0.5,\x)});
    \node at (-2.6,4.6) {\small $y = \left(\frac{1}{2}\right)^x$};
    \filldraw (0,1) circle (2pt);
  \end{tikzpicture}
\end{center}

Grafy $2^x$ a ${\left(\frac{1}{2}\right)}^x$ jsou souměrné podle osy $y$, protože ${\left(\frac{1}{2}\right)}^x = 2^{-x}$. To je transformace $f(-x)$ z kapitoly~\ref{kap:funkce}.

\begin{proof}[(monotonie pro $a > 1$)]
  Nechť $x_1 < x_2$. Pak $a^{x_2} = a^{x_1}\cdot a^{x_2 - x_1}$. Protože $x_2 - x_1 > 0$ a $a > 1$, je $a^{x_2 - x_1} > 1$.\footnote{Pro přirozený exponent je to jasné (součin čísel větších než $1$), pro racionální plyne z vlastností odmocnin a pro iracionální z aproximace. Detaily vynecháme.} Vynásobením kladného čísla $a^{x_1}$ číslem větším než $1$ dostaneme větší číslo, a tedy $a^{x_1} < a^{x_2}$.
\end{proof}

\subsection{Eulerovo číslo}

Mezi všemi možnými základy je jeden výjimečný. Představte si banku, která vám nabídne úrok $100\,\%$ ročně. Vložíte $1$ korunu a za rok máte $2$. Co když vám ale úrok připíše po půl roce (každé pololetí $50\,\%$)? Pak máte za rok ${(1 + \frac{1}{2})}^2 = 2{,}25$ korun. Při čtvrtletním připisování ${(1 + \frac{1}{4})}^4\approx 2{,}44$ korun. A co když se úrok připisuje každou vteřinu?
\begin{center}
  \renewcommand{\arraystretch}{1.3}
  \begin{tabular}{|c|c|c|c|c|c|c|}
    \hline
    $n$ & $1$ & $2$ & $4$ & $12$ & $365$ & $1\,000\,000$ \\
    \hline
    ${\left(1 + \frac{1}{n}\right)}^n$ & $2$ & $2{,}25$ & $2{,}441$ & $2{,}613$ & $2{,}7146$ & $2{,}71828$ \\
    \hline
  \end{tabular}
\end{center}
Hodnoty rostou, ale ne nade všechny meze. Blíží se k číslu
\begin{equation*}
  e = 2{,}718\,281\,828\,459\ldots,
\end{equation*}
kterému se říká \textbf{Eulerovo číslo}. Je iracionální (a dokonce transcendentní) a je to jedno z nejdůležitějších čísel matematiky, hned vedle $0$, $1$, $\pi$ a $i$. Exponenciální funkce $e^x$ se často nazývá prostě `exponenciála' a v kapitole o derivacích uvidíme, proč je tak výjimečná: je to jediná funkce (až na násobek), která je sama sobě derivací. Přesně řečeno je $e$ limitou posloupnosti ${\left(1 + \frac{1}{n}\right)}^n$, ale limity přijdou až později.

\section{Logaritmus}

Rovnici $2^x = 8$ vyřeší každý zpaměti: $x = 3$. Ale co rovnice $2^x = 5$? Řešení existuje (exponenciální funkce je prostá a nabývá všech kladných hodnot), ale nedá se vyjádřit pomocí operací, které zatím známe. Potřebujeme nový symbol.

\begin{definitionbox}
  \textbf{Logaritmus} čísla $x$ o základu $a$ (kde $a > 0$, $a\neq 1$ a $x > 0$) je číslo $y$, pro které platí $a^y = x$. Značíme ho $\log_a x$:
  \begin{equation*}
    \log_a x = y\quad\Leftrightarrow\quad a^y = x.
  \end{equation*}
  Slovy: logaritmus odpovídá na otázku `na kolikátou musím umocnit $a$, abych dostal $x$?'

  \vspace{0.2cm}

  \textbf{Speciální značení}:
  \begin{itemize}
    \item \textbf{dekadický logaritmus} (základ $10$): $\log x = \log_{10} x$,
    \item \textbf{přirozený logaritmus} (základ $e$): $\ln x = \log_e x$.
  \end{itemize}
\end{definitionbox}

\begin{example}
  Spočítejme podle definice:
  \begin{align*}
    \log_2 8 &= 3, &&\text{protože } 2^3 = 8,\\
    \log_3\tfrac{1}{9} &= -2, &&\text{protože } 3^{-2} = \tfrac{1}{9},\\
    \log_5\sqrt{5} &= \tfrac{1}{2}, &&\text{protože } 5^{\frac{1}{2}} = \sqrt{5},\\
    \log 0{,}001 &= -3, &&\text{protože } 10^{-3} = 0{,}001,\\
    \log_{\frac{1}{2}} 8 &= -3, &&\text{protože } {\left(\tfrac{1}{2}\right)}^{-3} = 8,\\
    \log_a 1 &= 0, &&\text{protože } a^0 = 1,\\
    \log_a a &= 1, &&\text{protože } a^1 = a.
  \end{align*}
  Logaritmus záporného čísla ani nuly neexistuje (v reálných číslech), protože $a^y$ je vždy kladné.
\end{example}

\begin{problem}
  Spočítejte: $\log_2 32$, $\log_4 2$, $\log_9 27$, $\log_{\frac{1}{3}} 81$, $\ln e^5$, $\log_7 1$, $\log 1000$, $\log_2\frac{\sqrt{2}}{4}$.
\end{problem}

\begin{solution}
  $5$, $\frac{1}{2}$, $\frac{3}{2}$ (protože $9^{\frac{3}{2}} = 27$), $-4$, $5$, $0$, $3$ a $\log_2 2^{\frac{1}{2} - 2} = -\frac{3}{2}$.
\end{solution}

\subsection{Logaritmická funkce}

\begin{definitionbox}
  \textbf{Logaritmická funkce} o základu $a$ ($a > 0$, $a\neq 1$) je funkce $f(x) = \log_a x$. Je to \textbf{inverzní funkce} k exponenciální funkci $a^x$. Proto:
  \begin{itemize}
    \item $D(f) = (0, \infty)$, $H(f) = \mathbb{R}$ (prohodily se oproti exponenciální funkci),
    \item pro $a > 1$ je rostoucí, pro $0 < a < 1$ klesající,
    \item je prostá, neomezená, nemá extrémy,
    \item graf prochází bodem $[1, 0]$ a bodem $[a, 1]$,
    \item osa $y$ je asymptota grafu,
    \item graf je souměrný s grafem $a^x$ podle přímky $y = x$.
  \end{itemize}
  Pro všechna přípustná $x$ platí
  \begin{equation*}
    a^{\log_a x} = x\qquad\text{a}\qquad\log_a a^x = x.
  \end{equation*}
\end{definitionbox}

\begin{center}
  \begin{tikzpicture}[scale=0.8]
    \draw[->] (-3,0) -- (5,0) node[right] {\small $x$};
    \draw[->] (0,-3) -- (0,5) node[above] {\small $y$};
    \foreach \x in {-2,-1,1,2,3,4} \draw (\x,0.08) -- (\x,-0.08) node[below] {\tiny $\x$};
    \draw[dashed, gray] (-2.8,-2.8) -- (4.6,4.6) node[above right] {\small $y = x$};
    \draw[thick, gray, domain=-2.9:2.25, samples=60] plot (\x, {pow(2,\x)}) node[above] {\small $y = 2^x$};
    \draw[thick, colorofwisdom, domain=0.13:4.8, samples=80] plot (\x, {ln(\x)/ln(2)}) node[right] {\small $y = \log_2 x$};
    \filldraw (1,0) circle (2pt) (0,1) circle (2pt);
  \end{tikzpicture}
\end{center}

\subsection{Pravidla pro počítání s logaritmy}

\begin{definitionbox}
  Pro $a, b > 0$, $a, b\neq 1$, $x, y > 0$ a $r\in\mathbb{R}$ platí:
  \begin{enumerate}
    \item $\log_a(x\cdot y) = \log_a x + \log_a y$,
    \item $\log_a\dfrac{x}{y} = \log_a x - \log_a y$,
    \item $\log_a x^r = r\cdot\log_a x$,
    \item $\log_a x = \dfrac{\log_b x}{\log_b a}$ \quad(\textbf{převod základu}).
  \end{enumerate}
\end{definitionbox}

Logaritmus tedy převádí násobení na sčítání a mocnění na násobení. To je přesně důvod, proč byly logaritmy po staletí tak cenné: sčítat se dá mnohem snáz než násobit.

\begin{proof}
  Všechna pravidla jsou jen přepisem pravidel pro mocniny. Označme $u = \log_a x$ a $v = \log_a y$, tedy $a^u = x$ a $a^v = y$.
  \begin{enumerate}
    \item $xy = a^u\cdot a^v = a^{u + v}$, a tedy podle definice $\log_a(xy) = u + v$.
    \item $\frac{x}{y} = \frac{a^u}{a^v} = a^{u - v}$, tedy $\log_a\frac{x}{y} = u - v$.
    \item $x^r = {(a^u)}^r = a^{ru}$, tedy $\log_a x^r = ru$.
    \item Zlogaritmujeme rovnost $x = a^u$ o základu $b$ a použijeme pravidlo 3: $\log_b x = \log_b a^u = u\cdot\log_b a$. Protože $a\neq 1$, je $\log_b a\neq 0$ a můžeme dělit.
  \end{enumerate}
\end{proof}

\begin{note}
  Pro součet žádné pravidlo neexistuje! $\log(x + y)\neq\log x + \log y$. Tohle je stejná chyba jako ${(a + b)}^2 = a^2 + b^2$, jen v jiném převleku. Pravidlo 1 platí pro logaritmus součinu, ne součtu.
\end{note}

\begin{example}
  Upravme:
  \begin{align*}
    \log 2 + \log 50 &= \log 100 = 2,\\
    \log_3 54 - \log_3 2 &= \log_3 27 = 3,\\
    2\log_5 10 - \log_5 4 &= \log_5 100 - \log_5 4 = \log_5 25 = 2,\\
    \log_2 3\cdot\log_3 8 &= \log_2 3\cdot\frac{\log_2 8}{\log_2 3} = \log_2 8 = 3.
  \end{align*}
\end{example}

\begin{note}
  Pravidlo pro převod základu je důvod, proč na kalkulačce stačí tlačítka $\log$ a $\ln$. Například $\log_2 5 = \frac{\log 5}{\log 2}\approx\frac{0{,}699}{0{,}301}\approx 2{,}32$.
\end{note}

\begin{problem}
  \begin{enumerate}[label=(\alph*)]
    \item Spočítejte $\log_6 4 + \log_6 9$, $\log_2 40 - \log_2 5$, $\frac{1}{2}\log_3 81$, $5^{\log_5 7}$, $10^{2\log 3}$.
    \item Vyjádřete pomocí $\log x$, $\log y$ a $\log z$: $\log\dfrac{x^3\sqrt{y}}{z^2}$.
    \item Zapište jako jeden logaritmus: $2\ln a + \ln b - 3\ln c$.
  \end{enumerate}
\end{problem}

\begin{solution}
  \begin{enumerate}[label=(\alph*)]
    \item $\log_6 36 = 2$, $\log_2 8 = 3$, $\frac{1}{2}\cdot 4 = 2$, $7$, $10^{\log 9} = 9$.
    \item $3\log x + \frac{1}{2}\log y - 2\log z$.
    \item $\ln\dfrac{a^2b}{c^3}$.
  \end{enumerate}
\end{solution}

\section{Exponenciální rovnice}

\begin{definitionbox}
  \textbf{Exponenciální rovnice} je rovnice, ve které se neznámá vyskytuje v exponentu. Základní metody:
  \begin{enumerate}
    \item \textbf{Převod na stejný základ}: Protože je exponenciální funkce prostá, platí
    \begin{equation*}
      a^{f(x)} = a^{g(x)}\quad\Leftrightarrow\quad f(x) = g(x).
    \end{equation*}
    \item \textbf{Vytýkání}: Pokud se v rovnici vyskytuje $a^{x + c}$ pro různá $c$, vytkneme $a^x$.
    \item \textbf{Substituce}: Pokud se v rovnici vyskytuje $a^x$ a $a^{2x} = {(a^x)}^2$, zavedeme $t = a^x$. Pozor, $t > 0$!
    \item \textbf{Logaritmování}: Pokud nelze převést na stejný základ, obě strany zlogaritmujeme.
  \end{enumerate}
\end{definitionbox}

\begin{example}
  Vyřešme:
  \begin{enumerate}[label=(\alph*)]
    \item $2^{x + 1} = 32$. Tedy $2^{x + 1} = 2^5$, $x + 1 = 5$ a $x = 4$.
    \item $9^x = 27$. Převedeme na základ $3$: $3^{2x} = 3^3$, tedy $x = \frac{3}{2}$.
    \item $5^{x^2 - 3x} = \frac{1}{25}$. Tedy $5^{x^2 - 3x} = 5^{-2}$, $x^2 - 3x + 2 = 0$, $x\in\{1, 2\}$.
    \item $2^x + 2^{x + 2} = 40$. Vytkneme $2^x$: $2^x(1 + 4) = 40$, tedy $2^x = 8$ a $x = 3$.
    \item $4^x - 3\cdot 2^x - 4 = 0$. Substituce $t = 2^x > 0$: $t^2 - 3t - 4 = 0$, tedy $t = 4$ nebo $t = -1$. Záporné $t$ zahodíme, protože $2^x$ je vždy kladné. Z $2^x = 4$ plyne $x = 2$.
    \item $3^x = 7$. Zlogaritmujeme: $x\log 3 = \log 7$, tedy $x = \frac{\log 7}{\log 3} = \log_3 7\approx 1{,}77$.
  \end{enumerate}
\end{example}

\begin{problem}
  Vyřešte v $\mathbb{R}$:
  \begin{enumerate}[label=(\alph*)]
    \item ${\left(\frac{1}{2}\right)}^{x} = 16$,
    \item $3^{2x - 1} = \sqrt{3}$,
    \item $9^x - 4\cdot 3^x + 3 = 0$,
    \item $5^{x + 1} - 5^{x - 1} = 120$,
    \item $2^{x} = 3^{x - 1}$.
  \end{enumerate}
\end{problem}

\begin{solution}
  \begin{enumerate}[label=(\alph*)]
    \item $2^{-x} = 2^4$, $x = -4$.
    \item $2x - 1 = \frac{1}{2}$, $x = \frac{3}{4}$.
    \item $t = 3^x$: $t^2 - 4t + 3 = 0$, $t\in\{1, 3\}$, tedy $x\in\{0, 1\}$.
    \item $5^{x - 1}(25 - 1) = 120$, $5^{x - 1} = 5$, $x = 2$.
    \item Zlogaritmujeme: $x\log 2 = (x - 1)\log 3$, tedy $x(\log 3 - \log 2) = \log 3$ a $x = \frac{\log 3}{\log 3 - \log 2} = \frac{\log 3}{\log 1{,}5}\approx 2{,}71$.
  \end{enumerate}
\end{solution}

\section{Logaritmické rovnice}

\begin{definitionbox}
  \textbf{Logaritmická rovnice} je rovnice, ve které se neznámá vyskytuje v argumentu logaritmu (případně v základu). Postup:
  \begin{enumerate}
    \item Určíme \textbf{podmínky}: argumenty logaritmů musí být kladné (a základy kladné a různé od $1$).
    \item Pomocí pravidel převedeme rovnici na tvar $\log_a f(x) = \log_a g(x)$, případně $\log_a f(x) = c$.
    \item Využijeme prostoty: $\log_a f(x) = \log_a g(x)\Leftrightarrow f(x) = g(x)$, případně z definice $\log_a f(x) = c\Leftrightarrow f(x) = a^c$.
    \item Ověříme podmínky (nebo provedeme zkoušku).
  \end{enumerate}
\end{definitionbox}

Pozor: úprava $\log_a f(x) + \log_a g(x) = \log_a(f(x)g(x))$ \textbf{rozšiřuje definiční obor}. Výraz vlevo vyžaduje $f(x) > 0$ \textbf{a} $g(x) > 0$, výraz vpravo jen $f(x)g(x) > 0$, což je splněno i když jsou obě záporné. Proto je nutné kontrolovat podmínky původní rovnice.

\begin{example}
  Vyřešme $\log_2 x + \log_2(x - 2) = 3$. Podmínky: $x > 0$ a $x > 2$, tedy $x > 2$.
  \begin{equation*}
    \log_2\left(x(x - 2)\right) = 3\quad\Leftrightarrow\quad x(x - 2) = 2^3\quad\Leftrightarrow\quad x^2 - 2x - 8 = 0\quad\Leftrightarrow\quad x\in\{4, -2\}.
  \end{equation*}
  Podmínku splňuje jen $x = 4$. $K = \{4\}$. Číslo $-2$ je falešný kořen, který vznikl přesně tím rozšířením definičního oboru: $(-2)\cdot(-4) = 8 > 0$, ale $\log_2(-2)$ neexistuje.
\end{example}

\begin{example}
  Vyřešme $\log^2 x - 3\log x + 2 = 0$ (kde $\log^2 x$ znamená ${(\log x)}^2$). Podmínka $x > 0$. Substituce $t = \log x$: $t^2 - 3t + 2 = 0$, tedy $t\in\{1, 2\}$ a $x\in\{10, 100\}$.
\end{example}

\begin{problem}
  Vyřešte v $\mathbb{R}$:
  \begin{enumerate}[label=(\alph*)]
    \item $\log_3(2x - 1) = 2$,
    \item $\log(x + 1) - \log(x - 1) = \log 2$,
    \item $\log_3(x^2 - 8) = 0$,
    \item $\ln x = 2\ln 3 - \ln 2$,
    \item $\log_2^2 x - \log_2 x^3 + 2 = 0$.
  \end{enumerate}
\end{problem}

\begin{solution}
  \begin{enumerate}[label=(\alph*)]
    \item $x > \frac{1}{2}$, $2x - 1 = 9$, $x = 5$.
    \item $x > 1$, $\log\frac{x + 1}{x - 1} = \log 2$, tedy $x + 1 = 2x - 2$ a $x = 3$.
    \item $x^2 > 8$, $x^2 - 8 = 3^0 = 1$, $x = \pm 3$. Obě splňují podmínku, $K = \{-3, 3\}$.
    \item $\ln x = \ln\frac{9}{2}$, $x = \frac{9}{2}$.
    \item $\log_2^2 x - 3\log_2 x + 2 = 0$, $t = \log_2 x\in\{1, 2\}$, $x\in\{2, 4\}$.
  \end{enumerate}
\end{solution}

\section{Exponenciální a logaritmické nerovnice}

Pro nerovnice je klíčová monotonie. Rostoucí funkce zachovává nerovnost, klesající ji otáčí.

\begin{definitionbox}
  Pro $a > 1$:
  \begin{equation*}
    a^{f(x)} < a^{g(x)}\Leftrightarrow f(x) < g(x),\qquad \log_a f(x) < \log_a g(x)\Leftrightarrow 0 < f(x) < g(x).
  \end{equation*}
  Pro $0 < a < 1$ se znak nerovnosti \textbf{otočí}:
  \begin{equation*}
    a^{f(x)} < a^{g(x)}\Leftrightarrow f(x) > g(x),\qquad \log_a f(x) < \log_a g(x)\Leftrightarrow f(x) > g(x) > 0.
  \end{equation*}
\end{definitionbox}

\begin{example}
  \begin{enumerate}[label=(\alph*)]
    \item $3^{x - 1} < 9$. Základ $3 > 1$: $x - 1 < 2$, tedy $K = (-\infty, 3)$.
    \item ${\left(\frac{1}{2}\right)}^x\geq 4$. Základ menší než $1$, znak se otočí: ${\left(\frac{1}{2}\right)}^x\geq{\left(\frac{1}{2}\right)}^{-2}\Leftrightarrow x\leq -2$. $K = (-\infty, -2]$.
    \item $\log_2(x + 3) < 2$. Podmínka $x > -3$. $x + 3 < 4$, tedy $x < 1$. $K = (-3, 1)$.
    \item $\log_{\frac{1}{3}}(x - 1)\geq -2$. Podmínka $x > 1$. Základ menší než $1$: $x - 1\leq{\left(\frac{1}{3}\right)}^{-2} = 9$, tedy $x\leq 10$. $K = (1, 10]$.
  \end{enumerate}
\end{example}

\section{Využití v praxi}

Exponenciální funkce popisují všechno, co roste nebo klesá úměrně tomu, kolik toho právě je. Logaritmy se pak používají všude, kde potřebujeme měřit veličiny, které se pohybují v obrovském rozsahu.

\begin{example}[ (rozpad)]
  Uhlík $^{14}$C má poločas rozpadu přibližně $5\,730$ let. To znamená, že za $5\,730$ let se rozpadne polovina atomů. Množství po čase $t$ je
  \begin{equation*}
    N(t) = N_0\cdot{\left(\frac{1}{2}\right)}^{\frac{t}{5730}}.
  \end{equation*}
  Archeologové najdou kost, ve které zbyla $\frac{1}{8}$ původního množství uhlíku $^{14}$C. Jak je stará? ${\left(\frac{1}{2}\right)}^{\frac{t}{5730}} = \frac{1}{8} = {\left(\frac{1}{2}\right)}^3$, tedy $t = 3\cdot 5730 = 17\,190$ let.
\end{example}

\begin{example}[ (růst)]
  Bakterie se dělí každých $20$ minut. Kolikrát se jejich počet zvětší za $2$ hodiny? Za $2$ hodiny proběhne $6$ dělení, takže počet se zvětší $2^6 = 64$krát. Za den by to bylo $2^{72}\approx 4{,}7\cdot 10^{21}$krát. Naštěstí jim dřív dojde jídlo.
\end{example}

\begin{note}
  Logaritmické stupnice jsou kolem nás:
  \begin{itemize}
    \item \textbf{pH}: $\mathrm{pH} = -\log[\mathrm{H}^+]$. Roztok s pH $3$ je desetkrát kyselejší než roztok s pH $4$.
    \item \textbf{Decibely}: zvýšení hlasitosti o $10\ \si{dB}$ znamená desetkrát větší intenzitu zvuku. Koncert ($110\ \si{dB}$) je tedy stomilionkrát intenzivnější než šeptání ($30\ \si{dB}$).
    \item \textbf{Richterova stupnice}: zemětřesení o síle $7$ uvolní zhruba $32$krát víc energie než zemětřesení o síle $6$.
  \end{itemize}
\end{note}

\begin{problem}
  Kolik let musí ležet $10\,000$ Kč v bance s úrokem $5\,\%$ ročně (úrok se připisuje jednou ročně a nechává se v bance), aby se částka zdvojnásobila?
\end{problem}

\begin{solution}
  $10000\cdot 1{,}05^n = 20000$, tedy $1{,}05^n = 2$ a $n = \frac{\log 2}{\log 1{,}05}\approx 14{,}2$. Částka se zdvojnásobí během 15. roku. Podrobně se úročením zabývá kapitola o finanční matematice.
\end{solution}

\section{Příběh --- Muž s černým kohoutem}

Logaritmy vymyslel skotský šlechtic John Napier a v roce 1614 je zveřejnil v knize s krásným názvem \textit{Mirifici Logarithmorum Canonis Descriptio}, tedy `Popis podivuhodného pravidla logaritmů'. Napier strávil přes dvacet let ručním počítáním tabulek. Astronomové je okamžitě začali používat, protože násobení dlouhých čísel se změnilo na sčítání a vyhledání v tabulce. Johannes Kepler díky nim dokončil výpočty drah planet. Odtud pochází i citát na začátku kapitoly. Pak přišla logaritmická pravítka, na kterých počítali inženýři až do 70. let 20. století. Na nich se mimo jiné počítaly i výpočty pro let na Měsíc.

Napier ale nebyl jen matematik. Mezi sousedy měl pověst čaroděje. Chodil prý v dlouhém černém plášti, mumlal si pro sebe a všude s sebou nosil černého kohouta, o kterém se říkalo, že je to jeho démonický pomocník. Nejznámější historka vypráví, jak Napier zjistil, kdo z jeho služebnictva krade. Zavřel kohouta do temné místnosti a oznámil, že kohout pozná zloděje. Každý sluha měl jít dovnitř a pohladit ho. Když se zloděj dotkne kohouta, ten zakokrhá. Kohout nezakokrhal ani jednou. Napier ale předtím kohouta natřel sazemi a pak si prohlédl ruce služebnictva. Jeden z nich měl ruce čisté, protože se kohouta bál dotknout.\footnote{Je to spíš legenda než doložený fakt, ale je to tak dobrá ukázka logického uvažování, že ji vyprávíme rádi. Napier mimochodem napsal i rozsáhlý výklad Apokalypsy, ve kterém spočítal, že konec světa nastane mezi lety 1688 a 1700. Tady se trochu přepočítal.}

\section{Shrnutí a otázky k ústní zkoušce}

\begin{problem}
  Odpovězte vlastními slovy:
  \begin{enumerate}
    \item Co je exponenciální funkce? Proč požadujeme $a > 0$ a $a\neq 1$? Popište její vlastnosti pro $a > 1$ a pro $0 < a < 1$.
    \item Co je Eulerovo číslo a jak ho můžeme získat?
    \item Jak je definován logaritmus? Proč neexistuje logaritmus záporného čísla?
    \item Jaký je vztah mezi exponenciální a logaritmickou funkcí? Jak to vypadá na grafu?
    \item Vyslovte a dokažte pravidla pro počítání s logaritmy, včetně převodu základu.
    \item Jaké znáte metody řešení exponenciálních rovnic?
    \item Jak se řeší logaritmické rovnice? Proč je nutné kontrolovat podmínky?
    \item Jak se řeší exponenciální a logaritmické nerovnice? Kdy se otáčí znak nerovnosti?
    \item Uveďte příklady využití exponenciálních funkcí a logaritmů v praxi.
  \end{enumerate}
\end{problem}

\koneckapitoly

\chapter{Goniometrie}
\label{kap:goniometrie}
\mathquote{Mathematics compares the most diverse phenomena and discovers the secret analogies that unite them.}{Joseph Fourier}

Goniometrie je část matematiky, která se zabývá úhly a funkcemi, které na nich závisí. Začíná nevinně u pravoúhlého trojúhelníku, ale končí u vln, kmitání, střídavého proudu, hudby a komprese obrázků v telefonu. Goniometrické funkce jsou totiž přesně ty funkce, které popisují všechno, co se opakuje. A to je ve fyzice skoro všechno.

\begin{note}
  Slovo \textit{sinus} má mimochodem jednu z nejzvláštnějších etymologií v matematice. Indičtí matematici mu říkali \textit{jyā}, tedy `tětiva luku'. Arabové slovo převzali jako \textit{jība}, ale v arabském písmu se nepíší samohlásky, takže pozdější překladatelé ho přečetli jako \textit{jaib}, což znamená `záhyb' nebo `ňadra'. Do latiny ho pak Gerard z Cremony v 12. století věrně přeložil jako \textit{sinus}, tedy také `záhyb, ohyb, ňadra'. Takže když počítáte sinus, počítáte vlastně překlep starý přes osm set let.
\end{note}

\section{Úhly a jejich míra}\label{sec:mira-uhlu}

\begin{definitionbox}
  \textbf{Stupňová míra}: Plný úhel má $360^\circ$, $1^\circ = 60'$ (minut), $1' = 60''$ (vteřin).

  \textbf{Oblouková míra}: Velikost úhlu v \textbf{radiánech} je délka oblouku, který úhel vytne na jednotkové kružnici (kružnici o poloměru $1$) se středem ve vrcholu úhlu. Plný úhel má tedy $2\pi$ radiánů (obvod jednotkové kružnice).

  \textbf{Převod}:
  \begin{equation*}
    180^\circ = \pi\ \text{rad},\qquad\alpha_{\text{rad}} = \frac{\pi}{180}\cdot\alpha^\circ,\qquad\alpha^\circ = \frac{180}{\pi}\cdot\alpha_{\text{rad}}.
  \end{equation*}
\end{definitionbox}

Proč radiány? Stupně jsou úplně libovolné, číslo $360$ si vymysleli Babyloňané, protože se dobře dělí a je blízko počtu dní v roce. Radiány jsou naopak `přirozené'. Vzorce v analýze (třeba pro derivaci sinu) fungují hezky jen v radiánech. Pokud v radiánech uvidíte úhel bez jednotky, je to v pořádku, radián se obvykle nepíše.

Co je tedy jeden radián? Představte si, že poloměr kružnice vezmete jako kus provázku a přiložíte ho na obvod kružnice. Úhel, který takový oblouk vytne, je právě $1$ radián, zhruba $57{,}3^\circ$. Na celý obvod se provázek vejde $2\pi\approx 6{,}28$krát, proto má plný úhel $2\pi$ radiánů.

\begin{center}
  \begin{tikzpicture}[scale=1.4]
    \begin{scope}
      \fill[colorofwisdom!15] (0,0) -- (1,0) arc (0:57.2958:1) -- cycle;
      \draw[thick] (0,0) circle (1);
      \draw[line width=2.2pt, colorofwisdom] (1,0) arc (0:57.2958:1);
      \draw[thick] (0,0) -- (1,0) node[midway, below] {\small $r$};
      \draw[thick] (0,0) -- (57.2958:1) node[midway, above left] {\small $r$};
      \draw (0.3,0) arc (0:57.2958:0.3);
      \node at (28.65:0.5) {\scriptsize $1$ rad};
      \node[colorofwisdom, right] at (28.65:1.05) {\scriptsize oblouk délky $r$};
      \fill (0,0) circle (1pt);
    \end{scope}
    \begin{scope}[shift={(4.9,0)}, scale=1.25]
      \draw[gray] (-1.12,0) -- (1.12,0);
      \draw[gray] (0,-1.12) -- (0,1.12);
      \draw[thick] (0,0) circle (1);
      \foreach \a/\l in {0/0, 30/\frac{\pi}{6}, 45/\frac{\pi}{4}, 60/\frac{\pi}{3}, 90/\frac{\pi}{2}, 120/\frac{2\pi}{3}, 135/\frac{3\pi}{4}, 150/\frac{5\pi}{6}, 180/\pi, 210/\frac{7\pi}{6}, 225/\frac{5\pi}{4}, 240/\frac{4\pi}{3}, 270/\frac{3\pi}{2}, 300/\frac{5\pi}{3}, 315/\frac{7\pi}{4}, 330/\frac{11\pi}{6}} {
        \draw[gray!60] (0,0) -- (\a:1);
        \fill[colorofwisdom] (\a:1) circle (1.1pt);
        \node at (\a:1.27) {\tiny $\l$};
      }
    \end{scope}
  \end{tikzpicture}
\end{center}

\begin{center}
  \renewcommand{\arraystretch}{1.4}
  \begin{tabular}{|c|c|c|c|c|c|c|c|c|}
    \hline
    stupně & $0^\circ$ & $30^\circ$ & $45^\circ$ & $60^\circ$ & $90^\circ$ & $180^\circ$ & $270^\circ$ & $360^\circ$ \\
    \hline
    radiány & $0$ & $\frac{\pi}{6}$ & $\frac{\pi}{4}$ & $\frac{\pi}{3}$ & $\frac{\pi}{2}$ & $\pi$ & $\frac{3\pi}{2}$ & $2\pi$ \\
    \hline
  \end{tabular}
\end{center}

\begin{definitionbox}
  \textbf{Orientovaný úhel}: Úhel měříme od kladné poloosy $x$. Proti směru hodinových ručiček je úhel \textbf{kladný}, po směru \textbf{záporný}. Úhly, které se liší o celistvý násobek $2\pi$, mají stejné rameno. Každý orientovaný úhel lze zapsat jako $x = x_0 + 2k\pi$, kde $x_0\in[0, 2\pi)$ je \textbf{základní velikost} a $k\in\mathbb{Z}$.
\end{definitionbox}

\begin{problem}
  \begin{enumerate}[label=(\alph*)]
    \item Převeďte na radiány: $15^\circ$, $120^\circ$, $135^\circ$, $210^\circ$, $-300^\circ$.
    \item Převeďte na stupně: $\frac{5\pi}{6}$, $\frac{7\pi}{4}$, $1$ rad.
    \item Určete základní velikost úhlů $\frac{17\pi}{3}$ a $-\frac{5\pi}{4}$.
  \end{enumerate}
\end{problem}

\begin{solution}
  \begin{enumerate}[label=(\alph*)]
    \item $\frac{\pi}{12}$, $\frac{2\pi}{3}$, $\frac{3\pi}{4}$, $\frac{7\pi}{6}$, $-\frac{5\pi}{3}$.
    \item $150^\circ$, $315^\circ$, $\frac{180^\circ}{\pi}\approx 57{,}3^\circ$.
    \item $\frac{17\pi}{3} = \frac{5\pi}{3} + 4\pi$, základní velikost $\frac{5\pi}{3}$. $-\frac{5\pi}{4} + 2\pi = \frac{3\pi}{4}$.
  \end{enumerate}
\end{solution}

\section{Goniometrické funkce v pravoúhlém trojúhelníku}\label{sec:gon-pravouhly}

\begin{definitionbox}
  V pravoúhlém trojúhelníku s ostrým úhlem $\alpha$, odvěsnou $a$ \textbf{protilehlou} k úhlu $\alpha$, odvěsnou $b$ \textbf{přilehlou} a přeponou $c$ definujeme
  \begin{equation*}
    \sin\alpha = \frac{a}{c},\qquad\cos\alpha = \frac{b}{c},\qquad\tg\alpha = \frac{a}{b},\qquad\cotg\alpha = \frac{b}{a}.
  \end{equation*}
\end{definitionbox}

\begin{center}
  \begin{tikzpicture}[scale=1.1]
    \draw[thick] (0,0) -- (4,0) -- (4,2.5) -- cycle;
    \draw (3.75,0) -- (3.75,0.25) -- (4,0.25);
    \draw[colorofwisdom, thick] (0.8,0) arc (0:32:0.8);
    \node[colorofwisdom] at (1.05,0.25) {$\alpha$};
    \node[below] at (2,0) {$b$ (přilehlá)};
    \node[right] at (4,1.25) {$a$ (protilehlá)};
    \node[above left] at (2,1.25) {$c$ (přepona)};
  \end{tikzpicture}
\end{center}

Proč tato definice dává smysl? Všechny pravoúhlé trojúhelníky se stejným úhlem $\alpha$ jsou si podobné, takže poměry jejich stran jsou stejné. Hodnota $\sin\alpha$ tedy závisí jen na úhlu, ne na velikosti trojúhelníku. Podobnosti se podrobně věnuje kapitola o planimetrii.

\begin{example}\label{ex:tabulka-hodnot}
  Odvoďme hodnoty pro $30^\circ$, $45^\circ$ a $60^\circ$.

  \textbf{$45^\circ$}: Vezmeme rovnoramenný pravoúhlý trojúhelník s odvěsnami $1$. Přepona je podle Pythagorovy věty $\sqrt{2}$, takže $\sin 45^\circ = \cos 45^\circ = \frac{1}{\sqrt{2}} = \frac{\sqrt{2}}{2}$ a $\tg 45^\circ = 1$.

  \textbf{$30^\circ$ a $60^\circ$}: Vezmeme rovnostranný trojúhelník o straně $2$ a rozpůlíme ho výškou. Dostaneme pravoúhlý trojúhelník s přeponou $2$, kratší odvěsnou $1$ a delší odvěsnou $\sqrt{4 - 1} = \sqrt{3}$. Úhly jsou $30^\circ$ a $60^\circ$. Tedy
  \begin{equation*}
    \sin 30^\circ = \frac{1}{2},\quad\cos 30^\circ = \frac{\sqrt{3}}{2},\quad\sin 60^\circ = \frac{\sqrt{3}}{2},\quad\cos 60^\circ = \frac{1}{2}.
  \end{equation*}
  \begin{center}
    \begin{tikzpicture}[scale=1.3]
      \begin{scope}
        \fill[colorofwisdom!15] (0,0) -- (1.8,0) -- (1.8,1.8) -- cycle;
        \draw[thick] (0,0) -- (1.8,0) node[midway, below] {\small $1$} -- (1.8,1.8) node[midway, right] {\small $1$} -- cycle;
        \node[above left] at (0.9,0.9) {\small $\sqrt{2}$};
        \draw (1.6,0) -- (1.6,0.2) -- (1.8,0.2);
        \draw (0.5,0) arc (0:45:0.5);
        \node at (22.5:0.75) {\scriptsize $45^\circ$};
      \end{scope}
      \begin{scope}[xshift=3.2cm]
        \fill[colorofwisdom!15] (0,0) -- (1.5,0) -- (1.5,2.598) -- cycle;
        \draw[thick] (0,0) -- (3,0) -- (1.5,2.598) -- cycle;
        \draw[thick, colorofwisdom] (1.5,0) -- (1.5,2.598);
        \draw (1.3,0) -- (1.3,0.2) -- (1.5,0.2);
        \node[below] at (0.75,0) {\small $1$};
        \node[below] at (2.25,0) {\small $1$};
        \node[above left] at (0.75,1.3) {\small $2$};
        \node[above right] at (2.25,1.3) {\small $2$};
        \node[colorofwisdom, right] at (1.5,1.0) {\small $\sqrt{3}$};
        \draw (0.45,0) arc (0:60:0.45);
        \node at (30:0.7) {\scriptsize $60^\circ$};
        \draw ($(1.5,2.598)+(240:0.55)$) arc (240:270:0.55);
        \node at ($(1.5,2.598)+(255:0.82)$) {\scriptsize $30^\circ$};
      \end{scope}
    \end{tikzpicture}
  \end{center}
\end{example}

\begin{definitionbox}
  \textbf{Tabulka významných hodnot}:
  \begin{center}
    \renewcommand{\arraystretch}{1.6}
    \begin{tabular}{|c|c|c|c|c|c|}
      \hline
      $x$ & $0$ & $\frac{\pi}{6}$ & $\frac{\pi}{4}$ & $\frac{\pi}{3}$ & $\frac{\pi}{2}$ \\
      \hline
      $\sin x$ & $0$ & $\frac{1}{2}$ & $\frac{\sqrt{2}}{2}$ & $\frac{\sqrt{3}}{2}$ & $1$ \\
      $\cos x$ & $1$ & $\frac{\sqrt{3}}{2}$ & $\frac{\sqrt{2}}{2}$ & $\frac{1}{2}$ & $0$ \\
      $\tg x$ & $0$ & $\frac{\sqrt{3}}{3}$ & $1$ & $\sqrt{3}$ & --- \\
      $\cotg x$ & --- & $\sqrt{3}$ & $1$ & $\frac{\sqrt{3}}{3}$ & $0$ \\
      \hline
    \end{tabular}
  \end{center}
  Mnemotechnická pomůcka pro sinus: $\frac{\sqrt{0}}{2}, \frac{\sqrt{1}}{2}, \frac{\sqrt{2}}{2}, \frac{\sqrt{3}}{2}, \frac{\sqrt{4}}{2}$. Kosinus je to samé pozpátku.
\end{definitionbox}

\begin{problem}
  \begin{enumerate}[label=(\alph*)]
    \item Žebřík dlouhý $5\ \si{m}$ je opřený o zeď a svírá se zemí úhel $60^\circ$. Jak vysoko na zdi končí?
    \item Ze vzdálenosti $50\ \si{m}$ od paty věže vidíme její vrchol pod úhlem $35^\circ$. Jak je věž vysoká?
  \end{enumerate}
\end{problem}

\begin{solution}
  \begin{enumerate}[label=(\alph*)]
    \item $h = 5\sin 60^\circ = \frac{5\sqrt{3}}{2}\approx 4{,}33\ \si{m}$.
    \item $v = 50\tg 35^\circ\approx 50\cdot 0{,}700\approx 35\ \si{m}$ (plus výška očí pozorovatele, pokud nekouká ze země).
  \end{enumerate}
\end{solution}

\section{Goniometrické funkce obecného úhlu}\label{sec:gon-obecny-uhel}

V trojúhelníku má smysl jen úhel mezi $0$ a $\frac{\pi}{2}$. Pro obecný úhel použijeme jednotkovou kružnici.

\begin{definitionbox}
  Nechť $x\in\mathbb{R}$. Na jednotkové kružnici se středem v počátku najdeme bod $M$, do kterého se dostaneme z bodu $[1, 0]$ otočením o orientovaný úhel $x$. Pak definujeme
  \begin{equation*}
    \cos x = \text{$x$-ová souřadnice bodu } M,\qquad\sin x = \text{$y$-ová souřadnice bodu } M,
  \end{equation*}
  tedy $M = [\cos x, \sin x]$. Dále
  \begin{equation*}
    \tg x = \frac{\sin x}{\cos x}\quad(\cos x\neq 0),\qquad\cotg x = \frac{\cos x}{\sin x}\quad(\sin x\neq 0).
  \end{equation*}
\end{definitionbox}

\begin{center}
  \begin{tikzpicture}[scale=2.2]
    \draw[->] (-1.3,0) -- (1.4,0) node[right] {\small $x$};
    \draw[->] (0,-1.2) -- (0,1.3) node[above] {\small $y$};
    \draw[thick] (0,0) circle (1);
    \coordinate (M) at (40:1);
    \draw[thick, colorofwisdom] (0,0) -- (M);
    \filldraw[colorofwisdom] (M) circle (0.8pt) node[above right] {$M = [\cos x, \sin x]$};
    \draw[dashed] (M) -- ({cos(40)},0);
    \draw[dashed] (M) -- (0,{sin(40)});
    \draw[line width=2pt, colorofwisdom!60] (0,0) -- ({cos(40)},0);
    \draw[line width=2pt, black!50] ({cos(40)},0) -- (M);
    \node[below] at ({cos(40)/2},0) {\small $\cos x$};
    \node[right] at ({cos(40)},{sin(40)/2}) {\small $\sin x$};
    \draw[->, colorofwisdom] (0.3,0) arc (0:40:0.3);
    \node[colorofwisdom] at (20:0.42) {\small $x$};
    \node[below right] at (1,0) {\small $1$};
    \node[above left] at (0,1) {\small $1$};
    \node at (0.75,0.95) {\small I};
    \node at (-0.75,0.95) {\small II};
    \node at (-0.75,-0.95) {\small III};
    \node at (0.75,-0.95) {\small IV};
  \end{tikzpicture}
\end{center}

Pro ostré úhly se nová definice shoduje se starou. Bod $M$, jeho průmět na osu $x$ a počátek tvoří pravoúhlý trojúhelník s přeponou $1$, takže $\sin x = \frac{\text{protilehlá}}{1}$.

I tangens má na jednotkové kružnici svůj obrázek. Představte si \textbf{tečnu} ke kružnici v bodě $[1, 0]$, tedy svislou přímku $x = 1$. Prodloužíme rameno úhlu, dokud tečnu neprotne v bodě $T$. Trojúhelník s vrcholy $[0, 0]$, $[1, 0]$ a $T$ je podobný trojúhelníku s vrcholy $[0, 0]$, $[\cos x, 0]$ a $M$, takže výška bodu $T$ je $\frac{\sin x}{\cos x}\cdot 1 = \tg x$. Tangens je tedy \textbf{výška bodu $T$ na tečně}. Teď je také vidět, proč pro $x\to\frac{\pi}{2}$ tangens utíká do nekonečna: rameno je skoro svislé a s tečnou se protne vysoko nahoře, pro $x = \frac{\pi}{2}$ jsou rovnoběžné a tangens neexistuje.

\begin{center}
  \begin{tikzpicture}[scale=2.2]
    \draw[->] (-1.2,0) -- (1.6,0) node[right] {\small $x$};
    \draw[->] (0,-1.2) -- (0,1.3) node[above] {\small $y$};
    \draw[thick] (0,0) circle (1);
    \draw[dashed, gray] (1,-1.15) -- (1,1.25) node[above] {\small $x = 1$};
    \draw[thick, colorofwisdom] (0,0) -- (1,0.70021);
    \draw[dashed] (35:1) -- (0.81915,0);
    \draw[line width=2pt, colorofwisdom!60] (0,0) -- (0.81915,0);
    \draw[line width=2pt, black!50] (0.81915,0) -- (35:1);
    \draw[line width=2pt, colorofwisdom] (1,0) -- (1,0.70021);
    \fill (35:1) circle (0.8pt) node[above left] {\small $M$};
    \fill[colorofwisdom] (1,0.70021) circle (0.8pt) node[right] {\small $T = [1, \tg x]$};
    \node[right, colorofwisdom] at (1,0.33) {\small $\tg x$};
    \node[below] at (0.41,0) {\scriptsize $\cos x$};
    \node[left] at (0.81915,0.29) {\scriptsize $\sin x$};
    \draw[->, colorofwisdom] (0.3,0) arc (0:35:0.3);
    \node[colorofwisdom] at (17:0.42) {\small $x$};
  \end{tikzpicture}
\end{center}

\begin{definitionbox}
  \textbf{Znaménka v kvadrantech}:
  \begin{center}
    \renewcommand{\arraystretch}{1.3}
    \begin{tabular}{|c|c|c|c|c|}
      \hline
      kvadrant & I & II & III & IV \\
      \hline
      $\sin x$ & $+$ & $+$ & $-$ & $-$ \\
      $\cos x$ & $+$ & $-$ & $-$ & $+$ \\
      $\tg x$, $\cotg x$ & $+$ & $-$ & $+$ & $-$ \\
      \hline
    \end{tabular}
  \end{center}

  \textbf{Převod do prvního kvadrantu} (pro $x\in\left(0, \frac{\pi}{2}\right)$):
  \begin{align*}
    \sin(\pi - x) &= \sin x, & \cos(\pi - x) &= -\cos x,\\
    \sin(\pi + x) &= -\sin x, & \cos(\pi + x) &= -\cos x,\\
    \sin(2\pi - x) &= -\sin x, & \cos(2\pi - x) &= \cos x.
  \end{align*}
  Všechny plynou ze souměrnosti jednotkové kružnice.
\end{definitionbox}

Tyto vzorce není potřeba se učit nazpaměť, stačí si je nakreslit. Úhly $x$, $\pi - x$, $\pi + x$ a $2\pi - x$ leží v rozích obdélníku vepsaného do kružnice. Úhel $\pi - x$ je zrcadlem úhlu $x$ podle osy $y$, takže má stejnou výšku (sinus) a opačnou vodorovnou souřadnici (kosinus). Úhel $\pi + x$ je otočením o půl kruhu, oběma souřadnicím se změní znaménko. Úhel $2\pi - x$, tedy $-x$, je zrcadlem podle osy $x$.

\begin{center}
  \begin{tikzpicture}[scale=2]
    \draw[->] (-1.3,0) -- (1.4,0) node[right] {\small $x$};
    \draw[->] (0,-1.25) -- (0,1.3) node[above] {\small $y$};
    \draw[thick] (0,0) circle (1);
    \draw[dashed, colorofwisdom] (30:1) -- (150:1) -- (210:1) -- (330:1) -- cycle;
    \foreach \a in {30,150,210,330} \draw[gray] (0,0) -- (\a:1);
    \fill[colorofwisdom] (30:1) circle (1pt) node[above right] {\small $x$};
    \fill (150:1) circle (1pt) node[above left] {\small $\pi - x$};
    \fill (210:1) circle (1pt) node[below left] {\small $\pi + x$};
    \fill (330:1) circle (1pt) node[below right] {\small $2\pi - x$};
    \node[right] at (0,0.58) {\scriptsize $\sin x$};
    \node[right] at (0,-0.58) {\scriptsize $-\sin x$};
    \node[below] at (0.6,0) {\scriptsize $\cos x$};
    \node[below] at (-0.6,0) {\scriptsize $-\cos x$};
    \draw[colorofwisdom, ->] (0.25,0) arc (0:30:0.25);
  \end{tikzpicture}
\end{center}

\begin{example}
  Spočítejme:
  \begin{align*}
    \sin\frac{5\pi}{6} &= \sin\left(\pi - \frac{\pi}{6}\right) = \sin\frac{\pi}{6} = \frac{1}{2},\\
    \cos\frac{4\pi}{3} &= \cos\left(\pi + \frac{\pi}{3}\right) = -\cos\frac{\pi}{3} = -\frac{1}{2},\\
    \tg\frac{7\pi}{4} &= \frac{\sin\frac{7\pi}{4}}{\cos\frac{7\pi}{4}} = \frac{-\frac{\sqrt{2}}{2}}{\frac{\sqrt{2}}{2}} = -1,\\
    \sin\frac{25\pi}{6} &= \sin\left(\frac{\pi}{6} + 4\pi\right) = \sin\frac{\pi}{6} = \frac{1}{2}.
  \end{align*}
\end{example}

\begin{problem}
  Spočítejte bez kalkulačky: $\cos\frac{5\pi}{6}$, $\sin\frac{7\pi}{6}$, $\tg\frac{2\pi}{3}$, $\cos\left(-\frac{\pi}{4}\right)$, $\sin\frac{3\pi}{2}$, $\cos 3\pi$, $\cotg\frac{5\pi}{3}$.
\end{problem}

\begin{solution}
  $-\frac{\sqrt{3}}{2}$, $-\frac{1}{2}$, $-\sqrt{3}$, $\frac{\sqrt{2}}{2}$, $-1$, $-1$, $-\frac{\sqrt{3}}{3}$.
\end{solution}

\section{Grafy a vlastnosti goniometrických funkcí}\label{sec:gon-grafy}

Pokud bod $M$ obíhá kolem jednotkové kružnice a my zapisujeme jeho výšku v závislosti na úhlu, dostaneme graf sinu. Vodorovná souřadnice dá graf kosinu.

Představte si, že úhel $x$ (délku oblouku) `odvineme' z kružnice na osu $x$ a do výšky nad ním naneseme výšku bodu $M$. Každý bod sinusoidy je tak jen výškou nějakého bodu kružnice.

\begin{center}
  \begin{tikzpicture}[scale=1.05]
    \draw[->] (-2.7,0) -- (6.9,0) node[right] {\small $x$};
    \draw[->] (0,-1.3) -- (0,1.45) node[above] {\small $y$};
    \draw[thick] (-1.5,0) circle (1);
    \draw (-1.5,-1.15) -- (-1.5,1.15);
    \foreach \t/\l in {1.5708/{\frac{\pi}{2}}, 3.1416/{\pi}, 4.7124/{\frac{3\pi}{2}}, 6.2832/{2\pi}} \draw (\t,0.06) -- (\t,-0.06) node[below] {\tiny $\l$};
    \foreach \a in {0,45,...,315} {
      \pgfmathsetmacro{\rr}{\a*pi/180}
      \draw[dotted, gray] ({-1.5 + cos(\a)},{sin(\a)}) -- (\rr,{sin(\a)});
      \fill[gray] ({-1.5 + cos(\a)},{sin(\a)}) circle (1.2pt);
      \fill[gray] (\rr,{sin(\a)}) circle (1.2pt);
    }
    \draw[thick, colorofwisdom, domain=0:6.2832, samples=120] plot (\x, {sin(\x r)});
    \draw[thick, colorofwisdom] (-1.5,0) -- (-1,0.866);
    \draw[->, colorofwisdom] (-1.2,0) arc (0:60:0.3);
    \node[colorofwisdom] at ($(-1.5,0)+(30:0.45)$) {\scriptsize $x$};
    \draw[line width=1.8pt, black!50] (-1,0) -- (-1,0.866);
    \draw[line width=1.8pt, black!50] (1.0472,0) -- (1.0472,0.866);
    \draw[dashed, colorofwisdom] (-1,0.866) -- (1.0472,0.866);
    \fill[colorofwisdom] (-1,0.866) circle (1.5pt) node[above] {\scriptsize $M$};
    \fill[colorofwisdom] (1.0472,0.866) circle (1.5pt);
    \draw[line width=1.8pt, colorofwisdom!50] (0,0) -- (1.0472,0) node[midway, below] {\tiny $x$};
    \node[colorofwisdom] at (3.9,1.05) {\small $y = \sin x$};
  \end{tikzpicture}
\end{center}

\begin{center}
  \begin{tikzpicture}[xscale=0.85, yscale=1.2]
    \draw[->] (-3.6,0) -- (7.2,0) node[right] {\small $x$};
    \draw[->] (0,-1.4) -- (0,1.5) node[above] {\small $y$};
    \draw[dashed, gray] (-3.5,1) -- (7,1);
    \draw[dashed, gray] (-3.5,-1) -- (7,-1);
    \node[left] at (0,1) {\tiny $1$};
    \node[left] at (0,-1) {\tiny $-1$};
    \foreach \t/\l in {-3.1416/{-\pi}, -1.5708/{-\frac{\pi}{2}}, 1.5708/{\frac{\pi}{2}}, 3.1416/{\pi}, 4.7124/{\frac{3\pi}{2}}, 6.2832/{2\pi}} \draw (\t,0.06) -- (\t,-0.06) node[below] {\tiny $\l$};
    \draw[thick, colorofwisdom, domain=-3.5:7, samples=150] plot (\x, {sin(\x r)});
    \draw[thick, domain=-3.5:7, samples=150] plot (\x, {cos(\x r)});
    \node[colorofwisdom] at (2.2,1.25) {\small $\sin x$};
    \node at (-0.7,1.25) {\small $\cos x$};
  \end{tikzpicture}
\end{center}

\begin{definitionbox}
  \textbf{Vlastnosti funkcí sinus a kosinus}:
  \begin{center}
    \renewcommand{\arraystretch}{1.4}
    \begin{tabular}{|l|l|l|}
      \hline
       & $\sin x$ & $\cos x$ \\
      \hline
      $D$, $H$ & $\mathbb{R}$, $[-1, 1]$ & $\mathbb{R}$, $[-1, 1]$ \\
      perioda & $2\pi$ & $2\pi$ \\
      parita & lichá & sudá \\
      nulové body & $k\pi$ & $\frac{\pi}{2} + k\pi$ \\
      maximum $1$ & $\frac{\pi}{2} + 2k\pi$ & $2k\pi$ \\
      minimum $-1$ & $\frac{3\pi}{2} + 2k\pi$ & $\pi + 2k\pi$ \\
      rostoucí na & $\left[-\frac{\pi}{2} + 2k\pi, \frac{\pi}{2} + 2k\pi\right]$ & $[-\pi + 2k\pi, 2k\pi]$ \\
      klesající na & $\left[\frac{\pi}{2} + 2k\pi, \frac{3\pi}{2} + 2k\pi\right]$ & $[2k\pi, \pi + 2k\pi]$ \\
      \hline
    \end{tabular}
  \end{center}
  ($k\in\mathbb{Z}$). Graf kosinu je graf sinu posunutý o $\frac{\pi}{2}$ doleva: $\cos x = \sin\left(x + \frac{\pi}{2}\right)$.
\end{definitionbox}

\begin{center}
  \begin{tikzpicture}[xscale=0.85, yscale=0.6]
    \draw[->] (-3.6,0) -- (5.2,0) node[right] {\small $x$};
    \draw[->] (0,-3.4) -- (0,3.6) node[above] {\small $y$};
    \foreach \t/\l in {-3.1416/{-\pi}, -1.5708/{-\frac{\pi}{2}}, 1.5708/{\frac{\pi}{2}}, 3.1416/{\pi}, 4.7124/{\frac{3\pi}{2}}} \draw (\t,0.1) -- (\t,-0.1) node[below] {\tiny $\l$};
    \foreach \a in {-1.5708, 1.5708, 4.7124} \draw[dashed, gray] (\a,-3.4) -- (\a,3.4);
    \draw[thick, colorofwisdom, domain=-1.27:1.27, samples=60] plot (\x, {tan(\x r)});
    \draw[thick, colorofwisdom, domain=1.87:4.41, samples=60] plot (\x, {tan(\x r)});
    \draw[thick, colorofwisdom, domain=-3.5:-1.87, samples=60] plot (\x, {tan(\x r)});
    \node[colorofwisdom] at (0.9,3.1) {\small $\tg x$};
  \end{tikzpicture}
\end{center}

\begin{center}
  \begin{tikzpicture}[xscale=0.85, yscale=0.6]
    \draw[->] (-3.6,0) -- (5.2,0) node[right] {\small $x$};
    \draw[->] (0,-3.4) -- (0,3.6) node[above] {\small $y$};
    \foreach \t/\l in {-3.1416/{-\pi}, -1.5708/{-\frac{\pi}{2}}, 1.5708/{\frac{\pi}{2}}, 3.1416/{\pi}, 4.7124/{\frac{3\pi}{2}}} \draw (\t,0.1) -- (\t,-0.1) node[below] {\tiny $\l$};
    \foreach \a in {-3.1416, 3.1416} \draw[dashed, gray] (\a,-3.4) -- (\a,3.4);
    \draw[thick, domain=0.29:2.85, samples=60] plot (\x, {cos(\x r)/sin(\x r)});
    \draw[thick, domain=-2.85:-0.29, samples=60] plot (\x, {cos(\x r)/sin(\x r)});
    \draw[thick, domain=3.43:5.15, samples=50] plot (\x, {cos(\x r)/sin(\x r)});
    \node at (1.15,2.6) {\small $\cotg x$};
  \end{tikzpicture}
\end{center}

Graf kotangentu je `zrcadlový' ke grafu tangentu: má asymptoty v nulových bodech sinu ($x = k\pi$), nulové body v $\frac{\pi}{2} + k\pi$ a na každém intervalu mezi asymptotami klesá. Opakuje se už po $\pi$, ne po $2\pi$, protože bod $M$ a bod k němu souměrný podle počátku (otočený o $\pi$) leží na stejné přímce procházející počátkem, a tedy dávají stejný poměr $\frac{\sin x}{\cos x}$.

\begin{definitionbox}
  \textbf{Vlastnosti funkce tangens}: $D = \mathbb{R}\setminus\left\{\frac{\pi}{2} + k\pi\right\}$, $H = \mathbb{R}$, perioda $\pi$, lichá, rostoucí na každém intervalu $\left(-\frac{\pi}{2} + k\pi, \frac{\pi}{2} + k\pi\right)$, nulové body $k\pi$, asymptoty $x = \frac{\pi}{2} + k\pi$. Kotangens má periodu $\pi$, je lichý a klesající na každém intervalu $(k\pi, \pi + k\pi)$.
\end{definitionbox}

\begin{example}
  Graf funkce $f(x) = 2\sin\left(x - \frac{\pi}{3}\right)$ dostaneme z grafu sinu posunem o $\frac{\pi}{3}$ doprava a roztažením dvakrát ve směru osy $y$ (transformace z kapitoly~\ref{kap:funkce}). Obor hodnot je $[-2, 2]$, perioda zůstává $2\pi$. Funkce $g(x) = \sin 2x$ má naopak periodu $\pi$, protože graf je stlačen dvakrát ve směru osy $x$. Obecně má $\sin(ax)$ periodu $\frac{2\pi}{|a|}$.
\end{example}

\begin{center}
  \begin{tikzpicture}[xscale=0.85, yscale=0.7]
    \draw[->] (-0.4,0) -- (7.2,0) node[right] {\small $x$};
    \draw[->] (0,-2.4) -- (0,2.6) node[above] {\small $y$};
    \foreach \t/\l in {1.0472/{\frac{\pi}{3}}, 3.1416/{\pi}, 6.2832/{2\pi}} \draw (\t,0.1) -- (\t,-0.1) node[below] {\tiny $\l$};
    \foreach \y in {-2,-1,1,2} \draw (0.08,\y) -- (-0.08,\y) node[left] {\tiny $\y$};
    \draw[thick, dashed, gray, domain=0:6.6, samples=100] plot (\x, {sin(\x r)});
    \draw[thick, colorofwisdom, domain=0:6.9, samples=100] plot (\x, {2*sin((\x - 1.0472) r)});
    \draw[thick, domain=0:6.6, samples=150] plot (\x, {sin(2*\x r)});
    \draw[thick, dashed, gray] (8,1.6) -- (8.6,1.6) node[right, black] {\scriptsize $\sin x$};
    \draw[thick, colorofwisdom] (8,0.8) -- (8.6,0.8) node[right, black] {\scriptsize $2\sin\left(x - \frac{\pi}{3}\right)$};
    \draw[thick] (8,0) -- (8.6,0) node[right] {\scriptsize $\sin 2x$};
  \end{tikzpicture}
\end{center}

\section{Goniometrické vzorce}\label{sec:gon-vzorce}

\begin{definitionbox}
  \textbf{Základní vztahy}: Pro všechna přípustná $x$ platí
  \begin{align*}
    \sin^2 x + \cos^2 x &= 1, & \tg x\cdot\cotg x &= 1,\\
    \sin(-x) &= -\sin x, & \cos(-x) &= \cos x,\\
    \sin\left(\tfrac{\pi}{2} - x\right) &= \cos x, & \cos\left(\tfrac{\pi}{2} - x\right) &= \sin x.
  \end{align*}

  \textbf{Součtové vzorce}:
  \begin{align*}
    \cos(x\pm y) &= \cos x\cos y\mp\sin x\sin y,\\
    \sin(x\pm y) &= \sin x\cos y\pm\cos x\sin y.
  \end{align*}

  \textbf{Vzorce pro dvojnásobný úhel}:
  \begin{equation*}
    \sin 2x = 2\sin x\cos x,\qquad\cos 2x = \cos^2 x - \sin^2 x = 1 - 2\sin^2 x = 2\cos^2 x - 1.
  \end{equation*}

  \textbf{Vzorce pro poloviční úhel}:
  \begin{equation*}
    \sin^2\frac{x}{2} = \frac{1 - \cos x}{2},\qquad\cos^2\frac{x}{2} = \frac{1 + \cos x}{2}.
  \end{equation*}
\end{definitionbox}

Vzorec $\sin^2 x + \cos^2 x = 1$ je Pythagorova věta na jednotkové kružnici. Bod $M = [\cos x, \sin x]$ leží na kružnici o poloměru $1$, takže součet čtverců jeho souřadnic je $1$. Značení $\sin^2 x$ znamená ${(\sin x)}^2$, ne $\sin(\sin x)$ ani $\sin(x^2)$.

\begin{proof}[(součtového vzorce pro $\cos(x - y)$)]
  Na jednotkové kružnici vezmeme body $A = [\cos x, \sin x]$ a $B = [\cos y, \sin y]$. Úhel mezi nimi je $x - y$. Druhou mocninu jejich vzdálenosti spočítáme dvakrát.

  Přímo ze souřadnic:
  \begin{equation*}
    |AB|^2 = {(\cos x - \cos y)}^2 + {(\sin x - \sin y)}^2 = 2 - 2(\cos x\cos y + \sin x\sin y),
  \end{equation*}
  kde jsme roznásobili a dvakrát použili $\sin^2 + \cos^2 = 1$.

  Nyní celý obrázek otočíme o úhel $-y$. Otočení zachovává vzdálenosti. Bod $B$ se dostane do $[1, 0]$ a bod $A$ do $[\cos(x - y), \sin(x - y)]$. Tedy
  \begin{equation*}
    |AB|^2 = {(\cos(x - y) - 1)}^2 + \sin^2(x - y) = 2 - 2\cos(x - y).
  \end{equation*}
  Porovnáním obou vyjádření dostaneme $\cos(x - y) = \cos x\cos y + \sin x\sin y$.
\end{proof}

\begin{center}
  \begin{tikzpicture}[scale=1.6]
    \begin{scope}
      \draw[->, gray] (-1.25,0) -- (1.35,0);
      \draw[->, gray] (0,-1.15) -- (0,1.3);
      \draw[thick] (0,0) circle (1);
      \draw (0,0) -- (110:1) (0,0) -- (40:1);
      \draw[line width=1.6pt, colorofwisdom] (110:1) -- (40:1);
      \fill (110:1) circle (1.2pt) node[above left] {\small $A$};
      \fill (40:1) circle (1.2pt) node[right] {\small $B$};
      \draw[->] (0.3,0) arc (0:40:0.3);
      \node at (20:0.45) {\scriptsize $y$};
      \draw[->] (0.6,0) arc (0:110:0.6);
      \node at (72:0.75) {\scriptsize $x$};
      \node at (0,-1.4) {\small původní poloha};
    \end{scope}
    \begin{scope}[xshift=3.4cm]
      \draw[->, gray] (-1.25,0) -- (1.35,0);
      \draw[->, gray] (0,-1.15) -- (0,1.3);
      \draw[thick] (0,0) circle (1);
      \draw (0,0) -- (70:1) (0,0) -- (1,0);
      \draw[line width=1.6pt, colorofwisdom] (70:1) -- (1,0);
      \fill (70:1) circle (1.2pt) node[above right] {\small $A'$};
      \fill (1,0) circle (1.2pt) node[below right] {\small $B' = [1, 0]$};
      \draw[->] (0.4,0) arc (0:70:0.4);
      \node at (35:0.58) {\scriptsize $x - y$};
      \node at (0,-1.4) {\small po otočení o úhel $-y$};
    \end{scope}
    \draw[->, thick, gray] (1.45,0.6) to[bend left=20] node[above] {\scriptsize otočení} (2.05,0.6);
  \end{tikzpicture}
\end{center}

Délka červené tětivy se otočením nezmění. Na levém obrázku ji ale počítáme z obecných souřadnic $A$ a $B$, na pravém z jednoduchých souřadnic $A'$ a $B'$. Obě čísla musí vyjít stejně, a z toho vypadne vzorec.

Ostatní vzorce z tohoto plynou. Dosazením $-y$ místo $y$ dostaneme $\cos(x + y)$. Ze vztahu $\sin x = \cos\left(\frac{\pi}{2} - x\right)$ dostaneme $\sin(x + y) = \cos\left(\left(\frac{\pi}{2} - x\right) - y\right)$ a po rozepsání vzorec pro sinus. Volbou $y = x$ dostaneme vzorce pro dvojnásobný úhel a z $\cos 2x = 1 - 2\sin^2 x$ po dosazení $\frac{x}{2}$ vzorce pro poloviční úhel. Celá tabulka vzorců tedy stojí na jednom obrázku.

\begin{example}
  Spočítejme přesně $\cos 15^\circ$:
  \begin{equation*}
    \cos 15^\circ = \cos(45^\circ - 30^\circ) = \frac{\sqrt{2}}{2}\cdot\frac{\sqrt{3}}{2} + \frac{\sqrt{2}}{2}\cdot\frac{1}{2} = \frac{\sqrt{6} + \sqrt{2}}{4}.
  \end{equation*}
\end{example}

\begin{example}
  Víme, že $\sin x = \frac{3}{5}$ a $x\in\left(\frac{\pi}{2}, \pi\right)$. Spočítejme ostatní funkce. Z $\sin^2 x + \cos^2 x = 1$ je $\cos^2 x = \frac{16}{25}$, a protože jsme ve druhém kvadrantu, je $\cos x = -\frac{4}{5}$. Pak
  \begin{equation*}
    \tg x = -\frac{3}{4},\qquad\sin 2x = 2\cdot\frac{3}{5}\cdot\left(-\frac{4}{5}\right) = -\frac{24}{25},\qquad\cos 2x = \frac{16}{25} - \frac{9}{25} = \frac{7}{25}.
  \end{equation*}
\end{example}

\begin{problem}
  \begin{enumerate}[label=(\alph*)]
    \item Spočítejte přesně $\sin 75^\circ$ a $\cos 105^\circ$.
    \item Víme, že $\cos x = -\frac{5}{13}$ a $x\in\left(\pi, \frac{3\pi}{2}\right)$. Určete $\sin x$, $\tg x$ a $\sin 2x$.
    \item Dokažte, že $\frac{\sin 2x}{1 + \cos 2x} = \tg x$ (pro přípustná $x$).
    \item Zjednodušte ${(\sin x + \cos x)}^2 - \sin 2x$.
  \end{enumerate}
\end{problem}

\begin{solution}
  \begin{enumerate}[label=(\alph*)]
    \item $\sin 75^\circ = \sin(45^\circ + 30^\circ) = \frac{\sqrt{6} + \sqrt{2}}{4}$. $\cos 105^\circ = \cos(60^\circ + 45^\circ) = \frac{\sqrt{2}}{4} - \frac{\sqrt{6}}{4} = \frac{\sqrt{2} - \sqrt{6}}{4}$.
    \item Ve třetím kvadrantu je sinus záporný: $\sin x = -\frac{12}{13}$, $\tg x = \frac{12}{5}$, $\sin 2x = 2\cdot\left(-\frac{12}{13}\right)\left(-\frac{5}{13}\right) = \frac{120}{169}$.
    \item $\frac{2\sin x\cos x}{1 + 2\cos^2 x - 1} = \frac{2\sin x\cos x}{2\cos^2 x} = \frac{\sin x}{\cos x} = \tg x$.
    \item $\sin^2 x + 2\sin x\cos x + \cos^2 x - 2\sin x\cos x = 1$.
  \end{enumerate}
\end{solution}

\section{Goniometrické rovnice}\label{sec:gon-rovnice}

Goniometrické funkce jsou periodické, a proto mají goniometrické rovnice (pokud mají řešení) nekonečně mnoho řešení. Najdeme je v jedné periodě a přičteme celistvé násobky periody.

\begin{definitionbox}
  \textbf{Základní goniometrické rovnice} ($k\in\mathbb{Z}$):
  \begin{itemize}
    \item $\sin x = a$, $|a|\leq 1$: najdeme $x_0\in\left[-\frac{\pi}{2}, \frac{\pi}{2}\right]$ se $\sin x_0 = a$. Řešení jsou $x = x_0 + 2k\pi$ a $x = \pi - x_0 + 2k\pi$.
    \item $\cos x = a$, $|a|\leq 1$: najdeme $x_0\in[0, \pi]$ s $\cos x_0 = a$. Řešení jsou $x = \pm x_0 + 2k\pi$.
    \item $\tg x = a$: najdeme $x_0\in\left(-\frac{\pi}{2}, \frac{\pi}{2}\right)$ s $\tg x_0 = a$. Řešení jsou $x = x_0 + k\pi$.
  \end{itemize}
  Pro $|a| > 1$ rovnice se sinem a kosinem nemají řešení.
\end{definitionbox}

Nejlepší je řešení vždy vidět na jednotkové kružnici. Rovnice $\sin x = \frac{1}{2}$ hledá body na kružnici ve výšce $\frac{1}{2}$. Takové body jsou dva: jeden v prvním kvadrantu ($\frac{\pi}{6}$) a jeden ve druhém ($\frac{5\pi}{6}$). Pak už jen přidáme $2k\pi$.

\begin{center}
  \begin{tikzpicture}[scale=1.7]
    \draw[->] (-1.3,0) -- (1.4,0);
    \draw[->] (0,-1.2) -- (0,1.3);
    \draw[thick] (0,0) circle (1);
    \draw[dashed, colorofwisdom] (-1.2,0.5) -- (1.2,0.5) node[right] {\small $y = \frac{1}{2}$};
    \filldraw[colorofwisdom] (30:1) circle (1pt) node[above right] {\small $\frac{\pi}{6}$};
    \filldraw[colorofwisdom] (150:1) circle (1pt) node[above left] {\small $\frac{5\pi}{6}$};
    \draw (0,0) -- (30:1) (0,0) -- (150:1);
  \end{tikzpicture}
\end{center}

Totéž na grafu: přímka $y = \frac{1}{2}$ protne sinusoidu v každé periodě dvakrát. Všechna řešení dostaneme z těchto dvou posouváním o celé násobky $2\pi$.

\begin{center}
  \begin{tikzpicture}[xscale=0.75]
    \draw[->] (-4,0) -- (9.9,0) node[right] {\small $x$};
    \draw[->] (0,-1.25) -- (0,1.6) node[above] {\small $y$};
    \foreach \t/\l in {-3.1416/{-\pi}, 3.1416/{\pi}, 6.2832/{2\pi}, 9.4248/{3\pi}} \draw (\t,0.06) -- (\t,-0.06) node[below] {\tiny $\l$};
    \draw[thick, colorofwisdom, domain=-3.9:9.7, samples=200] plot (\x, {sin(\x r)});
    \draw[dashed, gray] (-3.9,0.5) -- (9.7,0.5) node[right] {\small $y = \frac{1}{2}$};
    \foreach \p in {-3.6652, 0.5236, 2.6180, 6.8068, 8.9012} \fill (\p,0.5) circle (1.5pt);
    \node[above right] at (-3.67,0.5) {\tiny $-\frac{7\pi}{6}$};
    \node[above left] at (0.52,0.5) {\tiny $\frac{\pi}{6}$};
    \node[above right] at (2.62,0.5) {\tiny $\frac{5\pi}{6}$};
    \node[above left] at (6.81,0.5) {\tiny $\frac{13\pi}{6}$};
    \node[above right] at (8.90,0.5) {\tiny $\frac{17\pi}{6}$};
    \draw[<->, gray] (0.5236,1.3) -- (6.8068,1.3) node[midway, above] {\scriptsize $+2\pi$};
    \draw[dotted, gray] (0.5236,0.5) -- (0.5236,1.3) (6.8068,0.5) -- (6.8068,1.3);
  \end{tikzpicture}
\end{center}

\begin{example}
  Vyřešme:
  \begin{enumerate}[label=(\alph*)]
    \item $\sin x = \frac{1}{2}$: $x = \frac{\pi}{6} + 2k\pi$ nebo $x = \frac{5\pi}{6} + 2k\pi$, $k\in\mathbb{Z}$.
    \item $\cos x = -\frac{\sqrt{2}}{2}$: $x_0 = \frac{3\pi}{4}$, tedy $x = \pm\frac{3\pi}{4} + 2k\pi$.
    \item $\tg x = \sqrt{3}$: $x = \frac{\pi}{3} + k\pi$.
    \item $\sin 2x = 1$: $2x = \frac{\pi}{2} + 2k\pi$, tedy $x = \frac{\pi}{4} + k\pi$. Pozor, perioda se při dělení dvěma také vydělí!
  \end{enumerate}
\end{example}

\begin{definitionbox}
  \textbf{Složitější rovnice} převádíme na základní:
  \begin{itemize}
    \item \textbf{substitucí} (například $t = \sin x$, pokud se v rovnici vyskytuje jen $\sin x$ a jeho mocniny),
    \item \textbf{pomocí vzorců} (typicky $\sin^2 x = 1 - \cos^2 x$, abychom měli v rovnici jen jednu funkci),
    \item \textbf{rozkladem na součin} a větou o nulovém součinu.
  \end{itemize}
\end{definitionbox}

\begin{example}
  Vyřešme $2\sin^2 x - \sin x - 1 = 0$. Substituce $t = \sin x$: $2t^2 - t - 1 = 0$, tedy $t = 1$ nebo $t = -\frac{1}{2}$.
  \begin{itemize}
    \item $\sin x = 1$: $x = \frac{\pi}{2} + 2k\pi$,
    \item $\sin x = -\frac{1}{2}$: $x = -\frac{\pi}{6} + 2k\pi$ nebo $x = \frac{7\pi}{6} + 2k\pi$.
  \end{itemize}
\end{example}

\begin{example}
  Vyřešme $2\cos^2 x + 3\sin x - 3 = 0$. Nahradíme $\cos^2 x = 1 - \sin^2 x$:
  \begin{equation*}
    2 - 2\sin^2 x + 3\sin x - 3 = 0\quad\Leftrightarrow\quad 2\sin^2 x - 3\sin x + 1 = 0.
  \end{equation*}
  Tedy $\sin x = 1$ nebo $\sin x = \frac{1}{2}$, a tedy $x = \frac{\pi}{2} + 2k\pi$, $x = \frac{\pi}{6} + 2k\pi$ nebo $x = \frac{5\pi}{6} + 2k\pi$.
\end{example}

\begin{example}
  Vyřešme $\sin 2x = \sin x$. Podle vzorce $2\sin x\cos x - \sin x = 0$, tedy $\sin x(2\cos x - 1) = 0$.
  \begin{itemize}
    \item $\sin x = 0$: $x = k\pi$,
    \item $\cos x = \frac{1}{2}$: $x = \pm\frac{\pi}{3} + 2k\pi$.
  \end{itemize}
  Kdo by rovnici vydělil $\sin x$, přišel by o řešení $x = k\pi$. Tuto větu píšu v každé kapitole. Asi to něco znamená.
\end{example}

\begin{problem}
  Vyřešte v $\mathbb{R}$:
  \begin{enumerate}[label=(\alph*)]
    \item $2\cos x - 1 = 0$,
    \item $\sin^2 x = \frac{1}{4}$,
    \item $\cos 2x + \cos x = 0$,
    \item $\sin x + \cos x = 0$,
    \item $\tg^2 x - 3 = 0$.
  \end{enumerate}
\end{problem}

\begin{solution}
  \begin{enumerate}[label=(\alph*)]
    \item $\cos x = \frac{1}{2}$, $x = \pm\frac{\pi}{3} + 2k\pi$.
    \item $\sin x = \pm\frac{1}{2}$, $x = \pm\frac{\pi}{6} + k\pi$. Čtyři body na kružnici se dají zapsat dvěma soustavami s periodou $\pi$.
    \item $2\cos^2 x - 1 + \cos x = 0$, $\cos x\in\left\{\frac{1}{2}, -1\right\}$, tedy $x = \pm\frac{\pi}{3} + 2k\pi$ nebo $x = \pi + 2k\pi$.
    \item Pokud by $\cos x = 0$, bylo by $\sin x = \pm 1$ a rovnice by neplatila. Můžeme tedy dělit $\cos x$: $\tg x = -1$, $x = -\frac{\pi}{4} + k\pi$. Tady je dělení v pořádku, protože jsme předem ověřili, že nedělíme nulou.
    \item $\tg x = \pm\sqrt{3}$, $x = \pm\frac{\pi}{3} + k\pi$.
  \end{enumerate}
\end{solution}

\begin{example}
  Goniometrické nerovnice řešíme nejlépe pomocí jednotkové kružnice nebo grafu. $\sin x > \frac{1}{2}$ platí pro body na kružnici nad přímkou $y = \frac{1}{2}$, tedy na oblouku mezi $\frac{\pi}{6}$ a $\frac{5\pi}{6}$. $K = \bigcup_{k\in\mathbb{Z}}\left(\frac{\pi}{6} + 2k\pi, \frac{5\pi}{6} + 2k\pi\right)$.

  \begin{center}
    \begin{tikzpicture}[scale=1.5]
      \fill[colorofwisdom!12] (30:1) arc (30:150:1) -- cycle;
      \draw[->] (-1.3,0) -- (1.4,0);
      \draw[->] (0,-1.2) -- (0,1.3);
      \draw[thick] (0,0) circle (1);
      \draw[dashed, gray] (-1.25,0.5) -- (1.25,0.5) node[right] {\small $y = \frac{1}{2}$};
      \draw[line width=2.5pt, colorofwisdom] (30:1) arc (30:150:1);
      \draw[thick, colorofwisdom, fill=white] (30:1) circle (1.6pt) node[below right] {\small $\frac{\pi}{6}$};
      \draw[thick, colorofwisdom, fill=white] (150:1) circle (1.6pt) node[below left] {\small $\frac{5\pi}{6}$};
      \node[colorofwisdom, right] at (0.45,1.15) {\small $\sin x > \frac{1}{2}$};
    \end{tikzpicture}
  \end{center}
  Krajní body do řešení nepatří (prázdná kolečka), protože v nich je $\sin x = \frac{1}{2}$, ne větší.
\end{example}

\section{Řešení obecného trojúhelníku}\label{sec:obecny-trojuhelnik}

V obecném (ne nutně pravoúhlém) trojúhelníku značíme strany $a, b, c$ a protilehlé úhly $\alpha, \beta, \gamma$. Dvě věty nám umožní dopočítat všechno ostatní, pokud známe tři údaje (z nichž alespoň jeden je strana).

\begin{definitionbox}
  \textbf{Sinová věta}:
  \begin{equation*}
    \frac{a}{\sin\alpha} = \frac{b}{\sin\beta} = \frac{c}{\sin\gamma} = 2r,
  \end{equation*}
  kde $r$ je poloměr kružnice opsané.

  \textbf{Kosinová věta}:
  \begin{equation*}
    c^2 = a^2 + b^2 - 2ab\cos\gamma,
  \end{equation*}
  a obdobně pro ostatní strany.

  \textbf{Obsah trojúhelníku}:
  \begin{equation*}
    S = \frac{1}{2}ab\sin\gamma.
  \end{equation*}
\end{definitionbox}

\begin{proof}[(sinové věty bez části s $2r$)]
  Spusťme z vrcholu $C$ výšku $v_c$ na stranu $c$. Vzniknou dva pravoúhlé trojúhelníky. Z jednoho $v_c = b\sin\alpha$, z druhého $v_c = a\sin\beta$. Tedy $b\sin\alpha = a\sin\beta$, neboli $\frac{a}{\sin\alpha} = \frac{b}{\sin\beta}$. Pro třetí poměr stejně s jinou výškou. (Pro tupý úhel je důkaz stejný, jen pata výšky leží mimo stranu.)
\end{proof}

\begin{proof}[(kosinové věty)]
  Umístěme trojúhelník do souřadnic: $C = [0, 0]$, $B = [a, 0]$. Bod $A$ je ve vzdálenosti $b$ od $C$ pod úhlem $\gamma$, tedy $A = [b\cos\gamma, b\sin\gamma]$. Pak
  \begin{equation*}
    c^2 = |AB|^2 = {(b\cos\gamma - a)}^2 + b^2\sin^2\gamma = b^2(\cos^2\gamma + \sin^2\gamma) - 2ab\cos\gamma + a^2 = a^2 + b^2 - 2ab\cos\gamma.
  \end{equation*}
\end{proof}

\begin{center}
  \begin{tikzpicture}[scale=0.95]
    \begin{scope}
      \fill[colorofwisdom!8] (0,0) -- (5,0) -- (1.8,2.6) -- cycle;
      \draw[thick] (0,0) -- (5,0) -- (1.8,2.6) -- cycle;
      \draw[thick, dashed, colorofwisdom] (1.8,2.6) -- (1.8,0);
      \draw (1.8,0.2) -- (2,0.2) -- (2,0);
      \node[right, colorofwisdom] at (1.8,1.1) {\small $v_c$};
      \node[below left] at (0,0) {\small $A$};
      \node[below right] at (5,0) {\small $B$};
      \node[above] at (1.8,2.6) {\small $C$};
      \node[below] at (3.4,0) {\small $c$};
      \node[above left] at (0.9,1.3) {\small $b$};
      \node[above right] at (3.4,1.3) {\small $a$};
      \draw (0.6,0) arc (0:55.3:0.6);
      \node at (27:0.85) {\scriptsize $\alpha$};
      \draw ($(5,0)+(140.9:0.6)$) arc (140.9:180:0.6);
      \node at ($(5,0)+(160:0.85)$) {\scriptsize $\beta$};
      \node[align=center] at (2.5,-0.95) {\small $v_c = b\sin\alpha = a\sin\beta$};
    \end{scope}
    \begin{scope}[xshift=6.6cm]
      \fill[colorofwisdom!8] (0,0) -- (4.5,0) -- (1.7207,2.4575) -- cycle;
      \draw[thick] (0,0) -- (4.5,0) -- (1.7207,2.4575) -- cycle;
      \draw[dashed, gray] (1.7207,2.4575) -- (1.7207,0);
      \node[below left] at (0,0) {\small $C$};
      \node[below right] at (4.5,0) {\small $B = [a, 0]$};
      \node[above] at (1.7207,2.4575) {\small $A = [b\cos\gamma, b\sin\gamma]$};
      \node[above left] at (0.86,1.23) {\small $b$};
      \node[above right] at (3.11,1.23) {\small $c$};
      \node[above] at (3.1,0) {\small $a$};
      \draw (0.55,0) arc (0:55:0.55);
      \node at (27:0.8) {\scriptsize $\gamma$};
      \draw[decorate, decoration={brace, mirror, amplitude=4pt}] (0,-0.1) -- (1.7207,-0.1) node[midway, below=4pt] {\scriptsize $b\cos\gamma$};
      \node[right, gray] at (1.7207,1.0) {\scriptsize $b\sin\gamma$};
    \end{scope}
  \end{tikzpicture}
\end{center}

Kosinová věta je zobecněním Pythagorovy věty. Pro $\gamma = 90^\circ$ je $\cos\gamma = 0$ a dostaneme $c^2 = a^2 + b^2$. Člen $-2ab\cos\gamma$ je `oprava' pro trojúhelníky, které pravoúhlé nejsou.

\begin{example}
  V trojúhelníku je $a = 5$, $b = 7$ a $\gamma = 60^\circ$. Dopočítejme $c$ a obsah.
  \begin{equation*}
    c^2 = 25 + 49 - 2\cdot 5\cdot 7\cdot\frac{1}{2} = 39,\qquad c = \sqrt{39}\approx 6{,}24,\qquad S = \frac{1}{2}\cdot 5\cdot 7\cdot\frac{\sqrt{3}}{2} = \frac{35\sqrt{3}}{4}\approx 15{,}16.
  \end{equation*}
\end{example}

\begin{example}
  V trojúhelníku je $a = 10$, $\alpha = 30^\circ$ a $\beta = 45^\circ$. Pak $\gamma = 105^\circ$ a podle sinové věty
  \begin{equation*}
    b = \frac{a\sin\beta}{\sin\alpha} = \frac{10\cdot\frac{\sqrt{2}}{2}}{\frac{1}{2}} = 10\sqrt{2},\qquad c = \frac{a\sin\gamma}{\sin\alpha} = 20\sin 105^\circ = 20\cdot\frac{\sqrt{6} + \sqrt{2}}{4} = 5(\sqrt{6} + \sqrt{2}).
  \end{equation*}
  Hodnotu $\sin 105^\circ = \sin(60^\circ + 45^\circ)$ jsme spočítali součtovým vzorcem.
\end{example}

\begin{note}
  Pokud známe dvě strany a úhel \textbf{proti kratší} z nich, sinová věta může dát dvě řešení, protože $\sin\beta = \sin(180^\circ - \beta)$. Trojúhelník pak existuje ve dvou verzích (s ostrým a s tupým úhlem), nebo vůbec. Tomu se říká \textbf{nejednoznačný případ}. Kosinová věta tímto problémem netrpí, protože kosinus na $(0^\circ, 180^\circ)$ je prostý.
\end{note}

\begin{center}
  \begin{tikzpicture}[scale=1.1]
    \draw[thick, ->] (0,0) -- (6.2,0);
    \draw[dashed, gray] ($(3.464,2)+(-145:2.5)$) arc (-145:-35:2.5);
    \draw[thick] (0,0) -- (3.464,2) node[midway, above left] {\small $b$};
    \draw[thick, colorofwisdom] (3.464,2) -- (1.964,0) node[midway, left] {\small $a$};
    \draw[thick] (3.464,2) -- (4.964,0) node[midway, right] {\small $a$};
    \fill (0,0) circle (1.5pt) node[below left] {\small $A$};
    \fill (3.464,2) circle (1.5pt) node[above] {\small $C$};
    \fill[colorofwisdom] (1.964,0) circle (1.5pt) node[below] {\small $B_1$};
    \fill (4.964,0) circle (1.5pt) node[below] {\small $B_2$};
    \draw (0.7,0) arc (0:30:0.7);
    \node at (15:0.95) {\scriptsize $\alpha$};
  \end{tikzpicture}
\end{center}
Na obrázku známe $\alpha$, $b$ a $a < b$. Bod $B$ musí ležet na polopřímce z $A$ a zároveň ve vzdálenosti $a$ od $C$, tedy na kružnici se středem $C$. Ta polopřímku protne dvakrát, a proto vzniknou dva různé trojúhelníky $AB_1C$ (s tupým úhlem u $B_1$) a $AB_2C$. Kdyby byla strana $a$ příliš krátká, kružnice by polopřímku vůbec nezasáhla.

\begin{problem}
  \begin{enumerate}[label=(\alph*)]
    \item Trojúhelník má strany $a = 8$, $b = 5$, $c = 7$. Určete úhel $\gamma$.
    \item Dvě lodě vyplují ze stejného přístavu. Jedna jede $12\ \si{km}$ na sever, druhá $8\ \si{km}$ směrem, který svírá se severem úhel $60^\circ$. Jak daleko jsou od sebe?
    \item Určete obsah trojúhelníku se stranami $b = 6$, $c = 4$ a úhlem $\alpha = 150^\circ$.
  \end{enumerate}
\end{problem}

\begin{solution}
  \begin{enumerate}[label=(\alph*)]
    \item $\cos\gamma = \frac{a^2 + b^2 - c^2}{2ab} = \frac{64 + 25 - 49}{80} = \frac{1}{2}$, $\gamma = 60^\circ$.
    \item $d^2 = 144 + 64 - 2\cdot 12\cdot 8\cdot\frac{1}{2} = 112$, $d = 4\sqrt{7}\approx 10{,}6\ \si{km}$.
    \item $S = \frac{1}{2}\cdot 6\cdot 4\cdot\sin 150^\circ = 12\cdot\frac{1}{2} = 6$.
  \end{enumerate}
\end{solution}

\section{Příběh --- Muž, který miloval teplo}

Jean-Baptiste Joseph Fourier vyrůstal jako sirotek, za Francouzské revoluce byl zatčen a jen o vlásek unikl gilotině, a pak doprovázel Napoleona na výpravě do Egypta. Po návratu se začal zabývat šířením tepla a v roce 1807 přišel s tvrzením, které většina tehdejších matematiků považovala za nesmysl: \textit{každou} rozumnou periodickou funkci lze zapsat jako součet (nekonečně mnoha) sinů a kosinů různých frekvencí. Lagrange byl proti, práce byla roky odmítána a vyšla až v roce 1822.

Fourier měl pravdu. Dnes se tomu říká \textbf{Fourierova řada} (a obecněji Fourierova transformace) a stojí na tom značná část moderního světa. Váš telefon rozloží zvuk na sinusovky, ty nepodstatné zahodí a zbytek uloží. Tak funguje MP3. Fotka ve formátu JPEG se ukládá velice podobně, jen s kosiny ve dvou směrech. Na stejném principu fungují i rentgenová tomografie, rozpoznávání řeči a magnetická rezonance.

Teplo, kterým se Fourier celý život zabýval, se nakonec stalo i jeho posedlostí. Po návratu z Egypta byl přesvědčený, že horko je pro zdraví prospěšné. Svůj byt prý vytápěl tak, že návštěvníci nemohli dýchat, a sám chodil zabalený v několika vrstvách přikrývek. V roce 1830 spadl ze schodů a o několik týdnů později zemřel. Jeho jméno je jedním ze 72 jmen vědců vyrytých na Eiffelově věži.

\section{Shrnutí a otázky k ústní zkoušce}

\begin{problem}
  Odpovězte vlastními slovy:
  \begin{enumerate}
    \item Co je stupňová a oblouková míra úhlu? Jak se mezi nimi převádí? Proč se v matematice dává přednost radiánům?
    \item Jak jsou definovány goniometrické funkce v pravoúhlém trojúhelníku a na jednotkové kružnici? Proč se obě definice pro ostré úhly shodují?
    \item Odvoďte hodnoty goniometrických funkcí pro $30^\circ$, $45^\circ$ a $60^\circ$.
    \item Nakreslete grafy sinu, kosinu a tangentu a popište jejich vlastnosti.
    \item Odvoďte vztah $\sin^2 x + \cos^2 x = 1$. Dokažte součtový vzorec pro kosinus a odvoďte z něj vzorce pro dvojnásobný úhel.
    \item Jak se řeší základní goniometrické rovnice? Proč mají nekonečně mnoho řešení?
    \item Jaké znáte metody řešení složitějších goniometrických rovnic?
    \item Vyslovte a dokažte sinovou a kosinovou větu. Jak souvisí kosinová věta s Pythagorovou větou?
    \item Kdy může mít úloha o trojúhelníku dvě řešení?
  \end{enumerate}
\end{problem}

\koneckapitoly

\chapter{Planimetrie}
\label{kap:planimetrie}
\mathquote{Geometry is the art of correct reasoning from incorrectly drawn figures.}{George Pólya}

Planimetrie je geometrie v rovině. Je to nejstarší systematicky budovaná část matematiky. Eukleidovy \textit{Základy} ze 3. století před naším letopočtem jsou po Bibli pravděpodobně nejčastěji vydávanou knihou v dějinách a po dva tisíce let se z nich učilo, co je to důkaz. Právě v geometrii se poprvé objevila myšlenka, že z několika málo axiomů lze logicky odvodit všechno ostatní. Tu samou myšlenku jsme viděli u teorie množin v kapitole 2.

\section{Základní pojmy}

\begin{definitionbox}
  \textbf{Základní útvary}: Bod, přímka a rovina jsou základní pojmy, které nedefinujeme (podobně jako množinu). Značíme body velkými písmeny ($A, B, \ldots$) a přímky malými ($p, q, \ldots$).
  \begin{itemize}
    \item \textbf{Polopřímka} $\overrightarrow{AB}$: část přímky začínající v bodě $A$ a obsahující bod $B$.
    \item \textbf{Úsečka} $AB$: množina bodů přímky $AB$ ležících mezi $A$ a $B$ (včetně nich). Její délku značíme $|AB|$.
    \item \textbf{Polorovina} $pA$: část roviny ohraničená přímkou $p$ a obsahující bod $A$.
    \item \textbf{Úhel} $\angle AVB$: část roviny ohraničená dvěma polopřímkami $\overrightarrow{VA}$ a $\overrightarrow{VB}$ se společným počátkem $V$ (vrcholem).
  \end{itemize}

  \textbf{Vzájemná poloha dvou přímek v rovině}: \textbf{rovnoběžné různé} (žádný společný bod), \textbf{totožné} (všechny body společné), \textbf{různoběžné} (právě jeden společný bod, průsečík). Různoběžky jsou \textbf{kolmé}, pokud svírají pravý úhel.
\end{definitionbox}

\begin{definitionbox}
  \textbf{Dvojice úhlů}:
  \begin{itemize}
    \item \textbf{Vedlejší úhly} mají společné jedno rameno a další dvě ramena tvoří přímku. Jejich součet je $180^\circ$.
    \item \textbf{Vrcholové úhly} mají společný vrchol a jejich ramena tvoří dvě přímky. Jsou \textbf{shodné}.
    \item Pokud dvě \textbf{rovnoběžky} protneme třetí přímkou, jsou \textbf{souhlasné} úhly shodné a \textbf{střídavé} úhly shodné.
  \end{itemize}
  Úhel je \textbf{ostrý} ($< 90^\circ$), \textbf{pravý} ($= 90^\circ$), \textbf{tupý} ($> 90^\circ$ a $< 180^\circ$), \textbf{přímý} ($= 180^\circ$).
\end{definitionbox}

\begin{center}
  \begin{tikzpicture}[scale=0.9]
    \foreach \a/\n/\x in {50/ostrý/0, 90/pravý/2.6, 135/tupý/5.2, 180/přímý/8.6} {
      \begin{scope}[xshift=\x cm]
        \draw[thick] (1.2,0) -- (0,0) -- (\a:1.2);
        \ifnum\a=90 \draw (0.22,0) -- (0.22,0.22) -- (0,0.22); \else \draw[colorofwisdom, thick] (0.35,0) arc (0:\a:0.35); \fi
        \filldraw (0,0) circle (1pt);
        \node[below, font=\small] at (0.3,-0.25) {\n};
      \end{scope}
    }
  \end{tikzpicture}

  \vspace{0.3cm}

  \begin{tikzpicture}[scale=1.1]
    \draw[thick] (-2,-0.9) -- (2,0.9);
    \draw[thick] (-2,0.6) -- (2,-0.6);
    \filldraw (0,0) circle (1pt);
    \draw[colorofwisdom, thick] (0.45,0.2025) arc (24.2:-16.7:0.49);
    \draw[colorofwisdom, thick] (-0.45,-0.2025) arc (204.2:163.3:0.49);
    \draw[thick] (0.5,0.225) arc (24.2:163.3:0.55);
    \node[colorofwisdom] at (0.8,0.03) {\small $\alpha$};
    \node[colorofwisdom] at (-0.8,-0.03) {\small $\alpha$};
    \node at (0,0.75) {\small $\beta = 180^\circ - \alpha$};
    \node[font=\scriptsize, align=center] at (0,-1.3) {vrcholové úhly ($\alpha$ a $\alpha$) jsou shodné,\\ vedlejší ($\alpha$ a $\beta$) dávají $180^\circ$};
  \end{tikzpicture}
\end{center}

\begin{center}
  \begin{tikzpicture}[scale=0.9]
    \draw[thick] (-1,0) -- (6,0) node[right] {\small $p$};
    \draw[thick] (-1,2) -- (6,2) node[right] {\small $q$};
    \draw[thick] (0.5,-0.8) -- (4.3,2.8);
    \coordinate (A) at (1.35,0);
    \coordinate (B) at (3.46,2);
    \draw[colorofwisdom, thick] ($(A)+(0.5,0)$) arc (0:43.5:0.5);
    \draw[colorofwisdom, thick] ($(B)+(0.5,0)$) arc (0:43.5:0.5);
    \draw[thick] ($(B)+(-0.5,0)$) arc (180:223.5:0.5);
    \node[colorofwisdom] at ($(A)+(0.75,0.25)$) {\small $\alpha$};
    \node[colorofwisdom] at ($(B)+(0.75,0.25)$) {\small $\alpha$};
    \node at ($(B)+(-0.75,-0.25)$) {\small $\alpha$};
    \node[below] at (3,-0.9) {\small souhlasné (červeně) a střídavé úhly jsou shodné};
  \end{tikzpicture}
\end{center}

\section{Trojúhelník}

\begin{definitionbox}
  \textbf{Trojúhelník} $ABC$ je útvar určený třemi body, které neleží na jedné přímce. Strany značíme $a = |BC|$, $b = |AC|$, $c = |AB|$ (strana proti vrcholu), úhly $\alpha, \beta, \gamma$ při vrcholech $A, B, C$.

  \vspace{0.2cm}

  \textbf{Věta (o součtu úhlů)}: $\alpha + \beta + \gamma = 180^\circ$.

  \textbf{Trojúhelníková nerovnost}: Součet každých dvou stran je větší než strana třetí: $a + b > c$, $a + c > b$, $b + c > a$.

  \textbf{Vztah stran a úhlů}: Proti větší straně leží větší úhel a naopak.
\end{definitionbox}

\begin{proof}[(součtu úhlů)]
  Veďme vrcholem $C$ přímku $p$ rovnoběžnou se stranou $AB$. Přímka $p$ a ramena $CA$, $CB$ tvoří u vrcholu $C$ tři úhly, které dohromady dávají přímý úhel $180^\circ$. Prostřední z nich je $\gamma$. Levý je střídavý s úhlem $\alpha$ (rovnoběžky $p$ a $AB$ protnuté přímkou $AC$), a tedy má velikost $\alpha$. Pravý je ze stejného důvodu roven $\beta$. Tedy $\alpha + \gamma + \beta = 180^\circ$.
\end{proof}

\begin{center}
  \begin{tikzpicture}[scale=1]
    \coordinate (A) at (0,0); \coordinate (B) at (5,0); \coordinate (C) at (1.8,2.6);
    \draw[thick] (A) -- (B) -- (C) -- cycle;
    \draw[dashed] (-1,2.6) -- (6,2.6) node[right] {\small $p\parallel AB$};
    \node[below left] at (A) {$A$}; \node[below right] at (B) {$B$}; \node[above] at ($(C)+(0,0.1)$) {$C$};
    \draw[colorofwisdom, thick] (0.6,0) arc (0:55.3:0.6); \node[colorofwisdom] at (0.85,0.35) {\small $\alpha$};
    \draw[thick] (4.4,0) arc (180:140.9:0.6); \node at (4.15,0.25) {\small $\beta$};
    \draw[colorofwisdom, thick] ($(C)+(-0.6,0)$) arc (180:235.3:0.6); \node[colorofwisdom] at ($(C)+(-0.85,-0.3)$) {\small $\alpha$};
    \draw[thick] ($(C)+(0.6,0)$) arc (0:-39.1:0.6); \node at ($(C)+(0.9,-0.25)$) {\small $\beta$};
    \node at ($(C)+(0.1,-0.55)$) {\small $\gamma$};
  \end{tikzpicture}
\end{center}

\begin{note}
  Tento důkaz používá vlastnost rovnoběžek, která stojí na Eukleidově pátém postulátu. Bez něj věta o součtu úhlů neplatí! Na kouli může mít trojúhelník třeba tři pravé úhly. Více v odbočce na konci kapitoly.
\end{note}

\begin{definitionbox}
  \textbf{Rozdělení trojúhelníků}:
  \begin{itemize}
    \item podle stran: \textbf{různostranný}, \textbf{rovnoramenný} (dvě strany shodné, ramena; třetí je základna; úhly při základně jsou shodné), \textbf{rovnostranný} (všechny strany shodné, všechny úhly $60^\circ$),
    \item podle úhlů: \textbf{ostroúhlý}, \textbf{pravoúhlý} (strana proti pravému úhlu je \textbf{přepona}, ostatní \textbf{odvěsny}), \textbf{tupoúhlý}.
  \end{itemize}
\end{definitionbox}

\subsection{Významné přímky a body trojúhelníku}

\begin{definitionbox}
  \begin{itemize}
    \item \textbf{Těžnice} je úsečka spojující vrchol se středem protější strany. Tři těžnice se protínají v \textbf{těžišti} $T$, které dělí každou těžnici v poměru $2:1$ (delší část je u vrcholu).
    \item \textbf{Výška} je úsečka z vrcholu kolmá na protější stranu (nebo její prodloužení). Přímky, na kterých leží výšky, se protínají v \textbf{ortocentru} $V$.
    \item \textbf{Osy stran} (přímky kolmé na stranu v jejím středu) se protínají ve středu $S$ \textbf{kružnice opsané}. Bod $S$ má od všech vrcholů stejnou vzdálenost $r$.
    \item \textbf{Osy úhlů} se protínají ve středu $O$ \textbf{kružnice vepsané}. Bod $O$ má od všech stran stejnou vzdálenost $\rho$.
    \item \textbf{Střední příčka} spojuje středy dvou stran. Je rovnoběžná se třetí stranou a má poloviční délku.
  \end{itemize}
\end{definitionbox}

\begin{center}
  \begin{tikzpicture}[scale=0.62]
    \coordinate (A) at (0,0); \coordinate (B) at (5,0); \coordinate (C) at (1.5,3); \draw[thick] (A) -- (B) -- (C) -- cycle; \node[below left] at (A) {\scriptsize $A$}; \node[below right] at (B) {\scriptsize $B$}; \node[above] at (C) {\scriptsize $C$};
    \coordinate (Ma) at (3.25,1.5); \coordinate (Mb) at (0.75,1.5); \coordinate (Mc) at (2.5,0);
    \draw[colorofwisdom] (A) -- (Ma) (B) -- (Mb) (C) -- (Mc);
    \foreach \m in {Ma,Mb,Mc} \filldraw (\m) circle (1.5pt);
    \filldraw[colorofwisdom] (2.1667,1) circle (2.5pt) node[right, xshift=2pt] {\scriptsize $T$};
    \node[font=\small] at (2.5,-0.8) {těžnice a těžiště};
  \end{tikzpicture}
  \hspace{0.3cm}
  \begin{tikzpicture}[scale=0.62]
    \coordinate (A) at (0,0); \coordinate (B) at (5,0); \coordinate (C) at (1.5,3); \draw[thick] (A) -- (B) -- (C) -- cycle; \node[below left] at (A) {\scriptsize $A$}; \node[below right] at (B) {\scriptsize $B$}; \node[above] at (C) {\scriptsize $C$};
    \draw[colorofwisdom] (C) -- (1.5,0) (A) -- (2.1176,2.4706) (B) -- (1,2);
    \draw (1.5,0) rectangle (1.75,0.25);
    \filldraw[colorofwisdom] (1.5,1.75) circle (2.5pt) node[left, xshift=-2pt] {\scriptsize $V$};
    \node[font=\small] at (2.5,-0.8) {výšky a ortocentrum};
  \end{tikzpicture}

  \vspace{0.3cm}

  \begin{tikzpicture}[scale=0.62]
    \coordinate (A) at (0,0); \coordinate (B) at (5,0); \coordinate (C) at (1.5,3); \draw[thick] (A) -- (B) -- (C) -- cycle; \node[below left] at (A) {\scriptsize $A$}; \node[below right] at (B) {\scriptsize $B$}; \node[above] at (C) {\scriptsize $C$};
    \draw[colorofwisdom!80, dashed] (2.5,-0.4) -- (2.5,2.8);
    \draw[colorofwisdom!80, dashed] ($(0.75,1.5)!-1.1!(2.5,0.625)$) -- ($(0.75,1.5)!1.6!(2.5,0.625)$);
    \draw[colorofwisdom!80, dashed] ($(3.25,1.5)!-0.8!(2.5,0.625)$) -- ($(3.25,1.5)!1.9!(2.5,0.625)$);
    \draw[thick, colorofwisdom] (2.5,0.625) circle (2.577);
    \filldraw[colorofwisdom] (2.5,0.625) circle (2.5pt) node[below right] {\scriptsize $S$};
    \draw[gray] (2.5,0.625) -- (C) node[midway, left, font=\scriptsize] {$r$};
    \node[font=\small] at (2.5,-2.4) {osy stran a kružnice opsaná};
  \end{tikzpicture}
  \hspace{0.3cm}
  \begin{tikzpicture}[scale=0.62]
    \coordinate (A) at (0,0); \coordinate (B) at (5,0); \coordinate (C) at (1.5,3); \draw[thick] (A) -- (B) -- (C) -- cycle; \node[below left] at (A) {\scriptsize $A$}; \node[below right] at (B) {\scriptsize $B$}; \node[above] at (C) {\scriptsize $C$};
    \coordinate (I) at (1.8722,1.1571);
    \draw[colorofwisdom!80, dashed] (A) -- ($(A)!1.25!(I)$) (B) -- ($(B)!1.2!(I)$) (C) -- ($(C)!1.6!(I)$);
    \draw[thick, colorofwisdom] (I) circle (1.1571);
    \filldraw[colorofwisdom] (I) circle (2.5pt) node[above right] {\scriptsize $O$};
    \draw[gray] (I) -- (1.8722,0) node[midway, right, font=\scriptsize] {$\rho$};
    \node[font=\small] at (2.5,-0.8) {osy úhlů a kružnice vepsaná};
  \end{tikzpicture}
\end{center}

Všimněte si, že na prvních třech obrázcích je stejný trojúhelník a body $T$, $V$, $S$ opravdu leží na jedné přímce (Eulerově, viz poznámka níže). Střed kružnice vepsané na ní obecně neleží.

Proč se osy stran protínají v jednom bodě? Osa úsečky $AB$ je množina bodů, které mají od $A$ a $B$ stejnou vzdálenost. Průsečík os $AB$ a $BC$ má tedy stejnou vzdálenost od $A$ a $B$ a zároveň od $B$ a $C$. Má proto stejnou vzdálenost od $A$ a $C$ a leží i na ose $AC$. Pro osy úhlů je argument stejný, jen se místo vzdálenosti od bodů mluví o vzdálenosti od přímek.

\begin{note}
  Těžiště $T$, ortocentrum $V$ a střed kružnice opsané $S$ leží v každém trojúhelníku na jedné přímce, a navíc $|VT| = 2|TS|$. Této přímce se říká \textbf{Eulerova přímka}. Euler ji objevil v roce 1765, tedy přes dva tisíce let poté, co si Řekové s trojúhelníky začali hrát. Někdy jsou věci zajímavé i po dvou tisících letech.
\end{note}

\subsection{Shodnost a podobnost}

\begin{definitionbox}
  Dva trojúhelníky jsou \textbf{shodné}, pokud mají shodné všechny odpovídající si strany a úhly. Stačí ale ověřit jen některé z nich.

  \textbf{Věty o shodnosti trojúhelníků}:
  \begin{itemize}
    \item \textbf{sss}: shodují se ve všech třech stranách,
    \item \textbf{sus}: shodují se ve dvou stranách a úhlu jimi sevřeném,
    \item \textbf{usu}: shodují se v jedné straně a dvou úhlech k ní přilehlých,
    \item \textbf{Ssu}: shodují se ve dvou stranách a úhlu proti \textbf{větší} z nich.
  \end{itemize}

  Dva trojúhelníky jsou \textbf{podobné}, pokud existuje číslo $k > 0$ (\textbf{koeficient podobnosti}) takové, že strany jednoho jsou $k$-násobkem odpovídajících stran druhého. Podobné trojúhelníky mají shodné odpovídající úhly.

  \textbf{Věty o podobnosti trojúhelníků}:
  \begin{itemize}
    \item \textbf{uu}: shodují se ve dvou úhlech,
    \item \textbf{sus}: shodují se v poměru dvou stran a úhlu jimi sevřeném,
    \item \textbf{sss}: shodují se v poměrech všech tří stran.
  \end{itemize}
  Obvody podobných trojúhelníků jsou v poměru $k$, obsahy v poměru $k^2$.
\end{definitionbox}

Proč ve větě Ssu musí být úhel proti větší straně? Je to přesně nejednoznačný případ ze sinové věty v kapitole~\ref{kap:goniometrie}. Se dvěma stranami a úhlem proti kratší z nich lze sestrojit dva různé trojúhelníky.

\begin{example}
  Strom vrhá stín dlouhý $12\ \si{m}$. Ve stejnou chvíli vrhá metrová tyč stín dlouhý $1{,}5\ \si{m}$. Jak vysoký je strom?

  Paprsky slunce jsou rovnoběžné, takže dopadají na strom i na tyč pod stejným úhlem. Strom se stínem a tyč se stínem tvoří pravoúhlé trojúhelníky, které jsou podle věty uu podobné. Tedy $\frac{h}{12} = \frac{1}{1{,}5}$ a $h = 8\ \si{m}$. Podle legendy takhle změřil Thalés z Milétu výšku egyptských pyramid.
  \begin{center}
    \begin{tikzpicture}[scale=0.45]
      \draw[thick] (-1,0) -- (16,0);
      \fill[colorofwisdom!20] (0,0) -- (12,0) -- (0,8) -- cycle;
      \draw[very thick, black!70] (0,0) -- (0,8);
      \fill[black!60] (0,8.6) circle (1.1);
      \draw[dashed, colorofwisdom] (0,8) -- (12,0);
      \node[below] at (6,0) {\small stín $12\ \si{m}$};
      \node[left] at (0,4) {\small $h$};
      \fill[colorofwisdom!20] (13,0) -- (14.5,0) -- (13,1) -- cycle;
      \draw[very thick] (13,0) -- (13,1) node[left, font=\scriptsize] {$1$};
      \draw[dashed, colorofwisdom] (13,1) -- (14.5,0);
      \node[below] at (13.75,0) {\scriptsize $1{,}5$};
      \foreach \y in {2,4,6} \draw[->, gray!60] (-1+\y*1.5,\y+3) -- (\y*1.5+0.5,\y+2);
      \node[gray] at (11,8) {\small sluneční paprsky};
    \end{tikzpicture}
  \end{center}
\end{example}

\subsection{Pythagorova a Eukleidovy věty}

\begin{definitionbox}
  \textbf{Pythagorova věta}: V pravoúhlém trojúhelníku s odvěsnami $a, b$ a přeponou $c$ platí
  \begin{equation*}
    a^2 + b^2 = c^2.
  \end{equation*}
  Platí i \textbf{obrácená věta}: pokud pro strany trojúhelníku platí $a^2 + b^2 = c^2$, je trojúhelník pravoúhlý s pravým úhlem proti straně $c$.

  \vspace{0.2cm}

  \textbf{Eukleidovy věty}: Nechť $v$ je výška na přeponu, která ji rozdělí na úseky $c_a$ (přilehlý k odvěsně $a$) a $c_b$. Pak
  \begin{equation*}
    v^2 = c_a\cdot c_b\quad\text{(o výšce)},\qquad a^2 = c\cdot c_a,\quad b^2 = c\cdot c_b\quad\text{(o odvěsně)}.
  \end{equation*}
\end{definitionbox}

\begin{proof}[(Pythagorovy věty)]
  Vezměme čtverec o straně $a + b$ a umístěme do něj čtyři kopie našeho trojúhelníku tak, aby jejich přepony ohraničovaly vnitřní čtverec o straně $c$ (viz obrázek). Obsah velkého čtverce spočítáme dvakrát:
  \begin{equation*}
    {(a + b)}^2 = 4\cdot\frac{ab}{2} + c^2\quad\Leftrightarrow\quad a^2 + 2ab + b^2 = 2ab + c^2\quad\Leftrightarrow\quad a^2 + b^2 = c^2.
  \end{equation*}
  (To, že vnitřní útvar je opravdu čtverec, plyne z toho, že u každého jeho vrcholu se úhly trojúhelníků $\alpha$ a $\beta = 90^\circ - \alpha$ doplňují na $90^\circ$, takže zbývající úhel je také $90^\circ$.)
\end{proof}

\begin{center}
  \begin{tikzpicture}[scale=0.7]
    \draw[thick] (0,0) rectangle (5,5);
    \fill[colorofwisdom!20] (2,0) -- (5,2) -- (3,5) -- (0,3) -- cycle;
    \draw[thick, colorofwisdom] (2,0) -- (5,2) -- (3,5) -- (0,3) -- cycle;
    \node[below] at (1,0) {\small $a$}; \node[below] at (3.5,0) {\small $b$};
    \node[right] at (5,1) {\small $a$}; \node[right] at (5,3.5) {\small $b$};
    \node[above] at (4,5) {\small $a$}; \node[above] at (1.5,5) {\small $b$};
    \node[left] at (0,4) {\small $a$}; \node[left] at (0,1.5) {\small $b$};
    \node[colorofwisdom] at (2.5,2.5) {$c^2$};
    \node[colorofwisdom] at (3.7,0.75) {\small $c$};
  \end{tikzpicture}
\end{center}

\begin{center}
  \begin{tikzpicture}[scale=1.1]
    \coordinate (A) at (0,0); \coordinate (B) at (5,0); \coordinate (C) at (1.8,2.4); \coordinate (P) at (1.8,0);
    \fill[colorofwisdom!12] (A) -- (P) -- (C) -- cycle;
    \fill[black!8] (P) -- (B) -- (C) -- cycle;
    \draw[thick] (A) -- (B) -- (C) -- cycle;
    \draw[thick, colorofwisdom] (C) -- (P) node[midway, right] {\small $v$};
    \draw (P) rectangle (1.95,0.15);
    \draw ($(C)!0.13!(A)$) -- ($(C)!0.13!(A)!0.13!(B)$);
    \draw[decorate, decoration={brace, mirror, amplitude=4pt}] (0,-0.1) -- (1.8,-0.1) node[midway, below=4pt] {\small $c_b$};
    \draw[decorate, decoration={brace, mirror, amplitude=4pt}] (1.8,-0.1) -- (5,-0.1) node[midway, below=4pt] {\small $c_a$};
    \node[above left] at (0.9,1.2) {\small $b$};
    \node[above right] at (3.4,1.2) {\small $a$};
    \node[below left] at (A) {\small $A$}; \node[below right] at (B) {\small $B$}; \node[above] at (C) {\small $C$};
  \end{tikzpicture}
\end{center}

\begin{proof}[(Eukleidových vět)]
  Výška na přeponu rozdělí trojúhelník na dva menší pravoúhlé trojúhelníky. Oba mají s původním společný jeden ostrý úhel, a tedy jsou mu podle věty uu podobné, a tedy i sobě navzájem. Z podobnosti obou malých trojúhelníků plyne $\frac{v}{c_a} = \frac{c_b}{v}$, tedy $v^2 = c_ac_b$. Z podobnosti malého trojúhelníku u odvěsny $a$ s původním plyne $\frac{c_a}{a} = \frac{a}{c}$, tedy $a^2 = c\cdot c_a$.
\end{proof}

Pythagorovu větu znali Babyloňané přes tisíc let před Pythagorem. Existuje jí přes 370 různých důkazů, jeden z nich vymyslel i americký prezident James Garfield. Sám Pythagoras byl mimochodem spíš vůdce náboženské sekty než matematik. Jeho následovníci nesměli jíst fazole a podle jedné legendy Pythagoras zemřel, když utíkal před nepřáteli a odmítl proběhnout fazolovým polem.

\begin{example}
  Výška na přeponu dělí přeponu na úseky $c_a = 9$ a $c_b = 16$. Dopočítejme ostatní. Přepona je $c = 25$, výška $v = \sqrt{9\cdot 16} = 12$, odvěsny $a = \sqrt{25\cdot 9} = 15$ a $b = \sqrt{25\cdot 16} = 20$. Kontrola: $15^2 + 20^2 = 625 = 25^2$.
\end{example}

\begin{problem}
  \begin{enumerate}[label=(\alph*)]
    \item Je trojúhelník se stranami $7$, $24$, $25$ pravoúhlý? A se stranami $6$, $7$, $9$?
    \item Určete výšku rovnostranného trojúhelníku o straně $a$ a jeho obsah.
    \item Úhlopříčka obdélníku má délku $10$ a jedna strana $6$. Určete obsah obdélníku.
    \item V pravoúhlém trojúhelníku je $a = 6$ a $c_a = 3{,}6$. Určete $c$, $b$ a $v$.
  \end{enumerate}
\end{problem}

\begin{solution}
  \begin{enumerate}[label=(\alph*)]
    \item $49 + 576 = 625$, ano. $36 + 49 = 85\neq 81$, ne. Protože $85 > 81$, je úhel proti nejdelší straně ostrý a trojúhelník je ostroúhlý.
    \item $v = \sqrt{a^2 - \frac{a^2}{4}} = \frac{\sqrt{3}}{2}a$, $S = \frac{1}{2}a\cdot\frac{\sqrt{3}}{2}a = \frac{\sqrt{3}}{4}a^2$.
    \item Druhá strana $\sqrt{100 - 36} = 8$, obsah $48$.
    \item $c = \frac{a^2}{c_a} = 10$, $b = \sqrt{100 - 36} = 8$, $c_b = 6{,}4$, $v = \sqrt{3{,}6\cdot 6{,}4} = 4{,}8$.
  \end{enumerate}
\end{solution}

\subsection{Obsah trojúhelníku}

\begin{definitionbox}
  \begin{equation*}
    S = \frac{a\cdot v_a}{2} = \frac{1}{2}ab\sin\gamma = \rho\cdot s = \sqrt{s(s - a)(s - b)(s - c)},
  \end{equation*}
  kde $s = \frac{a + b + c}{2}$ je poloviční obvod a $\rho$ poloměr kružnice vepsané. Poslední vzorec se nazývá \textbf{Heronův vzorec}.
\end{definitionbox}

Vzorec $S = \rho s$ plyne z toho, že střed kružnice vepsané rozdělí trojúhelník na tři trojúhelníky s výškou $\rho$ a základnami $a$, $b$, $c$. Jejich obsahy dohromady dávají $\frac{\rho a + \rho b + \rho c}{2} = \rho s$.
\begin{center}
  \begin{tikzpicture}[scale=0.75]
    \coordinate (A) at (0,0); \coordinate (B) at (5,0); \coordinate (C) at (1.5,3); \coordinate (I) at (1.8722,1.1571);
    \fill[colorofwisdom!30] (A) -- (B) -- (I) -- cycle;
    \fill[colorofwisdom!12] (B) -- (C) -- (I) -- cycle;
    \fill[black!10] (C) -- (A) -- (I) -- cycle;
    \draw[thick] (A) -- (B) -- (C) -- cycle;
    \draw (I) -- (A) (I) -- (B) (I) -- (C);
    \draw[dashed] (I) circle (1.1571);
    \draw[colorofwisdom, thick] (I) -- (1.8722,0) node[midway, right, font=\scriptsize] {$\rho$};
    \node[below] at (2.5,0) {\small $c$}; \node[above right] at (3.25,1.5) {\small $a$}; \node[above left] at (0.75,1.5) {\small $b$};
  \end{tikzpicture}
\end{center}

\begin{example}
  Trojúhelník se stranami $13$, $14$, $15$ má $s = 21$ a podle Heronova vzorce
  \begin{equation*}
    S = \sqrt{21\cdot 8\cdot 7\cdot 6} = \sqrt{7056} = 84.
  \end{equation*}
  Poloměr kružnice vepsané je $\rho = \frac{S}{s} = 4$.
\end{example}

\section{Mnohoúhelníky}

\begin{definitionbox}
  \textbf{Konvexní $n$-úhelník} je mnohoúhelník, jehož každý vnitřní úhel je menší než $180^\circ$. \textbf{Úhlopříčka} spojuje dva nesousední vrcholy.
  \begin{itemize}
    \item Součet vnitřních úhlů konvexního $n$-úhelníku je $(n - 2)\cdot 180^\circ$.
    \item Počet úhlopříček je $\frac{n(n - 3)}{2}$.
    \item \textbf{Pravidelný} $n$-úhelník má všechny strany i úhly shodné. Každý jeho vnitřní úhel má velikost $\frac{(n - 2)\cdot 180^\circ}{n}$. Má kružnici opsanou i vepsanou se společným středem.
  \end{itemize}
\end{definitionbox}

\begin{proof}
  Z jednoho vrcholu vedeme všechny úhlopříčky. Ty rozdělí $n$-úhelník na $n - 2$ trojúhelníků, jejichž úhly dohromady tvoří přesně vnitřní úhly $n$-úhelníku. Součet je tedy $(n - 2)\cdot 180^\circ$.

  Z každého z $n$ vrcholů vede úhlopříčka do $n - 3$ vrcholů (ne do sebe ani do dvou sousedů). To dává $n(n - 3)$, ale každou úhlopříčku jsme započítali dvakrát (z obou konců).
\end{proof}

\begin{center}
  \begin{tikzpicture}[scale=1.2]
    \foreach \i/\c in {1/colorofwisdom!40, 2/colorofwisdom!20, 3/black!15, 4/black!5}
      \fill[\c] (90:1) -- ({90+60*\i}:1) -- ({90+60*(\i+1)}:1) -- cycle;
    \draw[thick] (90:1) \foreach \i in {1,...,5} { -- ({90+60*\i}:1) } -- cycle;
    \foreach \i in {2,3,4} \draw[colorofwisdom] (90:1) -- ({90+60*\i}:1);
    \node[font=\scriptsize, align=center] at (0,-1.4) {$6$-úhelník $=$ $4$ trojúhelníky,\\ $4\cdot 180^\circ = 720^\circ$};
  \end{tikzpicture}
  \hspace{1.2cm}
  \begin{tikzpicture}[scale=1.2]
    \draw[thick] (90:1) \foreach \i in {1,...,5} { -- ({90+60*\i}:1) } -- cycle;
    \foreach \i in {0,...,5} \foreach \j in {2,3} \draw[colorofwisdom!70] ({90+60*\i}:1) -- ({90+60*(\i+\j)}:1);
    \node[font=\scriptsize, align=center] at (0,-1.4) {všech $\frac{6\cdot 3}{2} = 9$ úhlopříček};
  \end{tikzpicture}
\end{center}

\begin{example}
  Pravidelný šestiúhelník má součet úhlů $720^\circ$, každý úhel $120^\circ$ a $9$ úhlopříček. Hlavní úhlopříčky ho rozdělí na $6$ rovnostranných trojúhelníků, takže pro stranu $a$ je jeho obsah $S = 6\cdot\frac{\sqrt{3}}{4}a^2 = \frac{3\sqrt{3}}{2}a^2$ a poloměr kružnice opsané je $r = a$. Proto mají včelí plástve tvar šestiúhelníků: ze všech pravidelných mnohoúhelníků, kterými lze vydláždit rovinu (trojúhelník, čtverec, šestiúhelník), uzavře šestiúhelník největší plochu při daném obvodu.
\end{example}

\subsection{Čtyřúhelníky}

\begin{definitionbox}
  \begin{center}
    \renewcommand{\arraystretch}{1.5}
    \begin{tabular}{|l|p{6.4cm}|l|}
      \hline
      Čtyřúhelník & Vlastnosti & Obsah \\
      \hline
      \textbf{rovnoběžník} & protější strany rovnoběžné a shodné, protější úhly shodné, úhlopříčky se půlí & $S = a\cdot v_a$ \\
      \textbf{obdélník} & rovnoběžník s pravými úhly, úhlopříčky shodné & $S = ab$ \\
      \textbf{kosočtverec} & rovnoběžník se shodnými stranami, úhlopříčky na sebe kolmé & $S = \frac{u_1u_2}{2}$ \\
      \textbf{čtverec} & obdélník i kosočtverec zároveň & $S = a^2 = \frac{u^2}{2}$ \\
      \textbf{lichoběžník} & právě jedna dvojice rovnoběžných stran (základny $a$, $c$) & $S = \frac{(a + c)v}{2}$ \\
      \textbf{deltoid} & dvě dvojice shodných sousedních stran, úhlopříčky na sebe kolmé & $S = \frac{u_1u_2}{2}$ \\
      \hline
    \end{tabular}
  \end{center}
\end{definitionbox}

\begin{center}
  \begin{tikzpicture}[scale=0.75]
    \begin{scope}
      \draw[thick, fill=colorofwisdom!12] (0,0) -- (2.2,0) -- (2.9,1.3) -- (0.7,1.3) -- cycle;
      \draw[gray, dashed] (0,0) -- (2.9,1.3) (2.2,0) -- (0.7,1.3);
      \node[font=\scriptsize] at (1.45,-0.4) {rovnoběžník};
    \end{scope}
    \begin{scope}[xshift=3.8cm]
      \draw[thick, fill=colorofwisdom!12] (0,0) rectangle (2.4,1.3);
      \draw[gray, dashed] (0,0) -- (2.4,1.3) (2.4,0) -- (0,1.3);
      \node[font=\scriptsize] at (1.2,-0.4) {obdélník};
    \end{scope}
    \begin{scope}[xshift=7.6cm]
      \draw[thick, fill=colorofwisdom!12] (0,0.75) -- (1.2,0) -- (2.4,0.75) -- (1.2,1.5) -- cycle;
      \draw[gray, dashed] (0,0.75) -- (2.4,0.75) (1.2,0) -- (1.2,1.5);
      \draw (1.2,0.75) rectangle (1.35,0.9);
      \node[font=\scriptsize] at (1.2,-0.4) {kosočtverec};
    \end{scope}
    \begin{scope}[yshift=-2.5cm]
      \draw[thick, fill=colorofwisdom!12] (0.4,0) rectangle (1.8,1.4);
      \draw[gray, dashed] (0.4,0) -- (1.8,1.4) (1.8,0) -- (0.4,1.4);
      \node[font=\scriptsize] at (1.1,-0.4) {čtverec};
    \end{scope}
    \begin{scope}[xshift=3.8cm, yshift=-2.5cm]
      \draw[thick, fill=colorofwisdom!12] (0,0) -- (2.6,0) -- (1.9,1.3) -- (0.6,1.3) -- cycle;
      \draw[gray, dashed] (0.6,1.3) -- (0.6,0) node[midway, left, font=\tiny] {$v$};
      \node[font=\tiny] at (1.3,0.15) {$a$}; \node[font=\tiny] at (1.25,1.45) {$c$};
      \node[font=\scriptsize] at (1.3,-0.4) {lichoběžník};
    \end{scope}
    \begin{scope}[xshift=7.6cm, yshift=-2.5cm]
      \draw[thick, fill=colorofwisdom!12] (1.2,1.5) -- (0.5,0.95) -- (1.2,0) -- (1.9,0.95) -- cycle;
      \draw[gray, dashed] (0.5,0.95) -- (1.9,0.95) (1.2,0) -- (1.2,1.5);
      \node[font=\scriptsize] at (1.2,-0.4) {deltoid};
    \end{scope}
  \end{tikzpicture}
\end{center}

Čtverec je zároveň obdélník, kosočtverec i rovnoběžník. Obdélník ale obecně čtverec není. Vztahy mezi čtyřúhelníky jsou vlastně inkluze množin: množina čtverců je průnikem množin obdélníků a kosočtverců. Opět se nám vrací kapitola 2.

\begin{problem}
  \begin{enumerate}[label=(\alph*)]
    \item Kosočtverec má úhlopříčky $6$ a $8$. Určete jeho stranu a obsah.
    \item Lichoběžník má základny $10$ a $4$ a výšku $5$. Jaký je jeho obsah?
    \item Rovnoramenný lichoběžník má základny $12$ a $6$ a ramena $5$. Určete výšku a obsah.
    \item Kolik úhlopříček má osmiúhelník a jaký je součet jeho vnitřních úhlů?
    \item Jaký pravidelný mnohoúhelník má vnitřní úhel $144^\circ$?
  \end{enumerate}
\end{problem}

\begin{solution}
  \begin{enumerate}[label=(\alph*)]
    \item Úhlopříčky se půlí a jsou kolmé, strana je přepona trojúhelníku s odvěsnami $3$ a $4$, tedy $5$. Obsah $\frac{6\cdot 8}{2} = 24$.
    \item $\frac{(10 + 4)\cdot 5}{2} = 35$.
    \item Rameno je přepona pravoúhlého trojúhelníku s odvěsnou $\frac{12 - 6}{2} = 3$, takže $v = 4$ a $S = \frac{18\cdot 4}{2} = 36$.
    \item $\frac{8\cdot 5}{2} = 20$ úhlopříček, součet úhlů $1080^\circ$.
    \item $\frac{(n - 2)180}{n} = 144$, tedy $180n - 360 = 144n$ a $n = 10$. Desetiúhelník.
  \end{enumerate}
\end{solution}

\section{Kružnice a kruh}

\begin{definitionbox}
  \textbf{Kružnice} $k(S, r)$ je množina všech bodů roviny, které mají od středu $S$ vzdálenost $r > 0$. \textbf{Kruh} $K(S, r)$ je množina bodů, které mají od $S$ vzdálenost nejvýše $r$.

  \vspace{0.2cm}

  \textbf{Vzájemná poloha přímky a kružnice} (kde $d$ je vzdálenost středu od přímky):
  \begin{itemize}
    \item $d > r$: \textbf{vnější přímka} (žádný společný bod),
    \item $d = r$: \textbf{tečna} (jeden společný bod, bod dotyku). Tečna je \textbf{kolmá} na poloměr v bodě dotyku.
    \item $d < r$: \textbf{sečna} (dva společné body). Úsečka mezi nimi je \textbf{tětiva}.
  \end{itemize}

  \textbf{Obvod a obsah}:
  \begin{equation*}
    o = 2\pi r,\qquad S = \pi r^2.
  \end{equation*}
  \textbf{Oblouk} příslušný středovému úhlu $\alpha$ (ve stupních) má délku $l = \frac{\pi r\alpha}{180^\circ}$ a \textbf{kruhová výseč} obsah $S_v = \frac{\pi r^2\alpha}{360^\circ}$. V radiánech prostě $l = r\alpha$ a $S_v = \frac{1}{2}r^2\alpha$.
\end{definitionbox}

\begin{center}
  \begin{tikzpicture}[scale=0.9]
    \foreach \d/\n/\x in {1.6/vnější přímka/0, 1/tečna/3.6, 0.45/sečna/7.2} {
      \begin{scope}[xshift=\x cm]
        \draw[thick] (0,0) circle (1);
        \filldraw (0,0) circle (1.2pt);
        \draw[very thick, colorofwisdom] (-1.5,-\d) -- (1.5,-\d);
        \draw[dashed] (0,0) -- (0,-\d) node[midway, right, font=\tiny] {$d$};
        \node[font=\small] at (0,-2.05) {\n};
      \end{scope}
    }
    \filldraw[colorofwisdom] (3.6,-1) circle (1.5pt);
    \draw (3.6,-1) rectangle (3.75,-0.85);
    \filldraw[colorofwisdom] ({7.2-0.893},-0.45) circle (1.5pt) ({7.2+0.893},-0.45) circle (1.5pt);
    \node[font=\tiny] at (7.2,-0.3) {tětiva};
  \end{tikzpicture}

  \vspace{0.3cm}

  \begin{tikzpicture}[scale=1.3]
    \fill[colorofwisdom!25] (0,0) -- (0:1) arc (0:70:1) -- cycle;
    \draw[thick] (0,0) circle (1);
    \draw[thick] (1,0) -- (0,0) -- (70:1);
    \draw[very thick, colorofwisdom] (0:1) arc (0:70:1);
    \draw (0.3,0) arc (0:70:0.3); \node at (35:0.45) {\scriptsize $\alpha$};
    \node[colorofwisdom] at (35:1.2) {\small $l$};
    \node at (35:0.72) {\scriptsize $S_v$};
    \node[below] at (0.5,0) {\scriptsize $r$};
    \node[font=\scriptsize, align=center] at (2.6,0.2) {výseč je stejný díl kruhu,\\ jako je $\alpha$ díl plného úhlu:\\ $S_v = \frac{\alpha}{360^\circ}\cdot\pi r^2$};
  \end{tikzpicture}
\end{center}

Číslo $\pi$ je definováno jako poměr obvodu kružnice k jejímu průměru. Že je tento poměr stejný pro všechny kružnice, plyne z toho, že všechny kružnice jsou si podobné. Číslo $\pi$ je iracionální (dokázal Lambert v roce 1761), takže jeho desetinný rozvoj nikdy nekončí a neopakuje se. Rekord v počítání jeho cifer je dnes přes sto bilionů míst. Pro všechny praktické účely přitom stačí asi čtyřicet: s nimi lze spočítat obvod pozorovatelného vesmíru s přesností na průměr atomu vodíku.

\subsection{Středový a obvodový úhel}

\begin{definitionbox}
  Nechť $AB$ je tětiva kružnice se středem $S$ a $V$ je bod kružnice, který neleží na oblouku $AB$.
  \begin{itemize}
    \item \textbf{Středový úhel} je $\angle ASB$.
    \item \textbf{Obvodový úhel} je $\angle AVB$.
  \end{itemize}

  \textbf{Věta (o obvodovém a středovém úhlu)}: Obvodový úhel je polovinou středového úhlu příslušného ke stejnému oblouku. Všechny obvodové úhly nad stejným obloukem jsou tedy shodné.

  \textbf{Thaletova věta}: Pokud je $AB$ průměr kružnice, je každý obvodový úhel nad ním pravý. Neboli: všechny body $V$, pro které je úhel $AVB$ pravý, leží na kružnici s průměrem $AB$ (\textbf{Thaletova kružnice}).
\end{definitionbox}

\begin{proof}
  Dokážeme speciální případ, kdy střed $S$ leží na rameni $VB$. Trojúhelník $AVS$ je rovnoramenný ($|SA| = |SV| = r$), takže úhly při jeho základně jsou shodné, oba označme $\varphi$. Úhel $ASB$ je vnějším úhlem trojúhelníku $AVS$ a je roven součtu dvou vnitřních úhlů, které k němu nepřiléhají\footnote{Vedlejší úhel k $\angle ASB$ je třetí úhel trojúhelníku $AVS$, tedy $180^\circ - 2\varphi$. Odtud $|\angle ASB| = 180^\circ - (180^\circ - 2\varphi) = 2\varphi$.}, tedy $2\varphi$. Obvodový úhel je $\varphi$, středový $2\varphi$. Obecný případ dostaneme rozdělením úhlu průměrem jdoucím z $V$ na dva takové speciální případy (a jejich sečtením nebo odečtením).

  Thaletova věta je speciální případ: středový úhel nad průměrem je přímý, tedy $180^\circ$, a obvodový je jeho polovina, $90^\circ$.
\end{proof}

\begin{center}
  \begin{tikzpicture}[scale=1.6]
    \draw[thick] (0,0) circle (1);
    \coordinate (S) at (0,0);
    \coordinate (A) at (-30:1); \coordinate (B) at (210:1); \coordinate (V) at (100:1);
    \filldraw (S) circle (0.8pt) node[below] {\small $S$};
    \draw[thick, colorofwisdom] (A) -- (S) -- (B);
    \draw[thick] (A) -- (V) -- (B);
    \node[right] at (A) {\small $A$}; \node[left] at (B) {\small $B$}; \node[above] at (V) {\small $V$};
    \draw[colorofwisdom] ($(S)+(-30:0.25)$) arc (-30:-150:0.25);
    \node[colorofwisdom] at (0,-0.4) {\small $2\varphi$};
    \node at ($(V)+(-0.03,-0.33)$) {\small $\varphi$};
  \end{tikzpicture}
  \hspace{1.5cm}
  \begin{tikzpicture}[scale=1.6]
    \draw[thick] (0,0) circle (1);
    \coordinate (A) at (-1,0); \coordinate (B) at (1,0); \coordinate (V) at (60:1);
    \draw[thick, colorofwisdom] (A) -- (B);
    \draw[thick] (A) -- (V) -- (B);
    \filldraw (0,0) circle (0.8pt);
    \node[left] at (A) {\small $A$}; \node[right] at (B) {\small $B$}; \node[above] at (V) {\small $V$};
    \coordinate (P1) at ($(V)!0.15!(A)$); \coordinate (P2) at ($(V)!0.15!(B)$);
    \draw (P1) -- ($(P1)+(P2)-(V)$) -- (P2);
    \node[below] at (0,-1.1) {\small Thaletova kružnice};
  \end{tikzpicture}
\end{center}

\begin{problem}
  \begin{enumerate}[label=(\alph*)]
    \item Kružnice má poloměr $6$. Určete délku oblouku a obsah výseče se středovým úhlem $60^\circ$.
    \item Bod $P$ má od středu kružnice o poloměru $5$ vzdálenost $13$. Jak dlouhé jsou tečny z bodu $P$ ke kružnici?
    \item Tětiva kružnice o poloměru $10$ má délku $16$. Jak daleko je od středu?
    \item Obvodový úhel nad tětivou $AB$ má velikost $35^\circ$. Jaká je velikost středového úhlu? A obvodového úhlu na druhé straně tětivy?
    \item Kolikrát se zvětší obsah kruhu, když zdvojnásobíme poloměr?
  \end{enumerate}
\end{problem}

\begin{solution}
  \begin{enumerate}[label=(\alph*)]
    \item $l = \frac{\pi\cdot 6\cdot 60}{180} = 2\pi$, $S_v = \frac{\pi\cdot 36\cdot 60}{360} = 6\pi$.
    \item Tečna je kolmá na poloměr, takže tvoří s poloměrem a úsečkou $SP$ pravoúhlý trojúhelník: $t = \sqrt{169 - 25} = 12$.
    \item Kolmice ze středu na tětivu ji půlí: $d = \sqrt{100 - 64} = 6$.
    \item Středový úhel $70^\circ$. Obvodové úhly na opačné straně tětivy příslušejí k doplňkovému oblouku se středovým úhlem $290^\circ$, mají tedy velikost $145^\circ = 180^\circ - 35^\circ$. (Odtud mimochodem plyne, že protější úhly čtyřúhelníku vepsaného do kružnice dávají dohromady $180^\circ$.)
    \item Čtyřikrát. Obsah roste s druhou mocninou, jako u všech podobných útvarů.
  \end{enumerate}
\end{solution}

\section{Odbočka --- Pátý postulát}

Eukleidés postavil celou geometrii na pěti postulátech. První čtyři jsou krátké a samozřejmé (například `všechny pravé úhly jsou si rovny'). Pátý je dlouhý a nepříjemný. V moderní podobě zní: \textit{bodem mimo přímku vede právě jedna rovnoběžka s touto přímkou}. Matematici měli dva tisíce let pocit, že tohle přece musí jít dokázat z ostatních čtyř, a pokoušeli se o to. Nikomu se to nepovedlo.

Maďarský matematik Farkas Bolyai se o důkaz pokoušel celý život. Když zjistil, že se do stejného problému pustil i jeho syn János, napsal mu v roce 1820 dopis:
\begin{center}
  \begin{minipage}{0.85\textwidth}
    \textit{`Pro Boha tě prosím, nech toho. Boj se toho stejně jako smyslných vášní, protože ti to také může vzít všechen čas a připravit tě o zdraví, klid mysli a štěstí v životě. Tahle bezedná tma by mohla pohltit tisíc newtonovských velikánů. Na zemi se nikdy nerozjasní.'}
  \end{minipage}
\end{center}
János neposlechl. Místo důkazu ale objevil, že pátý postulát dokázat nejde, protože existuje úplně jiná, stejně bezesporná geometrie, ve které neplatí: bodem mimo přímku v ní vede nekonečně mnoho rovnoběžek a součet úhlů v trojúhelníku je menší než $180^\circ$. Tutéž geometrii nezávisle objevil i Rus Nikolaj Lobačevskij. Když otec poslal synovu práci svému příteli Gaussovi, ten odpověděl, že ji nemůže pochválit, protože by tím chválil sám sebe. Na totéž prý přišel už před lety, jen se to bál zveřejnit. János se z toho nikdy nevzpamatoval a už nikdy nic dalšího nepublikoval.

Takovým geometriím říkáme \textbf{neeukleidovské}. Na povrchu koule (kde `přímkami' jsou hlavní kružnice, jako rovník) žádné rovnoběžky neexistují a součet úhlů v trojúhelníku je větší než $180^\circ$. Trojúhelník se dvěma vrcholy na rovníku a jedním na severním pólu může mít všechny tři úhly pravé. Einsteinova obecná teorie relativity pak ukázala, že náš vesmír sám není eukleidovský: hmota zakřivuje prostor. GPS ve vašem telefonu musí tohle zakřivení započítávat, jinak by se každý den mýlila o několik kilometrů.

\section{Shrnutí a otázky k ústní zkoušce}

\begin{problem}
  Odpovězte vlastními slovy:
  \begin{enumerate}
    \item Jaké jsou vzájemné polohy dvou přímek v rovině? Co jsou vedlejší, vrcholové, souhlasné a střídavé úhly?
    \item Dokažte větu o součtu úhlů v trojúhelníku. Na kterém předpokladu stojí?
    \item Co je těžnice, výška, osa strany a osa úhlu? Jaké významné body trojúhelníku znáte a proč se příslušné přímky protínají v jednom bodě?
    \item Vyslovte věty o shodnosti a podobnosti trojúhelníků. Proč musí být ve větě Ssu úhel proti větší straně?
    \item Vyslovte a dokažte Pythagorovu větu a Eukleidovy věty. Platí obrácená Pythagorova věta?
    \item Jaké znáte vzorce pro obsah trojúhelníku?
    \item Odvoďte vzorec pro součet vnitřních úhlů a počet úhlopříček konvexního $n$-úhelníku.
    \item Popište druhy čtyřúhelníků a jejich vlastnosti. Jak se počítají jejich obsahy?
    \item Jaké jsou vzájemné polohy přímky a kružnice? Jakou vlastnost má tečna?
    \item Vyslovte a dokažte větu o obvodovém a středovém úhlu. Jak z ní plyne Thaletova věta?
    \item Jak se počítá délka oblouku a obsah kruhové výseče?
  \end{enumerate}
\end{problem}

\koneckapitoly

\chapter{Zobrazení a konstrukční úlohy}
\label{kap:zobrazeni}
\mathquote{Symmetry, as wide or as narrow as you may define its meaning, is one idea by which man through the ages has tried to comprehend and create order, beauty, and perfection.}{Hermann Weyl}

V předchozí kapitole jsme zkoumali útvary samotné. Teď se podíváme na to, jak s nimi hýbat: překlápět je, posouvat, otáčet a zvětšovat. Takovým operacím se říká geometrická zobrazení. Ve druhé části kapitoly pak využijeme všechno, co zatím o geometrii víme, k řešení konstrukčních úloh, tedy úloh, ve kterých máme útvar sestrojit jen pomocí pravítka a kružítka.

\section{Geometrická zobrazení}\label{sec:geom-zobrazeni}

\begin{definitionbox}
  \textbf{Zobrazení v rovině} je předpis, který každému bodu $X$ roviny přiřadí právě jeden bod $X'$ roviny. Bod $X$ se nazývá \textbf{vzor}, bod $X'$ \textbf{obraz}. Je to vlastně funkce z kapitoly~\ref{kap:funkce}, jejímž definičním oborem i oborem hodnot je rovina.

  \vspace{0.2cm}

  \textbf{Samodružný bod} je bod, který se zobrazí sám na sebe: $X' = X$. \textbf{Samodružný útvar} je útvar, který se zobrazí sám na sebe (jeho jednotlivé body ale samodružné být nemusí).
\end{definitionbox}

\section{Shodná zobrazení}\label{sec:shodna-zobrazeni}

\begin{definitionbox}
  \textbf{Shodné zobrazení} (shodnost) je zobrazení, které zachovává vzdálenosti: pro každé dva body $X, Y$ a jejich obrazy $X', Y'$ platí $|X'Y'| = |XY|$.

  \vspace{0.2cm}

  Shodné zobrazení převádí přímku na přímku, úsečku na stejně dlouhou úsečku, úhel na shodný úhel, rovnoběžky na rovnoběžky a každý útvar na útvar s ním shodný.
\end{definitionbox}

\subsection{Osová souměrnost}\label{sec:osova-soumernost}

\begin{definitionbox}
  \textbf{Osová souměrnost} $O(o)$ s osou $o$ přiřadí bodu $X\notin o$ bod $X'$ tak, že přímka $o$ je osou úsečky $XX'$. Body osy $o$ jsou samodružné.

  \vspace{0.2cm}

  Osová souměrnost \textbf{mění orientaci} útvarů: pokud jsou vrcholy trojúhelníku $ABC$ seřazeny proti směru hodinových ručiček, jsou vrcholy obrazu $A'B'C'$ po směru. Je to jako zrcadlo, ve kterém se levá ruka stane pravou.
\end{definitionbox}

\begin{center}
  \begin{tikzpicture}[scale=0.9]
    \draw[thick, colorofwisdom] (2,-0.5) -- (2,3.5) node[above] {\small $o$};
    \coordinate (A) at (0,0); \coordinate (B) at (1.4,0.6); \coordinate (C) at (0.4,2.2);
    \coordinate (A') at (4,0); \coordinate (B') at (2.6,0.6); \coordinate (C') at (3.6,2.2);
    \draw[thick] (A) -- (B) -- (C) -- cycle;
    \draw[thick] (A') -- (B') -- (C') -- cycle;
    \draw[dashed, gray] (A) -- (A') (B) -- (B') (C) -- (C');
    \node[below left] at (A) {\small $A$}; \node[below] at (B) {\small $B$}; \node[above] at (C) {\small $C$};
    \node[below right] at (A') {\small $A'$}; \node[below] at (B') {\small $B'$}; \node[above] at (C') {\small $C'$};
  \end{tikzpicture}
\end{center}

Útvar je \textbf{osově souměrný}, pokud existuje osová souměrnost, ve které je samodružný. Rovnostranný trojúhelník má $3$ osy souměrnosti, čtverec $4$, pravidelný $n$-úhelník $n$ a kružnice nekonečně mnoho (každá přímka procházející středem). Lidský obličej má přibližně jednu, ale jen přibližně.\footnote{Pokud chcete zažít opravdový horor, vyfoťte se zepředu a vytvořte si dvě fotky: jednu složenou ze dvou levých polovin obličeje, druhou ze dvou pravých. Ani jedna nebudete vy.}

\subsection{Středová souměrnost}

\begin{definitionbox}
  \textbf{Středová souměrnost} $S(S)$ se středem $S$ přiřadí bodu $X$ bod $X'$ tak, že $S$ je středem úsečky $XX'$. Jediným samodružným bodem je střed $S$. Samodružné jsou všechny přímky procházející bodem $S$.

  \vspace{0.2cm}

  Středová souměrnost \textbf{zachovává orientaci} a přímku zobrazí na rovnoběžnou přímku. Je to totéž co otočení o $180^\circ$.
\end{definitionbox}

\begin{center}
  \begin{tikzpicture}[scale=0.9]
    \coordinate (S) at (0,0);
    \coordinate (A) at (-2.4,-0.6); \coordinate (B) at (-1,-1.2); \coordinate (C) at (-1.5,0.5);
    \draw[thick] (A) -- (B) -- (C) -- cycle;
    \draw[thick, colorofwisdom] ($(S)!-1!(A)$) -- ($(S)!-1!(B)$) -- ($(S)!-1!(C)$) -- cycle;
    \draw[dashed, gray] (A) -- ($(S)!-1!(A)$) (B) -- ($(S)!-1!(B)$) (C) -- ($(S)!-1!(C)$);
    \filldraw (S) circle (1.5pt) node[below] {\small $S$};
    \node[left] at (A) {\small $A$}; \node[below] at (B) {\small $B$}; \node[above] at (C) {\small $C$};
    \node[right] at ($(S)!-1!(A)$) {\small $A'$}; \node[above] at ($(S)!-1!(B)$) {\small $B'$}; \node[below] at ($(S)!-1!(C)$) {\small $C'$};
    \node[font=\scriptsize, align=left, right] at (3,0) {$S$ je středem každé\\ úsečky $XX'$; trojúhelník\\ se `otočí vzhůru nohama'};
  \end{tikzpicture}
\end{center}

Středově souměrný je například rovnoběžník (střed je průsečík úhlopříček), kružnice, elipsa nebo graf liché funkce (střed v počátku, viz kapitola~\ref{kap:funkce}). Trojúhelník středově souměrný není nikdy.

\subsection{Posunutí a otočení}\label{sec:posunuti-otoceni}

\begin{definitionbox}
  \textbf{Posunutí} (translace) $T(\vec{u})$ o vektor $\vec{u}$ posune každý bod o tentýž vektor: $X' = X + \vec{u}$. Pokud $\vec{u}\neq\vec{0}$, nemá žádný samodružný bod. Samodružné jsou přímky rovnoběžné s vektorem $\vec{u}$.

  \vspace{0.2cm}

  \textbf{Otočení} (rotace) $R(S, \varphi)$ se středem $S$ o orientovaný úhel $\varphi$ přiřadí bodu $X\neq S$ bod $X'$ tak, že $|SX'| = |SX|$ a orientovaný úhel $XSX'$ má velikost $\varphi$. Jediný samodružný bod je střed $S$ (pokud $\varphi$ není násobek $360^\circ$).
\end{definitionbox}

\begin{center}
  \begin{tikzpicture}[scale=0.85]
    \begin{scope}
      \coordinate (A) at (0,0); \coordinate (B) at (1.2,0.2); \coordinate (C) at (0.5,1.2);
      \draw[thick] (A) -- (B) -- (C) -- cycle;
      \draw[thick] ($(A)+(2.5,0.8)$) -- ($(B)+(2.5,0.8)$) -- ($(C)+(2.5,0.8)$) -- cycle;
      \draw[->, thick, colorofwisdom] (0.5,-0.6) -- (3,0.2) node[midway, below right] {\small $\vec{u}$};
      \draw[dashed, gray] (C) -- ($(C)+(2.5,0.8)$);
      \node at (1.9,-1.4) {\small posunutí};
    \end{scope}
    \begin{scope}[xshift=7.5cm]
      \coordinate (S) at (0,0);
      \filldraw[colorofwisdom] (S) circle (2pt) node[below left] {\small $S$};
      \coordinate (A) at (1.5,0.2); \coordinate (B) at (2.5,0.4); \coordinate (C) at (1.9,1.2);
      \draw[thick] (A) -- (B) -- (C) -- cycle;
      \draw[thick, rotate around={70:(S)}] (A) -- (B) -- (C) -- cycle;
      \draw[dashed, gray] (S) -- (B);
      \draw[dashed, gray, rotate around={70:(S)}] (S) -- (B);
      \draw[->, colorofwisdom] (0.8,0.13) arc (9:79:0.8);
      \node[colorofwisdom] at (0.85,0.85) {\small $\varphi$};
      \node at (1.2,-1.4) {\small otočení};
    \end{scope}
  \end{tikzpicture}
\end{center}

\begin{note}
  Všechna shodná zobrazení v rovině jsou buď identita, osová souměrnost, posunutí, otočení, nebo tzv. posunuté zrcadlení (osová souměrnost složená s posunutím ve směru osy). Navíc každé z nich lze složit nejvýše ze tří osových souměrností. Osová souměrnost je tedy základní stavební kámen všech shodností:
  \begin{itemize}
    \item složením dvou osových souměrností s \textbf{rovnoběžnými} osami vzdálenými $d$ vznikne \textbf{posunutí} o $2d$ kolmo na osy,
    \item složením dvou osových souměrností s osami \textbf{různoběžnými} svírajícími úhel $\alpha$ vznikne \textbf{otočení} kolem jejich průsečíku o $2\alpha$.
  \end{itemize}
  Speciálně složení dvou souměrností s kolmými osami je otočení o $180^\circ$, tedy středová souměrnost. Proto když se podíváte do dvou zrcadel postavených do pravého úhlu, vidíte se tak, jak vás vidí ostatní, a ne zrcadlově.
\end{note}

\begin{center}
  \begin{tikzpicture}[scale=0.85]
    \draw[thick, colorofwisdom] (1,-0.6) -- (1,2.4) node[above] {\small $o_1$};
    \draw[thick, colorofwisdom] (2.5,-0.6) -- (2.5,2.4) node[above] {\small $o_2$};
    \coordinate (A) at (0,0); \coordinate (B) at (0.7,0.3); \coordinate (C) at (0.2,1.3);
    \draw[thick, fill=black!10] (A) -- (B) -- (C) -- cycle;
    \draw[thick, dashed, fill=black!5] (2,0) -- (1.3,0.3) -- (1.8,1.3) -- cycle;
    \draw[thick, fill=colorofwisdom!20] (3,0) -- (3.7,0.3) -- (3.2,1.3) -- cycle;
    \draw[->, gray] (0.2,1.6) -- (1.8,1.6); \draw[->, gray] (1.8,1.9) -- (3.2,1.9);
    \draw[<->] (1,-0.45) -- (2.5,-0.45) node[midway, below, font=\scriptsize] {$d$};
    \draw[->, very thick, colorofwisdom] (0,-1) -- (3,-1) node[midway, below, font=\scriptsize] {posunutí o $2d$};
  \end{tikzpicture}
  \hspace{1cm}
  \begin{tikzpicture}[scale=0.85]
    \coordinate (P) at (0,0);
    \draw[thick, colorofwisdom] (P) -- (0:3) node[right] {\small $o_1$};
    \draw[thick, colorofwisdom] (P) -- (35:3) node[above right] {\small $o_2$};
    \draw (0.8,0) arc (0:35:0.8); \node at (17:1.05) {\scriptsize $\alpha$};
    \coordinate (X) at (-20:2.2); \coordinate (X1) at (20:2.2); \coordinate (X2) at (50:2.2);
    \filldraw (X) circle (1.5pt) node[below right] {\small $X$};
    \filldraw[gray] (X1) circle (1.5pt) node[right] {\small $X_1$};
    \filldraw[colorofwisdom] (X2) circle (1.5pt) node[above] {\small $X'$};
    \draw[dashed, gray] (X) -- (X1) -- (X2);
    \draw[->, colorofwisdom] (-20:1.6) arc (-20:50:1.6);
    \node[font=\scriptsize, colorofwisdom] at (100:0.8) {otočení o $2\alpha$};
    \filldraw (P) circle (1.5pt);
  \end{tikzpicture}
\end{center}

\begin{problem}
  \begin{enumerate}[label=(\alph*)]
    \item Kolik os souměrnosti mají písmena A, H, N, O, S, X? Která z nich jsou středově souměrná?
    \item Určete samodružné body a samodružné přímky osové souměrnosti.
    \item Co vznikne složením dvou středových souměrností se středy $S_1\neq S_2$?
    \item Které pravidelné mnohoúhelníky jsou středově souměrné?
  \end{enumerate}
\end{problem}

\begin{solution}
  \begin{enumerate}[label=(\alph*)]
    \item A: $1$, H: $2$, N: $0$, O: $2$ (v typickém písmu, které není dokonale kulaté), S: $0$, X: $2$. Středově souměrná jsou H, N, O, S, X.
    \item Samodružné body: všechny body osy. Samodružné přímky: osa sama a všechny přímky na ni kolmé.
    \item Posunutí o vektor $2\overrightarrow{S_1S_2}$. Každá středová souměrnost je složením dvou osových s kolmými osami, a pokud osy zvolíme vhodně, dvě z nich se vyruší a zbydou dvě rovnoběžné osy.
    \item Ty se sudým počtem vrcholů.
  \end{enumerate}
\end{solution}

\section{Podobná zobrazení}\label{sec:podobna-zobrazeni}

\begin{definitionbox}
  \textbf{Podobné zobrazení} s koeficientem $k > 0$ je zobrazení, pro které platí $|X'Y'| = k\cdot|XY|$ pro každé dva body. Pro $k = 1$ je to shodnost, pro $k > 1$ zvětšení, pro $k < 1$ zmenšení. Podobné zobrazení zachovává úhly a poměry délek.

  \vspace{0.2cm}

  \textbf{Stejnolehlost} (homotetie) $H(S, \kappa)$ se středem $S$ a koeficientem $\kappa\neq 0$ přiřadí bodu $X$ bod $X'$ na přímce $SX$ takový, že
  \begin{equation*}
    \overrightarrow{SX'} = \kappa\cdot\overrightarrow{SX}.
  \end{equation*}
  \begin{itemize}
    \item Pro $\kappa > 0$ leží $X'$ na polopřímce $\overrightarrow{SX}$, pro $\kappa < 0$ na polopřímce opačné.
    \item Stejnolehlost je podobné zobrazení s koeficientem $k = |\kappa|$.
    \item Přímku zobrazí na \textbf{rovnoběžnou} přímku.
    \item Pro $\kappa = -1$ je to středová souměrnost.
  \end{itemize}
\end{definitionbox}

\begin{center}
  \begin{tikzpicture}[scale=0.8]
    \coordinate (S) at (0,0);
    \coordinate (A) at (2,0.3); \coordinate (B) at (3,1.2); \coordinate (C) at (1.6,1.6);
    \filldraw[colorofwisdom] (S) circle (2pt) node[left] {\small $S$};
    \draw[thick] (A) -- (B) -- (C) -- cycle;
    \draw[thick, colorofwisdom] ($(S)!2!(A)$) -- ($(S)!2!(B)$) -- ($(S)!2!(C)$) -- cycle;
    \draw[dashed, gray] (S) -- ($(S)!2!(A)$) (S) -- ($(S)!2!(B)$) (S) -- ($(S)!2!(C)$);
    \node[below] at (A) {\small $A$}; \node[right] at (B) {\small $B$}; \node[above left] at (C) {\small $C$};
    \node[below] at ($(S)!2!(A)$) {\small $A'$}; \node[right] at ($(S)!2!(B)$) {\small $B'$}; \node[above left] at ($(S)!2!(C)$) {\small $C'$};
    \node at (3,-1) {\small $H(S, 2)$};
  \end{tikzpicture}
\end{center}

\begin{example}
  Stejnolehlost se skvěle hodí ke konstrukci útvarů vepsaných do jiných útvarů. Úloha: do trojúhelníku $ABC$ vepište čtverec tak, aby jedna jeho strana ležela na straně $AB$ a zbylé dva vrcholy na stranách $AC$ a $BC$.

  Řešení: Nejdřív sestrojíme \textit{pomocný} čtverec $K_0L_0M_0N_0$, který má stranu $K_0L_0$ na $AB$ a vrchol $N_0$ na $AC$. Stačí zvolit libovolný bod $N_0$ na $AC$, spustit z něj kolmici na $AB$ a čtverec dokreslit. Jeho vrchol $M_0$ ale obecně na straně $BC$ ležet nebude.

  Teď si všimneme, že všechny takové čtverce (se stranou na $AB$ a vrcholem na $AC$) jsou navzájem stejnolehlé se středem $A$. Jejich vrcholy $M$ proto leží na jedné polopřímce $\overrightarrow{AM_0}$. Hledaný vrchol $M$ je tedy průsečík polopřímky $\overrightarrow{AM_0}$ se stranou $BC$. Z něj už čtverec dokreslíme snadno. Pomocný čtverec jsme `nafoukli' stejnolehlostí přesně na správnou velikost.
  \begin{center}
    \begin{tikzpicture}[scale=1.15]
      \coordinate (A) at (0,0); \coordinate (B) at (5,0); \coordinate (C) at (1.5,3);
      \draw[thick] (A) -- (B) -- (C) -- cycle;
      \draw[gray, dashed, fill=black!8] (0.5,0) rectangle (1.5,1);
      \node[font=\tiny, gray] at (1,0.5) {pomocný};
      \draw[colorofwisdom!70, ->] (A) -- (3.3,2.2);
      \fill[colorofwisdom!25] (0.9375,0) rectangle (2.8125,1.875);
      \draw[very thick, colorofwisdom] (0.9375,0) rectangle (2.8125,1.875);
      \filldraw (1.5,1) circle (1.2pt) node[right, font=\tiny] {$M_0$};
      \filldraw[colorofwisdom] (2.8125,1.875) circle (1.5pt) node[right, font=\scriptsize] {$M$};
      \node[below left] at (A) {\small $A$}; \node[below right] at (B) {\small $B$}; \node[above] at (C) {\small $C$};
    \end{tikzpicture}
  \end{center}
\end{example}

\section{Konstrukční úlohy}\label{sec:konstrukcni-ulohy}

Konstrukční úloha je úloha, ve které máme sestrojit útvar daných vlastností pouze pomocí pravítka (bez stupnice) a kružítka. Je to tradice, která sahá až k Eukleidovi. Proč jen pravítko a kružítko? Protože to jsou nástroje odpovídající dvěma nejjednodušším křivkám, přímce a kružnici. A protože Řekové měli rádi omezení.

\begin{definitionbox}
  \textbf{Řešení konstrukční úlohy} má čtyři části:
  \begin{enumerate}
    \item \textbf{Rozbor}: Předpokládáme, že úloha je vyřešená, nakreslíme si náčrtek a hledáme vztahy, ze kterých lze útvar sestrojit. Zvýrazníme zadané prvky.
    \item \textbf{Postup konstrukce}: Krok za krokem zapíšeme, jak útvar sestrojíme, ideálně symbolicky.
    \item \textbf{Konstrukce}: Samotné narýsování.
    \item \textbf{Diskuze}: Kolik má úloha řešení a za jakých podmínek? (Například trojúhelník ze stran $1$, $2$, $5$ neexistuje.)
  \end{enumerate}
\end{definitionbox}

\begin{definitionbox}
  \textbf{Symbolický zápis}: Používají se tyto zkratky:
  \begin{itemize}
    \item $AB$; $|AB| = 5\ \si{cm}$ \quad (sestrojení úsečky),
    \item $k$; $k(S; 3\ \si{cm})$ \quad (kružnice se středem $S$ a poloměrem $3\ \si{cm}$),
    \item $\overrightarrow{AX}$; $|\angle BAX| = 60^\circ$ \quad (polopřímka svírající s $AB$ úhel $60^\circ$),
    \item $C$; $C\in k\cap\overrightarrow{AX}$ \quad (bod jako průsečík),
    \item $\triangle ABC$ \quad (výsledný trojúhelník).
  \end{itemize}
\end{definitionbox}

\subsection{Množiny bodů dané vlastnosti}\label{sec:mnoziny-bodu}

Jádrem většiny konstrukcí je hledání bodu, který má splňovat dvě podmínky zároveň. Každá podmínka určí nějakou množinu bodů a hledaný bod leží v jejich průniku. Opět množiny, opět průnik, opět kapitola 2.

\begin{definitionbox}
  \begin{center}
    \renewcommand{\arraystretch}{1.4}
    \begin{tabular}{|p{6.5cm}|p{6.5cm}|}
      \hline
      Množina bodů, které\ldots & je \\
      \hline
      mají od bodu $S$ vzdálenost $r$ & kružnice $k(S; r)$ \\
      mají stejnou vzdálenost od bodů $A$ a $B$ & osa úsečky $AB$ \\
      mají od přímky $p$ vzdálenost $d$ & dvojice rovnoběžek s $p$ ve vzdálenosti $d$ \\
      mají stejnou vzdálenost od dvou různoběžek & dvojice os úhlů, které přímky svírají \\
      z nichž je úsečka $AB$ vidět pod pravým úhlem & Thaletova kružnice nad $AB$ (bez bodů $A$, $B$) \\
      z nichž je úsečka $AB$ vidět pod úhlem $\alpha$ & dva kruhové oblouky nad $AB$ souměrné podle $AB$ (podle věty o obvodovém úhlu) \\
      \hline
    \end{tabular}
  \end{center}
\end{definitionbox}

\begin{center}
  \begin{tikzpicture}[scale=0.62]
    \begin{scope}
      \coordinate (A) at (0,0); \coordinate (B) at (2,0);
      \draw[thick] (A) -- (B); \filldraw (A) circle (1.5pt) node[below] {\tiny $A$}; \filldraw (B) circle (1.5pt) node[below] {\tiny $B$};
      \draw[thick, colorofwisdom] (1,-1.3) -- (1,1.7);
      \filldraw[colorofwisdom] (1,1.1) circle (1.5pt); \draw[dashed, gray] (A) -- (1,1.1) -- (B);
      \node[font=\scriptsize, align=center] at (1,-1.9) {osa úsečky};
    \end{scope}
    \begin{scope}[xshift=4cm]
      \draw[thick] (0,0) -- (2.4,0) (0,0) -- (50:2.4);
      \draw[thick, colorofwisdom] (0,0) -- (25:2.5);
      \coordinate (P) at (25:1.6); \filldraw[colorofwisdom] (P) circle (1.5pt);
      \draw[dashed, gray] (P) -- ({1.6*cos(25)},0) (P) -- ($(0,0)!(P)!(50:2.4)$);
      \node[font=\scriptsize] at (1.1,-1.9) {osa úhlu};
    \end{scope}
    \begin{scope}[xshift=8cm]
      \draw[thick] (0,0) -- (2.6,0) node[right, font=\tiny] {$p$};
      \draw[thick, colorofwisdom] (0,0.9) -- (2.6,0.9) (0,-0.9) -- (2.6,-0.9);
      \draw[<->, gray] (1.3,0) -- (1.3,0.9) node[midway, right, font=\tiny] {$d$};
      \node[font=\scriptsize] at (1.3,-1.9) {rovnoběžky ve vzd. $d$};
    \end{scope}
    \begin{scope}[xshift=13cm]
      \coordinate (A) at (-1.2,0); \coordinate (B) at (1.2,0); \coordinate (V) at (115:1.2);
      \draw[thick, colorofwisdom] (0,0) circle (1.2);
      \draw[thick] (A) -- (B); \draw (A) -- (V) -- (B);
      \filldraw (A) circle (1.5pt) (B) circle (1.5pt);
      \node[font=\scriptsize] at (0,-1.9) {Thaletova kružnice};
    \end{scope}
  \end{tikzpicture}
\end{center}

\begin{example}
  Sestrojme trojúhelník $ABC$, je-li dáno $c = 6\ \si{cm}$, $v_c = 3\ \si{cm}$ a $\gamma = 90^\circ$.

  \textbf{Rozbor}: Úsečku $AB$ sestrojíme rovnou. Bod $C$ musí splňovat dvě podmínky. Úhel $ACB$ je pravý, takže $C$ leží na Thaletově kružnici nad $AB$. Vzdálenost $C$ od přímky $AB$ je $v_c$, takže $C$ leží na rovnoběžce s $AB$ ve vzdálenosti $3\ \si{cm}$.

  \textbf{Postup}:
  \begin{enumerate}
    \item $AB$; $|AB| = 6\ \si{cm}$,
    \item $S$; $S$ je střed $AB$,
    \item $k$; $k(S; 3\ \si{cm})$ (Thaletova kružnice),
    \item $p$; $p\parallel AB$, $|pAB| = 3\ \si{cm}$,
    \item $C$; $C\in k\cap p$,
    \item $\triangle ABC$.
  \end{enumerate}

  \textbf{Diskuze}: Thaletova kružnice má poloměr $3\ \si{cm}$, tedy přesně $v_c$. Rovnoběžka $p$ je její tečnou a má s ní jediný společný bod. V jedné polorovině je tedy jedno řešení (v druhé polorovině je shodné, souměrné podle $AB$). Pokud by bylo $v_c > \frac{c}{2}$, úloha by neměla řešení. Pokud $v_c < \frac{c}{2}$, má v každé polorovině dvě (navzájem souměrná) řešení.
  \begin{center}
    \begin{tikzpicture}[scale=0.6]
      \coordinate (A) at (0,0); \coordinate (B) at (6,0); \coordinate (S) at (3,0); \coordinate (C) at (3,3);
      \draw[dashed, colorofwisdom] (S) circle (3);
      \draw[colorofwisdom!70] (-1,3) -- (7,3) node[right, font=\scriptsize] {$p$};
      \fill[colorofwisdom!15] (A) -- (B) -- (C) -- cycle;
      \draw[thick] (A) -- (B) -- (C) -- cycle;
      \draw ($(C)!0.12!(A)$) -- ($(C)!0.12!(A)!0.12!(B)$);
      \filldraw (S) circle (1.5pt) node[below] {\scriptsize $S$};
      \node[below left] at (A) {\small $A$}; \node[below right] at (B) {\small $B$}; \node[above] at (C) {\small $C$};
      \draw[<->, gray] (6.5,0) -- (6.5,3) node[midway, right, font=\scriptsize] {$v_c$};
      \node[font=\scriptsize, colorofwisdom] at (0.4,2.7) {$k$};
    \end{tikzpicture}
  \end{center}
\end{example}

\begin{example}\label{ex:tecny-z-bodu}
  Sestrojme tečny z bodu $P$ ke kružnici $k(S; r)$, kde $P$ leží vně kružnice.

  \textbf{Rozbor}: Tečna je kolmá na poloměr v bodě dotyku $T$. Úhel $STP$ je tedy pravý a bod $T$ leží na Thaletově kružnici nad $SP$.

  \textbf{Postup}:
  \begin{enumerate}
    \item $O$; $O$ je střed $SP$,
    \item $l$; $l(O; |OS|)$,
    \item $T_1, T_2$; $T_1, T_2\in k\cap l$,
    \item $t_1, t_2$; $t_1 = PT_1$, $t_2 = PT_2$.
  \end{enumerate}
  Úloha má vždy dvě řešení, pokud $P$ leží vně kružnice.
  \begin{center}
    \begin{tikzpicture}[scale=0.9]
      \coordinate (S) at (0,0); \coordinate (P) at (5,0); \coordinate (O) at (2.5,0);
      \draw[thick] (S) circle (1.5) node[above left, xshift=-0.9cm, yshift=0.9cm] {\small $k$};
      \draw[dashed, colorofwisdom] (O) circle (2.5);
      \filldraw (S) circle (1.5pt) node[below] {\small $S$};
      \filldraw (P) circle (1.5pt) node[right] {\small $P$};
      \filldraw (O) circle (1.5pt) node[below] {\small $O$};
      \coordinate (T1) at (0.45,1.431); \coordinate (T2) at (0.45,-1.431);
      \filldraw (T1) circle (1.5pt) node[above left] {\small $T_1$};
      \filldraw (T2) circle (1.5pt) node[below left] {\small $T_2$};
      \draw[thick] ($(P)!1.25!(T1)$) -- (P) -- ($(P)!1.25!(T2)$);
      \draw[gray] (S) -- (T1) (S) -- (T2) (S) -- (P);
      \node[colorofwisdom] at (2.5,2.75) {\small $l$};
    \end{tikzpicture}
  \end{center}
\end{example}

\begin{example}
  Sestrojme trojúhelník $ABC$, je-li dáno $c = 5\ \si{cm}$, $\gamma = 60^\circ$ a $t_c = 4\ \si{cm}$ (těžnice z vrcholu $C$).

  \textbf{Rozbor}: Úsečku $AB$ sestrojíme. Těžnice vede ze středu $S_c$ strany $AB$ do bodu $C$, takže $C$ leží na kružnici $k(S_c; 4\ \si{cm})$. Z bodu $C$ je úsečka $AB$ vidět pod úhlem $60^\circ$, takže $C$ leží na jednom ze dvou kruhových oblouků, ze kterých je $AB$ vidět pod úhlem $60^\circ$. Střed $O$ takového oblouku leží na ose $AB$ a středový úhel $AOB$ je dvojnásobkem obvodového, tedy $120^\circ$.

  \textbf{Postup} (v jedné polorovině):
  \begin{enumerate}
    \item $AB$; $|AB| = 5\ \si{cm}$,
    \item $S_c$; $S_c$ je střed $AB$,
    \item $o$; $o$ je osa $AB$,
    \item $O$; $O\in o$, $|\angle AOB| = 120^\circ$ (stačí sestrojit $|\angle OAB| = 30^\circ$),
    \item $m$; $m(O; |OA|)$,
    \item $k$; $k(S_c; 4\ \si{cm})$,
    \item $C$; $C\in m\cap k$,
    \item $\triangle ABC$.
  \end{enumerate}
  \textbf{Diskuze}: Kružnice $k$ a oblouk $m$ se mohou protnout ve dvou, jednom nebo žádném bodě. V našem případě dostaneme v jedné polorovině dva průsečíky, které jsou souměrné podle osy $o$, takže vzniknou dva shodné trojúhelníky. Až na shodnost je tedy řešení jedno.
  \begin{center}
    \begin{tikzpicture}[scale=0.75]
      \coordinate (A) at (0,0); \coordinate (B) at (5,0); \coordinate (Sc) at (2.5,0); \coordinate (O) at (2.5,1.443);
      \draw[thick, colorofwisdom] (A) arc (210:-30:2.887);
      \draw[dashed] (Sc) circle (4);
      \draw[gray] (2.5,-1) -- (2.5,4.8) node[above, font=\scriptsize] {$o$};
      \draw[gray] (A) -- (O) -- (B);
      \draw (O) ++(210:0.45) arc (210:330:0.45); \node[font=\tiny] at (2.5,0.75) {$120^\circ$};
      \coordinate (C1) at (4.642,3.378); \coordinate (C2) at (0.358,3.378);
      \fill[colorofwisdom!15] (A) -- (B) -- (C1) -- cycle;
      \draw[thick] (A) -- (B) -- (C1) -- cycle;
      \draw[gray] (Sc) -- (C1) node[midway, right, font=\scriptsize] {$t_c$};
      \filldraw (O) circle (1.5pt) node[left, font=\scriptsize] {$O$};
      \filldraw (Sc) circle (1.5pt) node[below, font=\scriptsize] {$S_c$};
      \filldraw (C1) circle (1.5pt) node[right, font=\scriptsize] {$C$};
      \filldraw[gray] (C2) circle (1.5pt) node[left, font=\scriptsize] {$C'$};
      \node[below left] at (A) {\small $A$}; \node[below right] at (B) {\small $B$};
      \node[colorofwisdom, font=\scriptsize] at (4.6,2.1) {$m$}; \node[font=\scriptsize] at (6.6,1) {$k$};
    \end{tikzpicture}
  \end{center}
\end{example}

\begin{problem}
  Proveďte rozbor a napište postup konstrukce:
  \begin{enumerate}[label=(\alph*)]
    \item trojúhelník $ABC$: $a = 5\ \si{cm}$, $b = 4\ \si{cm}$, $\gamma = 45^\circ$,
    \item trojúhelník $ABC$: $c = 7\ \si{cm}$, $v_c = 3\ \si{cm}$, $t_c = 4\ \si{cm}$,
    \item kružnice o poloměru $2\ \si{cm}$, která se dotýká dvou daných různoběžek,
    \item kosočtverec, je-li dána jeho strana $a = 4\ \si{cm}$ a úhlopříčka $u = 6\ \si{cm}$.
  \end{enumerate}
\end{problem}

\begin{solution}
  \begin{enumerate}[label=(\alph*)]
    \item Věta sus: $CB$; $|CB| = 5$; $\overrightarrow{CX}$; $|\angle BCX| = 45^\circ$; $A$; $A\in\overrightarrow{CX}$, $|CA| = 4$; $\triangle ABC$. Jedno řešení.
    \item $AB$; $|AB| = 7$; $S_c$ střed $AB$; $p\parallel AB$ ve vzdálenosti $3$; $k(S_c; 4)$; $C\in p\cap k$. Protože $t_c > v_c$, kružnice protne rovnoběžku ve dvou bodech, které dají dva shodné trojúhelníky (souměrné podle osy $AB$).
    \item Střed hledané kružnice má od obou přímek vzdálenost $2$, leží tedy na průsečíku rovnoběžek s danými přímkami ve vzdálenosti $2$ (což je zároveň na osách úhlů). Takové body jsou čtyři, úloha má čtyři řešení.
    \item Úhlopříčky kosočtverce se navzájem kolmo půlí. Sestrojíme úhlopříčku $AC$ délky $6$, její osu a na ní body $B, D$ ve vzdálenosti $4$ od $A$ (průsečíky kružnice $k(A; 4)$ s osou). Řešení existuje, protože $4 > 3$.
  \end{enumerate}
\end{solution}

\section{Odbočka --- Tři nemožné úlohy}\label{sec:nemozne-ulohy}

Řekové nám zanechali tři konstrukční úlohy, které nedokázali vyřešit, a na kterých si pak dva tisíce let lámaly zuby generace matematiků i amatérů:
\begin{enumerate}
  \item \textbf{Trisekce úhlu}: rozdělit libovolný úhel na tři shodné části.
  \item \textbf{Zdvojení krychle}: sestrojit hranu krychle, která má dvojnásobný objem než daná krychle. Neboli sestrojit délku $\sqrt[3]{2}$.
  \item \textbf{Kvadratura kruhu}: sestrojit čtverec se stejným obsahem jako daný kruh. Neboli sestrojit délku $\sqrt{\pi}$.
\end{enumerate}

Ke zdvojení krychle se váže temná legenda. Kolem roku 430 před naším letopočtem zachvátil Athény mor, na který zemřela asi čtvrtina obyvatel, včetně státníka Perikla. Athéňané poslali posly k věštírně na ostrov Délos, aby se zeptali, jak mor zastavit. Odpověď zněla: zdvojnásobte Apollónův oltář. Oltář měl tvar krychle. Athéňané nejdřív zdvojnásobili jeho hranu, ale tím objem vzrostl osmkrát a mor neustal. Úloze se proto dodnes říká \textbf{délský problém}. Podle Plútarcha pak Platón vysvětlil, že bůh oltář ve skutečnosti nepotřebuje. Chtěl jen, aby se Řekové víc věnovali geometrii.

Odpověď přišla až v 19. století a zněla: \textit{nejde to}. Pomocí pravítka a kružítka lze sestrojit jen délky, které vzniknou z daných délek sčítáním, odčítáním, násobením, dělením a druhou odmocninou. Číslo $\sqrt[3]{2}$ takhle získat nejde, a stejně tak (u obecného úhlu) trisekce vede na kubickou rovnici, kterou druhými odmocninami vyřešit nejde. Obojí dokázal v roce 1837 francouzský matematik Pierre Wantzel. Kvadraturu kruhu pohřbil v roce 1882 Ferdinand von Lindemann důkazem, že $\pi$ je transcendentní, tedy není kořenem žádného mnohočlenu s celočíselnými koeficienty.

Wantzel sám zemřel ve 33 letech. Jeho kolega Saint-Venant o něm napsal, že pracoval po večerech, málo spal, střídavě se uchyloval ke kávě a opiu a jedl, kdy si vzpomněl. A přestože je kvadratura kruhu dokázaně nemožná už přes sto let, matematické ústavy po celém světě dodnes dostávají dopisy od lidí, kteří ji `vyřešili'. Výraz `kvadratura kruhu' se v mnoha jazycích dokonce ujal jako synonym pro nemožný úkol.

\section{Shrnutí a otázky k ústní zkoušce}

\begin{problem}
  Odpovězte vlastními slovy:
  \begin{enumerate}
    \item Co je geometrické zobrazení? Co je samodružný bod a samodružný útvar?
    \item Co je shodné zobrazení? Popište osovou a středovou souměrnost, posunutí a otočení. Jaké mají samodružné body?
    \item Která shodná zobrazení mění orientaci?
    \item Co vznikne složením dvou osových souměrností?
    \item Co je podobné zobrazení a stejnolehlost? Jaké vlastnosti má stejnolehlost?
    \item Uveďte příklady osově a středově souměrných útvarů.
    \item Z jakých částí se skládá řešení konstrukční úlohy?
    \item Jaké znáte množiny bodů dané vlastnosti a jak se využívají v konstrukčních úlohách?
    \item Jak sestrojíte tečnu ke kružnici z bodu ležícího vně kružnice?
    \item Které klasické konstrukční úlohy nelze vyřešit pravítkem a kružítkem a proč?
  \end{enumerate}
\end{problem}

\koneckapitoly

\chapter{Stereometrie a tělesa}
\label{kap:stereometrie}
\mathquote{Let no one ignorant of geometry enter.}{nápis nad vchodem do Platónovy Akademie (podle legendy)}

Stereometrie je geometrie v prostoru. Všechno, co jsme se naučili v rovině, platí dál, jen se objeví jedna nová dimenze a s ní nové situace. Například dvě přímky, které se neprotínají, a přesto nejsou rovnoběžné. V první části kapitoly se podíváme na vzájemné polohy a vzdálenosti útvarů v prostoru, ve druhé na tělesa, jejich objemy a povrchy.

\section{Polohové vlastnosti}

\begin{definitionbox}
  \textbf{Vzájemná poloha dvou přímek v prostoru}:
  \begin{itemize}
    \item \textbf{totožné},
    \item \textbf{rovnoběžné různé}: leží v jedné rovině a nemají společný bod,
    \item \textbf{různoběžné}: mají právě jeden společný bod (a tedy leží v jedné rovině),
    \item \textbf{mimoběžné}: \textbf{neleží v jedné rovině}, nemají společný bod a nejsou rovnoběžné.
  \end{itemize}

  \textbf{Vzájemná poloha přímky a roviny}: přímka \textbf{leží v rovině}, je s ní \textbf{rovnoběžná} (nemá s ní společný bod), nebo je s ní \textbf{různoběžná} (protíná ji v jednom bodě, tzv. průsečíku).

  \textbf{Vzájemná poloha dvou rovin}: \textbf{totožné}, \textbf{rovnoběžné různé}, nebo \textbf{různoběžné} (protínají se v přímce, tzv. průsečnici).
\end{definitionbox}

Mimoběžky jsou ta novinka, kvůli které stereometrie stojí za samostatnou kapitolu. V rovině se dvě nerovnoběžné přímky vždy protnou. V prostoru se mohou minout. Představte si dálniční most přes jinou dálnici: silnice nahoře a dole se nikdy nepotkají a přitom nejsou rovnoběžné.

\begin{definitionbox}
  \textbf{Užitečná kritéria}:
  \begin{itemize}
    \item Přímka je rovnoběžná s rovinou, pokud je rovnoběžná s nějakou přímkou ležící v této rovině.
    \item Dvě roviny jsou rovnoběžné, pokud jedna obsahuje dvě různoběžky rovnoběžné s druhou rovinou.
    \item Pokud dvě rovnoběžné roviny protneme třetí rovinou, jsou průsečnice \textbf{rovnoběžné}.
    \item Přímka je kolmá k rovině, pokud je kolmá ke dvěma různoběžkám této roviny.
  \end{itemize}
\end{definitionbox}

V dalších příkladech budeme pracovat s krychlí $ABCDEFGH$, kde $ABCD$ je dolní podstava, $EFGH$ horní podstava a hrana $AE$ je svislá (bod $E$ je nad $A$, $F$ nad $B$ a tak dále).

\begin{center}
  \begin{tikzpicture}[scale=2, x={(1cm,0cm)}, y={(0.45cm,0.35cm)}, z={(0cm,1cm)}]
    \coordinate (A) at (0,0,0); \coordinate (B) at (1,0,0); \coordinate (C) at (1,1,0); \coordinate (D) at (0,1,0);
    \coordinate (E) at (0,0,1); \coordinate (F) at (1,0,1); \coordinate (G) at (1,1,1); \coordinate (H) at (0,1,1);
    \draw[thick] (A) -- (B) -- (C) -- (G) -- (H) -- (E) -- cycle;
    \draw[thick] (B) -- (F) -- (G) (E) -- (F);
    \draw[dashed] (A) -- (D) -- (C) (D) -- (H);
    \node[below left] at (A) {\small $A$}; \node[below right] at (B) {\small $B$}; \node[right] at (C) {\small $C$}; \node[left] at (D) {\small $D$};
    \node[left] at (E) {\small $E$}; \node[right] at (F) {\small $F$}; \node[above right] at (G) {\small $G$}; \node[above left] at (H) {\small $H$};
  \end{tikzpicture}
\end{center}

\begin{example}
  V krychli $ABCDEFGH$:
  \begin{itemize}
    \item $AB$ a $HG$ jsou rovnoběžné (obě jsou rovnoběžné s $DC$),
    \item $AB$ a $BF$ jsou různoběžné (průsečík $B$), dokonce kolmé,
    \item $AB$ a $CG$ jsou mimoběžné. Neprotínají se, nejsou rovnoběžné a žádná rovina je neobsahuje obě,
    \item $AB$ a rovina $EFG$ jsou rovnoběžné (protože $AB\parallel EF$),
    \item rovina $ABF$ (přední stěna) a rovina $BCG$ (pravá stěna) jsou různoběžné s průsečnicí $BF$.
  \end{itemize}
\end{example}

\begin{problem}
  Určete vzájemnou polohu v krychli $ABCDEFGH$:
  \begin{enumerate}[label=(\alph*)]
    \item přímek $AG$ a $BH$,
    \item přímek $AC$ a $FH$,
    \item přímek $BD$ a $EG$,
    \item přímek $BD$ a $CG$,
    \item přímky $AC$ a roviny $EFG$,
    \item rovin $ACG$ a $BDH$.
  \end{enumerate}
\end{problem}

\begin{solution}
  \begin{enumerate}[label=(\alph*)]
    \item Různoběžné, jsou to tělesové úhlopříčky a protínají se ve středu krychle.
    \item Mimoběžné. Na první pohled vypadají jako rovnoběžky, protože jsou to úhlopříčky protilehlých stěn. Jenže $FH$ je rovnoběžná s $BD$, ne s $AC$. A protože leží v různých rovnoběžných rovinách, nemohou se ani protnout. Tohle je přesně ten typ chyby, který se v prostoru dělá nejčastěji.
    \item Mimoběžné. Jejich průměty do podstavy jsou dvě různé úhlopříčky čtverce, které se protínají, ale přímky samotné jsou v různých výškách.
    \item Mimoběžné. Svislá přímka $CG$ protíná rovinu podstavy v bodě $C$, který na $BD$ neleží.
    \item Rovnoběžné, $AC\parallel EG$.
    \item Různoběžné, jejich průsečnicí je svislá přímka procházející středy podstav.
  \end{enumerate}
\end{solution}

\subsection{Řezy těles}

Řez tělesa rovinou je průnik tělesa s rovinou. U krychle je to mnohoúhelník, jehož strany leží ve stěnách krychle. Při konstrukci řezu se používají dvě pravidla:
\begin{enumerate}
  \item Dva body roviny řezu, které leží ve stejné stěně, spojíme úsečkou.
  \item Roviny řezu protínají \textbf{rovnoběžné} stěny v \textbf{rovnoběžných} úsečkách.
\end{enumerate}
Pokud to nestačí, prodloužíme úsečku řezu ve stěně až k průsečíku s přímkou hrany a tím získáme další bod roviny řezu v sousední stěně.

\begin{example}
  Řez krychle rovinou $BDE$ je rovnostranný trojúhelník $BDE$ se stranou $a\sqrt{2}$, protože všechny jeho strany jsou stěnové úhlopříčky. Řez rovinou procházející středy hran $AB$, $BC$, $CG$, $GH$, $HE$ a $EA$ je dokonce \textbf{pravidelný šestiúhelník}. Tato rovina je kolmá na tělesovou úhlopříčku $DF$ a prochází středem krychle.
  \begin{center}
    \begin{tikzpicture}[scale=2.1, x={(1cm,0cm)}, y={(0.45cm,0.35cm)}, z={(0cm,1cm)}]
      \coordinate (A) at (0,0,0); \coordinate (B) at (1,0,0); \coordinate (C) at (1,1,0); \coordinate (D) at (0,1,0);
      \coordinate (E) at (0,0,1); \coordinate (F) at (1,0,1); \coordinate (G) at (1,1,1); \coordinate (H) at (0,1,1);
      \fill[colorofwisdom!25] (0.5,0,0) -- (1,0.5,0) -- (1,1,0.5) -- (0.5,1,1) -- (0,0.5,1) -- (0,0,0.5) -- cycle;
      \draw[thick, colorofwisdom] (0.5,0,0) -- (1,0.5,0) -- (1,1,0.5) -- (0.5,1,1) -- (0,0.5,1) -- (0,0,0.5) -- cycle;
      \draw[thick] (A) -- (B) -- (C) -- (G) -- (H) -- (E) -- cycle;
      \draw[thick] (B) -- (F) -- (G) (E) -- (F);
      \draw[dashed] (A) -- (D) -- (C) (D) -- (H);
      \node[below left] at (A) {\small $A$}; \node[below right] at (B) {\small $B$}; \node[right] at (C) {\small $C$}; \node[left] at (D) {\small $D$};
      \node[left] at (E) {\small $E$}; \node[right] at (F) {\small $F$}; \node[above right] at (G) {\small $G$}; \node[above left] at (H) {\small $H$};
    \end{tikzpicture}
  \end{center}
  Ověřit to můžeme v souřadnicích. Pro krychli s hranou $1$ a $A = [0, 0, 0]$ leží všech šest středů hran v rovině $x - y + z = \frac{1}{2}$ (dosaďte si) a vzdálenost každých dvou sousedních je $\frac{\sqrt{2}}{2}$.
\end{example}

\section{Metrické vlastnosti}

\begin{definitionbox}
  \textbf{Vzdálenosti}:
  \begin{itemize}
    \item \textbf{Vzdálenost bodu od přímky} (roviny) je vzdálenost bodu od jeho kolmého průmětu (paty kolmice) na přímku (rovinu).
    \item \textbf{Vzdálenost rovnoběžných přímek} (rovin, přímky a s ní rovnoběžné roviny) je vzdálenost libovolného bodu jednoho útvaru od druhého útvaru.
    \item \textbf{Vzdálenost mimoběžek} je délka jejich \textbf{osy}, tedy úsečky, která je kolmá na obě mimoběžky a spojuje je.
  \end{itemize}

  \textbf{Odchylky}:
  \begin{itemize}
    \item \textbf{Odchylka dvou přímek} je velikost ostrého (nebo pravého) úhlu, který svírají. U mimoběžek je to odchylka různoběžek rovnoběžných s nimi, vedených jedním bodem.
    \item \textbf{Odchylka přímky a roviny} je odchylka přímky a jejího kolmého průmětu do roviny.
    \item \textbf{Odchylka dvou rovin} je odchylka přímek, které v obou rovinách leží a jsou kolmé na průsečnici v jednom jejím bodě.
  \end{itemize}
  Odchylky jsou vždy z intervalu $[0^\circ, 90^\circ]$.
\end{definitionbox}

Metrické úlohy v prostoru se skoro vždy řeší stejně. Najdeme vhodný pravoúhlý trojúhelník, který leží v nějaké rovině, a použijeme Pythagorovu větu nebo goniometrické funkce z kapitoly~\ref{kap:goniometrie}. Celý trik je v tom, ten trojúhelník najít.

\begin{example}
  V krychli o hraně $a$:
  \begin{itemize}
    \item \textbf{stěnová úhlopříčka}: $|AC| = \sqrt{a^2 + a^2} = a\sqrt{2}$,
    \item \textbf{tělesová úhlopříčka}: trojúhelník $ACG$ je pravoúhlý s pravým úhlem u $C$ (hrana $CG$ je kolmá na celou podstavu), takže $|AG| = \sqrt{2a^2 + a^2} = a\sqrt{3}$,
    \item \textbf{odchylka tělesové úhlopříčky od podstavy}: je to úhel $GAC$, pro který $\tg\varphi = \frac{|CG|}{|AC|} = \frac{a}{a\sqrt{2}} = \frac{\sqrt{2}}{2}$, tedy $\varphi\approx 35{,}26^\circ$,
    \item \textbf{odchylka přímek $AF$ a $AH$} (stěnové úhlopříčky ze stejného vrcholu): trojúhelník $AFH$ má všechny strany stěnové úhlopříčky, je tedy rovnostranný a odchylka je $60^\circ$,
    \item \textbf{vzdálenost mimoběžek $AB$ a $CG$}: úsečka $BC$ je kolmá na obě, takže vzdálenost je $a$.
  \end{itemize}
\end{example}

\begin{problem}
  \begin{enumerate}[label=(\alph*)]
    \item Kvádr má rozměry $3$, $4$ a $12$. Jak dlouhá je jeho tělesová úhlopříčka?
    \item V krychli o hraně $a$ určete vzdálenost bodu $A$ od přímky $FH$.
    \item V krychli o hraně $a$ určete odchylku rovin $ABC$ a $ABG$.
    \item V pravidelném čtyřbokém jehlanu má podstava hranu $6$ a výška je $4$. Určete délku boční hrany a odchylku boční stěny od podstavy.
  \end{enumerate}
\end{problem}

\begin{solution}
  \begin{enumerate}[label=(\alph*)]
    \item $\sqrt{9 + 16 + 144} = 13$.
    \item Trojúhelník $AFH$ je rovnostranný se stranou $a\sqrt{2}$, vzdálenost $A$ od $FH$ je jeho výška: $\frac{\sqrt{3}}{2}a\sqrt{2} = \frac{\sqrt{6}}{2}a$.
    \item Průsečnice je $AB$. V rovině $ABC$ je na ni kolmá $BC$, v rovině $ABG$ je na ni kolmá $BG$. Odchylka je úhel $CBG = 45^\circ$.
    \item Pata výšky je ve středu podstavy, jehož vzdálenost od vrcholu podstavy je $3\sqrt{2}$. Boční hrana je $\sqrt{16 + 18} = \sqrt{34}$. Pro odchylku boční stěny vezmeme trojúhelník z výšky jehlanu, spojnice středu podstavy se středem hrany podstavy (délka $3$) a stěnové výšky: $\tg\varphi = \frac{4}{3}$, $\varphi\approx 53{,}1^\circ$.
  \end{enumerate}
\end{solution}

\section{Tělesa}

\begin{definitionbox}
  \textbf{Mnohostěn} je těleso ohraničené konečně mnoha mnohoúhelníky (stěnami). Strany stěn jsou \textbf{hrany}, jejich koncové body \textbf{vrcholy}.
  \begin{itemize}
    \item \textbf{Hranol}: dvě shodné rovnoběžné podstavy (mnohoúhelníky) a boční stěny jsou rovnoběžníky. \textbf{Kolmý} hranol má boční stěny kolmé na podstavu (obdélníky). \textbf{Pravidelný} hranol je kolmý a má pravidelnou podstavu. \textbf{Kvádr} je kolmý hranol s obdélníkovou podstavou, \textbf{krychle} je kvádr se shodnými hranami.
    \item \textbf{Jehlan}: podstava je mnohoúhelník, boční stěny jsou trojúhelníky se společným vrcholem. \textbf{Pravidelný} jehlan má pravidelnou podstavu a vrchol nad jejím středem.
  \end{itemize}
  \textbf{Rotační tělesa}: \textbf{válec}, \textbf{kužel} a \textbf{koule} vzniknou rotací obdélníku, pravoúhlého trojúhelníku a půlkruhu kolem osy.
\end{definitionbox}

\begin{definitionbox}
  \textbf{Objemy a povrchy} ($S_p$ je obsah podstavy, $S_{pl}$ obsah pláště, $v$ výška, $r$ poloměr, $s$ délka strany kužele):
  \begin{center}
    \renewcommand{\arraystretch}{1.6}
    \begin{tabular}{|l|l|l|}
      \hline
      Těleso & Objem $V$ & Povrch $S$ \\
      \hline
      hranol & $S_p\cdot v$ & $2S_p + S_{pl}$ \\
      kvádr $a\times b\times c$ & $abc$ & $2(ab + bc + ac)$ \\
      krychle & $a^3$ & $6a^2$ \\
      jehlan & $\frac{1}{3}S_p\cdot v$ & $S_p + S_{pl}$ \\
      válec & $\pi r^2v$ & $2\pi r^2 + 2\pi rv$ \\
      kužel & $\frac{1}{3}\pi r^2v$ & $\pi r^2 + \pi rs$ \\
      koule & $\frac{4}{3}\pi r^3$ & $4\pi r^2$ \\
      komolý jehlan & $\frac{v}{3}\left(S_1 + \sqrt{S_1S_2} + S_2\right)$ & $S_1 + S_2 + S_{pl}$ \\
      \hline
    \end{tabular}
  \end{center}
\end{definitionbox}

Odkud se vzorce berou? U hranolu je to přirozené: objem je `podstava krát výška', jako když skládáte stejné vrstvy na sebe. U ostatních těles se hodí následující princip.

\begin{definitionbox}
  \textbf{Cavalieriho princip}: Pokud mají dvě tělesa stejnou výšku a každá rovina rovnoběžná s podstavou je protíná v řezech o stejném obsahu, mají obě tělesa stejný objem.
\end{definitionbox}

Představte si balíček karet. Pokud ho nakloníte tak, že karty se po sobě posunou, tvar se změní, ale objem ne. To je přesně Cavalieriho princip. Odtud například plyne, že šikmý hranol má stejný objem jako kolmý se stejnou podstavou a výškou.

\begin{example}
  Proč má jehlan třetinový objem? Krychli lze rozdělit na tři shodné jehlany se společným vrcholem $G$ a podstavami $ABCD$, $ABFE$ a $ADHE$ (tři stěny, které se nedotýkají vrcholu $G$). Každý z nich má podstavu $a^2$ a výšku $a$, a protože jsou tři a dohromady dávají krychli, je objem každého $\frac{1}{3}a^3 = \frac{1}{3}S_pv$. Pro obecný jehlan pak vzorec plyne z Cavalieriho principu a podobnosti řezů. Kužel je `jehlan s kruhovou podstavou', a proto má stejně třetinu objemu válce.
\end{example}

\begin{example}
  Pravidelný čtyřboký jehlan má hranu podstavy $a = 6$ a výšku $v = 4$. Objem je $V = \frac{1}{3}\cdot 36\cdot 4 = 48$. Stěnová výška (výška boční stěny) je přepona trojúhelníku s odvěsnami $v = 4$ a $\frac{a}{2} = 3$, tedy $w = 5$. Povrch je $S = 36 + 4\cdot\frac{6\cdot 5}{2} = 96$.
\end{example}

\begin{example}
  Kužel má poloměr podstavy $r = 3$ a výšku $v = 4$. Strana je $s = \sqrt{9 + 16} = 5$. Objem $V = \frac{1}{3}\pi\cdot 9\cdot 4 = 12\pi$ a povrch $S = 9\pi + 15\pi = 24\pi$. Plášť kužele se dá rozvinout do roviny jako kruhová výseč s poloměrem $s$, jejíž oblouk má délku $2\pi r$. Proto $S_{pl} = \frac{1}{2}\cdot s\cdot 2\pi r = \pi rs$ podle vzorce pro výseč z kapitoly~\ref{kap:planimetrie}.
\end{example}

\begin{example}\label{ex:archimedes}
  Vepíšeme kouli o poloměru $r$ do válce (válec má poloměr $r$ a výšku $2r$). Objem válce je $2\pi r^3$, objem koule $\frac{4}{3}\pi r^3$. Povrch válce je $2\pi r^2 + 2\pi r\cdot 2r = 6\pi r^2$, povrch koule $4\pi r^2$. V obou případech je poměr koule ku válci přesně
  \begin{equation*}
    \frac{2}{3}.
  \end{equation*}
  To je objev, který Archimédés považoval za svůj největší a který si nechal vytesat na hrob, jak jsme zmínili v kapitole~\ref{kap:elementarni}. Navíc objem kužele se stejnou podstavou a výškou je $\frac{2}{3}\pi r^3$, takže kužel, koule a válec mají objemy v poměru $1:2:3$.
\end{example}

\begin{note}
  Pokud těleso zvětšíme $k$-krát (podobností), jeho povrch se zvětší $k^2$-krát a objem $k^3$-krát. Proto obři z pohádek nemohou existovat. Kdyby byl člověk desetkrát vyšší, vážil by tisíckrát víc, ale průřez jeho kostí by byl jen stokrát větší. Zlomil by si nohy při prvním kroku. Tento argument použil už Galileo Galilei. A proto mají sloni tak tlusté nohy a mravenci tak tenké.
\end{note}

\begin{problem}
  \begin{enumerate}[label=(\alph*)]
    \item Kolik litrů vody se vejde do koule o průměru $6\ \si{dm}$?
    \item Válec má objem $100\pi\ \si{cm^3}$ a výšku $4\ \si{cm}$. Jaký má povrch?
    \item Krychli o hraně $a$ vepíšeme kouli. Kolik procent objemu krychle koule zabírá?
    \item Kolikrát se zvětší objem krychle, když její hranu zvětšíme o $50\,\%$?
    \item Z kruhové výseče o poloměru $10\ \si{cm}$ a středovém úhlu $216^\circ$ stočíme plášť kužele. Jaký má kužel poloměr podstavy a výšku?
  \end{enumerate}
\end{problem}

\begin{solution}
  \begin{enumerate}[label=(\alph*)]
    \item $r = 3\ \si{dm}$, $V = \frac{4}{3}\pi\cdot 27 = 36\pi\approx 113\ \si{dm^3} = 113$ litrů.
    \item $\pi r^2\cdot 4 = 100\pi$, $r = 5$. $S = 2\pi\cdot 25 + 2\pi\cdot 5\cdot 4 = 90\pi\ \si{cm^2}$.
    \item $r = \frac{a}{2}$, $V_k = \frac{4}{3}\pi\frac{a^3}{8} = \frac{\pi}{6}a^3$. Poměr $\frac{\pi}{6}\approx 0{,}524$, tedy asi $52{,}4\,\%$.
    \item $1{,}5^3 = 3{,}375$krát.
    \item Oblouk výseče je obvod podstavy: $\frac{\pi\cdot 10\cdot 216}{180} = 12\pi = 2\pi r$, tedy $r = 6\ \si{cm}$. Strana kužele je $10$, výška $\sqrt{100 - 36} = 8\ \si{cm}$.
  \end{enumerate}
\end{solution}

\subsection{Eulerova věta a pravidelné mnohostěny}

\begin{definitionbox}
  \textbf{Eulerova věta o mnohostěnech}: Pro každý konvexní mnohostěn s $V$ vrcholy, $H$ hranami a $S$ stěnami platí
  \begin{equation*}
    V - H + S = 2.
  \end{equation*}

  \textbf{Pravidelný mnohostěn} (platónské těleso) je konvexní mnohostěn, jehož stěny jsou shodné pravidelné mnohoúhelníky a v každém vrcholu se stýká stejný počet stěn. Existuje jich \textbf{právě pět}:
  \begin{center}
    \renewcommand{\arraystretch}{1.3}
    \begin{tabular}{|l|c|c|c|c|}
      \hline
      Těleso & stěny & $V$ & $H$ & $S$ \\
      \hline
      čtyřstěn & trojúhelníky & $4$ & $6$ & $4$ \\
      krychle (šestistěn) & čtverce & $8$ & $12$ & $6$ \\
      osmistěn & trojúhelníky & $6$ & $12$ & $8$ \\
      dvanáctistěn & pětiúhelníky & $20$ & $30$ & $12$ \\
      dvacetistěn & trojúhelníky & $12$ & $30$ & $20$ \\
      \hline
    \end{tabular}
  \end{center}
\end{definitionbox}

\begin{proof}[(že jich je nejvýš pět)]
  V každém vrcholu se stýkají alespoň tři stěny a součet jejich úhlů musí být menší než $360^\circ$ (jinak by se vrchol `rozplácl' do roviny). Pravidelný trojúhelník má úhly $60^\circ$, takže se jich ve vrcholu mohou stýkat $3$, $4$ nebo $5$ ($6\cdot 60^\circ = 360^\circ$ už nejde). Čtverec má úhly $90^\circ$, takže jen $3$. Pětiúhelník má úhly $108^\circ$, tedy jen $3$. Šestiúhelník má $120^\circ$ a $3\cdot 120^\circ = 360^\circ$, takže už nejde ani jeden. Dohromady $3 + 1 + 1 = 5$ možností a každá odpovídá jednomu tělesu v tabulce.
\end{proof}

Hráči deskových her znají všech pět pravidelných mnohostěnů jako kostky k4, k6, k8, k12 a k20. Je to tak proto, že jedině pravidelný mnohostěn dává každé stěně stejnou šanci. Pravděpodobnosti se věnujeme v pozdější kapitole.

\section{Odbočka --- Keplerův vesmír a proces s čarodějnicí}

Johannes Kepler věřil, že vesmír je postaven na geometrii. V roce 1596 publikoval model, ve kterém je šest tehdy známých planet uloženo na koulích, mezi které je vloženo právě pět platónských těles. Pět těles, šest koulí, šest planet. Bylo to tak krásné, že to přece musela být pravda. Nebyla. Kepler sám později z přesných pozorování Tychona Braheho (který mimochodem zemřel v Praze a je pohřbený v Týnském chrámu) zjistil, že planety obíhají po elipsách, a svůj vlastní nádherný model opustil. To je na vědě to nejlepší: i to nejkrásnější se zahodí, když to nesouhlasí s pozorováním.

V roce 1615 byla Keplerova matka Katharina obviněna z čarodějnictví. Sousedka tvrdila, že ji Katharina otrávila nápojem, a během několika let se přidali další svědci s historkami o uhynulém dobytku a o tom, jak Katharina prý prochází zavřenými dveřmi. V tehdejším Württembersku nebylo takové obvinění žádná legrace: v okolí bylo za čarodějnictví popraveno několik žen. Kepler přerušil práci, odcestoval domů a osobně vedl její obhajobu. Katharina strávila čtrnáct měsíců ve vězení v řetězech a byla vystavena takzvanému \textit{territio}, tedy výslechu zastrašováním, při kterém jí kat ukázal mučicí nástroje a podrobně popsal, co s ní udělá. Nepřiznala se. V roce 1621 byla propuštěna a o půl roku později zemřela.

\section{Shrnutí a otázky k ústní zkoušce}

\begin{problem}
  Odpovězte vlastními slovy:
  \begin{enumerate}
    \item Jaké jsou vzájemné polohy dvou přímek v prostoru? Co jsou mimoběžky? Uveďte příklad na krychli.
    \item Jaké jsou vzájemné polohy přímky a roviny a dvou rovin?
    \item Jak se konstruuje řez krychle rovinou? Jaká pravidla používáte?
    \item Jak je definována vzdálenost bodu od roviny a vzdálenost mimoběžek?
    \item Jak je definována odchylka dvou přímek, přímky a roviny a dvou rovin?
    \item Odvoďte délku stěnové a tělesové úhlopříčky krychle.
    \item Jaké znáte druhy těles? Uveďte vzorce pro jejich objemy a povrchy.
    \item Co říká Cavalieriho princip? Proč má jehlan třetinový objem hranolu?
    \item Jak se změní objem a povrch tělesa při podobnosti s koeficientem $k$?
    \item Vyslovte Eulerovu větu o mnohostěnech. Proč existuje jen pět pravidelných mnohostěnů?
  \end{enumerate}
\end{problem}

\koneckapitoly

\chapter{Vektory a analytická geometrie v rovině}
\label{kap:analyticka}
\mathquote{Each problem that I solved became a rule which served afterwards to solve other problems.}{René Descartes}

Analytická geometrie je místo, kde se setkává geometrie s algebrou. Body popíšeme čísly (souřadnicemi), přímky a kružnice rovnicemi a geometrické úlohy pak řešíme počítáním. Místo kružítka a pravítka používáme soustavy rovnic z kapitol~\ref{kap:rovnice} a~\ref{kap:kvadraticke}. Tahle myšlenka je tak silná, že na ní dnes stojí počítačová grafika, navigace i fyzika. Váš telefon totiž neumí rýsovat. Umí jen počítat.

\section{Soustava souřadnic}

\begin{definitionbox}
  \textbf{Kartézská soustava souřadnic} v rovině je tvořena dvěma navzájem kolmými číselnými osami $x$ a $y$ se společným počátkem $O$ a stejnou jednotkou. Každému bodu $A$ roviny odpovídá právě jedna uspořádaná dvojice $[a_1, a_2]$ jeho \textbf{souřadnic}. Rovina tak odpovídá kartézskému součinu $\mathbb{R}\times\mathbb{R}$ z kapitoly 2.

  \vspace{0.2cm}

  \textbf{Vzdálenost bodů} $A = [a_1, a_2]$ a $B = [b_1, b_2]$:
  \begin{equation*}
    |AB| = \sqrt{{(b_1 - a_1)}^2 + {(b_2 - a_2)}^2}.
  \end{equation*}
  \textbf{Střed úsečky} $AB$:
  \begin{equation*}
    S_{AB} = \left[\frac{a_1 + b_1}{2}, \frac{a_2 + b_2}{2}\right].
  \end{equation*}
\end{definitionbox}

Vzorec pro vzdálenost je jen Pythagorova věta. Úsečka $AB$ je přeponou pravoúhlého trojúhelníku, jehož odvěsny jsou rovnoběžné s osami a mají délky $|b_1 - a_1|$ a $|b_2 - a_2|$.
Představte si, že jdete po městě s pravoúhlými ulicemi z $A$ do $B$. Po ulicích ujdete $3$ bloky doprava a $4$ nahoru, ale pták letí přímo, po přeponě, a uletí jen $5$. Střed úsečky je zase jen `průměr': jeho souřadnice $x$ leží přesně v půli mezi $a_1$ a $b_1$, souřadnice $y$ v půli mezi $a_2$ a $b_2$.

\begin{center}
  \begin{tikzpicture}[scale=0.65, >=stealth]
    \draw[gray!30] (-0.5,-0.5) grid (5.5,6.5);
    \draw[->] (-0.5,0) -- (5.8,0) node[right] {\small $x$};
    \draw[->] (0,-0.5) -- (0,6.8) node[above] {\small $y$};
    \foreach \x in {1,...,5} \node[below] at (\x,0) {\scriptsize $\x$};
    \foreach \y in {1,...,6} \node[left] at (0,\y) {\scriptsize $\y$};
    \fill[colorofwisdom!15] (1,2) -- (4,2) -- (4,6) -- cycle;
    \draw[thick] (1,2) -- (4,2) node[midway, below] {\scriptsize $|b_1 - a_1| = 3$};
    \draw[thick] (4,2) -- (4,6) node[midway, right] {\scriptsize $|b_2 - a_2| = 4$};
    \draw (3.7,2) -- (3.7,2.3) -- (4,2.3);
    \draw[very thick, colorofwisdom] (1,2) -- (4,6) node[pos=0.8, left] {\small $5$};
    \draw[dotted, thick] (2.5,4) -- (2.5,0) (2.5,4) -- (0,4);
    \node[below] at (2.5,-0.3) {\scriptsize $\frac{5}{2}$};
    \filldraw (1,2) circle (2pt) node[below left] {\small $A$};
    \filldraw (4,6) circle (2pt) node[above right] {\small $B$};
    \filldraw[colorofwisdom] (2.5,4) circle (2pt) node[above left=-1pt] {\small $S_{AB}$};
  \end{tikzpicture}
\end{center}

\begin{example}
  Pro $A = [1, 2]$ a $B = [4, 6]$ je $|AB| = \sqrt{9 + 16} = 5$ a $S_{AB} = \left[\frac{5}{2}, 4\right]$.
\end{example}

\section{Vektory}

\begin{definitionbox}
  \textbf{Vektor} $\vec{u} = \overrightarrow{AB}$ je určen orientovanou úsečkou s počátečním bodem $A$ a koncovým bodem $B$. Dvě orientované úsečky určují stejný vektor, pokud mají stejnou délku, jsou rovnoběžné a stejně orientované. Vektor je tedy `posunutí', nezáleží na tom, odkud ho kreslíme.

  \vspace{0.2cm}

  \textbf{Souřadnice vektoru} $\overrightarrow{AB}$ jsou
  \begin{equation*}
    \vec{u} = B - A = (b_1 - a_1, b_2 - a_2).
  \end{equation*}
  \textbf{Velikost vektoru} $\vec{u} = (u_1, u_2)$ je $|\vec{u}| = \sqrt{u_1^2 + u_2^2}$. \textbf{Nulový vektor} $\vec{o} = (0, 0)$ má velikost $0$. \textbf{Jednotkový vektor} má velikost $1$.
\end{definitionbox}

Vektor si představte jako pokyn na mapě pokladu: `tři kroky na východ a jeden na sever'. Pokyn sám neříká, kde stojíte. Ať vyrazíte odkudkoli, posunete se stejně. Všechny šipky na obrázku jsou proto \textbf{tentýž} vektor $\vec{u} = (3, 1)$, jen nakreslený z různých míst.

\begin{center}
  \begin{tikzpicture}[scale=0.65, >=stealth]
    \draw[gray!30] (-0.5,-0.5) grid (9.5,4.5);
    \foreach \a/\b in {1/1, 0/3, 5/2.5, 6/0} \draw[->, very thick, colorofwisdom] (\a,\b) -- ++(3,1);
    \draw[dashed] (1,1) -- (4,1) node[midway, below] {\scriptsize $3$};
    \draw[dashed] (4,1) -- (4,2) node[midway, right] {\scriptsize $1$};
    \filldraw (1,1) circle (2pt) node[below left] {\small $A$};
    \filldraw (4,2) circle (2pt) node[above] {\small $B$};
    \node[above, colorofwisdom] at (1.5,3.5) {\small $\vec{u}$};
    \node[above, colorofwisdom] at (6.5,3) {\small $\vec{u}$};
    \node[above, colorofwisdom] at (7.5,0.5) {\small $\vec{u}$};
  \end{tikzpicture}
\end{center}

Souřadnice bodů píšeme do hranatých závorek, souřadnice vektorů do kulatých. Je to drobnost, ale zkoušející si na ní rádi zakládají, takže si na ni zvykněte.

\begin{definitionbox}
  \textbf{Operace s vektory}: Pro $\vec{u} = (u_1, u_2)$, $\vec{v} = (v_1, v_2)$ a $k\in\mathbb{R}$:
  \begin{itemize}
    \item \textbf{součet}: $\vec{u} + \vec{v} = (u_1 + v_1, u_2 + v_2)$. Geometricky: vektor $\vec{v}$ připojíme na konec $\vec{u}$ (nebo úhlopříčka rovnoběžníku),
    \item \textbf{násobek}: $k\vec{u} = (ku_1, ku_2)$. Geometricky: prodloužení $|k|$-krát, pro $k < 0$ i s otočením směru,
    \item \textbf{opačný vektor}: $-\vec{u} = (-u_1, -u_2)$.
  \end{itemize}
  Vektory $\vec{u}, \vec{v}$ jsou \textbf{kolineární} (lineárně závislé), pokud je jeden násobkem druhého. Nenulové vektory $(u_1, u_2)$ a $(v_1, v_2)$ jsou kolineární, právě když
  \begin{equation*}
    \begin{vmatrix} u_1 & v_1\\ u_2 & v_2 \end{vmatrix} = u_1v_2 - u_2v_1 = 0.
  \end{equation*}
\end{definitionbox}

\begin{center}
  \begin{tikzpicture}[scale=0.9, >=stealth]
    \coordinate (O) at (0,0); \coordinate (U) at (3,1); \coordinate (V) at (1,2);
    \draw[->, thick, colorofwisdom] (O) -- (U) node[midway, below right] {\small $\vec{u}$};
    \draw[->, thick] (O) -- (V) node[midway, left] {\small $\vec{v}$};
    \draw[->, dashed, colorofwisdom] (V) -- ($(U)+(V)$);
    \draw[->, dashed] (U) -- ($(U)+(V)$);
    \draw[->, very thick] (O) -- ($(U)+(V)$) node[above right] {\small $\vec{u} + \vec{v}$};
  \end{tikzpicture}
\end{center}

Násobek vektoru číslem si představte jako změnu rychlosti při stejném směru: $2\vec{u}$ je dvakrát delší krok, $\frac{1}{2}\vec{u}$ poloviční krok a $-\vec{u}$ krok stejně dlouhý, ale pozpátku.

\begin{center}
  \begin{tikzpicture}[scale=0.7, >=stealth]
    \draw[gray!30] (-0.5,-0.5) grid (11.5,2.5);
    \draw[->, very thick, colorofwisdom] (0,0) -- (2,1) node[above] {\small $\vec{u}$};
    \draw[->, very thick] (2.5,0) -- (6.5,2) node[above] {\small $2\vec{u}$};
    \draw[->, very thick] (7,0) -- (8,0.5) node[above] {\small $\frac{1}{2}\vec{u}$};
    \draw[->, very thick] (11,1) -- (9,0) node[below] {\small $-\vec{u}$};
  \end{tikzpicture}
\end{center}

Ano, ten determinant je ten samý jako v kapitole~\ref{kap:rovnice}. Tam jsme si řekli, že jeho absolutní hodnota je obsah rovnoběžníku, který vektory vytvoří. Kolineární vektory vytvoří `rovnoběžník' s nulovým obsahem, a proto je determinant nulový.

\begin{center}
  \begin{tikzpicture}[scale=0.8, >=stealth]
    \draw[gray!30] (-0.5,-0.5) grid (4.5,3.5);
    \fill[colorofwisdom!20] (0,0) -- (3,1) -- (4,3) -- (1,2) -- cycle;
    \draw[dashed] (3,1) -- (4,3) -- (1,2);
    \draw[->, very thick, colorofwisdom] (0,0) -- (3,1) node[below right] {\small $\vec{u} = (3, 1)$};
    \draw[->, very thick] (0,0) -- (1,2) node[above left] {\small $\vec{v} = (1, 2)$};
    \node at (2,1.5) {\small $S = 5$};
    \node at (2,-1) {\scriptsize $\begin{vmatrix} 3 & 1\\ 1 & 2 \end{vmatrix} = 6 - 1 = 5$};
  \end{tikzpicture}
  \hspace{1.5cm}
  \begin{tikzpicture}[scale=0.8, >=stealth]
    \draw[gray!30] (-0.5,-0.5) grid (4.5,3.5);
    \draw[->, very thick] (0,0) -- (4,2) node[above] {\small $\vec{v} = (4, 2)$};
    \draw[->, very thick, colorofwisdom] (0,0) -- (2,1) node[pos=0.6, below right=2pt] {\small $\vec{u} = (2, 1)$};
    \node at (2,-1) {\scriptsize $\begin{vmatrix} 2 & 4\\ 1 & 2 \end{vmatrix} = 4 - 4 = 0$};
  \end{tikzpicture}
\end{center}
Vpravo jsou vektory kolineární a rovnoběžník se `splácl' do úsečky.

\subsection{Skalární součin}

\begin{definitionbox}
  \textbf{Skalární součin} vektorů $\vec{u} = (u_1, u_2)$ a $\vec{v} = (v_1, v_2)$ je \textbf{číslo}
  \begin{equation*}
    \vec{u}\cdot\vec{v} = u_1v_1 + u_2v_2.
  \end{equation*}
  \textbf{Věta}: Pro úhel $\varphi\in[0^\circ, 180^\circ]$ mezi nenulovými vektory platí
  \begin{equation*}
    \vec{u}\cdot\vec{v} = |\vec{u}|\cdot|\vec{v}|\cdot\cos\varphi,\qquad\text{tedy}\qquad\cos\varphi = \frac{\vec{u}\cdot\vec{v}}{|\vec{u}|\cdot|\vec{v}|}.
  \end{equation*}
  Speciálně: nenulové vektory jsou \textbf{kolmé}, právě když $\vec{u}\cdot\vec{v} = 0$. Dále $\vec{u}\cdot\vec{u} = |\vec{u}|^2$.
\end{definitionbox}

Co skalární součin `měří'? Představte si, že vektor $\vec{v}$ je tyč a slunce svítí kolmo na přímku, na které leží $\vec{u}$. Stín tyče na této přímce má délku $|\vec{v}|\cos\varphi$, a skalární součin je tedy $|\vec{u}|$ krát délka stínu. Říká, \textit{jak moc míří $\vec{v}$ stejným směrem jako $\vec{u}$}. Je-li úhel tupý, stín padá dozadu a součin je záporný. Je-li pravý, stín je jen bod a součin je nula. Odtud ta nejdůležitější vlastnost: kolmost poznáme podle nulového skalárního součinu.

\begin{center}
  \begin{tikzpicture}[scale=0.75, >=stealth]
    \foreach \ang/\sx/\lab in {50/0/{$\vec{u}\cdot\vec{v} > 0$}, 90/5.2/{$\vec{u}\cdot\vec{v} = 0$}, 130/10.4/{$\vec{u}\cdot\vec{v} < 0$}} {
      \begin{scope}[xshift=\sx cm]
        \draw[gray] (-1.5,0) -- (3.3,0);
        \draw[dashed, gray] (\ang:2) -- ({2*cos(\ang)},0);
        \draw[line width=2.5pt, colorofwisdom!60] (0,0) -- ({2*cos(\ang)},0);
        \draw[->, very thick, colorofwisdom] (0,0) -- (3,0) node[below] {\small $\vec{u}$};
        \draw[->, very thick] (0,0) -- (\ang:2) node[above] {\small $\vec{v}$};
        \draw (0,0) ++(0:0.45) arc[start angle=0, end angle=\ang, radius=0.45];
        \node at ({\ang/2}:0.75) {\scriptsize $\varphi$};
        \node at (0.7,-0.9) {\small \lab};
      \end{scope}
    }
  \end{tikzpicture}
\end{center}
Tlustý úsek na přímce vektoru $\vec{u}$ je `stín' vektoru $\vec{v}$, tedy jeho kolmý průmět.

\begin{proof}
  Vektory $\vec{u}$, $\vec{v}$ umístíme do společného počátku. Třetí strana trojúhelníku, který vytvoří, je vektor $\vec{u} - \vec{v}$.
  \begin{center}
    \begin{tikzpicture}[scale=0.8, >=stealth]
      \coordinate (O) at (0,0); \coordinate (U) at (4,0.5); \coordinate (V) at (1,2);
      \fill[colorofwisdom!10] (O) -- (U) -- (V) -- cycle;
      \draw[->, very thick, colorofwisdom] (O) -- (U) node[midway, below] {\small $\vec{u}$};
      \draw[->, very thick] (O) -- (V) node[midway, left] {\small $\vec{v}$};
      \draw[->, very thick, gray] (V) -- (U) node[midway, above right] {\small $\vec{u} - \vec{v}$};
      \draw (O) ++(7.1:0.6) arc[start angle=7.1, end angle=63.4, radius=0.6];
      \node at (35:0.9) {\scriptsize $\varphi$};
    \end{tikzpicture}
  \end{center}
  Šedý vektor $\vec{u} - \vec{v}$ vede z konce $\vec{v}$ do konce $\vec{u}$, protože $\vec{v} + (\vec{u} - \vec{v}) = \vec{u}$. Podle kosinové věty z kapitoly~\ref{kap:goniometrie}:
  \begin{equation*}
    |\vec{u} - \vec{v}|^2 = |\vec{u}|^2 + |\vec{v}|^2 - 2|\vec{u}||\vec{v}|\cos\varphi.
  \end{equation*}
  Zároveň v souřadnicích:
  \begin{equation*}
    |\vec{u} - \vec{v}|^2 = {(u_1 - v_1)}^2 + {(u_2 - v_2)}^2 = |\vec{u}|^2 + |\vec{v}|^2 - 2(u_1v_1 + u_2v_2).
  \end{equation*}
  Porovnáním dostaneme $u_1v_1 + u_2v_2 = |\vec{u}||\vec{v}|\cos\varphi$.
\end{proof}

\begin{example}
  Vektory $\vec{u} = (3, 4)$ a $\vec{v} = (4, -3)$ mají skalární součin $12 - 12 = 0$, jsou tedy kolmé. Mimochodem, vektor kolmý k $(a, b)$ dostaneme vždy tak, že souřadnice prohodíme a jedné změníme znaménko: $(-b, a)$ nebo $(b, -a)$.

  Úhel mezi $\vec{u} = (1, 2)$ a $\vec{v} = (3, 1)$: $\cos\varphi = \frac{3 + 2}{\sqrt{5}\cdot\sqrt{10}} = \frac{5}{\sqrt{50}} = \frac{\sqrt{2}}{2}$, tedy $\varphi = 45^\circ$.
\end{example}

\begin{problem}
  \begin{enumerate}[label=(\alph*)]
    \item Pro $A = [-1, 3]$, $B = [2, -1]$ určete $\overrightarrow{AB}$, $|AB|$ a střed úsečky $AB$.
    \item Pro která $k$ jsou vektory $(2, k)$ a $(k, 8)$ kolineární? Pro které $k$ je $(2, k)$ kolmý na $(3, -6)$?
    \item Určete vnitřní úhly trojúhelníku $A = [0, 0]$, $B = [4, 0]$, $C = [1, 3]$.
    \item Doplňte bod $D$ tak, aby $ABCD$ byl rovnoběžník, kde $A = [1, 1]$, $B = [5, 2]$, $C = [6, 5]$.
  \end{enumerate}
\end{problem}

\begin{solution}
  \begin{enumerate}[label=(\alph*)]
    \item $\overrightarrow{AB} = (3, -4)$, $|AB| = 5$, $S = \left[\frac{1}{2}, 1\right]$.
    \item $2\cdot 8 - k\cdot k = 0$, $k = \pm 4$. Kolmost: $6 - 6k = 0$, $k = 1$.
    \item U $A$: $\overrightarrow{AB} = (4, 0)$, $\overrightarrow{AC} = (1, 3)$, $\cos\alpha = \frac{4}{4\sqrt{10}} = \frac{1}{\sqrt{10}}$, $\alpha\approx 71{,}6^\circ$. U $B$: $\overrightarrow{BA} = (-4, 0)$, $\overrightarrow{BC} = (-3, 3)$, $\cos\beta = \frac{12}{4\cdot 3\sqrt{2}} = \frac{\sqrt{2}}{2}$, $\beta = 45^\circ$. Pak $\gamma\approx 63{,}4^\circ$.
    \item $D = A + \overrightarrow{BC} = [1, 1] + (1, 3) = [2, 4]$.
  \end{enumerate}
  \begin{center}
    \begin{tikzpicture}[scale=0.6, >=stealth]
      \draw[gray!30] (-0.5,-0.5) grid (6.5,5.5);
      \fill[colorofwisdom!10] (1,1) -- (5,2) -- (6,5) -- (2,4) -- cycle;
      \draw[thick] (1,1) -- (5,2) (6,5) -- (2,4);
      \draw[->, very thick, colorofwisdom] (5,2) -- (6,5) node[midway, right] {\scriptsize $\overrightarrow{BC}$};
      \draw[->, very thick, colorofwisdom] (1,1) -- (2,4) node[midway, left] {\scriptsize $\overrightarrow{BC}$};
      \foreach \p/\c/\pos in {A/{1,1}/below left, B/{5,2}/below right, C/{6,5}/above right, D/{2,4}/above left} \filldraw (\c) circle (2pt) node[\pos] {\small $\p$};
    \end{tikzpicture}
  \end{center}
  Pozor na pořadí vrcholů: rovnoběžník $ABCD$ se obchází dokola, takže $D$ leží naproti $B$, ne naproti $C$.
\end{solution}

\section{Přímka v rovině}

Přímku můžeme popsat třemi způsoby. Každý se hodí na něco jiného, a proto je potřeba umět mezi nimi přecházet.

\begin{definitionbox}
  \textbf{Parametrické vyjádření}: Přímka procházející bodem $A$ se \textbf{směrovým vektorem} $\vec{u}\neq\vec{o}$ je množina bodů
  \begin{equation*}
    X = A + t\vec{u},\quad t\in\mathbb{R},\qquad\text{tedy}\qquad x = a_1 + tu_1,\quad y = a_2 + tu_2.
  \end{equation*}

  \textbf{Obecná rovnice}: Přímka s \textbf{normálovým vektorem} $\vec{n} = (a, b)\neq\vec{o}$ (kolmým na přímku) má rovnici
  \begin{equation*}
    ax + by + c = 0.
  \end{equation*}

  \textbf{Směrnicový tvar}: Pokud přímka není rovnoběžná s osou $y$, lze ji zapsat jako
  \begin{equation*}
    y = kx + q,
  \end{equation*}
  kde $k = \tg\varphi$ je \textbf{směrnice} ($\varphi$ je úhel s kladnou poloosou $x$) a $q$ je úsek na ose $y$. To je přesně graf lineární funkce z kapitoly~\ref{kap:elementarni}.
\end{definitionbox}

Parametrické vyjádření si můžete představit jako pohyb. V čase $t = 0$ jste v bodě $A$ a každou sekundu se posunete o vektor $\vec{u}$. Za všechny časy (i záporné, tedy i v minulosti) opíšete celou přímku. Obecná rovnice naopak přímku nepopisuje jako cestu, ale jako podmínku: bod leží na přímce, právě když splňuje rovnici.

\begin{example}
  Najděme všechna vyjádření přímky $AB$, kde $A = [1, 2]$ a $B = [4, 6]$.

  Směrový vektor $\vec{u} = (3, 4)$. Parametricky: $x = 1 + 3t$, $y = 2 + 4t$.

  Normálový vektor je kolmý na $\vec{u}$: $\vec{n} = (4, -3)$. Obecná rovnice $4x - 3y + c = 0$, dosazením bodu $A$: $4 - 6 + c = 0$, tedy $c = 2$ a $4x - 3y + 2 = 0$.

  Směrnicový tvar: $y = \frac{4}{3}x + \frac{2}{3}$.

  Z parametrického vyjádření dostaneme obecnou rovnici také tak, že z obou rovnic vyloučíme $t$: $t = \frac{x - 1}{3} = \frac{y - 2}{4}$, tedy $4x - 4 = 3y - 6$ a $4x - 3y + 2 = 0$.
\end{example}

Na obrázku je přímka z příkladu. Body s parametrem $t = 0$ a $t = 1$ jsou $A$ a $B$, polovina parametru dává střed úsečky $AB$ a záporné $t$ leží `za' bodem $A$. Směrový vektor $\vec{u}$ míří podél přímky, normálový vektor $\vec{n}$ je na ni kolmý. Obecná rovnice tedy říká: \textit{bod $X$ leží na přímce, právě když je vektor $\overrightarrow{AX}$ kolmý na $\vec{n}$}, tedy $\vec{n}\cdot(X - A) = 0$. Po roznásobení je to přesně $4x - 3y + 2 = 0$.

\begin{center}
  \begin{tikzpicture}[scale=0.55, >=stealth]
    \draw[gray!30] (-1.5,-1.5) grid (6.5,8.5);
    \draw[->] (-1.5,0) -- (6.8,0) node[right] {\small $x$};
    \draw[->] (0,-1.5) -- (0,8.8) node[above] {\small $y$};
    \draw[thick] (-1.25,-1) -- (6.25,9) node[below right] {\small $p$};
    \draw[->, very thick, colorofwisdom] (1,2) -- (4,6) node[pos=0.3, right=4pt] {\small $\vec{u} = (3, 4)$};
    \draw[->, very thick] (1,2) -- (5,-1) node[right] {\small $\vec{n} = (4, -3)$};
    \draw (1.24,2.32) -- (1.56,2.08) -- (1.32,1.76);
    \foreach \t/\lab in {-0.5/{-\frac{1}{2}}, 0/0, 0.5/{\frac{1}{2}}, 1/1, 1.5/{\frac{3}{2}}} {
      \filldraw[colorofwisdom] ({1+3*\t},{2+4*\t}) circle (2.5pt) node[left, xshift=-2pt, fill=white, inner sep=1pt] {\scriptsize $t = \lab$};
    }
    \node[above left=-1pt, xshift=-2pt, yshift=5pt] at (1,2) {\small $A$};
    \node[above left=-1pt, xshift=-2pt, yshift=5pt] at (4,6) {\small $B$};
  \end{tikzpicture}
\end{center}

\begin{definitionbox}
  \textbf{Vzájemná poloha dvou přímek}: Pro přímky $p: a_1x + b_1y + c_1 = 0$ a $q: a_2x + b_2y + c_2 = 0$:
  \begin{itemize}
    \item jsou \textbf{různoběžné}, pokud normálové vektory nejsou kolineární. Průsečík najdeme vyřešením soustavy dvou rovnic,
    \item jsou \textbf{rovnoběžné různé}, pokud jsou normálové vektory kolineární, ale rovnice nejsou násobkem jedna druhé,
    \item jsou \textbf{totožné}, pokud je jedna rovnice násobkem druhé.
  \end{itemize}
  To je přesně rozbor počtu řešení soustavy dvou lineárních rovnic z kapitoly~\ref{kap:rovnice}. Tam jsme řekli, že soustava popisuje dvě přímky, a teď to víme jistě.

  \vspace{0.2cm}

  \textbf{Odchylka přímek} $\varphi\in[0^\circ, 90^\circ]$:
  \begin{equation*}
    \cos\varphi = \frac{|\vec{n}_1\cdot\vec{n}_2|}{|\vec{n}_1|\cdot|\vec{n}_2|}\qquad\text{(nebo totéž se směrovými vektory)}.
  \end{equation*}
  Absolutní hodnota zajistí, že dostaneme ostrý úhel. Přímky v směrnicovém tvaru jsou rovnoběžné, právě když $k_1 = k_2$, a kolmé, právě když $k_1\cdot k_2 = -1$.

  \vspace{0.2cm}

  \textbf{Vzdálenost bodu} $M = [m_1, m_2]$ od přímky $ax + by + c = 0$:
  \begin{equation*}
    d = \frac{|am_1 + bm_2 + c|}{\sqrt{a^2 + b^2}}.
  \end{equation*}
\end{definitionbox}

\begin{center}
  \begin{tikzpicture}[scale=0.8]
    \draw[thick, colorofwisdom] (0,0) -- (3,2);
    \draw[thick] (0,2) -- (3,0.5);
    \filldraw (1.875,1.25) circle (2pt);
    \node at (1.5,-0.5) {\scriptsize různoběžné};
    \begin{scope}[xshift=4cm]
      \draw[thick, colorofwisdom] (0,0.3) -- (3,1.8);
      \draw[thick] (0,1) -- (3,2.5);
      \node at (1.5,-0.5) {\scriptsize rovnoběžné různé};
    \end{scope}
    \begin{scope}[xshift=8cm]
      \draw[line width=3pt, colorofwisdom!50] (0,0.3) -- (3,1.8);
      \draw[thick] (0,0.3) -- (3,1.8);
      \node at (1.5,-0.5) {\scriptsize totožné};
    \end{scope}
  \end{tikzpicture}
\end{center}

U odchylky můžeme klidně použít normálové vektory místo směrových. Normálový vektor je směrový vektor otočený o $90^\circ$, a když obě přímky otočíme o stejný pravý úhel, úhel mezi nimi se nezmění. Jediný háček je, že vektory mohou svírat tupý úhel $180^\circ - \varphi$ (stačí jeden z nich otočit opačným směrem). Proto je v čitateli absolutní hodnota.

\begin{center}
  \begin{tikzpicture}[scale=1, >=stealth]
    \draw[thick, colorofwisdom] (200:2) -- (20:2.6) node[right] {\small $p_1$};
    \draw[thick] (255:1.2) -- (75:2.2) node[above] {\small $p_2$};
    \draw (20:0.6) arc[start angle=20, end angle=75, radius=0.6];
    \node at (47.5:0.85) {\scriptsize $\varphi$};
    \draw[->, very thick, colorofwisdom] (0,0) -- (110:1.6) node[above] {\small $\vec{n}_1$};
    \draw[->, very thick] (0,0) -- (165:1.6) node[left] {\small $\vec{n}_2$};
    \draw (110:0.9) arc[start angle=110, end angle=165, radius=0.9];
    \node at (137.5:1.15) {\scriptsize $\varphi$};
    \draw[gray] (20:0.25) -- ($(20:0.25)+(110:0.25)$) -- (110:0.25);
  \end{tikzpicture}
\end{center}

Vzdálenost bodu od přímky je délka \textbf{nejkratší} cesty k ní, a ta vede kolmo, tedy ve směru normálového vektoru. Představte si, že stojíte na louce v bodě $M$ a přímka je řeka. K řece nejrychleji doběhnete kolmo k břehu, do bodu $Q$. Důkaz níže tu kolmou cestu nemusí vůbec hledat. Stačí vzít libovolný bod $P$ na břehu a z vektoru $\overrightarrow{PM}$ si nechat jen tu část, která míří ve směru $\vec{n}$ (jeho `stín' na normále, viz skalární součin). Na obrázku je to situace z příkladu níže: $M = [3, 1]$, přímka $4x - 3y + 2 = 0$ a $P = \left[\frac{5}{2}, 4\right]$. Body $P$, $P'$, $M$, $Q$ tvoří obdélník, takže $|PP'| = |QM| = d$.

\begin{center}
  \begin{tikzpicture}[scale=1.2, >=stealth]
    \draw[gray!30] (0,0.5) grid (5,4.5);
    \fill[colorofwisdom!12] (2.5,4) -- (4.26,2.68) -- (3,1) -- (1.24,2.32) -- cycle;
    \draw[thick] (0.25,1) -- (2.875,4.5) node[right] {\small $4x - 3y + 2 = 0$};
    \draw[->, thick, gray] (2.5,4) -- (3,1) node[pos=0.25, left] {\scriptsize $\overrightarrow{PM}$};
    \draw[very thick, colorofwisdom] (3,1) -- (1.24,2.32) node[midway, below left] {\small $d$};
    \draw[->, very thick, colorofwisdom!60] (2.5,4) -- (4.26,2.68) node[midway, above right] {\small $\vec{n}$};
    \draw[dashed] (4.26,2.68) -- (3,1);
    \draw ($(1.24,2.32)+(0.15,0.2)$) -- ++(0.2,-0.15) -- ++(-0.15,-0.2);
    \draw ($(4.26,2.68)+(-0.2,0.15)$) -- ++(-0.15,-0.2) -- ++(0.2,-0.15);
    \filldraw (3,1) circle (1.6pt) node[below right] {\small $M = [3, 1]$};
    \filldraw (2.5,4) circle (1.6pt) node[above left] {\small $P$};
    \filldraw (1.24,2.32) circle (1.6pt) node[left] {\small $Q$};
    \filldraw (4.26,2.68) circle (1.6pt) node[right] {\small $P'$};
  \end{tikzpicture}
\end{center}

\begin{proof}[(vzorce pro vzdálenost)]
  Nechť $P = [p_1, p_2]$ je libovolný bod přímky, tedy $ap_1 + bp_2 = -c$. Vzdálenost bodu $M$ od přímky je délka průmětu vektoru $\overrightarrow{PM}$ do směru kolmého na přímku, tedy do směru normálového vektoru $\vec{n}$. Z definice skalárního součinu je tato délka
  \begin{equation*}
    d = \frac{|\overrightarrow{PM}\cdot\vec{n}|}{|\vec{n}|} = \frac{|a(m_1 - p_1) + b(m_2 - p_2)|}{\sqrt{a^2 + b^2}} = \frac{|am_1 + bm_2 - (ap_1 + bp_2)|}{\sqrt{a^2 + b^2}} = \frac{|am_1 + bm_2 + c|}{\sqrt{a^2 + b^2}}.
  \end{equation*}
\end{proof}

\begin{proof}[(podmínky kolmosti $k_1k_2 = -1$)]
  Přímka $y = kx + q$ má směrový vektor $(1, k)$. Dvě přímky jsou kolmé, právě když jsou kolmé jejich směrové vektory, tedy $(1, k_1)\cdot(1, k_2) = 1 + k_1k_2 = 0$.
\end{proof}

\begin{example}
  \begin{itemize}
    \item Průsečík přímek $x + y - 3 = 0$ a $2x - y = 0$: sečtením $3x = 3$, tedy $x = 1$ a $y = 2$. Průsečík je $[1, 2]$.
    \item Vzdálenost bodu $M = [3, 1]$ od přímky $4x - 3y + 2 = 0$ je $\frac{|12 - 3 + 2|}{\sqrt{16 + 9}} = \frac{11}{5}$.
    \item Odchylka přímek $x - y = 0$ a $y = 0$: $\vec{n}_1 = (1, -1)$, $\vec{n}_2 = (0, 1)$, $\cos\varphi = \frac{1}{\sqrt{2}}$, $\varphi = 45^\circ$.
  \end{itemize}
\end{example}

\begin{definitionbox}
  \textbf{Polorovina}: Přímka $ax + by + c = 0$ rozdělí rovinu na dvě poloroviny. Body, pro které $ax + by + c > 0$, leží v jedné z nich, body s $ax + by + c < 0$ ve druhé. Normálový vektor $\vec{n} = (a, b)$ míří do té, kde je výraz kladný.
\end{definitionbox}

\begin{center}
  \begin{tikzpicture}[scale=0.7, >=stealth]
    \fill[colorofwisdom!12] (-0.5,4.5) -- (4.5,4.5) -- (4.5,-0.5) -- cycle;
    \draw[gray!30] (-0.5,-0.5) grid (4.5,4.5);
    \draw[->] (-0.5,0) -- (4.8,0) node[right] {\small $x$};
    \draw[->] (0,-0.5) -- (0,4.8) node[above] {\small $y$};
    \draw[thick] (-0.5,4.5) -- (4.5,-0.5) node[right] {\small $x + y - 4 = 0$};
    \draw[->, very thick, colorofwisdom] (1.5,2.5) -- (2.5,3.5) node[above] {\small $\vec{n} = (1, 1)$};
    \filldraw (0,0) circle (2pt) node[below left] {\scriptsize $[0, 0]$};
    \filldraw (3,3) circle (2pt) node[right] {\scriptsize $[3, 3]$};
    \node[colorofwisdom] at (3.9,1.5) {\scriptsize $> 0$};
    \node at (0.9,1) {\scriptsize $< 0$};
  \end{tikzpicture}
\end{center}
Polorovinu si představte jako jednu stranu plotu. Dosazením bodu do levé strany rovnice zjistíme, na které straně plotu stojí. Normálový vektor ukazuje `ven' na kladnou stranu.

\begin{problem}
  \begin{enumerate}[label=(\alph*)]
    \item Napište obecnou rovnici přímky, která prochází bodem $P = [2, -1]$ a je (i) rovnoběžná, (ii) kolmá k přímce $3x + 2y - 5 = 0$.
    \item Určete vzájemnou polohu přímek $p: x = 1 + 2t$, $y = 3 - t$ a $q: x + 2y - 7 = 0$.
    \item Určete vzdálenost rovnoběžek $3x - 4y + 1 = 0$ a $3x - 4y - 9 = 0$.
    \item Určete odchylku přímek $y = 2x + 1$ a $y = -3x + 4$.
    \item Leží body $[0, 0]$ a $[3, 3]$ ve stejné polorovině vzhledem k přímce $x + y - 4 = 0$?
  \end{enumerate}
\end{problem}

\begin{solution}
  \begin{enumerate}[label=(\alph*)]
    \item (i) $3x + 2y + c = 0$, $6 - 2 + c = 0$, tedy $3x + 2y - 4 = 0$. (ii) Normálový vektor kolmice je směrovým vektorem původní přímky, $(2, -3)$: $2x - 3y + c = 0$, $4 + 3 + c = 0$, tedy $2x - 3y - 7 = 0$.
    \item Směrový vektor $p$ je $(2, -1)$, normálový tedy $(1, 2)$, stejný jako u $q$. Bod $[1, 3]$ přímky $p$ dosadíme do $q$: $1 + 6 - 7 = 0$. Leží na ní, takže přímky jsou \textbf{totožné}.
    \item Vezmeme bod první přímky, například $[1, 1]$ ($3 - 4 + 1 = 0$), a spočítáme jeho vzdálenost od druhé: $\frac{|3 - 4 - 9|}{5} = 2$. Obecně pro rovnoběžky $ax + by + c_1 = 0$ a $ax + by + c_2 = 0$ je vzdálenost $\frac{|c_1 - c_2|}{\sqrt{a^2 + b^2}}$.
    \item Směrové vektory $(1, 2)$ a $(1, -3)$, $\cos\varphi = \frac{|1 - 6|}{\sqrt{5}\sqrt{10}} = \frac{5}{\sqrt{50}} = \frac{\sqrt{2}}{2}$, $\varphi = 45^\circ$.
    \item Dosadíme: $0 + 0 - 4 = -4 < 0$ a $3 + 3 - 4 = 2 > 0$. Body leží v opačných polorovinách.
  \end{enumerate}
\end{solution}

\section{Aplikace na trojúhelník}

\begin{definitionbox}
  Pro trojúhelník $ABC$ se souřadnicemi vrcholů:
  \begin{itemize}
    \item \textbf{těžiště}: $T = \left[\frac{a_1 + b_1 + c_1}{3}, \frac{a_2 + b_2 + c_2}{3}\right]$,
    \item \textbf{obsah}: $S = \frac{1}{2}\left|\det\left(\overrightarrow{AB}, \overrightarrow{AC}\right)\right| = \frac{1}{2}|(b_1 - a_1)(c_2 - a_2) - (b_2 - a_2)(c_1 - a_1)|$.
  \end{itemize}
\end{definitionbox}

Vzorec pro těžiště plyne z toho, že těžiště dělí těžnici v poměru $2:1$. Vzorec pro obsah z toho, že trojúhelník je polovina rovnoběžníku určeného vektory $\overrightarrow{AB}$ a $\overrightarrow{AC}$.

\begin{center}
  \begin{tikzpicture}[scale=0.85, >=stealth]
    \fill[colorofwisdom!10] (0,0) -- (3.5,0.5) -- (4.5,2.5) -- (1,2) -- cycle;
    \fill[colorofwisdom!30] (0,0) -- (3.5,0.5) -- (1,2) -- cycle;
    \draw[dashed] (3.5,0.5) -- (4.5,2.5) -- (1,2);
    \draw[thick] (3.5,0.5) -- (1,2);
    \draw[->, very thick] (0,0) -- (3.5,0.5) node[below right] {\small $B$};
    \draw[->, very thick] (0,0) -- (1,2) node[above left] {\small $C$};
    \node[below left] at (0,0) {\small $A$};
    \node at (1.5,0.85) {\small $S$};
    \node at (3,1.65) {\small $S$};
  \end{tikzpicture}
  \hspace{1.5cm}
  \begin{tikzpicture}[scale=0.85]
    \coordinate (A) at (0,0); \coordinate (B) at (4,0); \coordinate (C) at (1.2,2.7);
    \draw[thick] (A) -- (B) -- (C) -- cycle;
    \draw[gray] (A) -- ($(B)!0.5!(C)$) (B) -- ($(A)!0.5!(C)$);
    \draw[very thick, colorofwisdom] (C) -- ($(C)!0.6667!($(A)!0.5!(B)$)$) node[midway, right] {\scriptsize $2$};
    \draw[very thick, colorofwisdom!50] ($(C)!0.6667!($(A)!0.5!(B)$)$) -- ($(A)!0.5!(B)$) node[midway, right] {\scriptsize $1$};
    \filldraw ($(A)!0.5!(B)$) circle (1.5pt);
    \filldraw[colorofwisdom] (1.7333,0.9) circle (2pt) node[left=2pt] {\small $T$};
    \node[below left] at (A) {\small $A$}; \node[below right] at (B) {\small $B$}; \node[above] at (C) {\small $C$};
  \end{tikzpicture}
\end{center}
Vlevo: rovnoběžník je rozdělen úhlopříčkou na dva shodné trojúhelníky, takže trojúhelník má polovinu obsahu. Vpravo: těžiště $T$ dělí každou těžnici v poměru $2:1$, tedy leží ve třetině těžnice od strany. Představte si trojúhelník z kartonu: v těžišti ho udržíte na špičce tužky.

\begin{example}
  Trojúhelník $A = [0, 0]$, $B = [4, 0]$, $C = [1, 3]$ má těžiště $T = \left[\frac{5}{3}, 1\right]$ a obsah $S = \frac{1}{2}|4\cdot 3 - 0\cdot 1| = 6$. Kontrola: základna $AB$ má délku $4$ a výška z $C$ je $3$, tedy $S = \frac{4\cdot 3}{2} = 6$.

  Rovnice výšky z vrcholu $C$: je kolmá na $AB$, takže její normálový vektor je $\overrightarrow{AB} = (4, 0)$, a rovnice je $4x + c = 0$, po dosazení $C$: $x = 1$. Rovnice výšky z $A$: normálový vektor je $\overrightarrow{BC} = (-3, 3)$, neboli $(1, -1)$, rovnice $x - y = 0$. Průsečík obou výšek je ortocentrum $V = [1, 1]$.
\end{example}

\begin{center}
  \begin{tikzpicture}[scale=1.1]
    \draw[gray!30] (-0.5,-0.5) grid (4.5,3.5);
    \draw[->] (-0.5,0) -- (4.8,0) node[right] {\small $x$};
    \draw[->] (0,-0.5) -- (0,3.8) node[above] {\small $y$};
    \fill[colorofwisdom!10] (0,0) -- (4,0) -- (1,3) -- cycle;
    \draw[very thick] (0,0) -- (4,0) -- (1,3) -- cycle;
    \draw[thick, colorofwisdom, dashed] (1,3) -- (1,0) (0,0) -- (2,2) (4,0) -- (0.4,1.2);
    \draw[gray] (1,0) -- (1,0.2) -- (1.2,0.2) -- (1.2,0);
    \draw[gray] (2,2) -- (1.86,1.86) -- (2,1.72) -- (2.14,1.86);
    \draw[gray, dotted, thick] (1,3) -- (2,0) (0,0) -- (2.5,1.5) (4,0) -- (0.5,1.5);
    \filldraw[colorofwisdom] (1,1) circle (2pt) node[above left] {\small $V$};
    \filldraw (1.6667,1) circle (2pt) node[right=2pt] {\small $T$};
    \node[below left] at (0,0) {\small $A$}; \node[below right] at (4,0) {\small $B$}; \node[above] at (1,3) {\small $C$};
  \end{tikzpicture}
\end{center}
Čárkovaně jsou výšky (všechny tři se protnou v ortocentru $V$, třetí výška z $B$ má rovnici $x + 3y - 4 = 0$ a bodem $V$ opravdu prochází), tečkovaně těžnice protínající se v těžišti $T$.

\begin{problem}
  Je dán trojúhelník $A = [1, 1]$, $B = [7, 3]$, $C = [3, 7]$.
  \begin{enumerate}[label=(\alph*)]
    \item Určete jeho těžiště a obsah.
    \item Napište rovnici těžnice z vrcholu $C$.
    \item Napište rovnici osy strany $AB$.
    \item Rozhodněte, zda je trojúhelník rovnoramenný nebo pravoúhlý.
  \end{enumerate}
\end{problem}

\begin{solution}
  \begin{enumerate}[label=(\alph*)]
    \item $T = \left[\frac{11}{3}, \frac{11}{3}\right]$. $\overrightarrow{AB} = (6, 2)$, $\overrightarrow{AC} = (2, 6)$, $S = \frac{1}{2}|36 - 4| = 16$.
    \item Střed $AB$ je $[4, 2]$, směrový vektor těžnice $(1, -5)$, normálový $(5, 1)$: $5x + y + c = 0$, $20 + 2 + c = 0$, tedy $5x + y - 22 = 0$. (Kontrola: $C$ dává $15 + 7 - 22 = 0$.)
    \item Osa prochází středem $[4, 2]$ a má normálový vektor $\overrightarrow{AB} = (6, 2)$, tedy $(3, 1)$: $3x + y - 14 = 0$.
    \item $|AB| = |AC| = \sqrt{40}$, je rovnoramenný se základnou $BC$. Skalární součiny: $\overrightarrow{AB}\cdot\overrightarrow{AC} = 24\neq 0$, $\overrightarrow{BA}\cdot\overrightarrow{BC} = (-6)(-4) + (-2)(4) = 16\neq 0$, $\overrightarrow{CA}\cdot\overrightarrow{CB} = (-2)(4) + (-6)(-4) = 16\neq 0$. Pravoúhlý není.
  \end{enumerate}
  \begin{center}
    \begin{tikzpicture}[scale=0.55]
      \draw[gray!30] (-0.5,-0.5) grid (8.5,8.5);
      \draw[->] (-0.5,0) -- (8.8,0) node[right] {\small $x$};
      \draw[->] (0,-0.5) -- (0,8.8) node[above] {\small $y$};
      \fill[colorofwisdom!10] (1,1) -- (7,3) -- (3,7) -- cycle;
      \draw[very thick] (1,1) -- (7,3) -- (3,7) -- cycle;
      \draw[thick, colorofwisdom] (3,7) -- (4,2) node[pos=0.12, right=2pt] {\scriptsize těžnice};
      \draw[thick, dashed] (4.6667,0) -- (2,8) node[below left] {\scriptsize osa $AB$};
      \draw ($(4,2)+(0.2846,0.0949)$) -- ++(-0.0949,0.2846) -- ++(-0.2846,-0.0949);
      \filldraw (4,2) circle (2pt);
      \filldraw[colorofwisdom] (3.6667,3.6667) circle (2.5pt) node[right=2pt] {\small $T$};
      \node[below left] at (1,1) {\small $A$}; \node[right] at (7,3) {\small $B$}; \node[above right] at (3,7) {\small $C$};
      \draw[thick] (3.163,1.51) -- (3.037,1.89) (1.81,4.063) -- (2.19,3.937);
    \end{tikzpicture}
  \end{center}
  Čárky na stranách $AB$ a $AC$ naznačují, že jsou stejně dlouhé. Osa strany $AB$ prochází jejím středem $[4, 2]$ kolmo, těžnice z $C$ míří do téhož středu.
\end{solution}

\section{Příběh --- Lebka filozofa}

Souřadnice jsou pojmenované po Reném Descartovi (latinsky Cartesius), francouzském filozofovi a matematikovi, autorovi věty `Myslím, tedy jsem'. Podle legendy ho myšlenka souřadnic napadla, když jako nemocný ležel v posteli a pozoroval mouchu na stropě. Uvědomil si, že její polohu lze kdykoli popsat dvěma čísly: vzdáleností od dvou stěn. Descartes mimochodem celý život vstával pozdě, protože ráno v posteli prý myslel nejlépe. Čtenáři této knihy jistě chápou.

Právě to se mu stalo osudným. V roce 1649 ho švédská královna Kristina pozvala do Stockholmu, aby ji učil filozofii. Královna ale trvala na tom, že lekce budou v pět hodin ráno, v mrazivém zámku a uprostřed švédské zimy. Descartes po několika měsících dostal zápal plic a v únoru 1650 zemřel. Někteří historici dnes tvrdí, že ho otrávil arzenem katolický kněz, který se bál, že by Descartovy názory odradily protestantskou královnu od přestupu ke katolictví. Kristina mimochodem o několik let později opravdu konvertovala a vzdala se trůnu.

Tím ale Descartův příběh neskončil. Když po šestnácti letech převáželi jeho ostatky do Francie, hlava v rakvi chyběla. Lebka se pak víc než sto let potulovala Evropou, prodávala se a kupovala. Každý majitel na ni napsal své jméno a krátkou báseň, takže je dnes popsaná latinskými verši. Dnes ji najdete v Muzeu člověka v Paříži, zatímco zbytek těla leží v kostele Saint-Germain-des-Prés na druhém břehu Seiny. Zakladatel moderní filozofie, který tvrdil, že mysl a tělo jsou dvě oddělené věci, tak skončil přesně podle své vlastní teorie.

\section{Shrnutí a otázky k ústní zkoušce}

\begin{problem}
  Odpovězte vlastními slovy:
  \begin{enumerate}
    \item Jak se počítá vzdálenost dvou bodů a střed úsečky? Odkud se vzorec pro vzdálenost bere?
    \item Co je vektor? Jak se určí jeho souřadnice a velikost? Jak se vektory sčítají a násobí číslem?
    \item Kdy jsou dva vektory kolineární? Jak to souvisí s determinantem?
    \item Definujte skalární součin. Dokažte vztah se kosinem úhlu. Kdy jsou vektory kolmé?
    \item Jaké znáte způsoby vyjádření přímky v rovině? Co je směrový a normálový vektor? Jak se mezi vyjádřeními přechází?
    \item Jak určíte vzájemnou polohu dvou přímek a jejich průsečík?
    \item Jak se počítá odchylka dvou přímek a vzdálenost bodu od přímky? Odvoďte vzorec pro vzdálenost.
    \item Kdy jsou dvě přímky ve směrnicovém tvaru rovnoběžné a kdy kolmé?
    \item Jak spočítáte těžiště a obsah trojúhelníku ze souřadnic jeho vrcholů?
  \end{enumerate}
\end{problem}

\koneckapitoly

\chapter{Kuželosečky}
\label{kap:kuzelosecky}
\mathquote{Ubi materia, ibi geometria. (Kde je hmota, tam je geometrie.)}{Johannes Kepler}

Pokud dvojitý kužel (dva kužele spojené vrcholy) protneme rovinou, dostaneme podle sklonu roviny kružnici, elipsu, parabolu nebo hyperbolu. Těmto křivkám se proto říká \textbf{kuželosečky}. Zkoumali je už staří Řekové, hlavně Apollónios z Pergy, který jim dal jejich jména. Dva tisíce let to byla čistě abstraktní geometrie. Pak Kepler zjistil, že planety obíhají po elipsách, Newton vysvětlil proč, a ukázalo se, že kuželosečky jsou dráhy všeho, co se pohybuje v gravitačním poli: planet, komet, družic i hozeného míče.

\begin{center}
  \renewcommand{\arraystretch}{1.3}
  \begin{tabular}{|l|l|}
    \hline
    Rovina řezu & Kuželosečka \\
    \hline
    kolmá na osu kužele & kružnice \\
    šikmá, protíná jen jeden kužel celý & elipsa \\
    rovnoběžná s jednou povrchovou přímkou & parabola \\
    protíná oba kužele & hyperbola \\
    \hline
  \end{tabular}
\end{center}

V analytické geometrii je ale budeme definovat jako množiny bodů s určitou vlastností vzdálenosti. A jejich rovnice budou rovnice druhého stupně v $x$ a $y$.

\section{Kružnice}

\begin{definitionbox}
  \textbf{Kružnice} $k(S; r)$ je množina bodů, které mají od středu $S = [m, n]$ vzdálenost $r > 0$. Ze vzorce pro vzdálenost bodů dostaneme její \textbf{středovou rovnici}
  \begin{equation*}
    {(x - m)}^2 + {(y - n)}^2 = r^2.
  \end{equation*}
  Roznásobením dostaneme \textbf{obecnou rovnici} $x^2 + y^2 + ax + by + c = 0$. Naopak z obecné rovnice získáme střed a poloměr \textbf{doplněním na čtverec}.

  \vspace{0.2cm}

  \textbf{Tečna} ke kružnici v bodě $T = [x_0, y_0]$, který na ní leží:
  \begin{equation*}
    (x_0 - m)(x - m) + (y_0 - n)(y - n) = r^2.
  \end{equation*}
\end{definitionbox}

Rovnice tečny plyne z toho, že tečna je kolmá na poloměr $ST$. Vektor $\overrightarrow{ST} = (x_0 - m, y_0 - n)$ je tedy normálovým vektorem tečny. Rovnici pak dostaneme stejně jako u každé přímky z kapitoly~\ref{kap:analyticka}. (Zkuste si, že po úpravě vyjde přesně vzorec výše.)

\begin{example}
  Kružnice se středem $[2, -3]$ a poloměrem $5$ má rovnici ${(x - 2)}^2 + {(y + 3)}^2 = 25$. Po roznásobení $x^2 + y^2 - 4x + 6y - 12 = 0$.

  Opačně: rovnice $x^2 + y^2 - 4x + 6y - 12 = 0$. Doplníme na čtverec zvlášť pro $x$ a $y$:
  \begin{equation*}
    (x^2 - 4x + 4) - 4 + (y^2 + 6y + 9) - 9 - 12 = 0\quad\Leftrightarrow\quad{(x - 2)}^2 + {(y + 3)}^2 = 25.
  \end{equation*}

  Pozor, ne každá rovnice tvaru $x^2 + y^2 + ax + by + c = 0$ popisuje kružnici. Rovnice $x^2 + y^2 + 2x + 2 = 0$ dává ${(x + 1)}^2 + y^2 = -1$, a tu nesplňuje žádný bod. Rovnice $x^2 + y^2 - 2x + 1 = 0$ dává ${(x - 1)}^2 + y^2 = 0$, což je jediný bod $[1, 0]$.
\end{example}

\begin{example}
  Tečna ke kružnici $x^2 + y^2 = 25$ v bodě $[3, 4]$ je $3x + 4y = 25$.
\end{example}

\subsection{Vzájemná poloha přímky a kružnice}

Polohu přímky a kružnice zjistíme dvěma způsoby. Buď porovnáme vzdálenost středu od přímky s poloměrem (jako v kapitole~\ref{kap:planimetrie}), nebo vyřešíme soustavu rovnic. Dosazením přímky do rovnice kružnice dostaneme kvadratickou rovnici a o počtu průsečíků rozhodne diskriminant.

\begin{definitionbox}
  Po dosazení přímky do rovnice kuželosečky dostaneme kvadratickou rovnici s diskriminantem $D$:
  \begin{itemize}
    \item $D > 0$: dva průsečíky, přímka je \textbf{sečna},
    \item $D = 0$: jeden společný bod, přímka je \textbf{tečna},
    \item $D < 0$: žádný společný bod, přímka je \textbf{nesečna} (vnější přímka).
  \end{itemize}
\end{definitionbox}

\begin{example}
  Přímka $y = x + 1$ a kružnice $x^2 + y^2 = 25$: $x^2 + {(x + 1)}^2 = 25$, tedy $2x^2 + 2x - 24 = 0$, neboli $x^2 + x - 12 = 0$, a tedy $x\in\{3, -4\}$. Průsečíky jsou $[3, 4]$ a $[-4, -3]$, přímka je sečna.
\end{example}

\begin{example}
  Pro která $k$ je přímka $y = x + k$ tečnou kružnice $x^2 + y^2 = 8$? Dosazením $2x^2 + 2kx + k^2 - 8 = 0$, diskriminant $D = 4k^2 - 8(k^2 - 8) = 64 - 4k^2$. Tečna pro $D = 0$, tedy $k = \pm 4$.

  Kontrola vzdáleností: vzdálenost středu $[0, 0]$ od přímky $x - y + k = 0$ je $\frac{|k|}{\sqrt{2}}$ a má být rovna $r = 2\sqrt{2}$, tedy $|k| = 4$.
\end{example}

\begin{problem}
  \begin{enumerate}[label=(\alph*)]
    \item Určete střed a poloměr kružnice $x^2 + y^2 + 6x - 2y + 6 = 0$.
    \item Napište rovnici kružnice se středem $[1, 2]$, která prochází bodem $[4, 6]$.
    \item Napište rovnici kružnice, jejímž průměrem je úsečka $AB$, kde $A = [-2, 1]$ a $B = [4, 9]$.
    \item Určete vzájemnou polohu kružnic $x^2 + y^2 = 9$ a ${(x - 5)}^2 + y^2 = 4$.
    \item Určete průsečíky kružnice ${(x - 1)}^2 + y^2 = 5$ s osou $y$.
  \end{enumerate}
\end{problem}

\begin{solution}
  \begin{enumerate}[label=(\alph*)]
    \item ${(x + 3)}^2 + {(y - 1)}^2 = 9 + 1 - 6 = 4$, střed $[-3, 1]$, $r = 2$.
    \item $r = 5$, ${(x - 1)}^2 + {(y - 2)}^2 = 25$.
    \item Střed je střed $AB$, tedy $[1, 5]$, $r = \frac{|AB|}{2} = \frac{10}{2} = 5$: ${(x - 1)}^2 + {(y - 5)}^2 = 25$.
    \item Vzdálenost středů je $5 = 3 + 2$, kružnice mají \textbf{vnější dotyk} v bodě $[3, 0]$.
    \item $x = 0$: $1 + y^2 = 5$, $y = \pm 2$. Průsečíky $[0, 2]$ a $[0, -2]$.
  \end{enumerate}
\end{solution}

\section{Elipsa}

\begin{definitionbox}
  \textbf{Elipsa} je množina bodů roviny, které mají od dvou daných bodů $F_1$, $F_2$ (\textbf{ohnisek}) konstantní součet vzdáleností $2a$, kde $2a > |F_1F_2|$:
  \begin{equation*}
    |XF_1| + |XF_2| = 2a.
  \end{equation*}
  \begin{itemize}
    \item $S$ je \textbf{střed} (střed úsečky $F_1F_2$), $a$ je \textbf{hlavní poloosa}, $b$ \textbf{vedlejší poloosa},
    \item $e = |SF_1|$ je \textbf{výstřednost} (excentricita) a platí $a^2 = b^2 + e^2$,
    \item \textbf{středová rovnice} elipsy se středem $[m, n]$ a hlavní osou rovnoběžnou s osou $x$:
    \begin{equation*}
      \frac{{(x - m)}^2}{a^2} + \frac{{(y - n)}^2}{b^2} = 1.
    \end{equation*}
    Pokud je hlavní osa rovnoběžná s osou $y$, prohodí se $a$ a $b$ ve jmenovatelích.
  \end{itemize}
\end{definitionbox}

\begin{center}
  \begin{tikzpicture}[scale=0.9]
    \draw[->] (-5.8,0) -- (6,0) node[right] {\small $x$};
    \draw[->] (0,-3.5) -- (0,3.7) node[above] {\small $y$};
    \draw[thick, colorofwisdom] (0,0) ellipse (5 and 3);
    \filldraw (-4,0) circle (2pt) node[below] {\small $F_1$};
    \filldraw (4,0) circle (2pt) node[below] {\small $F_2$};
    \coordinate (X) at ({5*cos(60)},{3*sin(60)});
    \filldraw (X) circle (2pt) node[above right] {\small $X$};
    \draw[dashed] (-4,0) -- (X) -- (4,0);
    \draw[thick] (0,0) -- (0,3) node[midway, left] {\small $b$};
    \draw[thick] (0,0) -- (4,0) node[midway, below] {\small $e$};
    \draw[gray] (0,3) -- (4,0) node[midway, above right] {\small $a$};
    \node[below left] at (-5,0) {\small $A$}; \node[below right] at (5,0) {\small $B$};
    \node[above left] at (0,3) {\small $C$}; \node[below left] at (0,-3) {\small $D$};
  \end{tikzpicture}
\end{center}

Proč $a^2 = b^2 + e^2$? Vrchol $C$ na vedlejší ose má od obou ohnisek stejnou vzdálenost a jejich součet je $2a$. Každá z nich je tedy $a$. Z pravoúhlého trojúhelníku $SF_2C$ s odvěsnami $e$ a $b$ a přeponou $a$ už vzorec plyne z Pythagorovy věty.

Definice elipsy dává i praktický návod, jak ji nakreslit. Do dvou bodů zapíchněte připínáčky, uvažte mezi ně volný provázek a tužkou ho napínejte. Tužka opíše elipsu. Tak se to mimochodem dělá i při zakládání eliptických záhonů. Pokud dáte oba připínáčky do jednoho bodu, dostanete kružnici. Kružnice je tedy elipsa s $e = 0$.

\begin{example}
  Elipsa $\frac{x^2}{25} + \frac{y^2}{9} = 1$ má $a = 5$, $b = 3$, $e = \sqrt{25 - 9} = 4$. Ohniska jsou $[\pm 4, 0]$, hlavní vrcholy $[\pm 5, 0]$ a vedlejší vrcholy $[0, \pm 3]$. (To je přesně elipsa na obrázku.)

  Rovnice $9x^2 + 25y^2 - 18x - 100y - 116 = 0$ po doplnění na čtverec:
  \begin{equation*}
    9{(x - 1)}^2 - 9 + 25{(y - 2)}^2 - 100 = 116\quad\Leftrightarrow\quad\frac{{(x - 1)}^2}{25} + \frac{{(y - 2)}^2}{9} = 1.
  \end{equation*}
  Je to stejná elipsa, jen posunutá do středu $[1, 2]$.
\end{example}

\begin{note}
  Poměr $\varepsilon = \frac{e}{a}\in[0, 1)$ se nazývá \textbf{číselná výstřednost} a říká, jak moc je elipsa `zploštělá'. Oběžná dráha Země kolem Slunce má $\varepsilon\approx 0{,}017$, je tedy skoro kruhová. Halleyova kometa má $\varepsilon\approx 0{,}967$, její dráha je velice protáhlá. Slunce leží v jednom ohnisku dráhy, druhé ohnisko je prázdné. To je první Keplerův zákon.
\end{note}

\section{Hyperbola}

\begin{definitionbox}
  \textbf{Hyperbola} je množina bodů roviny, které mají od dvou daných bodů $F_1$, $F_2$ konstantní \textbf{absolutní hodnotu rozdílu} vzdáleností $2a$, kde $0 < 2a < |F_1F_2|$:
  \begin{equation*}
    \big||XF_1| - |XF_2|\big| = 2a.
  \end{equation*}
  \begin{itemize}
    \item $a$ je \textbf{hlavní poloosa}, $e = |SF_1|$ výstřednost, \textbf{vedlejší poloosa} $b$ je dána vztahem $e^2 = a^2 + b^2$,
    \item \textbf{středová rovnice} (hlavní osa rovnoběžná s $x$):
    \begin{equation*}
      \frac{{(x - m)}^2}{a^2} - \frac{{(y - n)}^2}{b^2} = 1,
    \end{equation*}
    \item hyperbola má dvě větve a dvě \textbf{asymptoty} $y - n = \pm\frac{b}{a}(x - m)$.
  \end{itemize}
\end{definitionbox}

\begin{center}
  \begin{tikzpicture}[scale=0.6]
    \draw[->] (-7,0) -- (7.2,0) node[right] {\small $x$};
    \draw[->] (0,-5) -- (0,5.2) node[above] {\small $y$};
    \draw[dashed, gray] (-6.6,-4.95) -- (6.6,4.95);
    \draw[dashed, gray] (-6.6,4.95) -- (6.6,-4.95) node[right] {\small asymptoty};
    \draw[gray] (-4,-3) rectangle (4,3);
    \draw[thick, colorofwisdom, domain=-1.1:1.1, samples=60, variable=\t] plot ({4*cosh(\t)}, {3*sinh(\t)});
    \draw[thick, colorofwisdom, domain=-1.1:1.1, samples=60, variable=\t] plot ({-4*cosh(\t)}, {3*sinh(\t)});
    \filldraw (5,0) circle (3pt) node[below right] {\small $F_2$};
    \filldraw (-5,0) circle (3pt) node[below left] {\small $F_1$};
    \node[below right] at (4,0) {\small $a$};
    \node[left] at (0,3) {\small $b$};
  \end{tikzpicture}
\end{center}

\begin{example}
  Hyperbola $\frac{x^2}{16} - \frac{y^2}{9} = 1$ má $a = 4$, $b = 3$, $e = \sqrt{16 + 9} = 5$. Ohniska $[\pm 5, 0]$, vrcholy $[\pm 4, 0]$, asymptoty $y = \pm\frac{3}{4}x$. Asymptoty jsou úhlopříčky obdélníku o stranách $2a$ a $2b$ se středem ve středu hyperboly (šedý obdélník na obrázku). To je nejrychlejší způsob, jak hyperbolu nakreslit.
\end{example}

\begin{note}
  Graf nepřímé úměrnosti $y = \frac{k}{x}$ z kapitoly~\ref{kap:elementarni} je také hyperbola, jen otočená o $45^\circ$. Asymptotami jsou osy souřadnic a $a = b$, proto se jí říká \textbf{rovnoosá}. Hyperboly se ve skutečném světě objevují méně často než elipsy, ale potkáte je: po hyperbole letí kometa, která má dost energie, aby ze Sluneční soustavy unikla, a na hyperbolách stál starší navigační systém LORAN, který určoval polohu lodí z rozdílu časů, za které dorazily signály ze dvou vysílačů.
\end{note}

\section{Parabola}

\begin{definitionbox}
  \textbf{Parabola} je množina bodů roviny, které mají stejnou vzdálenost od daného bodu $F$ (\textbf{ohniska}) a dané přímky $d$ (\textbf{řídicí přímky}), která bodem $F$ neprochází:
  \begin{equation*}
    |XF| = |Xd|.
  \end{equation*}
  \begin{itemize}
    \item $p = |Fd|$ je \textbf{parametr} paraboly, \textbf{vrchol} $V$ leží uprostřed mezi $F$ a $d$,
    \item \textbf{vrcholová rovnice} paraboly s vrcholem $[m, n]$:
    \begin{align*}
      {(y - n)}^2 &= \pm 2p(x - m) &&\text{(osa rovnoběžná s osou } x\text{, otevřená doprava/doleva)},\\
      {(x - m)}^2 &= \pm 2p(y - n) &&\text{(osa rovnoběžná s osou } y\text{, otevřená nahoru/dolů)}.
    \end{align*}
  \end{itemize}
\end{definitionbox}

\begin{example}
  Odvoďme rovnici paraboly s ohniskem $F = \left[0, \frac{p}{2}\right]$ a řídicí přímkou $y = -\frac{p}{2}$. Bod $X = [x, y]$ má od ohniska vzdálenost $\sqrt{x^2 + {\left(y - \frac{p}{2}\right)}^2}$ a od řídicí přímky $\left|y + \frac{p}{2}\right|$. Umocníme rovnost vzdáleností:
  \begin{equation*}
    x^2 + y^2 - py + \frac{p^2}{4} = y^2 + py + \frac{p^2}{4}\quad\Leftrightarrow\quad x^2 = 2py.
  \end{equation*}
  Graf funkce $y = x^2$ je tedy parabola s $2p = 1$, tedy $p = \frac{1}{2}$, ohniskem $\left[0, \frac{1}{4}\right]$ a řídicí přímkou $y = -\frac{1}{4}$. Parabola z kapitoly~\ref{kap:elementarni} a parabola z této kapitoly jsou opravdu tatáž křivka.
\end{example}

\begin{center}
  \begin{tikzpicture}[scale=1.2]
    \draw[->] (-2.5,0) -- (2.7,0) node[right] {\small $x$};
    \draw[->] (0,-0.8) -- (0,3.3) node[above] {\small $y$};
    \draw[thick, colorofwisdom, domain=-1.75:1.75, samples=60] plot (\x, {\x*\x});
    \draw[thick] (-2.5,-0.25) -- (2.5,-0.25) node[right] {\small $d$};
    \filldraw (0,0.25) circle (1.5pt) node[right] {\small $F$};
    \coordinate (X) at (1.3,1.69);
    \filldraw (X) circle (1.5pt) node[right] {\small $X$};
    \draw[dashed] (0,0.25) -- (X) -- (1.3,-0.25);
    \foreach \x in {-1.2,-0.6,0.6,1.2} \draw[gray, ->] (\x,3) -- (\x,{\x*\x}) -- (0,0.25);
  \end{tikzpicture}
\end{center}

Šedé paprsky na obrázku ukazují vlastnost, kterou jsme zmínili u Archimédových zrcadel v kapitole~\ref{kap:elementarni}: paprsky rovnoběžné s osou se po odrazu od paraboly setkají v ohnisku.

\begin{example}
  Rovnice $x^2 - 4x - 8y + 12 = 0$ po doplnění na čtverec dává ${(x - 2)}^2 - 4 - 8y + 12 = 0$, tedy ${(x - 2)}^2 = 8(y - 1)$. Vrchol $V = [2, 1]$, $2p = 8$, tedy $p = 4$. Parabola je otevřená nahoru, ohnisko je o $\frac{p}{2} = 2$ nad vrcholem, $F = [2, 3]$, a řídicí přímka $y = -1$.
\end{example}

\section{Rozpoznání kuželosečky a vzájemné polohy}

\begin{definitionbox}
  Rovnici $Ax^2 + Cy^2 + Dx + Ey + F = 0$ (bez členu $xy$) upravíme doplněním na čtverec. Podle koeficientů $A$ a $C$ jde o:
  \begin{itemize}
    \item \textbf{kružnici}, pokud $A = C\neq 0$,
    \item \textbf{elipsu}, pokud $A$ a $C$ mají stejné znaménko a $A\neq C$,
    \item \textbf{hyperbolu}, pokud $A$ a $C$ mají opačná znaménka,
    \item \textbf{parabolu}, pokud je právě jedno z $A$, $C$ nulové.
  \end{itemize}
  Ve výjimečných případech (tzv. \textbf{degenerované kuželosečky}) může rovnice popisovat jen bod, dvojici přímek nebo prázdnou množinu.
\end{definitionbox}

\begin{note}
  U paraboly a hyperboly může mít přímka s kuželosečkou jediný společný bod, a přesto nebýt tečnou. U paraboly je to přímka rovnoběžná s osou, u hyperboly přímka rovnoběžná s asymptotou. Po dosazení v tom případě nedostaneme kvadratickou, ale lineární rovnici, takže diskriminant nemá smysl.
\end{note}

\begin{problem}
  Rozhodněte, o jakou kuželosečku jde, a určete její základní parametry:
  \begin{enumerate}[label=(\alph*)]
    \item $4x^2 - 9y^2 = 36$,
    \item $x^2 + 4y^2 - 4 = 0$,
    \item $y^2 - 4y - 4x = 0$,
    \item $x^2 + y^2 - 10x = 0$,
    \item $4x^2 + y^2 + 8x - 6y + 9 = 0$.
  \end{enumerate}
\end{problem}

\begin{solution}
  \begin{enumerate}[label=(\alph*)]
    \item $\frac{x^2}{9} - \frac{y^2}{4} = 1$, hyperbola, $a = 3$, $b = 2$, $e = \sqrt{13}$, asymptoty $y = \pm\frac{2}{3}x$.
    \item $\frac{x^2}{4} + y^2 = 1$, elipsa, $a = 2$, $b = 1$, $e = \sqrt{3}$, ohniska $[\pm\sqrt{3}, 0]$.
    \item ${(y - 2)}^2 = 4x + 4 = 4(x + 1)$, parabola, $V = [-1, 2]$, $p = 2$, otevřená doprava, $F = [0, 2]$, řídicí přímka $x = -2$.
    \item ${(x - 5)}^2 + y^2 = 25$, kružnice se středem $[5, 0]$ a poloměrem $5$.
    \item $4{(x + 1)}^2 - 4 + {(y - 3)}^2 - 9 + 9 = 0$, tedy $\frac{{(x + 1)}^2}{1} + \frac{{(y - 3)}^2}{4} = 1$. Elipsa se středem $[-1, 3]$, $a = 2$ (ve směru osy $y$), $b = 1$.
  \end{enumerate}
\end{solution}

\begin{problem}
  \begin{enumerate}[label=(\alph*)]
    \item Určete průsečíky přímky $x - 2y + 2 = 0$ s elipsou $x^2 + 4y^2 = 20$.
    \item Pro která $q$ je přímka $y = 2x + q$ tečnou paraboly $y^2 = 8x$?
  \end{enumerate}
\end{problem}

\begin{solution}
  \begin{enumerate}[label=(\alph*)]
    \item $x = 2y - 2$, dosazením ${(2y - 2)}^2 + 4y^2 = 20$, $8y^2 - 8y - 16 = 0$, $y^2 - y - 2 = 0$, $y\in\{2, -1\}$. Průsečíky $[2, 2]$ a $[-4, -1]$.
    \item ${(2x + q)}^2 = 8x$, tedy $4x^2 + (4q - 8)x + q^2 = 0$. $D = {(4q - 8)}^2 - 16q^2 = -64q + 64 = 0$, tedy $q = 1$.
  \end{enumerate}
\end{solution}

\section{Odbočka --- Hypatia}

Apollóniovo dílo o kuželosečkách se nám dochovalo i díky komentářům, které k němu ve 4. a 5. století psala Hypatia z Alexandrie. Byla to jedna z prvních žen v dějinách matematiky, o kterých víme. Vedla v Alexandrii novoplatónskou filozofickou školu, učila matematiku a astronomii a byla tak uznávaná, že za ní přijížděli studenti z celého Středomoří. Radili se s ní i městští úředníci.

To se jí stalo osudným. V roce 415 se Alexandrie zmítala v mocenském sporu mezi římským prefektem Orestem a biskupem Kyrillem. Hypatia byla Orestovou přítelkyní a rádkyní a mezi Kyrillovými stoupenci se rozšířila pověst, že je to právě ona, kdo brání smíru, a že je čarodějnice. Dav fanatiků ji podle dobových zpráv vytáhl z vozu, odvlekl do kostela, svlékl ji a zaživa ji rozdrásal střepy z hliněných nádob. Její ostatky pak spálili.

Dílo Hypatie se nedochovalo celé. Víme, že psala komentáře k Apollóniovým \textit{Kuželosečkám}, k Diofantově \textit{Aritmetice} a k Ptolemaiovým astronomickým tabulkám. Pro mnoho historiků její smrt symbolizuje konec antické vědy. Další ženu, která by se do dějin matematiky zapsala podobně výrazně, známe až o více než tisíc let později.

\section{Shrnutí a otázky k ústní zkoušce}

\begin{problem}
  Odpovězte vlastními slovy:
  \begin{enumerate}
    \item Co jsou kuželosečky a jak vznikají řezem kužele?
    \item Definujte kružnici. Napište její středovou a obecnou rovnici. Jak z obecné rovnice určíte střed a poloměr?
    \item Jak zjistíte vzájemnou polohu přímky a kružnice? Jak vypadá rovnice tečny?
    \item Definujte elipsu. Jaký je vztah mezi $a$, $b$ a $e$? Napište středovou rovnici.
    \item Definujte hyperbolu. Co jsou asymptoty a jak je najdete?
    \item Definujte parabolu. Odvoďte její vrcholovou rovnici. Jak souvisí s grafem kvadratické funkce?
    \item Jak z obecné rovnice poznáte, o jakou kuželosečku jde?
    \item Kdy může mít přímka s kuželosečkou jediný společný bod, a přesto nebýt tečnou?
    \item Uveďte příklady kuželoseček v praxi.
  \end{enumerate}
\end{problem}

\koneckapitoly

\chapter{Analytická geometrie v prostoru}
\label{kap:prostor}
\mathquote{Geometry is knowledge of the eternally existent.}{Platón}

V kapitole~\ref{kap:stereometrie} jsme řešili prostorové úlohy tak, že jsme hledali vhodné pravoúhlé trojúhelníky. To funguje, ale vyžaduje to dobrou prostorovou představivost a trochu štěstí. Analytická geometrie v prostoru nabízí jinou cestu: zavedeme souřadnice a všechno spočítáme. Žádné hledání trojúhelníků, žádné kreslení. Jen vektory a rovnice. Pro ty z nás, kdo si v hlavě neumí otočit ani židli, je to vysvobození.

\section{Souřadnice a vektory v prostoru}

\begin{definitionbox}
  \textbf{Kartézská soustava souřadnic} v prostoru má tři navzájem kolmé osy $x$, $y$, $z$. Bod je určen trojicí $A = [a_1, a_2, a_3]$, vektor trojicí $\vec{u} = (u_1, u_2, u_3)$. Všechny vzorce z roviny se přenesou přidáním třetí souřadnice:
  \begin{align*}
    |AB| &= \sqrt{{(b_1 - a_1)}^2 + {(b_2 - a_2)}^2 + {(b_3 - a_3)}^2}, & S_{AB} &= \left[\tfrac{a_1 + b_1}{2}, \tfrac{a_2 + b_2}{2}, \tfrac{a_3 + b_3}{2}\right],\\
    |\vec{u}| &= \sqrt{u_1^2 + u_2^2 + u_3^2}, & \vec{u}\cdot\vec{v} &= u_1v_1 + u_2v_2 + u_3v_3.
  \end{align*}
  I zde platí $\vec{u}\cdot\vec{v} = |\vec{u}||\vec{v}|\cos\varphi$ a vektory jsou kolmé, právě když $\vec{u}\cdot\vec{v} = 0$.
\end{definitionbox}

Jak si bod v prostoru představit? Z počátku ujdeme $a_1$ ve směru osy $x$, pak $a_2$ ve směru osy $y$ a nakonec vystoupáme o $a_3$ nahoru. Bod je roh `krabice' o rozměrech $a_1\times a_2\times a_3$ a délka vektoru $\overrightarrow{OA}$ je její tělesová úhlopříčka. Proto vzorec pro délku vypadá jako Pythagorova věta použitá dvakrát: nejdřív úhlopříčka podstavy $\sqrt{a_1^2 + a_2^2}$, pak s výškou $a_3$.

\begin{center}
  \begin{tikzpicture}[krychle, scale=1.1]
    \draw[->] (0,0,0) -- (3.3,0,0) node[right] {\small $x$};
    \draw[->] (0,0,0) -- (0,4.2,0) node[above right] {\small $y$};
    \draw[->] (0,0,0) -- (0,0,2.9) node[above] {\small $z$};
    \coordinate (A) at (2,3,2);
    \draw[dashed, gray] (2,0,0) -- (2,3,0) -- (0,3,0);
    \draw[dashed, gray] (2,0,0) -- (2,0,2) (2,3,0) -- (2,3,2) (0,3,0) -- (0,3,2);
    \draw[dashed, gray] (0,0,2) -- (2,0,2) -- (2,3,2) -- (0,3,2) -- cycle;
    \draw[dotted, thick] (0,0,0) -- (2,3,0) node[midway, below right, font=\scriptsize] {$\sqrt{a_1^2 + a_2^2}$};
    \draw[->, thick, colorofwisdom] (0,0,0) -- (A);
    \fill (A) circle (1.5pt) node[above right] {\small $A = [2, 3, 2]$};
    \node[below left] at (0,0,0) {\small $O$};
    \node[below] at (2,0,0) {\scriptsize $a_1 = 2$};
    \node[right] at (0,3,0) {\scriptsize $a_2 = 3$};
    \node[left] at (0,0,2) {\scriptsize $a_3 = 2$};
  \end{tikzpicture}
\end{center}

\begin{example}
  Pro $A = [1, 2, 3]$ a $B = [3, -1, 9]$ je $\overrightarrow{AB} = (2, -3, 6)$, $|AB| = \sqrt{4 + 9 + 36} = 7$ a $S_{AB} = \left[2, \frac{1}{2}, 6\right]$.
\end{example}

\subsection{Vektorový součin}

V prostoru existuje ještě jeden součin vektorů, který v rovině nemá obdobu. Výsledkem není číslo, ale vektor.

\begin{definitionbox}
  \textbf{Vektorový součin} vektorů $\vec{u} = (u_1, u_2, u_3)$ a $\vec{v} = (v_1, v_2, v_3)$ je vektor
  \begin{equation*}
    \vec{u}\times\vec{v} = (u_2v_3 - u_3v_2,\ u_3v_1 - u_1v_3,\ u_1v_2 - u_2v_1).
  \end{equation*}
  \textbf{Vlastnosti}:
  \begin{itemize}
    \item $\vec{u}\times\vec{v}$ je \textbf{kolmý} na $\vec{u}$ i na $\vec{v}$,
    \item $|\vec{u}\times\vec{v}| = |\vec{u}||\vec{v}|\sin\varphi$, což je \textbf{obsah rovnoběžníku} určeného vektory $\vec{u}$, $\vec{v}$,
    \item $\vec{u}\times\vec{v} = -\vec{v}\times\vec{u}$ (není komutativní!),
    \item $\vec{u}\times\vec{v} = \vec{o}$, právě když jsou $\vec{u}$ a $\vec{v}$ kolineární.
  \end{itemize}
\end{definitionbox}


\begin{center}
  \begin{tikzpicture}[krychle, scale=1.25]
    \coordinate (O) at (0,0,0); \coordinate (U) at (2.5,0,0); \coordinate (V) at (0.8,1.8,0); \coordinate (UV) at (3.3,1.8,0);
    \draw[->, thick, gray] (O) -- (0,0,-1.1) node[below, font=\small] {$\vec{v}\times\vec{u}$};
    \fill[colorofwisdom!20] (O) -- (U) -- (UV) -- (V) -- cycle;
    \draw[dashed] (U) -- (UV) -- (V);
    \draw[->, thick] (O) -- (U) node[below right] {\small $\vec{u}$};
    \draw[->, thick] (O) -- (V) node[left] {\small $\vec{v}$};
    \node[font=\scriptsize, align=center] at (1.75,0.9,0) {obsah\\$|\vec{u}\times\vec{v}|$};
    \draw[->, thick, colorofwisdom] (O) -- (0,0,2.1) node[above, font=\small] {$\vec{u}\times\vec{v}$};
    \draw (0.25,0,0) -- (0.25,0,0.25) -- (0,0,0.25);
  \end{tikzpicture}
\end{center}
Na obrázku vidíte všechny vlastnosti najednou: výsledek trčí kolmo z roviny, ve které leží $\vec{u}$ a $\vec{v}$, jeho délka je obsah vybarveného rovnoběžníku a po prohození pořadí míří opačně (šedá šipka). Když jsou $\vec{u}$ a $\vec{v}$ rovnoběžné, rovnoběžník se `splácne' na úsečku s nulovým obsahem.

Vzorec vypadá hrozivě, ale dá se zapamatovat jako `determinant'. Každá složka je determinant $2\times 2$ z ostatních dvou souřadnic: první složka vznikne zakrytím prvního sloupce, druhá zakrytím druhého (s obráceným znaménkem), třetí zakrytím třetího:
\begin{equation*}
  \vec{u}\times\vec{v} = \left(\begin{vmatrix} u_2 & u_3\\ v_2 & v_3 \end{vmatrix},\ -\begin{vmatrix} u_1 & u_3\\ v_1 & v_3 \end{vmatrix},\ \begin{vmatrix} u_1 & u_2\\ v_1 & v_2 \end{vmatrix}\right).
\end{equation*}
Kolmost si ověříte přímo: $\vec{u}\cdot(\vec{u}\times\vec{v}) = u_1(u_2v_3 - u_3v_2) + u_2(u_3v_1 - u_1v_3) + u_3(u_1v_2 - u_2v_1) = 0$, protože se všechny členy po dvou vyruší. Směr výsledku určuje \textbf{pravidlo pravé ruky}: pokud prsty pravé ruky ukazují od $\vec{u}$ k $\vec{v}$, palec ukazuje směr $\vec{u}\times\vec{v}$. Fyzici tohle pravidlo používají u magnetické síly a točivého momentu tak často, že je u zkoušek z fyziky vidět sál plný studentů, kteří kroutí pravou rukou ve vzduchu.

\begin{example}
  Pro $\vec{u} = (1, 2, 3)$ a $\vec{v} = (4, 5, 6)$:
  \begin{equation*}
    \vec{u}\times\vec{v} = (2\cdot 6 - 3\cdot 5,\ 3\cdot 4 - 1\cdot 6,\ 1\cdot 5 - 2\cdot 4) = (-3, 6, -3).
  \end{equation*}
  Kontrola: $(1, 2, 3)\cdot(-3, 6, -3) = -3 + 12 - 9 = 0$ a $(4, 5, 6)\cdot(-3, 6, -3) = -12 + 30 - 18 = 0$.
\end{example}

\section{Přímka a rovina v prostoru}

\begin{definitionbox}
  \textbf{Přímka} procházející bodem $A$ se směrovým vektorem $\vec{u}$ má \textbf{parametrické vyjádření}
  \begin{equation*}
    X = A + t\vec{u},\quad t\in\mathbb{R},\qquad\text{tedy}\qquad x = a_1 + tu_1,\quad y = a_2 + tu_2,\quad z = a_3 + tu_3.
  \end{equation*}
  Přímka v prostoru \textbf{nemá} obecnou rovnici. Jedna lineární rovnice v prostoru totiž popisuje rovinu, ne přímku. Přímku lze popsat nejvýš jako průsečnici dvou rovin.

  \vspace{0.2cm}

  \textbf{Rovina} procházející bodem $A$ a určená dvěma nekolineárními vektory $\vec{u}$, $\vec{v}$ má \textbf{parametrické vyjádření}
  \begin{equation*}
    X = A + t\vec{u} + s\vec{v},\quad t, s\in\mathbb{R}.
  \end{equation*}
  \textbf{Obecná rovnice roviny} s normálovým vektorem $\vec{n} = (a, b, c)$ je
  \begin{equation*}
    ax + by + cz + d = 0.
  \end{equation*}
  Normálový vektor roviny určené vektory $\vec{u}$, $\vec{v}$ najdeme jako $\vec{n} = \vec{u}\times\vec{v}$.
\end{definitionbox}


\begin{center}
  \begin{tikzpicture}[scale=1.1]
    \begin{scope}[krychle]
      \draw[thick] (-0.94,-0.5,-0.4) -- (2.84,1.6,1.65) node[right] {\small $p$};
      \draw[->, very thick, colorofwisdom] (0.5,0.3,0.4) -- (1.4,0.8,0.9) node[midway, above left] {\small $\vec{u}$};
      \foreach \t in {-1,0,1,2} {
        \fill ({0.5+0.9*\t},{0.3+0.5*\t},{0.4+0.5*\t}) circle (1.5pt) node[below right, font=\scriptsize] {$t = \t$};
      }
      \node[above left] at (0.5,0.3,0.4) {\small $A$};
    \end{scope}
    \begin{scope}[xshift=5.2cm, krychle]
      \fill[colorofwisdom!15] (-0.8,-0.8,-0.4) -- (2.4,-0.8,0.24) -- (2.4,2.4,1.2) -- (-0.8,2.4,0.56) -- cycle;
      \draw[gray] (-0.8,-0.8,-0.4) -- (2.4,-0.8,0.24) -- (2.4,2.4,1.2) -- (-0.8,2.4,0.56) -- cycle;
      \draw[->, thick] (0,0,0) -- (2,0,0.4) node[below right] {\small $\vec{u}$};
      \draw[->, thick] (0,0,0) -- (0,2,0.6) node[above left] {\small $\vec{v}$};
      \draw[dashed] (0,0,0) -- (1.4,0,0.28) -- (1.4,1,0.58);
      \fill (1.4,1,0.58) circle (1.5pt) node[above right, font=\scriptsize] {$X = A + t\vec{u} + s\vec{v}$};
      \draw[->, thick, colorofwisdom] (0,0,0) -- (-0.24,-0.36,1.2) node[above] {\small $\vec{n} = \vec{u}\times\vec{v}$};
      \fill (0,0,0) circle (1.5pt) node[below left] {\small $A$};
      \node at (2.2,2.2,1.1) [below left] {\small $\rho$};
    \end{scope}
  \end{tikzpicture}
\end{center}
Parametr $t$ funguje jako `čas': v čase $t = 0$ stojíme v bodě $A$ a za každou jednotku času ujdeme vektor $\vec{u}$. Záporné $t$ znamená cestu pozpátku. U roviny máme dva nezávislé `ovladače' $t$ a $s$, a tedy dva rozměry, ve kterých se můžeme pohybovat. Normálový vektor $\vec{n}$ trčí z roviny kolmo jako stožár z podlahy.

Proč má rovina obecnou rovnici a přímka ne? Rovnice $\vec{n}\cdot\overrightarrow{AX} = 0$ říká, že vektor $\overrightarrow{AX}$ je kolmý na $\vec{n}$. V rovině jsou všechny takové body $X$ na jedné přímce. V prostoru tvoří všechny vektory kolmé na $\vec{n}$ celou rovinu. Jedna podmínka v prostoru `ubere' jen jeden rozměr ze tří.

\begin{example}
  Napišme obecnou rovnici roviny $ABC$, kde $A = [1, 0, 0]$, $B = [0, 2, 0]$ a $C = [0, 0, 3]$.

  Vektory $\vec{u} = \overrightarrow{AB} = (-1, 2, 0)$ a $\vec{v} = \overrightarrow{AC} = (-1, 0, 3)$. Normálový vektor
  \begin{equation*}
    \vec{n} = \vec{u}\times\vec{v} = (2\cdot 3 - 0\cdot 0,\ 0\cdot(-1) - (-1)\cdot 3,\ (-1)\cdot 0 - 2\cdot(-1)) = (6, 3, 2).
  \end{equation*}
  Rovnice $6x + 3y + 2z + d = 0$, dosazením $A$: $d = -6$. Tedy $6x + 3y + 2z - 6 = 0$. Kontrola s $B$: $6 - 6 = 0$, s $C$: $6 - 6 = 0$.

  Po vydělení šesti dostaneme hezký tvar $\frac{x}{1} + \frac{y}{2} + \frac{z}{3} = 1$, tzv. \textbf{úsekový tvar}, ze kterého jsou průsečíky s osami vidět okamžitě. A obsah trojúhelníku $ABC$ je polovina obsahu rovnoběžníku: $S = \frac{1}{2}|\vec{n}| = \frac{1}{2}\sqrt{36 + 9 + 4} = \frac{7}{2}$.
\end{example}

\begin{center}
  \begin{tikzpicture}[krychle, scale=1.15]
    \draw[->] (0,0,0) -- (1.8,0,0) node[right] {\small $x$};
    \draw[->] (0,0,0) -- (0,2.9,0) node[above right] {\small $y$};
    \draw[->] (0,0,0) -- (0,0,3.6) node[above] {\small $z$};
    \fill[colorofwisdom!30, opacity=0.8] (1,0,0) -- (0,2,0) -- (0,0,3) -- cycle;
    \draw[thick, colorofwisdom] (1,0,0) -- (0,2,0) -- (0,0,3) -- cycle;
    \fill (1,0,0) circle (1.5pt) node[below] {\small $A = [1, 0, 0]$};
    \fill (0,2,0) circle (1.5pt) node[right] {\small $B = [0, 2, 0]$};
    \fill (0,0,3) circle (1.5pt) node[left] {\small $C = [0, 0, 3]$};
    \node[font=\small] at (2.6,0,2.2) {$\dfrac{x}{1} + \dfrac{y}{2} + \dfrac{z}{3} = 1$};
  \end{tikzpicture}
\end{center}
Rovina `odřízne' z prvního oktantu trojúhelníkový kousek. Čísla ve jmenovatelích úsekového tvaru jsou přesně místa, kde rovina protíná osy.

\section{Vzájemné polohy}

\begin{definitionbox}
  \textbf{Dvě přímky} $p: X = A + t\vec{u}$ a $q: X = B + s\vec{v}$:
  \begin{itemize}
    \item pokud jsou $\vec{u}$ a $\vec{v}$ kolineární, jsou přímky \textbf{rovnoběžné} (různé, nebo totožné, pokud $B\in p$),
    \item jinak řešíme soustavu $A + t\vec{u} = B + s\vec{v}$ (tři rovnice pro dvě neznámé). Pokud má řešení, jsou \textbf{různoběžné}, pokud nemá, jsou \textbf{mimoběžné}.
  \end{itemize}

  \textbf{Přímka} $p: X = A + t\vec{u}$ a \textbf{rovina} $\rho: \vec{n}\cdot X + d = 0$:
  \begin{itemize}
    \item pokud $\vec{n}\cdot\vec{u}\neq 0$, jsou \textbf{různoběžné}. Průsečík najdeme dosazením parametrického vyjádření přímky do rovnice roviny,
    \item pokud $\vec{n}\cdot\vec{u} = 0$, je přímka s rovinou \textbf{rovnoběžná}, nebo v ní \textbf{leží} (pokud $A\in\rho$).
  \end{itemize}

  \textbf{Dvě roviny}: jsou \textbf{rovnoběžné}, pokud mají kolineární normálové vektory (a totožné, pokud jsou rovnice násobkem jedna druhé). Jinak jsou \textbf{různoběžné} a průsečnici najdeme vyřešením soustavy dvou rovnic se třemi neznámými (jedna neznámá bude parametr).
\end{definitionbox}

\begin{center}
  \begin{tikzpicture}[scale=1]
    \begin{scope}[krychle]
      \draw[thick] (0,0,0) -- (2,1,1) node[right] {\small $p$};
      \draw[thick] (2,0,0) -- (0,1,1) node[left] {\small $q$};
      \fill (1,0.5,0.5) circle (1.6pt);
      \node[font=\small] at (1,0,-0.6) {různoběžné};
    \end{scope}
    \begin{scope}[xshift=4cm, krychle]
      \draw[thick] (0,0,0) -- (2,0.5,1) node[right] {\small $p$};
      \draw[thick] (0,1,0.3) -- (2,1.5,1.3) node[right] {\small $q$};
      \node[font=\small] at (1,0,-0.6) {rovnoběžné};
    \end{scope}
    \begin{scope}[xshift=8cm, krychle]
      \draw[gray!60, dashed] (0,0,0) -- (2,0,0) -- (2,1.5,0) -- (0,1.5,0) -- cycle;
      \draw[gray!60, dashed] (0,0,1.2) -- (2,0,1.2) -- (2,1.5,1.2) -- (0,1.5,1.2) -- cycle;
      \draw[gray!60, dashed] (0,0,0) -- (0,0,1.2) (2,0,0) -- (2,0,1.2) (2,1.5,0) -- (2,1.5,1.2) (0,1.5,0) -- (0,1.5,1.2);
      \draw[thick] (-0.3,0.4,0) -- (2.3,0.4,0) node[right] {\small $p$};
      \draw[thick] (1,-0.3,1.2) -- (1,1.9,1.2) node[above right] {\small $q$};
      \node[font=\small] at (1,0,-0.6) {mimoběžné};
    \end{scope}
  \end{tikzpicture}
\end{center}
Mimoběžky jsou to, co v rovině nemáme: dvě přímky, které se nikdy neprotnou, a přesto nejsou rovnoběžné. Představte si přímku vedenou po podlaze podél jedné stěny a přímku po stropě podél stěny sousední. Krabice na obrázku pomáhá: kdyby tam nebyla, na papíře by se přímky zdánlivě křížily. U obrázků prostoru se proto vždy vyplatí přikreslit nějaký `rám'.

\begin{center}
  \begin{tikzpicture}[scale=1]
    \begin{scope}[krychle]
      \fill[colorofwisdom!15] (0,0,0) -- (2.4,0,0) -- (2.4,2,0) -- (0,2,0) -- cycle; \draw[gray] (0,0,0) -- (2.4,0,0) -- (2.4,2,0) -- (0,2,0) -- cycle;
      \draw[thick, dashed] (1.1,0.9,-0.8) -- (1.34,0.98,0);
      \draw[thick] (1.34,0.98,0) -- (1.7,1.1,1.2) node[right] {\small $p$};
      \fill (1.34,0.98,0) circle (1.6pt);
      \node[font=\small] at (1.2,0,-1.1) {různoběžné (průsečík)};
    \end{scope}
    \begin{scope}[xshift=4.3cm, krychle]
      \fill[colorofwisdom!15] (0,0,0) -- (2.4,0,0) -- (2.4,2,0) -- (0,2,0) -- cycle; \draw[gray] (0,0,0) -- (2.4,0,0) -- (2.4,2,0) -- (0,2,0) -- cycle;
      \draw[thick] (0.2,0.3,0.8) -- (2.4,1.5,0.8) node[right] {\small $p$};
      \draw[dashed, gray] (0.2,0.3,0) -- (2.4,1.5,0);
      \draw[dotted] (0.2,0.3,0) -- (0.2,0.3,0.8) (2.4,1.5,0) -- (2.4,1.5,0.8);
      \node[font=\small] at (1.2,0,-1.1) {rovnoběžné};
    \end{scope}
    \begin{scope}[xshift=8.6cm, krychle]
      \fill[colorofwisdom!15] (0,0,0) -- (2.4,0,0) -- (2.4,2,0) -- (0,2,0) -- cycle; \draw[gray] (0,0,0) -- (2.4,0,0) -- (2.4,2,0) -- (0,2,0) -- cycle;
      \draw[thick] (0.2,0.3,0) -- (2.2,1.7,0) node[above] {\small $p$};
      \node[font=\small] at (1.2,0,-1.1) {přímka leží v rovině};
    \end{scope}
  \end{tikzpicture}
\end{center}
Podmínka $\vec{n}\cdot\vec{u} = 0$ znamená, že směr přímky je kolmý na normálu, tedy `vodorovný' vzhledem k rovině. Taková přímka buď letí nad rovinou jako letadlo v konstantní výšce, nebo po ní jede jako auto po silnici.

\begin{example}
  Rozhodněme o poloze přímek $p: X = [1, 0, 0] + t(1, 1, 0)$ a $q: X = [0, 1, 1] + s(1, -1, 0)$. Směrové vektory kolineární nejsou. Soustava:
  \begin{align*}
    1 + t &= s,\\
    t &= 1 - s,\\
    0 &= 1.
  \end{align*}
  Třetí rovnice je spor, takže přímky jsou \textbf{mimoběžné}. Kdyby $q$ procházela bodem $[0, 1, 0]$, třetí rovnice by zněla $0 = 0$ a z prvních dvou bychom dostali $t = 0$, $s = 1$ a průsečík $[1, 0, 0]$.
\end{example}

\begin{example}
  Průsečík přímky $X = [1, 1, 1] + t(1, 2, 3)$ s rovinou $x + y + z = 12$: dosadíme $(1 + t) + (1 + 2t) + (1 + 3t) = 12$, tedy $6t = 9$ a $t = \frac{3}{2}$. Průsečík je $\left[\frac{5}{2}, 4, \frac{11}{2}\right]$.
\end{example}

\begin{example}
  Průsečnice rovin $x + y + z = 3$ a $x - y = 1$: zvolíme $y = t$, pak $x = 1 + t$ a $z = 3 - x - y = 2 - 2t$. Průsečnice je $X = [1, 0, 2] + t(1, 1, -2)$. Kontrola: směrový vektor $(1, 1, -2)$ je kolmý na oba normálové vektory $(1, 1, 1)$ a $(1, -1, 0)$, a tedy je rovnoběžný s $\vec{n}_1\times\vec{n}_2$.
\end{example}

\section{Metrické úlohy}

\begin{definitionbox}
  \textbf{Vzdálenost bodu} $M = [m_1, m_2, m_3]$ od roviny $ax + by + cz + d = 0$:
  \begin{equation*}
    v = \frac{|am_1 + bm_2 + cm_3 + d|}{\sqrt{a^2 + b^2 + c^2}}.
  \end{equation*}
  \textbf{Vzdálenost mimoběžek} $p: X = A + t\vec{u}$ a $q: X = B + s\vec{v}$:
  \begin{equation*}
    v = \frac{|\overrightarrow{AB}\cdot(\vec{u}\times\vec{v})|}{|\vec{u}\times\vec{v}|}.
  \end{equation*}
  \textbf{Odchylky}:
  \begin{align*}
    \text{dvou přímek:}\quad\cos\varphi &= \frac{|\vec{u}\cdot\vec{v}|}{|\vec{u}||\vec{v}|}, &
    \text{dvou rovin:}\quad\cos\varphi &= \frac{|\vec{n}_1\cdot\vec{n}_2|}{|\vec{n}_1||\vec{n}_2|},\\
    \text{přímky a roviny:}\quad\sin\varphi &= \frac{|\vec{u}\cdot\vec{n}|}{|\vec{u}||\vec{n}|}.
  \end{align*}
\end{definitionbox}

Všechny tyto vzorce vypadají jako něco, co je potřeba se naučit nazpaměť. Ve skutečnosti za nimi stojí jediná myšlenka: \textbf{průmět vektoru do vhodného směru}. Pojďme si je projít s obrázky.

\subsection*{Vzdálenost bodu od roviny}

Představte si, že $M$ je lampa visící nad podlahou $\rho$. Nejkratší cesta z lampy na podlahu vede svisle dolů, tedy ve směru normály $\vec{n}$. Vezmeme libovolný bod $P$ podlahy. Vektor $\overrightarrow{PM}$ míří šikmo, ale nás zajímá jen jeho `svislá' složka, tedy jeho průmět do směru $\vec{n}$.

\begin{center}
  \begin{tikzpicture}[krychle, scale=1.15]
    \fill[colorofwisdom!15] (0,0,0) -- (4,0,0) -- (4,3,0) -- (0,3,0) -- cycle;
    \draw[gray] (0,0,0) -- (4,0,0) -- (4,3,0) -- (0,3,0) -- cycle;
    \node at (3.7,2.7,0) {\small $\rho$};
    \coordinate (M) at (2.6,1.6,2); \coordinate (M') at (2.6,1.6,0); \coordinate (P) at (0.8,0.5,0);
    \draw[dashed] (P) -- (M');
    \draw[->, thick] (P) -- (M) node[midway, above left] {\small $\overrightarrow{PM}$};
    \draw[very thick, colorofwisdom] (M) -- (M') node[midway, right] {\small $v$};
    \draw (2.39,1.47,0) -- (2.39,1.47,0.2) -- (2.6,1.6,0.2);
    \draw[->, thick, colorofwisdom] (P) -- (0.8,0.5,1.1) node[left] {\small $\vec{n}$};
    \draw[dotted] (0.8,0.5,1.1) -- (0.8,0.5,2) -- (M);
    \fill (M) circle (1.6pt) node[above right] {\small $M$};
    \fill (M') circle (1.6pt) node[below right] {\small $M'$};
    \fill (P) circle (1.6pt) node[below left] {\small $P$};
  \end{tikzpicture}
\end{center}

Délka průmětu je $|\overrightarrow{PM}|\cdot|\cos\alpha|$, kde $\alpha$ je úhel mezi $\overrightarrow{PM}$ a $\vec{n}$. Ze skalárního součinu $\overrightarrow{PM}\cdot\vec{n} = |\overrightarrow{PM}||\vec{n}|\cos\alpha$ dostaneme
\begin{equation*}
  v = \frac{|\overrightarrow{PM}\cdot\vec{n}|}{|\vec{n}|}.
\end{equation*}
A teď rozepišme čitatel. Pro $P = [p_1, p_2, p_3]$ je
\begin{equation*}
  \overrightarrow{PM}\cdot\vec{n} = a(m_1 - p_1) + b(m_2 - p_2) + c(m_3 - p_3) = am_1 + bm_2 + cm_3 - \underbrace{(ap_1 + bp_2 + cp_3)}_{=\,-d} = am_1 + bm_2 + cm_3 + d,
\end{equation*}
protože $P$ leží v rovině, a tedy $ap_1 + bp_2 + cp_3 + d = 0$. To je celý vzorec. Všimněte si, že nezáleží na tom, který bod $P$ vybereme. Posouváním $P$ po podlaze se mění jen `vodorovná' část vektoru $\overrightarrow{PM}$, svislá zůstává stejná. Mimochodem, výraz $am_1 + bm_2 + cm_3 + d$ bez absolutní hodnoty má ještě jednu informaci: jeho znaménko říká, na \textit{které straně} roviny bod leží.

\subsection*{Vzdálenost mimoběžek}

Každé dvě mimoběžky leží ve dvou rovnoběžných rovinách, jako přímka na podlaze a přímka na stropě. Jejich vzdálenost je tedy prostě `výška místnosti'. Obě roviny mají stejnou normálu, která musí být kolmá na $\vec{u}$ i na $\vec{v}$. Takový vektor už známe: $\vec{u}\times\vec{v}$. Vzdálenost je tak délka průmětu libovolného vektoru $\overrightarrow{AB}$ (z podlahy na strop) do směru $\vec{u}\times\vec{v}$, přesně jako výše.

\begin{center}
  \begin{tikzpicture}[krychle, scale=1.15]
    \fill[colorofwisdom!15] (-0.5,-0.3,0) -- (3.5,-0.3,0) -- (3.5,2.5,0) -- (-0.5,2.5,0) -- cycle;
    \draw[gray] (-0.5,-0.3,0) -- (3.5,-0.3,0) -- (3.5,2.5,0) -- (-0.5,2.5,0) -- cycle;
    \draw[thick] (-0.5,1,0) -- (3.5,1,0) node[right] {\small $p$ (směr $\vec{u}$)};
    \draw[very thick, colorofwisdom] (1.5,1,0) -- (1.5,1,1.6) node[midway, right] {\small $v$};
    \draw[->, thick, colorofwisdom] (0.3,1,0) -- (0.3,1,1) node[left] {\small $\vec{u}\times\vec{v}$};
    \fill[colorofwisdom!8, opacity=0.7] (-0.5,-0.3,1.6) -- (3.5,-0.3,1.6) -- (3.5,2.5,1.6) -- (-0.5,2.5,1.6) -- cycle;
    \draw[gray] (-0.5,-0.3,1.6) -- (3.5,-0.3,1.6) -- (3.5,2.5,1.6) -- (-0.5,2.5,1.6) -- cycle;
    \draw[thick] (1.5,-0.3,1.6) -- (1.5,2.5,1.6) node[above right] {\small $q$ (směr $\vec{v}$)};
    \draw[->, thick] (0.3,1,0) -- (1.5,2.1,1.6);
    \node[right] at (1.1,1.65,1.05) {\small $\overrightarrow{AB}$};
    \fill (0.3,1,0) circle (1.6pt) node[below] {\small $A$};
    \fill (1.5,2.1,1.6) circle (1.6pt) node[above left] {\small $B$};
  \end{tikzpicture}
\end{center}

Vzorec má i krásný geometrický význam. Číslo $\overrightarrow{AB}\cdot(\vec{u}\times\vec{v})$ se jmenuje \textbf{smíšený součin} a jeho absolutní hodnota je \textbf{objem rovnoběžnostěnu} s hranami $\overrightarrow{AB}$, $\vec{u}$, $\vec{v}$. Proč? Objem je podstava krát výška, podstava je rovnoběžník s obsahem $|\vec{u}\times\vec{v}|$ a výška je průmět $\overrightarrow{AB}$ do kolmého směru. Vzorec pro vzdálenost mimoběžek tedy říká
\begin{equation*}
  v = \frac{\text{objem}}{\text{obsah podstavy}} = \text{výška}.
\end{equation*}

\subsection*{Odchylky}

Odchylka dvou rovin se měří jako úhel otevřené knihy: v rovině kolmé na hřbet. Normála každé stránky je od stránky otočená o $90^\circ$, takže když se stránky rozevřou o úhel $\varphi$, normály se rozevřou o stejný úhel. Proto stačí změřit úhel normál. Absolutní hodnota ve vzorci zajistí, že dostaneme úhel mezi $0^\circ$ a $90^\circ$. Normálu totiž můžeme vzít i opačně a odchylka nemá záviset na tom, jak se rozhodneme.

\begin{center}
  \begin{tikzpicture}[scale=1.3]
    \begin{scope}[krychle]
      \fill[colorofwisdom!15] (0,0,0) -- (3,0,0) -- (3,2,0) -- (0,2,0) -- cycle;
      \draw[gray] (0,0,0) -- (3,0,0) -- (3,2,0) -- (0,2,0) -- cycle;
      \coordinate (P) at (1,1,0);
      \coordinate (Q) at ({1+1.9*cos(35)},1,{1.9*sin(35)});
      \coordinate (Q') at ({1+1.9*cos(35)},1,0);
      \coordinate (N) at (1,1,1.4);
      \coordinate (R) at ({1-0.8*cos(35)},1,{-0.8*sin(35)});
    \end{scope}
    \draw[thick, dashed] (R) -- (P);
    \draw[thick] (P) -- (Q) node[right] {\small $p$};
    \draw[dashed] (P) -- (Q');
    \draw[dotted] (Q) -- (Q');
    \draw[->, thick, colorofwisdom] (P) -- (N) node[above] {\small $\vec{n}$};
    \draw ($(P)+(0:0.75)$) arc (0:35:0.75);
    \node at ($(P)+(17:1)$) {\small $\varphi$};
    \draw[colorofwisdom] ($(P)+(35:0.45)$) arc (35:90:0.45);
    \node at ($(P)+(64:0.85)$) {\scriptsize $90^\circ{-}\varphi$};
    \node[font=\small] at (1.8,-0.6) {přímka a rovina};
    \begin{scope}[xshift=5.2cm]
      \begin{scope}[krychle]
        \fill[colorofwisdom!15] (0,0,0) -- (2,0,0) -- (2,2.2,0) -- (0,2.2,0) -- cycle;
        \draw[gray] (0,0,0) -- (2,0,0) -- (2,2.2,0) -- (0,2.2,0) -- cycle;
        \fill[colorofwisdom!30, opacity=0.7] (0,0,0) -- ({2*cos(65)},0,{2*sin(65)}) -- ({2*cos(65)},2.2,{2*sin(65)}) -- (0,2.2,0) -- cycle;
        \draw[gray] (0,0,0) -- ({2*cos(65)},0,{2*sin(65)}) -- ({2*cos(65)},2.2,{2*sin(65)}) -- (0,2.2,0);
        \draw[very thick] (0,-0.2,0) -- (0,2.4,0);
        \coordinate (O) at (0,1.1,0);
        \coordinate (K1) at (1.7,1.1,0);
        \coordinate (K2) at ({1.7*cos(65)},1.1,{1.7*sin(65)});
        \coordinate (N1) at (0,1.1,1.2);
        \coordinate (N2) at ({-1.2*sin(65)},1.1,{1.2*cos(65)});
      \end{scope}
      \draw[thick] (O) -- (K1) (O) -- (K2);
      \draw ($(O)+(0:0.8)$) arc (0:65:0.8);
      \node at ($(O)+(32:1.05)$) {\small $\varphi$};
      \draw[->, thick, colorofwisdom] (O) -- (N1) node[above] {\small $\vec{n}_1$};
      \draw[->, thick, colorofwisdom] (O) -- (N2) node[above left] {\small $\vec{n}_2$};
      \draw[colorofwisdom] ($(O)+(90:0.55)$) arc (90:155:0.55);
      \node[colorofwisdom] at ($(O)+(122:0.8)$) {\small $\varphi$};
      \node[font=\small] at (1.2,-0.6) {dvě roviny};
    \end{scope}
  \end{tikzpicture}
\end{center}

Odchylka dvou přímek je jednodušší: je to úhel jejich směrových vektorů (opět s absolutní hodnotou). U mimoběžek si jednu přímku posuneme tak, aby druhou protínala, a měříme úhel tam.

U odchylky přímky a roviny je sinus, ne kosinus. Proč? Podívejte se na levý obrázek. Úhel přímky a roviny měříme mezi přímkou a jejím `stínem' (kolmým průmětem) v rovině. Vzorec s kosinem by nám ale dal úhel mezi přímkou a \textit{normálou} roviny. Úhel mezi přímkou a rovinou je doplněk do $90^\circ$ a $\cos(90^\circ - \varphi) = \sin\varphi$.

\begin{example}
  Vraťme se ke krychli z kapitoly~\ref{kap:stereometrie} a spočítejme všechno znovu, tentokrát bez hledání trojúhelníků. Krychle s hranou $1$ má vrcholy $A = [0, 0, 0]$, $B = [1, 0, 0]$, $C = [1, 1, 0]$, $D = [0, 1, 0]$, $E = [0, 0, 1]$, $F = [1, 0, 1]$, $G = [1, 1, 1]$, $H = [0, 1, 1]$.
  \begin{itemize}
    \item \textbf{Odchylka tělesové úhlopříčky od podstavy}: $\vec{u} = \overrightarrow{AG} = (1, 1, 1)$, normála podstavy $\vec{n} = (0, 0, 1)$. $\sin\varphi = \frac{1}{\sqrt{3}}$, $\varphi\approx 35{,}26^\circ$. Stejně jako minule.
    \item \textbf{Vzdálenost $A$ od roviny $BDE$}: rovina prochází body $[1, 0, 0]$, $[0, 1, 0]$, $[0, 0, 1]$, má tedy rovnici $x + y + z - 1 = 0$. Vzdálenost $\frac{|-1|}{\sqrt{3}} = \frac{\sqrt{3}}{3}$, tedy přesně třetina tělesové úhlopříčky.
    \item \textbf{Vzdálenost mimoběžek $AB$ a $CG$}: $\vec{u} = (1, 0, 0)$, $\vec{v} = (0, 0, 1)$, $\vec{u}\times\vec{v} = (0, -1, 0)$, $\overrightarrow{AC} = (1, 1, 0)$. Vzdálenost $\frac{|-1|}{1} = 1$.
    \item \textbf{Odchylka rovin $ABG$ a $ABC$}: rovina $ABG$ obsahuje vektory $(1, 0, 0)$ a $(1, 1, 1)$, normála $(0, -1, 1)$. Rovina $ABC$ má normálu $(0, 0, 1)$. $\cos\varphi = \frac{1}{\sqrt{2}}$, $\varphi = 45^\circ$.
  \end{itemize}
  Všechno sedí s výsledky z kapitoly o stereometrii. Kdo nemá prostorovou představivost, má souřadnice.
\end{example}

\begin{center}
  \begin{tikzpicture}[krychle, scale=2.7]
    \krychle
    \fill[colorofwisdom!25, opacity=0.7] (B) -- (D) -- (E) -- cycle;
    \draw[colorofwisdom] (B) -- (E) -- (D);
    \draw[colorofwisdom, dashed] (B) -- (D);
    \coordinate (T) at (1/3,1/3,1/3);
    \draw[very thick, colorofwisdom] (A) -- (T);
    \draw[dashed, colorofwisdom] (T) -- (G);
    \fill (T) circle (0.7pt) node[right, font=\scriptsize] {$T$};
    \node[font=\scriptsize, align=left, anchor=west] at (1.6,0,0.5) {rovina $BDE$ protíná\\úhlopříčku $AG$ v bodě $T$,\\$|AT| = \frac{1}{3}|AG| = \frac{\sqrt{3}}{3}$};
  \end{tikzpicture}
\end{center}
Na obrázku je vidět i to, proč vzdálenost $A$ od roviny $BDE$ měříme právě po úhlopříčce $AG$: ta má směr $(1, 1, 1)$, což je přesně normálový vektor roviny $x + y + z - 1 = 0$. Úhlopříčka tedy rovinu protíná kolmo.

\begin{problem}
  \begin{enumerate}[label=(\alph*)]
    \item Napište obecnou rovnici roviny, která prochází bodem $[2, -1, 3]$ a je kolmá na vektor $(1, 4, -2)$.
    \item Určete vzájemnou polohu přímky $X = [1, 0, 2] + t(2, 1, -1)$ a roviny $x - y + z - 4 = 0$. Pokud jsou rovnoběžné, určete jejich vzdálenost.
    \item Určete vzdálenost rovnoběžných rovin $2x - y + 2z + 3 = 0$ a $2x - y + 2z - 6 = 0$.
    \item Určete odchylku rovin $x + y + z = 1$ a $z = 0$.
    \item Určete obsah trojúhelníku $A = [1, 1, 0]$, $B = [3, 1, 0]$, $C = [1, 4, 4]$.
  \end{enumerate}
\end{problem}

\begin{solution}
  \begin{enumerate}[label=(\alph*)]
    \item $x + 4y - 2z + d = 0$, $2 - 4 - 6 + d = 0$, tedy $x + 4y - 2z + 8 = 0$.
    \item $\vec{n}\cdot\vec{u} = 2 - 1 - 1 = 0$, takže je přímka s rovinou rovnoběžná nebo v ní leží. Bod $[1, 0, 2]$: $1 - 0 + 2 - 4 = -1\neq 0$, rovnoběžná. Vzdálenost $\frac{1}{\sqrt{3}} = \frac{\sqrt{3}}{3}$.
    \item Bod první roviny, například $[0, 3, 0]$, má od druhé vzdálenost $\frac{|-3 - 6|}{\sqrt{4 + 1 + 4}} = \frac{9}{3} = 3$.
    \item $\cos\varphi = \frac{1}{\sqrt{3}}$, $\varphi\approx 54{,}7^\circ$.
    \item $\overrightarrow{AB} = (2, 0, 0)$, $\overrightarrow{AC} = (0, 3, 4)$, $\overrightarrow{AB}\times\overrightarrow{AC} = (0\cdot 4 - 0\cdot 3,\ 0\cdot 0 - 2\cdot 4,\ 2\cdot 3 - 0\cdot 0) = (0, -8, 6)$, $S = \frac{1}{2}\cdot 10 = 5$.
  \end{enumerate}
\end{solution}

\section{Příběh --- Most v Dublinu}

Irský matematik William Rowan Hamilton byl zázračné dítě. Ve třinácti letech prý mluvil třinácti jazyky a ve dvaadvaceti, ještě jako student, byl jmenován královským astronomem Irska. Přes deset let se pak snažil vyřešit problém, který ho pronásledoval: jak `násobit' trojice čísel, tedy body v prostoru, podobně jako komplexní čísla (kapitola~\ref{kap:komplexni}) umožňují násobit body v rovině. Každé ráno se ho prý děti u snídaně ptaly: `Tak co, tati, umíš už násobit trojice?' A on smutně odpovídal, že je umí jen sčítat a odčítat.

Šestnáctého října 1843 šel s manželkou podél kanálu v Dublinu a najednou mu došlo, že trojice nestačí. Potřeboval čtyři čísla. Byl tak nadšený, že nožem vyryl do kamene Broughamova mostu vzorec
\begin{equation*}
  i^2 = j^2 = k^2 = ijk = -1,
\end{equation*}
který definuje takzvané \textbf{kvaterniony}. Nápis už dávno zvětral, ale na mostě je dnes pamětní deska a matematici se k ní každý rok 16. října vydávají na pochod. Kvaterniony se dnes používají k počítání rotací v počítačových hrách, v robotice i v řízení vesmírných sond. Moderní vektorový součin, který jsme poznali v této kapitole, vznikl vlastně jako zjednodušený kus Hamiltonova násobení kvaternionů.

Hamiltonův osobní život tak šťastný nebyl. Celý život miloval ženu jménem Catherine Disneyová, kterou rodina provdala za jiného. V pozdějších letech propadl alkoholu a stal se samotářem. Když po jeho smrti v roce 1865 procházeli jeho pracovnu, našli prý mezi hromadami rukopisů i talíře se zbytky jídla, které tam ležely měsíce, protože odmítal, aby mu kdokoli sahal na papíry.

\section{Shrnutí a otázky k ústní zkoušce}

\begin{problem}
  Odpovězte vlastními slovy:
  \begin{enumerate}
    \item Jak se počítá vzdálenost bodů, velikost vektoru a skalární součin v prostoru?
    \item Co je vektorový součin? Jaké má vlastnosti a k čemu se používá?
    \item Jak vypadá parametrické vyjádření přímky a roviny v prostoru? Proč přímka v prostoru nemá obecnou rovnici?
    \item Jak najdete obecnou rovnici roviny procházející třemi body?
    \item Jak určíte vzájemnou polohu dvou přímek v prostoru? Jak poznáte mimoběžky?
    \item Jak určíte vzájemnou polohu přímky a roviny a dvou rovin?
    \item Jak se počítá vzdálenost bodu od roviny a vzdálenost mimoběžek?
    \item Jak se počítají odchylky přímek a rovin? Proč je u odchylky přímky a roviny sinus?
  \end{enumerate}
\end{problem}

\koneckapitoly

\chapter{Posloupnosti a řady}
\label{kap:posloupnosti}
\mathquote{The infinite! No other question has ever moved so profoundly the spirit of man.}{David Hilbert}

Posloupnost je seznam čísel, který (obvykle) nikdy nekončí. Na první pohled nic zajímavého. Jenže právě na posloupnostech se matematici poprvé naučili pořádně zacházet s nekonečnem. Naučíme se tu sčítat nekonečně mnoho čísel a dostat konečný výsledek a ukážeme si, proč $0{,}999\ldots$ je přesně jedna. Všechno, co přijde v kapitolách o limitách, derivacích a integrálech, stojí na tom, co se naučíme tady.

\section{Pojem posloupnosti}

\begin{definitionbox}
  \textbf{Posloupnost} je funkce, jejímž definičním oborem je množina přirozených čísel $\mathbb{N}$ (nebo její počáteční úsek $\{1, 2, \ldots, k\}$, pak mluvíme o \textbf{konečné} posloupnosti). Hodnotu v čísle $n$ značíme $a_n$ a nazýváme ji \textbf{$n$-tý člen}. Posloupnost zapisujeme ${(a_n)}_{n=1}^{\infty}$ nebo zkráceně $(a_n)$.

  \vspace{0.2cm}

  \textbf{Způsoby zadání}:
  \begin{itemize}
    \item \textbf{vzorcem pro $n$-tý člen}, například $a_n = \frac{1}{n}$,
    \item \textbf{rekurentně}: zadáme první člen (nebo několik prvních) a pravidlo, jak z předchozích členů spočítat další. Například $a_1 = 1$, $a_{n+1} = 2a_n + 1$,
    \item \textbf{výčtem} prvních členů (jen pokud je zřejmé, jak pokračovat, což u maturity bývá a v životě nebývá).
  \end{itemize}
\end{definitionbox}

Graf posloupnosti není křivka, ale izolované body $[1, a_1], [2, a_2], [3, a_3], \ldots$, protože definičním oborem jsou jen přirozená čísla.

\begin{center}
  \begin{tikzpicture}[scale=0.85]
    \begin{scope}
      \draw[->] (0,0) -- (5.4,0) node[right] {\small $n$};
      \draw[->] (0,0) -- (0,3.5) node[above] {\small $a_n$};
      \draw[dashed, gray] (0,3) -- (5.2,3);
      \node[left] at (0,3) {\small $1$};
      \foreach \n in {1,...,8} {
        \draw (0.6*\n,0.07) -- (0.6*\n,-0.07) node[below] {\tiny $\n$};
        \filldraw[colorofwisdom] (0.6*\n,{3*\n/(\n+1)}) circle (2pt);
      }
      \node at (2.7,-0.9) {\small $a_n = \frac{n}{n + 1}$: rostoucí, omezená};
    \end{scope}
    \begin{scope}[xshift=7cm, yshift=1.6cm]
      \draw[->] (0,0) -- (5.4,0) node[right] {\small $n$};
      \draw[->] (0,-1.6) -- (0,1.9) node[above] {\small $b_n$};
      \draw[dashed, gray] (0,1.2) -- (5.2,1.2);
      \draw[dashed, gray] (0,-1.2) -- (5.2,-1.2);
      \node[left] at (0,1.2) {\small $1$};
      \node[left] at (0,-1.2) {\small $-1$};
      \foreach \n in {1,...,8} {
        \draw (0.6*\n,0.07) -- (0.6*\n,-0.07);
        \filldraw[colorofwisdom] (0.6*\n,{1.2*ifthenelse(isodd(\n),-1,1)}) circle (2pt);
      }
      \node at (2.7,-2.5) {\small $b_n = {(-1)}^n$: omezená, není monotónní};
    \end{scope}
  \end{tikzpicture}
\end{center}

Body grafu nespojujeme čarou. Mezi $n = 1$ a $n = 2$ posloupnost žádnou hodnotu nemá, protože `první a půltý člen' neexistuje. Na obrázku jsou dvě posloupnosti z příkladu níže. Představte si posloupnost jako sérii snímků: na grafu hned vidíte, jestli body stoupají (rostoucí posloupnost), jestli se vejdou do vodorovného pruhu (omezená posloupnost), nebo jestli skáčou nahoru a dolů.

\begin{example}
  Slavná rekurentně zadaná posloupnost je \textbf{Fibonacciho posloupnost}: $F_1 = F_2 = 1$ a $F_{n+2} = F_{n+1} + F_n$. Prvních pár členů je
  \begin{equation*}
    1, 1, 2, 3, 5, 8, 13, 21, 34, 55, 89, \ldots
  \end{equation*}
  Leonardo Pisánský, zvaný Fibonacci, ji v roce 1202 popsal jako model rozmnožování králíků. Poměry sousedních členů $\frac{F_{n+1}}{F_n}$ se přitom blíží k číslu $\frac{1 + \sqrt{5}}{2}$, tedy ke zlatému řezu z kapitol~\ref{kap:vyrazy} a~\ref{kap:kvadraticke}. Zkuste si to: $\frac{89}{55}\approx 1{,}618$.
\end{example}

\begin{definitionbox}
  Posloupnost $(a_n)$ je
  \begin{itemize}
    \item \textbf{rostoucí}, pokud $a_{n+1} > a_n$ pro všechna $n$, \textbf{klesající}, pokud $a_{n+1} < a_n$ (a obdobně neklesající a nerostoucí),
    \item \textbf{shora omezená}, pokud existuje $h$ s $a_n\leq h$ pro všechna $n$, \textbf{zdola omezená}, pokud existuje $d$ s $a_n\geq d$, a \textbf{omezená}, pokud je obojí.
  \end{itemize}
  Je to totéž jako u funkcí v kapitole~\ref{kap:funkce}, jen stačí porovnávat sousední členy.
\end{definitionbox}

\begin{example}
  Posloupnost $a_n = \frac{n}{n + 1}$ je rostoucí, protože
  \begin{equation*}
    a_{n+1} - a_n = \frac{n + 1}{n + 2} - \frac{n}{n + 1} = \frac{{(n + 1)}^2 - n(n + 2)}{(n + 1)(n + 2)} = \frac{1}{(n + 1)(n + 2)} > 0.
  \end{equation*}
  Je omezená: $0 < a_n < 1$. Posloupnost $b_n = {(-1)}^n$ ($-1, 1, -1, 1, \ldots$) je omezená, ale není monotónní.
\end{example}

\section{Aritmetická posloupnost}

\begin{definitionbox}
  \textbf{Aritmetická posloupnost} je posloupnost, ve které je rozdíl každých dvou sousedních členů stejný. Tento rozdíl $d$ se nazývá \textbf{diference}:
  \begin{equation*}
    a_{n+1} = a_n + d.
  \end{equation*}
  Platí:
  \begin{align*}
    a_n &= a_1 + (n - 1)d,\\
    a_r - a_s &= (r - s)d,\\
    s_n &= a_1 + a_2 + \cdots + a_n = \frac{n(a_1 + a_n)}{2}.
  \end{align*}
\end{definitionbox}

\begin{proof}[(vzorce pro součet)]
  Napíšeme součet dvakrát, jednou popředu a jednou pozpátku, a sečteme:
  \begin{align*}
    s_n &= a_1 + a_2 + \cdots + a_{n-1} + a_n,\\
    s_n &= a_n + a_{n-1} + \cdots + a_2 + a_1,\\
    2s_n &= (a_1 + a_n) + (a_2 + a_{n-1}) + \cdots + (a_n + a_1).
  \end{align*}
  Každá závorka má stejnou hodnotu $a_1 + a_n$, protože když jeden člen o $d$ vzroste, druhý o $d$ klesne. Závorek je $n$, takže $2s_n = n(a_1 + a_n)$.
\end{proof}

Tohle je přesně ten trik, se kterým podle legendy přišel malý Gauss, když měl ve škole za trest sečíst čísla od $1$ do $100$. Za pár vteřin odevzdal výsledek $5050 = \frac{100\cdot 101}{2}$. Tutéž větu jsme dokázali matematickou indukcí v kapitole 1. Teď máte dva důkazy, takže si můžete vybrat ten, který se vám víc líbí.

Gaussův trik se dá i nakreslit. Součet $1 + 2 + \cdots + n$ si představte jako schodiště ze čtverečků. Vezmeme druhé, stejné schodiště, otočíme ho vzhůru nohama a přiložíme k prvnímu. Obě schodiště dohromady přesně vyplní obdélník $n\times(n + 1)$, takže jedno schodiště má $\frac{n(n + 1)}{2}$ čtverečků. Sčítání `popředu a pozpátku' v důkazu je přesně tohle přiložení.

\begin{center}
  \begin{tikzpicture}[scale=0.5]
    \foreach \k in {1,...,5} {
      \foreach \j in {1,...,\k} {
        \filldraw[fill=colorofwisdom!60, draw=white, thick] (\k-1,\j-1) rectangle (\k,\j);
      }
      \foreach \j in {\k,...,5} {
        \filldraw[fill=gray!30, draw=white, thick] (\k-1,\j) rectangle (\k,\j+1);
      }
    }
    \draw[thick] (0,0) rectangle (5,6);
    \draw[decorate, decoration={brace, mirror, amplitude=4pt}] (0,-0.2) -- (5,-0.2) node[midway, below=5pt] {\small $n = 5$};
    \draw[decorate, decoration={brace, amplitude=4pt}] (5.2,6) -- (5.2,0) node[midway, right=5pt] {\small $n + 1 = 6$};
    \node[anchor=west, align=left, font=\small] at (9.6,3) {červené schodiště: $1 + 2 + 3 + 4 + 5$\\šedé schodiště: totéž pozpátku\\dohromady: $5\cdot 6 = 30$ čtverečků\\jedno schodiště: $\frac{30}{2} = 15$};
  \end{tikzpicture}
\end{center}

Název `aritmetická' pochází z toho, že každý člen je \textbf{aritmetickým průměrem} svých sousedů: $a_n = \frac{a_{n-1} + a_{n+1}}{2}$.

Graficky: body aritmetické posloupnosti leží na \textbf{přímce}. Z každého bodu do dalšího uděláme krok o $1$ doprava a o $d$ nahoru (pro $d < 0$ dolů). Aritmetická posloupnost je tedy lineární funkce $y = a_1 + (x - 1)d$ vyhodnocená jen v přirozených číslech a diference $d$ je její směrnice. Odtud je vidět i vlastnost s průměrem: prostřední ze tří bodů na přímce leží přesně v polovině výšky mezi sousedy.

\begin{center}
  \begin{tikzpicture}[scale=0.85]
    \draw[->] (0,0) -- (7.4,0) node[right] {\small $n$};
    \draw[->] (0,0) -- (0,4.3) node[above] {\small $a_n$};
    \draw[dashed, gray, domain=0.3:7] plot (\x, {0.5 + 0.5*\x});
    \foreach \n in {1,...,6} {
      \draw (\n,0.07) -- (\n,-0.07) node[below] {\tiny $\n$};
    }
    \foreach \n in {1,...,5} {
      \draw[colorofwisdom!70] (\n,{0.5+0.5*\n}) -- (\n+1,{0.5+0.5*\n}) -- (\n+1,{1+0.5*\n});
    }
    \node[right, colorofwisdom] at (6,3.25) {\small $d$};
    \node[below, colorofwisdom] at (5.5,3) {\tiny $1$};
    \foreach \n in {1,...,6} {
      \filldraw[colorofwisdom] (\n,{0.5+0.5*\n}) circle (2.2pt);
    }
    \node[left] at (1,1.25) {\small $a_1$};
    \node[anchor=west, font=\small] at (7.6,2) {$a_1 = 1$, $d = \frac{1}{2}$};
  \end{tikzpicture}
\end{center}

\begin{example}
  \begin{itemize}
    \item Pro $a_1 = 3$, $d = 4$ je $a_{10} = 3 + 9\cdot 4 = 39$ a $s_{10} = \frac{10(3 + 39)}{2} = 210$.
    \item Je dáno $a_3 = 7$ a $a_8 = 22$. Pak $a_8 - a_3 = 5d = 15$, tedy $d = 3$ a $a_1 = a_3 - 2d = 1$. Dále $a_{20} = 1 + 19\cdot 3 = 58$ a $s_{20} = \frac{20(1 + 58)}{2} = 590$.
    \item Součet prvních $n$ lichých čísel: $1 + 3 + 5 + \cdots + (2n - 1) = \frac{n(1 + 2n - 1)}{2} = n^2$. To se dá i nakreslit: čtverec $n\times n$ se skládá z `háčků' o $1, 3, 5, \ldots$ čtverečcích.
  \end{itemize}

  \begin{center}
    \begin{tikzpicture}[scale=0.55]
      \foreach \i in {0,...,4} {
        \foreach \j in {0,...,4} {
          \pgfmathtruncatemacro{\m}{max(\i,\j)}
          \ifodd\m
            \filldraw[fill=gray!30, draw=white, thick] (\i,\j) rectangle (\i+1,\j+1);
          \else
            \filldraw[fill=colorofwisdom!55, draw=white, thick] (\i,\j) rectangle (\i+1,\j+1);
          \fi
        }
      }
      \foreach \m in {0,...,4} {
        \pgfmathtruncatemacro{\h}{2*\m+1}
        \node[font=\small] at (\m+0.5,5.5) {$\h$};
      }
      \draw[thick] (0,0) rectangle (5,5);
      \node[anchor=west, font=\small] at (5.8,2.5) {$1 + 3 + 5 + 7 + 9 = 5^2$};
    \end{tikzpicture}
  \end{center}
  Každý další háček obepne dosavadní čtverec zprava a shora, a tím z něj udělá čtverec o jedna větší. Ze čtverce $k\times k$ na $(k + 1)\times(k + 1)$ přibude právě $2k + 1$ čtverečků, tedy další liché číslo.
\end{example}

\begin{problem}
  \begin{enumerate}[label=(\alph*)]
    \item V aritmetické posloupnosti je $a_5 = 17$ a $a_{12} = 45$. Určete $a_1$, $d$ a $s_{12}$.
    \item Kolik členů posloupnosti $5, 9, 13, \ldots$ musíme sečíst, abychom dostali $230$?
    \item Určete součet všech dvojciferných čísel dělitelných třemi.
    \item V hledišti je $20$ řad, v první řadě je $15$ sedadel a v každé další o $2$ víc. Kolik sedadel má hlediště?
  \end{enumerate}
\end{problem}

\begin{solution}
  \begin{enumerate}[label=(\alph*)]
    \item $7d = 28$, $d = 4$, $a_1 = 17 - 16 = 1$, $s_{12} = \frac{12(1 + 45)}{2} = 276$.
    \item $s_n = \frac{n(10 + 4(n - 1))}{2} = n(2n + 3) = 230$, tedy $2n^2 + 3n - 230 = 0$, $D = 9 + 1840 = 43^2$, $n = \frac{-3 + 43}{4} = 10$.
    \item $12, 15, \ldots, 99$: $a_1 = 12$, $d = 3$, $99 = 12 + (n - 1)3$, $n = 30$. $s_{30} = \frac{30(12 + 99)}{2} = 1665$.
    \item $a_{20} = 15 + 38 = 53$, $s_{20} = \frac{20(15 + 53)}{2} = 680$.
  \end{enumerate}
\end{solution}

\section{Geometrická posloupnost}

\begin{definitionbox}
  \textbf{Geometrická posloupnost} je posloupnost, ve které je podíl každých dvou sousedních členů stejný. Tento podíl $q$ se nazývá \textbf{kvocient}:
  \begin{equation*}
    a_{n+1} = a_n\cdot q,\qquad a_1\neq 0,\ q\neq 0.
  \end{equation*}
  Platí:
  \begin{align*}
    a_n &= a_1\cdot q^{n-1},\\
    a_r &= a_s\cdot q^{r - s},\\
    s_n &= a_1\cdot\frac{q^n - 1}{q - 1}\quad\text{pro } q\neq 1,\qquad s_n = na_1\quad\text{pro } q = 1.
  \end{align*}
\end{definitionbox}

\begin{proof}[(vzorce pro součet)]
  Napíšeme $s_n$ a $q\cdot s_n$ pod sebe:
  \begin{align*}
    s_n &= a_1 + a_1q + a_1q^2 + \cdots + a_1q^{n-1},\\
    qs_n &= \phantom{a_1 + {}} a_1q + a_1q^2 + \cdots + a_1q^{n-1} + a_1q^n.
  \end{align*}
  Odečtením se skoro všechno vyruší: $qs_n - s_n = a_1q^n - a_1$, tedy $s_n(q - 1) = a_1(q^n - 1)$. Pro $q\neq 1$ můžeme dělit.
\end{proof}

Název `geometrická' pochází z toho, že (pro kladné členy) je každý člen \textbf{geometrickým průměrem} svých sousedů: $a_n = \sqrt{a_{n-1}a_{n+1}}$. Geometrické posloupnosti jsou diskrétní verzí exponenciálních funkcí z kapitoly~\ref{kap:exponencialni}: $a_n = a_1q^{n-1}$ je exponenciála $q^x$ vyhodnocená jen v celých číslech.

Zatímco body aritmetické posloupnosti leží na přímce, body geometrické posloupnosti leží na \textbf{exponenciále}. Pro $q > 1$ letí stále rychleji nahoru, protože každý krok je o stejný \textit{násobek} větší než ten předchozí. Pro $0 < q < 1$ se plíží k nule, každý krok je kratší. A pro záporné $q$ se znaménko střídá, takže body skáčou z jedné strany osy na druhou.

\begin{center}
  \begin{tikzpicture}[scale=0.75]
    \begin{scope}
      \draw[->] (0,0) -- (4.3,0) node[right] {\small $n$};
      \draw[->] (0,0) -- (0,3.6);
      \draw[dashed, gray, domain=0.3:3.7, samples=40] plot (\x, {0.1*pow(2,\x/0.6-1)});
      \foreach \n in {1,...,6} {
        \draw (0.6*\n,0.07) -- (0.6*\n,-0.07);
        \filldraw[colorofwisdom] (0.6*\n,{0.1*pow(2,\n-1)}) circle (2pt);
      }
      \node at (2,-0.7) {\small $q = 2$};
    \end{scope}
    \begin{scope}[xshift=5.5cm]
      \draw[->] (0,0) -- (4.3,0) node[right] {\small $n$};
      \draw[->] (0,0) -- (0,3.6);
      \draw[dashed, gray, domain=0.4:4, samples=40] plot (\x, {3*pow(0.5,\x/0.6-1)});
      \foreach \n in {1,...,6} {
        \draw (0.6*\n,0.07) -- (0.6*\n,-0.07);
        \filldraw[colorofwisdom] (0.6*\n,{3*pow(0.5,\n-1)}) circle (2pt);
      }
      \node at (2,-0.7) {\small $q = \frac{1}{2}$};
    \end{scope}
    \begin{scope}[xshift=11cm, yshift=1.2cm]
      \draw[->] (0,0) -- (4.3,0) node[right] {\small $n$};
      \draw[->] (0,-1.4) -- (0,2.4);
      \foreach \n in {1,...,6} {
        \draw (0.6*\n,0.07) -- (0.6*\n,-0.07);
        \draw[gray] (0.6*\n,0) -- (0.6*\n,{2*ifthenelse(isodd(\n),1,-1)*pow(0.5,\n-1)});
        \filldraw[colorofwisdom] (0.6*\n,{2*ifthenelse(isodd(\n),1,-1)*pow(0.5,\n-1)}) circle (2pt);
      }
      \node at (2,-1.9) {\small $q = -\frac{1}{2}$};
    \end{scope}
  \end{tikzpicture}
\end{center}

\begin{example}
  \begin{itemize}
    \item Pro $a_1 = 2$, $q = 3$ je $a_6 = 2\cdot 3^5 = 486$ a $s_6 = 2\cdot\frac{3^6 - 1}{3 - 1} = 728$.
    \item Je dáno $a_2 = 6$ a $a_5 = 162$. Pak $q^3 = 27$, tedy $q = 3$ a $a_1 = 2$.
    \item Podle legendy si vynálezce šachu vyžádal od krále odměnu: jedno zrnko rýže na první políčko šachovnice, dvě na druhé, čtyři na třetí a tak dál. Celkem $s_{64} = \frac{2^{64} - 1}{2 - 1} = 2^{64} - 1\approx 1{,}8\cdot 10^{19}$ zrnek. To je přes čtyři sta miliard tun, tedy zhruba tolik rýže, kolik se jí na celém světě vypěstuje za několik set let. Legenda pak má dvě verze konce: v jedné se vynálezce stal královým rádcem, ve druhé mu král za drzost nechal useknout hlavu.
  \end{itemize}
\end{example}

\begin{problem}
  \begin{enumerate}[label=(\alph*)]
    \item V geometrické posloupnosti je $a_1 = 3$ a $a_4 = -24$. Určete $q$ a $s_6$.
    \item Mezi čísla $2$ a $162$ vložte tři čísla tak, aby spolu s nimi tvořila geometrickou posloupnost.
    \item Město má $50\,000$ obyvatel a počet obyvatel každý rok klesne o $2\,\%$. Kolik obyvatel bude mít za $10$ let?
  \end{enumerate}
\end{problem}

\begin{solution}
  \begin{enumerate}[label=(\alph*)]
    \item $q^3 = -8$, $q = -2$. $s_6 = 3\cdot\frac{{(-2)}^6 - 1}{-2 - 1} = 3\cdot\frac{63}{-3} = -63$.
    \item $162 = 2q^4$, $q^4 = 81$, $q = \pm 3$. Buď $6, 18, 54$, nebo $-6, 18, -54$. Obě řešení jsou správně.
    \item $50000\cdot 0{,}98^{10}\approx 40\,854$ obyvatel.
  \end{enumerate}
\end{solution}

\section{Limita posloupnosti}

Co se děje s posloupností $a_n = \frac{1}{n}$, když $n$ roste? Členy $1, \frac{1}{2}, \frac{1}{3}, \ldots, \frac{1}{1000}, \ldots$ se blíží k nule. Nikdy nule nejsou rovny, ale dostanou se k ní libovolně blízko. Tuhle představu `blížení se' potřebujeme zachytit přesně.

\begin{definitionbox}
  Posloupnost $(a_n)$ má \textbf{limitu} $a\in\mathbb{R}$, pokud ke každému $\varepsilon > 0$ existuje $n_0\in\mathbb{N}$ takové, že pro všechna $n\geq n_0$ platí $|a_n - a| < \varepsilon$. Píšeme
  \begin{equation*}
    \lim_{n\to\infty} a_n = a.
  \end{equation*}
  Posloupnost s (konečnou) limitou se nazývá \textbf{konvergentní}, jinak \textbf{divergentní}. Pokud členy rostou nade všechny meze (ke každému $K$ existuje $n_0$ s $a_n > K$ pro $n\geq n_0$), píšeme $\lim a_n = \infty$.
\end{definitionbox}

Definice vypadá děsivě, ale říká jednoduchou věc. Řekněte mi libovolně malou toleranci $\varepsilon$. Já vám řeknu, od kterého členu dál jsou všechny členy od $a$ blíž než $\varepsilon$. Pokud to dokážu pro každou toleranci, je $a$ limita. Je to vlastně hra dvou hráčů s kvantifikátory $\forall\varepsilon\ \exists n_0\ \forall n$ z kapitoly 1.

Nakreslete si kolem přímky $y = a$ vodorovný pás o šířce $\varepsilon$ nahoru i dolů. Posloupnost má limitu $a$, pokud po nějakém čase do pásu \textit{vletí a už z něj nikdy nevyletí}. Na prvních pár členech nezáleží, ty mohou být kdekoli. Když pás zúžíme, bude to jen trvat déle (větší $n_0$), ale nakonec se do něj posloupnost vejde zase.

\begin{center}
  \begin{tikzpicture}[scale=0.9]
    \fill[colorofwisdom!12] (0,1.3) rectangle (9.2,2.7);
    \draw[dashed, colorofwisdom] (0,1.3) -- (9.2,1.3) node[right] {\small $a - \varepsilon$};
    \draw[dashed, colorofwisdom] (0,2.7) -- (9.2,2.7) node[right] {\small $a + \varepsilon$};
    \draw[colorofwisdom] (0,2) -- (9.2,2) node[right] {\small $a$};
    \draw[->] (0,-1.3) -- (0,4) node[above] {\small $a_n$};
    \draw[->] (0,0) -- (9.5,0) node[right] {\small $n$};
    \draw[dashed, gray] (3,-1.2) -- (3,3.8);
    \node[below right] at (3,0) {\small $n_0$};
    \foreach \n in {1,...,14} {
      \pgfmathsetmacro{\y}{2*(1+1.5*ifthenelse(isodd(\n),-1,1)/\n)}
      \ifnum\n<5
        \filldraw[black] (0.6*\n,\y) circle (2pt);
      \else
        \filldraw[colorofwisdom] (0.6*\n,\y) circle (2pt);
      \fi
    }
  \end{tikzpicture}
\end{center}
Na obrázku je $a_n = 1 + \frac{{(-1)}^n\cdot 1{,}5}{n}$ s limitou $a = 1$ a pás $\varepsilon = 0{,}35$. Čtvrtý člen $1{,}375$ je ještě venku, ale od pátého ($n_0 = 5$) jsou už všechny členy uvnitř.

\begin{example}
  Dokažme, že $\lim\frac{1}{n} = 0$. Nechť je dáno $\varepsilon > 0$. Chceme $\left|\frac{1}{n} - 0\right| = \frac{1}{n} < \varepsilon$, tedy $n > \frac{1}{\varepsilon}$. Zvolíme tedy $n_0$ jako libovolné přirozené číslo větší než $\frac{1}{\varepsilon}$. Pro $\varepsilon = 0{,}001$ například $n_0 = 1001$.
\end{example}

\begin{definitionbox}
  \textbf{Pravidla pro limity}: Pokud $\lim a_n = a$ a $\lim b_n = b$ (konečné), pak
  \begin{equation*}
    \lim(a_n\pm b_n) = a\pm b,\qquad\lim(a_n\cdot b_n) = a\cdot b,\qquad\lim\frac{a_n}{b_n} = \frac{a}{b}\quad(b\neq 0).
  \end{equation*}
  \textbf{Důležité limity}:
  \begin{equation*}
    \lim\frac{1}{n^k} = 0\ (k > 0),\qquad\lim q^n = 0\ (|q| < 1),\qquad\lim\sqrt[n]{a} = 1\ (a > 0),\qquad\lim{\left(1 + \frac{1}{n}\right)}^n = e.
  \end{equation*}
\end{definitionbox}

\begin{example}
  Typická úloha: limita podílu dvou mnohočlenů. Vydělíme čitatele i jmenovatele nejvyšší mocninou $n$ ve jmenovateli:
  \begin{align*}
    \lim\frac{2n^2 + 1}{n^2 - 3} &= \lim\frac{2 + \frac{1}{n^2}}{1 - \frac{3}{n^2}} = \frac{2 + 0}{1 - 0} = 2,\\
    \lim\frac{3n + 1}{n^2 + 1} &= \lim\frac{\frac{3}{n} + \frac{1}{n^2}}{1 + \frac{1}{n^2}} = \frac{0}{1} = 0,\\
    \lim\frac{n^2}{n + 1} &= \lim\frac{n}{1 + \frac{1}{n}} = \infty.
  \end{align*}
  Pravidlo: rozhoduje nejvyšší mocnina. Stejné stupně dají podíl koeficientů, vyšší stupeň ve jmenovateli dá nulu, vyšší stupeň v čitateli dá nekonečno.
\end{example}

\begin{example}
  Limita $\lim\left(\sqrt{n + 1} - \sqrt{n}\right)$ je typu `$\infty - \infty$', což nic neříká. Pomůže rozšíření výrazem $\sqrt{n + 1} + \sqrt{n}$ (podobně jako usměrňování v kapitole~\ref{kap:vyrazy}):
  \begin{equation*}
    \sqrt{n + 1} - \sqrt{n} = \frac{(n + 1) - n}{\sqrt{n + 1} + \sqrt{n}} = \frac{1}{\sqrt{n + 1} + \sqrt{n}}\to 0.
  \end{equation*}
\end{example}

\begin{problem}
  Spočítejte:
  \begin{enumerate}[label=(\alph*)]
    \item $\lim\dfrac{5n^3 - 2n}{1 - 2n^3}$,
    \item $\lim\dfrac{n + {(-1)}^n}{n}$,
    \item $\lim\dfrac{2^n + 3^n}{3^n}$,
    \item $\lim\left(\sqrt{n^2 + n} - n\right)$.
  \end{enumerate}
\end{problem}

\begin{solution}
  \begin{enumerate}[label=(\alph*)]
    \item $-\frac{5}{2}$.
    \item $1 + \frac{{(-1)}^n}{n}\to 1$.
    \item ${\left(\frac{2}{3}\right)}^n + 1\to 1$.
    \item $\frac{n}{\sqrt{n^2 + n} + n} = \frac{1}{\sqrt{1 + \frac{1}{n}} + 1}\to\frac{1}{2}$.
  \end{enumerate}
\end{solution}

\section{Nekonečná geometrická řada}

\begin{definitionbox}
  \textbf{Nekonečná řada} $\sum_{n=1}^{\infty} a_n = a_1 + a_2 + a_3 + \cdots$ má \textbf{součet} $s$, pokud posloupnost částečných součtů $s_n = a_1 + \cdots + a_n$ má limitu $s$. Pak řada \textbf{konverguje}, jinak \textbf{diverguje}.

  \vspace{0.2cm}

  \textbf{Věta}: Nekonečná geometrická řada $a_1 + a_1q + a_1q^2 + \cdots$ (s $a_1\neq 0$) konverguje, právě když $|q| < 1$, a její součet je
  \begin{equation*}
    s = \frac{a_1}{1 - q}.
  \end{equation*}
\end{definitionbox}

\begin{proof}
  Pro $q\neq 1$ je $s_n = a_1\frac{q^n - 1}{q - 1} = a_1\frac{1 - q^n}{1 - q}$. Pro $|q| < 1$ je $\lim q^n = 0$, a tedy $\lim s_n = \frac{a_1}{1 - q}$. Pro $|q|\geq 1$ členy řady nejdou k nule (mají stálou nebo rostoucí absolutní hodnotu) a součet nemůže existovat.
\end{proof}

\begin{example}
  \begin{itemize}
    \item $1 + \frac{1}{2} + \frac{1}{4} + \frac{1}{8} + \cdots = \frac{1}{1 - \frac{1}{2}} = 2$. Pokud ujdete půl cesty, pak polovinu zbytku, pak polovinu zbytku\ldots, dojdete přesně do cíle.

    \begin{center}
      \begin{tikzpicture}[scale=0.9]
        \filldraw[fill=colorofwisdom!70, draw=white, thick] (0,0) rectangle (2,4);
        \filldraw[fill=colorofwisdom!50, draw=white, thick] (2,2) rectangle (4,4);
        \filldraw[fill=colorofwisdom!35, draw=white, thick] (2,0) rectangle (3,2);
        \filldraw[fill=colorofwisdom!22, draw=white, thick] (3,0) rectangle (4,1);
        \filldraw[fill=colorofwisdom!14, draw=white, thick] (3,1) rectangle (3.5,2);
        \filldraw[fill=colorofwisdom!8, draw=white, thick] (3.5,1.5) rectangle (4,2);
        \filldraw[fill=colorofwisdom!4, draw=white, thick] (3.5,1) rectangle (3.75,1.5);
        \draw[thick] (0,0) rectangle (4,4);
        \node[white] at (1,2) {$\frac{1}{2}$};
        \node[white] at (3,3) {$\frac{1}{4}$};
        \node at (2.5,1) {$\frac{1}{8}$};
        \node at (3.5,0.5) {\small $\frac{1}{16}$};
        \node[anchor=west, align=left, font=\small] at (4.6,2) {čtverec o obsahu $1$:\\každý dílek je polovina toho, co zbylo,\\a dílky čtverec postupně celý vyplní:\\$\frac{1}{2} + \frac{1}{4} + \frac{1}{8} + \cdots = 1$};
      \end{tikzpicture}
    \end{center}
    Obrázek ukazuje, proč součet bez úvodní jedničky vychází $1$: nevyplněný zbytek čtverce je po každém kroku poloviční, a proto se zmenšuje k nule.
    \item $9 - 3 + 1 - \frac{1}{3} + \cdots = \frac{9}{1 + \frac{1}{3}} = \frac{27}{4}$.
    \item $0{,}\overline{3} = \frac{3}{10} + \frac{3}{100} + \cdots = \frac{\frac{3}{10}}{1 - \frac{1}{10}} = \frac{1}{3}$.
    \item $0{,}\overline{12} = \frac{12}{100} + \frac{12}{10000} + \cdots = \frac{\frac{12}{100}}{1 - \frac{1}{100}} = \frac{12}{99} = \frac{4}{33}$.
  \end{itemize}
\end{example}

\begin{example}\label{ex:0999}
  Kolik je $0{,}\overline{9} = 0{,}999\ldots$? Podle vzorce
  \begin{equation*}
    0{,}\overline{9} = \frac{\frac{9}{10}}{1 - \frac{1}{10}} = 1.
  \end{equation*}
  Ano, přesně jedna. Ne `skoro jedna', ne `nekonečně blízko jedné'. Jedna. Číslo $0{,}999\ldots$ je jen jiný zápis čísla $1$, podobně jako $\frac{1}{2}$ a $\frac{2}{4}$ jsou dva zápisy téhož čísla. Pokud tomu nevěříte, zkuste najít číslo, které leží mezi $0{,}999\ldots$ a $1$. Žádné takové není, a dvě různá reálná čísla mají mezi sebou vždy nějaké další (například jejich průměr). Tento příklad mimochodem na internetu spolehlivě rozpoutá hádku v jakékoli diskuzi.
\end{example}

\begin{note}
  Řecký filozof Zénón z Eleje tvrdil, že rychlonohý Achilles nikdy nedohoní želvu, která má náskok. Než Achilles doběhne tam, kde byla želva, ona se posune o kousek dál. Než doběhne tam, posune se zase dál. A tak donekonečna. Řekové nevěděli, co s tím, protože nevěděli, jak sečíst nekonečně mnoho čísel. Dnes víme, že časy, které Achilles potřebuje na jednotlivé úseky, tvoří geometrickou řadu s konečným součtem. Achilles želvu dohoní v konečném čase. Jen to bude trvat nekonečně mnoho kroků, ale každý bude kratší.
\end{note}

\begin{problem}
  \begin{enumerate}[label=(\alph*)]
    \item Převeďte na zlomek $0{,}\overline{27}$ a $1{,}2\overline{3}$.
    \item Pro která $x$ konverguje řada $1 + (x - 1) + {(x - 1)}^2 + \cdots$? Jaký má součet?
    \item Míček pustíme z výšky $10\ \si{m}$. Po každém odrazu vyskočí do $\frac{3}{4}$ předchozí výšky. Jakou celkovou dráhu urazí, než se zastaví?
    \item Vyřešte rovnici $x + x^2 + x^3 + \cdots = 3$.
  \end{enumerate}
\end{problem}

\begin{solution}
  \begin{enumerate}[label=(\alph*)]
    \item $0{,}\overline{27} = \frac{27}{99} = \frac{3}{11}$. $1{,}2\overline{3} = 1{,}2 + \frac{\frac{3}{100}}{1 - \frac{1}{10}} = \frac{6}{5} + \frac{1}{30} = \frac{37}{30}$.
    \item $|x - 1| < 1$, tedy $x\in(0, 2)$, součet $\frac{1}{1 - (x - 1)} = \frac{1}{2 - x}$.
    \item Dolů $10\ \si{m}$, pak nahoru a dolů $7{,}5$, pak $5{,}625$ a tak dál: $10 + 2\cdot\frac{7{,}5}{1 - 0{,}75} = 10 + 60 = 70\ \si{m}$.
    \item Pro $|x| < 1$ je součet $\frac{x}{1 - x} = 3$, tedy $x = \frac{3}{4}$. To splňuje podmínku $|x| < 1$, takže $K = \left\{\frac{3}{4}\right\}$.
  \end{enumerate}
\end{solution}

\begin{note}
  Pozor, z toho, že členy řady jdou k nule, ještě neplyne, že řada konverguje! Slavný protipříklad je \textbf{harmonická řada}
  \begin{equation*}
    1 + \frac{1}{2} + \frac{1}{3} + \frac{1}{4} + \cdots = \infty.
  \end{equation*}
  Důkaz pochází od středověkého biskupa Mikuláše Oresma ze 14. století. Členy seskupíme:
  \begin{equation*}
    1 + \frac{1}{2} + \underbrace{\left(\frac{1}{3} + \frac{1}{4}\right)}_{\geq\frac{1}{4} + \frac{1}{4} = \frac{1}{2}} + \underbrace{\left(\frac{1}{5} + \cdots + \frac{1}{8}\right)}_{\geq 4\cdot\frac{1}{8} = \frac{1}{2}} + \underbrace{\left(\frac{1}{9} + \cdots + \frac{1}{16}\right)}_{\geq\frac{1}{2}} + \cdots
  \end{equation*}
  Každá skupina je alespoň $\frac{1}{2}$ a skupin je nekonečně mnoho. Řada ale roste neuvěřitelně pomalu: abyste se dostali přes $100$, museli byste sečíst asi $10^{43}$ členů.

  \begin{center}
    \begin{tikzpicture}[scale=0.85]
      \draw[->] (0,0) -- (8.6,0) node[right] {\small $n$};
      \draw[->] (0,0) -- (0,3.4);
      \foreach \n in {1,...,16} {
        \filldraw[fill=gray!20, draw=gray] ({0.5*(\n-1)},0) rectangle ({0.5*\n},{3/\n});
      }
      \foreach \n/\b in {3/4, 4/4, 5/8, 6/8, 7/8, 8/8, 9/16, 10/16, 11/16, 12/16, 13/16, 14/16, 15/16, 16/16} {
        \fill[colorofwisdom!55] ({0.5*(\n-1)},0) rectangle ({0.5*\n},{3/\b});
      }
      \fill[colorofwisdom!55] (0,0) rectangle (0.5,3);
      \fill[colorofwisdom!55] (0.5,0) rectangle (1,1.5);
      \node[font=\scriptsize] at (0.25,3.2) {$1$};
      \node[font=\scriptsize] at (0.75,1.7) {$\frac{1}{2}$};
      \draw[decorate, decoration={brace, mirror, amplitude=3pt}] (1,-0.1) -- (2,-0.1) node[midway, below=3pt] {\scriptsize $\geq\frac{1}{2}$};
      \draw[decorate, decoration={brace, mirror, amplitude=3pt}] (2,-0.1) -- (4,-0.1) node[midway, below=3pt] {\scriptsize $\geq\frac{1}{2}$};
      \draw[decorate, decoration={brace, mirror, amplitude=3pt}] (4,-0.1) -- (8,-0.1) node[midway, below=3pt] {\scriptsize $\geq\frac{1}{2}$};
    \end{tikzpicture}
  \end{center}
  Sloupek nad číslem $n$ má výšku $\frac{1}{n}$. Tmavé části ukazují Oresmův odhad: v každé skupině zkrátíme všechny sloupky na výšku toho nejnižšího. Tmavá plocha každé skupiny je pak přesně $\frac{1}{2}$, a přitom skupiny jsou stále širší. Sloupky sice jdou k nule, ale ne dost rychle.
\end{note}

\section{Příběh --- Muž, který spočítal nekonečno}

Pojem nekonečna dělal matematikům problémy od Zénóna až do 19. století. Pak přišel Georg Cantor, zakladatel teorie množin z kapitoly 2, a ukázal, že nekonečna jsou různě velká. Přirozených čísel je stejně jako racionálních (dají se seřadit do posloupnosti), ale reálných čísel je víc. Tak moc víc, že se do posloupnosti seřadit nedají. Důkaz, kterému se říká \textit{diagonální metoda}, je dnes jedním z nejslavnějších v celé matematice.

Tehdejší matematici byli pobouření. Leopold Kronecker, Cantorův bývalý učitel a jeden z nejvlivnějších matematiků v Německu, ho veřejně nazýval `šarlatánem' a `kazitelem mládeže' a bránil mu v získání místa na prestižní univerzitě v Berlíně. Cantor celý život učil na malé univerzitě v Halle. Od roku 1884 trpěl záchvaty těžké deprese a opakovaně byl hospitalizován v psychiatrických zařízeních. Ve chvílích, kdy se mu nedařilo v matematice, se pokoušel dokázat, že autorem Shakespearových her byl ve skutečnosti Francis Bacon.

Uznání se nakonec dočkal. David Hilbert, autor citátů na začátku této kapitoly i kapitoly 2, prohlásil, že `z ráje, který pro nás Cantor vytvořil, nás nikdo nevyžene'. Cantor ale zemřel v lednu 1918 v psychiatrické léčebně v Halle, podvyživený, protože během první světové války nebylo dost jídla. Ráj, který vytvořil pro ostatní, si sám užít nemohl.

\section{Shrnutí a otázky k ústní zkoušce}

\begin{problem}
  Odpovězte vlastními slovy:
  \begin{enumerate}
    \item Co je posloupnost? Jakými způsoby ji lze zadat? Jak vypadá její graf?
    \item Kdy je posloupnost monotónní a kdy omezená?
    \item Definujte aritmetickou posloupnost. Odvoďte vzorce pro $n$-tý člen a součet prvních $n$ členů.
    \item Definujte geometrickou posloupnost. Odvoďte vzorce pro $n$-tý člen a součet prvních $n$ členů.
    \item Proč se posloupnostem říká aritmetická a geometrická?
    \item Co je limita posloupnosti? Vysvětlete definici. Jak se počítají limity podílů mnohočlenů?
    \item Co je nekonečná řada a její součet? Kdy konverguje geometrická řada a jaký má součet?
    \item Proč je $0{,}\overline{9} = 1$? Jak převedete periodické číslo na zlomek?
    \item Uveďte příklad řady, jejíž členy jdou k nule, ale která nekonverguje.
    \item Uveďte příklady využití posloupností v praxi.
  \end{enumerate}
\end{problem}

\koneckapitoly

\chapter{Finanční matematika}
\label{kap:financni}
\mathquote{Compound interest is the eighth wonder of the world. He who understands it, earns it; he who doesn't, pays it.}{připisováno Albertu Einsteinovi (nejspíš neprávem)}

Finanční matematika je ta část matematiky, která vás bude v životě pronásledovat nejdéle. Spořicí účty, hypotéky, splátky telefonu, kreditní karty, inflace. Banky a obchodníci s tím počítají každý den a spoléhají na to, že vy ne. Matematicky v této kapitole nic nového nebude: všechno jsou to procenta, geometrické posloupnosti z kapitoly~\ref{kap:posloupnosti} a exponenciála z kapitoly~\ref{kap:exponencialni}. Nové bude jen to, že tentokrát jde o vaše peníze.

\section{Procenta}\label{sec:procenta}

\begin{definitionbox}
  \textbf{Procento} je setina celku: $1\,\% = \frac{1}{100} = 0{,}01$. \textbf{Promile} je tisícina: $0{,}1\,\%$, tedy $0{,}001$.

  \vspace{0.2cm}

  \begin{itemize}
    \item $p\,\%$ ze základu $z$ je \textbf{procentová část} $c = \frac{p}{100}\cdot z$,
    \item zvětšení o $p\,\%$ znamená vynásobení číslem $1 + \frac{p}{100}$, zmenšení o $p\,\%$ vynásobení číslem $1 - \frac{p}{100}$,
    \item \textbf{procentní bod} je rozdíl dvou procentních údajů. Pokud úroková sazba vzroste ze $2\,\%$ na $3\,\%$, vzrostla o $1$ procentní bod, ale o $50\,\%$.
  \end{itemize}
\end{definitionbox}

\begin{example}
  Bunda stála $2\,000$ Kč. Obchod ji zdražil o $20\,\%$ a pak ji ve výprodeji zlevnil o $20\,\%$. Kolik stojí teď?
  \begin{equation*}
    2000\cdot 1{,}2\cdot 0{,}8 = 2000\cdot 0{,}96 = 1920\ \text{Kč}.
  \end{equation*}
  Bunda je o $4\,\%$ levnější než na začátku, ne stejně drahá. Zlevnění se počítá z vyšší ceny, takže ubere víc, než předtím přidalo zdražení. Obchody to vědí. Zákazníci většinou ne.
  \begin{center}
    \begin{tikzpicture}[xscale=0.9, yscale=0.6]
      \draw[->] (0,0) -- (7.2,0);
      \draw[->] (0,0) -- (0,5.3) node[above] {\small Kč};
      \foreach \y/\l in {1/500, 2/1000, 3/1500, 4/2000} \draw (0.07,\y) -- (-0.07,\y) node[left] {\tiny $\l$};
      \filldraw[fill=gray!35, draw=black] (0.6,0) rectangle (1.8,4);
      \filldraw[fill=gray!35, draw=black] (2.9,0) rectangle (4.1,4);
      \filldraw[fill=colorofwisdom!60, draw=black] (2.9,4) rectangle (4.1,4.8);
      \filldraw[fill=gray!35, draw=black] (5.2,0) rectangle (6.4,3.84);
      \draw[dashed, colorofwisdom] (5.2,3.84) rectangle (6.4,4.8);
      \draw[dashed, gray] (0,4) -- (7,4);
      \node[font=\scriptsize] at (3.5,4.4) {$+400$};
      \node[font=\scriptsize, colorofwisdom] at (5.8,4.32) {$-480$};
      \node[below, font=\small] at (1.2,0) {původně};
      \node[below, font=\small] at (3.5,0) {$+20\,\%$};
      \node[below, font=\small] at (5.8,0) {$-20\,\%$};
      \node[above, font=\scriptsize] at (1.2,4) {$2000$};
      \node[above, font=\scriptsize] at (3.5,4.8) {$2400$};
      \node[below, font=\scriptsize] at (5.8,3.84) {$1920$};
    \end{tikzpicture}
  \end{center}
  Na obrázku je to vidět: oba kroky jsou `$20\,\%$', ale pokaždé z jiného základu. Zdražení přidá $20\,\%$ ze $2\,000$, tedy $400$ Kč, zlevnění ubere $20\,\%$ ze $2\,400$, tedy $480$ Kč.
\end{example}

\begin{example}
  Po zdražení o $25\,\%$ stojí lístek do kina $250$ Kč. Kolik stál předtím? Ne $250\cdot 0{,}75 = 187{,}5$ Kč, ale $\frac{250}{1{,}25} = 200$ Kč. Abyste zrušili zdražení o $25\,\%$, musíte zlevnit o $20\,\%$.
\end{example}

\begin{problem}
  \begin{enumerate}[label=(\alph*)]
    \item Cena akcie vzrostla o $10\,\%$, pak klesla o $10\,\%$, pak vzrostla o $10\,\%$. Jaká je celková změna?
    \item Kolika procenty je $45$ z $180$? A kolik je $180$ procent z $45$?
    \item Volební strana měla $8\,\%$ a teď má $10\,\%$. O kolik procentních bodů a o kolik procent si polepšila?
    \item O kolik procent se zvětší obsah čtverce, když jeho stranu zvětšíme o $10\,\%$?
  \end{enumerate}
\end{problem}

\begin{solution}
  \begin{enumerate}[label=(\alph*)]
    \item $1{,}1\cdot 0{,}9\cdot 1{,}1 = 1{,}089$, tedy vzrostla o $8{,}9\,\%$.
    \item $\frac{45}{180} = 0{,}25 = 25\,\%$. $1{,}8\cdot 45 = 81$.
    \item O $2$ procentní body, což je $\frac{2}{8} = 25\,\%$.
    \item $1{,}1^2 = 1{,}21$, tedy o $21\,\%$.
  \end{enumerate}
\end{solution}

\section{Úročení}\label{sec:uroceni}

\begin{definitionbox}
  Pojmy:
  \begin{itemize}
    \item \textbf{jistina} $K_0$: počáteční vklad (nebo dluh),
    \item \textbf{úroková sazba} $i$: kolik procent z jistiny se připíše za jedno úrokovací období (obvykle rok, zkratka p.\,a., \textit{per annum}). Ve vzorcích ji píšeme jako desetinné číslo, tedy $3\,\% = 0{,}03$,
    \item \textbf{úrok}: částka, kterou za půjčení peněz dostaneme (nebo zaplatíme).
  \end{itemize}

  \textbf{Jednoduché úročení}: úrok se počítá stále jen z původní jistiny. Po $n$ obdobích
  \begin{equation*}
    K_n = K_0(1 + n\cdot i).
  \end{equation*}

  \textbf{Složené úročení}: úrok se připíše k jistině a v dalším období se úročí i on. Po $n$ obdobích
  \begin{equation*}
    K_n = K_0{(1 + i)}^n.
  \end{equation*}
\end{definitionbox}

Jednoduché úročení je aritmetická posloupnost (každý rok přibude stejná částka), složené úročení je geometrická posloupnost s kvocientem $1 + i$ (každý rok se částka vynásobí stejným číslem). Na krátkou dobu je rozdíl malý, na dlouhou obrovský, protože geometrická posloupnost roste exponenciálně.

Představte si to takhle. Při jednoduchém úročení dostáváte každý rok stejnou `výplatu' z původní jistiny. Při složeném úročení pracují i úroky z minulých let: úrok z úroku, pak úrok z úroku z úroku a tak dál. Na grafu je vývoj vkladu $10\,000$ Kč při $3\,\%$ ročně z následujícího příkladu, vždy po pěti letech. Šedé sloupce rostou rovnoměrně (přímka), červené stále rychleji (exponenciála).

\begin{center}
  \begin{tikzpicture}[yscale=1.1, xscale=0.95]
    \draw[->] (0,0) -- (11.3,0) node[below right] {\small roky};
    \draw[->] (0,0) -- (0,3.9) node[above] {\small tis. Kč};
    \foreach \y in {1,2,3} {
      \draw (0.07,\y) -- (-0.07,\y) node[left] {\tiny $\y0$};
      \draw[gray!30] (0.07,\y) -- (11,\y);
    }
    \foreach \k in {0,...,8} {
      \pgfmathsetmacro{\n}{5*\k}
      \pgfmathsetmacro{\s}{1 + 0.03*\n}
      \pgfmathsetmacro{\c}{pow(1.03,\n)}
      \filldraw[fill=gray!40, draw=black] ({1.2*\k+0.25},0) rectangle ({1.2*\k+0.6},\s);
      \filldraw[fill=colorofwisdom!70, draw=black] ({1.2*\k+0.6},0) rectangle ({1.2*\k+0.95},\c);
      \pgfmathtruncatemacro{\nn}{\n}
      \node[below, font=\tiny] at ({1.2*\k+0.6},0) {$\nn$};
    }
    \node[right, font=\scriptsize] at (10.6,2.2) {$22\,000$};
    \node[right, font=\scriptsize, colorofwisdom] at (10.6,3.26) {$32\,620$};
    \filldraw[fill=gray!40, draw=black] (0.5,3.55) rectangle (0.8,3.75);
    \node[right, font=\scriptsize] at (0.8,3.65) {jednoduché};
    \filldraw[fill=colorofwisdom!70, draw=black] (3.0,3.55) rectangle (3.3,3.75);
    \node[right, font=\scriptsize] at (3.3,3.65) {složené};
  \end{tikzpicture}
\end{center}

\begin{example}
  Vložíme $10\,000$ Kč na $4$ roky s úrokem $3\,\%$ ročně.
  \begin{itemize}
    \item Jednoduché úročení: $10000(1 + 4\cdot 0{,}03) = 11\,200$ Kč.
    \item Složené úročení: $10000\cdot 1{,}03^4\approx 11\,255{,}09$ Kč.
  \end{itemize}
  Za čtyři roky je rozdíl jen $55$ Kč. Za čtyřicet let by to bylo $22\,000$ Kč proti $32\,620$ Kč.
\end{example}

\begin{definitionbox}
  \textbf{Daň z úroků}: V České republice se z úroků na běžných účtech odvádí \textbf{srážková daň} $15\,\%$. Ze sazby $i$ nám tedy reálně zbude $0{,}85\cdot i$ a
  \begin{equation*}
    K_n = K_0{(1 + 0{,}85\cdot i)}^n.
  \end{equation*}

  \textbf{Častější připisování úroků}: Pokud je roční (nominální) sazba $i$ a úroky se připisují $m$-krát ročně, připíše se pokaždé $\frac{i}{m}$ a za $n$ let
  \begin{equation*}
    K_n = K_0{\left(1 + \frac{i}{m}\right)}^{mn}.
  \end{equation*}
  Skutečný roční výnos (\textbf{efektivní sazba}) je pak ${\left(1 + \frac{i}{m}\right)}^m - 1$, což je víc než $i$.
\end{definitionbox}

\begin{example}\label{ex:efektivni-sazba}
  Při sazbě $3\,\%$ a dani $15\,\%$ máme z $10\,000$ Kč za $4$ roky $10000\cdot 1{,}0255^4\approx 11\,059{,}68$ Kč.

  Při nominální sazbě $6\,\%$ s měsíčním připisováním je efektivní sazba ${(1 + 0{,}005)}^{12} - 1\approx 6{,}17\,\%$. Pokud bychom úroky připisovali nekonečně často (spojité úročení), dostali bychom efektivní sazbu $e^{0{,}06} - 1\approx 6{,}18\,\%$. Proč právě $e$? Připomeňme definici z kapitol~\ref{kap:exponencialni} a~\ref{kap:posloupnosti}:
  \begin{equation*}
    e = \lim_{n\to\infty}{\left(1 + \frac{1}{n}\right)}^n.
  \end{equation*}
  Pokud v našem výrazu označíme $m = n\cdot i$ (tedy počet připsání za rok je $n$-krát roční sazba), dostaneme
  \begin{equation*}
    {\left(1 + \frac{i}{m}\right)}^m = {\left(1 + \frac{1}{n}\right)}^{n\cdot i} = {\left({\left(1 + \frac{1}{n}\right)}^n\right)}^i\ \longrightarrow\ e^i\qquad\text{pro } m\to\infty.
  \end{equation*}
  Při spojitém úročení tedy jistina roste přesně jako exponenciála: $K(t) = K_0e^{it}$. Tady se Eulerovo číslo z kapitoly~\ref{kap:exponencialni} objevuje přesně v tom kontextu, ve kterém ho objevil Jacob Bernoulli.
\end{example}

\begin{note}
  \textbf{Pravidlo 72}: Doba (v letech), za kterou se vklad při složeném úročení zdvojnásobí, je přibližně $\frac{72}{p}$, kde $p$ je sazba v procentech. Při $7\,\%$ je to asi $10$ let (přesně $\frac{\ln 2}{\ln 1{,}07}\approx 10{,}24$). Platí to proto, že $\ln 2\approx 0{,}693$ a pro malá $i$ je $\ln(1 + i)\approx i$. Číslo $72$ se používá místo $69$ proto, že se lépe dělí (dvěma, třemi, čtyřmi, šesti, osmi, devíti a dvanácti).


  Pravidlo skrývá ještě jednu pěknou věc: doba zdvojnásobení \textit{nezávisí na tom, kolik zrovna máte}. Z $1$ na $2$ to trvá stejně dlouho jako ze $2$ na $4$ nebo ze $4$ na $8$. Exponenciála je jediná křivka s touto vlastností, a proto se po stejných časových krocích stále zdvojnásobuje:
  \begin{center}
    \begin{tikzpicture}[xscale=0.3, yscale=0.42]
      \draw[->] (0,0) -- (34,0) node[right] {\small roky};
      \draw[->] (0,0) -- (0,9.4) node[above] {\small $K_n/K_0$};
      \foreach \y in {1,2,4,8} {
        \draw[dashed, gray] (0,\y) -- ({ln(\y)/ln(1.07)},\y) -- ({ln(\y)/ln(1.07)},0);
        \node[left, font=\scriptsize] at (0,\y) {$\y$};
      }
      \foreach \t in {10,20,30} \node[below, font=\scriptsize] at (\t,0) {$\t$};
      \draw[thick, colorofwisdom, domain=0:32.5, samples=60] plot (\x, {pow(1.07,\x)});
      \foreach \a/\b in {0/10.24, 10.24/20.49, 20.49/30.73} {
        \draw[<->, colorofwisdom] (\a,-1.6) -- (\b,-1.6);
        \node[below, font=\tiny] at ({(\a+\b)/2},-1.6) {$\approx 10$ let};
      }
      \node[font=\small] at (8,7) {$1{,}07^n$};
    \end{tikzpicture}
  \end{center}
\end{note}

\begin{definitionbox}
  \textbf{Inflace} je růst cenové hladiny. Pokud je úroková sazba $i$ a inflace $\pi$, je \textbf{reálný výnos}
  \begin{equation*}
    r = \frac{1 + i}{1 + \pi} - 1\approx i - \pi.
  \end{equation*}
\end{definitionbox}

\begin{example}
  Banka nabízí $4\,\%$, inflace je $6\,\%$. Reálný výnos je $\frac{1{,}04}{1{,}06} - 1\approx -1{,}89\,\%$. Peníze na účtu se sice množí, ale koupíte si za ně každý rok méně. A to jsme ještě nezapočítali daň.
\end{example}

\begin{problem}
  \begin{enumerate}[label=(\alph*)]
    \item Jakou roční sazbu potřebujeme, aby se vklad za $10$ let zdvojnásobil?
    \item Auto za $500\,000$ Kč ztrácí každý rok $15\,\%$ své hodnoty. Kolik bude stát za $4$ roky?
    \item Kolik musíme dnes vložit při sazbě $5\,\%$ (bez daně), abychom za $18$ let měli $100\,000$ Kč?
    \item Babička vám v den vašeho narození uložila $10\,000$ Kč na $4\,\%$ ročně (bez daně). Kolik budete mít k maturitě (v $19$ letech)?
  \end{enumerate}
\end{problem}

\begin{solution}
  \begin{enumerate}[label=(\alph*)]
    \item ${(1 + i)}^{10} = 2$, $i = \sqrt[10]{2} - 1\approx 7{,}18\,\%$.
    \item $500000\cdot 0{,}85^4\approx 261\,003$ Kč. Auto ztratí skoro polovinu hodnoty.
    \item $K_0 = \frac{100000}{1{,}05^{18}}\approx 41\,552$ Kč. Tomuto se říká \textbf{současná hodnota} budoucí částky.
    \item $10000\cdot 1{,}04^{19}\approx 21\,068$ Kč. Za tu dobu se vklad zhruba zdvojnásobí, přesně podle pravidla 72 ($\frac{72}{4} = 18$).
  \end{enumerate}
\end{solution}

\section{Spoření}\label{sec:sporeni}

Při pravidelném spoření vkládáme každé období stejnou částku $a$. Každý vklad se pak úročí jinou dobu: ten první nejdéle, ten poslední nejkratší dobu nebo vůbec. Výsledná částka je součet geometrické posloupnosti.

\begin{center}
  \begin{tikzpicture}[scale=1]
    \draw[->] (-0.3,0) -- (10.6,0) node[right] {\small čas};
    \foreach \k in {0,...,5} {
      \draw (2*\k,0.1) -- (2*\k,-0.1) node[below] {\scriptsize $\k$};
    }
    \foreach \k/\e in {1/4, 2/3, 3/2, 4/1} {
      \filldraw[colorofwisdom] (2*\k,0) circle (2pt);
      \node[above, font=\scriptsize] at (2*\k,0.05) {$a$};
      \draw[->, colorofwisdom!80, looseness=0.9] (2*\k,0.35) to[out=60, in=180] (9.6,{0.35+0.45*\e});
      \ifnum\e=1
        \node[right, font=\scriptsize] at (10,{0.35+0.45*\e}) {$a(1 + i)$};
      \else
        \node[right, font=\scriptsize] at (10,{0.35+0.45*\e}) {$a{(1 + i)}^{\e}$};
      \fi
    }
    \filldraw[colorofwisdom] (10,0) circle (2pt);
    \node[right, font=\scriptsize] at (10,0.35) {$a$};
    \node[font=\scriptsize, align=center] at (0,-0.9) {začátek};
    \node[font=\scriptsize, align=center] at (10,-0.9) {konec};
  \end{tikzpicture}
\end{center}
Na obrázku je spoření po dobu $n = 5$ období. Každý vklad `jede' do konce jinak dlouho: první vklad se úročí čtyři období, poslední vůbec. Na konci sečteme pět čísel $a, a(1 + i), \ldots, a{(1 + i)}^4$, a to je geometrická posloupnost.

\begin{definitionbox}
  \textbf{Pravidelné spoření}: Pokud na konci každého z $n$ období vložíme částku $a$ a sazba za období je $i$, máme na konci
  \begin{equation*}
    S_n = a\cdot\frac{{(1 + i)}^n - 1}{i}.
  \end{equation*}
  Pokud vkládáme na \textbf{začátku} každého období, každý vklad se úročí o jedno období déle a
  \begin{equation*}
    S_n = a{(1 + i)}\cdot\frac{{(1 + i)}^n - 1}{i}.
  \end{equation*}
\end{definitionbox}

\begin{proof}
  Při vkladech na konci období se poslední vklad neúročí vůbec, předposlední jedno období a první $n - 1$ období. Celkem
  \begin{equation*}
    S_n = a + a(1 + i) + a{(1 + i)}^2 + \cdots + a{(1 + i)}^{n-1}.
  \end{equation*}
  To je součet $n$ členů geometrické posloupnosti s prvním členem $a$ a kvocientem $q = 1 + i$. Podle kapitoly~\ref{kap:posloupnosti}
  \begin{equation*}
    S_n = a\cdot\frac{{(1 + i)}^n - 1}{(1 + i) - 1} = a\cdot\frac{{(1 + i)}^n - 1}{i}.
  \end{equation*}
\end{proof}

\begin{example}
  Každý měsíc na konci měsíce vložíme $1\,000$ Kč na účet s roční sazbou $3\,\%$ (měsíčně tedy $0{,}25\,\%$, daň zanedbáme). Kolik budeme mít za dva roky?
  \begin{equation*}
    S_{24} = 1000\cdot\frac{1{,}0025^{24} - 1}{0{,}0025}\approx 24\,703\ \text{Kč}.
  \end{equation*}
  Vložili jsme $24\,000$ Kč, úroky dělají $703$ Kč. Na krátkou dobu se spoření moc nevyplatí. Na čtyřicet let by to ale bylo přes $925\,000$ Kč z vložených $480\,000$ Kč.
\end{example}

\section{Úvěry}\label{sec:uvery}

Úvěr je spoření obráceně. Banka nám půjčí částku $U$ a my ji splácíme pravidelnými splátkami. Každá splátka se skládá z \textbf{úroku} (z dosud nesplaceného dluhu) a \textbf{úmoru} (části, která skutečně snižuje dluh).

\begin{definitionbox}
  \textbf{Anuitní splácení}: Úvěr $U$ splácíme $n$ stejnými splátkami $a$ na konci každého období při sazbě $i$ za období. Pak
  \begin{equation*}
    a = U\cdot\frac{i{(1 + i)}^n}{{(1 + i)}^n - 1}.
  \end{equation*}
\end{definitionbox}

\begin{proof}
  Dluh po jednom období je $U(1 + i) - a$, po dvou $\left(U(1 + i) - a\right)(1 + i) - a = U{(1 + i)}^2 - a(1 + i) - a$, a tak dál. Po $n$ obdobích má být dluh nulový:
  \begin{equation*}
    U{(1 + i)}^n - a\left({(1 + i)}^{n-1} + \cdots + (1 + i) + 1\right) = 0.
  \end{equation*}
  Závorka je stejná geometrická řada jako u spoření, tedy $\frac{{(1 + i)}^n - 1}{i}$. Odtud
  \begin{equation*}
    U{(1 + i)}^n = a\cdot\frac{{(1 + i)}^n - 1}{i}\quad\Leftrightarrow\quad a = U\cdot\frac{i{(1 + i)}^n}{{(1 + i)}^n - 1}.
  \end{equation*}
  Rovnice má mimochodem pěknou interpretaci: to, na kolik by narostl dluh, kdybychom nesplácali (vlevo), se musí rovnat tomu, na kolik by narostly splátky, kdybychom je místo toho spořili (vpravo).
\end{proof}

\begin{example}\label{ex:splatkovy-kalendar}
  Půjčíme si $200\,000$ Kč na $5$ let s úrokem $6\,\%$ ročně a splácíme jednou ročně. Roční splátka je
  \begin{equation*}
    a = 200000\cdot\frac{0{,}06\cdot 1{,}06^5}{1{,}06^5 - 1}\approx 47\,479{,}28\ \text{Kč}.
  \end{equation*}
  Celkem zaplatíme $5a\approx 237\,396$ Kč, z toho úroky jsou asi $37\,396$ Kč. Začátek \textbf{splátkového kalendáře}:
  \begin{center}
    \renewcommand{\arraystretch}{1.3}
    \begin{tabular}{|c|r|r|r|r|}
      \hline
      rok & dluh na začátku & úrok ($6\,\%$) & úmor & dluh na konci \\
      \hline
      $1$ & $200\,000{,}00$ & $12\,000{,}00$ & $35\,479{,}28$ & $164\,520{,}72$ \\
      $2$ & $164\,520{,}72$ & $9\,871{,}24$ & $37\,608{,}04$ & $126\,912{,}68$ \\
      \vdots & \vdots & \vdots & \vdots & \vdots \\
      \hline
    \end{tabular}
  \end{center}
  Na začátku tvoří velkou část splátky úrok, ke konci skoro jen úmor. U třicetileté hypotéky to znamená, že prvních deset let splácíte hlavně bance, ne dluh.

  \begin{center}
    \begin{tikzpicture}[scale=0.8]
      \begin{scope}
        \draw[->] (0,0) -- (6.3,0) node[right] {\small rok};
        \draw[->] (0,0) -- (0,5.6) node[above] {\small tis. Kč};
        \foreach \y in {1,2,3,4} \draw (0.07,\y) -- (-0.07,\y) node[left] {\tiny $\y0$};
        \foreach \k/\u in {1/1.2, 2/0.987124, 3/0.761476, 4/0.522289, 5/0.268751} {
          \filldraw[fill=gray!40, draw=black] ({\k-0.35},0) rectangle ({\k+0.35},{4.747928-\u});
          \filldraw[fill=colorofwisdom!70, draw=black] ({\k-0.35},{4.747928-\u}) rectangle ({\k+0.35},4.747928);
          \node[below, font=\scriptsize] at (\k,0) {$\k$};
        }
        \filldraw[fill=colorofwisdom!70, draw=black] (0.3,5.15) rectangle (0.6,5.35);
        \node[right, font=\scriptsize] at (0.6,5.25) {úrok};
        \filldraw[fill=gray!40, draw=black] (2.4,5.15) rectangle (2.7,5.35);
        \node[right, font=\scriptsize] at (2.7,5.25) {úmor};
        \node[font=\small] at (3,-1) {složení splátky};
      \end{scope}
      \begin{scope}[xshift=8.3cm]
        \draw[->] (0,0) -- (6.2,0) node[right] {\small rok};
        \draw[->] (0,0) -- (0,5.6) node[above] {\small zbývající dluh};
        \foreach \y/\l in {1/25, 2/50, 3/75, 4/100} \draw (0.07,\y) -- (-0.07,\y) node[left] {\tiny $\l\,\%$};
        \foreach \t in {0,10,20,30} \draw (0.18*\t,0.07) -- (0.18*\t,-0.07) node[below] {\scriptsize $\t$};
        \draw[dashed, gray] (0,4) -- (5.4,0);
        \draw[thick, colorofwisdom, domain=0:30, samples=60] plot ({0.18*\x}, {4*(4.321942-pow(1.05,\x))/3.321942});
        \draw[dashed] (1.8,0) -- (1.8,3.243) -- (0,3.243);
        \filldraw[colorofwisdom] (1.8,3.243) circle (2pt) node[above right, font=\scriptsize] {po $10$ letech zbývá $81\,\%$};
        \node[font=\small] at (2.8,-1) {hypotéka na $30$ let při $5\,\%$};
      \end{scope}
    \end{tikzpicture}
  \end{center}
  Levý graf: všechny splátky jsou stejně vysoké, ale červená část (úrok) se zmenšuje a šedá (úmor) roste. Dluh proto neklesá rovnoměrně, ale zpočátku pomaleji a ke konci rychleji. U pětiletého úvěru je to skoro neznatelné, u dlouhé hypotéky už velmi. Pravý graf ukazuje zbývající dluh u hypotéky na $30$ let při $5\,\%$ ročně: po třetině doby splácení zbývá splatit ještě $81\,\%$ dluhu (dluh po $k$ letech je $U\cdot\frac{1{,}05^{30} - 1{,}05^k}{1{,}05^{30} - 1}$, což plyne ze stejného výpočtu jako v důkazu výše). Šedá přímka ukazuje, jak by dluh klesal, kdyby se splácel rovnoměrně.
\end{example}

\begin{note}\label{note:rpsn}
  Při porovnávání úvěrů se nedívejte jen na úrokovou sazbu, ale na \textbf{RPSN} (roční procentní sazba nákladů). Ta zahrnuje i poplatky za vedení úvěru, pojištění a další náklady. Půjčka `bez úroku' může mít RPSN klidně desítky procent. Nebankovní půjčky v minulosti nabízely RPSN v řádu stovek i tisíců procent, než zákon o spotřebitelském úvěru v roce 2016 výrazně zpřísnil pravidla pro jejich poskytovatele.
\end{note}

\begin{problem}
  \begin{enumerate}[label=(\alph*)]
    \item Telefon za $24\,000$ Kč koupíme na $12$ měsíčních splátek s měsíčním úrokem $1\,\%$. Jaká je splátka a kolik celkem přeplatíme?
    \item Kolik musíme spořit na konci každého roku při $4\,\%$, abychom za $10$ let měli $500\,000$ Kč?
    \item Úvěr $1\,000\,000$ Kč na $20$ let při $5\,\%$ ročně. Jaká je roční splátka a kolik zaplatíme celkem?
  \end{enumerate}
\end{problem}

\begin{solution}
  \begin{enumerate}[label=(\alph*)]
    \item $a = 24000\cdot\frac{0{,}01\cdot 1{,}01^{12}}{1{,}01^{12} - 1}\approx 2\,132{,}37$ Kč. Celkem $25\,588$ Kč, přeplatíme asi $1\,588$ Kč.
    \item $a = 500000\cdot\frac{0{,}04}{1{,}04^{10} - 1}\approx 41\,645$ Kč ročně.
    \item $a = 10^6\cdot\frac{0{,}05\cdot 1{,}05^{20}}{1{,}05^{20} - 1}\approx 80\,243$ Kč ročně, celkem asi $1\,604\,852$ Kč. Na úrocích zaplatíme přes $600\,000$ Kč.
  \end{enumerate}
\end{solution}

\section{Příběh --- Lichváři v sedmém kruhu}

Půjčovat peníze na úrok bylo po většinu evropských dějin hříchem. Už První nikajský koncil v roce 325 zakázal brát úroky duchovním a Třetí lateránský koncil v roce 1179 rozhodl, že lichváři budou exkomunikováni a nesmí být pohřbeni v posvěcené půdě. Za \textit{lichvu} se přitom považoval jakýkoli úrok, ne jen přemrštěný.

Dante Alighieri ve své \textit{Božské komedii} umístil lichváře do sedmého kruhu pekla, mezi ty, kdo se provinili násilím proti Bohu a přírodě. Sedí tam na rozžhaveném písku pod deštěm ohně a každý má na krku měšec se znakem svého rodu, na který neustále zírá. Dante je tam řadí vedle rouhačů a sodomitů. Podle středověké logiky totiž lichvář nechává peníze `rodit' další peníze, což je proti přirozenosti, protože peníze jsou neplodné.

Shakespearův \textit{Kupec benátský} pak vytvořil postavu Shylocka, lichváře, který si jako zástavu za půjčku vymíní libru masa z dlužníkova těla. Dnešní banky naštěstí berou jen nemovitosti. Ale principiálně se toho zas tolik nezměnilo: kdo neumí počítat složené úročení, ten platí. Možná proto je ten citát na začátku kapitoly tak oblíbený, i když ho Einstein nejspíš nikdy neřekl.

\section{Shrnutí a otázky k ústní zkoušce}

\begin{problem}
  Odpovězte vlastními slovy:
  \begin{enumerate}
    \item Co je procento a procentní bod? Proč zdražení o $20\,\%$ a následné zlevnění o $20\,\%$ nevrátí původní cenu?
    \item Jaký je rozdíl mezi jednoduchým a složeným úročením? Jak souvisí s aritmetickou a geometrickou posloupností?
    \item Jak se do výpočtu promítne daň z úroků a častější připisování úroků? Co je efektivní sazba?
    \item Co je inflace a reálný výnos?
    \item Odvoďte vzorec pro pravidelné spoření.
    \item Co je anuitní splácení? Odvoďte vzorec pro výši splátky.
    \item Z čeho se skládá splátka úvěru a jak se poměr mezi úrokem a úmorem v čase mění?
    \item Co je RPSN a proč je důležitější než samotná úroková sazba?
  \end{enumerate}
\end{problem}

\koneckapitoly

\chapter{Kombinatorika}
\label{kap:kombinatorika}
\mathquote{Counting is the religion of this generation. It is its hope and its salvation.}{Gertrude Stein}

Kombinatorika je umění počítat, aniž bychom počítali. Kolik je možných PINů? Kolika způsoby se může seřadit třída na fotku? Kolik je tiketů ve Sportce? Vypisovat všechny možnosti by trvalo věčnost, takže potřebujeme chytřejší metodu. Kombinatorika je také základem celé pravděpodobnosti, které se věnuje další kapitola. Bez počítání možností totiž nelze spočítat ani šance.

\section{Základní pravidla}\label{sec:kombinatorika-pravidla}

\begin{definitionbox}
  \textbf{Pravidlo součtu}: Pokud lze jednu věc vybrat $m$ způsoby a jinou (disjunktní, tedy se s ní nepřekrývající) $n$ způsoby, lze vybrat jednu z nich $m + n$ způsoby. Pro množiny: $|A\cup B| = |A| + |B|$, pokud $A\cap B = \emptyset$.

  \vspace{0.2cm}

  \textbf{Pravidlo součinu}: Pokud se výběr skládá ze dvou kroků, kde první krok lze provést $m$ způsoby a druhý (pro každou volbu prvního) $n$ způsoby, lze celý výběr provést $m\cdot n$ způsoby. Pro množiny: $|A\times B| = |A|\cdot|B|$.
\end{definitionbox}

Pravidlo součinu je jen počítání prvků kartézského součinu z kapitoly 2. Pokud máte $3$ trička a $4$ kalhoty, máte $3\cdot 4 = 12$ outfitů, protože outfit je uspořádaná dvojice (tričko, kalhoty).

Nejlepší pomůcka pro pravidlo součinu je \textbf{strom možností}. Každé rozhodnutí je jedno patro stromu a každá cesta od kořene k listu je jedna možnost. Na obrázku jsou $3$ trička a $2$ kalhoty: z kořene vedou $3$ větve, z každé z nich $2$, a listů je proto $3\cdot 2 = 6$.
\begin{center}
  \begin{tikzpicture}[xscale=1.6, yscale=0.55, every node/.style={font=\small}]
    \node[circle, fill=black, inner sep=1.5pt] (r) at (0,0) {};
    \foreach \t/\y/\c in {1/3/colorofwisdom, 2/0/black, 3/-3/gray} {
      \node[draw, fill=\c!20, rounded corners=1pt] (t\t) at (1.3,\y) {tričko \t};
      \draw (r) -- (t\t);
      \foreach \k/\dy in {1/0.8, 2/-0.8} {
        \node[draw, rounded corners=1pt] (k\t\k) at (2.9,\y+\dy) {kalhoty \k};
        \draw (t\t) -- (k\t\k);
        \node[right] at (3.45,\y+\dy) {$\to$ outfit (\t, \k)};
      }
    }
  \end{tikzpicture}
\end{center}

\begin{example}
  \begin{itemize}
    \item Čtyřmístných PINů (číslice $0$--$9$) je $10\cdot 10\cdot 10\cdot 10 = 10^4 = 10\,000$.
    \item Čtyřmístných čísel (první číslice nesmí být nula) je $9\cdot 10^3 = 9\,000$.
    \item Čtyřmístných čísel s různými číslicemi je $9\cdot 9\cdot 8\cdot 7 = 4\,536$. Na prvním místě $9$ možností (ne nula), na druhém zase $9$ (už i nula, ale ne ta první číslice), pak $8$ a $7$. Pomáhá si nakreslit \textbf{políčka} a do každého napsat, kolik možností na něm zbývá:
    \begin{center}
      \begin{tikzpicture}[scale=0.9]
        \foreach \i/\n/\l in {0/9/{ne $0$}, 1/9/{ne 1.}, 2/8/{ne 1., 2.}, 3/7/{ne 1., 2., 3.}} {
          \draw[thick, fill=colorofwisdom!12] (1.6*\i,0) rectangle (1.6*\i+1.2,1);
          \node at (1.6*\i+0.6,0.5) {\Large $\n$};
          \node[below, font=\tiny] at (1.6*\i+0.6,0) {\l};
        }
        \foreach \i in {0,1,2} \node at (1.6*\i+1.4,0.5) {$\cdot$};
        \node at (6.9,0.5) {$= 4\,536$};
      \end{tikzpicture}
    \end{center}
  \end{itemize}
\end{example}

\begin{definitionbox}
  \textbf{Faktoriál} přirozeného čísla $n$ je
  \begin{equation*}
    n! = n\cdot(n - 1)\cdot\ \cdots\ \cdot 2\cdot 1,\qquad 0! = 1.
  \end{equation*}
  Rekurentně: $0! = 1$ a $(n + 1)! = (n + 1)\cdot n!$.
\end{definitionbox}

Proč $0! = 1$? Ze stejného důvodu jako $a^0 = 1$ v kapitole~\ref{kap:vyrazy}: chceme, aby rekurentní vztah fungoval i pro $n = 0$, tedy $1! = 1\cdot 0!$. Navíc $n!$ bude počet způsobů, jak seřadit $n$ věcí, a nic se dá seřadit právě jedním způsobem. Faktoriály rostou neuvěřitelně rychle: $10! = 3\,628\,800$ a $52!$ (počet způsobů, jak zamíchat balíček karet) má $68$ cifer. Pokaždé, když pořádně zamícháte karty, s téměř jistotou vytvoříte pořadí, které v dějinách vesmíru ještě nikdy neexistovalo.

\section{Výběry bez opakování}\label{sec:vybery-bez-opakovani}

\begin{definitionbox}
  Mějme $n$ různých prvků a vybíráme z nich $k$ (každý nejvýš jednou).

  \vspace{0.2cm}

  \textbf{Variace} $k$-té třídy z $n$ prvků: uspořádané $k$-tice (\textbf{na pořadí záleží}). Jejich počet je
  \begin{equation*}
    V(k, n) = n\cdot(n - 1)\cdot\ \cdots\ \cdot(n - k + 1) = \frac{n!}{(n - k)!}.
  \end{equation*}

  \textbf{Permutace} z $n$ prvků: variace $n$-té třídy, tedy seřazení všech prvků. Jejich počet je
  \begin{equation*}
    P(n) = n!.
  \end{equation*}

  \textbf{Kombinace} $k$-té třídy z $n$ prvků: $k$-prvkové podmnožiny (\textbf{na pořadí nezáleží}). Jejich počet je
  \begin{equation*}
    C(k, n) = \binom{n}{k} = \frac{n!}{k!\,(n - k)!}.
  \end{equation*}
  Číslo $\binom{n}{k}$ čteme `$n$ nad $k$' a nazýváme \textbf{kombinační číslo}.
\end{definitionbox}

Vzorec pro variace plyne z pravidla součinu: na první místo vybíráme z $n$ prvků, na druhé z $n - 1$ zbylých a tak dál. Kombinace pak dostaneme z variací tak, že nebudeme rozlišovat pořadí. Každou $k$-prvkovou podmnožinu lze seřadit $k!$ způsoby, takže ve variacích je započítaná $k!$-krát:
\begin{equation*}
  C(k, n) = \frac{V(k, n)}{k!} = \frac{n!}{k!(n - k)!}.
\end{equation*}
Na obrázku je to vidět pro výběr tří písmen: šest různých pořadí (variací) písmen $A$, $B$, $C$ odpovídá jediné trojici (kombinaci) $\{A, B, C\}$. Kombinací je tedy $3! = 6$-krát méně než variací.
\begin{center}
  \begin{tikzpicture}[scale=0.9]
    \foreach \w/\i in {ABC/0, ACB/1, BAC/2, BCA/3, CAB/4, CBA/5} {
      \node[draw, rounded corners=2pt, fill=black!5, minimum width=1.1cm] (v\i) at (0,-0.6*\i) {\small $\w$};
    }
    \node[draw, thick, rounded corners=3pt, fill=colorofwisdom!20, minimum width=1.6cm, minimum height=0.8cm] (k) at (4,-1.5) {$\{A, B, C\}$};
    \foreach \i in {0,...,5} \draw[->, gray] (v\i.east) -- (k.west);
    \node[above, font=\small] at (0,0.35) {variace (záleží na pořadí)};
    \node[above, font=\small] at (4,-1) {kombinace};
  \end{tikzpicture}
\end{center}

\begin{definitionbox}
  \textbf{Jak poznat, co použít?} Položte si dvě otázky:
  \begin{enumerate}
    \item \textbf{Záleží na pořadí?} Pokud ano (stupně vítězů, hesla, pořadí ve frontě), jde o variace nebo permutace. Pokud ne (výběr týmu, karty v ruce, čísla ve Sportce), jde o kombinace.
    \item \textbf{Mohou se prvky opakovat?} Pokud ano, jde o výběry s opakováním (viz dále).
  \end{enumerate}
\end{definitionbox}

\begin{example}
  \begin{itemize}
    \item Pět lidí se může postavit do řady $5! = 120$ způsoby.
    \item Na stupních vítězů z $10$ běžců (zlato, stříbro, bronz) může stát $V(3, 10) = 10\cdot 9\cdot 8 = 720$ různých trojic.
    \begin{center}
      \begin{tikzpicture}[scale=0.8]
        \draw[thick, fill=black!15] (0,0) rectangle (1.4,0.9);
        \draw[thick, fill=colorofwisdom!35] (1.4,0) rectangle (2.8,1.5);
        \draw[thick, fill=black!30] (2.8,0) rectangle (4.2,0.6);
        \node at (0.7,0.45) {\small $2.$};
        \node at (2.1,0.75) {\small $1.$};
        \node at (3.5,0.3) {\small $3.$};
        \node[above] at (2.1,1.5) {\small $10$ možností};
        \node[above] at (0.7,0.9) {\small $9$};
        \node[above] at (3.5,0.6) {\small $8$};
      \end{tikzpicture}
    \end{center}
    \item Tříčlenný výbor z $10$ lidí (bez funkcí) lze vybrat $\binom{10}{3} = \frac{720}{6} = 120$ způsoby. Je to $3! = 6$-krát méně než stupně vítězů, protože tady nerozlišujeme, kdo je `první'.
    \item Ve Sportce se vybírá $6$ čísel ze $49$. Různých tiketů je $\binom{49}{6} = 13\,983\,816$. Pokud si vsadíte jeden tiket týdně, vyhrajete první cenu v průměru jednou za zhruba $270\,000$ let.
  \end{itemize}
\end{example}

\begin{problem}
  \begin{enumerate}[label=(\alph*)]
    \item Kolika způsoby lze rozesadit $8$ hostů kolem kulatého stolu, pokud nás zajímá jen to, kdo sedí vedle koho (otočení stolu nic nemění)?
    \item Kolik je podání rukou, pokud si na večírku s $10$ lidmi podá ruku každý s každým?
    \item Kolik úhlopříček má konvexní $n$-úhelník? Odvoďte to kombinatoricky.
    \item Z $6$ mužů a $4$ žen máme vybrat tříčlennou komisi, ve které je alespoň jedna žena. Kolika způsoby to jde?
    \item Kolik je pěticiferných čísel s různými číslicemi, která jsou dělitelná pěti?
  \end{enumerate}
\end{problem}

\begin{hint}
  U (d) je jednodušší spočítat doplněk: všechny komise minus komise bez žen.
\end{hint}

\begin{solution}
  \begin{enumerate}[label=(\alph*)]
    \item $\frac{8!}{8} = 7! = 5\,040$. Každé rozesazení jsme započítali osmkrát (pro každé otočení). Na obrázku: když celý stůl pootočíte o jedno místo, každý má pořád stejné sousedy, takže jde o totéž rozesazení.
    \begin{center}
      \begin{tikzpicture}[scale=0.75]
        \draw[thick, fill=black!8] (0,0) circle (1);
        \foreach \i in {1,...,8} \node[draw, circle, fill=white, inner sep=1.5pt, font=\scriptsize] at ({90-45*(\i-1)}:1.45) {\i};
        \draw[->, colorofwisdom, thick] (60:0.55) arc (60:20:0.55);
        \node[font=\scriptsize, colorofwisdom] at (0,0) {otočení};
      \end{tikzpicture}
    \end{center}
    \item $\binom{10}{2} = 45$.
    \item Dvojic vrcholů je $\binom{n}{2}$, z toho $n$ jsou strany. Úhlopříček je $\binom{n}{2} - n = \frac{n(n - 1)}{2} - n = \frac{n(n - 3)}{2}$. Stejný výsledek jako v kapitole~\ref{kap:planimetrie}.
    \item $\binom{10}{3} - \binom{6}{3} = 120 - 20 = 100$.
    \item Končí nulou: $9\cdot 8\cdot 7\cdot 6 = 3\,024$. Končí pětkou: na prvním místě nesmí být nula ani pětka, tedy $8\cdot 8\cdot 7\cdot 6 = 2\,688$. Celkem $5\,712$.
  \end{enumerate}
\end{solution}

\section{Výběry s opakováním}\label{sec:vybery-s-opakovanim}

\begin{definitionbox}
  \textbf{Variace s opakováním} $k$-té třídy z $n$ prvků (uspořádané $k$-tice, prvky se mohou opakovat):
  \begin{equation*}
    V'(k, n) = n^k.
  \end{equation*}

  \textbf{Permutace s opakováním}: seřazení $n$ prvků, mezi kterými je $n_1$ stejných prvků prvního druhu, $n_2$ druhého, \ldots, $n_r$ druhu $r$ ($n_1 + \cdots + n_r = n$):
  \begin{equation*}
    P'(n_1, \ldots, n_r) = \frac{n!}{n_1!\,n_2!\cdots n_r!}.
  \end{equation*}

  \textbf{Kombinace s opakováním} $k$-té třídy z $n$ druhů (neuspořádané výběry, prvky se mohou opakovat):
  \begin{equation*}
    C'(k, n) = \binom{n + k - 1}{k}.
  \end{equation*}
\end{definitionbox}

\begin{example}
  \begin{itemize}
    \item PIN je variace s opakováním čtvrté třídy z $10$ číslic: $10^4$.
    \item Slovo MATEMATIKA má $10$ písmen: M dvakrát, A třikrát, T dvakrát a E, I, K po jednom. Různých přesmyček je
    \begin{equation*}
      \frac{10!}{2!\,3!\,2!} = \frac{3\,628\,800}{24} = 151\,200.
    \end{equation*}
    Kdybychom všechna písmena rozlišovali, bylo by jich $10!$. Prohozením obou M, tří A mezi sebou nebo obou T ale dostaneme stejné slovo, a proto dělíme.
    \item Ve zmrzlinárně mají $5$ druhů zmrzliny a chceme kornout se třemi kopečky (na pořadí kopečků nezáleží, druhy se mohou opakovat). Možností je $\binom{5 + 3 - 1}{3} = \binom{7}{3} = 35$.
  \end{itemize}
\end{example}

Odkud se bere vzorec pro kombinace s opakováním? Výběr $k$ kopečků z $n$ druhů zakódujeme jako řadu $k$ hvězdiček a $n - 1$ přepážek. Hvězdičky před první přepážkou jsou kopečky prvního druhu, mezi první a druhou kopečky druhého a tak dál. Například pro $3$ kopečky z $5$ druhů znamená $\star\star\,|\,|\,\star\,|\,|$ dva kopečky prvního druhu a jeden třetího. Každý výběr odpovídá právě jedné takové řadě a řad je tolik, kolika způsoby lze mezi $n + k - 1$ místy vybrat $k$ míst pro hvězdičky. Tomuto triku se v angličtině říká \textit{stars and bars}.
\begin{center}
  \begin{tikzpicture}[scale=0.75]
    \foreach \x/\t in {0/s, 1/s, 2/b, 3/b, 4/s, 5/b, 6/b} {
      \ifx\t s \node[colorofwisdom] at (\x,0) {\LARGE $\star$}; \else \draw[very thick] (\x,-0.45) -- (\x,0.45); \fi
    }
    \draw[decorate, decoration={brace, mirror, amplitude=4pt}] (-0.3,-0.6) -- (1.3,-0.6) node[midway, below=4pt, font=\scriptsize] {2$\times$ druh 1};
    \node[below, font=\scriptsize] at (2.5,-0.65) {druh 2: 0};
    \node[below, font=\scriptsize] at (4,-0.65) {1$\times$ druh 3};
    \node[below, font=\scriptsize] at (5.5,-0.65) {druhy 4, 5: 0};
  \end{tikzpicture}
\end{center}

\begin{problem}
  \begin{enumerate}[label=(\alph*)]
    \item Kolik je SPZ ve tvaru `číslice, písmeno, číslice, číslice, číslice, číslice, číslice', pokud máme $26$ písmen?
    \item Kolik různých slov (i nesmyslných) lze sestavit přeskládáním písmen slova BANANA?
    \item Kolika způsoby lze rozdělit $10$ stejných bonbonů mezi $4$ děti?
    \item Kolik je cest z levého dolního do pravého horního rohu mřížky $5\times 3$, pokud smíme chodit jen doprava a nahoru?
  \end{enumerate}
\end{problem}

\begin{solution}
  \begin{enumerate}[label=(\alph*)]
    \item $10\cdot 26\cdot 10^5 = 26\,000\,000$.
    \item $\frac{6!}{1!\,3!\,2!} = 60$.
    \item $\binom{4 + 10 - 1}{10} = \binom{13}{10} = \binom{13}{3} = 286$.
    \item Každá cesta je posloupnost $5$ kroků doprava a $3$ nahoru. Je to permutace s opakováním: $\frac{8!}{5!\,3!} = \binom{8}{3} = 56$. Červená cesta na obrázku odpovídá slovu PPNPNPPN (P = doprava, N = nahoru).
    \begin{center}
      \begin{tikzpicture}[scale=0.7]
        \draw[gray!50] (0,0) grid (5,3);
        \draw[ultra thick, colorofwisdom, ->] (0,0) -- (2,0) -- (2,1) -- (3,1) -- (3,2) -- (5,2) -- (5,3);
        \filldraw (0,0) circle (2pt) node[below left] {\small start};
        \filldraw (5,3) circle (2pt) node[above right] {\small cíl};
      \end{tikzpicture}
    \end{center}
  \end{enumerate}
\end{solution}

\section{Kombinační čísla a binomická věta}\label{sec:binomicka-veta}

\begin{definitionbox}
  \textbf{Vlastnosti kombinačních čísel}:
  \begin{enumerate}
    \item $\binom{n}{0} = \binom{n}{n} = 1$, $\binom{n}{1} = n$,
    \item $\binom{n}{k} = \binom{n}{n - k}$ \quad(souměrnost),
    \item $\binom{n}{k} + \binom{n}{k + 1} = \binom{n + 1}{k + 1}$ \quad(Pascalův vztah),
    \item $\binom{n}{0} + \binom{n}{1} + \cdots + \binom{n}{n} = 2^n$.
  \end{enumerate}

  \textbf{Binomická věta}: Pro všechna $a, b\in\mathbb{R}$ a $n\in\mathbb{N}$ platí
  \begin{equation*}
    {(a + b)}^n = \sum_{k=0}^{n}\binom{n}{k}a^{n-k}b^k = \binom{n}{0}a^n + \binom{n}{1}a^{n-1}b + \cdots + \binom{n}{n}b^n.
  \end{equation*}
\end{definitionbox}

Všechny vlastnosti se dají dokázat počítáním se vzorcem, ale mnohem hezčí jsou \textbf{kombinatorické důkazy}, kde stejnou věc spočítáme dvěma způsoby.

\begin{proof}[(vlastností 2--4 a binomické věty)]
  \textbf{Souměrnost}: Vybrat $k$ prvků, které \textit{vezmeme}, je totéž jako vybrat $n - k$ prvků, které \textit{nevezmeme}.

  \textbf{Pascalův vztah}: Kolika způsoby lze vybrat $k + 1$ prvků z $n + 1$? Označme jeden prvek jako `speciální'. Buď ho vybereme, a pak zbylých $k$ vybíráme z $n$ ostatních ($\binom{n}{k}$ způsobů), nebo ho nevybereme a všech $k + 1$ vybíráme z $n$ ostatních ($\binom{n}{k + 1}$ způsobů). Sečtením dostaneme $\binom{n + 1}{k + 1}$.

  \textbf{Součet řádku}: Levá strana počítá všechny podmnožiny $n$-prvkové množiny podle počtu prvků. Pravá strana počítá totéž jinak: u každého prvku se rozhodneme, zda ho vezmeme, nebo ne, což dává $2^n$ možností. Je to počet prvků potenční množiny z kapitoly 2.

  \textbf{Binomická věta}: Roznásobíme ${(a + b)}^n = (a + b)(a + b)\cdots(a + b)$. Z každé závorky vybereme buď $a$, nebo $b$. Člen $a^{n-k}b^k$ dostaneme, kdykoli vybereme $b$ z právě $k$ závorek. Těch závorek lze vybrat $\binom{n}{k}$ způsoby, a proto je to koeficient u $a^{n-k}b^k$.
\end{proof}

Pascalův vztah je přesně pravidlo, podle kterého se staví Pascalův trojúhelník z kapitoly~\ref{kap:vyrazy}: každé číslo je součtem dvou čísel nad ním. Čísla v $n$-tém řádku jsou $\binom{n}{0}, \binom{n}{1}, \ldots, \binom{n}{n}$. Tam jsme slíbili důkaz binomické věty a teď jsme ho dodali.
\begin{center}
  \begin{tikzpicture}[xscale=0.95, yscale=0.75]
    \foreach \n in {0,...,5} {
      \foreach \k in {0,...,\n} {
        \pgfmathtruncatemacro{\c}{round(factorial(\n)/(factorial(\k)*factorial(\n-\k)))}
        \node at ({\k - \n/2}, -\n) {\small $\c$};
      }
    }
    \draw[colorofwisdom, thick, rounded corners] (-0.85,-3.35) rectangle (0.85,-2.65);
    \draw[colorofwisdom, thick, ->] (-0.5,-3.35) -- (-0.1,-3.7);
    \draw[colorofwisdom, thick, ->] (0.5,-3.35) -- (0.1,-3.7);
    \draw[colorofwisdom, thick] (0,-4) circle (0.32);
    \node[right, font=\small, colorofwisdom] at (2.6,-3.4) {$\binom{3}{1} + \binom{3}{2} = \binom{4}{2}$, tedy $3 + 3 = 6$};
  \end{tikzpicture}
\end{center}

\begin{example}
  \begin{itemize}
    \item ${(x + 2)}^4 = x^4 + 4x^3\cdot 2 + 6x^2\cdot 4 + 4x\cdot 8 + 16 = x^4 + 8x^3 + 24x^2 + 32x + 16$.
    \item Koeficient u $x^3$ v ${(2x - 1)}^5$: člen je $\binom{5}{2}{(2x)}^3{(-1)}^2 = 10\cdot 8x^3 = 80x^3$.
    \item $1{,}01^{10} = {(1 + 0{,}01)}^{10}\approx 1 + 10\cdot 0{,}01 + 45\cdot 0{,}0001 = 1{,}1045$. Další členy jsou zanedbatelné (přesně $1{,}10462\ldots$). Takhle se počítalo před kalkulačkami.
  \end{itemize}
\end{example}

\begin{problem}
  \begin{enumerate}[label=(\alph*)]
    \item Vyřešte v $\mathbb{N}$: $\binom{n}{2} = 28$ a $\frac{(n + 1)!}{(n - 1)!} = 30$.
    \item Zjednodušte $\binom{n}{2} + \binom{n}{3}$.
    \item Určete člen, který neobsahuje $x$, v rozvoji ${\left(x + \frac{1}{x}\right)}^6$.
    \item Kolik podmnožin má množina $\{1, 2, \ldots, 10\}$? Kolik z nich má sudý počet prvků?
  \end{enumerate}
\end{problem}

\begin{solution}
  \begin{enumerate}[label=(\alph*)]
    \item $\frac{n(n - 1)}{2} = 28$, $n^2 - n - 56 = 0$, $n = 8$ (záporný kořen zahodíme). Dále $(n + 1)n = 30$, $n = 5$.
    \item Podle Pascalova vztahu $\binom{n + 1}{3}$.
    \item Obecný člen je $\binom{6}{k}x^{6-k}x^{-k} = \binom{6}{k}x^{6 - 2k}$. Bez $x$ pro $k = 3$: $\binom{6}{3} = 20$.
    \item $2^{10} = 1024$. Polovina z nich, tedy $512$, má sudý počet prvků. Z binomické věty pro $a = 1$, $b = -1$ totiž $0 = {(1 - 1)}^{10} = \binom{10}{0} - \binom{10}{1} + \binom{10}{2} - \cdots$, takže součet `sudých' kombinačních čísel je roven součtu `lichých'.
  \end{enumerate}
\end{solution}

\section{Odbočka --- Faktoriál z polovičky}

Faktoriál je definovaný jen pro přirozená čísla: $n!$ je počet způsobů, jak seřadit $n$ věcí. Co by mohlo znamenat $\left(\frac{1}{2}\right)!$? Seřadit půl věci? To nedává smysl. Ale matematici se nenechají odradit tím, že něco nedává smysl. V kapitole~\ref{kap:vyrazy} jsme takhle `nesmyslně' rozšířili mocniny na záporné a racionální exponenty, a pokaždé tak, aby dál platila stará pravidla. Pro faktoriál je to pravidlo $(n + 1)! = (n + 1)\cdot n!$.

Na obrázku vidíte body $[n, (n-1)!]$ pro $n = 1, 2, 3, 4$, tedy $1, 1, 2, 6$. Chceme jimi proložit hladkou křivku, která dál splňuje to rekurentní pravidlo. Takovou funkci v roce 1729 našel Leonhard Euler.

\begin{definitionbox}
  \textbf{Funkce gama}: Pro $x > 0$ definujeme
  \begin{equation*}
    \Gamma(x) = \int_0^{\infty} t^{x-1}e^{-t}\,\mathrm{d}t.
  \end{equation*}
  \textbf{Vlastnosti}:
  \begin{itemize}
    \item $\Gamma(x + 1) = x\,\Gamma(x)$ (stejné pravidlo jako u faktoriálu),
    \item $\Gamma(1) = 1$, a tedy $\Gamma(n + 1) = n!$ pro všechna $n\in\mathbb{N}_0$,
    \item $\Gamma\left(\frac{1}{2}\right) = \sqrt{\pi}$, a tedy `$\left(\frac{1}{2}\right)! = \Gamma\left(\frac{3}{2}\right) = \frac{\sqrt{\pi}}{2}\approx 0{,}886$'.
  \end{itemize}
\end{definitionbox}

\begin{proof}[(vztahu $\Gamma(x + 1) = x\Gamma(x)$)]
  Integrací per partes z kapitoly~\ref{kap:integraly} s $u = t^x$ a $v' = e^{-t}$:
  \begin{equation*}
    \Gamma(x + 1) = \int_0^{\infty}t^xe^{-t}\,\mathrm{d}t = {\big[-t^xe^{-t}\big]}_0^{\infty} + x\int_0^{\infty}t^{x-1}e^{-t}\,\mathrm{d}t = 0 + x\Gamma(x).
  \end{equation*}
  (Člen v hranaté závorce je nulový v nule i v nekonečnu, protože exponenciála $e^{-t}$ klesá rychleji, než jakákoli mocnina roste.) A $\Gamma(1) = \int_0^\infty e^{-t}\,\mathrm{d}t = 1$. Indukcí pak $\Gamma(n + 1) = n\Gamma(n) = \cdots = n!$.
\end{proof}

\begin{center}
  \begin{tikzpicture}[scale=0.85]
    \draw[->] (-2.3,0) -- (5,0) node[right] {\small $x$};
    \draw[->] (0,-5.4) -- (0,5.6) node[above] {\small $y$};
    \foreach \x in {-2,-1,1,2,3,4} \draw (\x,0.08) -- (\x,-0.08) node[below] {\tiny $\x$};
    \foreach \y in {-4,-2,2,4} \draw (0.08,\y) -- (-0.08,\y) node[left] {\tiny $\y$};
    \foreach \a in {0,-1,-2} \draw[dashed, gray] (\a,-5.3) -- (\a,5.4);
    \draw[thick, colorofwisdom] plot[smooth] coordinates {(0.200,4.591) (0.273,3.300) (0.347,2.572) (0.420,2.110) (0.493,1.796) (0.567,1.571) (0.640,1.404) (0.713,1.277) (0.787,1.180) (0.860,1.103) (0.933,1.043) (1.007,0.996) (1.080,0.960) (1.153,0.932) (1.227,0.912) (1.300,0.897) (1.373,0.889) (1.447,0.886) (1.520,0.887) (1.593,0.893) (1.667,0.903) (1.740,0.917) (1.813,0.935) (1.887,0.957) (1.960,0.984) (2.033,1.015) (2.107,1.050) (2.180,1.090) (2.253,1.135) (2.327,1.186) (2.400,1.242) (2.473,1.305) (2.547,1.374) (2.620,1.451) (2.693,1.537) (2.767,1.631) (2.840,1.734) (2.913,1.849) (2.987,1.976) (3.060,2.115) (3.133,2.270) (3.207,2.440) (3.280,2.629) (3.353,2.837) (3.427,3.068) (3.500,3.323) (3.573,3.607) (3.647,3.921) (3.720,4.269) (3.793,4.657) (3.867,5.088)};
    \draw[thick, colorofwisdom] plot[smooth] coordinates {(-0.770,-5.143) (-0.747,-4.800) (-0.725,-4.522) (-0.703,-4.296) (-0.680,-4.111) (-0.657,-3.961) (-0.635,-3.839) (-0.613,-3.741) (-0.590,-3.666) (-0.568,-3.610) (-0.545,-3.571) (-0.522,-3.550) (-0.500,-3.545) (-0.477,-3.556) (-0.455,-3.583) (-0.432,-3.627) (-0.410,-3.689) (-0.388,-3.771) (-0.365,-3.874) (-0.343,-4.002) (-0.320,-4.159) (-0.297,-4.350) (-0.275,-4.582) (-0.252,-4.866)};
    \draw[thick, colorofwisdom] plot[smooth] coordinates {(-1.883,4.839) (-1.860,4.181) (-1.837,3.715) (-1.815,3.370) (-1.792,3.108) (-1.770,2.905) (-1.748,2.747) (-1.725,2.622) (-1.702,2.523) (-1.680,2.447) (-1.657,2.390) (-1.635,2.348) (-1.613,2.320) (-1.590,2.306) (-1.567,2.303) (-1.545,2.312) (-1.522,2.332) (-1.500,2.363) (-1.478,2.407) (-1.455,2.463) (-1.433,2.532) (-1.410,2.616) (-1.388,2.718) (-1.365,2.838) (-1.343,2.981) (-1.320,3.151) (-1.297,3.353) (-1.275,3.594) (-1.252,3.885) (-1.230,4.241) (-1.208,4.681)};
    \foreach \x/\y in {1/1, 2/1, 3/2} \filldraw (\x,\y) circle (2.5pt);
    \filldraw (1.5,0.886) circle (2.5pt) node[below right, font=\tiny] {$\Gamma(\frac{3}{2}) = (\frac{1}{2})!$};
    \node[colorofwisdom, font=\small] at (3.4,4.9) {$y = \Gamma(x)$};
  \end{tikzpicture}
\end{center}

Černé body jsou obyčejné faktoriály $0! = 1! = 1$ a $2! = 2$ (bod $3! = 6$ v $x = 4$ už se do obrázku nevešel). Funkce gama mezi nimi vede hladce a vztah $\Gamma(x) = \frac{\Gamma(x + 1)}{x}$ ji rozšíří i na záporná čísla, kromě nuly a záporných celých čísel, kde má svislé asymptoty. Odtud mimochodem plyne další důvod, proč je $0! = 1$: je to prostě $\Gamma(1)$.

Že $\Gamma\left(\frac{1}{2}\right) = \sqrt{\pi}$, je jedna z těch matematických náhod, které nejsou náhodou. Po substituci $t = s^2$ se integrál změní na $\int_{-\infty}^{\infty}e^{-s^2}\,\mathrm{d}s$, tedy obsah pod Gaussovou křivkou, kterou uvidíme u normálního rozdělení v kapitole o pravděpodobnosti. A ten je roven $\sqrt{\pi}$. Kruh a pravděpodobnost se tu potkávají v jednom čísle.

\begin{note}
  Pro velká $n$ je faktoriál neprakticky velký, ale dá se dobře odhadnout \textbf{Stirlingovým vzorcem}
  \begin{equation*}
    n!\approx\sqrt{2\pi n}{\left(\frac{n}{e}\right)}^n.
  \end{equation*}
  Už pro $n = 10$ se liší jen o $0{,}8\,\%$ ($3\,598\,696$ místo $3\,628\,800$) a s rostoucím $n$ je relativní chyba čím dál menší. Ve vzorci jsou zase $\pi$ i $e$, přestože faktoriál je jen součin celých čísel.
\end{note}

\section{Odbočka --- Josephův problém}

Roku 67 našeho letopočtu obléhali Římané galilejské město Jotapata. Když město padlo, ukryl se židovský velitel Josephus se čtyřiceti svými muži v jeskyni. Muži se odmítli vzdát a rozhodli se raději zemřít. Podle Josephova vlastního vyprávění se domluvili, že se postaví do kruhu a budou se navzájem zabíjet podle losu, aby nikdo nemusel spáchat sebevraždu, kterou jejich víra zakazovala. Josephus a jeden další muž zůstali poslední, a místo aby splnili dohodu, vzdali se Římanům. Josephus pak přežil válku, stal se římským občanem a sepsal dějiny židovské války, ze kterých celý příběh známe.

Pozdější verze legendy (od matematiků, ne od historiků) z toho udělaly hádanku: $41$ lidí stojí v kruhu a počítá se dokola. Každý třetí je vyřazen, dokud nezůstane jediný. Na které místo se musí postavit Josephus, aby přežil? Odpověď je: na $31.$ místo, a jeho přítel na $16.$ místo, aby byli poslední dva.
\begin{center}
  \begin{tikzpicture}[scale=1.6]
    \foreach \i in {1,...,41} {
      \pgfmathsetmacro{\ang}{90-360*(\i-1)/41}
      \ifnum\i=31 \filldraw[colorofwisdom] (\ang:1) circle (2.2pt); \node[colorofwisdom, font=\scriptsize] at (\ang:1.22) {31};
      \else\ifnum\i=16 \filldraw[colorofwisdom] (\ang:1) circle (2.2pt); \node[colorofwisdom, font=\scriptsize] at (\ang:1.22) {16};
      \else \filldraw[black!40] (\ang:1) circle (1.2pt); \fi\fi
    }
    \node[font=\scriptsize] at (90:1.22) {1};
    \draw[->, gray] (80:0.75) arc (80:20:0.75);
    \node[font=\scriptsize, gray] at (0,0) {každý třetí ven};
  \end{tikzpicture}
\end{center}

Obecně, pokud vyřazujeme každého druhého z $n$ lidí, dá se poloha přeživšího spočítat přesně. Napíšeme $n = 2^m + l$, kde $2^m$ je největší mocnina dvojky nepřevyšující $n$ a $0\leq l < 2^m$. Přeživší stojí na místě $2l + 1$. Pro $41 = 32 + 9$ je to místo $19$. Zkuste si to ověřit na malém příkladu, například pro $n = 5$ (přežije třetí) nebo $n = 8$ (přežije první). Úloha se dnes používá jako klasické cvičení v programování, kde se jí říká \textit{Josephus problem}. Je to asi jediná programovací úloha, jejímž historickým základem je hromadná sebevražda v jeskyni.

\section{Shrnutí a otázky k ústní zkoušce}

\begin{problem}
  Odpovězte vlastními slovy:
  \begin{enumerate}
    \item Vysvětlete pravidlo součtu a součinu. Jak souvisí s množinami?
    \item Co je faktoriál? Proč je $0! = 1$?
    \item Co jsou variace, permutace a kombinace bez opakování? Odvoďte vzorce pro jejich počet.
    \item Jak poznáte, který typ výběru v úloze použít?
    \item Co jsou variace, permutace a kombinace s opakováním? Vysvětlete vzorec pro permutace s opakováním.
    \item Jaké vlastnosti mají kombinační čísla? Dokažte Pascalův vztah.
    \item Vyslovte a dokažte binomickou větu. Jak souvisí s Pascalovým trojúhelníkem?
    \item Kolik podmnožin má $n$-prvková množina a proč?
  \end{enumerate}
\end{problem}

\koneckapitoly

\chapter{Pravděpodobnost}
\label{kap:pravdepodobnost}
\mathquote{Probability theory is nothing but common sense reduced to calculation.}{Pierre-Simon Laplace}

Laplace tvrdí, že pravděpodobnost je jen zdravý rozum převedený na počty. Dobrá zpráva je, že má pravdu. Špatná zpráva je, že lidský zdravý rozum je v pravděpodobnosti překvapivě špatný. V této kapitole si ukážeme, jak pravděpodobnost počítat, a narazíme na několik příkladů, ve kterých intuice selhává tak spolehlivě, že kvůli nim vznikaly hádky mezi profesory matematiky.

\section{Náhodný pokus a jevy}

\begin{definitionbox}
  \begin{itemize}
    \item \textbf{Náhodný pokus} je děj, jehož výsledek není předem jednoznačně určen (hod kostkou, losování, měření).
    \item \textbf{Prostor výsledků} $\Omega$ je množina všech možných výsledků pokusu.
    \item \textbf{Jev} $A$ je libovolná podmnožina $\Omega$. Jev \textbf{nastal}, pokud výsledek pokusu patří do $A$.
    \item \textbf{Jistý jev} je $\Omega$, \textbf{nemožný jev} je $\emptyset$.
  \end{itemize}

  \textbf{Operace s jevy} jsou množinové operace z kapitoly 2:
  \begin{itemize}
    \item $A\cup B$: nastal $A$ \textbf{nebo} $B$ (sjednocení),
    \item $A\cap B$: nastal $A$ \textbf{a} $B$ (průnik),
    \item $A' = \Omega\setminus A$: \textbf{nenastal} $A$ (opačný jev),
    \item $A$ a $B$ jsou \textbf{neslučitelné}, pokud $A\cap B = \emptyset$ (nemohou nastat současně).
  \end{itemize}
\end{definitionbox}

\begin{example}
  Hod kostkou: $\Omega = \{1, 2, 3, 4, 5, 6\}$. Jev `padne sudé číslo' je $A = \{2, 4, 6\}$, jev `padne víc než $4$' je $B = \{5, 6\}$. Pak $A\cup B = \{2, 4, 5, 6\}$, $A\cap B = \{6\}$ a $A' = \{1, 3, 5\}$.
\end{example}

\section{Klasická pravděpodobnost}

\begin{definitionbox}
  \textbf{Klasická definice pravděpodobnosti}: Pokud má pokus konečně mnoho výsledků a všechny jsou \textbf{stejně pravděpodobné}, je pravděpodobnost jevu $A$
  \begin{equation*}
    P(A) = \frac{|A|}{|\Omega|} = \frac{\text{počet příznivých výsledků}}{\text{počet všech výsledků}}.
  \end{equation*}

  \textbf{Vlastnosti}:
  \begin{enumerate}
    \item $0\leq P(A)\leq 1$, $P(\emptyset) = 0$, $P(\Omega) = 1$,
    \item $P(A') = 1 - P(A)$,
    \item $P(A\cup B) = P(A) + P(B) - P(A\cap B)$, speciálně pro neslučitelné jevy $P(A\cup B) = P(A) + P(B)$.
  \end{enumerate}
\end{definitionbox}

Vlastnost 3 je stejný princip jako u počítání prvků sjednocení: když sečteme $|A|$ a $|B|$, prvky průniku započítáme dvakrát, a jednou je proto musíme odečíst. Vlastnost 2 je nejužitečnější trik celé pravděpodobnosti. Kdykoli úloha obsahuje slovo `alespoň', zkuste nejdřív spočítat opačný jev.

Podmínka stejně pravděpodobných výsledků je důležitá. Výroky jako `buď vyhraju v loterii, nebo ne, takže je to padesát na padesát' selhávají přesně na ní.

\begin{example}
  \begin{itemize}
    \item Hodíme dvěma kostkami. Výsledků je $6\cdot 6 = 36$ (uspořádané dvojice). Součet $7$ padne v $6$ případech ($1 + 6, 2 + 5, \ldots, 6 + 1$), tedy $P = \frac{6}{36} = \frac{1}{6}$. Součet $2$ jen v jednom případě, $P = \frac{1}{36}$. Proto má v deskové hře Osadníci z Katanu políčko se sedmičkou zvláštní pravidla.
    \item Z balíčku $52$ karet vytáhneme jednu. $P(\text{eso}) = \frac{4}{52} = \frac{1}{13}$. $P(\text{eso nebo srdce}) = \frac{4}{52} + \frac{13}{52} - \frac{1}{52} = \frac{16}{52} = \frac{4}{13}$, protože srdcové eso jsme započítali dvakrát.
    \item Pravděpodobnost výhry v prvním pořadí Sportky je $\frac{1}{\binom{49}{6}} = \frac{1}{13\,983\,816}$.
  \end{itemize}
\end{example}

\begin{example}
  Jaká je pravděpodobnost, že při čtyřech hodech kostkou padne alespoň jedna šestka? Opačný jev je `nepadne žádná šestka' s pravděpodobností ${\left(\frac{5}{6}\right)}^4$. Tedy
  \begin{equation*}
    P = 1 - {\left(\frac{5}{6}\right)}^4 = 1 - \frac{625}{1296}\approx 0{,}518.
  \end{equation*}
  Sázet na alespoň jednu šestku ve čtyřech hodech se tedy vyplatí. To je přesně sázka, kterou ve Francii v 17. století s oblibou nabízel hazardní hráč Chevalier de Méré a vydělával na ní. Viz odbočka.
\end{example}

\begin{problem}
  \begin{enumerate}[label=(\alph*)]
    \item Hodíme třemi mincemi. Jaká je pravděpodobnost, že padnou alespoň dva líce?
    \item Z osudí s $5$ bílými a $3$ černými koulemi vytáhneme najednou dvě. Jaká je pravděpodobnost, že budou obě bílé? A že bude každá jiné barvy?
    \item Jaká je pravděpodobnost, že ve Sportce uhodnete právě $3$ čísla?
    \item Číslo z množiny $\{1, 2, \ldots, 100\}$ vybereme náhodně. Jaká je pravděpodobnost, že je dělitelné dvěma nebo třemi?
  \end{enumerate}
\end{problem}

\begin{solution}
  \begin{enumerate}[label=(\alph*)]
    \item $|\Omega| = 8$, příznivé: LLR, LRL, RLL, LLL, tedy $\frac{4}{8} = \frac{1}{2}$.
    \item $\frac{\binom{5}{2}}{\binom{8}{2}} = \frac{10}{28} = \frac{5}{14}$. Každá jiné barvy: $\frac{5\cdot 3}{28} = \frac{15}{28}$.
    \item Tři ze šesti tažených a tři ze $43$ netažených: $\frac{\binom{6}{3}\binom{43}{3}}{\binom{49}{6}} = \frac{20\cdot 12341}{13983816}\approx 0{,}0177$, tedy asi $1{,}8\,\%$.
    \item Dělitelných dvěma je $50$, třemi $33$, šesti $16$. $P = \frac{50 + 33 - 16}{100} = \frac{67}{100}$.
  \end{enumerate}
\end{solution}

\section{Podmíněná pravděpodobnost a nezávislost}

\begin{definitionbox}
  \textbf{Podmíněná pravděpodobnost} jevu $A$ za podmínky, že nastal jev $B$ (s $P(B) > 0$), je
  \begin{equation*}
    P(A\,|\,B) = \frac{P(A\cap B)}{P(B)}.
  \end{equation*}
  Jevy $A$ a $B$ jsou \textbf{nezávislé}, pokud
  \begin{equation*}
    P(A\cap B) = P(A)\cdot P(B),
  \end{equation*}
  tedy pokud $P(A\,|\,B) = P(A)$: to, že nastal $B$, nijak nemění šanci na $A$.
\end{definitionbox}

Podmíněná pravděpodobnost znamená, že jsme dostali nějakou informaci a tím se zmenšil prostor výsledků. Už nepočítáme ze všech možností, jen z těch, ve kterých nastal $B$.

\begin{example}
  Hodíme kostkou a víme, že padlo sudé číslo. Jaká je pravděpodobnost, že padla šestka? $P(6\,|\,\text{sudé}) = \frac{\frac{1}{6}}{\frac{1}{2}} = \frac{1}{3}$. Jasné: ze tří sudých čísel je jedno šestka.
\end{example}

\begin{example}
  Rodina má dvě děti a víme, že alespoň jedno je kluk. Jaká je pravděpodobnost, že jsou to dva kluci? Intuice říká $\frac{1}{2}$, ale prostor je $\{KK, KD, DK, DD\}$ a informace vyřadí jen $DD$. Ze tří zbylých stejně pravděpodobných možností je příznivá jedna: $P = \frac{1}{3}$. Pokud bychom ale věděli, že \textit{starší} dítě je kluk, bylo by to $\frac{1}{2}$. Na tom, jak přesně informaci dostaneme, záleží.
\end{example}

\begin{note}
  Pozor na \textbf{gamblerův omyl}. Pokud na ruletě padla pětkrát po sobě červená, šance na černou v dalším kole není větší. Jednotlivá kola rulety jsou nezávislá a ruleta si nepamatuje, co padlo. V roce 1913 v kasinu v Monte Carlu padla černá šestadvacetkrát po sobě a hráči během té série prosázeli miliony franků, protože byli přesvědčení, že červená `už musí přijít'.
\end{note}

\subsection{Bayesova věta}

\begin{definitionbox}
  \textbf{Věta o úplné pravděpodobnosti}: Pokud jevy $B$ a $B'$ rozdělí prostor na dvě části, platí
  \begin{equation*}
    P(A) = P(A\,|\,B)P(B) + P(A\,|\,B')P(B').
  \end{equation*}
  \textbf{Bayesova věta}:
  \begin{equation*}
    P(B\,|\,A) = \frac{P(A\,|\,B)\cdot P(B)}{P(A)}.
  \end{equation*}
\end{definitionbox}

\begin{example}
  Nemoc má $1\,\%$ populace. Test ji odhalí u nemocného s pravděpodobností $99\,\%$, ale u zdravého člověka vyjde falešně pozitivní v $5\,\%$ případů. Váš test vyšel pozitivní. Jaká je pravděpodobnost, že jste nemocní?

  Označme $N$ = nemocný, $+$ = pozitivní test. $P(+) = 0{,}99\cdot 0{,}01 + 0{,}05\cdot 0{,}99 = 0{,}0099 + 0{,}0495 = 0{,}0594$. Podle Bayese
  \begin{equation*}
    P(N\,|\,+) = \frac{0{,}99\cdot 0{,}01}{0{,}0594} = \frac{0{,}0099}{0{,}0594} = \frac{1}{6}\approx 16{,}7\,\%.
  \end{equation*}
  Ne $99\,\%$, ale jen asi $17\,\%$! Zdravých lidí je tolik, že i malé procento falešných poplachů jich vyprodukuje víc než skutečně nemocných. Studie ukázaly, že velká část lékařů na takovou otázku odpoví špatně. Proto se pozitivní výsledky screeningových testů vždy ověřují dalším testem.
\end{example}

\begin{example}[ (Monty Hall)]
  V televizní soutěži jsou tři dveře. Za jedněmi je auto, za ostatními kozy. Vyberete si dveře, třeba číslo $1$. Moderátor, který ví, kde je auto, otevře jiné dveře, třeba číslo $3$, za kterými je koza. Pak se vás zeptá, jestli chcete změnit volbu na dveře číslo $2$. Vyplatí se to?

  Ano! Pokud volbu změníte, vyhrajete s pravděpodobností $\frac{2}{3}$. Vaše původní volba měla šanci $\frac{1}{3}$ a moderátor na tom nic nezměnil, protože koza za nějakými dalšími dveřmi je vždycky. Zbylé $\frac{2}{3}$ se ale `soustředily' na jediné zbývající dveře. Nevěříte? Představte si stejnou hru se sto dveřmi, ze kterých moderátor otevře $98$ s kozou. Nechali byste si svou původní volbu?

  Když tuto úlohu v roce 1990 v časopise vyřešila Marilyn vos Savant, dostala přes deset tisíc dopisů, z toho asi tisíc od lidí s doktorátem, že se mýlí. Přesvědčit se prý nenechal ani slavný matematik Paul Erdős, dokud neviděl počítačovou simulaci.
\end{example}

\section{Opakované pokusy}

\begin{definitionbox}
  \textbf{Bernoulliho schéma}: Pokus, ve kterém nastane `úspěch' s pravděpodobností $p$, opakujeme $n$-krát nezávisle. Pravděpodobnost, že úspěch nastane \textbf{právě $k$-krát}, je
  \begin{equation*}
    P_n(k) = \binom{n}{k}p^k{(1 - p)}^{n - k}.
  \end{equation*}
\end{definitionbox}

Odkud vzorec? Konkrétní pořadí s $k$ úspěchy a $n - k$ neúspěchy má pravděpodobnost $p^k{(1 - p)}^{n-k}$ (z nezávislosti). A takových pořadí je tolik, kolika způsoby lze vybrat $k$ pozic pro úspěchy z $n$, tedy $\binom{n}{k}$.

\begin{example}
  Test má $10$ otázek, každá má $4$ možnosti. Student nic neumí a tipuje. Jaká je pravděpodobnost, že odpoví správně na právě $5$ otázek?
  \begin{equation*}
    P_{10}(5) = \binom{10}{5}{\left(\frac{1}{4}\right)}^5{\left(\frac{3}{4}\right)}^5 = 252\cdot\frac{243}{1048576}\approx 0{,}058.
  \end{equation*}
  Alespoň jednu správnou odpověď má s pravděpodobností $1 - {\left(\frac{3}{4}\right)}^{10}\approx 0{,}944$. Aspoň něco.
\end{example}

\begin{definitionbox}
  \textbf{Geometrická pravděpodobnost}: Pokud vybíráme náhodně bod z oblasti $\Omega$ (úsečky, obrazce, tělesa) a každé místo je `stejně pravděpodobné', je pravděpodobnost, že bod padne do části $A$,
  \begin{equation*}
    P(A) = \frac{\text{míra } A}{\text{míra } \Omega},
  \end{equation*}
  kde mírou je délka, obsah nebo objem.
\end{definitionbox}

\begin{example}
  Do čtverce o straně $2$ náhodně hodíme šipku. Jaká je pravděpodobnost, že padne do vepsaného kruhu? $P = \frac{\pi\cdot 1^2}{4} = \frac{\pi}{4}\approx 0{,}785$. Takto se dá $\pi$ dokonce \textit{změřit}: hodíme velké množství náhodných bodů a spočítáme, kolik jich padlo do kruhu. Této metodě se říká \textbf{Monte Carlo} a dnes se používá ve fyzice, ve financích i při vývoji her.
\end{example}

\begin{problem}
  \begin{enumerate}[label=(\alph*)]
    \item Kolik lidí musí být v místnosti, aby pravděpodobnost, že dva z nich mají narozeniny ve stejný den, byla větší než $\frac{1}{2}$? Odhadněte nejdřív bez počítání.
    \item Střelec zasáhne terč s pravděpodobností $0{,}8$. Jaká je pravděpodobnost, že z $5$ výstřelů zasáhne alespoň čtyřikrát?
    \item Dva kamarádi se domluví, že se sejdou mezi 18:00 a 19:00. Každý přijde v náhodnou chvíli a čeká $15$ minut. Jaká je pravděpodobnost, že se potkají?
  \end{enumerate}
\end{problem}

\begin{hint}
  U (a) spočítejte pravděpodobnost opačného jevu: všichni mají narozeniny v různé dny. U (c) si nakreslete čtverec $60\times 60$, kde osa $x$ je příchod prvního a osa $y$ druhého.
\end{hint}

\begin{solution}
  \begin{enumerate}[label=(\alph*)]
    \item Pro $n$ lidí je pravděpodobnost různých narozenin $\frac{365\cdot 364\cdots(365 - n + 1)}{365^n}$. Pro $n = 23$ vyjde asi $0{,}493$, takže pravděpodobnost shody je asi $0{,}507$. Stačí $23$ lidí! Většina lidí tipuje kolem $180$. Tomu se říká \textbf{narozeninový paradox}. Ve třídě o $30$ žácích je pravděpodobnost shody už asi $70\,\%$.
    \item $P_5(4) + P_5(5) = 5\cdot 0{,}8^4\cdot 0{,}2 + 0{,}8^5 = 0{,}4096 + 0{,}32768\approx 0{,}737$.
    \item Potkají se, pokud $|x - y|\leq 15$. Mimo tento pás leží dva trojúhelníky s odvěsnami $45$, dohromady s obsahem $45^2 = 2025$. $P = 1 - \frac{2025}{3600} = \frac{1575}{3600} = \frac{7}{16}\approx 0{,}44$.
  \end{enumerate}
\end{solution}

\section{Odbočka --- Hráč a noc ohně}

Moderní teorie pravděpodobnosti začala hazardem. Chevalier de Méré, francouzský šlechtic a vášnivý hráč, roky vydělával na sázce, že ve čtyřech hodech kostkou padne alespoň jedna šestka (jak jsme spočítali, pravděpodobnost je asi $0{,}518$). Když ho to přestalo bavit, vymyslel si `vylepšenou' verzi: ve $24$ hodech dvěma kostkami padne alespoň jednou dvojitá šestka. Jeho úvaha: dvojitá šestka je šestkrát méně pravděpodobná, a proto potřebuji šestkrát víc hodů. Jenže začal prohrávat. Skutečná pravděpodobnost je $1 - {\left(\frac{35}{36}\right)}^{24}\approx 0{,}491$, tedy těsně pod polovinou.

Zmatený de Méré se v roce 1654 obrátil na Blaise Pascala. Pascal si o problému začal dopisovat s Pierrem de Fermatem a v jejich korespondenci se zrodila teorie pravděpodobnosti. Řešili mimo jiné, jak spravedlivě rozdělit sázku, když se hra musí přerušit před koncem.

Ve stejném roce, v noci 23. listopadu 1654, prožil Pascal svou mystickou zkušenost, o které jsme mluvili v kapitole~\ref{kap:vyrazy}. Trvala zhruba od půl jedenácté do půl jedné v noci a Pascal ji popsal na kousek pergamenu, který začíná slovem \textit{OHEŇ}. Ten pergamen si pak nechal zašít do podšívky kabátu a nosil ho u sebe každý den až do smrti. Celých osm let o něm nikdo nevěděl. Našel ho až sluha, který po Pascalově smrti procházel jeho oblečení. Pascal se po té noci matematice skoro přestal věnovat. Tak se stalo, že jeden z otců teorie pravděpodobnosti opustil hazard i vědu ve stejném roce, ve kterém je spojil.

Pascal ale ještě stihl jeden pravděpodobnostní argument, tentokrát o Bohu, takzvanou \textbf{Pascalovu sázku}. Pokud Bůh existuje a věříte v něj, získáte nekonečnou odměnu. Pokud neexistuje, ztratíte jen málo. Očekávaná hodnota víry je tedy nekonečná, takže se vyplatí věřit. Filozofové se o tento argument přou dodnes.

\section{Shrnutí a otázky k ústní zkoušce}

\begin{problem}
  Odpovězte vlastními slovy:
  \begin{enumerate}
    \item Co je náhodný pokus, prostor výsledků a jev? Jak souvisí operace s jevy s množinovými operacemi?
    \item Definujte klasickou pravděpodobnost. Za jakých předpokladů ji lze použít?
    \item Jaké vlastnosti má pravděpodobnost? Proč se u úloh se slovem `alespoň' vyplatí počítat opačný jev?
    \item Co je podmíněná pravděpodobnost? Kdy jsou jevy nezávislé?
    \item Vysvětlete Bayesovu větu na příkladu lékařského testu.
    \item Co je Bernoulliho schéma? Odvoďte vzorec pro pravděpodobnost $k$ úspěchů v $n$ pokusech.
    \item Co je geometrická pravděpodobnost?
    \item Vysvětlete narozeninový paradox a úlohu Montyho Halla.
  \end{enumerate}
\end{problem}

\koneckapitoly

\chapter{Statistika}
\label{kap:statistika}
\mathquote{There are three kinds of lies: lies, damned lies, and statistics.}{Mark Twain (který to připsal Benjaminu Disraelimu, který to nejspíš nikdy neřekl)}

Pravděpodobnost se ptá: když víme, jak funguje náhoda, co uvidíme? Statistika se ptá obráceně: když něco vidíme, co to říká o světě? V této kapitole se podíváme na \textbf{popisnou statistiku}, tedy na to, jak velké množství dat shrnout několika čísly a grafy. Je to ta část matematiky, kterou nejčastěji uvidíte ve zprávách, a také ta, se kterou se nejčastěji manipuluje.

\section{Základní pojmy}\label{sec:stat-zakladni-pojmy}

\begin{definitionbox}
  \begin{itemize}
    \item \textbf{Statistický soubor} je množina objektů, které zkoumáme (žáci třídy, obyvatelé města, výrobky z linky).
    \item \textbf{Statistická jednotka} je jeden prvek souboru. Počet jednotek je \textbf{rozsah souboru} $n$.
    \item \textbf{Statistický znak} je vlastnost, kterou u jednotek sledujeme. Může být \textbf{kvantitativní} (číselný: výška, známka, plat) nebo \textbf{kvalitativní} (slovní: barva očí, oblíbený předmět).
    \item \textbf{Absolutní četnost} $n_i$ hodnoty $x_i$ je počet jednotek, které mají tuto hodnotu. \textbf{Relativní četnost} je $\frac{n_i}{n}$ (často v procentech). \textbf{Kumulativní četnost} je součet četností všech hodnot menších nebo rovných $x_i$.
  \end{itemize}
\end{definitionbox}

Často nezkoumáme celý soubor (\textbf{základní soubor}, například všechny voliče), ale jen jeho část (\textbf{výběrový soubor}, například tisíc dotázaných v průzkumu). Aby výsledky něco znamenaly, musí být výběr \textbf{reprezentativní}. Slavný případ opaku: v roce 1936 americký časopis \textit{Literary Digest} obeslal přes deset milionů lidí a předpověděl drtivé vítězství republikána Landona nad Rooseveltem. Roosevelt vyhrál ve 46 státech ze 48. Časopis totiž vybíral adresy z telefonních seznamů a seznamů majitelů aut, a ty si uprostřed hospodářské krize mohli dovolit hlavně bohatší voliči. O dva roky později časopis zanikl.

\begin{example}\label{ex:znamky}
  Známky z písemky v malé skupině $10$ studentů: $1, 2, 2, 3, 1, 4, 2, 3, 5, 2$. Tabulka četností:
  \begin{center}
    \renewcommand{\arraystretch}{1.3}
    \begin{tabular}{|c|c|c|c|c|c|}
      \hline
      známka $x_i$ & $1$ & $2$ & $3$ & $4$ & $5$ \\
      \hline
      absolutní četnost $n_i$ & $2$ & $4$ & $2$ & $1$ & $1$ \\
      relativní četnost & $20\,\%$ & $40\,\%$ & $20\,\%$ & $10\,\%$ & $10\,\%$ \\
      kumulativní četnost & $2$ & $6$ & $8$ & $9$ & $10$ \\
      \hline
    \end{tabular}
  \end{center}
\end{example}

\section{Grafické znázornění}\label{sec:graficke-znazorneni}

\begin{definitionbox}
  \begin{itemize}
    \item \textbf{Sloupcový diagram} (histogram): výška sloupce odpovídá četnosti. Vhodný pro kvantitativní znaky.
    \item \textbf{Kruhový (koláčový) diagram}: velikost výseče odpovídá relativní četnosti. Úhel výseče je $\frac{n_i}{n}\cdot 360^\circ$. Vhodný pro podíly celku.
    \item \textbf{Spojnicový graf}: vhodný pro vývoj v čase.
    \item \textbf{Krabicový graf} (box plot): zobrazuje minimum, kvartily a maximum, viz dále.
  \end{itemize}
\end{definitionbox}

\begin{center}
  \begin{tikzpicture}[scale=0.8]
    \draw[->] (0,0) -- (6.2,0) node[right] {\small známka};
    \draw[->] (0,0) -- (0,4.8) node[above] {\small četnost};
    \foreach \y in {1,2,3,4} \draw (0.08,\y) -- (-0.08,\y) node[left] {\tiny $\y$};
    \foreach \x/\h in {1/2, 2/4, 3/2, 4/1, 5/1} {
      \fill[colorofwisdom!60] (\x-0.35,0) rectangle (\x+0.35,\h);
      \draw (\x-0.35,0) rectangle (\x+0.35,\h);
      \node[below] at (\x,0) {\small $\x$};
    }
  \end{tikzpicture}
  \hspace{1.5cm}
  \begin{tikzpicture}[scale=1.3]
    \fill[colorofwisdom!80] (0,0) -- (0:1) arc (0:72:1) -- cycle;
    \fill[colorofwisdom!40] (0,0) -- (72:1) arc (72:216:1) -- cycle;
    \fill[black!30] (0,0) -- (216:1) arc (216:288:1) -- cycle;
    \fill[black!60] (0,0) -- (288:1) arc (288:324:1) -- cycle;
    \fill[black!85] (0,0) -- (324:1) arc (324:360:1) -- cycle;
    \draw (0,0) circle (1);
    \node at (36:0.65) {\small $1$};
    \node at (144:0.65) {\small $2$};
    \node at (252:0.65) {\small $3$};
    \node[white] at (306:0.7) {\small $4$};
    \node[white] at (342:0.7) {\small $5$};
  \end{tikzpicture}
\end{center}

\begin{note}
  Grafy jsou nejsnadnější způsob, jak se statistikou lhát, aniž byste řekli jedinou nepravdu. Oblíbené triky: osa $y$, která nezačíná nulou (malý rozdíl pak vypadá obrovsky), trojrozměrné koláčové grafy (přední výseče vypadají větší) a vybírání časového úseku tak, aby trend vypadal, jak se nám hodí. Když vidíte graf ve zprávách, nejdřív se podívejte na osy.

  \begin{center}
    \begin{tikzpicture}[scale=0.8]
      \begin{scope}
        \draw[->] (0,0) -- (3.4,0);
        \draw[->] (0,0) -- (0,3.4);
        \foreach \y/\l in {1/20, 2/40} \draw (0.07,\y) -- (-0.07,\y) node[left] {\tiny $\l\,\%$};
        \filldraw[fill=gray!40, draw=black] (0.5,0) rectangle (1.4,2.5);
        \filldraw[fill=colorofwisdom!70, draw=black] (2,0) rectangle (2.9,2.6);
        \node[below, font=\scriptsize] at (0.95,0) {strana A};
        \node[below, font=\scriptsize] at (2.45,0) {strana B};
        \node[above, font=\scriptsize] at (0.95,2.5) {$50\,\%$};
        \node[above, font=\scriptsize] at (2.45,2.6) {$52\,\%$};
        \node[font=\small] at (1.7,-1) {osa od nuly};
      \end{scope}
      \begin{scope}[xshift=6cm]
        \draw[->] (0,0) -- (3.4,0);
        \draw[->] (0,0) -- (0,3.4);
        \draw (-0.12,0.15) -- (0.12,0.25);
        \draw (-0.12,0.25) -- (0.12,0.35);
        \foreach \y/\l in {0.9/50, 2.6/52} \draw (0.07,\y) -- (-0.07,\y) node[left] {\tiny $\l\,\%$};
        \filldraw[fill=gray!40, draw=black] (0.5,0) rectangle (1.4,0.9);
        \filldraw[fill=colorofwisdom!70, draw=black] (2,0) rectangle (2.9,2.6);
        \node[below, font=\scriptsize] at (0.95,0) {strana A};
        \node[below, font=\scriptsize] at (2.45,0) {strana B};
        \node[font=\small] at (1.7,-1) {osa od $49\,\%$};
      \end{scope}
    \end{tikzpicture}
  \end{center}
  Oba grafy ukazují stejná data. V levém je vidět, že strany jsou skoro vyrovnané. V pravém to vypadá, že strana B vede trojnásobně, protože osa začíná těsně pod $50\,\%$.
\end{note}

\section{Charakteristiky polohy}\label{sec:charakteristiky-polohy}

Charakteristiky polohy odpovídají na otázku `jaká je typická hodnota?'.

\begin{definitionbox}
  Pro hodnoty $x_1, x_2, \ldots, x_n$:
  \begin{itemize}
    \item \textbf{Aritmetický průměr}:
    \begin{equation*}
      \bar{x} = \frac{x_1 + x_2 + \cdots + x_n}{n} = \frac{1}{n}\sum_{i=1}^{n}x_i.
    \end{equation*}
    Pokud se hodnoty $x_i$ opakují s četnostmi $n_i$, je to \textbf{vážený průměr} $\bar{x} = \frac{\sum n_ix_i}{\sum n_i}$.
    \item \textbf{Medián} $\tilde{x}$: prostřední hodnota po seřazení dat podle velikosti. Pro sudé $n$ je to průměr dvou prostředních hodnot.
    \item \textbf{Modus} $\hat{x}$: nejčastější hodnota.
    \item \textbf{Geometrický průměr} (pro kladné hodnoty): $\bar{x}_G = \sqrt[n]{x_1x_2\cdots x_n}$.
    \item \textbf{Harmonický průměr}: $\bar{x}_H = \dfrac{n}{\frac{1}{x_1} + \cdots + \frac{1}{x_n}}$.
  \end{itemize}
\end{definitionbox}

\begin{example}
  V datech z příkladu~\ref{ex:znamky} je průměr $\bar{x} = \frac{1\cdot 2 + 2\cdot 4 + 3\cdot 2 + 4 + 5}{10} = \frac{25}{10} = 2{,}5$. Seřazená data: $1, 1, 2, 2, 2, 2, 3, 3, 4, 5$. Medián je průměr pátého a šestého prvku, tedy $\tilde{x} = 2$. Modus je $\hat{x} = 2$.

  Průměr má krásný fyzikální význam: je to \textbf{těžiště}. Představte si číselnou osu jako lehké prkno a na každou hodnotu položte jednu stejně těžkou kuličku. Prkno bude v rovnováze, právě když ho podepřete v průměru. Odchylky nalevo ($-1{,}5$ dvakrát a $-0{,}5$ čtyřikrát, celkem $-5$) se přesně vyváží s odchylkami napravo ($0{,}5$ dvakrát, $1{,}5$ a $2{,}5$, celkem $+5$).
  \begin{center}
    \begin{tikzpicture}[scale=1]
      \draw[thick] (0.3,0) -- (8.7,0);
      \foreach \x in {1,...,5} \draw (1.6*\x,0.08) -- (1.6*\x,-0.08) node[below=1pt] {\small $\x$};
      \foreach \x/\n in {1/2, 2/4, 3/2, 4/1, 5/1} {
        \foreach \k in {1,...,\n} {
          \filldraw[fill=colorofwisdom!60, draw=black] (1.6*\x,{0.28*\k-0.1}) circle (0.12);
        }
      }
      \fill[colorofwisdom] (4,-0.42) -- (3.8,-0.8) -- (4.2,-0.8) -- cycle;
      \node[right, font=\small, colorofwisdom] at (4.25,-0.65) {$\bar{x} = 2{,}5$};
      \draw[dashed, gray] (3.2,-0.3) -- (3.2,1.4) node[above, black, font=\small] {$\tilde{x} = 2$};
    \end{tikzpicture}
  \end{center}
  Medián naopak dělí kuličky na dvě stejně \textit{početné} poloviny, ať jsou krajní kuličky jakkoli daleko.
\end{example}

\begin{example}
  Proč potřebujeme medián? Firma má $9$ zaměstnanců s platem $30\,000$ Kč a ředitele s platem $300\,000$ Kč. Průměrný plat je
  \begin{equation*}
    \bar{x} = \frac{9\cdot 30000 + 300000}{10} = 57\,000\ \text{Kč}.
  \end{equation*}
  Ale nikdo kromě ředitele tolik nebere! Medián je $30\,000$ Kč a mnohem lépe popisuje `typického' zaměstnance. Průměr je citlivý na extrémní hodnoty, medián ne. Proto se ve zprávách o mzdách vždy vyplatí zjistit, zda jde o průměr, nebo o medián. Medián mezd v Česku bývá zhruba o šestinu nižší než průměr.

  \begin{center}
    \begin{tikzpicture}[scale=1]
      \draw[->] (0,0) -- (11,0) node[right] {\small tis. Kč};
      \foreach \x in {0,50,...,300} \draw ({0.035*\x},0.08) -- ({0.035*\x},-0.08) node[below] {\scriptsize $\x$};
      \foreach \k in {1,...,9} \filldraw[fill=gray!50, draw=black] (1.05,{0.24*\k-0.08}) circle (0.1);
      \filldraw[fill=colorofwisdom!70, draw=black] (10.5,0.16) circle (0.1);
      \node[above, font=\scriptsize] at (10.5,0.3) {ředitel};
      \draw[dashed, gray] (1.05,-0.5) -- (1.05,2.5) node[above, black, font=\small] {medián $30$};
      \fill[colorofwisdom] (1.995,-0.45) -- (1.82,-0.8) -- (2.17,-0.8) -- cycle;
      \node[right, font=\small, colorofwisdom] at (2.2,-0.65) {průměr $57$};
    \end{tikzpicture}
  \end{center}
  V řeči těžiště: jedna kulička daleko vpravo má dlouhé `rameno páky', takže přetáhne podpěru (průměr) doprava, přestože nalevo je kuliček devět. Medián se o vzdálenost nestará, jen o počet.
\end{example}

\begin{example}\label{ex:jine-prumery}
  Kdy použít jiné průměry?
  \begin{itemize}
    \item \textbf{Harmonický průměr} pro průměrnou rychlost. Pojedete do Brna rychlostí $60\ \si{km/h}$ a zpátky rychlostí $40\ \si{km/h}$. Průměrná rychlost není $50$, ale $\frac{2}{\frac{1}{60} + \frac{1}{40}} = 48\ \si{km/h}$, protože pomalejší cestou strávíte víc času. Pokud pochybujete, spočítejte to pro konkrétní vzdálenost, třeba $120\ \si{km}$: tam $2$ hodiny, zpátky $3$ hodiny, celkem $240\ \si{km}$ za $5$ hodin.
    \item \textbf{Geometrický průměr} pro průměrný růst. Akcie vzrostla první rok o $10\,\%$ a druhý o $60\,\%$. Průměrný roční růst není $35\,\%$, ale $\sqrt{1{,}1\cdot 1{,}6} - 1 = \sqrt{1{,}76} - 1\approx 32{,}7\,\%$. Kontrola: $1{,}327^2\approx 1{,}76$.
  \end{itemize}
\end{example}

\begin{note}
  Pro kladná čísla vždy platí $\bar{x}_H\leq\bar{x}_G\leq\bar{x}$, a rovnost nastává jen tehdy, když jsou všechna čísla stejná. Pro dvě čísla $a, b$ se nerovnost $\sqrt{ab}\leq\frac{a + b}{2}$ dokáže jednoduše: je ekvivalentní s $0\leq{(\sqrt{a} - \sqrt{b})}^2$. Geometricky to znamená, že ze všech obdélníků s daným obvodem má největší obsah čtverec, což jsme viděli v kapitole~\ref{kap:elementarni}.
\end{note}

\section{Charakteristiky variability}\label{sec:charakteristiky-variability}

Dva soubory mohou mít stejný průměr a přitom vypadat úplně jinak. Data $\{4, 5, 6\}$ a $\{0, 5, 10\}$ mají oba průměr $5$, ale ty druhé jsou mnohem víc `rozházené'. Charakteristiky variability měří, jak moc se data od průměru odchylují.

\begin{center}
  \begin{tikzpicture}[scale=0.85]
    \foreach \row/\vals/\s/\lab in {0/{4,5,6}/0.816/{$\{4, 5, 6\}$: $s\approx 0{,}82$}, -1.6/{0,5,10}/4.082/{$\{0, 5, 10\}$: $s\approx 4{,}08$}} {
      \fill[colorofwisdom!12] ({5-\s},\row-0.35) rectangle ({5+\s},\row+0.35);
      \draw (-0.3,\row) -- (10.3,\row);
      \foreach \x in {0,...,10} \draw (\x,\row+0.06) -- (\x,\row-0.06);
      \foreach \v in \vals {
        \ifnum\v=5\relax\else
          \draw[->, colorofwisdom, thick] (5,\row+0.2) -- (\v,\row+0.2);
        \fi
        \filldraw[fill=colorofwisdom!60, draw=black] (\v,\row) circle (0.12);
      }
      \node[anchor=west, font=\small] at (10.6,\row) {\lab};
    }
    \foreach \x in {0,5,10} \node[below, font=\scriptsize] at (\x,-1.7) {$\x$};
    \fill[colorofwisdom] (5,-2.15) -- (4.85,-2.45) -- (5.15,-2.45) -- cycle;
    \node[right, font=\scriptsize] at (5.15,-2.3) {$\bar{x} = 5$};
  \end{tikzpicture}
\end{center}
Šipky jsou odchylky od průměru, světlý pruh sahá od $\bar{x} - s$ do $\bar{x} + s$. Směrodatnou odchylku si můžete představit jako `typickou vzdálenost' hodnoty od průměru. U horních dat je to necelá jednička, u dolních přes čtyři.

\begin{definitionbox}
  \begin{itemize}
    \item \textbf{Variační rozpětí}: $R = x_{\max} - x_{\min}$.
    \item \textbf{Rozptyl} (průměrná druhá mocnina odchylky od průměru):
    \begin{equation*}
      s^2 = \frac{1}{n}\sum_{i=1}^{n}{(x_i - \bar{x})}^2 = \overline{x^2} - \bar{x}^2.
    \end{equation*}
    \item \textbf{Směrodatná odchylka}: $s = \sqrt{s^2}$. Má stejné jednotky jako data.
    \item \textbf{Variační koeficient}: $v = \frac{s}{\bar{x}}$ (relativní míra variability, obvykle v procentech).
  \end{itemize}
\end{definitionbox}

Proč druhé mocniny? Kdybychom průměrovali přímo odchylky $x_i - \bar{x}$, vyšla by vždy nula, protože kladné a záporné odchylky se přesně vyruší (zkuste si to dokázat). Druhá mocnina udělá ze všech odchylek kladná čísla. Směrodatná odchylka pak vrací výsledek do původních jednotek. Vzorec $s^2 = \overline{x^2} - \bar{x}^2$ (průměr druhých mocnin minus druhá mocnina průměru) se počítá rychleji a odvodí se roznásobením.
 Druhá mocnina má i geometrický smysl: nad každou odchylkou si nakreslete čtverec. Rozptyl je \textbf{průměrný obsah těchto čtverců} a směrodatná odchylka je strana čtverce s tímto průměrným obsahem. Vzdálené hodnoty mají obrovské čtverce, takže rozptyl na ně reaguje mnohem citlivěji než na hodnoty blízko průměru.

\begin{example}
  Pro známky z příkladu~\ref{ex:znamky} s průměrem $2{,}5$:
  \begin{align*}
    s^2 &= \frac{2\cdot{(1 - 2{,}5)}^2 + 4\cdot{(2 - 2{,}5)}^2 + 2\cdot{(3 - 2{,}5)}^2 + {(4 - 2{,}5)}^2 + {(5 - 2{,}5)}^2}{10}\\
    &= \frac{4{,}5 + 1 + 0{,}5 + 2{,}25 + 6{,}25}{10} = 1{,}45.
  \end{align*}
  Směrodatná odchylka $s = \sqrt{1{,}45}\approx 1{,}20$ a variační koeficient $v = \frac{1{,}20}{2{,}5}\approx 48\,\%$. Kontrola druhým vzorcem: $\overline{x^2} = \frac{2 + 16 + 18 + 16 + 25}{10} = 7{,}7$ a $7{,}7 - 6{,}25 = 1{,}45$.
\end{example}

\begin{note}
  Ve statistice výběrových souborů se rozptyl často počítá s $n - 1$ místo $n$ ve jmenovateli (tzv. výběrový rozptyl), protože to dává lepší odhad rozptylu celé populace. U maturity se obvykle používá verze s $n$, ale pokud vám kalkulačka nebo tabulkový procesor vyhodí jiné číslo, je to pravděpodobně tohle.
\end{note}

\begin{definitionbox}
  \textbf{Kvantily}: \textbf{Dolní kvartil} $Q_1$ je hodnota, pod kterou leží čtvrtina dat (medián dolní poloviny), \textbf{horní kvartil} $Q_3$ hodnota, pod kterou leží tři čtvrtiny dat (medián horní poloviny). Medián je druhý kvartil. \textbf{Mezikvartilové rozpětí} je $Q_3 - Q_1$.

  \textbf{Krabicový graf} znázorňuje minimum, $Q_1$, medián, $Q_3$ a maximum. Krabice obsahuje prostřední polovinu dat.
\end{definitionbox}

\begin{center}
  \begin{tikzpicture}[scale=1.5]
    \draw[->] (0.5,0) -- (5.5,0);
    \foreach \x in {1,2,3,4,5} \draw (\x,0.06) -- (\x,-0.06) node[below] {\small $\x$};
    \draw[thick] (1,0.6) -- (2,0.6);
    \draw[thick] (3,0.6) -- (5,0.6);
    \draw[thick, fill=colorofwisdom!30] (2,0.35) rectangle (3,0.85);
    \draw[very thick, colorofwisdom] (2,0.35) -- (2,0.85);
    \draw[thick] (1,0.45) -- (1,0.75);
    \draw[thick] (5,0.45) -- (5,0.75);
    \node[above] at (1,0.8) {\tiny min};
    \node[above] at (2,0.9) {\tiny $Q_1 = \tilde{x}$};
    \node[above] at (3,0.9) {\tiny $Q_3$};
    \node[above] at (5,0.8) {\tiny max};
  \end{tikzpicture}
\end{center}

Pro naše známky je dolní polovina $1, 1, 2, 2, 2$ s mediánem $Q_1 = 2$ a horní polovina $2, 3, 3, 4, 5$ s mediánem $Q_3 = 3$. Medián je také $2$, takže splývá s dolním kvartilem. To je u malých souborů s mnoha stejnými hodnotami běžné.

\begin{problem}
  Výšky (v cm) osmi hráčů basketbalového týmu jsou $185, 192, 178, 201, 188, 195, 183, 190$.
  \begin{enumerate}[label=(\alph*)]
    \item Určete průměr, medián a variační rozpětí.
    \item Určete směrodatnou odchylku.
    \item Do týmu přibude hráč s výškou $230$ cm. Jak se změní průměr a medián?
  \end{enumerate}
\end{problem}

\begin{solution}
  \begin{enumerate}[label=(\alph*)]
    \item Součet $1512$, průměr $\bar{x} = 189$. Seřazeno: $178, 183, 185, 188, 190, 192, 195, 201$, medián $\frac{188 + 190}{2} = 189$. Rozpětí $201 - 178 = 23$.
    \item Odchylky: $-11, -6, -4, -1, 1, 3, 6, 12$. Čtverce: $121, 36, 16, 1, 1, 9, 36, 144$, součet $364$. $s^2 = 45{,}5$, $s\approx 6{,}75\ \si{cm}$.
    \item Průměr $\frac{1742}{9}\approx 193{,}6$, vzroste skoro o pět centimetrů. Medián je teď pátá hodnota, tedy $190$, vzroste jen o centimetr. Opět vidíme, jak je průměr citlivý na extrémní hodnoty.
  \end{enumerate}
\end{solution}

\section{Závislost dvou znaků}\label{sec:zavislost-znaku}

Často nás zajímá, jestli spolu dva znaky souvisí. Mají vyšší lidé větší boty? Mají studenti, kteří víc spí, lepší známky? Dvojice hodnot $(x_i, y_i)$ můžeme zakreslit jako body do roviny (\textbf{bodový graf}) a podívat se, zda tvoří nějaký trend.

\begin{definitionbox}
  \textbf{Korelační koeficient} $r\in[-1, 1]$ měří míru \textbf{lineární} závislosti dvou znaků:
  \begin{equation*}
    r = \frac{\sum{(x_i - \bar{x})(y_i - \bar{y})}}{\sqrt{\sum{(x_i - \bar{x})}^2}\cdot\sqrt{\sum{(y_i - \bar{y})}^2}}.
  \end{equation*}
  \begin{itemize}
    \item $r$ blízko $1$: silná přímá závislost (s rostoucím $x$ roste $y$),
    \item $r$ blízko $-1$: silná nepřímá závislost,
    \item $r$ blízko $0$: žádná lineární závislost (ale může existovat jiná, třeba kvadratická).
  \end{itemize}
\end{definitionbox}

\begin{center}
  \begin{tikzpicture}[scale=0.8]
    \foreach \sx/\pts/\lab in {
      0/{0.2/0.28, 0.44/0.68, 0.67/0.74, 0.91/1.1, 1.15/1.31, 1.38/1.1, 1.62/1.26, 1.85/2.02, 2.09/1.8, 2.33/1.98, 2.56/2.69, 2.8/2.52}/{$r\approx 0{,}96$},
      3.7/{0.2/2.62, 0.44/2.26, 0.67/2.17, 0.91/1.74, 1.15/1.8, 1.38/1.75, 1.62/1.4, 1.85/1.33, 2.09/1.12, 2.33/0.63, 2.56/0.81, 2.8/0.55}/{$r\approx -0{,}98$},
      7.4/{0.2/2.05, 0.44/2.11, 0.67/1.37, 0.91/0.88, 1.15/0.3, 1.38/1.76, 1.62/1.33, 1.85/1.97, 2.09/1.12, 2.33/1.99, 2.56/0.9, 2.8/2.06}/{$r\approx 0$},
      11.1/{0.2/0.73, 0.44/1.17, 0.67/1.65, 0.91/2.05, 1.15/2.48, 1.38/2.47, 1.62/2.51, 1.85/2.32, 2.09/2.12, 2.33/1.71, 2.56/1.23, 2.8/0.66}/{$r\approx 0$}} {
      \begin{scope}[xshift=\sx cm]
        \draw[->] (0,0) -- (3.1,0) node[right] {\scriptsize $x$};
        \draw[->] (0,0) -- (0,3) node[above] {\scriptsize $y$};
        \foreach \x/\y in \pts \filldraw[colorofwisdom] (\x,\y) circle (2.2pt);
        \node[font=\small] at (1.5,-0.45) {\lab};
      \end{scope}
    }
    \node[font=\scriptsize] at (1.5,-0.95) {přímá závislost};
    \node[font=\scriptsize] at (5.2,-0.95) {nepřímá závislost};
    \node[font=\scriptsize] at (8.9,-0.95) {žádná závislost};
    \node[font=\scriptsize, align=center] at (12.6,-1.1) {silná, ale ne lineární!};
  \end{tikzpicture}
\end{center}
Korelační koeficient si můžete představit takhle. Bodovým grafem proložíme svislou a vodorovnou přímku průměry $\bar{x}$ a $\bar{y}$, čímž ho rozdělíme na čtyři kvadranty. Body vpravo nahoře a vlevo dole přispívají do čitatele kladně (obě odchylky mají stejné znaménko), body v ostatních dvou kvadrantech záporně. Pokud body tvoří stoupající pruh, převažují kladné příspěvky a $r > 0$. Na posledním obrázku se příspěvky z levé a pravé poloviny paraboly přesně vyruší, takže $r\approx 0$, přestože $y$ je na $x$ silně závislé. Proto je dobré se vždy nejdřív podívat na graf.

Pozorný čtenář si všimne, že vzorec je přesně kosinus úhlu mezi vektory odchylek $(x_1 - \bar{x}, \ldots, x_n - \bar{x})$ a $(y_1 - \bar{y}, \ldots, y_n - \bar{y})$ podle kapitoly~\ref{kap:analyticka}, jen v $n$-rozměrném prostoru. Proto je vždy mezi $-1$ a $1$.

\begin{note}
  \textbf{Korelace neznamená kauzalitu.} V létě roste prodej zmrzliny i počet utonulých. Zmrzlina ale k utonutí nevede. Obojí způsobuje teplé počasí. Na internetu najdete stránky se sbírkami absurdních korelací. Například počet filmů, ve kterých hrál Nicolas Cage, v letech 1999 až 2009 překvapivě dobře koreloval s počtem lidí, kteří utonuli v bazénu. Pokud někdo tvrdí, že `$A$ způsobuje $B$', protože spolu korelují, vždy se zeptejte, jestli za obojím nestojí nějaké $C$.
\end{note}

\section{Příběh --- Účty smrti a růže z Krymu}

Moderní statistika se narodila mezi mrtvými. V Londýně se od 16. století každý týden vydávaly takzvané \textit{Bills of Mortality}, tedy seznamy úmrtí podle farností a příčin. Pomáhaly sledovat, kdy se blíží další vlna moru. Mezi příčinami smrti v nich najdete například `zemřel leknutím', `planetou zasažen' nebo `šílenství'. Obchodník John Graunt je v roce 1662 poprvé systematicky zpracoval. Zjistil, že se rodí trochu víc chlapců než dívek, že úmrtnost dětí je děsivě vysoká a že čísla se dají použít k odhadu počtu obyvatel. Jeho kniha je považována za počátek demografie i statistiky vůbec.

O dvě století později se do Krymské války vydala ošetřovatelka Florence Nightingaleová. Ve vojenských nemocnicích zjistila, že vojáci umírají mnohem častěji na nemoci jako cholera nebo tyfus než na zranění z bitev. Aby o tom přesvědčila politiky, kteří čísla nečetli, vymyslela nový druh grafu, takzvaný \textit{coxcomb} (polární diagram, někdy nazývaný Nightingaleova růže). Každý měsíc tvořil jednu výseč a plocha barevných částí ukazovala, kolik vojáků zemřelo na nemoci, na zranění a z jiných příčin. Graf byl tak výmluvný, že vedl k reformě vojenských nemocnic a hygieny. Florence Nightingaleová se v roce 1858 stala první ženou zvolenou do britské Královské statistické společnosti. Lidé si ji pamatují jako `dámu s lampou'. Měli by si ji pamatovat i jako dámu s grafem.

\section{Shrnutí a otázky k ústní zkoušce}

\begin{problem}
  Odpovězte vlastními slovy:
  \begin{enumerate}
    \item Co je statistický soubor, jednotka a znak? Jaké druhy znaků znáte?
    \item Co je absolutní, relativní a kumulativní četnost?
    \item Jaké znáte způsoby grafického znázornění dat? Jak lze grafem manipulovat?
    \item Definujte aritmetický průměr, medián a modus. Kdy je lepší použít medián?
    \item Kdy se používá geometrický a harmonický průměr? Uveďte příklad.
    \item Co je rozptyl a směrodatná odchylka? Proč se v rozptylu odchylky umocňují?
    \item Co jsou kvartily a krabicový graf?
    \item Co měří korelační koeficient? Proč korelace neznamená kauzalitu?
  \end{enumerate}
\end{problem}

\koneckapitoly

\chapter{Komplexní čísla}
\label{kap:komplexni}
\mathquote{The shortest path between two truths in the real domain passes through the complex domain.}{Jacques Hadamard}

V kapitole 2 jsme slíbili, že se ke komplexním číslům vrátíme. V kapitole~\ref{kap:kvadraticke} jsme viděli, že rovnice $x^2 + 2x + 5 = 0$ nemá reálné řešení, a v příběhu o Cardanovi, že se odmocninám ze záporných čísel nedalo vyhnout ani při řešení rovnic, které reálné řešení mají. Komplexní čísla dlouho měla pověst něčeho podezřelého. Descartes je pohrdavě nazval \textit{imaginárními}, tedy pomyslnými, a to jméno jim už zůstalo. Dnes bez nich nefunguje elektrotechnika, kvantová fyzika ani zpracování signálů. Jsou totiž jediným rozumným způsobem, jak počítat s otáčením.

\section{Algebraický tvar}\label{sec:algebraicky-tvar}

\begin{definitionbox}
  \textbf{Imaginární jednotka} $i$ je číslo, pro které platí $i^2 = -1$.

  \textbf{Komplexní číslo} je číslo tvaru
  \begin{equation*}
    z = a + bi,\qquad a, b\in\mathbb{R}.
  \end{equation*}
  Číslo $a = \mathrm{Re}\,z$ je \textbf{reálná část}, $b = \mathrm{Im}\,z$ \textbf{imaginární část}. Množinu všech komplexních čísel značíme $\mathbb{C}$. Pro $b = 0$ dostáváme reálná čísla, takže $\mathbb{R}\subset\mathbb{C}$. Čísla s $a = 0$ se nazývají \textbf{ryze imaginární}.

  Dvě komplexní čísla se rovnají, právě když se rovnají jejich reálné i imaginární části.
\end{definitionbox}

V kapitole 2 jsme psali $i = \sqrt{-1}$. To je dobrá intuice, ale přesnější je definovat $i$ rovností $i^2 = -1$. Odmocnina ze záporného čísla totiž v $\mathbb{C}$ není jednoznačná (číslo $-i$ splňuje ${(-i)}^2 = -1$ úplně stejně) a pravidla jako $\sqrt{a}\sqrt{b} = \sqrt{ab}$ přestávají platit: $\sqrt{-1}\cdot\sqrt{-1} = i\cdot i = -1$, ale $\sqrt{(-1)(-1)} = \sqrt{1} = 1$.

\begin{definitionbox}
  \textbf{Operace}: Pro $z_1 = a + bi$ a $z_2 = c + di$:
  \begin{align*}
    z_1 + z_2 &= (a + c) + (b + d)i,\\
    z_1\cdot z_2 &= (ac - bd) + (ad + bc)i.
  \end{align*}
  Násobíme tedy jako mnohočleny a pak nahradíme $i^2 = -1$.

  \textbf{Komplexně sdružené číslo} k $z = a + bi$ je $\bar{z} = a - bi$. Platí $z\cdot\bar{z} = a^2 + b^2$ (reálné číslo!).

  \textbf{Dělení}: zlomek rozšíříme číslem komplexně sdruženým se jmenovatelem:
  \begin{equation*}
    \frac{z_1}{z_2} = \frac{z_1\cdot\bar{z}_2}{z_2\cdot\bar{z}_2} = \frac{z_1\cdot\bar{z}_2}{c^2 + d^2}.
  \end{equation*}
\end{definitionbox}

Dělení je přesně stejný trik jako usměrňování zlomků v kapitole~\ref{kap:vyrazy}. Tam jsme se zbavovali odmocniny ve jmenovateli pomocí $(a + \sqrt{b})(a - \sqrt{b}) = a^2 - b$, tady se zbavujeme $i$ pomocí $(c + di)(c - di) = c^2 + d^2$.

\begin{example}
  Pro $z_1 = 2 + 3i$ a $z_2 = 1 - i$:
  \begin{align*}
    z_1 + z_2 &= 3 + 2i,\\
    z_1\cdot z_2 &= 2 - 2i + 3i - 3i^2 = 2 + i + 3 = 5 + i,\\
    \frac{z_1}{z_2} &= \frac{(2 + 3i)(1 + i)}{(1 - i)(1 + i)} = \frac{2 + 2i + 3i + 3i^2}{1 + 1} = \frac{-1 + 5i}{2} = -\frac{1}{2} + \frac{5}{2}i.
  \end{align*}
\end{example}

\begin{note}
  Mocniny $i$ se opakují s periodou $4$:
  \begin{equation*}
    i^1 = i,\quad i^2 = -1,\quad i^3 = -i,\quad i^4 = 1,\quad i^5 = i,\ \ldots
  \end{equation*}
  Stačí tedy vydělit exponent čtyřmi a podívat se na zbytek. Například $i^{2026} = i^{4\cdot 506 + 2} = i^2 = -1$.
  Proč zrovna perioda $4$? Za chvíli v Gaussově rovině uvidíme, že násobení číslem $i$ znamená otočení o $90^\circ$. Každá další mocnina je tedy o čtvrt otáčky dál a po čtyřech čtvrtotáčkách jsme zpátky. Vpravo je totéž otáčení použité na číslo $z = 2 + i$: čísla $z, iz, i^2z, i^3z$ tvoří čtverec.
  \begin{center}
    \begin{tikzpicture}[scale=0.95]
      \begin{scope}
        \draw[->] (-1.7,0) -- (1.8,0) node[right] {\tiny Re};
        \draw[->] (0,-1.6) -- (0,1.7) node[above] {\tiny Im};
        \draw[dashed, gray] (0,0) circle (1.2);
        \foreach \a/\l/\p in {0/{$i^0 = 1$}/above left, 90/{$i$}/below right, 180/{$i^2 = -1$}/below right, 270/{$i^3 = -i$}/above left} {
          \filldraw[colorofwisdom] (\a:1.2) circle (2pt) node[\p, font=\scriptsize] {\l};
        }
        \foreach \a in {0,90,180,270} \draw[->, colorofwisdom, thick] (\a+12:1.45) arc (\a+12:\a+78:1.45);
        \node[font=\scriptsize, colorofwisdom] at (45:1.8) {$\cdot i$};
      \end{scope}
      \begin{scope}[xshift=5.5cm, scale=0.6]
        \draw[->] (-2.8,0) -- (2.9,0) node[right] {\tiny Re};
        \draw[->] (0,-2.7) -- (0,2.8) node[above] {\tiny Im};
        \draw[gray, dashed] (2,1) -- (-1,2) -- (-2,-1) -- (1,-2) -- cycle;
        \foreach \x/\y/\l/\p in {2/1/{$z = 2 + i$}/right, -1/2/{$iz = -1 + 2i$}/above left, -2/-1/{$i^2z$}/left, 1/-2/{$i^3z$}/below right} {
          \draw[colorofwisdom] (0,0) -- (\x,\y);
          \filldraw[colorofwisdom] (\x,\y) circle (3pt) node[\p, font=\scriptsize] {\l};
        }
      \end{scope}
    \end{tikzpicture}
  \end{center}
\end{note}

\begin{note}\label{note:komplexni-usporadani}
  Komplexní čísla tvoří \textbf{těleso} z kapitoly 2: platí všechny axiomy sčítání a násobení a každé nenulové číslo má inverzi $\frac{1}{z} = \frac{\bar{z}}{|z|^2}$. Na rozdíl od reálných čísel ale \textbf{nejsou uspořádaná}. Nedává smysl říct, že $i > 0$ nebo $i < 0$. Kdyby totiž $i > 0$, pak $i\cdot i > 0$, tedy $-1 > 0$, a kdyby $i < 0$, pak $(-i)(-i) > 0$, tedy zase $-1 > 0$. Nerovnice v komplexních číslech proto řešit nemůžeme.
\end{note}

\begin{problem}
  Vypočítejte:
  \begin{enumerate}[label=(\alph*)]
    \item $(3 - 2i)(4 + i)$,
    \item ${(1 + i)}^2$ a ${(1 + i)}^4$,
    \item $\dfrac{5}{2 - i}$,
    \item $\dfrac{1 + i}{1 - i}$,
    \item $i + i^2 + i^3 + \cdots + i^{100}$.
  \end{enumerate}
\end{problem}

\begin{solution}
  \begin{enumerate}[label=(\alph*)]
    \item $12 + 3i - 8i - 2i^2 = 14 - 5i$.
    \item ${(1 + i)}^2 = 1 + 2i - 1 = 2i$, ${(1 + i)}^4 = {(2i)}^2 = -4$.
    \item $\frac{5(2 + i)}{4 + 1} = 2 + i$.
    \item $\frac{{(1 + i)}^2}{2} = \frac{2i}{2} = i$.
    \item Každá čtveřice $i + (-1) + (-i) + 1 = 0$, a čtveřic je $25$. Součet je $0$.
  \end{enumerate}
\end{solution}

\section{Gaussova rovina}\label{sec:gaussova-rovina}

Komplexní číslo $a + bi$ je určeno dvojicí reálných čísel $(a, b)$. Můžeme ho tedy zakreslit jako bod v rovině. Vodorovná osa je \textbf{reálná}, svislá \textbf{imaginární}. Takové rovině se říká \textbf{Gaussova} (nebo komplexní) rovina.

\begin{definitionbox}
  \textbf{Absolutní hodnota} komplexního čísla $z = a + bi$ je jeho vzdálenost od počátku:
  \begin{equation*}
    |z| = \sqrt{a^2 + b^2} = \sqrt{z\bar{z}}.
  \end{equation*}
  \textbf{Argument} $\arg z$ je orientovaný úhel, který svírá spojnice počátku s bodem $z$ s kladnou reálnou poloosou.

  \vspace{0.2cm}

  \textbf{Geometrický význam operací}:
  \begin{itemize}
    \item sčítání komplexních čísel je sčítání vektorů z kapitoly~\ref{kap:analyticka},
    \item komplexně sdružené číslo je obraz v osové souměrnosti podle reálné osy,
    \item $|z_1 - z_2|$ je vzdálenost bodů $z_1$ a $z_2$.
  \end{itemize}
\end{definitionbox}

\begin{center}
  \begin{tikzpicture}[scale=1.1]
    \draw[->] (-1,0) -- (4.2,0) node[right] {\small Re};
    \draw[->] (0,-2.5) -- (0,3) node[above] {\small Im};
    \coordinate (Z) at (3,2);
    \draw[thick, colorofwisdom, ->] (0,0) -- (Z);
    \filldraw[colorofwisdom] (Z) circle (1.5pt) node[above right] {\small $z = a + bi$};
    \filldraw (3,-2) circle (1.5pt) node[below right] {\small $\bar{z} = a - bi$};
    \draw[dashed, gray] (3,-2) -- (Z);
    \draw[dashed] (Z) -- (3,0) node[below right] {\small $a$};
    \draw[dashed] (Z) -- (0,2) node[left] {\small $b$};
    \node[colorofwisdom, above left] at (1.5,1) {\small $|z|$};
    \draw[->] (0.6,0) arc (0:33.7:0.6);
    \node at (0.85,0.25) {\small $\varphi$};
  \end{tikzpicture}
\end{center}

Komplexní číslo si tedy můžeme představit dvojím způsobem: jako \textbf{bod} $[a, b]$, nebo jako \textbf{šipku} (vektor) z počátku do tohoto bodu. Sčítání je pak skládání šipek, přesně jako u sil ve fyzice. Šipku $z_2$ přiložíme na konec šipky $z_1$, a kde skončíme, tam je $z_1 + z_2$. Obě šipky a jejich součet tvoří rovnoběžník. Na obrázku jsou čísla z příkladu výše, $z_1 = 2 + 3i$ a $z_2 = 1 - i$.

\begin{center}
  \begin{tikzpicture}[scale=0.85]
    \draw[->] (-0.8,0) -- (4.3,0) node[right] {\small Re};
    \draw[->] (0,-1.6) -- (0,3.7) node[above] {\small Im};
    \foreach \x in {1,2,3} \draw (\x,0.06) -- (\x,-0.06);
    \foreach \y in {-1,1,2,3} \draw (0.06,\y) -- (-0.06,\y);
    \fill[colorofwisdom!8] (0,0) -- (2,3) -- (3,2) -- (1,-1) -- cycle;
    \draw[->, thick, colorofwisdom] (0,0) -- (2,3) node[above left, font=\small] {$z_1 = 2 + 3i$};
    \draw[->, thick] (0,0) -- (1,-1) node[below right, font=\small] {$z_2 = 1 - i$};
    \draw[->, dashed] (2,3) -- (3,2);
    \draw[->, dashed, colorofwisdom] (1,-1) -- (3,2);
    \draw[->, very thick, colorofwisdom!60!black] (0,0) -- (3,2) node[right, font=\small] {$z_1 + z_2 = 3 + 2i$};
  \end{tikzpicture}
  \hspace{1cm}
  \begin{tikzpicture}[scale=0.85]
    \draw[->] (-0.8,0) -- (4.3,0) node[right] {\small Re};
    \draw[->] (0,-1.6) -- (0,3.7) node[above] {\small Im};
    \draw[->, thick, colorofwisdom] (0,0) -- (2,3) node[above left, font=\small] {$z_1$};
    \draw[->, thick] (0,0) -- (1,-1) node[below right, font=\small] {$z_2$};
    \draw[very thick, colorofwisdom!60!black] (1,-1) -- (2,3) node[midway, right, font=\small] {$|z_1 - z_2| = \sqrt{17}$};
    \draw[dashed, gray] (0,0) circle (1.414);
    \node[gray, font=\scriptsize] at (-0.3,-1.65) {$|z| = \sqrt{2}$};
  \end{tikzpicture}
\end{center}
Vpravo je vzdálenost dvou čísel: $z_1 - z_2 = 1 + 4i$, a tedy $|z_1 - z_2| = \sqrt{1 + 16} = \sqrt{17}$. Šedá kružnice je množina všech čísel se stejnou absolutní hodnotou jako $z_2$, tedy $\{z : |z| = \sqrt{2}\}$. Rovnice $|z - z_0| = r$ obecně popisuje kružnici se středem $z_0$ a poloměrem $r$.

\section{Goniometrický tvar}\label{sec:goniometricky-tvar}

\begin{definitionbox}
  Každé nenulové komplexní číslo lze zapsat v \textbf{goniometrickém tvaru}
  \begin{equation*}
    z = |z|(\cos\varphi + i\sin\varphi),
  \end{equation*}
  kde $\varphi = \arg z$. Pro $z = a + bi$ platí
  \begin{equation*}
    \cos\varphi = \frac{a}{|z|},\qquad\sin\varphi = \frac{b}{|z|}.
  \end{equation*}

  \textbf{Násobení a dělení v goniometrickém tvaru}: Pro $z_1 = r_1(\cos\varphi_1 + i\sin\varphi_1)$ a $z_2 = r_2(\cos\varphi_2 + i\sin\varphi_2)$:
  \begin{align*}
    z_1\cdot z_2 &= r_1r_2\left(\cos(\varphi_1 + \varphi_2) + i\sin(\varphi_1 + \varphi_2)\right),\\
    \frac{z_1}{z_2} &= \frac{r_1}{r_2}\left(\cos(\varphi_1 - \varphi_2) + i\sin(\varphi_1 - \varphi_2)\right).
  \end{align*}
  Při násobení se tedy \textbf{absolutní hodnoty násobí} a \textbf{argumenty sčítají}.
\end{definitionbox}

\begin{proof}[(násobení)]
  Roznásobíme a použijeme součtové vzorce z kapitoly~\ref{kap:goniometrie}:
  \begin{align*}
    z_1z_2 &= r_1r_2\left(\cos\varphi_1\cos\varphi_2 + i\cos\varphi_1\sin\varphi_2 + i\sin\varphi_1\cos\varphi_2 + i^2\sin\varphi_1\sin\varphi_2\right)\\
    &= r_1r_2\left((\cos\varphi_1\cos\varphi_2 - \sin\varphi_1\sin\varphi_2) + i(\sin\varphi_1\cos\varphi_2 + \cos\varphi_1\sin\varphi_2)\right)\\
    &= r_1r_2\left(\cos(\varphi_1 + \varphi_2) + i\sin(\varphi_1 + \varphi_2)\right).
  \end{align*}
\end{proof}

Tohle je ten nejdůležitější poznatek celé kapitoly. Násobení komplexním číslem s absolutní hodnotou $1$ je \textbf{otočení} roviny. Násobení číslem $i$ (které má $|i| = 1$ a $\arg i = \frac{\pi}{2}$) otočí každý bod o $90^\circ$ proti směru hodinových ručiček. Dvojí otočení o $90^\circ$ je otočení o $180^\circ$, tedy násobení $-1$. A to je geometrický důvod, proč $i^2 = -1$! Nic mystického, jen dvě otočení za sebou.

\begin{center}
  \begin{tikzpicture}[scale=1.15]
    \draw[->] (-1.4,0) -- (2.6,0) node[right] {\small Re};
    \draw[->] (0,-0.3) -- (0,3.1) node[above] {\small Im};
    \coordinate (Z1) at (1,1);
    \coordinate (Z2) at (1,1.732);
    \coordinate (P) at (-0.732,2.732);
    \draw[->, thick, gray] (0,0) -- (Z1) node[right, font=\small, black] {$z_1 = 1 + i$};
    \draw[->, thick, black] (0,0) -- (Z2) node[right, font=\small] {$z_2 = 1 + \sqrt{3}\,i$};
    \draw[->, very thick, colorofwisdom] (0,0) -- (P) node[above left, font=\small] {$z_1z_2$};
    \draw[gray] (0.55,0) arc (0:45:0.55);
    \node[font=\scriptsize, gray] at (22:0.75) {$45^\circ$};
    \draw (60:0.95) arc (60:105:0.95);
    \node[font=\scriptsize] at (83:1.15) {$45^\circ$};
    \draw (1.25,0) arc (0:60:1.25);
    \node[font=\scriptsize] at (30:1.48) {$60^\circ$};
    \draw[->, colorofwisdom] (1.8,0) arc (0:105:1.8);
    \node[font=\scriptsize, colorofwisdom] at (12:2.1) {$105^\circ$};
    \node[anchor=west, align=left, font=\small] at (3,1.6) {$|z_1| = \sqrt{2}$, $\arg z_1 = 45^\circ$\\$|z_2| = 2$, $\arg z_2 = 60^\circ$\\[2pt]$|z_1z_2| = 2\sqrt{2}$, $\arg z_1z_2 = 105^\circ$\\[2pt]$z_1z_2 = (1 - \sqrt{3}) + (1 + \sqrt{3})i$};
  \end{tikzpicture}
\end{center}
Na obrázku násobíme $z_2 = 1 + \sqrt{3}\,i$ číslem $z_1 = 1 + i$. Šipka $z_2$ se otočí o $\arg z_1 = 45^\circ$ a natáhne se $|z_1| = \sqrt{2}$-krát. Zkuste si ověřit roznásobením, že $(1 + i)(1 + \sqrt{3}\,i)$ vyjde opravdu tak, jak je napsáno vpravo. Násobení komplexních čísel je tedy vždy kombinace dvou jednoduchých pohybů: \textbf{otočení a natažení}.

\begin{example}
  Převeďme do goniometrického tvaru:
  \begin{itemize}
    \item $z = 1 + i$: $|z| = \sqrt{2}$, $\cos\varphi = \sin\varphi = \frac{1}{\sqrt{2}}$, tedy $\varphi = \frac{\pi}{4}$ a $z = \sqrt{2}\left(\cos\frac{\pi}{4} + i\sin\frac{\pi}{4}\right)$.
    \item $z = -1 + i\sqrt{3}$: $|z| = 2$, $\cos\varphi = -\frac{1}{2}$, $\sin\varphi = \frac{\sqrt{3}}{2}$, tedy $\varphi = \frac{2\pi}{3}$ (druhý kvadrant!).
    \item $z = -5$: $|z| = 5$, $\varphi = \pi$, $z = 5(\cos\pi + i\sin\pi)$.
    \item $z = -3i$: $|z| = 3$, $\varphi = \frac{3\pi}{2}$.
  \end{itemize}
  Při určování argumentu se vždy podívejte, ve kterém kvadrantu číslo leží. Samotný tangens nestačí, protože $\tg\varphi = \frac{b}{a}$ je stejný pro $z$ i $-z$.
\end{example}

\begin{definitionbox}\label{def:moivre}
  \textbf{Moivreova věta}: Pro každé $n\in\mathbb{Z}$ platí
  \begin{equation*}
    {\left(r(\cos\varphi + i\sin\varphi)\right)}^n = r^n(\cos n\varphi + i\sin n\varphi).
  \end{equation*}
\end{definitionbox}

\begin{proof}
  Pro přirozená $n$ matematickou indukcí. Pro $n = 1$ je to zřejmé. Pokud tvrzení platí pro $n$, pak podle pravidla o násobení
  \begin{equation*}
    z^{n+1} = z^n\cdot z = r^n(\cos n\varphi + i\sin n\varphi)\cdot r(\cos\varphi + i\sin\varphi) = r^{n+1}\left(\cos(n + 1)\varphi + i\sin(n + 1)\varphi\right).
  \end{equation*}
  Pro záporná $n$ to plyne z pravidla pro dělení.
\end{proof}

Moivreova věta říká graficky, že mocniny $z, z^2, z^3, \ldots$ leží na \textbf{spirále}. Každé další vynásobení číslem $z$ otočí o stejný úhel $\varphi$ a natáhne $|z|$-krát. Pro $|z| > 1$ se spirála rozvíjí ven, pro $|z| < 1$ se stáčí do počátku a pro $|z| = 1$ body obíhají po jednotkové kružnici. Na obrázku jsou mocniny čísla $1 + i$, každá otočená o $45^\circ$ a $\sqrt{2}$-krát delší než předchozí.

\begin{center}
  \begin{tikzpicture}[scale=0.6]
    \draw[->] (-5,0) -- (2.6,0) node[right] {\small Re};
    \draw[->] (0,-4.6) -- (0,2.9) node[above] {\small Im};
    \draw[thick, colorofwisdom!50, domain=0:225, samples=80] plot ({pow(2,\x/90)*cos(\x)}, {pow(2,\x/90)*sin(\x)});
    \foreach \x/\y/\l/\p in {1/0/{$z^0 = 1$}/below right, 1/1/{$z = 1 + i$}/right, 0/2/{$z^2 = 2i$}/above right, -2/2/{$z^3 = -2 + 2i$}/above left, -4/0/{$z^4 = -4$}/above left, -4/-4/{$z^5 = -4 - 4i$}/below right} {
      \draw[gray] (0,0) -- (\x,\y);
      \filldraw[colorofwisdom] (\x,\y) circle (3.5pt) node[\p, font=\scriptsize] {\l};
    }
  \end{tikzpicture}
\end{center}

\begin{example}
  \begin{itemize}
    \item ${(1 + i)}^8 = {\left(\sqrt{2}\right)}^8\left(\cos 2\pi + i\sin 2\pi\right) = 16$. Zkuste to roznásobit osmkrát.
    \item ${(1 + i\sqrt{3})}^6 = 2^6(\cos 2\pi + i\sin 2\pi) = 64$, protože $1 + i\sqrt{3} = 2\left(\cos\frac{\pi}{3} + i\sin\frac{\pi}{3}\right)$.
    \item Moivreova věta dá zadarmo vzorce pro násobné úhly. Pro $n = 2$ a $r = 1$: ${(\cos\varphi + i\sin\varphi)}^2 = \cos^2\varphi - \sin^2\varphi + 2i\sin\varphi\cos\varphi$, a to se musí rovnat $\cos 2\varphi + i\sin 2\varphi$. Porovnáním reálných a imaginárních částí dostaneme oba vzorce pro dvojnásobný úhel z kapitoly~\ref{kap:goniometrie}.
  \end{itemize}
\end{example}

\begin{problem}
  \begin{enumerate}[label=(\alph*)]
    \item Převeďte do goniometrického tvaru $\sqrt{3} - i$ a $-2 - 2i$.
    \item Vypočítejte ${(\sqrt{3} - i)}^{12}$.
    \item Vypočítejte $\dfrac{{(1 + i)}^{10}}{{(1 - i)}^{8}}$.
  \end{enumerate}
\end{problem}

\begin{solution}
  \begin{enumerate}[label=(\alph*)]
    \item $\sqrt{3} - i = 2\left(\cos\frac{11\pi}{6} + i\sin\frac{11\pi}{6}\right)$, $-2 - 2i = 2\sqrt{2}\left(\cos\frac{5\pi}{4} + i\sin\frac{5\pi}{4}\right)$.
    \item $2^{12}\left(\cos 22\pi + i\sin 22\pi\right) = 4096$.
    \item $1 + i = \sqrt{2}\left(\cos\frac{\pi}{4} + i\sin\frac{\pi}{4}\right)$ a $1 - i = \sqrt{2}\left(\cos\left(-\frac{\pi}{4}\right) + i\sin\left(-\frac{\pi}{4}\right)\right)$. Podíl má absolutní hodnotu ${\left(\sqrt{2}\right)}^{10 - 8} = 2$ a argument $10\cdot\frac{\pi}{4} - 8\cdot\left(-\frac{\pi}{4}\right) = \frac{9\pi}{2}$, tedy výsledek je $2\left(\cos\frac{9\pi}{2} + i\sin\frac{9\pi}{2}\right) = 2i$.
  \end{enumerate}
\end{solution}

\section{Rovnice v komplexních číslech}

\subsection{Kvadratické rovnice}\label{sec:komplexni-kvadraticke}

\begin{definitionbox}
  Kvadratická rovnice $ax^2 + bx + c = 0$ s reálnými koeficienty a $D < 0$ má v $\mathbb{C}$ dva \textbf{komplexně sdružené} kořeny
  \begin{equation*}
    x_{1,2} = \frac{-b\pm i\sqrt{-D}}{2a}.
  \end{equation*}
  Vietovy vzorce i rozklad $a(x - x_1)(x - x_2)$ platí dál.
\end{definitionbox}

\begin{example}
  Rovnice $x^2 + 2x + 5 = 0$ z kapitoly~\ref{kap:kvadraticke}: $D = -16$, $\sqrt{-D} = 4$, tedy $x_{1,2} = \frac{-2\pm 4i}{2} = -1\pm 2i$. Kontrola Vietem: součet $-2$, součin ${(-1)}^2 + 2^2 = 5$. Sedí.

  \begin{center}
    \begin{tikzpicture}[scale=0.8]
      \begin{scope}
        \draw[->] (-3.4,0) -- (1.6,0) node[right] {\small $x$};
        \draw[->] (0,-0.4) -- (0,6) node[above] {\small $y$};
        \draw[thick, colorofwisdom, domain=-3.2:1.2, samples=60] plot (\x, {(\x*\x + 2*\x + 5)/1.6});
        \filldraw (-1,2.5) circle (2pt) node[below, font=\scriptsize] {$[-1, 4]$};
        \node[font=\small] at (-1,-1) {$y = x^2 + 2x + 5$};
      \end{scope}
      \begin{scope}[xshift=6.5cm, yshift=2.6cm]
        \draw[->] (-2.4,0) -- (1.6,0) node[right] {\small Re};
        \draw[->] (0,-2.6) -- (0,2.7) node[above] {\small Im};
        \draw[dashed, gray] (-1,2) -- (-1,-2);
        \filldraw[colorofwisdom] (-1,2) circle (2.5pt) node[left, font=\scriptsize] {$-1 + 2i$};
        \filldraw[colorofwisdom] (-1,-2) circle (2.5pt) node[left, font=\scriptsize] {$-1 - 2i$};
        \draw (-1,0.06) -- (-1,-0.06) node[below right, font=\scriptsize] {$-1$};
        \node[font=\small] at (-0.4,-3.6) {kořeny v Gaussově rovině};
      \end{scope}
    \end{tikzpicture}
  \end{center}
  Parabola osu $x$ neprotíná, takže reálné kořeny nejsou. Komplexní kořeny leží symetricky podle reálné osy (jsou komplexně sdružené) a jejich reálná část $-1$ je přesně $x$-ová souřadnice vrcholu paraboly. To platí obecně: $-\frac{b}{2a}$ je střed mezi kořeny, ať jsou reálné, nebo komplexní.
\end{example}

\begin{definitionbox}
  \textbf{Základní věta algebry}: Každý mnohočlen stupně $n\geq 1$ s komplexními koeficienty má v $\mathbb{C}$ alespoň jeden kořen. Důsledek: lze ho rozložit na součin $n$ lineárních činitelů, takže má (počítáno s násobností) právě $n$ kořenů.
\end{definitionbox}

Tuto větu dokázal (ne úplně kompletně) Carl Friedrich Gauss ve své disertaci v roce 1799, když mu bylo dvaadvacet. Je to důvod, proč u komplexních čísel nemusíme vymýšlet žádná další čísla. U reálných jsme museli přidat $i$, aby měla řešení rovnice $x^2 + 1 = 0$. U komplexních čísel už mají řešení všechny polynomiální rovnice. Komplexní čísla jsou v tomto smyslu `hotová'.

\subsection{Binomické rovnice}\label{sec:binomicke-rovnice}

\begin{definitionbox}
  \textbf{Binomická rovnice} $z^n = a$, kde $a = r(\cos\alpha + i\sin\alpha)\neq 0$, má v $\mathbb{C}$ právě $n$ řešení
  \begin{equation*}
    z_k = \sqrt[n]{r}\left(\cos\frac{\alpha + 2k\pi}{n} + i\sin\frac{\alpha + 2k\pi}{n}\right),\qquad k = 0, 1, \ldots, n - 1.
  \end{equation*}
  V Gaussově rovině tvoří tato řešení vrcholy \textbf{pravidelného $n$-úhelníku} se středem v počátku.
\end{definitionbox}

Proč $+2k\pi$? Podle Moivreovy věty má $z^n$ argument $n\varphi$. Ten se musí rovnat $\alpha$, ale jen až na násobek $2\pi$, protože úhly lišící se o $2\pi$ jsou stejné. Tedy $n\varphi = \alpha + 2k\pi$. Pro $k = 0, \ldots, n - 1$ dostaneme různé body, pro $k = n$ se vrátíme na začátek.

\begin{example}
  Rovnice $z^3 = 1$ má řešení $\cos\frac{2k\pi}{3} + i\sin\frac{2k\pi}{3}$ pro $k = 0, 1, 2$:
  \begin{equation*}
    z_0 = 1,\qquad z_1 = -\frac{1}{2} + \frac{\sqrt{3}}{2}i,\qquad z_2 = -\frac{1}{2} - \frac{\sqrt{3}}{2}i.
  \end{equation*}
  Tvoří rovnostranný trojúhelník vepsaný do jednotkové kružnice. Číslo $1$ má tedy tři třetí odmocniny, jen jedna z nich je reálná.

  Rovnice $z^4 = -16$: $r = 16$, $\alpha = \pi$, tedy $z_k = 2\left(\cos\frac{\pi + 2k\pi}{4} + i\sin\frac{\pi + 2k\pi}{4}\right)$. Úhly $\frac{\pi}{4}, \frac{3\pi}{4}, \frac{5\pi}{4}, \frac{7\pi}{4}$ dávají čtyři vrcholy čtverce $\pm\sqrt{2}\pm\sqrt{2}i$.
  \begin{center}
    \begin{tikzpicture}[scale=1.1]
      \begin{scope}
        \draw[->] (-1.5,0) -- (1.6,0) node[right] {\tiny Re};
        \draw[->] (0,-1.4) -- (0,1.5) node[above] {\tiny Im};
        \draw[dashed, gray] (0,0) circle (1);
        \draw[thick, colorofwisdom] (0:1) -- (120:1) -- (240:1) -- cycle;
        \foreach \a in {0,120,240} \filldraw[colorofwisdom] (\a:1) circle (1.5pt);
        \node at (0,-1.8) {\small $z^3 = 1$};
      \end{scope}
      \begin{scope}[xshift=5cm]
        \draw[->] (-1.8,0) -- (1.9,0) node[right] {\tiny Re};
        \draw[->] (0,-1.7) -- (0,1.8) node[above] {\tiny Im};
        \draw[dashed, gray] (0,0) circle (1.3);
        \draw[thick, colorofwisdom] (45:1.3) -- (135:1.3) -- (225:1.3) -- (315:1.3) -- cycle;
        \foreach \a in {45,135,225,315} \filldraw[colorofwisdom] (\a:1.3) circle (1.5pt);
        \node at (0,-1.8) {\small $z^4 = -16$};
      \end{scope}
    \end{tikzpicture}
  \end{center}
\end{example}

\begin{problem}
  Vyřešte v $\mathbb{C}$:
  \begin{enumerate}[label=(\alph*)]
    \item $x^2 - 6x + 13 = 0$,
    \item $x^2 + 9 = 0$,
    \item $z^2 = i$,
    \item $z^6 = 64$,
    \item $x^3 + 8 = 0$.
  \end{enumerate}
\end{problem}

\begin{solution}
  \begin{enumerate}[label=(\alph*)]
    \item $D = 36 - 52 = -16$, $x = 3\pm 2i$.
    \item $x = \pm 3i$.
    \item $i = \cos\frac{\pi}{2} + i\sin\frac{\pi}{2}$, řešení $\cos\frac{\pi}{4} + i\sin\frac{\pi}{4} = \frac{\sqrt{2}}{2}(1 + i)$ a $\cos\frac{5\pi}{4} + i\sin\frac{5\pi}{4} = -\frac{\sqrt{2}}{2}(1 + i)$.
    \item $z_k = 2\left(\cos\frac{k\pi}{3} + i\sin\frac{k\pi}{3}\right)$ pro $k = 0, \ldots, 5$: $\pm 2$, $\pm 1\pm\sqrt{3}i$. Pravidelný šestiúhelník.
    \item $x^3 + 8 = (x + 2)(x^2 - 2x + 4)$ podle vzorce z kapitoly~\ref{kap:vyrazy}. Kořeny $-2$ a $1\pm\sqrt{3}i$.
  \end{enumerate}
\end{solution}

\section{Příběh --- Muž, který předpověděl svou smrt}

Abraham de Moivre, autor Moivreovy věty, byl francouzský protestant. Když v roce 1685 Ludvík XIV. zrušil edikt nantský, který protestantům zaručoval náboženskou svobodu, byl de Moivre uvězněn a teprve po propuštění mohl uprchnout do Anglie. Tam strávil zbytek života. Navzdory tomu, že se přátelil s Newtonem a Halleym a byl členem Královské společnosti, nikdy nezískal univerzitní místo. Živil se doučováním matematiky, poradenstvím pro hazardní hráče a pojišťovny a tím, že v londýnské kavárně Slaughter's hrál o peníze šachy.

V pozdním věku si prý všiml, že každý den spí o čtvrt hodiny déle než předchozí den. Jako správný matematik z toho sestavil aritmetickou posloupnost z kapitoly~\ref{kap:posloupnosti} a spočítal, kdy dosáhne dvaceti čtyř hodin. Toho dne, 27. listopadu 1754, opravdu zemřel. Lékař jako příčinu smrti uvedl somnolenci, tedy chorobnou spavost. Je to pravděpodobně jediný případ v dějinách, kdy někdo předpověděl datum své smrti pomocí aritmetické posloupnosti a vyšlo mu to. Matematika je někdy až příliš přesná.

A ještě jedna věc, kterou nemůžeme vynechat. Pokud se v goniometrickém tvaru podíváme na $\cos\varphi + i\sin\varphi$, Euler v 18. století ukázal, že se to rovná $e^{i\varphi}$, kde $e$ je Eulerovo číslo z kapitoly~\ref{kap:exponencialni}. Moivreova věta je pak jen pravidlo ${(e^{i\varphi})}^n = e^{in\varphi}$. A pro $\varphi = \pi$ dostaneme
\begin{equation*}
  e^{i\pi} + 1 = 0,
\end{equation*}
rovnici, která v jednom řádku spojuje pět nejdůležitějších čísel matematiky: $0$, $1$, $\pi$, $e$ a $i$. Mnoho matematiků ji považuje za nejkrásnější rovnici, jaká kdy byla napsána. Fyzik Richard Feynman ji ve svém deníku ve čtrnácti letech nazval `nejpozoruhodnějším vzorcem v matematice'.

\section{Shrnutí a otázky k ústní zkoušce}

\begin{problem}
  Odpovězte vlastními slovy:
  \begin{enumerate}
    \item Proč potřebujeme komplexní čísla? Jak je definována imaginární jednotka?
    \item Co je algebraický tvar komplexního čísla? Jak se komplexní čísla sčítají, násobí a dělí?
    \item Co je komplexně sdružené číslo a absolutní hodnota? Jaký mají geometrický význam?
    \item Proč komplexní čísla nelze uspořádat?
    \item Co je Gaussova rovina? Jak v ní vypadá sčítání?
    \item Co je goniometrický tvar? Jak se v něm násobí a dělí? Dokažte to.
    \item Vyslovte a dokažte Moivreovu větu.
    \item Jak se řeší kvadratické rovnice se záporným diskriminantem? Co říká základní věta algebry?
    \item Jak se řeší binomické rovnice? Proč jejich řešení tvoří pravidelný mnohoúhelník?
  \end{enumerate}
\end{problem}

\koneckapitoly

\chapter{Limity a spojitost}
\label{kap:limity}
\mathquote{If I have seen further it is by standing on the shoulders of Giants.}{Isaac Newton}

Tři poslední kapitoly této knihy jsou o \textbf{matematické analýze}, tedy o limitách, derivacích a integrálech. Je to matematika, kterou v 17. století vymysleli Newton a Leibniz, aby mohli popsat pohyb, a která od té doby pohání skoro celou fyziku a techniku. Všechno začíná limitou. Limitu posloupnosti už známe z kapitoly~\ref{kap:posloupnosti}. Teď se podíváme na limitu funkce: co se děje s hodnotami $f(x)$, když se $x$ blíží k nějakému bodu?

\section{Limita funkce}

Podívejme se na funkci
\begin{equation*}
  f(x) = \frac{x^2 - 4}{x - 2}.
\end{equation*}
V bodě $x = 2$ není definovaná, protože bychom dělili nulou. Ale co se děje \textit{blízko} dvojky? Pro $x\neq 2$ můžeme zkrátit: $f(x) = \frac{(x - 2)(x + 2)}{x - 2} = x + 2$. Když se tedy $x$ blíží ke $2$, hodnoty $f(x)$ se blíží ke $4$. Graf je přímka $y = x + 2$, ze které chybí jeden bod. Říkáme, že limita funkce v bodě $2$ je $4$, přestože funkční hodnota v tom bodě vůbec neexistuje.

\begin{center}
  \begin{tikzpicture}[scale=0.9]
    \draw[->] (-1,0) -- (4,0) node[right] {\small $x$};
    \draw[->] (0,-0.5) -- (0,5.5) node[above] {\small $y$};
    \draw[thick, colorofwisdom] (-0.8,1.2) -- (3.4,5.4);
    \filldraw[fill=white, draw=colorofwisdom, thick] (2,4) circle (2.5pt);
    \draw[dashed] (2,0) node[below] {\small $2$} -- (2,3.9);
    \draw[dashed] (0,4) node[left] {\small $4$} -- (1.9,4);
  \end{tikzpicture}
\end{center}

\begin{definitionbox}
  \textbf{Okolí bodu} $a$ s poloměrem $\delta > 0$ je interval $(a - \delta, a + \delta)$. \textbf{Prstencové okolí} je totéž bez samotného bodu $a$.

  \vspace{0.2cm}

  Funkce $f$ má v bodě $a$ \textbf{limitu} $L$, pokud ke každému $\varepsilon > 0$ existuje $\delta > 0$ takové, že pro všechna $x$ z prstencového okolí $a$ s poloměrem $\delta$ (kde je $f$ definovaná) platí $|f(x) - L| < \varepsilon$. Píšeme
  \begin{equation*}
    \lim_{x\to a} f(x) = L.
  \end{equation*}
\end{definitionbox}

Slovy: hodnoty $f(x)$ budou libovolně blízko $L$ ($\varepsilon$), pokud budeme $x$ volit dost blízko $a$ ($\delta$). Je to stejná hra jako u posloupností, jen místo `od nějakého $n_0$ dál' říkáme `dost blízko bodu $a$'. Důležité je, že \textbf{hodnota v bodě $a$ samotném vůbec nehraje roli}. Proto je v definici prstencové okolí.

\begin{definitionbox}
  \textbf{Jednostranné limity}: Pokud se k $a$ blížíme jen zprava ($x > a$), mluvíme o \textbf{limitě zprava} $\lim_{x\to a^+} f(x)$, pokud jen zleva, o \textbf{limitě zleva} $\lim_{x\to a^-} f(x)$. Limita existuje, právě když existují obě jednostranné limity a jsou si rovny.

  \vspace{0.2cm}

  \textbf{Nevlastní limity a limity v nekonečnu}: $\lim_{x\to a} f(x) = \infty$ znamená, že hodnoty rostou nade všechny meze, když se $x$ blíží k $a$. $\lim_{x\to\infty} f(x) = L$ znamená, že se hodnoty blíží k $L$, když $x$ roste nade všechny meze.
\end{definitionbox}

\begin{example}
  Funkce $f(x) = \frac{1}{x}$:
  \begin{equation*}
    \lim_{x\to 0^+}\frac{1}{x} = \infty,\qquad\lim_{x\to 0^-}\frac{1}{x} = -\infty,\qquad\lim_{x\to\infty}\frac{1}{x} = \lim_{x\to-\infty}\frac{1}{x} = 0.
  \end{equation*}
  Limita v nule neexistuje, protože jednostranné limity jsou různé. Nekonečné limity odpovídají svislým asymptotám grafu, limity v nekonečnu vodorovným. To jsou přesně asymptoty hyperboly z kapitoly~\ref{kap:elementarni}.
\end{example}

\begin{example}
  Funkce $\mathrm{sgn}(x)$ (signum: $-1$ pro záporná, $0$ pro nulu, $1$ pro kladná $x$) má v nule limitu zprava $1$ a zleva $-1$. Limita neexistuje, i když jsou obě jednostranné limity konečné.
\end{example}

\section{Počítání limit}

\begin{definitionbox}
  \textbf{Pravidla}: Pokud $\lim_{x\to a} f(x) = A$ a $\lim_{x\to a} g(x) = B$ (konečné), pak
  \begin{equation*}
    \lim(f\pm g) = A\pm B,\qquad\lim(f\cdot g) = A\cdot B,\qquad\lim\frac{f}{g} = \frac{A}{B}\ (B\neq 0).
  \end{equation*}

  \textbf{Typické postupy}:
  \begin{enumerate}
    \item \textbf{Dosazení}: pro `hezké' funkce (spojité, viz dále) stačí dosadit: $\lim_{x\to a} f(x) = f(a)$.
    \item \textbf{Výraz typu $\frac{0}{0}$}: rozložíme na součin a zkrátíme, nebo rozšíříme výrazem s odmocninou.
    \item \textbf{Limita v nekonečnu u podílu mnohočlenů}: vydělíme nejvyšší mocninou jmenovatele (stejně jako u posloupností).
  \end{enumerate}
  Výrazy jako $\frac{0}{0}$, $\frac{\infty}{\infty}$, $\infty - \infty$ nebo $0\cdot\infty$ jsou \textbf{neurčité}. Neznamená to, že limita neexistuje, ale že z nich samotných ji nepoznáme a musíme výraz upravit.
\end{definitionbox}

\begin{example}
  \begin{align*}
    \lim_{x\to 3}(x^2 - 2x + 1) &= 9 - 6 + 1 = 4,\\
    \lim_{x\to 1}\frac{x^3 - 1}{x^2 - 1} &= \lim_{x\to 1}\frac{(x - 1)(x^2 + x + 1)}{(x - 1)(x + 1)} = \lim_{x\to 1}\frac{x^2 + x + 1}{x + 1} = \frac{3}{2},\\
    \lim_{x\to 0}\frac{\sqrt{x + 4} - 2}{x} &= \lim_{x\to 0}\frac{(x + 4) - 4}{x(\sqrt{x + 4} + 2)} = \lim_{x\to 0}\frac{1}{\sqrt{x + 4} + 2} = \frac{1}{4},\\
    \lim_{x\to\infty}\frac{3x^2 + 1}{x^2 - x} &= \lim_{x\to\infty}\frac{3 + \frac{1}{x^2}}{1 - \frac{1}{x}} = 3.
  \end{align*}
\end{example}

\subsection{Limita $\frac{\sin x}{x}$}

\begin{definitionbox}
  \textbf{Věta}:
  \begin{equation*}
    \lim_{x\to 0}\frac{\sin x}{x} = 1.
  \end{equation*}
\end{definitionbox}

\begin{proof}
  Pro $x\in\left(0, \frac{\pi}{2}\right)$ porovnáme tři obsahy na jednotkové kružnici: trojúhelník s vrcholy $O$, $[1, 0]$, $[\cos x, \sin x]$ (obsah $\frac{1}{2}\sin x$), kruhovou výseč s úhlem $x$ (obsah $\frac{1}{2}x$, kapitola~\ref{kap:planimetrie}) a pravoúhlý trojúhelník s odvěsnami $1$ a $\tg x$ (obsah $\frac{1}{2}\tg x$). Každý obrazec obsahuje ten předchozí, a tedy
  \begin{equation*}
    \sin x < x < \tg x.
  \end{equation*}
  Vydělíme $\sin x > 0$ a vezmeme převrácené hodnoty (které otočí nerovnosti):
  \begin{equation*}
    \cos x < \frac{\sin x}{x} < 1.
  \end{equation*}
  Pro $x\to 0^+$ jde $\cos x\to 1$, a výraz $\frac{\sin x}{x}$, sevřený mezi dvě funkce s limitou $1$, musí mít limitu $1$ také. Pro záporná $x$ to plyne ze sudosti funkce $\frac{\sin x}{x}$.
\end{proof}

Tomuto argumentu se říká \textbf{věta o dvou policajtech} (nebo o sevření): pokud je funkce mezi dvěma funkcemi se stejnou limitou, nezbývá jí nic jiného, než jít s nimi. Jako opilec mezi dvěma policisty, kteří ho vedou na stanici.\footnote{V angličtině se jí říká \textit{squeeze theorem} nebo, v Itálii, \textit{teorema dei carabinieri}. Opilec je mezinárodní.} Limita $\frac{\sin x}{x}$ je důvodem, proč pro malé úhly (v radiánech!) platí $\sin x\approx x$. Pro $x = 0{,}1$ je $\sin x\approx 0{,}0998$. Fyzici toho využívají u kyvadla.

\begin{example}
  $\lim_{x\to 0}\frac{\sin 3x}{x} = \lim_{x\to 0} 3\cdot\frac{\sin 3x}{3x} = 3\cdot 1 = 3$ (substituce $t = 3x$).
\end{example}

\begin{problem}
  Vypočítejte:
  \begin{enumerate}[label=(\alph*)]
    \item $\lim\limits_{x\to -2}\dfrac{x^2 + 5x + 6}{x + 2}$,
    \item $\lim\limits_{x\to 9}\dfrac{x - 9}{\sqrt{x} - 3}$,
    \item $\lim\limits_{x\to\infty}\dfrac{2x^3 - x}{5x^3 + 4}$,
    \item $\lim\limits_{x\to 0}\dfrac{\sin 5x}{\sin 2x}$,
    \item $\lim\limits_{x\to 2^+}\dfrac{1}{x - 2}$ a $\lim\limits_{x\to 2^-}\dfrac{1}{x - 2}$.
  \end{enumerate}
\end{problem}

\begin{solution}
  \begin{enumerate}[label=(\alph*)]
    \item $\frac{(x + 2)(x + 3)}{x + 2} = x + 3\to 1$.
    \item $\frac{(\sqrt{x} - 3)(\sqrt{x} + 3)}{\sqrt{x} - 3} = \sqrt{x} + 3\to 6$.
    \item $\frac{2}{5}$.
    \item $\frac{5\cdot\frac{\sin 5x}{5x}}{2\cdot\frac{\sin 2x}{2x}}\to\frac{5}{2}$.
    \item $\infty$ a $-\infty$.
  \end{enumerate}
\end{solution}

\section{Spojitost}

\begin{definitionbox}
  Funkce $f$ je \textbf{spojitá v bodě} $a$, pokud
  \begin{equation*}
    \lim_{x\to a} f(x) = f(a).
  \end{equation*}
  Musí tedy platit tři věci: $f(a)$ existuje, limita existuje a obě se rovnají. Funkce je \textbf{spojitá na intervalu}, pokud je spojitá v každém jeho bodě (v krajních bodech uzavřeného intervalu jen jednostranně).
\end{definitionbox}

Intuitivně: spojitou funkci na intervalu nakreslíte jedním tahem, aniž byste zvedli tužku. Všechny elementární funkce z předchozích kapitol (mnohočleny, lomené funkce, odmocniny, exponenciála, logaritmus, goniometrické funkce) jsou spojité ve všech bodech svého definičního oboru. Spojité jsou i jejich součty, součiny, podíly a složení. Proto u nich limitu spočítáme prostým dosazením.

\begin{definitionbox}
  \textbf{Druhy nespojitosti} v bodě $a$:
  \begin{itemize}
    \item \textbf{odstranitelná}: limita existuje, ale $f(a)$ neexistuje nebo je jiná (díra v grafu, jako u $\frac{x^2 - 4}{x - 2}$),
    \item \textbf{skok}: jednostranné limity existují, ale jsou různé (jako u signa),
    \item \textbf{nekonečná} (podstatná): alespoň jedna jednostranná limita je nekonečná nebo neexistuje (jako u $\frac{1}{x}$).
  \end{itemize}
\end{definitionbox}

\begin{example}
  Pro jaké $a$ je funkce
  \begin{equation*}
    f(x) = \begin{cases} x^2 & \text{pro } x\leq 1,\\ ax + 2 & \text{pro } x > 1\end{cases}
  \end{equation*}
  spojitá? Obě části jsou spojité, problém může být jen v bodě $1$. Limita zleva je $1$, limita zprava je $a + 2$ a $f(1) = 1$. Spojitost nastane pro $a + 2 = 1$, tedy $a = -1$.
\end{example}

\subsection{Věty o spojitých funkcích}

\begin{definitionbox}
  \textbf{Weierstrassova věta}: Funkce spojitá na \textbf{uzavřeném} intervalu $[a, b]$ je na něm omezená a nabývá na něm svého maxima i minima.

  \textbf{Bolzanova věta}: Pokud je funkce $f$ spojitá na $[a, b]$ a $f(a)$ a $f(b)$ mají opačná znaménka, existuje $c\in(a, b)$ takové, že $f(c) = 0$.
\end{definitionbox}

Obě věty zní samozřejmě a obě potřebují předpoklady, které nejsou samozřejmé. Funkce $\frac{1}{x}$ na otevřeném intervalu $(0, 1]$ maximum nemá (Weierstrass potřebuje uzavřený interval). Funkce $\frac{1}{x}$ má v $-1$ a $1$ opačná znaménka, ale nulu nikde nemá (Bolzano potřebuje spojitost na celém intervalu, a nula tam chybí).

Bolzanova věta má velice praktický důsledek: dá se podle ní hledat kořen rovnice, i když ji neumíme vyřešit. Této metodě se říká \textbf{metoda půlení intervalu} (bisekce).

\begin{example}
  Rovnice $x^3 + x - 1 = 0$ nemá celočíselný kořen a vzorec pro kubickou rovnici nechceme. Ale $f(0) = -1 < 0$ a $f(1) = 1 > 0$, takže podle Bolzanovy věty má kořen v $(0, 1)$. Půlíme:
  \begin{itemize}
    \item $f(0{,}5) = -0{,}375 < 0$, kořen je v $(0{,}5; 1)$,
    \item $f(0{,}75)\approx 0{,}172 > 0$, kořen je v $(0{,}5; 0{,}75)$,
    \item $f(0{,}625)\approx -0{,}131 < 0$, kořen je v $(0{,}625; 0{,}75)$,
    \item a tak dál.
  \end{itemize}
  Každým krokem se interval zmenší na polovinu. Po deseti krocích známe kořen na tři desetinná místa: $x\approx 0{,}682$. Takhle (jen chytřeji) počítá kalkulačka i počítač.
\end{example}

\begin{note}
  Bolzanova věta má i jeden zábavný důsledek, kterému se říká \textit{věta o mnichovi}. Mnich vyjde v 6:00 ráno z kláštera na vrchol hory a dorazí tam v 18:00. Druhý den ve stejný čas sejde stejnou cestou dolů (jakkoli rychle, klidně se cestou zastaví). Existuje místo, kde byl v obou dnech ve stejnou hodinu? Ano! Představte si oba dny současně jako dva mnichy, jednoho jdoucího nahoru a druhého dolů. Někde se nutně potkají. Formálně: rozdíl jejich výšek je spojitá funkce času, na začátku je záporný a na konci kladný.
\end{note}

\begin{problem}
  \begin{enumerate}[label=(\alph*)]
    \item Určete body nespojitosti funkce $f(x) = \frac{x^2 - 1}{x^2 - 3x + 2}$ a jejich druh.
    \item Dokažte, že rovnice $\cos x = x$ má řešení.
    \item Dokažte, že každý mnohočlen lichého stupně má reálný kořen.
  \end{enumerate}
\end{problem}

\begin{solution}
  \begin{enumerate}[label=(\alph*)]
    \item $\frac{(x - 1)(x + 1)}{(x - 1)(x - 2)}$. V bodě $1$ je odstranitelná nespojitost (limita je $\frac{2}{-1} = -2$), v bodě $2$ nekonečná nespojitost.
    \item $g(x) = \cos x - x$ je spojitá, $g(0) = 1 > 0$, $g\left(\frac{\pi}{2}\right) = -\frac{\pi}{2} < 0$. Podle Bolzana má kořen v $\left(0, \frac{\pi}{2}\right)$. (Je přibližně $0{,}739$. Když na kalkulačce v radiánech opakovaně mačkáte tlačítko $\cos$, skončíte přesně u něj.)
    \item Pro mnohočlen lichého stupně s kladným vedoucím koeficientem je $\lim_{x\to\infty} P(x) = \infty$ a $\lim_{x\to-\infty} P(x) = -\infty$. Existují tedy body, kde je záporný a kde je kladný, a mezi nimi podle Bolzana leží kořen. (Pro záporný vedoucí koeficient obráceně.)
  \end{enumerate}
\end{solution}

\section{Odbočka --- Strážce mince}

Isaac Newton byl génius, ale ne příjemný člověk. Citát na začátku kapitoly napsal v roce 1675 v dopise Robertu Hookovi, se kterým se celý život nesnášel. Někteří historici ho čtou jako skromnou poklonu předchůdcům, jiní jako skrytou urážku: Hooke byl prý malé postavy a shrbený, takže těžko mohl být `obrem'. Po Hookově smrti zmizel jediný jeho známý portrét. Kdo za to může, se neví.

Newton strávil desítky let alchymií a pokusy o výrobu kamene mudrců a napsal o ní víc než o fyzice. Analýzy jeho vlasů ve 20. století ukázaly vysoký obsah rtuti, což by vysvětlovalo jeho období paranoie a nespavosti kolem roku 1693.

V roce 1696 se stal správcem a později ředitelem Královské mincovny. Většina lidí by tu funkci brala jako čestnou synekuru. Newton ne. Osobně vyšetřoval padělatele mincí, převlečený navštěvoval londýnské hospody, aby sbíral svědectví, a vyslýchal podezřelé. Padělání tehdy bylo velezradou, za kterou se věšelo. Jeho nejslavnějším protivníkem byl William Chaloner, padělatel, který Newtona veřejně obviňoval z neschopnosti. Newton ho s neúprosnou pečlivostí usvědčil a Chaloner byl v roce 1699 oběšen, přestože Newtonovi z vězení psal zoufalé dopisy s prosbou o milost. Newton na ně neodpověděl.

\section{Shrnutí a otázky k ústní zkoušce}

\begin{problem}
  Odpovězte vlastními slovy:
  \begin{enumerate}
    \item Co je limita funkce v bodě? Vysvětlete definici. Proč nezáleží na hodnotě funkce v samotném bodě?
    \item Co jsou jednostranné limity? Kdy limita existuje?
    \item Co jsou nevlastní limity a limity v nekonečnu? Jak souvisí s asymptotami?
    \item Jak se počítají limity typu $\frac{0}{0}$? Co jsou neurčité výrazy?
    \item Dokažte, že $\lim_{x\to 0}\frac{\sin x}{x} = 1$.
    \item Co je spojitost funkce v bodě a na intervalu? Jaké znáte druhy nespojitosti?
    \item Vyslovte Weierstrassovu a Bolzanovu větu. Proč jsou důležité jejich předpoklady?
    \item Jak se dá pomocí Bolzanovy věty přibližně najít kořen rovnice?
  \end{enumerate}
\end{problem}

\koneckapitoly

\chapter{Derivace}
\label{kap:derivace}
\mathquote{It is unworthy of excellent men to lose hours like slaves in the labour of calculation which could safely be relegated to anyone else if machines were used.}{Gottfried Wilhelm Leibniz}

Derivace měří, jak rychle se něco mění. Rychlost je derivace dráhy podle času, zrychlení je derivace rychlosti. Sklon tečny ke grafu je derivace funkce. A když chceme najít maximum nebo minimum čehokoli, od zisku firmy po tvar plechovky, hledáme místa, kde je derivace nulová. Je to pravděpodobně nejužitečnější nástroj, který středoškolská matematika nabízí. Pokud vám z celé analýzy zůstane jen jedna věc, ať je to tahle.

\section{Pojem derivace}

Představte si auto, které ujede za čas $t$ dráhu $s(t)$. Průměrná rychlost mezi časy $t_0$ a $t_0 + h$ je
\begin{equation*}
  \frac{s(t_0 + h) - s(t_0)}{h}.
\end{equation*}
Okamžitou rychlost (tu, kterou ukazuje tachometr) dostaneme, když necháme $h$ jít k nule. Geometricky je průměrná rychlost směrnice \textbf{sečny} grafu procházející body $[t_0, s(t_0)]$ a $[t_0 + h, s(t_0 + h)]$. Když se druhý bod blíží k prvnímu, sečna se změní v \textbf{tečnu}.

\begin{center}
  \begin{tikzpicture}[scale=1]
    \draw[->] (-0.3,0) -- (4.6,0) node[right] {\small $x$};
    \draw[->] (0,-0.3) -- (0,4.3) node[above] {\small $y$};
    \draw[thick, colorofwisdom, domain=0.2:4.2, samples=60] plot (\x, {0.22*\x*\x + 0.3});
    \coordinate (A) at (1.5,0.795); \coordinate (B) at (3.5,2.995);
    \filldraw (A) circle (1.5pt); \filldraw (B) circle (1.5pt);
    \draw[gray] ($(A)!-0.4!(B)$) -- ($(A)!1.3!(B)$) node[right] {\small sečna};
    \draw[thick] (0.3,0.003) -- (4,2.445) node[right] {\small tečna};
    \draw[dashed] (A) -- (1.5,0) node[below] {\small $x_0$};
    \draw[dashed] (B) -- (3.5,0) node[below] {\small $x_0 + h$};
  \end{tikzpicture}
\end{center}

\begin{definitionbox}
  \textbf{Derivace funkce} $f$ v bodě $x_0$ je limita
  \begin{equation*}
    f'(x_0) = \lim_{h\to 0}\frac{f(x_0 + h) - f(x_0)}{h},
  \end{equation*}
  pokud existuje a je konečná. Pak říkáme, že $f$ je v bodě $x_0$ \textbf{diferencovatelná}. Funkce, která každému $x$ přiřadí $f'(x)$, je \textbf{derivace} $f'$.

  \vspace{0.2cm}

  \textbf{Geometrický význam}: $f'(x_0)$ je směrnice tečny ke grafu v bodě $[x_0, f(x_0)]$. \textbf{Rovnice tečny}:
  \begin{equation*}
    y - f(x_0) = f'(x_0)(x - x_0).
  \end{equation*}
  \textbf{Fyzikální význam}: derivace je okamžitá rychlost změny.
\end{definitionbox}

Leibniz derivaci značil $\frac{\mathrm{d}y}{\mathrm{d}x}$, tedy `nekonečně malá změna $y$ ku nekonečně malé změně $x$'. Newton psal tečku nad písmeno, $\dot{s}$. Značení $f'$ zavedl až Lagrange. Fyzici dodnes používají všechna tři, podle nálady.

\begin{example}
  Spočítejme derivaci $f(x) = x^2$ podle definice:
  \begin{equation*}
    f'(x) = \lim_{h\to 0}\frac{{(x + h)}^2 - x^2}{h} = \lim_{h\to 0}\frac{2xh + h^2}{h} = \lim_{h\to 0}(2x + h) = 2x.
  \end{equation*}
  Tečna ke grafu $y = x^2$ v bodě $x_0 = 3$ má směrnici $f'(3) = 6$ a rovnici $y - 9 = 6(x - 3)$, tedy $y = 6x - 9$.
\end{example}

\begin{note}
  Každá diferencovatelná funkce je spojitá. Obráceně to neplatí: funkce $|x|$ je v nule spojitá, ale derivaci tam nemá. Zprava je sklon $1$, zleva $-1$ a graf má `zlom'. Weierstrassova funkce z odbočky v kapitole~\ref{kap:funkce} je spojitá, ale nemá derivaci v žádném bodě. Je to funkce, která se skládá jen ze zlomů.
\end{note}

\section{Pravidla pro derivování}

\begin{definitionbox}
  \textbf{Derivace elementárních funkcí}:
  \begin{center}
    \renewcommand{\arraystretch}{2.3}
    \begin{tabular}{|c|c||c|c|}
      \hline
      $f(x)$ & $f'(x)$ & $f(x)$ & $f'(x)$ \\
      \hline
      $c$ & $0$ & $\sin x$ & $\cos x$ \\
      $x^n$ & $nx^{n-1}$ & $\cos x$ & $-\sin x$ \\
      $e^x$ & $e^x$ & $\tg x$ & $\dfrac{1}{\cos^2 x}$ \\
      $a^x$ & $a^x\ln a$ & $\cotg x$ & $-\dfrac{1}{\sin^2 x}$ \\
      $\ln x$ & $\dfrac{1}{x}$ & $\log_a x$ & $\dfrac{1}{x\ln a}$ \\
      \hline
    \end{tabular}
  \end{center}
  Vzorec pro $x^n$ platí pro libovolný reálný exponent $n$ (na vhodném definičním oboru), například $(\sqrt{x})' = {\left(x^{\frac{1}{2}}\right)}' = \frac{1}{2}x^{-\frac{1}{2}} = \frac{1}{2\sqrt{x}}$.

  \vspace{0.2cm}

  \textbf{Pravidla}:
  \begin{align*}
    (c\cdot f)' &= c\cdot f', & (f\pm g)' &= f'\pm g',\\
    (f\cdot g)' &= f'g + fg', & \left(\frac{f}{g}\right)' &= \frac{f'g - fg'}{g^2},\\
    \left(f(g(x))\right)' &= f'(g(x))\cdot g'(x) & &\text{(derivace složené funkce)}.
  \end{align*}
\end{definitionbox}

\begin{proof}[(několika vzorců)]
  \textbf{$(x^n)' = nx^{n-1}$ pro $n\in\mathbb{N}$}: Podle binomické věty z kapitoly~\ref{kap:kombinatorika} je ${(x + h)}^n = x^n + nx^{n-1}h + \binom{n}{2}x^{n-2}h^2 + \cdots$. Po odečtení $x^n$ a vydělení $h$ zbyde $nx^{n-1} + (\text{členy s } h)$, a ty pro $h\to 0$ zmizí.

  \textbf{$(\sin x)' = \cos x$}: Podle součtového vzorce
  \begin{equation*}
    \frac{\sin(x + h) - \sin x}{h} = \sin x\cdot\frac{\cos h - 1}{h} + \cos x\cdot\frac{\sin h}{h}.
  \end{equation*}
  Podle kapitoly~\ref{kap:limity} je $\lim\frac{\sin h}{h} = 1$. Dále $\frac{\cos h - 1}{h} = \frac{-2\sin^2\frac{h}{2}}{h} = -\sin\frac{h}{2}\cdot\frac{\sin\frac{h}{2}}{\frac{h}{2}}\to 0\cdot 1 = 0$. Celkem tedy $\cos x$. Tady je vidět, proč radiány: limita $\frac{\sin h}{h} = 1$ platí jen v nich. Ve stupních by derivace sinu byla $\frac{\pi}{180}\cos x$, což nechce nikdo.

  \textbf{Součin}: $f(x + h)g(x + h) - f(x)g(x) = \left(f(x + h) - f(x)\right)g(x + h) + f(x)\left(g(x + h) - g(x)\right)$. Po vydělení $h$ a limitě dostaneme $f'g + fg'$. (Použili jsme, že $g$ je spojitá, takže $g(x + h)\to g(x)$.)
\end{proof}

Teď dokážeme všechno ostatní z tabulky. Nejdřív zbylá pravidla, protože z nich pak vzorce pro jednotlivé funkce vyjdou skoro samy.

\begin{proof}[(zbylých pravidel)]
  \textbf{Konstanta, násobek a součet}: Pro $f(x) = c$ je $\frac{c - c}{h} = 0$, takže $c' = 0$. Pro $c\cdot f$ a $f\pm g$ stačí vytknout $c$, respektive rozdělit zlomek na dva, a použít, že limita násobku a součtu je násobek a součet limit (kapitola~\ref{kap:limity}).

  \textbf{Složená funkce}: Označme $k = g(x + h) - g(x)$. Pro $h\to 0$ jde i $k\to 0$, protože $g$ je spojitá. Pokud je $k\neq 0$, můžeme rozšířit:
  \begin{equation*}
    \frac{f(g(x + h)) - f(g(x))}{h} = \frac{f(g(x) + k) - f(g(x))}{k}\cdot\frac{g(x + h) - g(x)}{h}\ \longrightarrow\ f'(g(x))\cdot g'(x).
  \end{equation*}
  První zlomek je přesně definice derivace $f$ v bodě $g(x)$ (s přírůstkem $k$), druhý je definice $g'(x)$. (Pokud by $k$ bylo nulové pro libovolně malá $h$, musí se důkaz udělat trochu opatrněji, ale výsledek platí také. To si nechte na vysokou školu.)

  \textbf{Podíl}: Nejdřív $\left(\frac{1}{g}\right)'$. Funkce $\frac{1}{u} = u^{-1}$ má derivaci $-\frac{1}{u^2}$ (vzorec pro $x^n$ níže s $n = -1$), takže podle pravidla pro složenou funkci $\left(\frac{1}{g(x)}\right)' = -\frac{g'(x)}{g^2(x)}$. Pak podle pravidla pro součin
  \begin{equation*}
    {\left(\frac{f}{g}\right)}' = {\left(f\cdot\frac{1}{g}\right)}' = \frac{f'}{g} - \frac{fg'}{g^2} = \frac{f'g - fg'}{g^2}.
  \end{equation*}
\end{proof}

\begin{proof}[(zbylých vzorců z tabulky)]
  \textbf{$e^x$}: Podle definice
  \begin{equation*}
    \frac{e^{x + h} - e^x}{h} = e^x\cdot\frac{e^h - 1}{h},
  \end{equation*}
  takže stačí ukázat, že $\lim_{h\to 0}\frac{e^h - 1}{h} = 1$. Označme $t = e^h - 1$, tedy $h = \ln(1 + t)$, a pro $h\to 0$ jde i $t\to 0$. Pak
  \begin{equation*}
    \frac{e^h - 1}{h} = \frac{t}{\ln(1 + t)} = \frac{1}{\frac{1}{t}\ln(1 + t)} = \frac{1}{\ln{(1 + t)}^{\frac{1}{t}}}\ \longrightarrow\ \frac{1}{\ln e} = 1,
  \end{equation*}
  protože ${(1 + t)}^{\frac{1}{t}}\to e$, což je (pro $t = \frac{1}{n}$) přesně definice Eulerova čísla z kapitoly~\ref{kap:exponencialni}. Geometricky to znamená, že graf $e^x$ má v bodě $[0, 1]$ tečnu se směrnicí přesně $1$. Žádný jiný základ tuto vlastnost nemá, a proto je $e$ tak výjimečné.

  \textbf{$a^x$}: Protože $a = e^{\ln a}$, je $a^x = e^{x\ln a}$. Podle pravidla pro složenou funkci $(a^x)' = e^{x\ln a}\cdot\ln a = a^x\ln a$.

  \textbf{$\ln x$}: Pro $x > 0$ platí $e^{\ln x} = x$. Zderivujeme obě strany (vlevo podle pravidla pro složenou funkci): $e^{\ln x}\cdot(\ln x)' = 1$, tedy $x\cdot(\ln x)' = 1$ a $(\ln x)' = \frac{1}{x}$. Tomuto triku se říká \textbf{derivace inverzní funkce}: graf $\ln x$ je zrcadlový obraz grafu $e^x$ podle přímky $y = x$, takže i sklony tečen se `převrátí'.

  \textbf{$\log_a x$}: Podle převodu základu je $\log_a x = \frac{\ln x}{\ln a}$, a tedy $(\log_a x)' = \frac{1}{x\ln a}$.

  \textbf{$x^n$ pro libovolné reálné $n$} (pro $x > 0$): $x^n = e^{n\ln x}$, takže $(x^n)' = e^{n\ln x}\cdot\frac{n}{x} = x^n\cdot\frac{n}{x} = nx^{n-1}$.

  \textbf{$\cos x$}: Protože $\cos x = \sin\left(\frac{\pi}{2} - x\right)$, je $(\cos x)' = \cos\left(\frac{\pi}{2} - x\right)\cdot(-1) = -\sin x$.

  \textbf{$\tg x$ a $\cotg x$}: Podle pravidla pro podíl a vztahu $\sin^2 x + \cos^2 x = 1$:
  \begin{equation*}
    (\tg x)' = {\left(\frac{\sin x}{\cos x}\right)}' = \frac{\cos x\cos x - \sin x(-\sin x)}{\cos^2 x} = \frac{1}{\cos^2 x},\qquad(\cotg x)' = \frac{-\sin^2 x - \cos^2 x}{\sin^2 x} = -\frac{1}{\sin^2 x}.
  \end{equation*}
\end{proof}

Celou tabulku tedy stačí odvodit z definice derivace, ze dvou limit ($\frac{\sin h}{h}\to 1$ a definice $e$) a ze tří pravidel (součin, složená funkce a součet). Všechno ostatní je jen šikovné kombinování. Za zapamatování opravdu stojí jen pár řádků.

\begin{example}
  \begin{align*}
    (3x^4 - 5x^2 + 7)' &= 12x^3 - 10x,\\
    (x^2\sin x)' &= 2x\sin x + x^2\cos x,\\
    {\left(\frac{x + 1}{x - 1}\right)}' &= \frac{1\cdot(x - 1) - (x + 1)\cdot 1}{{(x - 1)}^2} = \frac{-2}{{(x - 1)}^2},\\
    \left(\sin(x^2)\right)' &= \cos(x^2)\cdot 2x,\\
    \left(e^{3x}\right)' &= e^{3x}\cdot 3,\\
    \left(\sqrt{x^2 + 1}\right)' &= \frac{1}{2\sqrt{x^2 + 1}}\cdot 2x = \frac{x}{\sqrt{x^2 + 1}},\\
    \left(\ln(\cos x)\right)' &= \frac{1}{\cos x}\cdot(-\sin x) = -\tg x.
  \end{align*}
  U složených funkcí derivujeme `zvenku dovnitř': nejdřív vnější funkci (vnitřek necháme, jak je) a pak vynásobíme derivací vnitřku. Rozklad na vnější a vnitřní funkci jsme si procvičili v kapitole~\ref{kap:funkce}.
\end{example}

\begin{problem}
  Zderivujte:
  \begin{enumerate}[label=(\alph*)]
    \item $f(x) = 2x^3 - \frac{1}{x} + \sqrt{x}$,
    \item $f(x) = xe^x$,
    \item $f(x) = \dfrac{x^2}{x + 1}$,
    \item $f(x) = {(2x - 1)}^5$,
    \item $f(x) = \ln(x^2 + 1)$,
    \item $f(x) = \sin^2 x$.
  \end{enumerate}
\end{problem}

\begin{solution}
  \begin{enumerate}[label=(\alph*)]
    \item $6x^2 + \frac{1}{x^2} + \frac{1}{2\sqrt{x}}$.
    \item $e^x + xe^x = (1 + x)e^x$.
    \item $\frac{2x(x + 1) - x^2}{{(x + 1)}^2} = \frac{x^2 + 2x}{{(x + 1)}^2}$.
    \item $5{(2x - 1)}^4\cdot 2 = 10{(2x - 1)}^4$. Roznásobovat pátou mocninu by byl hřích.
    \item $\frac{2x}{x^2 + 1}$.
    \item $2\sin x\cos x = \sin 2x$.
  \end{enumerate}
\end{solution}

\section{Průběh funkce}

\begin{definitionbox}
  \textbf{Monotonie}: Pokud je na intervalu $f'(x) > 0$, je $f$ na něm \textbf{rostoucí}. Pokud $f'(x) < 0$, je \textbf{klesající}.

  \textbf{Lokální extrémy}: Pokud má diferencovatelná funkce v bodě $x_0$ lokální extrém, je $f'(x_0) = 0$ (takovým bodům se říká \textbf{stacionární}). Pokud $f'$ v $x_0$ mění znaménko z $+$ na $-$, je tam \textbf{lokální maximum}, pokud z $-$ na $+$, \textbf{lokální minimum}. Ekvivalentně: pokud $f'(x_0) = 0$ a $f''(x_0) < 0$, je tam maximum, pokud $f''(x_0) > 0$, minimum.

  \textbf{Konvexnost a konkávnost}: Druhá derivace $f'' = (f')'$ udává, jak se mění sklon. Pokud $f'' > 0$, je funkce \textbf{konvexní} (graf `drží vodu', $\cup$). Pokud $f'' < 0$, je \textbf{konkávní} ($\cap$). Bod, ve kterém se konvexnost mění v konkávnost, je \textbf{inflexní bod}.
\end{definitionbox}

Pozor: $f'(x_0) = 0$ je jen \textit{nutná} podmínka, ne postačující. Funkce $x^3$ má $f'(0) = 0$, ale v nule extrém nemá, je všude rostoucí. Tečna je tam vodorovná, ale graf jde dál nahoru. Proto je třeba ověřit změnu znaménka derivace nebo druhou derivaci.

\begin{definitionbox}
  \textbf{Postup vyšetření průběhu funkce}:
  \begin{enumerate}
    \item definiční obor, sudost/lichost, periodicita,
    \item průsečíky s osami,
    \item limity v krajních bodech definičního oboru a asymptoty,
    \item první derivace: intervaly monotonie a lokální extrémy,
    \item druhá derivace: konvexnost, konkávnost a inflexní body,
    \item obor hodnot a graf.
  \end{enumerate}
\end{definitionbox}

\begin{example}
  Vyšetřeme $f(x) = x^3 - 3x$.
  \begin{enumerate}
    \item $D(f) = \mathbb{R}$. Funkce je lichá: $f(-x) = -x^3 + 3x = -f(x)$.
    \item $f(0) = 0$. Nulové body: $x(x^2 - 3) = 0$, tedy $x\in\{0, \pm\sqrt{3}\}$.
    \item $\lim_{x\to\infty} f(x) = \infty$, $\lim_{x\to-\infty} f(x) = -\infty$. Asymptoty nemá.
    \item $f'(x) = 3x^2 - 3 = 3(x - 1)(x + 1)$. Nulové body $\pm 1$. Na $(-\infty, -1)$ je $f' > 0$ (rostoucí), na $(-1, 1)$ je $f' < 0$ (klesající), na $(1, \infty)$ je $f' > 0$ (rostoucí). Lokální maximum $f(-1) = 2$, lokální minimum $f(1) = -2$.
    \item $f''(x) = 6x$. Konkávní na $(-\infty, 0)$, konvexní na $(0, \infty)$, inflexní bod $[0, 0]$.
    \item $H(f) = \mathbb{R}$.
  \end{enumerate}
  \begin{center}
    \begin{tikzpicture}[scale=0.8]
      \draw[->] (-3,0) -- (3.2,0) node[right] {\small $x$};
      \draw[->] (0,-3.2) -- (0,3.4) node[above] {\small $y$};
      \foreach \x in {-2,-1,1,2} \draw (\x,0.08) -- (\x,-0.08) node[below] {\tiny $\x$};
      \foreach \y in {-2,2} \draw (0.08,\y) -- (-0.08,\y) node[left] {\tiny $\y$};
      \draw[thick, colorofwisdom, domain=-2.2:2.2, samples=80] plot (\x, {\x*\x*\x - 3*\x});
      \filldraw (-1,2) circle (2pt) node[above] {\small max};
      \filldraw (1,-2) circle (2pt) node[below] {\small min};
      \filldraw (0,0) circle (2pt);
      \node[above right] at (0.05,0.05) {\tiny inflexe};
      \draw[dashed, gray] (-1.8,2) -- (-0.2,2);
      \draw[dashed, gray] (0.2,-2) -- (1.8,-2);
    \end{tikzpicture}
  \end{center}
\end{example}

\section{Užití derivací}

\subsection{Optimalizační úlohy}

V kapitole~\ref{kap:elementarni} jsme hledali maxima jen u kvadratických funkcí, kde stačilo najít vrchol. S derivacemi umíme najít extrém téměř čehokoli.

\begin{example}
  Plechovka tvaru válce má mít objem $1$ litr ($1000\ \si{cm^3}$). Jaké rozměry má mít, aby se na ni spotřebovalo co nejméně plechu?

  Z objemu $\pi r^2v = 1000$ máme $v = \frac{1000}{\pi r^2}$. Povrch je
  \begin{equation*}
    S(r) = 2\pi r^2 + 2\pi rv = 2\pi r^2 + \frac{2000}{r}.
  \end{equation*}
  Derivace $S'(r) = 4\pi r - \frac{2000}{r^2} = 0$, tedy $r^3 = \frac{500}{\pi}$ a $r\approx 5{,}42\ \si{cm}$. Výška vychází $v = \frac{1000}{\pi r^2}\approx 10{,}84\ \si{cm}$, tedy přesně $v = 2r$. Druhá derivace $S''(r) = 4\pi + \frac{4000}{r^3} > 0$, takže jde opravdu o minimum.

  Optimální plechovka má výšku rovnou průměru. Až půjdete příště do obchodu, všimněte si, že většina plechovek je vyšší. Na vině je prý marketing, pohodlný úchop a to, že víčka bývají z pevnějšího, a tedy dražšího plechu než plášť.
\end{example}

\begin{problem}
  \begin{enumerate}[label=(\alph*)]
    \item Z obdélníkového kartonu $30\times 20\ \si{cm}$ vystřihneme v rozích čtverce o straně $x$ a ohneme krabici bez víka. Pro jaké $x$ má největší objem?
    \item Najděte bod na parabole $y = x^2$, který je nejblíže bodu $[0, 2]$.
    \item Kámen vyhozený vzhůru je v čase $t$ ve výšce $h(t) = 20t - 5t^2$. Určete jeho rychlost v čase $t = 1$ a kdy dosáhne vrcholu.
  \end{enumerate}
\end{problem}

\begin{solution}
  \begin{enumerate}[label=(\alph*)]
    \item $V(x) = x(30 - 2x)(20 - 2x) = 4x^3 - 100x^2 + 600x$ pro $x\in(0, 10)$. $V'(x) = 12x^2 - 200x + 600 = 0$, tedy $3x^2 - 50x + 150 = 0$ a $x = \frac{50\pm\sqrt{2500 - 1800}}{6} = \frac{50\pm\sqrt{700}}{6}$. V intervalu je jen $x\approx 3{,}92\ \si{cm}$.
    \item Druhá mocnina vzdálenosti je $d^2 = x^2 + {(x^2 - 2)}^2 = x^4 - 3x^2 + 4$. Derivace $4x^3 - 6x = 2x(2x^2 - 3) = 0$. V bodě $0$ je lokální maximum ($d^2 = 4$), pro $x = \pm\sqrt{\frac{3}{2}}$ minimum s $d^2 = \frac{9}{4} - \frac{9}{2} + 4 = \frac{7}{4}$. Nejbližší body jsou $\left[\pm\sqrt{\frac{3}{2}}, \frac{3}{2}\right]$.
    \item $v(t) = h'(t) = 20 - 10t$, $v(1) = 10\ \si{m/s}$. Vrchol pro $v(t) = 0$, tedy $t = 2\ \si{s}$. Stejně jako v kapitole~\ref{kap:elementarni}, jen bez doplňování na čtverec. Zrychlení $h''(t) = -10\ \si{m/s^2}$ je tíhové zrychlení.
  \end{enumerate}
\end{solution}

\subsection{L'Hospitalovo pravidlo}

\begin{definitionbox}
  \textbf{L'Hospitalovo pravidlo}: Pokud je limita $\lim_{x\to a}\frac{f(x)}{g(x)}$ typu $\frac{0}{0}$ nebo $\frac{\infty}{\infty}$ a existuje $\lim_{x\to a}\frac{f'(x)}{g'(x)}$, pak
  \begin{equation*}
    \lim_{x\to a}\frac{f(x)}{g(x)} = \lim_{x\to a}\frac{f'(x)}{g'(x)}.
  \end{equation*}
\end{definitionbox}

\begin{example}
  \begin{equation*}
    \lim_{x\to 0}\frac{e^x - 1}{x} = \lim_{x\to 0}\frac{e^x}{1} = 1,\qquad\lim_{x\to\infty}\frac{\ln x}{x} = \lim_{x\to\infty}\frac{\frac{1}{x}}{1} = 0.
  \end{equation*}
  Druhá limita říká, že logaritmus roste pomaleji než jakákoli přímka. Pravidlo pochází od markýze de l'Hospital, který ho v roce 1696 publikoval v první učebnici diferenciálního počtu vůbec. Ve skutečnosti ho ale objevil jeho učitel Johann Bernoulli, kterému markýz platil za to, že mu bude posílat své objevy a nechá ho je vydávat pod svým jménem. Pravidlo tedy dodnes nese jméno člověka, který si ho koupil.
\end{example}

Na závěr malý výhled. Derivace odpovídá na otázku `znám funkci, jak rychle se mění?'. Často ale potřebujeme opak: `znám rychlost, kde jsem?' Na tuto obrácenou otázku odpovídá integrál. Že jde opravdu o obrácení derivace, a že tatáž operace zároveň počítá plochy pod grafem, říká \textbf{základní věta integrálního počtu} v sekci~\ref{sec:zvip} na straně~\pageref{sec:zvip}. Právě tato věta je to, co Newton a Leibniz objevili. Derivace samotné znali v nějaké podobě už Fermat a další.

\section{Příběh --- Válka o kalkulus}

Newton vymyslel diferenciální počet kolem roku 1666, během moru, kdy byla univerzita v Cambridge zavřená a on se uchýlil na rodinnou farmu. Nic ale nezveřejnil. Leibniz na totéž přišel nezávisle kolem roku 1675 a v roce 1684 to publikoval, s mnohem lepším značením, které používáme dodnes. Pak začala jedna z nejošklivějších hádek v dějinách vědy.

Newtonovi stoupenci obvinili Leibnize, že myšlenku ukradl. Leibniz se v roce 1711 obrátil na Královskou společnost v Londýně, aby spor rozhodla. Společnost jmenovala komisi, která po roce vydala zprávu, že prvenství patří Newtonovi a Leibniz je podezřelý z plagiátorství. Prezidentem Královské společnosti byl v té době Isaac Newton. Zprávu, jak se později ukázalo, napsal z velké části on sám. A pak o ní ještě anonymně napsal pochvalnou recenzi.

Leibniz zemřel v roce 1716 v Hannoveru, zapomenutý a v nemilosti u svého zaměstnavatele, který se mezitím stal anglickým králem Jiřím I. Na jeho pohřbu byl podle očitého svědka jediný člověk, jeho tajemník. Newton se o jeho smrti dozvěděl a prý prohlásil, že `zlomil Leibnizovi srdce'. Spor otrávil vztahy mezi britskými a kontinentálními matematiky na celé století. Britové z vlastenectví trvali na Newtonově nepraktickém značení s tečkami, a britská matematika proto za kontinentem desítky let zaostávala. Dnes víme, že oba přišli na kalkulus nezávisle, a celý svět používá Leibnizovo $\frac{\mathrm{d}y}{\mathrm{d}x}$.

\section{Shrnutí a otázky k ústní zkoušce}

\begin{problem}
  Odpovězte vlastními slovy:
  \begin{enumerate}
    \item Definujte derivaci funkce v bodě. Jaký má geometrický a fyzikální význam?
    \item Spočítejte derivaci $x^2$ podle definice. Jak vypadá rovnice tečny?
    \item Jaký je vztah mezi spojitostí a diferencovatelností? Uveďte příklad spojité funkce bez derivace.
    \item Uveďte derivace elementárních funkcí. Odvoďte derivaci $x^n$ a $\sin x$.
    \item Uveďte pravidla pro derivaci součtu, součinu, podílu a složené funkce.
    \item Jak souvisí znaménko první derivace s monotonií? Jak se hledají lokální extrémy? Je $f'(x_0) = 0$ postačující podmínka?
    \item Co říká druhá derivace? Co je inflexní bod?
    \item Popište postup vyšetření průběhu funkce.
    \item Uveďte příklad optimalizační úlohy a vyřešte ji.
    \item Co říká l'Hospitalovo pravidlo?
  \end{enumerate}
\end{problem}

\koneckapitoly

\chapter{Integrály}
\label{kap:integraly}
\mathquote{Give me a place to stand, and I shall move the earth.}{Archimédés}

Integrál je opak derivace, a zároveň je to obsah plochy pod grafem. To, že tyto dvě věci, které spolu na první pohled vůbec nesouvisí, jsou ve skutečnosti totéž, je obsahem \textbf{základní věty integrálního počtu} a je to jeden z nejkrásnějších objevů v dějinách matematiky. Díky integrálům spočítáme obsahy, objemy, dráhy, práci a pravděpodobnosti. Tahle kapitola je posledním dílkem skládačky. Konečně dokážeme vzorce pro objem koule a kužele, které jsme v kapitole~\ref{kap:stereometrie} museli jen přijmout.

\section{Primitivní funkce a neurčitý integrál}

\begin{definitionbox}
  Funkce $F$ je \textbf{primitivní funkcí} k funkci $f$ na intervalu $I$, pokud pro všechna $x\in I$ platí
  \begin{equation*}
    F'(x) = f(x).
  \end{equation*}
  Pokud je $F$ primitivní funkce k $f$, je primitivní funkcí i $F + C$ pro každou konstantu $C$, protože derivace konstanty je nula. A jiné primitivní funkce (na intervalu) neexistují. Množinu všech primitivních funkcí nazýváme \textbf{neurčitý integrál} a značíme
  \begin{equation*}
    \int f(x)\,\mathrm{d}x = F(x) + C.
  \end{equation*}
  Funkci $f$ nazýváme \textbf{integrand}, $C$ \textbf{integrační konstanta}.
\end{definitionbox}

Integrování je tedy derivování pozpátku: hledáme funkci, jejíž derivace je to, co máme. Každý integrál si můžeme snadno ověřit tím, že výsledek zderivujeme. A pozor na konstantu $C$! Zapomenout ji je u ústní maturity spolehlivý způsob, jak zkoušejícího přimět ke smutnému výrazu.

\begin{definitionbox}
  \textbf{Základní integrály} (dostaneme je obrácením tabulky derivací z kapitoly~\ref{kap:derivace}):
  \begin{center}
    \renewcommand{\arraystretch}{1.8}
    \begin{tabular}{|c|c||c|c|}
      \hline
      $f(x)$ & $\int f(x)\,\mathrm{d}x$ & $f(x)$ & $\int f(x)\,\mathrm{d}x$ \\
      \hline
      $x^n\ (n\neq -1)$ & $\dfrac{x^{n+1}}{n + 1} + C$ & $\sin x$ & $-\cos x + C$ \\
      $\dfrac{1}{x}$ & $\ln|x| + C$ & $\cos x$ & $\sin x + C$ \\
      $e^x$ & $e^x + C$ & $\dfrac{1}{\cos^2 x}$ & $\tg x + C$ \\
      $a^x$ & $\dfrac{a^x}{\ln a} + C$ & $\dfrac{1}{\sin^2 x}$ & $-\cotg x + C$ \\
      \hline
    \end{tabular}
  \end{center}

  \textbf{Linearita}:
  \begin{equation*}
    \int\left(af(x) + bg(x)\right)\mathrm{d}x = a\int f(x)\,\mathrm{d}x + b\int g(x)\,\mathrm{d}x.
  \end{equation*}
\end{definitionbox}

Proč $\ln|x|$ a ne $\ln x$? Funkce $\frac{1}{x}$ je definovaná i pro záporná $x$, ale $\ln x$ ne. Pro $x < 0$ je ${(\ln(-x))}' = \frac{1}{-x}\cdot(-1) = \frac{1}{x}$, takže absolutní hodnota pokryje oba případy.

\begin{example}
  \begin{align*}
    \int(3x^2 + 2x - 1)\,\mathrm{d}x &= x^3 + x^2 - x + C,\\
    \int\left(\sqrt{x} + \frac{1}{x^2}\right)\mathrm{d}x &= \int\left(x^{\frac{1}{2}} + x^{-2}\right)\mathrm{d}x = \frac{2}{3}x^{\frac{3}{2}} - \frac{1}{x} + C,\\
    \int(2\cos x - 3e^x)\,\mathrm{d}x &= 2\sin x - 3e^x + C.
  \end{align*}
  Kontrola derivací: ${\left(\frac{2}{3}x^{\frac{3}{2}} - \frac{1}{x}\right)}' = x^{\frac{1}{2}} + \frac{1}{x^2}$.
\end{example}

\begin{note}
  Integrovat je mnohem těžší než derivovat. Derivovat se dá podle pravidel skoro cokoli. Integrál ale nemá pravidlo pro součin ani pro složenou funkci a mnoho funkcí, které vypadají nevinně, nemá primitivní funkci vyjádřitelnou elementárními funkcemi. Například $\int e^{-x^2}\,\mathrm{d}x$, integrál klíčový pro statistiku, se elementárně spočítat nedá. To se dokonce dá dokázat. Derivování je řemeslo, integrování umění.
\end{note}

\subsection{Substituce a integrace per partes}

\begin{definitionbox}
  \textbf{Substituční metoda} (obrácené pravidlo pro derivaci složené funkce):
  \begin{equation*}
    \int f(g(x))\cdot g'(x)\,\mathrm{d}x = \int f(t)\,\mathrm{d}t,\qquad\text{kde } t = g(x),\ \mathrm{d}t = g'(x)\,\mathrm{d}x.
  \end{equation*}

  \textbf{Integrace per partes} (obrácené pravidlo pro derivaci součinu):
  \begin{equation*}
    \int u(x)v'(x)\,\mathrm{d}x = u(x)v(x) - \int u'(x)v(x)\,\mathrm{d}x.
  \end{equation*}
\end{definitionbox}

Per partes plyne ze vzorce $(uv)' = u'v + uv'$. Zintegrováním obou stran dostaneme $uv = \int u'v + \int uv'$ a stačí jeden integrál převést na druhou stranu. Metoda se vyplatí, pokud je $\int u'v$ jednodušší než původní integrál. Typicky za $u$ volíme mnohočlen (derivováním se zjednoduší) nebo logaritmus (jeho derivace je hezčí než on sám).

\begin{example}[ (substituce)]
  \begin{align*}
    \int 2x\cos(x^2)\,\mathrm{d}x &= \left|\begin{array}{c} t = x^2\\ \mathrm{d}t = 2x\,\mathrm{d}x \end{array}\right| = \int\cos t\,\mathrm{d}t = \sin t + C = \sin(x^2) + C,\\
    \int xe^{x^2}\,\mathrm{d}x &= \left|\begin{array}{c} t = x^2\\ \mathrm{d}t = 2x\,\mathrm{d}x \end{array}\right| = \frac{1}{2}\int e^t\,\mathrm{d}t = \frac{1}{2}e^{x^2} + C,\\
    \int\tg x\,\mathrm{d}x = \int\frac{\sin x}{\cos x}\,\mathrm{d}x &= \left|\begin{array}{c} t = \cos x\\ \mathrm{d}t = -\sin x\,\mathrm{d}x \end{array}\right| = -\int\frac{\mathrm{d}t}{t} = -\ln|\cos x| + C.
  \end{align*}
\end{example}

\begin{example}[ (per partes)]
  \begin{align*}
    \int xe^x\,\mathrm{d}x &= \left|\begin{array}{cc} u = x & v' = e^x\\ u' = 1 & v = e^x \end{array}\right| = xe^x - \int e^x\,\mathrm{d}x = xe^x - e^x + C,\\
    \int\ln x\,\mathrm{d}x &= \left|\begin{array}{cc} u = \ln x & v' = 1\\ u' = \frac{1}{x} & v = x \end{array}\right| = x\ln x - \int 1\,\mathrm{d}x = x\ln x - x + C,\\
    \int x\sin x\,\mathrm{d}x &= \left|\begin{array}{cc} u = x & v' = \sin x\\ u' = 1 & v = -\cos x \end{array}\right| = -x\cos x + \int\cos x\,\mathrm{d}x = -x\cos x + \sin x + C.
  \end{align*}
  U $\ln x$ je trik v tom, že integrand chápeme jako $\ln x\cdot 1$. Matematici jsou někdy líní až vynalézaví.
\end{example}

\begin{problem}
  Vypočítejte:
  \begin{enumerate}[label=(\alph*)]
    \item $\int(4x^3 - 6x + 5)\,\mathrm{d}x$,
    \item $\int\frac{x^2 + 1}{x}\,\mathrm{d}x$,
    \item $\int{(2x + 1)}^{5}\,\mathrm{d}x$,
    \item $\int\sin x\cos x\,\mathrm{d}x$,
    \item $\int x^2\ln x\,\mathrm{d}x$,
    \item $\int\frac{2x}{x^2 + 1}\,\mathrm{d}x$.
  \end{enumerate}
\end{problem}

\begin{solution}
  \begin{enumerate}[label=(\alph*)]
    \item $x^4 - 3x^2 + 5x + C$.
    \item $\int\left(x + \frac{1}{x}\right)\mathrm{d}x = \frac{x^2}{2} + \ln|x| + C$.
    \item $t = 2x + 1$, $\mathrm{d}t = 2\,\mathrm{d}x$: $\frac{1}{2}\cdot\frac{t^6}{6} = \frac{{(2x + 1)}^6}{12} + C$.
    \item $t = \sin x$: $\frac{\sin^2 x}{2} + C$. (Jiná cesta přes $\sin x\cos x = \frac{1}{2}\sin 2x$ dá $-\frac{\cos 2x}{4} + C$. Oba výsledky jsou správně, liší se jen o konstantu. Ověřte si to.)
    \item $u = \ln x$, $v' = x^2$: $\frac{x^3}{3}\ln x - \int\frac{x^3}{3}\cdot\frac{1}{x}\,\mathrm{d}x = \frac{x^3}{3}\ln x - \frac{x^3}{9} + C$.
    \item $t = x^2 + 1$: $\ln(x^2 + 1) + C$. Obecně $\int\frac{f'(x)}{f(x)}\,\mathrm{d}x = \ln|f(x)| + C$.
  \end{enumerate}
\end{solution}

\section{Určitý integrál}

Jak spočítat obsah plochy pod grafem funkce $f$ mezi $a$ a $b$? Rozdělíme interval $[a, b]$ na $n$ malých kousků, nad každým postavíme obdélník s výškou rovnou hodnotě funkce a obsahy obdélníků sečteme. Čím jemnější dělení, tím přesnější odhad. V limitě dostaneme přesný obsah.

\begin{center}
  \begin{tikzpicture}[scale=1.6]
    \draw[->] (-0.2,0) -- (2.3,0) node[right] {\small $x$};
    \draw[->] (0,-0.2) -- (0,2.2) node[above] {\small $y$};
    \foreach \i in {0,...,7} {
      \pgfmathsetmacro{\x}{0.25*\i}
      \pgfmathsetmacro{\h}{0.5*(\x+0.25)*(\x+0.25)}
      \fill[colorofwisdom!25] (\x,0) rectangle ({\x+0.25},\h);
      \draw[colorofwisdom!70] (\x,0) rectangle ({\x+0.25},\h);
    }
    \draw[thick, colorofwisdom, domain=0:2.05, samples=60] plot (\x, {0.5*\x*\x});
    \node[below] at (0,0) {\small $a$};
    \node[below] at (2,0) {\small $b$};
  \end{tikzpicture}
\end{center}

\begin{definitionbox}
  \textbf{Určitý (Riemannův) integrál} funkce $f$ na intervalu $[a, b]$ je limita součtů obsahů obdélníků
  \begin{equation*}
    \int_a^b f(x)\,\mathrm{d}x = \lim_{n\to\infty}\sum_{i=1}^{n}f(x_i)\cdot\Delta x,
  \end{equation*}
  kde $\Delta x = \frac{b - a}{n}$ je šířka obdélníků a $x_i$ jsou body v jednotlivých dílcích. Pro spojitou funkci tato limita vždy existuje. Pro nezápornou funkci je to \textbf{obsah plochy} pod grafem. Části pod osou $x$ se počítají záporně.

  \vspace{0.2cm}

  \textbf{Newtonova--Leibnizova věta} (základní věta integrálního počtu): Pokud je $f$ spojitá na $[a, b]$ a $F$ je k ní primitivní funkce, pak
  \begin{equation*}
    \int_a^b f(x)\,\mathrm{d}x = F(b) - F(a) = {\big[F(x)\big]}_a^b.
  \end{equation*}
\end{definitionbox}

Znak $\int$ je mimochodem protažené písmeno S, jako \textit{summa}. Zavedl ho Leibniz a připomíná, že integrál je součet nekonečně mnoha nekonečně tenkých obdélníků o šířce $\mathrm{d}x$.

Proč platí Newtonova--Leibnizova věta? Označme $S(x)$ obsah plochy pod grafem od $a$ do $x$. Když $x$ zvětšíme o malé $h$, přibude proužek o šířce $h$ a výšce přibližně $f(x)$, tedy $S(x + h) - S(x)\approx f(x)\cdot h$. Po vydělení $h$ a limitě dostaneme $S'(x) = f(x)$. Funkce `obsah plochy' je tedy primitivní funkcí k $f$! A protože $S(a) = 0$, je $S(b) = F(b) - F(a)$ pro libovolnou primitivní funkci $F$ (konstanta se odečte). Celá věta je tedy důsledkem toho, že rychlost, jakou roste plocha, je rovna výšce grafu.

\begin{example}
  \begin{itemize}
    \item $\int_0^1 x^2\,\mathrm{d}x = {\left[\frac{x^3}{3}\right]}_0^1 = \frac{1}{3}$. Obsah pod parabolou od $0$ do $1$ je přesně třetina jednotkového čtverce. Archimédés tento výsledek (v trochu jiné podobě) odvodil už ve 3. století před naším letopočtem, devatenáct století před Newtonem a Leibnizem, takzvanou metodou vyčerpávání.
    \item $\int_0^{\pi}\sin x\,\mathrm{d}x = {\big[-\cos x\big]}_0^{\pi} = 1 - (-1) = 2$. Obsah pod jedním `obloukem' sinusoidy je přesně $2$.
    \item $\int_1^e\frac{1}{x}\,\mathrm{d}x = {\big[\ln x\big]}_1^e = 1 - 0 = 1$. Tohle je mimochodem jedna z možných definic čísla $e$: je to číslo, pro které má plocha pod hyperbolou $\frac{1}{x}$ od $1$ do $e$ obsah $1$.
  \end{itemize}
\end{example}

\begin{definitionbox}
  \textbf{Vlastnosti určitého integrálu}:
  \begin{align*}
    \int_a^b(f + g)\,\mathrm{d}x &= \int_a^b f\,\mathrm{d}x + \int_a^b g\,\mathrm{d}x, & \int_a^b c\cdot f\,\mathrm{d}x &= c\int_a^b f\,\mathrm{d}x,\\
    \int_a^b f\,\mathrm{d}x &= \int_a^c f\,\mathrm{d}x + \int_c^b f\,\mathrm{d}x, & \int_a^b f\,\mathrm{d}x &= -\int_b^a f\,\mathrm{d}x.
  \end{align*}
\end{definitionbox}

\section{Užití určitého integrálu}

\subsection{Obsah rovinného útvaru}

\begin{definitionbox}
  Obsah plochy mezi grafy funkcí $f$ a $g$, kde $f(x)\geq g(x)$ na $[a, b]$, je
  \begin{equation*}
    S = \int_a^b\left(f(x) - g(x)\right)\mathrm{d}x.
  \end{equation*}
  Pokud funkce na části intervalu klesne pod osu $x$, musíme interval rozdělit a části pod osou vzít v absolutní hodnotě.
\end{definitionbox}

\begin{example}
  Obsah plochy mezi přímkou $y = x$ a parabolou $y = x^2$. Průsečíky: $x = x^2$, tedy $x\in\{0, 1\}$. Na $[0, 1]$ je $x\geq x^2$, a tedy
  \begin{equation*}
    S = \int_0^1(x - x^2)\,\mathrm{d}x = {\left[\frac{x^2}{2} - \frac{x^3}{3}\right]}_0^1 = \frac{1}{2} - \frac{1}{3} = \frac{1}{6}.
  \end{equation*}
  \begin{center}
    \begin{tikzpicture}[scale=2.5]
      \draw[->] (-0.1,0) -- (1.3,0) node[right] {\small $x$};
      \draw[->] (0,-0.1) -- (0,1.3) node[above] {\small $y$};
      \fill[colorofwisdom!30, domain=0:1, samples=40] plot (\x, {\x}) -- plot[domain=1:0] (\x, {\x*\x}) -- cycle;
      \draw[thick, domain=-0.05:1.2] plot (\x, \x) node[right] {\small $y = x$};
      \draw[thick, colorofwisdom, domain=-0.05:1.1, samples=40] plot (\x, {\x*\x}) node[right] {\small $y = x^2$};
      \node[below] at (1,0) {\small $1$};
    \end{tikzpicture}
  \end{center}
\end{example}

\begin{example}
  Pozor na části pod osou. $\int_0^{2\pi}\sin x\,\mathrm{d}x = 0$, protože kladná a záporná část se vyruší. Obsah plochy mezi sinusoidou a osou $x$ na $[0, 2\pi]$ je ale $2 + 2 = 4$.
\end{example}

\subsection{Objem rotačního tělesa}

\begin{definitionbox}
  Objem tělesa, které vznikne rotací plochy pod grafem $f$ (na $[a, b]$) kolem osy $x$, je
  \begin{equation*}
    V = \pi\int_a^b f^2(x)\,\mathrm{d}x.
  \end{equation*}
\end{definitionbox}

Vzorec plyne ze stejné myšlenky jako určitý integrál: těleso rozřežeme na tenké kotoučky. Kotouček v bodě $x$ je válec s poloměrem $f(x)$ a výškou $\mathrm{d}x$, a má tedy objem $\pi f^2(x)\,\mathrm{d}x$. Sečtením dostaneme integrál. Je to vlastně Cavalieriho princip z kapitoly~\ref{kap:stereometrie} dovedený do konce.

\begin{example}\label{ex:objem-koule}
  Konečně dokážeme vzorce z kapitoly o stereometrii.

  \textbf{Koule}: vznikne rotací půlkružnice $f(x) = \sqrt{r^2 - x^2}$ na $[-r, r]$:
  \begin{equation*}
    V = \pi\int_{-r}^{r}(r^2 - x^2)\,\mathrm{d}x = \pi{\left[r^2x - \frac{x^3}{3}\right]}_{-r}^{r} = \pi\left(2r^3 - \frac{2r^3}{3}\right) = \frac{4}{3}\pi r^3.
  \end{equation*}

  \textbf{Kužel}: vznikne rotací úsečky $f(x) = \frac{r}{v}x$ na $[0, v]$:
  \begin{equation*}
    V = \pi\int_0^v\frac{r^2}{v^2}x^2\,\mathrm{d}x = \pi\frac{r^2}{v^2}\cdot\frac{v^3}{3} = \frac{1}{3}\pi r^2v.
  \end{equation*}
  Odtud pochází ta záhadná třetina. Je to $\int x^2 = \frac{x^3}{3}$.
\end{example}

\begin{problem}
  \begin{enumerate}[label=(\alph*)]
    \item Vypočítejte $\int_1^2(3x^2 - 2x)\,\mathrm{d}x$ a $\int_0^{\frac{\pi}{2}}\cos x\,\mathrm{d}x$.
    \item Určete obsah plochy ohraničené parabolou $y = 4 - x^2$ a osou $x$.
    \item Určete obsah plochy mezi grafy $y = x^2$ a $y = 2x + 3$.
    \item Určete objem tělesa, které vznikne rotací plochy pod grafem $y = \sqrt{x}$ na $[0, 4]$ kolem osy $x$.
    \item Auto zrychluje a jeho rychlost v čase $t$ je $v(t) = 3t\ \si{m/s}$. Jakou dráhu ujede za prvních $10$ sekund?
  \end{enumerate}
\end{problem}

\begin{solution}
  \begin{enumerate}[label=(\alph*)]
    \item ${\big[x^3 - x^2\big]}_1^2 = (8 - 4) - (1 - 1) = 4$. ${\big[\sin x\big]}_0^{\frac{\pi}{2}} = 1$.
    \item Průsečíky s osou $\pm 2$. $S = \int_{-2}^{2}(4 - x^2)\,\mathrm{d}x = {\left[4x - \frac{x^3}{3}\right]}_{-2}^{2} = 2\left(8 - \frac{8}{3}\right) = \frac{32}{3}$.
    \item Průsečíky: $x^2 = 2x + 3$, $x\in\{-1, 3\}$. Na $[-1, 3]$ je přímka nad parabolou: $S = \int_{-1}^{3}(2x + 3 - x^2)\,\mathrm{d}x = {\left[x^2 + 3x - \frac{x^3}{3}\right]}_{-1}^{3} = 9 - \left(-\frac{5}{3}\right) = \frac{32}{3}$.
    \item $V = \pi\int_0^4 x\,\mathrm{d}x = \pi\cdot 8 = 8\pi$. Těleso je rotační paraboloid, tvarem podobný satelitní anténě z kapitoly~\ref{kap:kuzelosecky}.
    \item Dráha je integrál rychlosti: $s = \int_0^{10}3t\,\mathrm{d}t = {\left[\frac{3t^2}{2}\right]}_0^{10} = 150\ \si{m}$. Derivace dráhy je rychlost, integrál rychlosti je dráha. Kruh se uzavírá.
  \end{enumerate}
\end{solution}

\begin{note}
  Integrovat se dá i přes nekonečný interval. Takovému integrálu se říká \textbf{nevlastní}. Například $\int_1^{\infty}\frac{1}{x^2}\,\mathrm{d}x = \lim_{b\to\infty}{\left[-\frac{1}{x}\right]}_1^b = 1$. Nekonečně dlouhá plocha má konečný obsah! Ještě podivnější je \textbf{Gabrielův roh}, který vznikne rotací grafu $\frac{1}{x}$ na $[1, \infty)$. Jeho objem je $\pi\int_1^{\infty}\frac{1}{x^2}\,\mathrm{d}x = \pi$, ale jeho povrch je nekonečný. Do rohu se tedy vejde konečně mnoho barvy, ale nikdy nenatřete jeho vnitřní stěnu. Objevil ho v roce 1643 Evangelista Torricelli a pojmenován je po archandělu Gabrielovi, který podle tradice zatroubí na roh v den posledního soudu.
\end{note}

\section{Odbočka --- Spálené papíry}

Bernhard Riemann, po kterém je pojmenovaný určitý integrál, byl plachý a nemocný muž z chudé rodiny luteránského pastora. Původně měl studovat teologii, ale matematika ho nepustila. Za svůj krátký život publikoval jen hrstku prací, ale skoro každá z nich změnila matematiku. Jeho geometrie zakřivených prostorů se o šedesát let později stala jazykem Einsteinovy obecné teorie relativity, o které jsme mluvili u pátého postulátu v kapitole~\ref{kap:planimetrie}.

V roce 1859 publikoval jediný článek o prvočíslech, osm stran dlouhý. Na jeho konci mimochodem zmínil domněnku o tom, kde leží nulové body jedné komplexní funkce, a poznamenal, že ji zatím nedokázal, ale pro jeho účely ji nepotřebuje. \textbf{Riemannova hypotéza} je dnes považována za nejdůležitější nevyřešený problém matematiky. Clayův matematický institut za její důkaz v roce 2000 vypsal odměnu jeden milion dolarů. Dosud nikdo nepřišel na to, jak ji dokázat.

Riemann onemocněl tuberkulózou a odjel se léčit do Itálie, kde v roce 1866 ve věku 39 let zemřel. Po jeho smrti začala jeho hospodyně uklízet pracovnu a velkou část jeho poznámek spálila, dřív než ji kolegové stačili zastavit. Nikdo neví, co v těch papírech bylo. Možná nic. A možná tam byl ten důkaz.

\section{Shrnutí a otázky k ústní zkoušce}

\begin{problem}
  Odpovězte vlastními slovy:
  \begin{enumerate}
    \item Co je primitivní funkce a neurčitý integrál? Proč se přidává konstanta $C$?
    \item Uveďte základní integrály. Proč je $\int\frac{1}{x}\,\mathrm{d}x = \ln|x| + C$?
    \item Popište substituční metodu a integraci per partes. Z jakých pravidel pro derivace vycházejí?
    \item Co je určitý integrál? Jak souvisí s obsahem plochy?
    \item Vyslovte Newtonovu--Leibnizovu větu a vysvětlete, proč platí.
    \item Jak spočítáte obsah plochy mezi dvěma grafy? Na co si dát pozor, když je funkce záporná?
    \item Jak spočítáte objem rotačního tělesa? Odvoďte vzorec pro objem koule a kužele.
    \item Uveďte příklady využití integrálů ve fyzice.
  \end{enumerate}
\end{problem}

\koneckapitoly

\end{document}